% ############################################### PREAMBLE AND CONTROL PANEL #

% ########################################################### DOCUMENT CLASS #
\documentclass[  %   draft
             % , demo  % black boxes instead of figures
              , 8pt, twoside, openright]{extbook}

% ########################################################## AND CONTROL PANEL #

% true - hmarginration 1:1, urlcolor - darkgreen, version - for searching
\newif\ifscreen 	   \screentrue
% true - no apendix, no author of grammer, no backmatter ISBN, 
\newif\ifPDF               \PDFtrue
% true - rotate pages of morph verb tables
\newif\ifPDFrotate         \PDFrotatetrue
% true - adds white searchable value !
\newif\ifPDFsearchvalue    \PDFsearchvaluetrue
\newif\ifmakecovers        \makecoverstrue
\newif\ifmakefrontmatter   \makefrontmattertrue
\newif\ifinputintroduction \inputintroductiontrue
\newif\ifinputletters      \inputletterstrue
\newif\ifinputphon         \inputphontrue
\newif\ifinputmorpho       \inputmorphotrue
\newif\ifmakebackmatter    \makebackmattertrue
\newif\ifshowcrop           \showcroptrue
% corrections
\newif\ifC                  \Ctrue

\newif\ifPRINTHOUSE 	   \PRINTHOUSEfalse
\newif\ifPDFSEARCH  	   \PDFSEARCHfalse
\newif\ifPDFPRINT 	   \PDFPRINTfalse
\newif\ifPDFtest	   \PDFtesttrue

%versions
%PRINTHOUSE
\ifPRINTHOUSE
\makecoversfalse
\screenfalse
\PDFfalse
\PDFrotatefalse
\PDFsearchvaluefalse
\showcroptrue
\Ctrue
\fi

%PDFSEARCH
\ifPDFSEARCH
\screentrue
\PDFtrue
\PDFrotatetrue
\PDFsearchvaluetrue
\showcropfalse
\Ctrue
\fi

%PDFPRINT
\ifPDFPRINT
\screenfalse
\PDFtrue
\PDFrotatefalse
\PDFsearchvaluefalse
\showcropfalse
\Ctrue
\fi

%PDFtest
\ifPDFtest
\makecoversfalse
\makefrontmatterfalse
\inputintroductionfalse
\inputletterstrue
\inputphonfalse
\inputmorphofalse
\makebackmatterfalse
\screentrue
\PDFtrue
\PDFrotatefalse
\PDFsearchvaluefalse
\showcropfalse
\Cfalse
\fi

% 1. PDF (for printing)  
% a) searchable value, b) no apendix, c) version= PDF d) no ISBN and no ISBN bar
% e) automatic rotation of landscape pages, f) no author of Icelandic Grammer
% 2. PDF (for viewing offline and searching) + screen (horizontal ratio 1:1) 
% 3. PRINTHOUSE
% a) no searchable value, b) apendix, c) version ISBN d) ISBN bar
% e) no rotation of landscape pages, f)urllinkcolor black
% f) A4 paper dimentions, C5 layout, shows cropping marks

% ########################################################## PAGE GEOMETRIES #

\usepackage%[showframe]
           {geometry}
           

\ifshowcrop
% Basic geometry of document
\geometry
  { headsep    =   \baselineskip
  , textwidth  = 42\baselineskip
  , textheight = 60\baselineskip
  , hmarginratio = \ifscreen 1:1\else 2:3\fi
  , vmarginratio = 2:3
  , bindingoffset = 0cm
  , onecolumn
  , a4paper 
  , layout=c5paper,
  , layouthoffset=\dimexpr(\paperwidth-\csname Gm@layoutwidth\endcsname)/2\relax
   , layoutvoffset=\dimexpr(\paperheight-\csname Gm@layoutheight\endcsname)/2\relax
  , showcrop
  }
\else
\geometry
  { headsep    =   \baselineskip
  , textwidth  = 42\baselineskip
  , textheight = 60\baselineskip
  , hmarginratio = \ifscreen 1:1\else 2:3\fi
  , vmarginratio = 2:3
  , bindingoffset = 0cm
  , onecolumn,
  , c5paper
  }
\fi

\ifshowcrop
% Geometry of the dictionary section
\newcommand\dictionarygeometry{\newgeometry
  { textwidth  = 42\baselineskip
  , textheight = 60\baselineskip
  , hmarginratio = \ifscreen 1:1\else 2:3\fi
  , vmarginratio = 3:2
  , bindingoffset = 0cm
  , twocolumn
  , a4paper 
  , layout=c5paper
  , layouthoffset=\dimexpr(\paperwidth-\csname Gm@layoutwidth\endcsname)/2\relax
   , layoutvoffset=\dimexpr(\paperheight-\csname Gm@layoutheight\endcsname)/2\relax
  , showcrop
  }}
   % Geometry of the cover pages
\newcommand\covergeometry{\newgeometry
  { textwidth  = 42\baselineskip
  , textheight = 60\baselineskip
  , hmarginratio = 1:1
  , vmarginratio = 2:3
  , bindingoffset = 0cm
  , onecolumn
  , layout=c5paper
 , layouthoffset=\dimexpr(\paperwidth-\csname Gm@layoutwidth\endcsname)/2\relax
   , layoutvoffset=\dimexpr(\paperheight-\csname Gm@layoutheight\endcsname)/2\relax
  , showcrop
  }} 
\else
% Geometry of the dictionary section
\newcommand\dictionarygeometry{\newgeometry
  { textwidth  = 42\baselineskip
  , textheight = 60\baselineskip
  , hmarginratio = \ifscreen 1:1\else 2:3\fi
  , vmarginratio = 3:2
  , bindingoffset = 0cm
  , twocolumn,
  , c5paper
  }}
   % Geometry of the cover pages
\newcommand\covergeometry{\newgeometry
  { textwidth  = 42\baselineskip
  , textheight = 60\baselineskip
  , hmarginratio = 1:1
  , vmarginratio = 2:3
  , bindingoffset = 0cm
  , onecolumn
  , c5paper
  }} 
\fi

% ################################################################# LANGUAGES #

\usepackage[icelandic, latin, czech]{babel}
\usepackage{csquotes}


% ################################################################## ENCODING #

\usepackage[utf8]{inputenc}

% Smashing uppercase letters with diacritics
% ------------------ THIS IS EVIL AND SHOULD NEVER BE DONE -------------------
% Don't smash anything if we are on the first run to produce the index
\IfFileExists{figures.idx}{
  % Czech uppercase letters with diacritics:
  \DeclareUnicodeCharacter{00C1}{\strut\smash{\' A}} % Á
  \DeclareUnicodeCharacter{010C}{\strut\smash{\v C}} % Č
  \DeclareUnicodeCharacter{010E}{\strut\smash{\v D}} % Ď
  \DeclareUnicodeCharacter{00C9}{\strut\smash{\' E}} % É
  \DeclareUnicodeCharacter{011A}{\strut\smash{\v E}} % Ě
  \DeclareUnicodeCharacter{00CD}{\strut\smash{\' I}} % Í
  \DeclareUnicodeCharacter{0147}{\strut\smash{\v N}} % Ň
  \DeclareUnicodeCharacter{00D3}{\strut\smash{\' O}} % Ó
  \DeclareUnicodeCharacter{0158}{\strut\smash{\v R}} % Ř
  \DeclareUnicodeCharacter{0160}{\strut\smash{\v S}} % Š
  \DeclareUnicodeCharacter{0164}{\strut\smash{\v T}} % Ť
  \DeclareUnicodeCharacter{00DA}{\strut\smash{\' U}} % Ú
  \DeclareUnicodeCharacter{016E}{\strut\smash{\r U}} % Ů
  \DeclareUnicodeCharacter{00DD}{\strut\smash{\' Y}} % Ý
  \DeclareUnicodeCharacter{017D}{\strut\smash{\v Z}} % Ž
  % Icelandic uppercase letters with diacritics: ÁÉÍÓÚÝ (ÐÞÆ) and
  \DeclareUnicodeCharacter{00D6}{\strut\smash{\" O}} % Ö
}{}
% ------------------ THAT WAS EVIL AND SHOULD NEVER BE DONE ------------------

% Non-breakable space
\DeclareUnicodeCharacter{00A0}{~}

% Non-breakable hyphen
\newcommand*{\nobreakhyphen}{\mbox{-}}
\DeclareUnicodeCharacter{2011}{\nobreakhyphen}


% ############################################################ FONTS, SYMBOLS #

\usepackage[T1]{fontenc}

\usepackage{tgpagella}
\usepackage[scaled=0.88]{helvet}      % relative scale of the two fonts

\def\phvfamily{\fontfamily{phv}\selectfont}

%\usepackage{stmaryrd}
% Let's keep this minimal instead of using the package:
\DeclareSymbolFont{stmry}{U}{stmry}{m}{n}
%\SetSymbolFont{stmry}{bold}{U}{stmry}{b}{n}
\DeclareMathSymbol\shortrightarrow\mathrel{stmry}{"01}
\DeclareMathSymbol\shortuparrow\mathrel{stmry}{"02}

% Methinks these are not necessary
%\usepackage{dingbat}
%\usepackage{manfnt}
%\usepackage{latexsym}

\usepackage{pifont}




% #################################################################### LAYOUT #

% \normalsize should be {8pt}{9.6pt}
\def\HUGE {\fontsize{23.887872pt}{3\baselineskip}\selectfont}
\def\Huge {\fontsize{19.906560pt}{3\baselineskip}\selectfont}
\def\huge {\fontsize{16.588800pt}{3\baselineskip}\selectfont}
\def\LARGE{\fontsize{13.824000pt}{2\baselineskip}\selectfont}
\def\Large{\fontsize{11.520000pt}{2\baselineskip}\selectfont}
\def\large{\fontsize{ 9.600000pt}{2\baselineskip}\selectfont}

% Two columns layout ruler
\setlength\columnsep    {2\baselineskip}
\setlength\columnseprule{0.4pt}

% Necessary for baseline alignment
\topskip=\baselineskip
\raggedbottom
\setlength\parskip{0pt} % it's better to avoid glue

% Temporarily suppress warnings
\hbadness=5000
\vbadness=5000

\setlength\emergencystretch{17pt}

% Allow smaller emergencystretch in several cases
\newenvironment{xtolerant}[2]{%
  \par
  \ifx\empty#1\empty\else\tolerance=#1\relax\fi
  \ifx\empty#2\empty\else\emergencystretch=#2\relax\fi
}{%
  \par
}

% ######################################################### ASSORTED PACKAGES #

\usepackage{tikz}
\usetikzlibrary{calc}

\usepackage{tipa}                      % IPA phonetics

\usepackage{fourier-orns}              % ornaments
\usepackage{amsmath}                   % non-breakable dash
\usepackage{hyphenat}                  % no hyphen in abbreviations

\usepackage{hanging}
\usepackage{fix2col}
\usepackage{fixltx2e}                  % subscript
\usepackage{pagecolor}
\usepackage{afterpage}                 % Used in coloring trick for covers
\usepackage{pdfpages}

\usepackage{multirow}
\usepackage{multicol}
\usepackage{rotating}

\usepackage{etoolbox} % http://ctan.org/pkg/etoolbox
%\usepackage{newunicodechar} 		% new definition of middle dot
\usepackage[ISBN=978-80-260-2385-2,SC0]{ean13isbn}

\usepackage{pgffor}

\usepackage{newunicodechar}

\newunicodechar{·}{\chejnikcdot}
\makeatletter
\newcommand{\chejnikcdot}{%
  \ifnum\lastpenalty=10042
    \kern-0.06667em\relax % dots should always be kerned
  \fi
  \textperiodcentered
  \@ifnextchar|{\kern\chejnikkern\relax}{\penalty10042 }%
}

\newcommand{\closeupdotbar}{%
  \renewcommand\chejnikkern{-0.06667em}%
}
\newcommand\chejnikkern{0pt} % default

% Middle dot + | little bit closer, nicer
%\newunicodechar{·}{\advancedcdot}
%\makeatletter
%\newcommand*\advancedcdot{\textperiodcentered
%  \@ifnextchar|{\hspace{-0.03667em}}{\checkifcdotnext}}
%\newcommand*\checkifcdotnext[2]
%  {\def\tmpa{·}\def\tmpb{#1#2}\ifx\tmpa\tmpb\hspace{-0.06667em}\fi#1#2}
%\makeatother

% ############################################################ BASELINE GRID #

\usepackage{eso-pic}
\usepackage{tikzpagenodes}

% Command to draw a baseline grid
\newcommand\drawbaselinegrid{%
  \begin{tikzpicture}[overlay,remember picture]
    \draw [red!30!white, ultra thin, dashed]
      (0,0) grid [ ystep  = \baselineskip, xstep = \textwidth
                 , shift  = (current page text area.north west)
                 , yshift = -\dp\strutbox
                 ] ++(\textwidth,-\textheight);
    \draw [red!30!white, thin]
      (0,0) grid [ step = \baselineskip, xstep = \textwidth
                 , shift=(current page text area.north)
                 ] ++(0.5\textwidth,-\textheight)
      (0,0) grid [ step = \baselineskip, xstep = \textwidth
                 , shift=(current page text area.north)
                 ] ++(-0.5\textwidth,-\textheight);
    \draw[black!10!white]
      (current page text area.north west)
        rectangle (current page text area.south east);
  \end{tikzpicture}}

% Draw a baseline grid on every page
% \AddToShipoutPicture{\drawbaselinegrid}


% ################################################################## METADATA #

\title{%
  \textbf{Islandsko-český studijní slovník}%
  \thanks{Tato kniha byla vytvořena v \LaTeX{}u pod Ubuntu 14.04.%
    Poděkování patří všem autorům, kteří publikují pod svobodnými licencemi.}}
\author{Aleš Chejn, Ján Zaťko, Jón Gíslason}
\date{říjen 2011}


% ############################################################## BIBLIOGRAPHY #

\usepackage[style=authoryear]{biblatex}
\addbibresource{slovnik.bib}


% ##################################################################### INDEX #

\usepackage{imakeidx}

% Index of authors of photographs
\makeindex
  [ name = figures
  , title = {Seznam autorů fotografií a ilustrací}
  , intoc
  , columnseprule
  , columnsep = \columnsep
  , noautomatic                       % We handle this in make.sh
  ]

% icelandic rules start
\usepackage{filecontents}
\def\mygroup#1{{\bfseries#1}}
\begin{filecontents*}{style.xdy}
;; style.xdy
(markup-letter-group :open-head "~n\mygroup{" :close-head "}")
\end{filecontents*}
% icelandic rules end


% ############################################################### PAGE STYLES #

\usepackage{fancyhdr}

% -------------------------------------- BASIC PAGE STYLE
\fancypagestyle{plain}{
  \fancyhf\relax                       % Clear header/footer
  \renewcommand \headrulewidth {0.0pt} % No header rule
  \renewcommand \footrulewidth {0.0pt} % No footer rule
  \fancyfoot[C]{\thepage}              % Page number in footer, centred
}

% -------------------------------------- INDEX PAGE STYLE  
%\fancypagestyle{indexstyle}{
%  \fancyhf\relax                       % Clear header/footer
%  \renewcommand \headrulewidth {0.4pt} % Thin header rule
%  \renewcommand \footrulewidth {0.0pt} % No footer rule
%  \fancyhead[C]{\thepage}              % Page number in footer, centred
%}

% -------------------------------------- DICTIONARY PAGE STYLE
\fancypagestyle{myheadings}{
  \fancyhf\relax                       % Clear header/footer
  \renewcommand \headrulewidth {0.4pt} % Thin header rule
  \fancyhead[CO,CE]{\thepage}          % Page number in header, centred
  % NOTE: the page numbers will be printed when the dictionary is ready
  \fancyhead[LE,LO]{\phvfamily\bfseries\rightmark}
  \fancyhead[RE,RO]{\phvfamily\bfseries\leftmark}
}

\pagestyle{plain}  

% ##################################################################### LISTS #

\usepackage{expdlist}
\usepackage{enumitem}

\setlist{nosep, topsep=\baselineskip}
\setlist[itemize]{leftmargin=*,label=\scalebox{.6}{\textbullet}}


% #################################################################### FLOATS #

\usepackage{caption}
\captionsetup{labelformat=empty}

% Alter some LaTeX defaults for better treatment of figures:
% See p.105 of "TeX Unbound" for suggested values. 
% See pp. 199-200 of Lamport's "LaTeX" book for details.
% ------------------------------------- Parameters for ALL pages:
\renewcommand \topfraction    {0.9}   %   max fraction of floats at top
\renewcommand \bottomfraction {0.8}   %   max fraction of floats at bottom
% ------------------------------------- Parameters for TEXT (non float) pages:
\setcounter{topnumber}    {1}
\setcounter{bottomnumber} {1}
\setcounter{totalnumber}  {1}         %   2 may work better
\setcounter{dbltopnumber} {2}         %   for 2-column pages
\renewcommand \dbltopfraction   {0.9} %   fit big float above 2-col. text
\renewcommand \textfraction     {0.07}%   allow minimal text w. figs
% ------------------------------------- Parameters for FLOAT (non text) pages:
\renewcommand \floatpagefraction{0.7} %   require fuller float pages
% N.B.: floatpagefraction MUST be less than topfraction !!
\renewcommand{\dblfloatpagefraction}{0.7} % require fuller float pages
% remember to use [htp] or [htpb] for placement

\makeatletter
\setlength{\@fptop}{0pt}
\setlength{\@fpbot}{0pt plus 1fil}
\makeatother

% This kills spacing around floats
\setlength \intextsep        {0pt}
\setlength \floatsep         {0pt}
\setlength \textfloatsep     {0pt}
\setlength \dbltextfloatsep  {0pt}
\setlength \dblfloatsep      {0pt}

% This kills spacing around captions
\setlength \abovecaptionskip {0pt}
\setlength \belowcaptionskip {0pt}


% #################################################################### TABLES #

\usepackage{booktabs}
\usepackage{tabularx}
\usepackage{longtable}
\usepackage{ltxtable}
\usepackage{lscape}

\ifPDFrotate
\else
\makeatletter			       
\renewcommand\PLS@Rotate[1]{}          % this command prevents pdflscape 
\makeatother                           % to rotate pages
\fi

%\LetLtxMacro\oldtabularx\tabularx
%\LetLtxMacro\oldendtabularx\endtabularx
%\renewenvironment{tabularx}[1]
%  {\oldtabularx{\columnwidth}{#2}}
%  {\oldendtabularx}

\newcommand\LTXfw[1]{\LTXtable{\columnwidth}{#1}}

%%% \LetLtxMacro\oldtabularx\tabularx
%%% \LetLtxMacro\oldendtabularx\endtabularx
%%% \newenvironment{longtabularx}[2]
%%%   {\oldtabularx{#1}{#2}}
%%%   {\oldendtabularx}

% ######################################################### czech babel vs cmidrule
% ###################################### http://tex.stackexchange.com/a/112001/8380
\begingroup
  \makeatletter
  \catcode`\-=\active
  \AtBeginDocument{
  \def\@@@cmidrule[#1-#2]#3#4{\global\@cmidla#1\relax
    \global\advance\@cmidla\m@ne
    \ifnum\@cmidla>0\global\let\@gtempa\@cmidrulea\else
    \global\let\@gtempa\@cmidruleb\fi
    \global\@cmidlb#2\relax
    \global\advance\@cmidlb-\@cmidla
    \global\@thisrulewidth=#3
    \@setrulekerning{#4}
    \ifnum\@lastruleclass=\z@\vskip \aboverulesep\fi
    \ifnum0=`{\fi}\@gtempa
    \noalign{\ifnum0=`}\fi\futurenonspacelet\@tempa\@xcmidrule}
  }
\endgroup

% macro that supress cmidrule at page brake
% http://tex.stackexchange.com/questions/110548/using-longtable-with-cmidrule-page-breaking-issues
\makeatletter
\def\LT@output{%
  \ifnum\outputpenalty <-\@Mi
    \ifnum\outputpenalty > -\LT@end@pen
      \LT@err{floats and marginpars not allowed in a longtable}\@ehc
    \else
      \setbox\z@\vbox{\unvbox\@cclv}%
      \ifdim \ht\LT@lastfoot>\ht\LT@foot
        \dimen@\pagegoal
        \advance\dimen@-\ht\LT@lastfoot
        \ifdim\dimen@<\ht\z@
          \setbox\@cclv\vbox{\unvbox\z@
\copy\LT@foot\vss}%
          \@makecol
          \@outputpage
          \setbox\z@\vbox{\box\LT@head}%
        \fi
      \fi
      \global\@colroom\@colht
      \global\vsize\@colht
      \vbox
        {\unvbox\z@
\box\ifvoid\LT@lastfoot\LT@foot\else\LT@lastfoot\fi}%
    \fi
  \else
    \setbox\@cclv\vbox{\unvbox\@cclv

\setbox0\lastbox
\ifdim\ht0=0.29999pt \unskip % hope that was a cline we threw away
\else
\nointerlineskip
\box0     % put it back, whatever it was
\fi

\copy\LT@foot\vss}%
    \@makecol
    \@outputpage
      \global\vsize\@colroom
    \copy\LT@head\nobreak
  \fi}

\makeatother

% These values ensure that for top, mid and bottom rules
% space above + rule width + space below = baselineskip
\setlength \abovetopsep    {0.43\baselineskip}
\setlength \heavyrulewidth {0.10\baselineskip}
\setlength \belowbottomsep {0.43\baselineskip}
\setlength \aboverulesep   {0.47\baselineskip}
\setlength \lightrulewidth {0.06\baselineskip}
\setlength \belowrulesep   {0.47\baselineskip}

%\usepackage{letltxmacro}

%%% \LetLtxMacro\oldtabular\tabular
%%% \LetLtxMacro\oldendtabular\endtabular
%%% \renewenvironment{tabular}[2][t]
%%%   {\oldtabular[#1]{#2}}
%%%   {\oldendtabular}

%%% \LetLtxMacro\oldtable\table
%%% \LetLtxMacro\oldendtable\endtable
%%% \renewenvironment{table}[1][h]
%%%   {\leavevmode\oldtable[#1]\vspace{\dp\strutbox}}
%%%   {\vspace{-\dp\strutbox}\oldendtable}
%%%   % This aligns a [t] aligned tab-envo to baseline grid


%%% % declention and conjugation tables
%%% \usepackage{floatrow}
%%% \DeclareFloatFont{footnotesize}{\footnotesize}
%%% % "scriptsize" is defined by floatrow, "tiny" not
%%% \floatsetup[table]{font=footnotesize}

% ################################################################### FIGURES #

\usepackage{adjustbox} % loads also graphicx

\graphicspath{%
  {/home/chejnik/Dokumenty/hvalur.org/images/biolib/full/}%
  {/home/chejnik/Dokumenty/hvalur.org/images/uploaded_files/}}

\newlength\dicfigheight
\newlength\dicfigskip

\newcommand\dicFigure[3]{ % {filename}{caption}{index entry}
  \setlength\fboxsep{0pt}\setlength\fboxrule{0.5pt}
  \def\dicfigbox{\fbox{%
    \adjincludegraphics [ max height = 0.8\columnwidth
                        , max width  = 0.8\columnwidth ] {#1}}}
  \setlength\abovecaptionskip {\dp\strutbox}
  \setlength\belowcaptionskip {\ht\strutbox} 
  \setlength\dicfigheight {\heightof\dicfigbox+\fboxrule}
  \setlength\dicfigskip
    {\baselineskip*\numexpr1+\dicfigheight/\baselineskip\relax-\dicfigheight}
  \begin{figure}[ht]
    \centering\vspace*{\dicfigskip}\dicfigbox\caption{#2}\index[figures]{#3}
  \end{figure}
}


% #################################################################### COLORS #

\usepackage{color}

% http://www.colorschemer.com/online.html
% main color - grammar color
% #66023C
\definecolor {darkgreen}          {rgb} {0.40, 0.01, 0.24}

\definecolor {royalazure}         {rgb} {0.00, 0.22, 0.66}
\definecolor {brown}              {rgb} {0.40, 0.01, 0.24}

% COLORS FOR THUMB INDEXES
\definecolor {babyblueeyes}       {rgb} {0.63, 0.79, 0.95}
\definecolor {unitednationsblue}  {rgb} {0.36, 0.57, 0.90}
\definecolor {blue(ryb)}          {rgb} {0.01, 0.28, 1.00}
\definecolor {darkblue}           {rgb} {0.00, 0.00, 0.55}
\definecolor {screamingreen}      {rgb} {0.46, 1.00, 0.44}
\definecolor {limegreen}          {rgb} {0.20, 0.80, 0.20}
\definecolor {islamicgreen}       {rgb} {0.00, 0.56, 0.00}
\definecolor {upforestgreen}      {rgb} {0.00, 0.27, 0.13}
\definecolor {icterine}           {rgb} {0.99, 0.97, 0.37}
\definecolor {orange(colorwheel)} {rgb} {1.00, 0.50, 0.00}
\definecolor {orange-red}         {rgb} {1.00, 0.27, 0.00}
\definecolor {oucrimsonred}       {rgb} {0.60, 0.00, 0.00}
\definecolor {cottoncandy}        {rgb} {1.00, 0.74, 0.85}
\definecolor {orchid}             {rgb} {0.85, 0.44, 0.84}
\definecolor {vividcerise}        {rgb} {0.85, 0.11, 0.51}
\definecolor {patriarch}          {rgb} {0.50, 0.00, 0.50}

\definecolor {title}              {RGB} {  16,   13,   32}
\definecolor {golden}             {RGB} { 241,  184,   45}
\definecolor {freqstar}           {RGB} {  91,    3,   99}
\definecolor {newdev}             {RGB} {   3,   11,   99}

\definecolor {color1}             {RGB} { 182,   86,    0}
\definecolor {color2}             {RGB} { 143,    9,    6}
\definecolor {color3}             {RGB} {   3,   23,  118}
\definecolor {color4}             {RGB} {   0,   82,  168}
\definecolor {color5}             {RGB} {   0,   85,  142}
\definecolor {color6}             {RGB} {   0,  115,  162}
\definecolor {color7}             {RGB} {  34,  146,  186}
\definecolor {color8}             {RGB} {  40,  159,  153}
\definecolor {color9}             {RGB} {   0,  125,  111}
\definecolor {color10}            {RGB} {   4,  107,   60}
\definecolor {color11}            {RGB} {  71,  134,   81}
\definecolor {color12}            {RGB} { 109,  134,   42}
\definecolor {color13}            {RGB} { 205,  194,   18}
\definecolor {color14}            {RGB} { 204,  162,   24}

%\definecolor {color1}             {RGB} { 239,  222,  205}
%\definecolor {color2}             {RGB} { 251,  206,  177}
%\definecolor {color3}             {RGB} { 251,  185,  142}
%\definecolor {color4}             {RGB} { 248,  166,  112}
%\definecolor {color5}             {RGB} { 216,  231,  246}
%\definecolor {color6}             {RGB} { 186,  215,  244}
%\definecolor {color7}             {RGB} { 152,  197,  243}
%\definecolor {color8}             {RGB} { 118,  181,  245}
%\definecolor {color9}             {RGB} { 181,  249,  185}
%\definecolor {color10}            {RGB} { 160,  248,  167}
%\definecolor {color11}            {RGB} { 128,  247,  136}
%\definecolor {color12}            {RGB} {  87,  248,   98}
%\definecolor {color13}            {RGB} { 250,  241,  138}
%\definecolor {color14}            {RGB} { 250,  240,  111}
%\definecolor {color15}            {RGB} { 244,  239,   85}
%\definecolor {color16}            {RGB} { 245,  239,   57}

\definecolor {lightgrey}          {RGB} { 105,  105,  105}
\definecolor {darkgreen_real}     {rgb} {0.00, 0.50, 0.00}


% ############################################################# THUMB INDEXES #

% Thumb indexes' colors
%\newcommand\BoxColor{%
%  \ifcase\theletternum darkgreen!30\or babyblueeyes\or unitednationsblue\or blue(ryb)\or screamingreen\or limegreen\or islamicgreen\or upforestgreen\or icterine\or orange(colorwheel)\or orange-red%
%  \or oucrimsonred\or cottoncandy\or orchid\or vividcerise\or patriarch\or babyblueeyes\or unitednationsblue\or blue(ryb)\or screamingreen\or limegreen\or islamicgreen\or upforestgreen\or icterine\or orange(colorwheel)\or orange-red%
%  \or oucrimsonred\or cottoncandy\or orchid\or vividcerise\or patriarch\else darkgreen!30\fi}
  
% Thumb indexes' colors
%\newcommand\BoxColor{%
%  \ifcase\theletternum darkgreen!30\or color1\or color2\or color3\or color4\or color5\or color6\or color7\or color8\or color9\or color10%
%  \or color11\or color12\or color13\or color14\or color1\or color2\or color3\or color4\or color5\or color6\or color7\or color8\or color9%
%  \or color10\or color11\or color12\or color13\or color14\or color1\or color2\or color3\or color4\or color5\or color6\else color7\or color8\or color9\or color10\or color11\or color12\fi}

%\def\BoxColor#1{darkgreen!\the\numexpr103-#1-#1-#1\relax!blue}

\newcommand\BoxColor[1]{%
\ifcase#1 darkgreen!30\or color13\or color14\or color1\or color2\or color3\or color4\or color5\or color6\or color7\or color8\or color9\or color10%
\or color11\or color12\or color13\or color14\or color1\or color2\or color3\or color4\or color5\or color6\or color7\or color8\or color9%
\or color10\or color11\or color12\or color13\or color14\or color1\or color2\or color3\or color4\or color5\or color6\else color7\or color8\or color9\or color10\or color11\or color12\fi}

\ifshowcrop
%%%%%%%%%%%%%%%%%%%%%%%%%%%%%%%%%%%%%%%%%%%%%%%%%%%%%%%%%%%%%%%%%%%%%%%%%%%%%%%%

\newcounter{letternum}
\newcounter{lettersum}                 \setcounter{lettersum}{34}
% these margins separate top thumb from top of page and last thumb from
% bottom of page. BEWARE, it's from top of the (cut out) page, not the folio!
\newlength{\thumbtopmargin}            \setlength{\thumbtopmargin}{1cm}
\newlength{\thumbbottommargin}         \setlength{\thumbbottommargin}{1cm}
\newlength{\thumbwidth}                \setlength{\thumbwidth}{2.20cm}
\newlength{\thumbheight}
\pgfmathsetlength{\thumbheight}%
  {(\csname Gm@layoutheight\endcsname-\thumbtopmargin-\thumbbottommargin)/\value{lettersum}}
% inner xsep of letter from box margin
\newlength{\thumbsep}                  \setlength{\thumbsep}{.30cm}
% overlap over cutting line
\newlength{\thumboverlap}              \setlength{\thumboverlap}{3cm}

\tikzset{
  thumb/.style={
    text=white,
    outer sep=0pt,
    inner xsep=\thumbsep,
    minimum height=\thumbheight,
    text width=\thumbwidth-2\thumbsep,
    font=\sffamily\bfseries,},
  thumb east/.style={thumb,
    xshift=\dimexpr(+\thumboverlap-\paperwidth+\csname Gm@layoutwidth\endcsname)/2\relax,
    align=left,anchor=north east,},
  thumb west/.style={thumb,
    xshift=\dimexpr(-\thumboverlap+\paperwidth-\csname Gm@layoutwidth\endcsname)/2\relax,
    align=right,anchor=north west,},
}

\def\thumbnew{}
\def\thumbold{}
\usepackage{everypage}
\AddEverypageHook{\if\relax\thumbnew\relax\xdef\thumbnew{\thumbold}\fi}

\def\ethumbs#1,#2\relax%
  {\if\relax#1\relax\else\eventhumb{#1}\fi%
   \if\relax#2\relax\else\ethumbs#2\relax\fi%
   \gdef\thumbnew{}%
   \gdef\thumbold{#1,}}

\def\othumbs#1,#2\relax%
  {\if\relax#1\relax\else\oddthumb{#1}\fi%
   \if\relax#2\relax\else\othumbs#2\relax\fi%
   \gdef\thumbold{#1,}%
   \gdef\thumbnew{}}

\newcommand{\drawthumb}[2]{%
  \begin{tikzpicture}[remember picture, overlay]
    \node [thumb #2, fill=\BoxColor{#1}]
       at ($(current page.north #2)-%
            (0,\csname Gm@layoutvoffset\endcsname+\thumbtopmargin-\thumbheight+%
               #1*\thumbheight)$) {\csname Let#1\endcsname};
   \end{tikzpicture}}



%%%%%%%%%%%%%%%%%%%%%%%%%%%%%%%%%%%%%%%%%%%%%%%%%%%%%%%%%%%%%%%%%%%%%%%%%%%%%%%%
\else
% THUMB INDEXES
% new counter to hold the current number of the letter to determine the vertical position
\newcounter{letternum}
% newcounter for the sum of all letters to get the right height of a box
\newcounter{lettersum}
\setcounter{lettersum}{34}
% some margin settings
\newlength{\thumbtopmargin}
\setlength{\thumbtopmargin}{\ifshowcrop 3cm\else 1.5cm\fi} %1cm
\newlength{\thumbbottommargin}
\setlength{\thumbbottommargin}{\ifshowcrop 6cm\else 2cm\fi} %2.5cm
% calculate the box height by dividing the page height
\newlength{\thumbheight}
\pgfmathsetlength{\thumbheight}{%
(\paperheight-\thumbtopmargin-\thumbbottommargin)%
/%
\value{lettersum}
}
% box width
\newlength{\thumbwidth}
\setlength{\thumbwidth}{\ifshowcrop 2cm\else 0.5cm\fi} %0.5cm
% style the boxes
\tikzset{
thumb/.style={
   text=white,
   minimum height=\thumbheight,
   text width=\thumbwidth,
   outer sep=0pt,
   font=\sffamily\bfseries,
 }
 }

\def\thumbnew{}
\def\thumbold{}
\usepackage{everypage}
\AddEverypageHook{\if\relax\thumbnew\relax\xdef\thumbnew{\thumbold}\fi}

\def\ethumbs#1,#2\relax{\if\relax#1\relax\else\eventhumb{#1}\fi%
                        \if\relax#2\relax\else\ethumbs#2\relax\fi%
                        \gdef\thumbnew{}%
                        \gdef\thumbold{#1,}%
}

\def\othumbs#1,#2\relax{\if\relax#1\relax\else\oddthumb{#1}\fi%
                        \if\relax#2\relax\else\othumbs#2\relax\fi%
                        \gdef\thumbold{#1,}%
                        \gdef\thumbnew{}%
}

\newcommand{\drawthumb}[2]{%
  % see pgfmanual.pdf for more information about this part
  \begin{tikzpicture}[remember picture, overlay]
    \node [thumb, fill=\BoxColor{#1}, text centered, anchor=north #2]
       at ($(current page.north #2)-%
%            (0,\thumbtopmargin+\value{letternum}*\thumbheight)%
            (0,\thumbtopmargin+#1*\thumbheight)$) {\csname Let#1\endcsname};
   \end{tikzpicture}}
\fi

\newcommand{\oddthumb} [1]{\drawthumb{#1}{west}}
\newcommand{\eventhumb}[1]{\drawthumb{#1}{east}}

\newcommand{\lettergroup}[1]%
  {\refstepcounter{letternum}%
   \expandafter\gdef\csname Let\theletternum\endcsname{#1}%
   \xdef\thumbnew{\theletternum,\thumbnew}%
   \fancyhead[LO]{\phvfamily\bfseries\rightmark%
     \expandafter\ethumbs\thumbnew\relax\relax\relax}%
   \fancyhead[RE]{\phvfamily\bfseries\leftmark%
     \expandafter\othumbs\thumbnew\relax\relax\relax}}
     
% ############################################################## LOCALIZATION #

% renames the index name
%\addto\captionsczech{%
%  \renewcommand\indexname{{Seznam autorů fotografií a ilustrací}}}

% renames the contents name
\addto\captionsczech{%
  \renewcommand\contentsname{Obsah}}

\renewcommand \contentsname {Obsah}
\renewcommand \chaptername  {Kapitola}


% ################################################################ SECTIONING #
\usepackage[explicit]{titlesec}

\setlength{\beforetitleunit}{\baselineskip}
\setlength{\aftertitleunit} {\baselineskip}

\titlespacing*{\chapter}   {0em}{*0}{*3} % Remembere there's a topskip too
\titlespacing*{\section}   {0em}{*1}{*2}
\titlespacing*{\subsection}{0em}{*1}{*2}

\newcommand\ghostdrop[2]{%
  \raisebox{-#1\baselineskip}[0pt][0pt]{#2}}

\titleformat {\chapter} {} {} {0pt} % {#1}
  {\ghostdrop{1.0}{\parbox[t]{\columnwidth}
     {\LARGE\bfseries\centering\MakeUppercase{#1}}}}

\def\longchapterskip{\vspace{2\baselineskip}}
\def\longsectionskip{\vspace{2\baselineskip}}

\titleformat {\section} {\bfseries} {} {0pt}
  {\ghostdrop{1.0}{\parbox[t]{\columnwidth}
     {\Large\bfseries\centering#1}}}

\titleformat {\subsection} {\bfseries} {} {0pt}
  {\ghostdrop{1.0}{\hspace{1em}\parbox[t]{\columnwidth-1em}
     {\large\bfseries#1}}}


\setlength\multicolsep{0pt}
%\BeforeBeginEnvironment{multicols}{\vspace{-\dp\strutbox}}


% ######################################################## TABLE OF CONTENTS #
\usepackage{titletoc}

\titlecontents {chapter} [0pt] {\addvspace\baselineskip\bfseries}
  {} {} {\titlerule*[2\baselineskip]{}\contentspage}

\titlecontents {section} [\baselineskip] {} {} {}
  {\titlerule*[1\baselineskip]{.}\contentspage}


% #################################################### ASSORTED CUSTOM MACROS #

\def\textIS#1{\foreignlanguage{icelandic}{#1}}
\def\textCS#1{\foreignlanguage{czech}{#1}}
\def\textLA#1{\foreignlanguage{latin}{#1}}

\usepackage{accsupp}

\makeatletter                                                                                                          
  \def\ACCSUPP@bdc{\special{pdf:code \ACCSUPP@span BDC}}                                                               
  \def\ACCSUPP@emc{\special{pdf:code EMC}}                                                                             
\makeatother

\usepackage{pgffor}

\newcommand\blspace[1][1]{\vspace{#1\baselineskip}}

\usepackage{xparse}


\newcommand \dicsymFrequent  {\raisebox{-.2ex}{\color{darkgreen}\ding{167}}}
\newcommand \dicsymSee       {$\shortrightarrow$}
\newcommand \dicsymCompare   {$\shortuparrow$}
\newcommand \dicsymIdiom     {$\diamondsuit$}
\newcommand \dicsymExampleIS {$\triangleright$}
\newcommand \dicsymExampleCS {\guilsinglright}
\newcommand \dicsymProverb   {{\color{newdev}\ding{96}}}
\newcommand \dicsymPhoto
  {\includegraphics[keepaspectratio,height=0.65em]{photo_icon}}

%\NewDocumentCommand\dicTerm{mo}
 % {{\phvfamily\bfseries\textIS{#1}}\IfValueT{#2}{, \phvfamily\bfseries\textIS{#2}}}

\NewDocumentCommand\dicEntry{o}
  {\par\hangpara{0.25\baselineskip}{1}%
   \IfValueT{#1}{\markboth{#1}{#1}}%
    \ifPDFsearchvalue
    \settowidth{\eyja}{#1}%
  \makebox[\eyja]{\color{white}#1}\hspace{-\eyja}%
  \fi
   \ignorespaces}
   
   \NewDocumentCommand\dicEntryGuide{o}
  {\par\hangpara{0.29\baselineskip}{1}%
   \IfValueT{#1}{\markboth{#1}{#1}}%
    \ifPDFsearchvalue
    \settowidth{\eyja}{#1}%
  \makebox[\eyja]{\color{white}#1}\hspace{-\eyja}%
  \fi
   \ignorespaces}
   
\newcommand\dicSubEntry
  {\\\hspace*{-\hangindent}}

\NewDocumentCommand\textAlt{mm}
  {\BeginAccSupp{method=pdfstringdef,ActualText={#2}}#1\EndAccSupp{}}

\newcounter{terms}

\NewDocumentCommand\dicStyleIntTerm{m}{\textIS{{\phvfamily\bfseries#1}}}
\NewDocumentCommand\dicStyleExtTerm{m}{\textIS{\textit{#1}}}

   \NewDocumentCommand\dicTerm{mo}
  {{\closeupdotbar \phvfamily\bfseries\textIS{#1}}\IfValueT{#2}{, \phvfamily\bfseries\textIS{#2}}}

  \newcommand\dicTermList[1]
  {\foreach \t [count=\n] in {#1}{\ifnum\n=1\else, \fi\dicTerm{\t}}}
  
\def\dicIPA#1{\textipa{[#1]}}

\NewDocumentCommand\dicPos{smo}
  {{\color{darkgreen}\textbf{\small#2}%
   \IfValueT{#3}{\nobreak\textsubscript{#3}}}}

%\NewDocumentCommand\dicFlx{smo}
 % {{\color{darkgreen}\footnotesize#1\IfValueT{#2}{\textsubscript{#2}}}}

  \NewDocumentCommand\dicFlx{smo}
  {{\color{darkgreen}\footnotesize#2%
   \IfValueT{#3}{\nobreak\textsubscript{#3}}}}
  
\newcommand\dicExampleCS[1]{\dicsymExampleCS~\textit{#1}}
\newcommand\dicExampleIS[1]{\dicsymExampleIS~\textIS{\textit{#1}}}

\NewDocumentCommand\dicLink{sm}
  {\IfBooleanF#1{\dicsymSee\nobreak\hspace{.08667em plus .03333em}}\nobreak\hspace{.03667em}\dicTerm{#2}}
\NewDocumentCommand\dicSynonym{sm}
  {(\IfBooleanF#1{\dicsymSee\nobreak\hspace{.08667em plus .03333em}}\nobreak\hspace{.03667em}\textIS{\textit{#2}})}
\NewDocumentCommand\dicAntonym{sm}
  {(\IfBooleanF#1{\dicsymCompare\nobreak\hspace{.08667em plus .03333em}}\nobreak\hspace{.03667em}\textIS{\textit{#2}})}
\NewDocumentCommand\dicCompare{sm}
  {(\IfBooleanF#1{\dicsymCompare\nobreak\hspace{.08667em plus .03333em}}\nobreak\hspace{.03667em}\textIS{\textit{#2}})}


\newcommand\dicPhraseIS[1]{\textIS{\textbf{#1}}}

\NewDocumentCommand\dicIdiom{mo}
  {\dicSubEntry{\color{newdev}\dicStyleIntTerm{#1} %
     \IfValueTF{#2}{+ \dicStyleIntTerm{#2} \markboth{#1 + #2}{#1 + #2}}{\markboth{#1}{#1}}}\dicsymIdiom}
     
\newcommand\dicProverb
  {\dicSubEntry\dicsymProverb}



\newcommand\dicLangCat[1]{{\footnotesize#1}}
\newcommand\dicFieldCat[1]{{\footnotesize#1}}


\newcommand\dicDirectTranslationCS[1]{#1}
\newcommand\dicIndirectTranslationCS[1]{{\footnotesize#1}}

\usepackage{placeins}

\newcommand*{\dicLetter}[2]{%
  \FloatBarrier\newpage\phantomsection\pdfbookmark[1]{#1}{#2}%
  \noindent\parbox[b][9\baselineskip][c]{\columnwidth}
    {\centering\HUGE\strut\smash{\MakeUppercase{#1}~\MakeLowercase{#1}}}%
      %\par\xdef\tpd{\the\prevdepth}
      %\par\prevdepth\tpd%
  \expandafter\lettergroup{\MakeLowercase{#1}}}

  
%tabulky
\newcommand{\specialcell}[2][l]{%
\begin{tabular}[#1]{@{}l@{}}#2\end{tabular}}


%\newcommand{\devisionguide}[1]{\hspace*{0.3cm}{{{{\foreignlanguage{icelandic}{\color{newdev}{\phvfamily{\textbf{#1}}}}}}}}}
%\newcommand{\devisionguideindent}[1]{\hspace*{-0.1cm}{{{{\foreignlanguage{icelandic}{\color{newdev}{\phvfamily{\textbf{#1}}}}}}}}}

%\newcommand\n{$n$\nobreakdash-\hspace{0pt}} %nonbreakable dash in grammar endings
%\def \nobreakseq {\nobreak \hskip 0pt \hbox} %nolinebreak in case like (-s, -)

% rule line
\newcommand{\HRule}{\rule{\linewidth}{0.1mm}}

% width of pdf found box
\newlength\eyja

%\newcommand\strutsmash[1]{\strut\smash{#1}}


% Dictionary symbols


% Command synonyms
%\let\freqstar\dicsymFrequent
%\let\photoicon\dicsymPhoto

\newcommand{\vstretch}[1]{\vspace*{\stretch{#1}}}

\newcommand*{\addthin}{\hskip0.16667em\relax}
\newcommand*{\addthinS}{\hskip0.06667em\relax}
\newcommand*{\addthinM}{\hskip0.10007em\relax}
\newcommand*{\addthinSS}{\hskip0.00007em\relax}
% #################################################################### COVERS #

\def\dictnameCZ{islandsko-český studijní slovník}
\def\dictnameIS{íslensk-tékknesk stúdentaorðabók}

\newcommand*{\frontcoverimages}{%
  ds_image_blaberjalyng_0_2.jpg,
  ds_image_lundi_0_1.jpg,
  ds_image_hestur_0_1.jpg,
  ds_image_hrafnafifa_0_2.jpg,
  ds_image_baldursbra_0_1.jpg,
  20948.jpg,
  ds_image_hvalur_0_1.jpg,
  ds_image_blodberg_0_1.jpg}

\newcommand*{\backcoverimages}{%
  ds_image_spoi_0_1.jpg,
  ds_image_mosalyng_0_1.jpg,
  ds_image_blaklukkulyng_0_1.jpg,
  20787.jpg,
  ds_image_ufsi_0_2.jpg,
  %ds_image_gullbra_0_1.jpg,
  ds_image_vatnsberi_0_1.jpg,
  ds_image_islandsfifill_0_2.jpg,
  21137.jpg}

\newcommand\makecoverwith[1]{
 \pagecolor{title}\afterpage\nopagecolor
 \begin{center}
 \vspace*{0.3cm}%{\fill}
 {\Huge\scshape\color{white}%
   \dictnameCZ\\\dictnameIS}
 \vspace*{1em}                        % I think this is necessary
 \begin{figure}[ht]\centering%
   \setlength\fboxsep{0pt}\setlength\fboxrule{0.5pt}%
   \foreach \file in #1{
     \fbox{\includegraphics[width=5.5cm]{\file}}}
 \end{figure}
 \vspace*{\fill}
 \end{center}
}


% ###################################################### HYPERREF & PDF INFO #

% unicode neccessary so that national characters in hypersetup appear ok
\usepackage[pdftex, unicode, hyperfootnotes=false]{hyperref}

% hyperlinks in black
\makeatletter
\let\Hy@linktoc\Hy@linktoc@none
\makeatother

\hypersetup
  { pdftitle={Islandsko-český studijní slovník}
  , pdfauthor={Aleš Chejn, Jón Gíslason, Ján Zaťko}
  , pdfsubject={Islandsko-český studijní slovník}
  , pdfkeywords={Islandsko-český studijní slovník%
               , islandština, čeština, slovník}
  , bookmarks=true
  , colorlinks=true
  , citecolor=black
  , urlcolor=\ifscreen darkgreen\else black\fi
  , linkcolor=black}


% ################################################################# MICROTYPE #

\usepackage[final]{microtype}

% ############################################################### HYPHENATION #

% this allows to include brackets inside hyphenation rules f.e (pro)-dis-ku-to-vat
\lccode`\(`\(
\lccode`\)`\)
\lccode`\[`\[
\lccode`\]`\]

\begin{hyphenrules}{icelandic}
\hyphenation{
  frum-semja
  burm-neskur
  ó-breytan-le-gur
  pla-tó-nskt
  skatta-fram-tal
  skrúf-gang-ur
  strand-rauð-viður
  prent-sver-ta
  grun-nur
  uppi-skrop-pa
  ger-vi-mennska
  sö-lu-stað-ur
  sý-na-gó-ga
  mark-mið
  trau-st-ur
  ís-kyggi-leg-ur
  brjó-sta-höld
  kné-krjú-pa
  nel-li-ka
  kaffi-kvörn
  margs-konar
  skrúf-gang-ur
  uppi-skrop-pa
  veð-rað-ur
  vest-ur-átt
  mis-mun-andi
  Bar-ba-do-si
  tung-u-mál
  au-ga-brún
  há-sker-pu-sjón-varp
  rúnn-stykki
  ger-vi-tungl
  yfir-borðs-menn-ska
  gæ-fu-mun-ur
  hlí-fi-skil-di
  á-vís-un
  ó-um-flýj-an-leg-ur
  san-mar-ín-sk-ur
  gæ-ða-stað-all
}
\end{hyphenrules}
\begin{hyphenrules}{czech}
\hyphenation{
  teo-rie
  mi-liardy
  (z)-od-po-věd-nost
  tan-tié-ma
  alian-ce
  pe-dia-tr(ička)
  pa-cienty
  zá-chra-nář(ka)
  dia-man-to-vá
  poe-zie
  ma-te-riál-ní
  egoi-s-ta
  egoist-ka
  nai-vi-ta
  ne-upřím-nost
  pě-ti-úhel-ník
  še-s-ti-úhel-ník
  za-šmo-dr-chat
  ar-cheo-lo-gic-ký
  ne-omlu-vi-tel-ný
  (vy)-pro-du-ko-vat
  dia-gno-s-ti-ko-vat
  teo-lož-ka
  ateist-ka
  oceá-nu
  im-pe-ria-li-s-ta
  (vy)-zdvih-nout
  ne-úměr-ný
  kon-ti-nuál-ní
  neu-t-rum
  ne-opráv-ně-ně
  vig-vam
  pro-zaic-ký
  geo-lož-ka
  seiz-mic-ká
  teo-rém
  olean-dr
  poe-ta
  (vy)-fo-to-gra-fo-vat
  (vy)-ven-ti-lo-vat
  pneu-mo-nie
  Lu-cem-bur-čan-(ka)
  ma-te-riály
  am-bi-cióz-ní
  Myan-ma(r)
  Myan-ma-řan-(ka)
  mol-dav-šti-na
  si-tua-cí
  ne-ukáz-ně-ný
  he-te-ro-trof-ní
  ne-omlu-vi-tel-ná
  ne-upřím-ný
  dvou-bliz-ný
  ne-úmy-sl-ný
  ne-ustá-lý-mi
  ne-ohle-du-pl-ný
  ne-osob-ní
  ne-oso-le-ný
  ne-uvě-ři-tel-ný
  ne-ohro-že-ný
  ne-oby-čej-ný
  ne-útul-ný
  ne-únav-ný
  ne-ože-nit
  po-ukáz-ka
  za-úto-čit
  aero-dy-na-mic-ký
  hru-bián-ství
  (pro)-di-s-ku-to-vat
  (vy)-ča-lou-nit
  (na)-oč-ko-vá-ní
  Bhú-tán-ka
  (za)-šně-ro-vat
  soud-ní
  re-vo-lu-cio-nář(ka)
  á-zer-báj-d-žán
  čtyř-hran-ný
  ze-m-dle-lý
  klou-zač-ka
  Vest-mann-ské
  tří-stran-ný
  (za)-bu-rá-ce-ní
  os-tud-ný
  pro-du-cent-(ka)
  chvá-s-tal-(ka)
  apar-theid
  asis-ten-ce
  asis-tent-(ka)
  (u)-či-nit
  (u)-za-mknout
  past-vi-na
  (za)-zá-lo-ho-vat
  Af-ri-čan-(ka)
  (za)-an-ga-žo-vá-ní
  se-rióz-ní
  opo-nent-(ka)
  pro-tiv-ník
  zvlášt-ní
  (po)-div-ný
  ho-mo-se-xuál-ní
  ko-alič-ní
  si-tua-ce
  si-tua-cí
  egoi-s-mus
  spe-ciál-ní
  spe-cia-li-zo-va-ný
  spe-cia-li-zo-va-né
  sobě-stač-ný
  sa-mo-zvo-le-ný
  sa-mo-čin-ně
 chvá-s-ta-vost
 klient-(ka)
 ku-rióz-ní
 (na)-šrou-bo-vat
 (za)-šrou-bo-vat
 krea-ti-vi-ta
 pneu-ma-ti-ka
 toa-le-ta
 so-cia-list-ka
 (vy)-dez-in-fi-ko-va-ný
 vše-uměl
 vše-uměl-(kyně)
 (pře)-škrt-nout
 (vy)-svlék-nout
 vy-učo-va-cí
 vy-učo-va-cích
 (vy)-chla-s-tat
 (vy)-te-to-vat
 pa-cient-(ka)
 ne-opod-s-tat-ně-ný
 de-se-ti-úhel-ník
 pe-rio-di-kum
 (za)-mum-lá-ní
 ne-smě-lá
 ne-dů-vě-ři-vá
 tee-na-ger-(ka)
 (z)-de-ge-ne-ro-vá-ní
 mo-ci-chti-vý
 bez-úhon-nost
 zne-uctít
 vo-dě-odol-ný
 we-bo-vý
 web-de-sign
 in-že-nýr-(ka)
 zá-pa-do-is-land-ský
 klient-(ka)
 za-uja-tý
 vy-uží-va-né
 vy-uží-vat
 (za)-re-gi-s-t-ro-vat
 (z)-kom-pri-mo-va-ný
 na-cio-na-li-s-mus
 ma-k-ro-eko-no-mic-ká
 bez-bo-lest-ný
 viol-ka
 troj-úhel-ník
 (za)-bu-rá-ce-ní
 ne-ústup-nost
 (vy)-cvi-če-ný
 zprá-va
 ge-ria-tr-(ička)
 osud-ný
 ná-pra-vy
 vstup-ní
 obec-ně
 re-pub-li-ka
 (po)-div-ný
 pa-slun-ce
 vliv-ný
 ba-ham-ský
 pe-dia-tr-(ička)
 ha-sič-(ka)
 vnuč-ka
 bro-s-kvoň
 lo-di-vod-(ka)
 mi-s-tr-(yně)
 in-že-nýr-(ka)
 dr-ži-tel-(ka)
 to-pič-(ka)
 pro-mi-nent-ní
 takt-ní
 frag-men-tár-ní
 po-zo-ru-hod-ný
 úspě-š-ný
 trans-fu-ze
 bur-žoa-zie
 pře-stro-je-ní
 oh-ni-vzdor-ný
 dis-tan-ční
 de-frau-da-ce
 fi-s-kál-ní
 zběh-lý
 kvap-ný
 střed-ní
 sa-mo-hlá-s-ka
 pře-vo-di-tel-ný
 zban-k-ro-to-va-ný
 zá-clo-na
 zná-mý
 zá-klad-ní
 na-uka
 ro-š-ťa-č-ka
 lí-bán-ky
 sou-du
 zkou-š-ka
 mlý-nek
 prá-šiv-ka
 kul-ma
 la-hod-ný
 zá-meč-ni-ce
 ne-vý-znam-ný
 ovzdu-ší
 mno-ho-úhel-ní-ko-vý
 žra-vý
 zcest-ný
 esen-ce
 (při)-nu-ce-ný
 skr-b-lic-ký
 skle-s-lý
 bez-bož-ný
 ne-opráv-ně-ný
 ošu-mě-lý
 pru-der-ní
 chvá-s-ta-vý
 sou-drž-ný
 van-dal-(ka)
 sa-bo-tér-(ka)
 (po)-div-ně
 mo-k-řad
 plyt-ký
 na-spo-du
 izo-la-ce
 le-toun
 za-han-be-ný
 ří-ze-ní
 kom-plex-ní
 pr-kén-ko
 ba-vič-(ka)
 plyt-ký
 ne-mo-tor-ně
 šéf-(ka)
 přes-čas
 zji-s-ti-tel-ný
 ta-ko-vý
 kňou-ral-(ka)
 se-me-s-tr
 po-pra-s-ka-ný
 han-ba
 vzpří-me-ný
 chys-ta-ný
 lik-vi-da-ce
 or-ga-ni-za-ce
 ri-s-kant-ní
 (za)-pís-kat
 (za)-hvíz-dat
 vrou-cí
  }
\end{hyphenrules}


% ############################################################# CZECH TYPOGRAPHY #
% https://math.feld.cvut.cz/tkadlec/ftp/text/textypo.pdf



% ################################################################# DOCUMENT #

\begin{document}
\pagenumbering{gobble}

% -------------------------------------------------------------- Front cover -

\ifmakecovers
\covergeometry
  \makecoverwith\frontcoverimages
\restoregeometry  
\fi


% ============================================================= FRONT MATTER =

\ifmakefrontmatter

% ------------------------------------------------------------------ Patitul -

\cleardoublepage

\begin{center}
  %\vspace*{2cm}{\fontsize{20}{20}\selectfont\scshape\dictnameCZ\\\dictnameIS}
  \vspace*{\stretch 1}
  {\Huge\scshape\dictnameCZ\\\dictnameIS}
  \vspace{\stretch 4}
\end{center}

% -------------------------------------------------------------------- Titul -

\cleardoublepage

\begin{center}
  %\vspace*{2cm}{\fontsize{24}{24}\selectfont\scshape\dictnameCZ\\\dictnameIS}
  \vspace*{\stretch 1}
  {\HUGE\scshape\dictnameCZ\\\dictnameIS}
  \vspace{\stretch 2}
\end{center}

% ----------------------------------------------------------------- Impresum -

\clearpage

\noindent\begin{minipage}[t][\textheight][t]{\textwidth}
Autoři \textit{Islandsko-českého studijního slovníku}:\\
  \textbf{Aleš Chejn, Ján Zaťko, Jón Gíslason}\\[\baselineskip]
\ifPDF
\else
Autor \textit{Islandské gramatiky}:\\
  \textbf{Vojtěch Kupča}\\[\baselineskip]
\fi
Poděkování:\\
  \textbf{Dorotě Nierychlewské-Chejn, Jiřině Hanzlové, Jiřině Chejnové,
          Simoně Petrásové,\\ Petru Mikešovi, Vojtěchu Kupčovi, 
          Renatě Peškové Emilsson, Amiru Mulahumici, \\ Kristýně Antonové, 
          Martině Kašparové, Lucii Brabcové, Samoriele, Adamu Kožouškovi}\\[\baselineskip]
Finanční příspěvek na výtisk exemplářů určených knihovnám\\ v České republice,
Slovenské republice a na Islandu poskytli:\\
  \textbf{Petr Mikeš, Dorota Nierychlewska-Chejn, Jiřina Chejnová, Otta Chejn}
  \ifPDF
  \\[39\baselineskip]
  \else
  \\[36\baselineskip]
  \fi
\textbf{Hvalur.org - 2016}\\
Licence: \textbf{GNU Free Documentation License}\\
Sazba: \textbf{Aleš Chejn} v prostředí  \textbf\LaTeX\\
  Zvláštní poděkování patří \textbf{Paolu al-Imk\=an\=\i\ Brasolin} a uživatelům  \textbf{\TeX\ - \LaTeX\ Stack Exchange}\\[\baselineskip]
\ifPDF
Verze: \textbf{PDF verze} \textit{Islandsko-českého studijního slovníku} vhodná pro \textbf{\ifscreen offline prohlížení a vyhledávání\else vlastní tisk\fi}\\ volně dostupná na \url{www.hvalur.org}\\[\baselineskip]
\else
\textbf\ISBN\\[\baselineskip]
\fi
\ifPDF
\else
Tisk: \textbf{Tiskárna Melmen, Pardubice}\\
\fi
Heslová část vysázena patkovým fontem \TeX\ Gyre Pagella,
hesla bezpatkovým fontem {\phvfamily Helvetica}.
\end{minipage}

% -------------------------------------------------------- Table of contents -

\cleardoublepage 

\pdfbookmark\contentsname{toc}    
\tableofcontents
\thispagestyle{empty}
\cleardoublepage 

\fi

\pagenumbering{arabic}
\setcounter{page}{3}
% ============================================================= INTRODUCTION =

\ifinputintroduction

% ------------------------------------- The study Icelandic-Czech dictionary -

\chapter{O Islandsko-českém studijním slovníku}
Cílem vytvoření\textit{ Islandsko-českého studijního slovníku} je zpřístupnění islandského jazyka všem zájemcům o~tento jazyk. 

\textit{Islandsko-český studijní slovník} je určen především českým a~slovenským studentům islandštiny, začátečníkům i~pokročilým, kteří zde naleznou ucelené fonetické, morfologické a~syntaktické informace, jakož
i~slovní zásobu potřebnou k~aktivnímu i~pasivnímu užití s~rozsahem vhodným ke studiu islandštiny na vysoké škole.
Slovník je rovněž určen všem zájemcům o~islandštinu a~Island, jimž slovník nabízí encyklopedické informace, úvod do fonetiky a~gramatiky a~velké množství fotografií a~ilustrací. 
Třebaže se jedná o~slovník primárně zaměřený na potřeby českých a~slovenských uživatelů, obsahuje také řadu informací, které usnadní práci islandskému uživateli (např. použití více jak 10 000 synonym a~antonym, překlady příkladů ap.).

Slovní zásoba původně vychází z~\textit{Concise Icelandic-English Dictionary} (\cite {ic_en}), značně však byla rozšířena díky přístupu k~webovému projektu \textit{ISLEX} (\cite {int1}). Velké množství hesel bylo
přidáno na podnět uživatelů online verze \textit{Islandsko-českého studijního slovníku}.

Slovník obsahuje 30~575 hesel s~46~725 významy nebo slovními spojeními. Z~celkového počtu 30~575 je 28~097 svébytných hesel obsahujících alespoň jeden význam nebo slovní spojení. 
Zbylá hesla (v~počtu 2~478) jsou varianty, dublety nebo nepravidelné tvary.

Slovník je primárně online elektronický slovník, volně přístupný na internetových stránkách \url{www.hvalur.org}  (\cite {int14}). 
Tištěná verze slovníku je publikována na vlastní náklady v~počtu 30 kusů, přičemž většina exemplářů je určena knihovnám a~institucím v~České republice, Slovenské republice a~na Islandu
a~několik exemplářů je určeno autorům nebo zájemcům z~řad veřejnosti. Dvě PDF verze (vyhledávání\,/\addthin prohlížení offline a~vlastní tisk) jsou dostupné zdarma na internetových stránkách slovníku (viz výše).

\textit{Islandsko-český studijní slovník} je otevřený svobodný projekt, jehož se v~průběhu tvorby zúčastnilo velké množství spoluautorů (viz Historie tvorby IČSS), kterým tímto jménem autorů velmi děkuji.
Děkuji také všem uživatelům online verze slovníku za trpělivost, dobrou vůli a~veškeré připomínky a~návrhy, jmenovitě především Petru Mikešovi, Adamu Kožouškovi a~Samoriele za neutuchající víru v~náš projekt.
Poděkování patří také našim rodinám a~blízkým, díky kterým jsme se mohli soustavně věnovat tvorbě slovníku.
Srdečné poděkování patří mé ženě Dorotce za podporu a~nadšení potřebné k~tvorbě slovníku.

\blspace[5]

{\centering Aleš Chejn, Brwinów, 24. prosince 2015\par}


% ------------ The history of the formation Icelandic-Czech dictionary study -

\chapter{Historie tvorby Islandsko-českého studijního slovníku}
\longchapterskip
Projekt \textit{Islandsko-českého studijního slovníku} (IČSS) vznikl v~roce 2001 na nepřímý podnět doc. PhDr. Heleny Kadečkové, CSc. z~Univerzity Karlovy,
kdy byl vypracován seznam slov. 
Tento seznam obsahoval islandská slova a~jejich překlady, později byly přidány základní morfologické informace o~slovním druhu a~deklinační\,/\addthin konjugační koncovky podle vzoru z~\textit{Concise Icelandic-English Dictionary} (\cite {ic_en}). 
Seznam v~této době obsahoval okolo 1 200 slov. Byl zveřejněn pomocí statických webových stránek a~později byl umístěn na webu Univerzity Karlovy. 

V~roce 2006 byl slovník rozšířen o~nových 1 500 heslových slov a~stávající hesla byla rozšířena o~nové významy, zvláště významy frázových sloves. 
Ke konci roku 2006  byl slovník převeden z programu Excel do databáze MySQL a~byly zahájeny práce na vytvoření aplikace \textit{Dictionary System} (DS) (\cite {int19}), která by podporovala tvorbu IČSS. 
Obsah slovníku byl zveřejněn pod svobodnou licencí na webových stránkách \url{www.hvalur.org} (\cite {int14}), kde se nachází dodnes. 

V~roce 2007 byla vytvořena první verze aplikace DS. V~této době se k~projektu přidala Renata Pešková Emilsson, která se nyní věnuje podpoře výuky češtiny na Islandu. 
Ta slovník obohatila rozsáhlým seznamem slov z~praktického života na Islandu. Dále se slovník rozšiřoval o~heslová slova podle \textit{Concise Icelandic-English Dictionary} (\cite {ic_en}). 
V průběhu následujících třech let bylo doplněno přibližně 15~000 slov. V této době byla hesla rozšířena o~synonyma, příklady, frekvenci, oborovou a~stylovou charakteristiku. K projektu se na Islandu připojila na krátkou dobu Kristýna Antonová, studentka islandštiny na Islandské univerzitě, která 
nás upozornila na nesrovnalosti v~uživatelském rozhraní aplikace DS.

Během léta 2009 byl vytvořen skript na generování skloňování a~časování heslových slov. Výsledky generování byly porovnány s~\textit{Beygingarlýsing íslensks nútímamáls} (\cite {DalvikVM}) a~během následujících let konzultovány
s~Jónem Gíslasonem. Kristín Bjarnadóttir, vedoucí projektu \textit{Beygingarlýsing íslensks nútímamáls}, zodpověděla řadu našich dotazů týkajících se deklinací a~konjugací.

V~roce 2009 začala spolupráce s~\textit{Biolib.cz} (\cite {int5}), jmenovitě s~Ondřejem Zichou. Obdrželi jsme  databázi s~náhledy obrázků, autory a~licencemi a~latinskými názvy botanických a~zoologických druhů.

V~zimním semestru 2009 jsme požádali s~Amirem Mulahumicem, bosenským studentem islandštiny na Islandské univerzitě, Studentský fond Islandské univerzity o~grant na tvorbu zvukové databáze islandských slov. 
Studentský fond návrh přijal. Peníze z~grantu jsme určili pro islandského rodilého mluvčího, Jóna Gíslasona, lektora islandského jazyka na Islandské univerzitě v~Reykjavíku, který nahrál více jak 22~000 islandských slov do elektronické formy. 
S~Dorotou Nierychlewskou-Chejn jsme připravili dva seznamy heslových slov (jeden pro nahrávajícího a~druhý pro zpracování nahrávek). Na podporu nahrávek jsme vytvořili webové stránky nazvané \textit{Databáze s~výslovností islandských slov} (\cite {int13}) a~tyto stránky jsme přeložili do 6~jazyků (Jón Gíslason do islandštiny, Ján Zaťko do slovenštiny, francouzštiny a~angličtiny, Amir Mulamuhic do srbochorvatštiny a~Dorota Nierychlewska-Chejn do polštiny).

V~roce 2010 jsme přidali do databáze pravidla islandské výslovnosti, která jsou uvedena v~knize \textit{Handbók um íslenskan framburð} (\cite {is2}). Z~těchto pravidel jsme automaticky vygenerovali fonetický zápis výslovnosti všech hesel ve slovníku. 
Jón Gíslason ověřil správnost fonetického zápisu u~všech heslových slov.

Na konci roku 2011 jsme spojili webové stránky slovníku se stránkami aplikace a~sjednotili jsme a~vylepšili vzhled webových stránek. Cennými připomínkami se na vytvoření aplikace podílela
RNDr. Ingrid Nagyová, PhD., vysokoškolská pedagožka Ostravské Univerzity. Ján Zaťko přeložil nové webové stránky, Průvodce po slovníku do slovenštiny, francouzštiny a~angličtiny a~celý manuál aplikace do angličtiny. 
Dorota Nierychlewska-Chejn přeložila stránky do polštiny, Jón Gíslason do islandštiny a~Marcos Helena do portugalštiny.

Od prosince 2011 do března 2014 jsme realizovali první čtení, ve kterém jsme doplnili chybějící informace, přidali nová hesla a~opravili pravopisné chyby. 
Obsah slovníku jsme porovnávali především s~\textit{Íslensk orðabók} (\mbox{\cite {is}}) a~\textit{Islandsko-českým slovníkem} (\cite {is7}). S~Vojtěchem Kupčou jsme konzultovali mnohé nesnáze při překladech.

V~letech 2013 až 2015 byl vytvořen morfologický klíč, který rozšiřuje klasický systém zobrazení morfologické informace v~islandských slovnících o~přesné zařazení daného hesla do určité deklinační\,/\addthin konjugační skupiny, jejíž
skloňování\,/\addthin časování je ukázáno na příkladu. V~roce 2014 jsme obdrželi roční grant od Ostravské univerzity v~Ostravě na tvorbu morfologického klíče. S~realizací grantu nám velmi pomohla doc. PhDr. Diana Svobodová, Ph.D. a~Ing. Zuzana Havrlantová.

Od března 2014 do srpna 2015 jsme realizovali druhé finální čtení, ve kterém jsme porovnávali informace s~webovým slovníkem \textit{ISLEX} (\cite {int1}). 
Na požádání jsme obdrželi od Þórdís Úlfarsdóttir (vedoucí projektu \textit{ISLEX}) seznam víceslovných příslovečných frází. Takto poskytnutý seznam nám usnadnil opravy slovníku. 
S~Þórdís Úlfarsdóttir jsme konzultovali různé gramatické jevy a~záležitosti týkající se deklinací\,/\addthin konjugací.
Martina Kašparová, studentka islandštiny na Islandské univerzitě, prošla celý slovník, přičemž se soustředila mimo jiné na stylistické a~pravopisné korektury překladů. 
Lucie Brabcová udělala jazykové a~typografické korektury celého slovníku.

Vojtěch Kupča vydal v~roce 2014 \textit{Islandskou gramatiku} (\cite {is77}) a~souhlasil s~připojením této publikace ke~knižní verzi IČSS.
V~roce 2015 jsme s~přispěním Paola al-Imk\=an\=\i\ Brasolin vysázeli slovník v~typografickém prostředí {\LaTeX}. V~únoru 2016 byl slovník vytištěn v~Tiskárně MELMEN, s.r.o v~Pardubicích a svázán 
v~Knihařství Vojtíšková Hrubá v~Pardubicích. 

\blspace[5]

{\centering Všem spoluautorům děkujeme za spolupráci.\par}


% -------------------------------------------------  Guide to the dictionary -

\chapter{Průvodce po slovníku}
\begin{multicols}{2}
\def\tableA{\noindent\begin{tabularx}{\columnwidth}{>{\bfseries}lll}
  \toprule
  \textit{přípona} & \textit{přípona s~koncovkou} & \textit{příklad} \\
  \midrule
  -ing-   & -ingur   & rugl··ingur \\
          & -ing     & hreyf··ing \\
          & -ingi    & erf··ingi \\
          & -ingar   & lækn··ingar \\
  -ling-  & -lingur  & vett··lingur \\
  -ning-  & -ningur  & flut··ningur \\
          & -ning    & set··ning \\
          & -ningar  & kos··ningar \\
  -að-    & -aður    & mark··aður \\
  -nað-   & -naður   & skil··naður \\
  -ar-    & -ari     & dóm··ari \\
  -skap-  & -skapur  & kveð··skapur \\
  -dóm-   & -dómur   & lær··dómur \\
  -leik-  & -leikur  & sann··leikur \\
          & -leiki   & mögu··leiki \\
  -lát-   & -látur   & laus··látur \\
  -átt-   & -átta    & víð··átta \\
  -sam-   & -samur   & starf··samur \\
  -sem-   & -semi    & hamingju··semi \\
  -nesk-  & -neskur  & tékk··neskur \\
          & -neska   & tékk··neska \\
  -leg-   & -legur   & blíð··legur \\
          & -lega    & blíð··lega \\
  -ung-   & -ungur   & gler··ungur \\
          & -unga    & snopp··unga \\
          & -ungi    & sveit··ungi \\
          & -ung     & nýj··ung \\
  -ótt-   & -óttur   & blett··óttur \\
  -ræn-   & -rænn    & suð··rænn \\
          & -ræna    & nor··ræna \\
  -ern-   & -erni    & fað··erni \\
  -ind-   & -indi    & vís··indi \\
  -and-   & -andi    & gef··andi \\
  -ug-    & -ugur    & göf··ugur \\
          & -ugt     & stöð··ugt \\
  -ust-   & -usta    & þjón··usta \\
  -ald-   & -aldur   & far··aldur \\
  -uð-    & -uður    & söfn··uður \\
  -úð-    & -úð      & dul··úð \\
  -versk- & -verskur & ung··verskur \\
          & -verska  & ung··verska \\
  -ist-   & -isti    & hass··isti \\
  -ism-   & -ismi    & húman··ismi \\
  -ang-   & -angur   & far··angur \\
  \bottomrule
  \end{tabularx}
  \captionof{table}{Seznam přípon označených ve slovníku}
  \label{table1}}

\def\tableB{\noindent\begin{tabularx}{\columnwidth}{>{\bfseries}lX}
  \toprule
  \textbf{symbol} & \textbf{vysvětlivka} \\
  \midrule
  m                   & označuje slovní druh \\
  \textsubscript{1}   & označuje třídu \\
  (-a, -ar)           & označuje deklinační / konjugační koncovky hesla \\
  \textsubscript{8}   & označuje podtřídu\\
  \bottomrule
  \end{tabularx}
  \captionof{table}{Morfologická informace}
  \label{table2}}

\def\tableC{\noindent\begin{tabularx}{\columnwidth}
    {>{\bfseries}l>{\footnotesize\itshape}lllll}
  \toprule
  \multirow{3}{*}{{{\textbf{f{\textsubscript{1}}} \Large{\textbf{7}}}}}
    & nom & kisa & kisan    & kisur & kisurnar \\
    & acc & kisu & kisuna   & kisur & kisurnar \\
    & dat & kisu & kisunni  & kisum & kisunum \\
    & gen & kisu & kisunnar & kisa  & kisanna \\
  \bottomrule
  \end{tabularx}
  \captionof{table}{Skloňování hesla \textit{kisa}}
  \label{table3}}

\def\tableD{\noindent\begin{tabularx}{\columnwidth}{>{\bfseries}lllll}
  \toprule
  \textit{zkratka} & \textit{slovo} & \textit{pád} & \textit{číslo}
    & \specialcell{\textit{předmět}\\\textit{podmět}} \\
  \midrule
  e-að  & eitthvað   & 1. & sg & neživotný \\
  e-ð   & eitthvað   & 4. & sg & neživotný \\
  e-ir  & einhverjir & 1. & pl & životný \\
  e-ja  & einhverja  & 4. & pl & životný \\
  e-jum & einhverjum & 3. & pl & životný \\
  e-m   & einhverjum & 3. & sg & životný \\
  e-n   & einhvern   & 4. & sg & životný \\
  e-r   & einhver    & 1. & sg & životný \\
  e-rra & einhverra  & 2. & pl & životný \\
  e-rs  & einhvers   & 2. & sg & životný \\
  e-s   & einhvers  & 2. & sg & neživotný \\
  e-u   & einhverju  & 3. & sg & neživotný \\
  \bottomrule
  \end{tabularx}
  \captionof{table}{Seznam použitých zkratek neurčitých zájmen}
  \label{table4}}

%#############################################################################

\section{Heslo}

\subsection*{Řazení slov}

Heslová slova jsou zobrazena tučně a~jsou seřazena podle islandské abecedy:

\blspace
{\bfseries\centering
  a~á b c d ð e é f g h i~í j k~l m n\\
  o~ó p q r s~t u~ú v~w x y ý z~þ æ ö\par}
\blspace

K~snadnější orientaci v~heslové části je použit u~každého písmene tzv. palcový index, který se nachází na hřbetě slovníku. Každé písmeno je v~palcovém indexu barevně odlišeno.

\subsection*{Homonyma}

Homonyma (heslová slova stejného reprezentativního tvaru) jsou označena povýšenými arabskými číslicemi uváděnými za heslovým slovem.

\blspace
  \dicEntry \dicTerm{vor\textsuperscript{1}} \dicPos{n} \dicFlx{(-s, -)}
  \dicEntry \dicTerm{vor\textsuperscript{2}} \dicPos{pron} \dicFlx{pers gen pl}
\blspace

Homonyma jsou seřazena podle českého ustáleného řazení slovních druhů. Poté následují předpony a~zkratky.
V~případě dvou podstatných jmen je pořadí mužský rod, ženský rod, střední rod. V~případě dvou podstatných jmen stejného rodu má přednost slovo se slabým skloňováním před slovem se silným skloňováním. V~případě dvou sloves stejného reprezentativního tvaru má přednost sloveso se slabým časováním před slovesem se silným časováním.

\subsection*{Varianty a~dublety}

U~některých heslových slov jsou v~záhlaví hesla uvedeny a) varianty (obměny)

\blspace
  \dicEntry \dicTerm{ærsla··fenginn}, \dicTerm{ærslafullur} \dicPos{adj}
\blspace

nebo b) dublety (dvojtvary).

\blspace
  \dicEntry \dicTerm{{ævin··týri}, \dicTerm{ævintýr}} \dicPos{n}
\blspace

Varianty a~dublety jsou méně frekventované než samotné heslo. Lze je vyhledat ve slovníku jako hesla obsahující fonetické a~morfologické informace s~odkazem na původní heslové slovo.

\blspace
  \dicEntry \dicTerm{ærsla··|fullur} \dicIPA{{a}{i}{\textsubring{r}}{s}{\textsubring{d}}{l}{a}{f}{\textscy}{\textsubring{d}}{l}{\textscy}{\textsubring{r}}} \dicPos{adj} \dicFlx{(comp -fyllri, sup -fyllstur)} \dicLink{ærslafenginn}
  \dicEntry \dicTerm{ævin··týr} \dicIPA{{a}{i}{\textlengthmark}{v}{\textsci}{n}{t\textsuperscript{h}}{i}{\textsubring{r}}} \dicPos{n} \dicFlx{(-s, -)} \dicLink{ævintýri}
\blspace

\subsection*{Vybrané tvary jako hesla}

Ve slovníku jsou uvedeny vybrané tvary heslových slov, u~kterých by mohl mít uživatel potíže s~identifikací. Jedná se o~a) genitiv jednotného čísla ({\textbf{sg}}) a~nominativ množného čísla ({\textbf{pl}}) podstatných jmen, b) ženský a~střední rod přídavných jmen a~2. a~3. stupeň rodu mužského přídavných jmen, c) oznamovací způsob ({\textbf{ind}}) přítomného času ({\textbf{praes}}) a~minulého času ({\textbf{pf}}) 1. osoby ({\textbf{pers}}) sg a~pl a~spojovací způsob ({\textbf{con}}) minulého času 1. osoby sg a~jiné nepravidelné tvary.

\blspace
  \dicEntry \dicTerm{öttum} \dicPos{v} \dicFlx{ind pf pl 1 pers} \dicLink{etja}
\blspace

\subsection*{Dělení heslových slov}

K~dělení heslových slov je použita střední tečka \textbf{'·'}, která odděluje části složeného slova, přípony a~předpony. Platí následující pravidla:

\renewcommand{\labelenumi}{\alph{enumi})}
\begin{enumerate}
\item jediné dělení heslového slova je označeno dvěma středními tečkami \textbf{'··'}. Může se jednat o~dělení částí složeného slova
\blspace
\dicEntry \dicTerm{ráð··hús}
\blspace

nebo o~dělení mezi kořenem a~předponou / příponou.

\blspace
\dicEntry \dicTerm{þjón··usta}
\blspace
\item V~případě, že se v~hesle vyskytuje více dělení, jsou použity dvě střední tečky \textbf{'··'} pro významově poslední dělení a~jedna střední tečka \textbf{'·'} pro zbylé / zbylá dělení.
\blspace
\dicEntry \dicTerm{við·bótar··líf·eyris·sparn·aður}
\blspace
\item V~hesle s~předponou je označena jedna nebo i~více předpon (každá zvlášť).
\blspace
\dicEntry \dicTerm{ó··að·finnan·legur}
\blspace
\item V~hesle s~příponou je označena jedna nebo i~více přípon (každá zvlášť).
\blspace
\dicEntry \dicTerm{létt·úð··ugur}
\blspace
\item V~hesle s~příponou je označena pouze přípona / přípony poslední části složeného slova. Důvodem pro neoznačení přípon jiných částí složených slov je lepší čitelnost hesla. Příklad označení pouze přípony poslední části složeniny:
\blspace
\dicEntry \dicTerm{flutninga··miðl·ari}
\end{enumerate}

Následuje seznam přípon, které jsou ve slovníku označeny:

\tableA

\section{Frekvence slov}

Symbolem \dicsymFrequent\ je označeno 2~825 hesel patřících mezi velmi frekventovaná islandská slova. Tento seznam vychází z~korpusu \textit{Íslenskur Orðasjóður} (\cite{int7}) a~knihy \textit{Íslensk orðtíðnibók} (\cite{is2}).

\blspace
  \dicEntry \dicTerm{koma}\textsuperscript{2} \dicsymFrequent
\blspace

\section{Fonetický zápis}

V~hranatých závorkách v~záhlaví hesla je uveden fonetický zápis výslovnosti hesla v~mezinárodní fonetické abecedě (IPA, angl. International Phonetic Alphabet).

\blspace
  \dicEntry \dicTerm{öldu··rót} \dicIPA{{\oe}{l}{\textsubring{d}}{\textscy}{r}{ou}{\textsubring{d}}}
\blspace

Znak '{\textlengthmark}' označuje dlouhou samohlásku nebo souhlásku.

\blspace
  \dicEntry \dicTerm{ör}\textsuperscript{2} \dicIPA{{\oe}{\textlengthmark}{\textsubring{r}}}
\blspace

\section{Slovní druhy, předpony a~zkratky}

V~záhlaví hesla je uvedena morfologická informace popisující a) slovní druh hesla, b) předponu nebo c) zkratku. Pokud se heslo řadí do více slovních druhů, jsou jednotlivé zkratky pro slovní druhy odděleny lomítkem.

\blspace
  \dicEntry \dicTerm{af} \dicPos{prep / adv}
\blspace

Informace o~slovním druhu, předponě či zkratce je vysázena barevně a~tučně.

\subsection*{Podstatná jména}

Podstatná jména jsou uvedena v~1. pádě (\textbf{nom}) jednotného čísla (\textbf{sg}) nebo v~1. pádě množného čísla (\textbf{pl}), pokud se heslo v~jednotném čísle nevyskytuje. Podstatná jména jsou označena následovně: \textbf{m} (l. masculinum) - podstatné jméno rodu mužského,
\textbf{f} (l. femininum) - podstatné jméno rodu ženského, \textbf{n} (l. neutrum) - podstatné jméno rodu středního, \textbf{subs} (l. substantivum) - podstatné jméno bez jednoznačného mluvnického rodu.

\blspace
  \dicEntry \dicTerm{hestur} \dicPos{m}
  \dicEntry \dicTerm{Hanoí} \dicPos{subs}
\blspace

\subsection*{Přídavná jména}

Přídavná jména jsou uvedena v~1. pádě (\textbf{nom}) rodu mužského (\textbf{m}) jednotného čísla (\textbf{sg}) 1. stupně (\textbf{pos}) (l. positivus) a~jsou označena zkratkou \textbf{adj} (l. adjectivum).

\blspace
  \dicEntry \dicTerm{fallegur} \dicPos{adj}
\blspace

V~případě, že u~přídavného jména dochází v~rodě ženském (\textbf{f}) nebo středním (\textbf{n}) k~přehlásce anebo je tvar druhého stupně (\textbf{comp} (l. comparativus)) či třetího stupně (\textbf{sup} (l. superlativus)) nepravidelný, pak jsou tyto tvary uvedeny v~nominativu jako heslová slova s~odkazem na hlavní heslový tvar

\blspace
  \dicEntry \dicTerm{yngri} \dicPos{adj} \dicFlx{comp m} \dicLink{ungur}
\blspace

nebo jako svébytná hesla s~významy.

\blspace
  \dicEntry \dicTerm{fremri} \dicPos{adj} \dicFlx{comp (sup fremstur)}
\blspace

\subsection*{Zájmena}

Zájmena jsou označena zkratkou \textbf{pron} (l. pronomen) a~jsou tříděna na zájmena ukazovací \textbf{dem} (angl. demonstrative), přivlastňovací \textbf{poss} (angl. possessive), osobní \textbf{pers} (angl. personal), neurčitá \textbf{indef} (angl. indefinite), zvratná \textbf{refl} (angl. reflexive) a~vztažná \textbf{rel} (angl. relative).

\blspace
  \dicEntry \dicTerm{þessi} \dicPos{pron} \dicFlx{dem}
\blspace

U~zájmen přivlastňovacích je v~3. osobě označen rod.

\blspace
  \dicEntry \dicTerm{hans} \dicPos{pron} \dicFlx{poss m}
\blspace

U~zájmen osobních je označen pád a~číslo.

\blspace
  \dicEntry \dicTerm{mig} \dicPos{pron} \dicFlx{pers acc sg}
\blspace

\subsection*{Číslovky}

Číslovky jsou označeny zkratkou \textbf{num} (l. numerale). V~případě vybraných číslovek (1 až 4) jsou ve slovníku uvedeny jako heslová slova také tvary ženského i~středního rodu. V~takovém případě jsou označeny zkratkou \textbf{m} pro mužský rod, zkratkou \textbf{f} pro ženský rod a~zkratkou \textbf{n} pro střední rod. Zkratkou \textbf{ord} (angl. ordinal number) jsou označeny číslovky řadové.

\blspace
  \dicEntry \dicTerm{tut··tugu} \dicPos{num}
  \dicEntry \dicTerm{tveir} \dicPos{num} \dicFlx{pl m}
  \dicEntry \dicTerm{fyrsti} \dicPos{num} \dicFlx{ord}
\blspace

\subsection*{Slovesa}

Slovesa jsou uvedena v~infinitivu a~jsou označena zkratkou \textbf{v} (l. verbum).

\blspace
  \dicEntry \dicTerm{fara} \dicPos{v}
\blspace

Slovesa ve středním rodě jsou označena jako \textbf{refl} (angl. reflexive).

\blspace
  \dicEntry \dicTerm{nálgast} \dicPos{v} \dicFlx{refl}
\blspace

Slovesa, která jsou vždy neosobní, jsou označena \textbf{impers} (angl. impersonal).

\blspace
  \dicEntry \dicTerm{svima} \dicPos{v} \dicFlx{impers}
\blspace

U~sloves je dále uvedeno, s~jakým pádem se pojí. Zkratka \textbf{acc} označuje 4. pád, \textbf{dat} 3. pád, \textbf{gen} 2. pád a~\textbf{nom} 1. pád. V~případě, že se sloveso pojí s~různými pády, je použito lomítko.

\blspace
  \dicEntry \dicTerm{klóra} \dicPos{v} \dicFlx{acc / dat}
\blspace

V~případě, že se sloveso pojí s~více pády, je použito znaménko plus.

\blspace
  \dicEntry \dicTerm{gefa} \dicPos{v} \dicFlx{dat + acc}
\blspace

\subsection*{Příslovce}

Příslovce jsou označena zkratkou \textbf{adv} (l. adverbium). Pokud chybí informace o~stupni, jedná se o~1. stupeň stupňování (\textbf{pos}). 2. stupeň je vždy označen zkratkou \textbf{comp} a~3. stupeň zkratkou \textbf{sup}.

\blspace
  \dicEntry \dicTerm{þar} \dicPos{adv}
  \dicEntry \dicTerm{ofar} \dicPos{adv} \dicFlx{comp (pos uppi, sup efst)}
  \dicEntry \dicTerm{síðast} \dicPos{adv} \dicFlx{sup (comp síðar)}
\blspace

\subsection*{Předložky}

Předložky jsou označeny zkratkou \textbf{prep} (l. praepositio). Pokud se předložka pojí výhradně s~jedním pádem, jsou použity zkratky pro 4. pád (\textbf{acc}), 3. pád (\textbf{dat}), 2. pád (\textbf{gen}).

\blspace
  \dicEntry \dicTerm{af} \dicPos{prep / adv} \dicFlx{dat}
\blspace

V~případě, že se předložka pojí s~různými pády, je použito lomítko.

\blspace
  \dicEntry \dicTerm{fyrir} \dicPos{prep / adv} \dicFlx{acc / dat}
\blspace

\subsection*{Spojky}

Spojky jsou označeny zkratkou \textbf{conj} (l. conjunctio).

\subsection*{Částice}

Částice jsou označeny zkratkou \textbf{part} (l. particula).

\subsection*{Citoslovce}

Citoslovce jsou označeny zkratkou \textbf{inter} (l. interjectio).

\subsection*{Předpony}

Předpony jsou označeny zkratkou \textbf{predp}.

\subsection*{Zkratky}

Zkratky jsou označeny zkratkou \textbf{zkr}.

\section{Skloňování, časování a~stupňování}

Každé heslo obsahuje zápis morfologické informace, která přesně informuje, jak se heslo skloňuje / časuje nebo stupňuje.

V~slovníku jsou použity dva odlišné systémy zápisu morfologické informace, které se vzájemně doplňují:

\begin{enumerate}
\item tradiční systém - do závorky za slovním druhem se zapisují deklinační / konjugační koncovky. Počet koncovek závisí na slovním druhu (viz dále).
\item numerický systém (vychází ze systému Helgi Haraldssona (\cite {is6})) - číslice odkazující na morfologický klíč, kde jsou podrobně vypsány deklinační / konjugační třídy.
Další informace o~numerickém systému lze nalézt na straně \pageref{sec:dec}.
\end{enumerate}

Deklinační / konjugační tvary nebo koncovky podstatných jmen, přídavných jmen, zájmen, číslovek, sloves a~příslovcí jsou uvedeny v~závorkách za slovním druhem.

Hnízdovací šev \textbf{'|'} označuje místo ve slově, kde se deklinační koncovka připojuje ke slovu.

\blspace
  \dicEntry \dicTerm{heils|a{\textsuperscript{1}}} \dicPos{f} \dicFlx{(-u)}
\blspace

Pokud není hnízdovací šev \textbf{'|'} uveden, tak se koncovka připojuje k~samotnému heslu.

\blspace
  \dicEntry \dicTerm{mynd} \dicPos{f} \dicFlx{(-ar, -ir)}
\blspace

Pokud dochází v~ohýbání k~přehlásce nebo je ohýbání nepravidelné, pak jsou uvedeny celé tvary hesla.

\blspace
  \dicEntry \dicTerm{maður} \dicPos{m} \dicFlx{(manns, menn)}
\blspace

V~případě, že heslové slovo má více tvarů, jsou varianty koncovky odděleny lomítkem.

\blspace
  \dicEntry \dicTerm{beð|ur} \dicPos{m} \dicFlx{(-s / -jar, -ir)}
\blspace

\subsection*{Podstatná jména}

U~podstatných jmen je uveden tvar hesla pro 2.~pád jednotného čísla a~1.~pád množného čísla.

\blspace
  \dicEntry \dicTerm{hval|ur} \dicPos{m} \dicFlx{(-s, -ir)}
\blspace

kde '\textit{hvals}' je 2.~pád jednotného čísla a~'\textit{hvalir}' 1.~pád množného čísla.
Pokud je uveden pouze jeden tvar, jedná se o~2.~pád jednotného čísla a~znamená to rovněž, že heslové slovo se neohýbá v~množném čísle.

\blspace
  \dicEntry \dicTerm{heils|a{\textsuperscript{1}}} \dicPos{f} \dicFlx{(-u)}
\blspace

Pokud je použita zkratka pro množné číslo (\textbf{pl}), znamená to, že se heslo ohýbá pouze v~množném čísle.

\blspace
  \dicEntry \dicTerm{afar··kostir} \dicPos{m} \dicFlx{pl}
\blspace

Pokud je použita zkratka \textbf{indecl}, znamená to, že podstatné jméno je nesklonné.

\blspace
  \dicEntry \dicTerm{október} \dicPos{m} \dicFlx{indecl}
\blspace

\subsection*{Přídavná jména}

U~přídavných jmen jsou uvedeny v~závorkách tvary, které poukazují na nepravidelnosti ve skloňování.

Pokud přídavné jméno v~ženském rodě neztrácí koncové '\textit{-r}', je uveden tvar hesla v~ženském rodě.

\blspace
  \dicEntry \dicTerm{hýr} \dicPos{adj} \dicFlx{(f -)}
\blspace

Dále jsou uvedeny tvary ženského rodu, ve~kterých dochází k~přehlásce:

\blspace
  \dicEntry \dicTerm{reglu··|samur} \dicPos{adj} \dicFlx{(f -söm)}
\blspace

Pokud je přídavné jméno nesklonné, je uvedena zkratka \textbf{indecl}.

\blspace
  \dicEntry \dicTerm{tví··mála} \dicPos{adj} \dicFlx{indecl}
\blspace

Pokud má přídavné jméno nepravidelné tvary 2.~a~3.~stupně stupňování, jsou tyto tvary uvedeny v~závorce spolu se zkratkou pro 2.~stupeň (\textbf{comp}) a~3.~stupeň (\textbf{sup}).

\blspace
  \dicEntry \dicTerm{þunnur} \dicPos{adj} \dicFlx{(comp þynnri, sup þynnstur)}
\blspace

V~případě, že je uveden pouze jeden ze dvou zbývajících stupňů, znamená to, že heslo se stupňuje pouze v~uvedených stupních.

\blspace
  \dicEntry \dicTerm{syðri} \dicPos{adj} \dicFlx{comp (sup syðstur)}
\blspace

\subsection*{Slovesa}

Islandská slovesa se dělí na slabá, silná a~nepravidelná. Vzhledem k~tomuto dělení se liší počet slovních tvarů uvedených v~závorce za zkratkou \textbf{v}. Princip uvádění slovních tvarů spočívá v~tom, že uživatel
je schopen (více méně) z~těchto tvarů odvodit časování slovesa. V~tomto slovníku jsme
záměrně umístili některá slabá slovesa mezi silná (místo dvou slovesných tvarů je uvedeno pět tvarů) kvůli tomu, aby byl uživatel přesněji informován o~časování slovesa.

Slabá slovesa, která tvoří 1.~osobu jednotného čísla minulého času oznamovacího způsobu koncovkou '\textit{-aði}', mají v~závorce uvedenu jen jednu koncovku, jmenovitě '\textit{(-aði)}'.

\blspace
  \dicEntry \dicTerm{ætl|a} \dicPos{v} \dicFlx{(-aði)}
\blspace

U~zbývajících skupin slabých sloves se v~závorce nachází koncovka 1.~osoby jednotného čísla minulého času oznamovacího způsobu a~supinum.

\blspace
  \dicEntry \dicTerm{kenn|a} \dicPos{v} \dicFlx{(-di, -t)}
\blspace

Silná nebo nepravidelná slovesa mají v~závorce vždy pět tvarů, jmenovitě 1.~osobu jednotného čísla přítomného času oznamovacího způsobu, 1.~osobu jednotného čísla minulého času oznamovacího způsobu, 1.~osobu množného čísla minulého času oznamovacího způsobu, 1.~osobu jednotného čísla minulého času spojovacího způsobu a~supinum.

\blspace
  \dicEntry \dicTerm{grípa} \dicPos{v} \dicFlx{(gríp, greip, gripum, gripi, gripið)}
\blspace

Silná nebo nepravidelná slovesa vyskytující se pouze v~neosobním užití mají v~závorce 4 tvary, jmenovitě 3.~osobu jednotného čísla přítomného času oznamovacího způsobu, 3.~osobu jednotného čísla minulého času oznamovacího způsobu, 3.~osobu jednotného čísla minulého času spojovacího způsobu a~supinum.

\blspace
  \dicEntry \dicTerm{kross··|bregða} \dicPos{v} \dicFlx{(-bregður, -brá, -brygði, -brugðið)}
\blspace

U~některých významů v~heslovém slově se vyskytují gramatické informace, které popisují chování daného významu.

Zkratka \textbf{impers} označuje, že se daný význam slovesa vyskytuje v~neosobním užití. Pomocí slovního spojení je vyjádřeno užití slovesa.

\blspace
  \dicEntry \dicTerm{batn|a} \dicPos{v} \dicFlx{(-aði)} \dicTerm{e-m batnar} \dicFlx{impers} (kdo) se uzdravuje, (kdo) se zotavuje, (komu) je lépe
\blspace

Zkratka \textbf{refl} označuje mediopasivní tvary slovesa.

\blspace
  \dicEntry \dicTerm{að··lag|a} \dicPos{v} \dicFlx{(-aði)} \dicIdiom{aðlagast} {\dicPhraseIS{aðlagast e-u}} \dicFlx{refl} {přizpůsobit se (čemu)}
\blspace

\subsection*{Příslovce}

U~příslovcí se v~záhlaví nachází tvary stupňování. Tyto tvary jsou uvedeny v~závorce spolu se zkratkou pro 2.~stupeň (\textbf{comp}) a~3.~stupeň (\textbf{sup}).

\blspace
  \dicEntry \dicTerm{vel\textsuperscript{2}} \dicPos{adv} \dicFlx{(comp betur, sup best)}
\blspace

Nepravidelné tvary 2.~a~3.~stupně jsou uvedeny buď jako hesla odkazující na základní tvar v~1.~stupni (\textbf{pos}), anebo jako svébytná hesla. V~obou případech jsou uvedeny v~závorce zbývající stupně.

\blspace
  \dicEntry \dicTerm{síst} \dicPos{adv} \dicFlx{sup (pos varla, comp síður)}
\blspace

Pokud některý ze stupňů uveden není, znamená to, že pro dané slovo neexistuje.

\blspace
  \dicEntry \dicTerm{síðar} \dicPos{adv} \dicFlx{comp (sup síðast)}
\blspace

\section{Morfologická informace} \label{sec:dec}

V~záhlaví hesla se kromě deklinačních / konjugačních koncovek nachází také morfologické numerické informace, které odkazují na morfologický klíč (viz strana \pageref{sec:morpho}), kde jsou podrobně vypsány kompletní deklinační / konjugační tabulky.

\blspace
  \dicEntry \dicTerm{ap|i} \dicPos{m}[1] \dicFlx{(-a, -ar)}[8]
\blspace

Morfologická informace je zapsána podle následujícího schématu: \textbf{m{\textsubscript{1}} (-a, -ar){\textsubscript{8}}}

\blspace
\tableB
\blspace

Deklinační / konjugační třídy a~jejich počet nejsou v~islandštině ustálené a~slouží pouze pro rozlišení odlišných skloňování / časování.

Názvy podtříd slouží pro rozlišení odlišných skloňování / časování.

\subsection*{Práce s~morfologickým klíčem}

Pro vyhledání kompletních deklinačních / konjugačních tabulek slouží následující postup, který demonstrujeme na ukázce:

\blspace
  \dicEntry \dicTerm{kis|a} \dicPos{f}[1] \dicFlx{(-u, -ur)}[7]
\blspace

Nejdříve určíme slovní druh (v~případě podstatných jmen také rod), v~tomto případě podstatné jméno rodu ženského, a~navštívíme Morfologický klíč (str. \pageref{sec:morpho}) pro podstatná jména rodu ženského (str. \pageref{sec:morpho_f}).
Následně dohledáme třídu (v~tomto případě číslo {\textsubscript{\textbf{1}}}) a~následně dohledáme dodatečné numerické informace za závorkou (v~tomto případě číslo {\textsubscript{\textbf{7}}}).

\tableC

\section{Definice}

\subsection*{Základní tvar}

Základním tvarem definice je, že jednomu islandskému slovu odpovídá jeden nebo více českých významů.

\blspace
  \dicEntry \dicTerm{landa··fræð|i} \dicPos{f} \dicFlx{(-i)} \textCS{zeměpis, geografie}
\blspace

V~některých případech je v~české části použita závorka. V~prvním případě je závorka užívaná ke~zkrácení zápisu dvou slov ('\textit{pomoc}', '\textit{výpomoc}').

\blspace
  \dicEntry \dicTerm{að··stoð} \dicPos{f} \dicFlx{(-ar)} \textCS{(vý)pomoc, asistence}
\blspace

V~druhém případě se v~závorce nachází upřesňující informace, která zabraňuje dvojznačnosti.

\blspace
  \dicEntry \dicTerm{fíkj|a} \dicPos{f} \dicFlx{(-u, -ur)} \textCS{fík} {\footnotesize {\textCS{(plod)}}}
\blspace

V~mnoha případech je závorka v~české části použita pro příklad užití českého slova a~tím vymezení jeho významu.

\blspace
  \dicEntry \dicTerm{að·gengi··legur} \dicPos{adj} {přístupný, dostupný (informace ap.)}
\blspace

Lomítko \textbf{'/'} slouží v~české části významu k~oddělení více významů pojících se s~jedním nebo více slovy a~je použito v~české části následujícím způsobem:

\blspace
  \dicEntry \dicTerm{aðal··|gata} \dicPos{f} \dicFlx{(-götu, -götur)} \textCS{hlavní třída / ulice}
\blspace

Význam přečteme jako '\textit{hlavní třída, hlavní ulice}'.

\subsection*{Opis heslového slova}

V~případech, kdy islandské slovo nemá ekvivalent v~českém jazyce, jsou významy islandských slov opsány delším opisem a~zobrazeny menším písmem. Příkladem jsou slova z~lidových pověstí, kulinářské speciality nebo botanické pojmy.

\blspace
  \dicEntry \dicTerm{til··ber|i} \dicPos{m} \dicFlx{(-a, -ar)} {\footnotesize{pov.}} {\footnotesize{stvoření, které dojí krávy a~ovce hospodářů}}
\blspace

\subsection*{Slovní spojení}

V~definici hesla jsou kromě samotných významů uváděna také slovní spojení, která jsou sázena tučně. Po slovním spojení následuje překlad slovního spojení.

\blspace
  \dicEntry \dicTerm{lær|a} \dicPos{v} \dicFlx{(-ði, -t)} \textbf{læra e-ð utanbókar} {učit se (co) zpaměti, memorovat (co)}
\blspace

Závorky jsou použity na tu část slovního spojení, kterou je možné vypustit.

\blspace
  \dicEntry \dicTerm{ár\textsuperscript{1}} \dicPos{f} \dicFlx{(-ar, -ar)} \textbf{taka (of) djúpt í árinni} {\footnotesize{přen.}} {přehnat (co v~tvrzení)}
\blspace

Příklad můžeme přečíst jako '\textit{taka djúpt í árinni}' nebo '\textit{taka of djúpt í árinni}'.

Mezi slovní spojení v~tomto slovníku řadíme a) tvary heslového slova, b) příkladová spojení, c) ustálená slovní spojení a~d) rčení nebo přísloví.

\begin{enumerate}
\item Tvary heslového slova\\
Do této skupiny řadíme tvar množného čísla hesla (zkratka \textbf{pl}).
\blspace
\dicEntry \dicTerm{bót} \dicPos{f} \dicFlx{(-ar, bætur)} \textbf{bætur} \dicFlx{pl} {dávky, přídavky, příspěvky}\\

Dále do této skupiny řadíme tvar slovesa v~středním slovesném rodě (zkratka \textbf{refl}).\\

\dicEntry \dicTerm{finna} \dicPos{v} \dicFlx{(finn, fann, fundum, fyndi, fundið)} \dicIdiom{finnast} {\dicPhraseIS{finnast}} \dicFlx{refl} {být nalezen, najít se, nalézt se, nacházet se}
\blspace
\item Příkladová spojení\\
Příkladová spojení ilustrují komunikační chování heslových slov v~příslušných významech a~v~různých syntagmatech (např. sloveso s~určitou předložkou ap.). Jsou označena tučně modrou barvou a~odsazená pro snadnou orientaci jako svébytné heslo.
\blspace
\dicEntry \dicTerm{koma\textsuperscript{2}} \dicPos{v} \dicFlx{(kem, kom, komum, kæmi, komið)} \dicIdiom{koma}[að] {\dicPhraseIS{koma að e-m}} {přiblížit se ke (komu), přijít ke (komu)}
\blspace
\item Ustálená slovní spojení\\
Tento druh slovních spojení je charakterizován většinou přeneseným významem. Tato spojení jsou uváděna za posledním významem nebo příkladovým spojením a~jsou označená tučně modrou barvou a~odsazená pro snadnou orientaci jako svébytné heslo.
\blspace
\dicEntry \dicTerm{kast} \dicPos{n} \dicFlx{(-s, köst)} \dicIdiom{kast} {\dicPhraseIS{þegar til kastanna kemur}} {\footnotesize{přen.}} {když jde do tuhého}
\blspace
\item Rčení a~přísloví\\
Rčení a~přísloví se nacházejí vždy na konci definice a~jsou odlišená symbolem \dicsymProverb\ a~odsazená pro snadnou orientaci.
\blspace
\dicEntry \dicTerm{æfing} \dicPos{f} \dicFlx{(-ar, -ar)} \dicProverb{} \dicPhraseIS{Æfing skapar meistara.} {\footnotesize{přís.}} \textCS{Cvičení dělá mistra.}
\end{enumerate}

\section{Vícevýznamovost}

\subsection*{Vícevýznamová hesla}

Značné množství hesel má více než jeden význam. Významy jsou odlišeny arabskými číslicemi a~v~některých případech i~islandskými synonymy.

\blspace
  \dicEntry \dicTerm{land} \dicPos{n} \dicFlx{(-s, lönd)} {\textbf{1.}} \dicSynonym{þurrlendi} {souš, pevnina, země} {\textbf{2.}} \dicSynonym{árbakki} {břeh} {\textbf{3.}} \dicSynonym{ríki} {země, stát, vlast} {\textbf{4.}} \dicSynonym{landareign} {země, půda, pozemek}
\blspace

\subsection*{Vícevýznamová slovní spojení}

V~případě, že slovní spojení uvnitř definice má více významů, jsou jednotlivé významy odlišeny písmeny {\textbf{a., b., c.}} atd.

\blspace
  \dicEntry \dicTerm{drag} {\dicPos{n} \dicFlx{(-s, drög)} {\textbf{1.}} {mokřina, podmoklý terén;} {\textbf{2.}} \textbf{drög} {\footnotesize{pl}} {\textbf{a.}} \dicSynonym{uppsprettur} {prameny (řeky ap.)} {\textbf{b.}} \dicSynonym{undirbúningur} {náčrt, návrh.}}
\blspace

\section{Řazení významů a~slovních spojení v~definici}
\longsectionskip

V~definici jsou nejdříve seřazeny významy slova a~odlišeny arabskými číslicemi.

\subsection*{V rámci významu}

Slovní spojení jsou uváděna v~rámci jednotlivého významu, pokud se s~ním významově vážou.

\blspace
  \dicEntry \dicTerm{tím|i} \dicPos{m} \dicFlx{(-a, -ar)} {\textbf{1.}} \dicSynonym{tíð} {čas, doba}; \textbf{í þann tíma} \dicFlx{adv} \dicSynonym{þá\textsuperscript{2}} {pak} {\textbf{2.}} \dicSynonym{rétt stund} {(správný) čas / moment}
\blspace

\subsection*{Za posledním očíslovaným významem}

Pokud se slovní spojení neváže k~žádnému významu, je řazeno za posledním očíslovaným významem. Ustálená slovní spojení jsou uvedena symbolem \dicsymIdiom.

\blspace
  \dicEntry \dicTerm{finna} \dicPos{v} \dicFlx{(finn, fann, fundum, fyndi, fundið)} {\textbf{3.}} \dicSynonym{skynja} {cítit, vnímat} \dicIdiom{finna}[að] {\dicPhraseIS{finna að e-u}} {kritizovat (co), nacházet chyby na (čem)}
\blspace

\subsection*{Způsob řazení slovních spojení}

Slovní spojení, která se nachází za posledním očíslovaným významem, jsou řazena podle předložky / příslovce / zájmena, se kterou se pojí. Pro snadnou orientaci v~rozsáhlých definicích je uvedeno heslové slovo nebo tvar heslového slova plus předložka / příslovce / zájmeno a~barevně odlišeno modrou barvou.

\blspace
  \dicEntry \dicTerm{koma\textsuperscript{2}} \dicPos{v} \dicFlx{(kem, kom, komum, kæmi, komið)} \dicIdiom{koma}[út] {\dicPhraseIS{koma út}} {vyjít (kniha ap.), objevit se (na trhu ap.)}
\blspace

Poté následují slovní spojení v~slovesném středním rodě řazená opět podle předložky / příslovce / zájmena, se kterou se pojí:

\blspace
  \dicEntry \dicTerm{koma\textsuperscript{2}} \dicPos{v} \dicFlx{(kem, kom, komum, kæmi, komið)} \dicIdiom{komast} {\dicPhraseIS{komast}} \dicFlx{refl} {dostat se, dostávat se}
\blspace

Poté následují ustálená slovní spojení, která nebylo možné zařadit do předchozích položek:

\blspace
  \dicEntry \dicTerm{koma\textsuperscript{2}} \dicPos{v} \dicFlx{(kem, kom, komum, kæmi, komið)} \dicIdiom{koma} {\dicPhraseIS{(komdu) sæll (og blessaður)}} {ahoj, nazdar, dobrý den}
\blspace

Nakonec jsou uvedena přísloví a~rčení označená symbolem \dicsymProverb\:

\blspace
  \dicEntry \dicTerm{eik} \dicPos{n} \dicFlx{(-ar / -ur, -ur)} \dicProverb{} \dicPhraseIS{Ekki / eigi fellur eik við fyrsta högg.} {\footnotesize{přís.}} {Jednou ranou dub nespadne.}
\blspace

\section{Synonyma a~antonyma / hesla ke~srovnání}
\longsectionskip

\subsection*{Synonyma}

Synonyma nebo synonymní slovní spojení jsou uvedeny před opisem významu hesla. Jsou umístěna v~závorkách. Islandská synonyma slouží a) islandskému uživateli k~rozlišení českých významů b) českému uživateli jako dodatečná informace o~významu slova.

\blspace
  \dicEntry \dicTerm{staða} \dicPos{f} \dicFlx{(stöðu, stöður)} {\textbf{1.}} \dicSynonym{ástand} {situace, stav, poměry} {\textbf{2.}} \dicSynonym{embætti} {post, (pracovní) místo / pozice (ve firmě ap.)}
\blspace

V~případě, že synonymum lze nalézt jako heslo ve slovníku, je použita šipka '\dicsymSee\ '.

\blspace
  \dicEntry \dicTerm{bygging} {\dicPos{f} \dicFlx{(-ar, -ar)} {\textbf{1.}} \dicSynonym{hús} {budova, stavba, stavení} {\textbf{2.}} \textit{(það að byggja)} {(vý)stavba} }
\blspace

\subsection*{Antonyma / Hesla ke~srovnání}

Antonyma či hesla ke~srovnání jsou uvedena za opisem významu heslového slova a~jsou označena šipkou '$\shortuparrow$'.

\blspace
  \dicEntry \dicTerm{kaldur} \dicPos{adj} \dicFlx{(f köld)} {\textbf{1.}} {studený, chladný} \textit{($\shortuparrow$ heitur)}
\blspace

\section{Oborové a~stylové charakteristiky}

K~upřesnění významu heslového slova slouží oborové a~stylové charakteristiky.

\subsection*{Oborové charakteristiky}

Oborové charakteristiky zařazují islandské slovo do určitého oboru a~naznačují použití daného slova v~jazyce

\blspace
  \dicEntry \dicTerm{berg} \dicPos{n} \dicFlx{(-s, -)} {\textbf{1.}} {\footnotesize{geol.}} {hornina, kámen} {\textbf{2.}} \dicSynonym{klöpp\textsuperscript{1}} {skalní stěna, skalisko}
\blspace

a~vymezují význam českého slova.

\blspace
  \dicEntry \dicTerm{mús} \dicPos{f} \dicFlx{(-ar, mýs)} {\textbf{1.}} {\footnotesize{zool.}} {myš} (\textit{l. Mus}) {\textbf{2.}} {\footnotesize{poč.}} {myš}
\blspace

\subsection*{Stylové charakteristiky}

Stylové charakteristiky jsou uváděny v~případě, kdy význam heslového slova je možné použít pouze při určité příležitosti nebo je význam jazykově zabarven.

\blspace
  \dicEntry \dicTerm{gaur} \dicPos{m} \dicFlx{(-s, -ar)} {\textbf{1.}} {\footnotesize{hovor.}} {chlápek, maník, člověk}
\blspace

\section{Syntax}

Syntaktické informace popisují:

\begin{enumerate}
\item zdali je sloveso tranzitivní nebo intranzitivní a~s~jakým(i) pádem (pády) se pojí, pokud je tranzitivní,
\item zdali je podmět / předmět životný či neživotný,
\item zdali je sloveso osobní či neosobní a~v~jakém pádu je podmět, pokud je sloveso neosobní,
\item zdali je podmět gramatický,
\item zdali je sloveso použito ve středním rodě,
\item nejčastější příklady, které se pojí se slovesem,
\item se kterými příslovci nebo předložkami se sloveso pojí.
\end{enumerate}

Syntaktické informace se ve slovníku vyskytují ve třech formách: a) informace uvedené zkratkou, b) teoretické příklady a~c) praktické příklady.

\subsection*{Syntaktické informace uvedené zkratkou}

Zde máme na mysli informace o~pádě (\textbf{nom, acc, dat, gen}), slovesu v~středním slovesném rodě (\textbf{refl}), neosobním slovesu (\textbf{impers}), čísle (\textbf{sg, pl}), rozkazovacím způsobu (\textbf{imper}) a~spojovacím způsobu (\textbf{con}).

\subsection*{Teoretické příklady}

Teoretickým příkladem nazýváme příklad, kdy podmět nebo předmět je zastoupen zájmenem neurčitým v~pádě, se~kterým se slovo pojí a~sloveso je zpravidla v~infinitivu. V~případě, že podmět není v~1.~pádě (např. u~sloves neosobních), je sloveso uvedeno v~3.~osobě jednotného čísla přítomného času.
Ve slovních spojeních používáme zkratky pro tvary neurčitého zájmena \textbf{einhver} (někdo) a~\textbf{eitthvað} (něco). Tyto zkratky označují a) pád, b) číslo a~c) životnost nebo neživotnost podmětu / předmětu.\\
Následuje seznam zkratek neurčitých zájmen použitých ve slovníku: \\

\tableD
\blspace
Příklad použití zkratek s~neurčitými zájmeny:

\blspace
  \dicEntry \dicTerm{leit|a} \dicPos{v} \dicFlx{(-aði)} \textbf{leita að e-u / e-m} {hledat (co / koho)} \dicExampleIS{leita að lyklunum} \dicExampleCS{hledat klíče}
\blspace

Příklad tedy přečteme jako '\textit{leita að einhverju}' (\textit{hledat co}) a~'\textit{leita að einhverjum}' (\textit{hledat koho}). V~tomto případě se sloveso pojí s~3.~pádem a~předmět může být životný i~neživotný.

\subsection*{Praktické příklady}

Ve slovníku se vyskytuje velké množství praktických příkladů a~jejich překladů. Praktický příklad se liší od příkladu teoretického tím, že místo zájmen neurčitých jsou uvedena frekventovaná slova z~běžného projevu. Praktické příklady slouží k~ilustraci použití
heslového slova.\\
Praktickým příkladem může být a) (složené) slovo, b) slovní spojení nebo c) věta.

\begin{enumerate}
\item Slovo jako praktický příklad:
\blspace
\dicEntry \dicTerm{hett|a} \dicPos{f} \dicFlx{(-u, -ur)} \textbf{1.} {kapuce, kápě} \dicExampleIS{hettuúlpa} \dicExampleCS{bunda s~kapucí}
\blspace
\item Slovní spojení jako praktický příklad:
\blspace
\dicEntry \dicTerm{tím|i} \dicPos{m} \dicFlx{(-a, -ar)} \textbf{3.} \dicSynonym{klukkustund} {hodina} {\footnotesize{(šedesát minut)}} \dicExampleIS{tveggja tíma gangur} \dicExampleCS{dvouhodinový pochod}
\blspace
\item Věta jako praktický příklad:
\blspace
\dicEntry \dicTerm{mæt|a} \dicPos{v} \dicFlx{(-ti, -t)} \dicIdiom{mætast} \dicPhraseIS{mætast} \dicSynonym*{ná saman} \dicFlx{refl} {setkat se, potkat se} \dicExampleIS{Þau mættust á gangstéttinni.} \dicExampleCS{Potkali se na chodníku.}
\end{enumerate}

\section{Fotografie a~ilustrace}

Ve slovníku je publikováno 630 fotografií rostlin, živočichů a~fotografií, která dokumentují místopisně heslová slova (např. kulinářské speciality).
Fotografie jsou publikovány pod veřejnými licencemi. Jména autorů a~názvy licencí lze nalézt v~Seznamu autorů fotografií a~ilustrací na straně \pageref{sec:photo}. U~hesla a~konkrétně u~daného významu, ke kterému se fotografie vztahuje, je použit symbol fotoaparátu \dicsymPhoto. Fotografie je umístěna buď pod heslem, anebo v~horní části stránky, kde se heslo nachází, anebo na stránce (stránkách) následující(ch).

%\par\begin{center}\setlength\fboxsep{0pt}\setlength\fboxrule{0.5pt}
%\fbox{\includegraphics[width=6cm]{ds_image_posthus_0_2.jpg}}\end{center}
%\par\begin{center}\footnotesize {Pósthús}\end{center}

\end{multicols}

% ---------------------------------------- List of abbreviations and symbols -

\chapter{Seznam použitých zkratek a symbolů}
\begin{multicols}{2}
\begin{description}[leftmargin=!,labelwidth=1.5cm]
%[\breaklabel\setleftmargin{40pt}\setlabelstyle{\bfseries}
%\addtolength{\itemsep}{-0.5\baselineskip}]
\item[{abb}] {zkrácený}
\item[{acc}] {4. pád, akuzativ}
\item[{adj}] {přídavné jméno}
\item[{adv}] {příslovce}
\item[{akt}] {rod činný, aktivum}
\item[{anat.}] {anatomie}
\item[{angl.}] {anglicky}
\item[{astro.}] {astronomie, výzkum vesmíru}
\item[{ap.}] {a podobně}
\item[{básn.}] {básnický, poetický výraz}
\item[{biol.}] {biologie}
\item[{bot.}] {botanika}

\item[{comp}] {2. stupeň stupňování, komparativ}
\item[{con}] {spojovací způsob, konjunktiv}
\item[{conj}] {spojka}
\item[{dat}] {3. pád, dativ}
\item[{def}] {určitý tvar}
\item[{dem}] {zájmeno ukazovací}
\item[{dět.}] {dětsky}

\item[{e-að}] {co (sg nom neživotný) (\textit{eitthvað})}
\item[{e-ð}] {co (sg acc neživotný) (\textit{eitthvað})}
\item[{e-ir}] {kdo (pl nom životný) (\textit{einhverjir})}
\item[{e-ja}] {koho (pl acc životný) (\textit{einhverja})}
\item[{e-jum}] {komu (pl dat životný) (\textit{einhverjum})}
\item[{e-m}] {komu (sg dat životný) (\textit{einhverjum})}
\item[{e-n}] {koho (sg acc životný) (\textit{einhvern})}
\item[{e-r}] {kdo (sg nom životný) (\textit{einhver})}
\item[{e-rra}] {koho (pl gen životný) (\textit{einhverra})}
\item[{e-rs}] {koho (sg gen životný) (\textit{einhvers})}
\item[{e-s}] {čeho (sg gen neživotný) (\textit{einhvers})}
\item[{e-u}] {čemu (sg dat neživotný) (\textit{einhverju})}

\item[{ekon.}] {ekonomika, obchod}
\item[{elek.}] {elektřina}
\item[{f}] {rod ženský}
\item[{filos.}] {filosofie, logika}
\item[{form.}] {formálně}
\item[{fyz.}] {fyzika}
\item[{gen}] {2. pád, genitiv}
\item[{geog.}] {zeměpis, geografie}
\item[{geol.}] {geologie}
\item[{han.}] {hanlivě, pejorativně}
\item[{hist.}] {historie}
\item[{hovor.}] {hovorově}
\item[{hrub.}] {hrubě}
\item[{hud.}] {hudba}
\item[{imper}] {rozkazovací způsob, imperativ}
\item[{impers}] {sloveso neosobní}
\item[{ind}] {oznamovací způsob, indikativ}
\item[{indecl}] {nesklonný}
\item[{indef}] {neurčitý tvar, (pron +) zájmeno neurčité}
\item[{inf}] {infinitiv}
\item[{int}] {tázací zájmeno}
\item[{inter}] {citoslovce}
\item[{chem.}] {chemie}
\item[{jaz.}] {jazykověda}
\item[{kulin.}] {kulinářství, vaření}
\item[{l.}] {latinsky}
\item[{let.}] {letectví}
\item[{lit.}] {literatura, vydavatelství}
\item[{m}] {rod mužský}
\item[{mat.}] {matematika}
\item[{med. }] {lékařství}
\item[{meteo.}] {meteorologie}
\item[{myt.}] {mytologie}
\item[{n}] {rod střední}
\item[{náb.}] {náboženství}
\item[{nám.}] {námořnictví, rybolov}
\item[{nom}] {1. pád, nominativ}
\item[{num}] {číslovka}
\item[{ord}] {řadová číslovka}
\item[{p}] {osoba}
\item[{part}] {částice}
\item[{pers}] {(v +) osoba, (pron +) zájmeno osobní}
\item[{pf}] {minulý čas}
\item[{pl}] {množné číslo, plurál}
\item[{poč.}] {informatika}
\item[{pol.}] {politika, politologie}
\item[{poss}] {zájmeno přivlastňovací}
\item[{pos}] {1. stupeň stupňování, pozitiv}
\item[{pov.}] {lidové pověsti, folkloristika}
\item[{pp}] {příčestí minulé}
\item[{praes}] {přítomný čas}
\item[{práv.}] {právnictví, soudnictví}
\item[{predp}] {předpona}
\item[{prep}] {předložka}
\item[{presp}] {příčestí přítomné}
\item[{pron}] {zájmeno}
\item[{prop}] {vlastní jméno, proprium}
\item[{přen.}] {přeneseně}
\item[{přís.}] {přísloví}
\item[{psych.}] {psychologie}
\item[{refl}] {(v +) médium, (pron +) zvratné zájmeno}
\item[{rel}] {vztažné zájmeno}
\item[{sg}] {jednotné číslo, singulár}
\item[{slang.}] {slang}
\item[{sport.}] {sport}
\item[{stav.}] {stavebnictví, architektura}
\item[{subs}] {podstatné jméno (bez rodu)}
\item[{sup}] {3. stupeň stupňování, superlativ}
\item[{supin}] {supinum}
\item[{škol.}] {školství}
\item[{techn.}] {technika, mechanika}
\item[{v}] {sloveso}
\item[{voj.}] {vojenství}
\item[{zast.}] {zastarale}
\item[{zkr}] {zkratka}
\item[{zool.}] {zoologie}

\item[{;}] {středník uprostřed definice odděluje věty; středník na konci definice odděluje významy}

\item[{(...)}] {v kulatých závorkách jsou uvedeny doplňkové informace}
\item[{[...]}] {v hranatých závorkách je uveden fonetický zápis výslovnosti}\footnote{Vysvětlení fonetických symbolů je uvedeno v kapitole Fonetika \ref{sec:phon_phonems} na straně \pageref{sec:phon_phonems}}
\item[{|}] {naznačuje místo, kde se deklinační koncovka připojuje ke slovnímu základu}
\item[{·}] {vedlejší dělení složeného slova}
\item[{··}] {hlavní dělení složeného slova}
\item[{+}] {plus}
\item[{/}] {lomítko označuje více možností}
\item[{\dicsymSee}] {šipka odkazuje na heslo ve slovníku}
\item[{\dicsymCompare}] {šipka odkazuje na heslo ve slovníku ke srovnání}
\item[{\dicsymExampleIS}] {uvozuje islandský příklad}
\item[{\dicsymExampleCS}] {uvozuje český překlad příkladu}
\item[{\dicsymIdiom}] {uvozuje slovní spojení}
\item[{\dicsymFrequent}] {heslo patří mezi frekventované výrazy}

\item[{\dicsymPhoto}] {fotografie, ilustrace}
\item[{\dicsymProverb}] {přísloví, rčení}

\end{description}
\end{multicols}

\chapter{Skrá yfir skammstafanir og tákn}
\begin{multicols}{2}
\begin{description}[leftmargin=!,labelwidth=1.5cm]
%[\breaklabel\setleftmargin{40pt}\setlabelstyle{\bfseries}
%\addtolength{\itemsep}{-0.5\baselineskip}]
\item[{abb}] {stýfður}
\item[{acc}] {þolfall}
\item[{adj}] {lýsingarorð}
\item[{adv}] {atviksorð}
\item[{akt}] {germynd}
\item[{anat.}] {líffærafræði}
\item[{angl.}] {enska}
\item[{astro.}] {geimvísindi, stjörnufræði}
\item[{ap.}] {e.þ.h.}
\item[{básn.}] {ljóðrænn}
\item[{biol.}] {líffræði}
\item[{bot.}] {grasafræði}

\item[{comp}] {miðstig}
\item[{con}] {viðtengingarháttur}
\item[{conj}] {samtenging}
\item[{dat}] {þágufall}
\item[{def}] {með greini}
\item[{dem}] {ábendingarfornafn}
\item[{dět.}] {barnamál}

\item[{e-að}] {(\textit{eitthvað})}
\item[{e-ð}] {(\textit{eitthvað})}
\item[{e-ir}] {(\textit{einhverjir})}
\item[{e-ja}] {(\textit{einhverja})}
\item[{e-jum}] {(\textit{einhverjum})}
\item[{e-m}] {(\textit{einhverjum})}
\item[{e-n}] {(\textit{einhvern})}
\item[{e-r}] { (\textit{einhver})}
\item[{e-rra}] {(\textit{einhverra})}
\item[{e-rs}] {(\textit{einhvers})}
\item[{e-s}] {(\textit{einhvers})}
\item[{e-u}] { (\textit{einhverju})}

\item[{ekon.}] {hagfræði, viðskipti}
\item[{elek.}] {rafmagn}
\item[{f}] {kvenkyn}
\item[{filos.}] {heimspeki, rökfræði}
\item[{form.}] {formlegt mál}
\item[{fyz.}] {eðlisfræði}
\item[{gen}] {eignarfall}
\item[{geog.}] {landafræði}
\item[{geol.}] {jarðfræði}
\item[{han.}] {niðrandi}
\item[{hist.}] {sagnfræði}
\item[{hovor.}] {talmál}
\item[{hrub.}] {gróft}
\item[{hud.}] {tónlist}
\item[{imper}] {boðháttur}
\item[{impers}] {ópersónuleg sögn}
\item[{ind}] {framsöguháttur}
\item[{indecl}] {óbeygjanlegur}
\item[{indef}] {óákveðinn, (pron +) óákveðið fornafn}
\item[{inf}] {nafnháttur}
\item[{int}] {spurnigarfornafn}
\item[{inter}] {upphrópun}
\item[{chem.}] {efnafræði}
\item[{jaz.}] {málfræði, málvísindi}
\item[{kulin.}] {eldhús}
\item[{l.}] {latína}
\item[{let.}] {flug, flugmál}
\item[{lit.}] {skáldskapur}
\item[{m}] {nafnorð, karlkyn}
\item[{mat.}] {stærðfræði}
\item[{med. }] {læknisfræði}
\item[{meteo.}] {veðurfræði}
\item[{myt.}] {goðafræði}
\item[{n}] {nafnorð, hvorugkyn}
\item[{náb.}] {trúfræði}
\item[{nám.}] {sjómennska}
\item[{nom}] {nefnifall}
\item[{num}] {töluorð}
\item[{ord}] {raðtala}
\item[{p}] {persóna}
\item[{part}] {ögn}
\item[{pers}] {(v +) persóna, (pron +) persónufornafn}
\item[{pf}] {þátíð}
\item[{pl}] {fleirtala}
\item[{poč.}] {tölvufræði, tölvur}
\item[{pol.}] {stjórnmál, stjórnmálafræði}
\item[{poss}] {eignarfornafn}
\item[{pos}] {frumstig}
\item[{pov.}] {þjóðtrú}
\item[{pp}] {lýsingarháttur þátíðar}
\item[{praes}] {nútíð}
\item[{práv.}] {lögfræði, dómsmál}
\item[{predp}] {forskeyti}
\item[{prep}] {forsetning}
\item[{presp}] {lýsingarháttur nútíðar}
\item[{pron}] {fornafn}
\item[{prop}] {sérnafn}
\item[{přen.}] {afleidd merking}
\item[{přís.}] {málsháttur}
\item[{psych.}] {sálfræði}
\item[{refl}] {(v +) miðmynd, (pron +) afturbeygt fornafn}
\item[{rel}] {tilvísunarfornafn}
\item[{sg}] {eintala}
\item[{slang.}] {slangur}
\item[{sport.}] {íþróttir}
\item[{stav.}] {húsasmíði, byggingarlist}
\item[{subs}] {nafnorð (án kyns)}
\item[{sup}] {efstastig}
\item[{supin}] {sagnbót}
\item[{škol.}] {skólamál}
\item[{techn.}] {tækni, verkfræði}
\item[{v}] {sagnorð}
\item[{voj.}] {her}
\item[{zast.}] {úrelt mál}
\item[{zkr}] {skammstöfun}
\item[{zool.}] {dýrafræði}

\item[{;}] {semíkomma í miðri skýringu aðgreinir setningu; semíkomma í enda skýringar aðgreinir mismunandi merkingarsvið}

\item[{(...)}] {innan sviga eru viðbótarupplýsingar}
\item[{[...]}] {innan hornklofa er hljóðritaður framburður}
\item[{|}] {merkir stað, þar sem beygingarending bætist við stofn}
\item[{·}] {aukasamskeyti samsetts orðs}
\item[{··}] {aðalsamskeyti samsetts orðs}
\item[{+}] {plús}
\item[{/}] {skástrik táknar fleiri möguleika}
\item[{\dicsymSee}] {ör vísar til uppflettiorðs í orðabókinni}
\item[{\dicsymCompare}] {ör vísar til uppflettiorðs í bókinni til samanburðar}
\item[{\dicsymExampleIS}] {kemur á undan dæmi á íslensku}
\item[{\dicsymExampleCS}] {kemur á undan tékkneskri þýðingu á dæminu}
\item[{\dicsymIdiom}] {orð með háa tíðni}

\item[{\dicsymPhoto}] {ljósmynd}
\item[{\dicsymProverb}] {málsháttur}

\end{description}
\end{multicols}

% ------------------------------------------------- Headword part and alphabet

\chapter{Heslová část}
{\renewcommand{\arraystretch}{2}%
%\topskip0pt
\vstretch{1}
{\setlength{\tabcolsep}{1.5em}
\begin{table}[htpb]
\centering

\begin{tabular}{ccccccc} \hline
A a &  & G g &  & O o &  & V v \\
Á á  &  & H h &  & Ó ó &  & W w \\
B b &  & I i &  & P p &  & X x \\
C c &  & Í í &  & Q q &  & Y y \\
D d &  & J j &  & R r &  & Ý ý \\
Ð ð &  & K k &  & S s &  & Z z \\
E e &  & L l &  & T t &  & Þ þ \\
É é  &  & M m &  & U u &  & Æ æ \\
F f &  & N n &  & Ú ú &  & Ö ö \\
\hline
\end{tabular}

%\caption{\textbf{Islandská abeceda}}
\label{my-islandska-abeceda}
\end{table}
}
\vstretch{1}
{\renewcommand{\arraystretch}{1}%

\clearpage

\fi

% =========================================================== THE DICTIONARY =

\ifinputletters

\cleardoublepage

\dictionarygeometry
\pagestyle{myheadings}

\newcommand*{\alphabetF}{%
  a,aa,b,c,d,e,ee,f,g,h,i,ii,j,k,l,m,n,o,oo,p,r,s1,s2,t,u,uu,v,w,x,y,yy,th,ae,oe}
\ifPDFtest
\dicLetter{t}{letter23}
\dicEntry[t.a.m.] \dicTerm{t.a.m.} \dicPos{zkr} \dicPhraseIS{til að mynda} \dicFlx{adv} \dicDirectTranslationCS{pro ilustraci\,/\addthin názornost}
\dicEntry[t.d.] \dicTerm{t.d.} \dicPos{zkr} \dicPhraseIS{til dæmis} \dicFlx{adv} \dicDirectTranslationCS{například}
\dicEntry[t.h.] \dicTerm{t.h.} \dicPos{zkr} \dicPhraseIS{til hægri} \dicFlx{adv} \dicDirectTranslationCS{vpravo, doprava} \dicAntonym{t.v.}
\dicEntry[t.v.] \dicTerm{t.v.} \dicPos{zkr} \dicPhraseIS{til vinstri} \dicFlx{adv} \dicDirectTranslationCS{vlevo, doleva} \dicAntonym{t.h.}
\dicEntry[tabú] \dicTerm{tabú} \dicIPA{{t\smash{\textsuperscript{h}}}{a}{\textlengthmark}{\textsubring{b}}{u}} \dicPos{n}[2] \dicFlx{(‑s)}[2] \dicSynonym{bannhelgi} \dicDirectTranslationCS{tabu}
\dicEntry[Tadsjiki] \dicTerm{Tadsjik|i} \dicIPA{{t\smash{\textsuperscript{h}}}{a}{\textsubring{d}}{s}{j}{\textsci}{\r{\textObardotlessj}}{\textsci}} \dicPos{m}[1] \dicFlx{(‑a, ‑ar)}[1] \dicDirectTranslationCS{Tádžik, Tádžička}
\dicEntry[Tadsjikistan] \dicTerm{Tadsjikistan} \dicIPA{{t\smash{\textsuperscript{h}}}{a}{\textsubring{d}}{s}{j}{\textsci}{\r{\textObardotlessj}}{\textsci}{s}{\textsubring{d}}{a}{\textsubring{n}}} \dicPos{n}[2] \dicFlx{(‑\,/\addthin ‑s)}[35] \dicFieldCat{geog.} \dicDirectTranslationCS{Tádžikistán}
\dicEntry[tadsjikska] \dicTerm{tadsjiksk|a} \dicIPA{{t\smash{\textsuperscript{h}}}{a}{\textsubring{d}}{s}{j}{\textsci}{\r{g}}{s}{\r{g}}{a}} \dicPos{f}[1] \dicFlx{(‑u)}[5] \dicDirectTranslationCS{tádžičtina}
\dicEntry[tadsjikskur] \dicTerm{tadsjikskur} \dicIPA{{t\smash{\textsuperscript{h}}}{a}{\textsubring{d}}{s}{j}{\textsci}{\r{g}}{s}{\r{g}}{\textscy}{\textsubring{r}}} \dicPos{adj}[1]\dicFlx{}[-6] \dicDirectTranslationCS{tádžický}
\dicEntry[tað] \dicTerm{tað} \dicIPA{{t\smash{\textsuperscript{h}}}{a}{\textlengthmark}{\texttheta}} \dicPos{n}[2] \dicFlx{(‑s, töð)}[8] \dicSynonym{skán} \dicDirectTranslationCS{vysušený trus} \dicIndirectTranslationCS{(používaný jako palivo)}
\dicEntry[taða] \dicTerm{taða} \dicIPA{{t\smash{\textsuperscript{h}}}{a}{\textlengthmark}{ð}{a}} \dicPos{f}[1] \dicFlx{(töðu, töður)}[8] \dicSynonym*{hey af túni} \dicDirectTranslationCS{seno (z~louky)}
\dicEntry[tafar] \dicTerm{tafar} \dicIPA{{t\smash{\textsuperscript{h}}}{a}{\textlengthmark}{v}{a}{\textsubring{r}}} \dicPos{f} \dicFlx{sg gen} \dicLink{töf}
\dicEntry[tafarlaus] \dicTerm{tafar··laus} \dicIPA{{t\smash{\textsuperscript{h}}}{a}{\textlengthmark}{v}{a}{r}{l}{\oe i}{s}} \dicPos{adj}[5]\dicFlx{}[-1] \dicDirectTranslationCS{okamžitý, bezodkladný, neprodlený}
\dicEntry[tafarlaust] \dicTerm{tafar··laust} \dicIPA{{t\smash{\textsuperscript{h}}}{a}{\textlengthmark}{v}{a}{r}{l}{\oe i}{s}{\textsubring{d}}} \dicPos{adv} \dicDirectTranslationCS{okamžitě, bezodkladně, neprodleně}
\dicEntry[tafði] \dicTerm{tafði} \dicIPA{{t\smash{\textsuperscript{h}}}{a}{v}{ð}{\textsci}} \dicPos{v} \dicFlx{ind pf sg 1 pers} \dicLink{tefja}
\dicEntry[tafið] \dicTerm{tafið} \dicIPA{{t\smash{\textsuperscript{h}}}{a}{\textlengthmark}{v}{\textsci}{\texttheta}} \dicPos{v} \dicFlx{supin} \dicLink{tefja}
\dicEntry[tafir] \dicTerm{tafir} \dicIPA{{t\smash{\textsuperscript{h}}}{a}{\textlengthmark}{v}{\textsci}{\textsubring{r}}} \dicPos{f} \dicFlx{pl nom} \dicLink{töf}
\dicEntry[tafl] \dicTerm{tafl} \dicIPA{{t\smash{\textsuperscript{h}}}{a}{\textsubring{b}}{\textsubring{l}}} \dicPos{n}[2] \dicFlx{(‑s, töfl)}[8] \dicSynonym{skák} \dicDirectTranslationCS{šachy};  \dicPhraseIS{hér eru brögð í tafli} \dicLangCat{přen.} \dicDirectTranslationCS{je v~tom nějaký háček}
\dicEntry[tafla] \dicTerm{tafla} \dicsymFrequent\  \dicIPA{{t\smash{\textsuperscript{h}}}{a}{\textsubring{b}}{l}{a}} \dicPos{f}[1] \dicFlx{(töflu, töflur)}[20] \textbf{1.} \dicDirectTranslationCS{tabule (školní ap.)} \dicExampleIS{skrifa orðið upp á töflu} \dicExampleCS{napsat slovo na tabuli}  \textbf{2.} \dicSynonym{listi} \dicDirectTranslationCS{tabulka} \dicIndirectTranslationCS{(přehledný graficky členěný seznam)} \dicExampleIS{töflur og línurit} \dicExampleCS{tabulky a~grafy}  \textbf{3.} \dicSynonym{pilla\smash{\textsuperscript{1}}} \dicDirectTranslationCS{tabletka, pilulka, prášek} \dicExampleIS{taka töflu við höfuðverknum} \dicExampleCS{vzít si tabletku při bolesti hlavy}  \textbf{4.} \dicFieldCat{poč.} \dicDirectTranslationCS{tabulka}
\dicEntry[taflborð] \dicTerm{tafl··borð} \dicIPA{{t\smash{\textsuperscript{h}}}{a}{\textsubring{b}}{\textsubring{l}}{\textsubring{b}}{\textopeno}{r}{\texttheta}} \dicPos{n}[2] \dicFlx{(‑s, ‑)}[5] \dicSynonym{skákborð} \dicDirectTranslationCS{šachovnice}
\dicEntry[taflfélag] \dicTerm{tafl··fé·|lag} \dicIPA{{t\smash{\textsuperscript{h}}}{a}{\textsubring{b}}{\textsubring{l}}{f}{j}{\textepsilon}{l}{a}{x}} \dicPos{n}[2] \dicFlx{(‑lags, ‑lög)}[8] \dicDirectTranslationCS{šachový klub}
\dicEntry[taflmaður] \dicTerm{tafl··|maður} \dicIPA{{t\smash{\textsuperscript{h}}}{a}{\textsubring{b}}{l}{m}{a}{ð}{\textscy}{\textsubring{r}}} \dicPos{m}[13] \dicFlx{(‑manns, ‑menn)}[2] \textbf{1.} \dicDirectTranslationCS{šachista, šachistka}  \textbf{2.} \dicDirectTranslationCS{šachová figurka}
\dicEntry[tafsamur] \dicTerm{taf··|samur} \dicIPA{{t\smash{\textsuperscript{h}}}{a}{f}{s}{a}{m}{\textscy}{\textsubring{r}}} \dicPos{adj}[1] \dicFlx{(f ‑söm)}[2] \dicSynonym{tímafrekur} \dicDirectTranslationCS{časově náročný, zdlouhavý} \dicExampleIS{tafsamt verk} \dicExampleCS{časově náročná práce}
\dicEntry[tag] \dicTerm{tag} \dicsymFrequent\  \dicIPA{{t\smash{\textsuperscript{h}}}{a}{\textlengthmark}{x}} \dicPos{n}[2] \dicFlx{(‑s)}[2] \dicSynonym{tegund} \dicDirectTranslationCS{druh} \dicExampleIS{af ýmsu tagi} \dicExampleCS{různého druhu}
\dicEntry[tagl] \dicTerm{tagl} \dicIPA{{t\smash{\textsuperscript{h}}}{a}{\r{g}}{\textsubring{l}}} \dicPos{n}[2] \dicFlx{(‑s, tögl)}[8] \textbf{1.} \dicDirectTranslationCS{ohon, oháňka, ocas} \dicIndirectTranslationCS{(koně)}  \textbf{2.} \dicSynonym*{hár í teygju} \dicDirectTranslationCS{koňský ohon} \dicIndirectTranslationCS{(vlasy sepnuté na temeni hlavy)};  \dicPhraseIS{hafa tögl(in) og hagldir(nar)} \dicLangCat{přen.} \dicDirectTranslationCS{držet otěže pevně ve svých rukou}
\dicEntry[taglhnýtingur] \dicTerm{tagl··hnýt·ing|ur} \dicIPA{{t\smash{\textsuperscript{h}}}{a}{\r{g}}{\textsubring{l}}{\textsubring{n}}{i}{\textsubring{d}}{i}{\ng}{\r{g}}{\textscy}{\textsubring{r}}} \dicPos{m}[6] \dicFlx{(‑s, ‑ar)}[8] \textbf{1.} \dicDirectTranslationCS{kůň přivázaný za ocas jiného koně}  \textbf{2.} \dicSynonym*{ósjálfstæður fylgismaður} \dicDirectTranslationCS{loutka, ocásek} \dicIndirectTranslationCS{(člověk ovládaný cizí vůlí)}
\dicEntry[Taíland] \dicTerm{Taí··land}\dicTerm{, Tæland} \dicIPA{{t\smash{\textsuperscript{h}}}\-{a}\-{i}\-{\textlengthmark}\-{l}\-{a}\-{n}\-{\textsubring{d}}\-} \dicPos{n}[2] \dicFlx{(‑s)}[4] \dicFieldCat{geog.} \dicDirectTranslationCS{Thajsko}
\dicEntry[Taílendingur] \dicTerm{Taí·lend··ing|ur}\dicTerm{, Tælendingur} \dicIPA{{t\smash{\textsuperscript{h}}}\-{a}\-{i}\-{\textlengthmark}\-{l}\-{\textepsilon}\-{n}\-{\textsubring{d}}\-{i}\-{\ng}\-{\r{g}}\-{\textscy}\-{\textsubring{r}}\-} \dicPos{m}[6] \dicFlx{(‑s, ‑ar)}[8] \dicDirectTranslationCS{Thajec, Thajka}
\dicEntry[taílenska] \dicTerm{taí··lensk|a}\dicTerm{, tælenska} \dicIPA{{t\smash{\textsuperscript{h}}}\-{a}\-{i}\-{\textlengthmark}\-{l}\-{\textepsilon}\-{n}\-{s}\-{\r{g}}\-{a}\-} \dicPos{f}[1] \dicFlx{(‑u)}[5] \dicDirectTranslationCS{thajština}
\dicEntry[taílenskur] \dicTerm{taí··lenskur}\dicTerm{, tælenskur} \dicIPA{{t\smash{\textsuperscript{h}}}\-{a}\-{i}\-{\textlengthmark}\-{l}\-{\textepsilon}\-{n}\-{s}\-{\r{g}}\-{\textscy}\-{\textsubring{r}}\-} \dicPos{adj}[1]\dicFlx{}[-6] \dicDirectTranslationCS{thajský}
\dicEntry[Taívan] \dicTerm{Taívan} \dicIPA{{t\smash{\textsuperscript{h}}}{a}{i}{\textlengthmark}{v}{a}{\textsubring{n}}} \dicPos{n}[2] \dicFlx{(‑\,/\addthin ‑s)}[35] \dicFieldCat{geog.} \dicDirectTranslationCS{Tchaj-wan}
\dicEntry[Taívani] \dicTerm{Taívan|i} \dicIPA{{t\smash{\textsuperscript{h}}}{a}{i}{\textlengthmark}{v}{a}{n}{\textsci}} \dicPos{m}[1] \dicFlx{(‑a, ‑ar)}[8] \dicDirectTranslationCS{obyvatel(ka) Tchaj-wanu}
\dicEntry[taívanskur] \dicTerm{taívanskur} \dicIPA{{t\smash{\textsuperscript{h}}}{a}{i}{\textlengthmark}{v}{a}{n}{s}{\r{g}}{\textscy}{\textsubring{r}}} \dicPos{adj}[1] \dicFlx{(f taívönsk)}[3] \dicDirectTranslationCS{tchajwanský}
\dicEntry[tak] \dicTerm{tak} \dicsymFrequent\  \dicIPA{{t\smash{\textsuperscript{h}}}{a}{\textlengthmark}{\r{g}}} \dicPos{n}[2] \dicFlx{(‑s, tök)}[8] \textbf{1.} \dicDirectTranslationCS{sevření, stisk, úchop};  \dicPhraseIS{hafa tak á e‑u} \dicDirectTranslationCS{svírat (co), (pevně) držet (co)};  \dicPhraseIS{losa um takið} \dicDirectTranslationCS{povolit sevření};  \dicPhraseIS{missa takið á e‑u} \dicDirectTranslationCS{pustit (co)} \dicIndirectTranslationCS{(neúmyslně)};  \dicPhraseIS{ná taki á e‑u} \dicDirectTranslationCS{sevřít (co), uchopit (co)};  \dicPhraseIS{sleppa takinu á e‑u} \dicDirectTranslationCS{pustit (co)} \dicIndirectTranslationCS{(úmyslně)}  \textbf{2.} \dicPhraseIS{tök} \dicFlx{pl} \dicDirectTranslationCS{osvojení si, ovládnutí};  \dicPhraseIS{kunna tökin á e‑u} \dicDirectTranslationCS{zvládat (co), mít (co) osvojené};  \dicPhraseIS{missa tökin á e‑u} \dicDirectTranslationCS{ztratit kontakt s~(čím) (projektem ap.)};  \dicPhraseIS{ná tökum á e‑u} \dicDirectTranslationCS{dostat se do (čeho), osvojit si (co)}  \textbf{3.} \dicSynonym{verkur} \dicDirectTranslationCS{bolest} \dicExampleIS{fá tak í bakið} \dicExampleCS{dostat bolest zad};  \dicIdiom{tak}{ \dicPhraseIS{e‑að er til taks}} \dicDirectTranslationCS{(co) je k~dispozici} \dicExampleIS{Verkfærin eru til taks.} \dicExampleCS{Nářadí je k~dispozici.}; { \dicPhraseIS{hafa\,/\addthin vera með (margt) í takinu}} \dicLangCat{přen.} \dicDirectTranslationCS{dělat (více věcí) najednou}; { \dicPhraseIS{hafa engin tök á að (gera e‑ð)}} \dicDirectTranslationCS{nemít šanci (udělat (co))};  \dicPhraseIS{taka e‑ð föstum tökum} \dicLangCat{přen.} \dicDirectTranslationCS{zaměřit se intenzivně\,/\addthin zcela na (co) (problém ap.)};  \dicPhraseIS{taka sér tak} \dicLangCat{přen.} \dicDirectTranslationCS{sebrat se, dát se dohromady};  \dicPhraseIS{það eru engin tök á að (gera e‑ð)} \dicDirectTranslationCS{není možné (udělat (co)), nedá se (udělat (co))}
\dicEntry[taka] \dicTerm{taka\smash{\textsuperscript{1}}} \dicIPA{{t\smash{\textsuperscript{h}}}{a}{\textlengthmark}{\r{g}}{a}} \dicPos{f}[1] \dicFlx{(töku, tökur)}[20] \textbf{1.} \dicSynonym{valdataka} \dicDirectTranslationCS{(pře)vzetí, obsazení, zabrání}  \textbf{2.} \dicSynonym*{kvikmyndataka} \dicDirectTranslationCS{filmování, natáčení}  \textbf{3.} \dicFieldCat{práv.} \dicDirectTranslationCS{okupace} \dicIndirectTranslationCS{(nabytí vlastnického práva k~věci bez vlastníka)}
\dicEntry[taka] \dicTerm{taka\smash{\textsuperscript{2}}} \dicsymFrequent\  \dicIPA{{t\smash{\textsuperscript{h}}}{a}{\textlengthmark}{\r{g}}{a}} \dicPos{v}[6] \dicFlx{(tek, tók, tókum, tæki, tekið)}[125] \dicFlx{acc\,/\addthin dat} \textbf{1.} \dicFlx{acc} \dicDirectTranslationCS{vzít, sebrat, brát} \dicExampleIS{taka bókina af honum} \dicExampleCS{vzít mu knížku}  \textbf{2.} \dicFlx{acc} \dicDirectTranslationCS{zmocnit se, získat, ovládnout, dobýt} \dicExampleIS{taka borg} \dicExampleCS{zmocnit se města}  \textbf{3.} \dicFlx{acc} \dicDirectTranslationCS{sebrat, zadržet, zatknout};  \dicPhraseIS{taka e‑n fastan} \dicDirectTranslationCS{zadržet (koho), zatknout (koho)}  \textbf{4.} \dicFlx{acc} \dicDirectTranslationCS{udělat, zhotovit (snímek ap.)};  \dicPhraseIS{taka mynd} \dicDirectTranslationCS{udělat fotku\,/\addthin snímek, (vy)fotit};  \dicPhraseIS{taka dæmi} \dicDirectTranslationCS{uvést\,/\addthin dát příklad};  \dicPhraseIS{taka afrit} \dicDirectTranslationCS{udělat kopii, okopírovat}  \textbf{5.} \dicFlx{acc} \dicDirectTranslationCS{pojmout, obsáhnout} \dicExampleIS{Ílátið tekur tvo lítra.} \dicExampleCS{Nádoba pojme dva litry.}  \textbf{6.} \dicFlx{dat} \dicDirectTranslationCS{přijmout, přijímat, vzít, brát} \dicExampleIS{taka e‑m opnum örmum} \dicExampleCS{přijmout (koho) s~otevřenou náručí}  \textbf{7.} \dicFlx{dat} \dicDirectTranslationCS{přijmout, přijímat, (za)reagovat} \dicExampleIS{taka e‑u vel} \dicExampleCS{přijmout (co) dobře}  \textbf{8.} \dicFlx{acc} \dicDirectTranslationCS{brát, chápat, pojímat} \dicExampleIS{taka orð hans alvarlega} \dicExampleCS{brát jeho slova vážně}  \textbf{9.} \dicDirectTranslationCS{schylovat se, začínat (pršet ap.)} \dicExampleIS{Það tekur að rigna.} \dicExampleCS{Schyluje se k~dešti.}  \textbf{10.} \dicPhraseIS{taka sér e‑ð} \dicDirectTranslationCS{vzít si (co)} \dicExampleIS{taka sér frí} \dicExampleCS{vzít si volno}  \textbf{11.} \dicPhraseIS{það tekur því ekki} \dicFlx{impers} \dicDirectTranslationCS{nevyplatí se, nemá cenu};  \dicIdiom{taka}[að]{ \dicPhraseIS{taka e‑ð að sér}} \dicDirectTranslationCS{vzít na sebe (co), ujmout se (čeho) (úkolu ap.)}; { \dicPhraseIS{taka e‑n að sér}} \dicDirectTranslationCS{ujmout se (koho), vzít (koho) do opatrování (dítě ap.)} \dicExampleIS{taka barn að sér} \dicExampleCS{ujmout se dítěte};  \dicIdiom{taka}[af]{ \dicPhraseIS{e‑ð tekur fljótt af}} \dicFlx{impers} \dicDirectTranslationCS{(co) brzy skončí\,/\addthin pomine (bolest ap.)}; { \dicPhraseIS{taka af sér e‑ð}} \dicDirectTranslationCS{odložit (co), sundat si (co) (klobouk ap.)} \dicExampleIS{taka af sér húfuna} \dicExampleCS{sundat si čepici}; { \dicPhraseIS{taka e‑n af lífi}} \dicDirectTranslationCS{popravit (koho)};  \dicIdiom{taka}[aftur]{ \dicPhraseIS{taka e‑ð aftur}} \dicDirectTranslationCS{vzít (co) zpět} \dicExampleIS{taka loforðið aftur} \dicExampleCS{vzít slib zpět};  \dicIdiom{taka}[á]{ \dicPhraseIS{taka á e‑u}} \dicDirectTranslationCS{chopit se (čeho), zaměřit se na (co)}; { \dicPhraseIS{taka á rás}} \dicDirectTranslationCS{dát se do běhu};  \dicPhraseIS{taka á öllu sínu} \dicDirectTranslationCS{udělat vše, co je v~lidských silách}; { \dicPhraseIS{taka e‑n á orðinu}} \dicLangCat{přen.} \dicDirectTranslationCS{vzít (koho) za slovo}; { \dicPhraseIS{taka sig á}} \dicDirectTranslationCS{sebrat se, polepšit se};  \dicIdiom{taka}[á móti]{ \dicPhraseIS{taka á móti e‑u}} \dicDirectTranslationCS{přijmout (co), přijímat (co)} \dicExampleIS{taka á móti gjöf} \dicExampleCS{přijmout dárek}; { \dicPhraseIS{taka á móti e‑m}} \dicDirectTranslationCS{přijmout (koho), přijímat (koho)};  \dicIdiom{taka}[eftir]{ \dicPhraseIS{taka eftir e‑u\,/\addthin e‑m}} \dicDirectTranslationCS{všimnout si (čeho\,/\addthin koho), zaznamenat (co\,/\addthin koho)} \dicExampleIS{Hún tók ekki eftir mér þótt ég kallaði til hennar.} \dicExampleCS{Nevšimla si mě, i~když jsem na ni volal.};  \dicIdiom{taka}[fram]{ \dicPhraseIS{taka fram e‑ð}} \dicDirectTranslationCS{vzít (co), sebrat (co), posbírat (co)} \dicExampleIS{taka fram dót} \dicExampleCS{posbírat věci}; { \dicPhraseIS{taka e‑u\,/\addthin e‑m fram}} \dicDirectTranslationCS{převýšit (co\,/\addthin koho), převyšovat (co\,/\addthin koho)} \dicExampleIS{taka honum fram} \dicExampleCS{převyšovat ho}; { \dicPhraseIS{taka e‑ð fram}} \dicDirectTranslationCS{podtrhnout (co), zdůraznit (co)};  \dicIdiom{taka}[fram yfir]{ \dicPhraseIS{taka e‑ð fram yfir e‑ð}} \dicSynonym*{kjósa heldur} \dicDirectTranslationCS{preferovat (co) před (čím), upřednostňovat (co) před (čím)} \dicExampleIS{taka fiskmeti fram yfir kjöt} \dicExampleCS{preferovat ryby před masem};  \dicIdiom{taka}[frá]{ \dicPhraseIS{taka frá e‑ð}} \dicDirectTranslationCS{rezervovat (co), zamluvit (co), odložit (co) stranou} \dicExampleIS{taka frá vöru fyrir hana} \dicExampleCS{zamluvit pro ni zboží}; { \dicPhraseIS{taka e‑ð frá e‑m}} \dicDirectTranslationCS{vzít (co komu), brát (co komu)} \dicExampleIS{taka bókina frá honum} \dicExampleCS{vzít mu knížku};  \dicIdiom{taka}[fyrir]{ \dicPhraseIS{taka (borgun) fyrir e‑ð}} \dicDirectTranslationCS{přijmout (platbu) za (co)}; { \dicPhraseIS{taka fyrir e‑ð}} {\textbf{a.}} \dicDirectTranslationCS{udělat (čemu) přítrž, přerušit (co)};  {\textbf{b.}} \dicDirectTranslationCS{projednat (co), prohovořit (co)};  \dicIdiom{taka}[inn]{ \dicPhraseIS{taka e‑n inn í e‑ð}} \dicDirectTranslationCS{přijmout (koho) do (čeho) (žáka do školy ap.)}; { \dicPhraseIS{taka inn e‑ð}} \dicDirectTranslationCS{přijmout (co), vzít (co) (lék ap.)} \dicExampleIS{taka inn lyf} \dicExampleCS{přijmout lék};  \dicIdiom{taka}[inn á]{ \dicPhraseIS{taka e‑ð inn á sig}} \dicLangCat{přen.} \dicDirectTranslationCS{brát si (co) k~srdci};  \dicIdiom{taka}[í]{ \dicPhraseIS{taka í e‑n, taka í lurginn á e‑m}} \dicDirectTranslationCS{potrestat (koho) tělesně}; { \dicPhraseIS{taka í nefið}} \dicDirectTranslationCS{šňupnout si, šňupat (tabák ap.)}; { \dicPhraseIS{taka vel\,/\addthin illa í e‑ð}} \dicDirectTranslationCS{stavět se pozitivně\,/\addthin negativně k~(čemu), přijmout\,/\addthin brát (co) dobře\,/\addthin špatně};  \dicIdiom{taka}[í sundur]{ \dicPhraseIS{taka e‑ð í sundur}} \dicDirectTranslationCS{rozebrat (co), demontovat (co)};  \dicIdiom{taka}[með]{ \dicPhraseIS{taka e‑n með sér}} \dicDirectTranslationCS{vzít (koho) s~sebou} \dicExampleIS{taka hana með sér í vinnuna} \dicExampleCS{vzít ji s~sebou do práce};  \dicPhraseIS{taka með (sér) e‑ð} \dicDirectTranslationCS{vzít si s~sebou (co)} \dicExampleIS{taka með nesti í vinnuna} \dicExampleCS{vzít si s~sebou do práce svačinu};  \dicIdiom{taka}[nærri]{ \dicPhraseIS{taka e‑ð nærri sér}} \dicLangCat{přen.} \dicDirectTranslationCS{brát si (co) (k~srdci ap.)};  \dicIdiom{taka}[ofan]{ \dicPhraseIS{taka ofan fyrir e‑m}} {\textbf{a.}} \dicDirectTranslationCS{smeknout\,/\addthin smekat před (kým)} \dicIndirectTranslationCS{(sejmout pokrývku hlavy na pozdrav)} \dicExampleIS{taka ofan hattinn fyrir henni} \dicExampleCS{smeknout před ní klobouk};  {\textbf{b.}} \dicLangCat{přen.} \dicDirectTranslationCS{smeknout\,/\addthin smekat před (kým)} \dicIndirectTranslationCS{(uznat něčí kvality)};  \dicIdiom{taka}[saman]{ \dicPhraseIS{taka saman e‑ð}} {\textbf{a.}} \dicDirectTranslationCS{(po)sbírat (co)} \dicExampleIS{taka saman fötin} \dicExampleCS{posbírat šaty};  {\textbf{b.}} \dicDirectTranslationCS{složit, sepsat (článek ap.)} \dicExampleIS{taka saman grein} \dicExampleCS{sepsat článek}; { \dicPhraseIS{taka sig saman um e‑ð}} \dicDirectTranslationCS{dát se dohromady v~(čem), spolčit se kvůli (čemu)} \dicExampleIS{Þau tóku sig saman um að kaupa gjöf handa honum.} \dicExampleCS{Dali se dohromady, aby mu společně koupili dárek.};  \dicIdiom{taka}[til]{ \dicPhraseIS{taka til}} \dicDirectTranslationCS{uklidit, uklízet} \dicExampleIS{taka til eftir sig í herberginu} \dicExampleCS{uklidit po sobě v~pokoji}; { \dicPhraseIS{taka til e‑ð}} \dicDirectTranslationCS{připravit (co), připravovat (co)} \dicExampleIS{taka til nesti handa henni} \dicExampleCS{připravit pro ni svačinu}; { \dicPhraseIS{taka til fótanna}} \dicDirectTranslationCS{rozběhnout se, dát se do běhu, začít běžet}; { \dicPhraseIS{láta e‑ð til sín taka}} \dicDirectTranslationCS{angažovat se v~(čem)}; { \dicPhraseIS{taka til starfa}} \dicDirectTranslationCS{dát\,/\addthin pustit se do práce}; { \dicPhraseIS{taka e‑ð til sín}} \dicDirectTranslationCS{brát si (co) (k~srdci ap.)} \dicExampleIS{taka orð hans til sín} \dicExampleCS{brát si jeho slova k~srdci};  \dicIdiom{taka}[til við]{ \dicPhraseIS{taka til við að (gera e‑ð)}} \dicDirectTranslationCS{dát\,/\addthin pustit se do (dělání (čeho))} \dicExampleIS{taka til við að syngja} \dicExampleCS{pustit se do zpívání};  \dicIdiom{taka}[undir]{ \dicPhraseIS{taka undir e‑ð}} \dicDirectTranslationCS{přidat\,/\addthin připojit se k~(čemu) (zpěvu ap.)} \dicExampleIS{taka hressilega undir sönginn} \dicExampleCS{přidat se nadšeně ke zpěvu};  \dicIdiom{taka}[undir í]{ \dicPhraseIS{það tekur undir í e‑u}} \dicFlx{impers} \dicDirectTranslationCS{v~(čem) se rozléhá ozvěna} \dicExampleIS{Það tekur undir í fjöllunum.} \dicExampleCS{V~horách se rozléhá ozvěna.};  \dicIdiom{taka}[upp]{ \dicPhraseIS{taka e‑n upp}} \dicDirectTranslationCS{vyvolat (koho), vyzkoušet (koho)}; { \dicPhraseIS{taka upp e‑ð}} {\textbf{a.}} \dicDirectTranslationCS{zvednout (co), zvedat (co)} \dicExampleIS{taka upp stein} \dicExampleCS{zvednout kámen};  {\textbf{b.}} \dicDirectTranslationCS{vyhrabat (co), vykopat (co)} \dicExampleIS{taka upp kartöflur} \dicExampleCS{vykopat brambory};  {\textbf{c.}} \dicDirectTranslationCS{nahrát (co), nahrávat (co)} \dicExampleIS{taka lagið upp á segulband} \dicExampleCS{nahrát písničku na kazetu};  {\textbf{d.}} \dicDirectTranslationCS{otevřít (co), otevírat (co) (láhev ap.)};  {\textbf{e.}} \dicDirectTranslationCS{zavést (co), přijmout (co) (nové metody ap.)};  \dicIdiom{taka}[upp á]{ \dicPhraseIS{taka upp á e‑u}} \dicDirectTranslationCS{začít (co), přijmout (co) (nový zvyk ap.)};  \dicIdiom{taka}[upp]{ \dicPhraseIS{snjóinn tekur upp}} \dicFlx{impers} \dicDirectTranslationCS{sníh taje};  \dicIdiom{taka}[utan af]{ \dicPhraseIS{taka utan af e‑u}} \dicDirectTranslationCS{otevřít (co), otevírat (co) (balík ap.)} \dicExampleIS{taka utan af pakkanum} \dicExampleCS{otevřít balík};  \dicIdiom{taka}[utan um]{ \dicPhraseIS{taka utan um e‑n}} \dicDirectTranslationCS{obejmout (koho), objímat (koho)} \dicExampleIS{taka þétt utan um hana} \dicExampleCS{pevně ji obejmout};  \dicIdiom{taka}[út]{ \dicPhraseIS{taka út e‑ð}} {\textbf{a.}} \dicDirectTranslationCS{vybrat (co), vybírat (co) (peníze ap.)} \dicExampleIS{taka þúsund krónur út úr bankanum} \dicExampleCS{vybrat tisíc korun z~banky};  {\textbf{b.}} \dicDirectTranslationCS{odhadnout (co), ohodnotit (co) (znaleckým posudkem ap.)} \dicExampleIS{taka út húsið} \dicExampleCS{odhadnout dům}; { \dicPhraseIS{taka út refsingu}} \dicDirectTranslationCS{být potrestán};  \dicIdiom{taka}[við]{ \dicPhraseIS{taka við e‑u}} \dicDirectTranslationCS{přijmout (co), přijímat (co) (dopis ap.)}; { \dicPhraseIS{taka við af e‑m}} \dicDirectTranslationCS{vystřídat (koho), vyměnit (koho) (v~práci ap.)};  \dicIdiom{taka}[yfir]{ \dicPhraseIS{taka yfir e‑ð}} {\textbf{a.}} \dicDirectTranslationCS{převzít (co) (řízení ap.)};  {\textbf{b.}} \dicDirectTranslationCS{podmanit si (co), dobýt (co)};  \dicIdiom{takast}{ \dicPhraseIS{e‑m tekst að (gera e‑ð)}} \dicFlx{refl impers} \dicDirectTranslationCS{(komu) se podaří (udělat (co)), (kdo) uspěje v~(dělání (čeho))} \dicExampleIS{Honum tókst að ljúka verkinu.} \dicExampleCS{Podařilo se mu dokončit práci.}; { \dicPhraseIS{e‑að tekst}} \dicFlx{refl} \dicSynonym{heppnast} \dicDirectTranslationCS{(co) se (po)daří} \dicExampleIS{Það tókst sem betur fer.} \dicExampleCS{Naštěstí se to podařilo.};  \dicIdiom{takast}[á]{ \dicPhraseIS{takast e‑ð á hendur}} \dicFlx{refl} \dicDirectTranslationCS{vzít na sebe (co) (práci ap.)}; { \dicPhraseIS{takast á loft}} \dicFlx{refl} \dicDirectTranslationCS{být vržen do vzduchu};  \dicIdiom{takast}[á um]{ \dicPhraseIS{takast á um e‑ð}} \dicFlx{refl} \dicDirectTranslationCS{soupeřit o~(co), zápolit o~(co)};  \dicIdiom{takast}[á við]{ \dicPhraseIS{takast á við e‑ð}} \dicFlx{refl} \dicDirectTranslationCS{potýkat se s~(čím), vypořádat se s~(čím)};  \dicIdiom{takast}[í]{ \dicPhraseIS{takast í hendur}} \dicFlx{refl} \dicDirectTranslationCS{potřást si rukama};  \dicIdiom{takast}[með]{ \dicPhraseIS{það tekst með þeim e‑að}} \dicFlx{refl impers} \dicDirectTranslationCS{(co) vzniká mezi nimi (přátelství ap.)};  \dicIdiom{takast}[til]{ \dicPhraseIS{það tókst (þannig) til}} \dicFlx{refl impers} \dicDirectTranslationCS{(tak) se zdařilo, (tak) se vydařilo}
\dicEntry[takk] \dicTerm{takk} \dicsymFrequent\  \dicIPA{{t\smash{\textsuperscript{h}}}{a}{h}{\r{g}}} \dicPos{inter} \dicDirectTranslationCS{děkuju, děkuji, díky};  \dicPhraseIS{takk fyrir} \dicDirectTranslationCS{děkuju, děkuji, díky};  \dicPhraseIS{takk fyrir e‑ð} \dicDirectTranslationCS{děkuji\,/\addthin díky za (co)};  \dicPhraseIS{takk fyrir mig} \dicDirectTranslationCS{děkuji (za večeři ap.)} \dicIndirectTranslationCS{(užívané po dojedení jídla)};  \dicPhraseIS{takk fyrir síðast} \dicDirectTranslationCS{děkuji za posledně}
\dicEntry[takki] \dicTerm{takk|i} \dicIPA{{t\smash{\textsuperscript{h}}}{a}{h}{\r{\textObardotlessj}}{\textsci}} \dicPos{m}[1] \dicFlx{(‑a, ‑ar)}[8] \dicSynonym{hnappur} \dicDirectTranslationCS{tlačítko, knoflík} \dicExampleIS{ýta á takkann} \dicExampleCS{zmáčknout tlačítko}
\dicEntry[takmark] \dicTerm{tak··|mark} \dicsymFrequent\  \dicIPA{{t\smash{\textsuperscript{h}}}{a}{\textlengthmark}{\r{g}}{m}{a}{\textsubring{r}}{\r{g}}} \dicPos{n}[2] \dicFlx{(‑marks, ‑mörk)}[8] \textbf{1.} \dicDirectTranslationCS{cíl, záměr, úmysl, meta} \dicExampleIS{setja sér ákveðið takmark} \dicExampleCS{stanovit si určitý cíl}  \textbf{2.} \dicPhraseIS{takmörk} \dicFlx{pl} \dicSynonym{mörk\smash{\textsuperscript{1}}} \dicDirectTranslationCS{mez, hranice, omezení, limit};  \dicPhraseIS{e‑að á sér engin takmörk} \dicLangCat{přen.} \dicDirectTranslationCS{(co) přesahuje všechny meze, (co) nezná hranic};  \dicPhraseIS{þekkja takmörk sín} \dicLangCat{přen.} \dicDirectTranslationCS{znát svoje možnosti\,/\addthin hranice}
\dicEntry[takmarka] \dicTerm{tak··mark|a} \dicIPA{{t\smash{\textsuperscript{h}}}{a}{\textlengthmark}{\r{g}}{m}{a}{\textsubring{r}}{\r{g}}{a}} \dicPos{v}[1] \dicFlx{(‑aði)}[13] \dicFlx{acc} \dicDirectTranslationCS{omezit, omezovat, ohraničit, ohraničovat, limitovat} \dicExampleIS{takmarka þátttökuna} \dicExampleCS{omezit účast};  \dicPhraseIS{takmarka e‑ð við e‑ð} \dicDirectTranslationCS{omezit (co) na (co)};  \dicIdiom{takmarkast}{ \dicPhraseIS{takmarkast}} \dicFlx{refl} \dicDirectTranslationCS{být omezen\,/\addthin ohraničen, omezit se, ohraničit se};  \dicPhraseIS{e‑að takmarkast af e‑u} \dicFlx{refl} \dicDirectTranslationCS{(co) je omezeno\,/\addthin ohraničeno (čím)};  \dicPhraseIS{e‑að takmarkast við e‑ð} \dicFlx{refl} \dicDirectTranslationCS{(co) je omezeno\,/\addthin ohraničeno na (co)}
\dicEntry[takmarkaður] \dicTerm{tak··|mark·aður} \dicIPA{{t\smash{\textsuperscript{h}}}{a}{\textlengthmark}{\r{g}}{m}{a}{\textsubring{r}}{\r{g}}{a}{ð}{\textscy}{\textsubring{r}}} \dicPos{adj}[3] \dicFlx{(f ‑mörkuð)}[1] \textbf{1.} \dicDirectTranslationCS{omezený, vymezený, ohraničený, limitovaný}  \textbf{2.} \dicLangCat{přen.} \dicDirectTranslationCS{omezený, hloupý}
\dicEntry[takmarkalaus] \dicTerm{tak·marka··laus} \dicIPA{{t\smash{\textsuperscript{h}}}{a}{\textlengthmark}{\r{g}}{m}{a}{\textsubring{r}}{\r{g}}{a}{l}{\oe i}{s}} \dicPos{adj}[5]\dicFlx{}[-1] \dicSynonym{endalaus} \dicDirectTranslationCS{neomezený, neohraničený, bezmezný} \dicExampleIS{takmarkalaus forvitni} \dicExampleCS{bezmezná zvědavost}
\dicEntry[takmörkun] \textls[15]{\dicTerm{tak··|mörkun} \dicIPA{{t\smash{\textsuperscript{h}}}{a}{\textlengthmark}{\r{g}}{m}{\oe}{\textsubring{r}}{\r{g}}{\textscy}{\textsubring{n}}} \dicPos{f}[7] \dicFlx{(‑mörkunar, ‑markanir)}[11]} \dicDirectTranslationCS{omezení, ohraničení} \dicExampleIS{takmörkun á stærð} \dicExampleCS{omezení na velikost}
\dicEntry[taktfast] \dicTerm{takt··fast} \dicIPA{{t\smash{\textsuperscript{h}}}{a}{x}{\textsubring{d}}{f}{a}{s}{\textsubring{d}}} \dicPos{adv} \dicDirectTranslationCS{rytmicky}
\dicEntry[taktfastur] \dicTerm{takt··|fastur} \dicIPA{{t\smash{\textsuperscript{h}}}{a}{x}{\textsubring{d}}{f}{a}{s}{\textsubring{d}}{\textscy}{\textsubring{r}}} \dicPos{adj}[1] \dicFlx{(f ‑föst)}[11] \dicDirectTranslationCS{rytmický}
\dicEntry[taktík] \dicTerm{taktík} \dicIPA{{t\smash{\textsuperscript{h}}}{a}{x}{\textsubring{d}}{i}{\r{g}}} \dicPos{f}[10] \dicFlx{(‑ur)}[2] \dicDirectTranslationCS{taktika, strategie}
\dicEntry[taktískur] \dicTerm{taktískur} \dicIPA{{t\smash{\textsuperscript{h}}}{a}{x}{\textsubring{d}}{i}{s}{\r{g}}{\textscy}{\textsubring{r}}} \dicPos{adj}[1]\dicFlx{}[-1] \dicDirectTranslationCS{taktický, strategický}
\dicEntry[taktlaus] \dicTerm{takt··laus} \dicIPA{{t\smash{\textsuperscript{h}}}{a}{x}{\textsubring{d}}{l}{\oe i}{s}} \dicPos{adj}[5]\dicFlx{}[-1] \dicDirectTranslationCS{netaktní, neohleduplný}
\dicEntry[taktleysi] \dicTerm{takt··leysi} \dicIPA{{t\smash{\textsuperscript{h}}}{a}{x}{\textsubring{d}}{l}{ei}{s}{\textsci}} \dicPos{n}[2] \dicFlx{(‑s)}[20] \dicDirectTranslationCS{netaktnost, neohleduplnost}
\dicEntry[taktmælir] \dicTerm{takt··mæl|ir} \dicIPA{{t\smash{\textsuperscript{h}}}{a}{x}{\textsubring{d}}{m}{a}{i}{l}{\textsci}{\textsubring{r}}} \dicPos{m}[7] \dicFlx{(‑is, ‑ar)}[1] \dicFieldCat{hud.} \dicDirectTranslationCS{metronom}
\dicEntry[taktur] \dicTerm{takt|ur} \dicsymFrequent\  \dicIPA{{t\smash{\textsuperscript{h}}}{a}{x}{\textsubring{d}}{\textscy}{\textsubring{r}}} \dicPos{m}[6] \dicFlx{(‑s, ‑ar)}[20] \textbf{1.} \dicFieldCat{hud.} \dicSynonym*{hljóðfall} \dicDirectTranslationCS{takt, rytmus} \dicExampleIS{slá takt} \dicExampleCS{odbíjet takt}  \textbf{2.} \dicSynonym{háttalag} \dicDirectTranslationCS{takt, chování} \dicIndirectTranslationCS{(nejčastěji v~pl)} \dicExampleIS{Hann sýndi góða takta.} \dicExampleCS{Projevil takt.}  \textbf{3.} \dicDirectTranslationCS{krok, takt} \dicExampleIS{ganga í takt} \dicExampleCS{držet krok}
\dicEntry[tal] \dicTerm{tal} \dicsymFrequent\  \dicIPA{{t\smash{\textsuperscript{h}}}{a}{\textlengthmark}{\textsubring{l}}} \dicPos{n}[2] \dicFlx{(‑s, töl)}[8] \textbf{1.} \dicSynonym*{það að tala} \dicDirectTranslationCS{mluvení, povídání, konverzování}  \textbf{2.} \dicSynonym*{samræða} \dicDirectTranslationCS{(roz)hovor, konverzace};  \dicPhraseIS{beina talinu að e‑u} \dicDirectTranslationCS{zavést hovor na (co)};  \dicPhraseIS{e‑að berst í tal} \dicFlx{refl} \dicDirectTranslationCS{(co) je zmíněno, (co) se dostane do řečí};  \dicPhraseIS{eiga tal við e‑n} \dicDirectTranslationCS{promluvit si s~(kým)};  \dicPhraseIS{e‑ir taka tal saman} \dicDirectTranslationCS{(kdo) spolu mluví\,/\addthin hovoří};  \dicPhraseIS{færa e‑ð í tal} \dicDirectTranslationCS{přivést řeč na (co)};  \dicPhraseIS{hafa tal af e‑m} \dicDirectTranslationCS{promluvit si s~(kým), kontaktovat se s~(kým)};  \dicPhraseIS{ná tali af e‑m} \dicDirectTranslationCS{kontaktovat se s~(kým)};  \dicPhraseIS{síminn er á tali} \dicDirectTranslationCS{je obsazeno, telefonní linka je obsazená};  \dicPhraseIS{taka e‑n tali} \dicDirectTranslationCS{promluvit si s~(kým)}  \textbf{3.} \dicSynonym{fjöldi} \dicDirectTranslationCS{počet, množství} \dicExampleIS{tíu talsins} \dicExampleCS{v~počtu deseti}
\dicEntry[tala] \dicTerm{tala\smash{\textsuperscript{1}}} \dicsymFrequent\  \dicIPA{{t\smash{\textsuperscript{h}}}{a}{\textlengthmark}{l}{a}} \dicPos{f}[1] \dicFlx{(tölu, tölur)}[14] \textbf{1.} \dicSynonym{erindi} \dicDirectTranslationCS{řeč, proslov, promluva};  \dicPhraseIS{halda tölu} \dicDirectTranslationCS{mít proslov}  \textbf{2.} \dicFieldCat{mat.} \dicDirectTranslationCS{číslo, číslice} \dicExampleIS{leggja tölurnar saman} \dicExampleCS{sčítat čísla};  \dicPhraseIS{blendin tala} \dicFieldCat{mat.} \dicDirectTranslationCS{smíšené číslo};  \dicPhraseIS{heil tala} \dicFieldCat{mat.} \dicDirectTranslationCS{celé číslo};  \dicPhraseIS{jöfn\,/\addthin slétt tala} \dicFieldCat{mat.} \dicDirectTranslationCS{sudé číslo};  \dicPhraseIS{ójöfn tala} \dicFieldCat{mat.} \dicDirectTranslationCS{liché číslo};  \dicPhraseIS{óræð tala} \dicFieldCat{mat.} \dicDirectTranslationCS{iracionální číslo};  \dicPhraseIS{ræð tala} \dicFieldCat{mat.} \dicDirectTranslationCS{racionální číslo}  \textbf{3.} \dicSynonym{talning} \dicDirectTranslationCS{spočítání, počítání};  \dicPhraseIS{hafa tölu á e‑ð} \dicDirectTranslationCS{spočítat (co), mít (co) spočítané};  \dicPhraseIS{kasta tölu á e‑ð} \dicDirectTranslationCS{odhadnout počet (čeho)}  \textbf{4.} \dicFieldCat{jaz.} \dicDirectTranslationCS{číslo}  \textbf{5.} \dicSynonym{hnappur} \dicDirectTranslationCS{knoflík} \dicExampleIS{buxnatala} \dicExampleCS{knoflík od kalhot}
\dicEntry[tala] \dicTerm{tal|a\smash{\textsuperscript{2}}} \dicsymFrequent\  \dicIPA{{t\smash{\textsuperscript{h}}}{a}{\textlengthmark}{l}{a}} \dicPos{v}[1] \dicFlx{(‑aði)}[10] \dicDirectTranslationCS{mluvit, povídat, hovořit} \dicExampleIS{tala um e‑ð} \dicExampleCS{mluvit o~(čem)};  \dicPhraseIS{tala á e‑u} \dicDirectTranslationCS{mluvit (čím)} \dicExampleIS{tala á íslensku} \dicExampleCS{mluvit islandsky};  \dicPhraseIS{tala e‑ð} \dicDirectTranslationCS{mluvit (čím)} \dicExampleIS{tala tékknesku} \dicExampleCS{mluvit česky};  \dicPhraseIS{tala við e‑n} \dicDirectTranslationCS{mluvit s~(kým)};  \dicIdiom{tala}[af]{ \dicPhraseIS{tala af sér}} \dicDirectTranslationCS{podřeknout se};  \dicIdiom{tala}[saman]{ \dicPhraseIS{tala saman}} \dicDirectTranslationCS{mluvit spolu, povídat si};  \dicIdiom{tala}[til]{ \dicPhraseIS{tala e‑n til}} \dicDirectTranslationCS{domluvit (komu), domlouvat (komu)}; { \dicPhraseIS{tala til e‑rs}} \dicDirectTranslationCS{mluvit ke (komu)};  \dicIdiom{tala}[um]{ \dicPhraseIS{tala vel\,/\addthin illa um e‑n}} \dicDirectTranslationCS{mluvit o~(kom) dobře\,/\addthin špatně};  \dicIdiom{tala}[um fyrir]{ \dicPhraseIS{tala um fyrir e‑m}} \dicDirectTranslationCS{rozmluvit\,/\addthin rozmlouvat (komu), vymluvit\,/\addthin vymlouvat (komu)};  \dicIdiom{tala}[út]{ \dicPhraseIS{tala út}} \dicDirectTranslationCS{dopovědět, domluvit};  \dicIdiom{talast}[við]{ \dicPhraseIS{talast við}} \dicFlx{refl} \dicDirectTranslationCS{mluvit spolu, povídat si}
\dicEntry[talaður] \dicTerm{tal··aður} \dicIPA{{t\smash{\textsuperscript{h}}}{a}{\textlengthmark}{l}{a}{ð}{\textscy}{\textsubring{r}}} \dicPos{adj}[3] \dicFlx{(f töluð)}[1] \dicPhraseIS{í þeim töluðum orðum} \dicFlx{adv} \dicDirectTranslationCS{v~tom okamžiku, v~téže chvíli}
\dicEntry[talandi] \dicTerm{tal··|andi\smash{\textsuperscript{1}}} \dicIPA{{t\smash{\textsuperscript{h}}}{a}{\textlengthmark}{l}{a}{n}{\textsubring{d}}{\textsci}} \dicPos{m}[2] \dicFlx{(‑anda, ‑endur)}[1] \dicPhraseIS{hafa liðugan talanda} \dicDirectTranslationCS{mít výřečný jazyk}
\dicEntry[talandi] \dicTerm{tal··andi\smash{\textsuperscript{2}}} \dicIPA{{t\smash{\textsuperscript{h}}}{a}{\textlengthmark}{l}{a}{n}{\textsubring{d}}{\textsci}} \dicPos{adj}[13] \dicFlx{indecl}[1] \dicDirectTranslationCS{umějící mluvit};  \dicPhraseIS{talandi dæmi} \dicDirectTranslationCS{výmluvný příklad}
\dicEntry[taldi] \dicTerm{taldi} \dicIPA{{t\smash{\textsuperscript{h}}}{a}{l}{\textsubring{d}}{\textsci}} \dicPos{v} \dicFlx{ind pf sg 1 pers} \dicLink{telja}
\dicEntry[talfæri] \dicTerm{tal··færi} \dicIPA{{t\smash{\textsuperscript{h}}}{a}{l}{f}{a}{i}{r}{\textsci}} \dicPos{n}[2] \dicFlx{(‑s, ‑)}[14] \dicFieldCat{anat.} \dicDirectTranslationCS{hlasové\,/\addthin fonační ústrojí}
\dicEntry[talgervill] \dicTerm{tal··gerv|ill} \dicIPA{{t\smash{\textsuperscript{h}}}{a}{l}{\r{\textObardotlessj}}{\textepsilon}{r}{v}{\textsci}{\textsubring{d}}{\textsubring{l}}} \dicPos{m}[6] \dicFlx{(‑ils, ‑lar)}[35] \dicFieldCat{poč.} \dicDirectTranslationCS{syntetizátor\,/\addthin syntetizér řeči}
\dicEntry[talhlýðinn] \dicTerm{tal··hlýðinn} \dicIPA{{t\smash{\textsuperscript{h}}}{a}{l}{\textsubring{l}}{i}{ð}{\textsci}{\textsubring{n}}} \dicPos{adj}[6]\dicFlx{}[-2] \dicDirectTranslationCS{(snadno) ovlivnitelný\,/\addthin přemluvitelný (člověk ap.)}
\dicEntry[talhólf] \dicTerm{tal··hólf} \dicIPA{{t\smash{\textsuperscript{h}}}{a}{\textlengthmark}{\textsubring{l}}{h}{ou}{l}{f}} \dicPos{n}[2] \dicFlx{(‑s, ‑)}[5] \dicDirectTranslationCS{hlasová schránka}
\dicEntry[talið] \dicTerm{talið} \dicIPA{{t\smash{\textsuperscript{h}}}{a}{\textlengthmark}{l}{\textsci}{\texttheta}} \dicPos{v} \dicFlx{supin} \dicLink{telja}
\dicEntry[talía] \dicTerm{talí|a} \dicIPA{{t\smash{\textsuperscript{h}}}{a}{\textlengthmark}{l}{i}{j}{a}} \dicPos{f}[1] \dicFlx{(‑u, ‑ur)}[7] \dicFieldCat{techn.} \dicDirectTranslationCS{kladka, kladkostroj}
\dicEntry[Talíbanar] \dicTerm{Talíbanar} \dicIPA{{t\smash{\textsuperscript{h}}}{a}{\textlengthmark}{l}{i}{\textsubring{b}}{a}{n}{a}{\textsubring{r}}} \dicPos{m}[1] \dicFlx{pl}[9] \dicFieldCat{pol.} \dicDirectTranslationCS{Tálibán}
\dicEntry[talk] \dicTerm{talk} \dicIPA{{t\smash{\textsuperscript{h}}}{a}{\textsubring{l}}{\r{g}}} \dicPos{n}[2] \dicFlx{(‑s)}[2] \dicFieldCat{geol.} \dicDirectTranslationCS{mastek, talek}
\dicEntry[talkennari] \dicTerm{tal··kenn·ar|i} \dicIPA{{t\smash{\textsuperscript{h}}}{a}{l}{c\smash{\textsuperscript{h}}}{\textepsilon}{n}{a}{r}{\textsci}} \dicPos{m}[1] \dicFlx{(‑a, ‑ar)}[13] \dicDirectTranslationCS{odborník\,/\addthin odbornice na výuku řeči, logoped(ka)}
\dicEntry[talkúm] \dicTerm{talkúm} \dicIPA{{t\smash{\textsuperscript{h}}}{a}{\textsubring{l}}{\r{g}}{u}{\textsubring{m}}} \dicPos{n}[2] \dicFlx{(‑s)}[2] \dicDirectTranslationCS{mletý mastek}
\dicEntry[Tallinn] \dicTerm{Tallinn} \dicIPA{{t\smash{\textsuperscript{h}}}{a}{l}{\textlengthmark}{\textsci}{\textsubring{n}}} \dicPos{subs} \dicFlx{indecl} \dicDirectTranslationCS{Tallinn} \dicIndirectTranslationCS{(hlavní město Estonska)}
\dicEntry[talmál] \dicTerm{tal··mál} \dicIPA{{t\smash{\textsuperscript{h}}}{a}{l}{m}{au}{\textsubring{l}}} \dicPos{n}[2] \dicFlx{(‑s, ‑)}[5] \dicFieldCat{jaz.} \dicDirectTranslationCS{mluvený\,/\addthin hovorový jazyk}
\dicEntry[talmeinafræði] \dicTerm{tal·meina··fræð|i} \dicIPA{{t\smash{\textsuperscript{h}}}{a}{l}{m}{ei}{n}{a}{f}{r}{a}{i}{ð}{\textsci}} \dicPos{f}[3] \dicFlx{(‑i)}[3] \dicDirectTranslationCS{logopedie}
\dicEntry[talmeinafræðingur] \dicTerm{tal·meina·fræð··ing|ur} \dicIPA{{t\smash{\textsuperscript{h}}}\-{a}\-{l}\-{m}\-{ei}\-{n}\-{a}\-{f}\-{r}\-{a}\-{i}\-{ð}\-{i}\-{\ng}\-{\r{g}}\-{\textscy}\-{\textsubring{r}}\-} \dicPos{m}[6] \dicFlx{(‑s, ‑ar)}[8] \dicDirectTranslationCS{logoped(ka)}
\dicEntry[talnaband] \dicTerm{talna··|band} \dicIPA{{t\smash{\textsuperscript{h}}}{a}{l}{n}{a}{\textsubring{b}}{a}{n}{\textsubring{d}}} \dicPos{n}[2] \dicFlx{(‑bands, ‑bönd)}[8] \dicFieldCat{náb.} \dicDirectTranslationCS{růženec} \dicIndirectTranslationCS{(modlitební pomůcka)}
\dicEntry[talnafræði] \dicTerm{talna··fræð|i} \dicIPA{{t\smash{\textsuperscript{h}}}{a}{l}{n}{a}{f}{r}{a}{i}{ð}{\textsci}} \dicPos{f}[3] \dicFlx{(‑i)}[3] \dicFieldCat{mat.} \dicDirectTranslationCS{aritmetika}
\dicEntry[talnagrind] \dicTerm{talna··grind} \dicIPA{{t\smash{\textsuperscript{h}}}{a}{l}{n}{a}{\r{g}}{r}{\textsci}{n}{\textsubring{d}}} \dicPos{f}[8] \dicFlx{(‑ar, ‑ur)}[1] \dicDirectTranslationCS{(ruční) počítadlo}
\dicEntry[talning] \dicTerm{tal··ning} \dicIPA{{t\smash{\textsuperscript{h}}}{a}{l}{n}{i}{\ng}{\r{g}}} \dicPos{f}[4] \dicFlx{(‑ar, ‑ar)}[5] \dicDirectTranslationCS{počítání, sčítání} \dicExampleIS{talning á birgðum} \dicExampleCS{počítání zásob}
\dicEntry[talsamband] \dicTerm{tal··sam·band} \dicIPA{{t\smash{\textsuperscript{h}}}{a}{l}{s}{a}{m}{\textsubring{b}}{a}{n}{\textsubring{d}}} \dicPos{n}[2] \dicFlx{(‑s)}[2] \dicDirectTranslationCS{telefonické\,/\addthin telefonní spojení}
\dicEntry[talset] \dicTerm{tal··set} \dicIPA{{t\smash{\textsuperscript{h}}}{a}{l}{s}{\textepsilon}{\textsubring{d}}} \dicPos{v} \dicFlx{ind praes sg 1 pers} \dicLink{talsetja}
\dicEntry[talsetja] \dicTerm{tal··|setja} \dicIPA{{t\smash{\textsuperscript{h}}}{a}{l}{s}{\textepsilon}{\textsubring{d}}{j}{a}} \dicPos{v}[4] \dicFlx{(‑set, ‑setti, ‑settum, ‑setti, ‑sett)}[26] \dicFlx{acc} \dicDirectTranslationCS{(na)dabovat} \dicExampleIS{talsetja kvikmyndina á íslensku} \dicExampleCS{(na)dabovat film do islandštiny}
\dicEntry[talsetning] \textls[10]{\dicTerm{tal··set·ning} \dicIPA{{t\smash{\textsuperscript{h}}}{a}{l}{s}{\textepsilon}{h}{\textsubring{d}}{n}{i}{\ng}{\r{g}}} \dicPos{f}[4] \dicFlx{(‑ar)}[7] \dicDirectTranslationCS{dabing, dabink}} \dicExampleIS{Myndin er sýnd með íslenskri talsetningu.} \dicExampleCS{Film je (promítán) s~islandským dabingem.}
\dicEntry[talsett] \dicTerm{tal··sett} \dicIPA{{t\smash{\textsuperscript{h}}}{a}{l}{s}{\textepsilon}{h}{\textsubring{d}}} \dicPos{v} \dicFlx{supin} \dicLink{talsetja}
\dicEntry[talsetti] \dicTerm{tal··setti} \dicIPA{{t\smash{\textsuperscript{h}}}{a}{l}{s}{\textepsilon}{h}{\textsubring{d}}{\textsci}} \dicPos{v} \dicFlx{ind\,/\addthin con pf sg 1 pers} \dicLink{talsetja}
\dicEntry[talsettum] \dicTerm{tal··settum} \dicIPA{{t\smash{\textsuperscript{h}}}{a}{l}{s}{\textepsilon}{h}{\textsubring{d}}{\textscy}{\textsubring{m}}} \dicPos{v} \dicFlx{ind pf pl 1 pers} \dicLink{talsetja}
\dicEntry[talsháttur] \dicTerm{tals··|háttur} \dicIPA{{t\smash{\textsuperscript{h}}}{a}{l}{s}{h}{au}{h}{\textsubring{d}}{\textscy}{\textsubring{r}}} \dicPos{m}[12] \dicFlx{(‑háttar, ‑hættir)}[7] \dicSynonym{orðtak} \dicDirectTranslationCS{slovní obrat, idiom}
\dicEntry[talsins] \dicTerm{talsins} \dicIPA{{t\smash{\textsuperscript{h}}}{a}{l}{s}{\textsci}{n}{s}} \dicPos{adv} \dicDirectTranslationCS{celkem, dohromady, úhrnem}
\dicEntry[talsmaður] \dicTerm{tals··|maður} \dicIPA{{t\smash{\textsuperscript{h}}}{a}{l}{s}{m}{a}{ð}{\textscy}{\textsubring{r}}} \dicPos{m}[13] \dicFlx{(‑manns, ‑menn)}[2] \dicDirectTranslationCS{mluvčí}
\dicEntry[talsmáti] \dicTerm{tals··mát|i} \dicIPA{{t\smash{\textsuperscript{h}}}{a}{l}{s}{m}{au}{\textsubring{d}}{\textsci}} \dicPos{m}[1] \dicFlx{(‑a)}[3] \dicSynonym{orðbragð} \dicDirectTranslationCS{způsob řeči\,/\addthin vyjadřování}
\dicEntry[talstöð] \dicTerm{tal··stöð} \dicIPA{{t\smash{\textsuperscript{h}}}{a}{l}{s}{\textsubring{d}}{\oe}{\texttheta}} \dicPos{f}[6] \dicFlx{(‑var, ‑var)}[1] \dicDirectTranslationCS{vysílačka}
\dicEntry[talsverður] \dicTerm{tals··verður} \dicsymFrequent\  \dicIPA{{t\smash{\textsuperscript{h}}}{a}{l}{s}{v}{\textepsilon}{r}{ð}{\textscy}{\textsubring{r}}} \dicPos{adj}[2]\dicFlx{}[-4] \dicSynonym{töluverður} \dicDirectTranslationCS{značný, podstatný, zásadní} \dicExampleIS{vekja talsverða furðu} \dicExampleCS{vzbudit značný úžas}
\dicEntry[talsvert] \dicTerm{tals··vert} \dicsymFrequent\  \dicIPA{{t\smash{\textsuperscript{h}}}{a}{l}{s}{v}{\textepsilon}{\textsubring{r}}{\textsubring{d}}} \dicPos{adv} \dicSynonym{býsna} \dicDirectTranslationCS{značně, podstatně, zásadně} \dicExampleIS{Þessi aðferð er talsvert tímafrek.} \dicExampleCS{Tento postup je značně zdlouhavý.}
\dicEntry[talþjálfun] \dicTerm{tal··þjálf|un} \dicIPA{{t\smash{\textsuperscript{h}}}{a}{l}{\texttheta}{j}{au}{l}{v}{\textscy}{\textsubring{n}}} \dicPos{f}[7] \dicFlx{(‑unar)}[9] \dicDirectTranslationCS{mluvní cvičení, cvičení výslovnosti}
\dicEntry[tambúrína] \dicTerm{tambúrín|a} \dicIPA{{t\smash{\textsuperscript{h}}}{a}{m}{\textsubring{b}}{u}{r}{i}{n}{a}} \dicPos{f}[1] \dicFlx{(‑u, ‑ur)}[7] \dicFieldCat{hud.} \dicDirectTranslationCS{tamburína}
\dicEntry[tamdi] \dicTerm{tamdi} \dicIPA{{t\smash{\textsuperscript{h}}}{a}{m}{\textsubring{d}}{\textsci}} \dicPos{v} \dicFlx{ind pf sg 1 pers} \dicLink{temja}
\dicEntry[tamið] \dicTerm{tamið} \dicIPA{{t\smash{\textsuperscript{h}}}{a}{\textlengthmark}{m}{\textsci}{\texttheta}} \dicPos{v} \dicFlx{supin} \dicLink{temja}
\dicEntry[taminn] \dicTerm{taminn} \dicIPA{{t\smash{\textsuperscript{h}}}{a}{\textlengthmark}{m}{\textsci}{\textsubring{n}}} \dicPos{adj}[6]\dicFlx{}[-8] \dicDirectTranslationCS{ochočený, krotký, zkrocený}
\dicEntry[tamning] \dicTerm{tam··ning} \dicIPA{{t\smash{\textsuperscript{h}}}{a}{m}{n}{i}{\ng}{\r{g}}} \dicPos{f}[4] \dicFlx{(‑ar, ‑ar)}[5] \dicDirectTranslationCS{(z)krocení, ochočení, ochočování} \dicExampleIS{tamning hesta} \dicExampleCS{krocení koní}
\dicEntry[tamningamaður] \dicTerm{tamninga··|maður} \dicIPA{{t\smash{\textsuperscript{h}}}{a}{m}{n}{i}{\ng}{\r{g}}{a}{m}{a}{ð}{\textscy}{\textsubring{r}}} \dicPos{m}[13] \dicFlx{(‑manns, ‑menn)}[2] \dicDirectTranslationCS{krotitel(ka) (koní ap.)}
\dicEntry[tamur] \dicTerm{tamur} \dicIPA{{t\smash{\textsuperscript{h}}}{a}{\textlengthmark}{m}{\textscy}{\textsubring{r}}} \dicPos{adj}[1] \dicFlx{(f töm)}[2] \dicSynonym{vanur\smash{\textsuperscript{2}}} \dicDirectTranslationCS{přivyklý, zvyklý, zažitý};  \dicPhraseIS{e‑m er tamt að (gera e‑ð)} \dicFlx{impers} \dicDirectTranslationCS{(kdo) má zažité (dělání (čeho)), (kdo) je zvyklý (dělat (co))};  \dicProverb\  \dicPhraseIS{Það er tungunni tamast sem hjartanu er kærast.} \dicLangCat{přís.} \dicDirectTranslationCS{Co na srdci, to na jazyku.}
\dicEntry[tandurhreinn] \dicTerm{tandur··hreinn} \dicIPA{{t\smash{\textsuperscript{h}}}{a}{n}{\textsubring{d}}{\textscy}{\textsubring{r}}{ei}{\textsubring{d}}{\textsubring{n}}} \dicPos{adj}[7]\dicFlx{}[-2] \dicDirectTranslationCS{naprosto\,/\addthin absolut\-ně čistý} \dicExampleIS{tandurhreinn bíll} \dicExampleCS{absolutně čisté auto}
\dicEntry[tangar] \dicTerm{tangar} \dicIPA{{t\smash{\textsuperscript{h}}}{au}{\ng}{\r{g}}{a}{\textsubring{r}}} \dicPos{f} \dicFlx{sg gen} \dicLink{töng}
\dicEntry[tangarfæðing] \dicTerm{tangar··fæð·ing} \dicIPA{{t\smash{\textsuperscript{h}}}{au}{\ng}{\r{g}}{a}{\textsubring{r}}{f}{a}{i}{ð}{i}{\ng}{\r{g}}} \dicPos{f}[4] \dicFlx{(‑ar, ‑ar)}[5] \dicFieldCat{med.} \dicDirectTranslationCS{klešťový porod, porod kleštěmi}
\dicEntry[tangarhald] \dicTerm{tangar··hald} \dicIPA{{t\smash{\textsuperscript{h}}}{au}{\ng}{\r{g}}{a}{\textsubring{r}}{h}{a}{l}{\textsubring{d}}} \dicPos{n}[2] \dicFlx{(‑s)}[2] \dicPhraseIS{hafa tangarhald á e‑m} \dicLangCat{přen.} \dicDirectTranslationCS{mít\,/\addthin držet (koho) v~hrsti};  \dicPhraseIS{ná tangarhaldi á e‑u} \dicDirectTranslationCS{zmocnit se (čeho), uchvátit (co), uzurpovat (co)}
\dicEntry[tangens] \dicTerm{tangens} \dicIPA{{t\smash{\textsuperscript{h}}}{a}{n}{\r{\textObardotlessj}}{\textepsilon}{n}{s}} \dicPos{m}[4] \dicFlx{(‑, ‑ar)}[31] \dicFieldCat{mat.} \dicDirectTranslationCS{tangens}
\dicEntry[tangi] \dicTerm{tang|i} \dicIPA{{t\smash{\textsuperscript{h}}}{au}{\textltailn}{\r{\textObardotlessj}}{\textsci}} \dicPos{m}[1] \dicFlx{(‑a, ‑ar)}[8] \dicSynonym{oddi} \dicDirectTranslationCS{výběžek (země), ostroh}
\dicEntry[tangó] \dicTerm{tangó} \dicIPA{{t\smash{\textsuperscript{h}}}{au}{\ng}{\r{g}}{ou}} \dicPos{m}[4] \dicFlx{(‑s, ‑ar)}[14] \dicFieldCat{hud.} \dicDirectTranslationCS{tango}
\dicEntry[tankbíll] \dicTerm{tank··bíl|l} \dicIPA{{t\smash{\textsuperscript{h}}}{au}{\r{\ng}}{\r{g}}{\textsubring{b}}{i}{\textsubring{d}}{\textsubring{l}}} \dicPos{m}[6] \dicFlx{(‑s, ‑ar)}[48] \dicDirectTranslationCS{cisterna} \dicIndirectTranslationCS{(vůz s~nádrží na kapaliny nebo sypké látky)}
\dicEntry[tankskip] \dicTerm{tank··skip} \dicIPA{{t\smash{\textsuperscript{h}}}{au}{\r{\ng}}{\r{g}}{s}{\r{\textObardotlessj}}{\textsci}{\textsubring{b}}} \dicPos{n}[2] \dicFlx{(‑s, ‑)}[5] \dicDirectTranslationCS{tanker, tanková\,/\addthin cisternová loď}
\dicEntry[tankur] \dicTerm{tank|ur} \dicIPA{{t\smash{\textsuperscript{h}}}{au}{\r{\ng}}{\r{g}}{\textscy}{\textsubring{r}}} \dicPos{m}[6] \dicFlx{(‑s, ‑ar)}[11] \dicDirectTranslationCS{nádrž, cisterna, tank} \dicExampleIS{mjólkurtankur} \dicExampleCS{cisterna na mléko}
\dicEntry[tannar] \dicTerm{tannar} \dicIPA{{t\smash{\textsuperscript{h}}}{a}{n}{\textlengthmark}{a}{\textsubring{r}}} \dicPos{f} \dicFlx{sg gen} \dicLink{tönn}
\dicEntry[tannbergshljóð] \dicTerm{tann·bergs··hljóð} \dicIPA{{t\smash{\textsuperscript{h}}}{a}{n}{\textsubring{b}}{\textepsilon}{r}{\r{g}}{s}{\textsubring{l}}{j}{ou}{\texttheta}} \dicPos{n}[2] \dicFlx{(‑s, ‑)}[5] \dicFieldCat{jaz.} \dicDirectTranslationCS{alveolára, předodásňová\,/\addthin dásňová souhláska}
\dicEntry[tannbergsmæltur] \dicTerm{tann·bergs··mæltur} \dicIPA{{t\smash{\textsuperscript{h}}}{a}{n}{\textsubring{b}}{\textepsilon}{r}{\r{g}}{s}{m}{a}{i}{\textsubring{l}}{\textsubring{d}}{\textscy}{\textsubring{r}}} \dicPos{adj}[1]\dicFlx{}[-13] \dicFieldCat{jaz.} \dicDirectTranslationCS{předodásňový, alveolární}
\dicEntry[tannbursta] \dicTerm{tann··burst|a} \dicIPA{{t\smash{\textsuperscript{h}}}{a}{n}{\textsubring{b}}{\textscy}{\textsubring{r}}{s}{\textsubring{d}}{a}} \dicPos{v}[1] \dicFlx{(‑aði)}[34] \dicFlx{acc} \dicDirectTranslationCS{(vy)čistit zuby};  \dicPhraseIS{tannbursta sig} \dicDirectTranslationCS{(vy)čistit si zuby}
\dicEntry[tannbursti] \dicTerm{tann··burst|i} \dicIPA{{t\smash{\textsuperscript{h}}}{a}{n}{\textsubring{b}}{\textscy}{\textsubring{r}}{s}{\textsubring{d}}{\textsci}} \dicPos{m}[1] \dicFlx{(‑a, ‑ar)}[1] \dicDirectTranslationCS{zubní kartáček, kartáček na zuby}
\dicEntry[tannburstun] \dicTerm{tann··burst|un} \dicIPA{{t\smash{\textsuperscript{h}}}{a}{n}{\textsubring{b}}{\textscy}{\textsubring{r}}{s}{\textsubring{d}}{\textscy}{\textsubring{n}}} \dicPos{f}[7] \dicFlx{(‑unar)}[9] \dicDirectTranslationCS{(vy)čištění zubů}
\dicEntry[tannfylling] \dicTerm{tann··fyll·ing} \dicIPA{{t\smash{\textsuperscript{h}}}{a}{n}{f}{\textsci}{\textsubring{d}}{l}{i}{\ng}{\r{g}}} \dicPos{f}[4] \dicFlx{(‑ar, ‑ar)}[5] \dicDirectTranslationCS{zubní výplň\,/\addthin plomba}
\dicEntry[tanngómmæltur] \dicTerm{tann·góm··mæltur} \dicIPA{{t\smash{\textsuperscript{h}}}{a}{n}{\r{g}}{ou}{m}{a}{i}{\textsubring{l}}{\textsubring{d}}{\textscy}{\textsubring{r}}} \dicPos{adj}[1]\dicFlx{}[-13] \dicFieldCat{jaz.} \dicDirectTranslationCS{zadodásňový, postalveolární}
\dicEntry[tannhjól] \dicTerm{tann··hjól} \dicIPA{{t\smash{\textsuperscript{h}}}{a}{n}{\c{c}}{ou}{\textsubring{l}}} \dicPos{n}[2] \dicFlx{(‑s, ‑)}[5] \dicDirectTranslationCS{ozubené kolo}
\dicEntry[tannhljóð] \dicTerm{tann··hljóð} \dicIPA{{t\smash{\textsuperscript{h}}}{a}{n}{\textsubring{l}}{j}{ou}{\texttheta}} \dicPos{n}[2] \dicFlx{(‑s, ‑)}[5] \dicFieldCat{jaz.} \dicDirectTranslationCS{dentála, zubodásňová\,/\addthin zubní souhláska}
\dicEntry[tannhold] \dicTerm{tann··hold} \dicIPA{{t\smash{\textsuperscript{h}}}{a}{n}{h}{\textopeno}{l}{\textsubring{d}}} \dicPos{n}[2] \dicFlx{(‑s)}[2] \dicFieldCat{anat.} \dicDirectTranslationCS{dáseň}
\dicEntry[tannholdsbólga] \dicTerm{tann·holds··bólg|a} \dicIPA{{t\smash{\textsuperscript{h}}}{a}{n}{h}{\textopeno}{l}{\textsubring{d}}{s}{\textsubring{b}}{ou}{l}{\r{g}}{a}} \dicPos{f}[1] \dicFlx{(‑u)}[5] \dicFieldCat{med. } \dicDirectTranslationCS{zánět dásní, gingivitida}
\dicEntry[tannhvalur] \dicTerm{tann··hval|ur} \dicIPA{{t\smash{\textsuperscript{h}}}{a}{n}{k\smash{\textsuperscript{h}}}{v}{a}{l}{\textscy}{\textsubring{r}}} \dicPos{m}[9] \dicFlx{(‑s, ‑ir)}[9] \dicFieldCat{zool.} \dicDirectTranslationCS{kytovec z~podřádu ozubených} \textit{(l.~{\textLA{Odontoceti}})}  \dicsymPhoto\ 
\dicFigure{ds_image_tannhvalur_0_1.jpg}{Tannhvalur}{Tannhvalur - Jessica Redfern, PD}
\dicEntry[tannkrem] \dicTerm{tann··krem} \dicIPA{{t\smash{\textsuperscript{h}}}{a}{n}{k\smash{\textsuperscript{h}}}{r}{\textepsilon}{\textsubring{m}}} \dicPos{n}[2] \dicFlx{(‑s, ‑)}[5] \dicDirectTranslationCS{zubní pasta}
\dicEntry[tannkróna] \dicTerm{tann··krón|a} \dicIPA{{t\smash{\textsuperscript{h}}}{a}{n}{k\smash{\textsuperscript{h}}}{r}{ou}{n}{a}} \dicPos{f}[1] \dicFlx{(‑u, ‑ur)}[7] \dicDirectTranslationCS{zubní korunka}
\dicEntry[tannkýli] \dicTerm{tann··kýli} \dicIPA{{t\smash{\textsuperscript{h}}}{a}{n}{c\smash{\textsuperscript{h}}}{i}{l}{\textsci}} \dicPos{n}[2] \dicFlx{(‑s, ‑)}[14] \dicFieldCat{med.} \dicDirectTranslationCS{otok dásní, parulis}
\dicEntry[tannlaus] \dicTerm{tann··laus} \dicIPA{{t\smash{\textsuperscript{h}}}{a}{n}{l}{\oe i}{s}} \dicPos{adj}[5]\dicFlx{}[-1] \dicDirectTranslationCS{bezzubý, (jsoucí) bez zubů} \dicExampleIS{tannlaust gamalmenni} \dicExampleCS{bezzubý stařík}
\dicEntry[tannlæknadeild] \dicTerm{tann·lækna··deild} \dicIPA{{t\smash{\textsuperscript{h}}}{a}{n}{l}{a}{i}{h}{\r{g}}{n}{a}{\textsubring{d}}{ei}{l}{\textsubring{d}}} \dicPos{f}[7] \dicFlx{(‑ar, ‑ir)}[1] \dicDirectTranslationCS{fakulta zubního lékařství}
\dicEntry[tannlæknakostnaður] \dicTerm{tann·lækna··kost·nað|ur} \dicIPA{{t\smash{\textsuperscript{h}}}\-{a}\-{n}\-{l}\-{a}\-{i}\-{h}\-{\r{g}}\-{n}\-{a}\-{k\smash{\textsuperscript{h}}}\-{\textopeno}\-{s}\-{\textsubring{d}}\-{n}\-{a}\-{ð}\-{\textscy}\-{\textsubring{r}}\-} \dicPos{m}[10] \dicFlx{(‑ar)}[9] \dicDirectTranslationCS{náklady na ošetření zubním lékařem}
\dicEntry[tannlæknastofa] \dicTerm{tann·lækna··stof|a} \dicIPA{{t\smash{\textsuperscript{h}}}{a}{n}{l}{a}{i}{h}{\r{g}}{n}{a}{s}{\textsubring{d}}{\textopeno}{v}{a}} \dicPos{f}[1] \dicFlx{(‑u, ‑ur)}[7] \dicDirectTranslationCS{zubní ordinace}
\dicEntry[tannlæknavakt] \dicTerm{tann·lækna··vakt} \dicIPA{{t\smash{\textsuperscript{h}}}{a}{n}{l}{a}{i}{h}{\r{g}}{n}{a}{v}{a}{x}{\textsubring{d}}} \dicPos{f}[7] \dicFlx{(‑ar, ‑ir)}[1] \dicDirectTranslationCS{zubní pohotovost}
\dicEntry[tannlækning] \dicTerm{tann··lækn·ing} \dicIPA{{t\smash{\textsuperscript{h}}}{a}{n}{l}{a}{i}{h}{\r{g}}{n}{i}{\ng}{\r{g}}} \dicPos{f}[4] \dicFlx{(‑ar, ‑ar)}[5] \textbf{1.} \dicDirectTranslationCS{zubní léčení}  \textbf{2.} \dicPhraseIS{tannlækningar} \dicFlx{pl} \dicSynonym{tannlæknisfræði} \dicDirectTranslationCS{zubní lékařství, stomatologie}
\dicEntry[tannlæknir] \dicTerm{tann··lækn|ir} \dicIPA{{t\smash{\textsuperscript{h}}}{a}{n}{l}{a}{i}{h}{\r{g}}{n}{\textsci}{\textsubring{r}}} \dicPos{m}[7] \dicFlx{(‑is, ‑ar)}[1] \dicDirectTranslationCS{zubař(ka), stomatolog, stomatoložka}
\dicEntry[tannlæknisfræði] \dicTerm{tann·læknis··fræð|i} \dicIPA{{t\smash{\textsuperscript{h}}}{a}{n}{l}{a}{i}{h}{\r{g}}{n}{\textsci}{s}{f}{r}{a}{i}{ð}{\textsci}} \dicPos{f}[3] \dicFlx{(‑i)}[3] \dicDirectTranslationCS{stomatologie, zubní lékařství}
\dicEntry[tannmæltur] \dicTerm{tann··mæltur} \dicIPA{{t\smash{\textsuperscript{h}}}{a}{n}{m}{a}{i}{\textsubring{l}}{\textsubring{d}}{\textscy}{\textsubring{r}}} \dicPos{adj}[1]\dicFlx{}[-13] \dicFieldCat{jaz.} \dicDirectTranslationCS{dentální, zubodásňový, zubní}
\dicEntry[tannpína] \dicTerm{tann··pín|a} \dicIPA{{t\smash{\textsuperscript{h}}}{a}{n}{p\smash{\textsuperscript{h}}}{i}{n}{a}} \dicPos{f}[1] \dicFlx{(‑u)}[5] \dicDirectTranslationCS{bolest zubu}
\dicEntry[tannrétting] \dicTerm{tann··rétt·ing} \dicIPA{{t\smash{\textsuperscript{h}}}{a}{n}{r}{j}{\textepsilon}{h}{\textsubring{d}}{i}{\ng}{\r{g}}} \dicPos{f}[4] \dicFlx{(‑ar, ‑ar)}[5] \dicFieldCat{med.} \dicDirectTranslationCS{ortodoncie, ortodontie}
\dicEntry[tannrót] \dicTerm{tann··|rót} \dicIPA{{t\smash{\textsuperscript{h}}}{a}{n}{r}{ou}{\textsubring{d}}} \dicPos{f}[8] \dicFlx{(‑rótar, ‑rætur)}[5] \dicDirectTranslationCS{zubní kořen}
\dicEntry[tannskán] \dicTerm{tann··skán} \dicIPA{{t\smash{\textsuperscript{h}}}{a}{n}{s}{\r{g}}{au}{\textsubring{n}}} \dicPos{f}[7] \dicFlx{(‑ar)}[3] \dicDirectTranslationCS{zubní plak}
\dicEntry[tannskemmd] \dicTerm{tann··skemmd} \dicIPA{{t\smash{\textsuperscript{h}}}{a}{n}{s}{\r{\textObardotlessj}}{\textepsilon}{m}{\textsubring{d}}} \dicPos{f}[7] \dicFlx{(‑ar, ‑ir)}[1] \dicDirectTranslationCS{zubní kaz}
\dicEntry[tannsmiður] \dicTerm{tann··smið|ur} \dicIPA{{t\smash{\textsuperscript{h}}}{a}{n}{s}{m}{\textsci}{ð}{\textscy}{\textsubring{r}}} \dicPos{m}[9] \dicFlx{(‑s, ‑ir)}[8] \dicDirectTranslationCS{zubní technik\,/\addthin technička, zubní laborant(ka)}
\dicEntry[tannsmíði] \dicTerm{tann··smíð|i} \dicIPA{{t\smash{\textsuperscript{h}}}{a}{n}{s}{m}{i}{ð}{\textsci}} \dicPos{f}[2] \dicFlx{(‑i, ‑ar)}[1] \dicDirectTranslationCS{zubní technika}
\dicEntry[tannsteinn] \dicTerm{tann··stein|n} \dicIPA{{t\smash{\textsuperscript{h}}}{a}{n}{s}{\textsubring{d}}{ei}{\textsubring{d}}{\textsubring{n}}} \dicPos{m}[6] \dicFlx{(‑s, ‑ar)}[42] \dicFieldCat{med.} \dicDirectTranslationCS{zubní kámen}
\dicEntry[tannstöngull] \dicTerm{tann··stöng|ull} \dicIPA{{t\smash{\textsuperscript{h}}}{a}{n}{s}{\textsubring{d}}{\oe i}{\ng}{\r{g}}{\textscy}{\textsubring{d}}{\textsubring{l}}} \dicPos{m}[6] \dicFlx{(‑uls, ‑lar)}[31] \dicDirectTranslationCS{párátko (do zubů)}
\dicEntry[tannvarahljóð] \dicTerm{tann·vara··hljóð} \dicIPA{{t\smash{\textsuperscript{h}}}{a}{n}{v}{a}{r}{a}{\textsubring{l}}{j}{ou}{\texttheta}} \dicPos{n}[2] \dicFlx{(‑s, ‑)}[5] \dicFieldCat{jaz.} \dicDirectTranslationCS{labiodentála, retozubní souhláska}
\dicEntry[tannvaramæltur] \dicTerm{tann·vara··mæltur} \dicIPA{{t\smash{\textsuperscript{h}}}{a}{n}{v}{a}{r}{a}{m}{a}{i}{\textsubring{l}}{\textsubring{d}}{\textscy}{\textsubring{r}}} \dicPos{adj}[1]\dicFlx{}[-13] \dicFieldCat{jaz.} \dicDirectTranslationCS{labiodentální, retozubní}
\dicEntry[tannviðgerðir] \dicTerm{tann··við·gerðir} \dicIPA{{t\smash{\textsuperscript{h}}}{a}{n}{v}{\textsci}{ð}{\r{\textObardotlessj}}{\textepsilon}{r}{ð}{\textsci}{\textsubring{r}}} \dicPos{f}[7] \dicFlx{pl}[2] \dicDirectTranslationCS{ošetření zubním lékařem}
\dicEntry[Tansani] \dicTerm{Tansan|i} \dicIPA{{t\smash{\textsuperscript{h}}}{a}{n}{s}{a}{n}{\textsci}} \dicPos{m}[3] \dicFlx{(‑a, ‑ar\,/\addthin ‑ir)}[4] \dicDirectTranslationCS{Tanzaňan(ka), Tanzanec, Tanzanka}
\dicEntry[Tansanía] \dicTerm{Tansaní|a} \dicIPA{{t\smash{\textsuperscript{h}}}{a}{n}{s}{a}{n}{i}{j}{a}} \dicPos{f}[1] \dicFlx{(‑u)}[6] \dicFieldCat{geog.} \dicDirectTranslationCS{Tanzanie}
\dicEntry[tansanískur] \dicTerm{tansanískur} \dicIPA{{t\smash{\textsuperscript{h}}}{a}{n}{s}{a}{n}{i}{s}{\r{g}}{\textscy}{\textsubring{r}}} \dicPos{adj}[1]\dicFlx{}[-6] \dicDirectTranslationCS{tanzanský}
\dicEntry[tap] \dicTerm{tap} \dicIPA{{t\smash{\textsuperscript{h}}}{a}{\textlengthmark}{\textsubring{b}}} \dicPos{n}[2] \dicFlx{(‑s, töp)}[8] \textbf{1.} \dicSynonym*{fjárhagslegur skaði} \dicDirectTranslationCS{ztráta, prodělek, deficit};  \dicPhraseIS{reka e‑ð með tapi} \dicDirectTranslationCS{vést (co) se ztrátou (firmu ap.)};  \dicPhraseIS{verða fyrir tapi} \dicDirectTranslationCS{prodělat, ztratit};  \dicPhraseIS{það er tap á e‑u} \dicFlx{impers} \dicDirectTranslationCS{(co) je ztrátové\,/\addthin prodělečné}  \textbf{2.} \dicSynonym{ósigur} \dicDirectTranslationCS{prohra, porážka} \dicExampleIS{vinningur og tap} \dicExampleCS{výhra a~prohra}
\dicEntry[tapa] \dicTerm{tap|a} \dicsymFrequent\  \dicIPA{{t\smash{\textsuperscript{h}}}{a}{\textlengthmark}{\textsubring{b}}{a}} \dicPos{v}[1] \dicFlx{(‑aði)}[13] \dicFlx{dat} \textbf{1.} \dicSynonym{missa} \dicDirectTranslationCS{ztratit, ztrácet, přijít\,/\addthin přicházet o~(peníze ap.)} \dicExampleIS{tapa stórfé á viðskiptunum} \dicExampleCS{ztratit obrovské množství peněz při obchodování}  \textbf{2.} \dicSynonym*{bíða ósigur} \dicDirectTranslationCS{prohrát, prohrávat, být poražen} \dicExampleIS{tapa leiknum} \dicExampleCS{prohrát hru}  \textbf{3.} \dicPhraseIS{tapa af e‑u} \dicDirectTranslationCS{nestihnout (co), zpozdit se na (co) (autobus ap.)};  \dicIdiom{tapast}{ \dicPhraseIS{tapast}} \dicFlx{refl} \dicDirectTranslationCS{ztratit se, zapomenout se};  \dicIdiom{tapa}{ \dicPhraseIS{tapað\,/\addthin fundið}} \dicDirectTranslationCS{ztráty a~nálezy}
\dicEntry[tapír] \dicTerm{tapír} \dicIPA{{t\smash{\textsuperscript{h}}}{a}{\textlengthmark}{\textsubring{b}}{i}{\textsubring{r}}} \dicPos{m}[4] \dicFlx{(‑s, ‑ar)}[14] \dicFieldCat{zool.} \dicDirectTranslationCS{tapír} \textit{(l.~{\textLA{Tapirus}})}
\dicEntry[tappa] \dicTerm{tapp|a} \dicIPA{{t\smash{\textsuperscript{h}}}{a}{h}{\textsubring{b}}{a}} \dicPos{v}[1] \dicFlx{(‑aði)}[13] \textbf{1.} \dicPhraseIS{tappa á e‑ð} \dicDirectTranslationCS{načepovat (co), napustit (co), napouštět (co), stočit (co) (víno do láhve ap.)}  \textbf{2.} \dicPhraseIS{tappa af e‑u} \dicDirectTranslationCS{vyčepovat (co), vypustit (co) (benzín z~nádrže ap.)}
\dicEntry[tappatogari] \dicTerm{tappa··tog·ar|i} \dicIPA{{t\smash{\textsuperscript{h}}}{a}{h}{\textsubring{b}}{a}{t\smash{\textsuperscript{h}}}{\textopeno}{\textbabygamma}{a}{r}{\textsci}} \dicPos{m}[1] \dicFlx{(‑a, ‑ar)}[13] \dicDirectTranslationCS{vývrtka} \dicIndirectTranslationCS{(nástroj na vytahování zátek)}
\dicEntry[tappi] \dicTerm{tapp|i} \dicIPA{{t\smash{\textsuperscript{h}}}{a}{h}{\textsubring{b}}{\textsci}} \dicPos{m}[1] \dicFlx{(‑a, ‑ar)}[8] \dicDirectTranslationCS{zátka, špunt, uzávěr} \dicExampleIS{skrúfa tappann af flöskunni} \dicExampleCS{odšroubovat zátku láhve}
\dicEntry[taprekstur] \dicTerm{tap··rekst|ur} \dicIPA{{t\smash{\textsuperscript{h}}}{a}{\textlengthmark}{\textsubring{b}}{r}{\textepsilon}{x}{s}{\textsubring{d}}{\textscy}{\textsubring{r}}} \dicPos{m}[5] \dicFlx{(‑urs\,/\addthin ‑rar)}[10] \dicFieldCat{ekon.} \dicDirectTranslationCS{provozní ztráta, ztrátová operace}
\dicEntry[tarantúla] \dicTerm{tarantúl|a} \dicIPA{{t\smash{\textsuperscript{h}}}{a}{\textlengthmark}{r}{a}{\textsubring{n}}{\textsubring{d}}{u}{l}{a}} \dicPos{f}[1] \dicFlx{(‑u, ‑ur)}[19] \dicFieldCat{zool.} \dicDirectTranslationCS{tarantule, slíďák tarentský} \textit{(l.~{\textLA{Lycosa tarantula}})}
\dicEntry[tarfur] \dicTerm{tarf|ur} \dicIPA{{t\smash{\textsuperscript{h}}}{a}{r}{v}{\textscy}{\textsubring{r}}} \dicPos{m}[6] \dicFlx{(‑s, ‑ar)}[5] \dicSynonym{naut\smash{\textsuperscript{1}}} \dicDirectTranslationCS{býk}
\dicEntry[tarnar] \dicTerm{tarnar} \dicIPA{{t\smash{\textsuperscript{h}}}{a}{r}{\textsubring{d}}{n}{a}{\textsubring{r}}} \dicPos{f} \dicFlx{sg gen} \dicLink{törn}
\dicEntry[taska] \dicTerm{taska} \dicsymFrequent\  \dicIPA{{t\smash{\textsuperscript{h}}}{a}{s}{\r{g}}{a}} \dicPos{f}[1] \dicFlx{(tösku, töskur)}[8] \textbf{1.} \dicSynonym{ferðataska} \dicDirectTranslationCS{kufr} \dicExampleIS{í stórum töskum} \dicExampleCS{ve velkých kufrech}  \textbf{2.} \dicSynonym{handtaska} \dicDirectTranslationCS{taška, kabelka}  \textbf{3.} \dicSynonym{skjalataska} \dicDirectTranslationCS{kufřík}  \textbf{4.} \dicSynonym{skólataska} \dicDirectTranslationCS{(školní) baťoh\,/\addthin brašna}
\dicEntry[Tasmanía] \dicTerm{Tasmaní|a} \dicIPA{{t\smash{\textsuperscript{h}}}{a}{s}{m}{a}{n}{i}{j}{a}} \dicPos{f}[1] \dicFlx{(‑u)}[5] \dicFieldCat{geog.} \dicDirectTranslationCS{Tasmánie}
\dicEntry[tattóvera] \dicTerm{tattó··ver|a}\dicTerm{, tattúera} \dicIPA{{t\smash{\textsuperscript{h}}}\-{a}\-{h}\-{\textsubring{d}}\-{ou}\-{v}\-{\textepsilon}\-{r}\-{a}\-} \dicPos{v}[1] \dicFlx{(‑aði)}[1] \dicFlx{acc} \dicDirectTranslationCS{(vy)tetovat}
\dicEntry[tattóvering] \dicTerm{tattó··ver·ing} \dicIPA{{t\smash{\textsuperscript{h}}}{a}{h}{\textsubring{d}}{ou}{v}{\textepsilon}{r}{i}{\ng}{\r{g}}} \dicPos{f}[4] \dicFlx{(‑ar, ‑ar)}[5] \textbf{1.} \dicSynonym{húðflúr} \dicDirectTranslationCS{tetování} \dicIndirectTranslationCS{(obrazec vzniklý tetováním)}  \textbf{2.} \dicSynonym*{það að tattóvera} \dicDirectTranslationCS{tetování} \dicIndirectTranslationCS{(vpichování obrazců do kůže)}
\dicEntry[tattú] \dicTerm{tattú} \dicIPA{{t\smash{\textsuperscript{h}}}{a}{h}{\textsubring{d}}{u}} \dicPos{n}[2] \dicFlx{(‑s, ‑)}[5] \dicSynonym{húðflúr} \dicDirectTranslationCS{tetování} \dicIndirectTranslationCS{(obrazec vzniklý tetováním)}
\dicEntry[tattúera] \dicTerm{tattú··er|a} \dicIPA{{t\smash{\textsuperscript{h}}}{a}{h}{\textsubring{d}}{u}{\textepsilon}{r}{a}} \dicPos{v}[1] \dicFlx{(‑aði)}[1] \dicLink{tattóvera}
\dicEntry[tau] \dicTerm{tau} \dicIPA{{t\smash{\textsuperscript{h}}}{\oe i}{\textlengthmark}} \dicPos{n}[2] \dicFlx{(‑s, ‑)}[5] \textbf{1.} \dicDirectTranslationCS{(lněné) plátno} \dicExampleIS{Kjóllinn er úr taui.} \dicExampleCS{Šaty jsou z~plátna.}  \textbf{2.} \dicSynonym{þvottur} \dicDirectTranslationCS{prádlo (připravené na praní)}
\dicEntry[taug] \dicTerm{taug} \dicsymFrequent\  \dicIPA{{t\smash{\textsuperscript{h}}}{\oe i}{\textlengthmark}{x}} \dicPos{f}[4] \dicFlx{(‑ar, ‑ar)}[1] \textbf{1.} \dicFieldCat{anat.} \dicDirectTranslationCS{nerv}  \textbf{2.} \dicSynonym{strengur} \dicDirectTranslationCS{vlákno, lanko, lano};  \dicIdiom{taug}{ \dicPhraseIS{e‑að fer í taugarnar á e‑m}} \dicLangCat{přen.} \dicDirectTranslationCS{(co komu) leze na nervy};  \dicPhraseIS{fara á taugum} \dicLangCat{přen.} \dicDirectTranslationCS{mít nervy nadranc};  \dicPhraseIS{hafa taugar til e‑s} \dicLangCat{přen.} \dicDirectTranslationCS{mít silné pouto k~(čemu)};  \dicPhraseIS{hafa sterkar taugar} \dicLangCat{přen.} \dicDirectTranslationCS{mít silné nervy};  \dicPhraseIS{taka e‑n á taugum} \dicDirectTranslationCS{znervóznit (koho)}
\dicEntry[taugaáfall] \dicTerm{tauga··á·|fall} \dicIPA{{t\smash{\textsuperscript{h}}}{\oe i}{\textlengthmark}{\textbabygamma}{a}{au}{f}{a}{\textsubring{d}}{\textsubring{l}}} \dicPos{n}[2] \dicFlx{(‑falls, ‑föll)}[8] \dicFieldCat{med.} \dicDirectTranslationCS{nervové zhroucení}
\dicEntry[taugabilaður] \dicTerm{tauga··bil·|aður} \dicIPA{{t\smash{\textsuperscript{h}}}{\oe i}{\textlengthmark}{\textbabygamma}{a}{\textsubring{b}}{\textsci}{l}{a}{ð}{\textscy}{\textsubring{r}}} \dicPos{adj}[3] \dicFlx{(f ‑uð)}[3] \dicDirectTranslationCS{mající sla\-bé\,/\addthin pocuchané nervy, mající nervy nadranc}
\dicEntry[taugabilun] \dicTerm{tauga··bil|un} \dicIPA{{t\smash{\textsuperscript{h}}}{\oe i}{\textlengthmark}{\textbabygamma}{a}{\textsubring{b}}{\textsci}{l}{\textscy}{\textsubring{n}}} \dicPos{f}[7] \dicFlx{(‑unar, ‑anir)}[8] \dicDirectTranslationCS{nervozita, slabé nervy}
\dicEntry[taugadeild] \dicTerm{tauga··deild} \dicIPA{{t\smash{\textsuperscript{h}}}{\oe i}{\textlengthmark}{\textbabygamma}{a}{\textsubring{d}}{ei}{l}{\textsubring{d}}} \dicPos{f}[7] \dicFlx{(‑ar, ‑ir)}[1] \dicDirectTranslationCS{neurologické oddělení}
\dicEntry[taugafruma] \dicTerm{tauga··frum|a} \dicIPA{{t\smash{\textsuperscript{h}}}{\oe i}{\textlengthmark}{\textbabygamma}{a}{f}{r}{\textscy}{m}{a}} \dicPos{f}[1] \dicFlx{(‑u, ‑ur)}[19] \dicFieldCat{anat.} \dicDirectTranslationCS{neuron, nervová buňka}
\dicEntry[taugagripla] \dicTerm{tauga··gripl|a} \dicIPA{{t\smash{\textsuperscript{h}}}{\oe i}{\textlengthmark}{\textbabygamma}{a}{\r{g}}{r}{\textsci}{h}{\textsubring{b}}{l}{a}} \dicPos{f}[1] \dicFlx{(‑u, ‑ur)}[13] \dicFieldCat{anat.} \dicDirectTranslationCS{dendrit}
\dicEntry[taugakerfi] \dicTerm{tauga··kerfi} \dicIPA{{t\smash{\textsuperscript{h}}}{\oe i}{\textlengthmark}{\textbabygamma}{a}{c\smash{\textsuperscript{h}}}{\textepsilon}{r}{v}{\textsci}} \dicPos{n}[2] \dicFlx{(‑s, ‑)}[14] \dicFieldCat{anat.} \dicDirectTranslationCS{nervová soustava, nervový systém}
\dicEntry[taugalæknir] \dicTerm{tauga··lækn|ir} \dicIPA{{t\smash{\textsuperscript{h}}}{\oe i}{\textlengthmark}{\textbabygamma}{a}{l}{a}{i}{h}{\r{g}}{n}{\textsci}{\textsubring{r}}} \dicPos{m}[7] \dicFlx{(‑is, ‑ar)}[1] \dicDirectTranslationCS{neurolog, neuroložka}
\dicEntry[taugaóstyrkur] \dicTerm{tauga··ó·styrkur} \dicIPA{{t\smash{\textsuperscript{h}}}{\oe i}{\textlengthmark}{\textbabygamma}{a}{ou}{s}{\textsubring{d}}{\textsci}{\textsubring{r}}{\r{g}}{\textscy}{\textsubring{r}}} \dicPos{adj}[1]\dicFlx{}[-1] \dicDirectTranslationCS{nervózní, napjatý, podrážděný}
\dicEntry[taugasjúkdómur] \textls[10]{\dicTerm{tauga··sjúk·dóm|ur} \dicIPA{{t\smash{\textsuperscript{h}}}{\oe i}{\textlengthmark}{\textbabygamma}{a}{s}{j}{u}{\r{g}}{\textsubring{d}}{ou}{m}{\textscy}{\textsubring{r}}} \dicPos{m}[6] \dicFlx{(‑s, ‑ar)}[10] \dicFieldCat{med.} \dicDirectTranslationCS{neurologické onemocnění, nemoc nervového ústrojí}}
\dicEntry[taugasjúklingur] \dicTerm{tauga··sjúk·ling|ur} \dicIPA{{t\smash{\textsuperscript{h}}}{\oe i}{\textlengthmark}{\textbabygamma}{a}{s}{j}{u}{\r{g}}{l}{i}{\ng}{\r{g}}{\textscy}{\textsubring{r}}} \dicPos{m}[6] \dicFlx{(‑s, ‑ar)}[8] \dicDirectTranslationCS{pacient(ka) s~neurologickým onemocněním}
\dicEntry[taugaspenna] \dicTerm{tauga··spenn|a} \dicIPA{{t\smash{\textsuperscript{h}}}{\oe i}{\textlengthmark}{\textbabygamma}{a}{s}{\textsubring{b}}{\textepsilon}{n}{a}} \dicPos{f}[1] \dicFlx{(‑u)}[5] \dicDirectTranslationCS{nervozita, napětí}
\dicEntry[taugastríð] \dicTerm{tauga··stríð} \dicIPA{{t\smash{\textsuperscript{h}}}{\oe i}{\textlengthmark}{\textbabygamma}{a}{s}{\textsubring{d}}{r}{i}{\texttheta}} \dicPos{n}[2] \dicFlx{(‑s, ‑)}[5] \dicDirectTranslationCS{válka nervů, psychologická válka}
\dicEntry[taugaveiki] \dicTerm{tauga··veik|i} \dicIPA{{t\smash{\textsuperscript{h}}}{\oe i}{\textlengthmark}{\textbabygamma}{a}{v}{ei}{\r{\textObardotlessj}}{\textsci}} \dicPos{f}[3] \dicFlx{(‑i)}[3] \dicFieldCat{med.} \dicDirectTranslationCS{tyfus}
\dicEntry[taugaveiklaður] \dicTerm{tauga··veikl·|aður} \dicIPA{{t\smash{\textsuperscript{h}}}{\oe i}{\textlengthmark}{\textbabygamma}{a}{v}{ei}{h}{\r{g}}{l}{a}{ð}{\textscy}{\textsubring{r}}} \dicPos{adj}[3] \dicFlx{(f ‑uð)}[3] \textbf{1.} \dicFieldCat{med.} \dicDirectTranslationCS{neurózní}  \textbf{2.} \dicSynonym{taugaóstyrkur} \dicDirectTranslationCS{nervózní, napjatý} \dicExampleIS{Hann gaf frá sér taugaveiklaðan hlátur.} \dicExampleCS{Nervózně se zasmál.}
\dicEntry[taugaveiklun] \dicTerm{tauga··veikl|un} \dicIPA{{t\smash{\textsuperscript{h}}}{\oe i}{\textlengthmark}{\textbabygamma}{a}{v}{ei}{h}{\r{g}}{l}{\textscy}{\textsubring{n}}} \dicPos{f}[7] \dicFlx{(‑unar)}[9] \textbf{1.} \dicFieldCat{med.} \dicDirectTranslationCS{neuróza}  \textbf{2.} \dicSynonym{óróleiki} \dicDirectTranslationCS{nervozita, slabé nervy}
\dicEntry[taugungur] \dicTerm{taug··ung|ur} \dicIPA{{t\smash{\textsuperscript{h}}}{\oe i}{\textlengthmark}{\textbabygamma}{u}{\ng}{\r{g}}{\textscy}{\textsubring{r}}} \dicPos{m}[6] \dicFlx{(‑s, ‑ar)}[8] \dicFieldCat{anat.} \dicDirectTranslationCS{neuron, nervová buňka}
\dicEntry[taumlaus] \dicTerm{taum··laus} \dicIPA{{t\smash{\textsuperscript{h}}}{\oe i}{m}{l}{\oe i}{s}} \dicPos{adj}[5]\dicFlx{}[-1] \dicSynonym{stjórnlaus} \dicDirectTranslationCS{nekontrolovatelný, nespoutaný, divoký} \dicExampleIS{taumlaus ímyndun} \dicExampleCS{nekontrolovatelná představivost}
\dicEntry[taumur] \dicTerm{taum|ur} \dicsymFrequent\  \dicIPA{{t\smash{\textsuperscript{h}}}{\oe i}{\textlengthmark}{m}{\textscy}{\textsubring{r}}} \dicPos{m}[6] \dicFlx{(‑s, ‑ar)}[22] \textbf{1.} \dicSynonym*{beislistaumur} \dicDirectTranslationCS{otěž, oprať, uzda} \dicExampleIS{leðurtaumur} \dicExampleCS{kožená otěž}  \textbf{2.} \dicSynonym{taug} \dicDirectTranslationCS{šňůra, provaz}  \textbf{3.} \dicSynonym{rák} \dicDirectTranslationCS{proužek, pruh};  \dicIdiom{taumur}{ \dicPhraseIS{draga taum e‑rs}} \dicDirectTranslationCS{podpořit (koho), podporovat (koho)}; { \dicPhraseIS{gefa e‑u lausan tauminn}} \dicLangCat{přen.} \dicDirectTranslationCS{popustit (čemu) uzdu (fantazii ap.)};  \dicPhraseIS{grípa\,/\addthin taka í taumana} \dicDirectTranslationCS{zasáhnout, zakročit, vložit se (do sporu ap.)}
\dicEntry[taut] \dicTerm{taut} \dicIPA{{t\smash{\textsuperscript{h}}}{\oe i}{\textlengthmark}{\textsubring{d}}} \dicPos{n}[2] \dicFlx{(‑s)}[2] \dicSynonym{muldur} \dicDirectTranslationCS{mrmlání, brblání};  \dicPhraseIS{koma engu tauti við e‑n} \dicLangCat{přen.} \dicDirectTranslationCS{nemoct přivést (koho) k~rozumu}
\dicEntry[tauta] \dicTerm{taut|a} \dicsymFrequent\  \dicIPA{{t\smash{\textsuperscript{h}}}{\oe i}{\textlengthmark}{\textsubring{d}}{a}} \dicPos{v}[1] \dicFlx{(‑aði)}[1] \dicFlx{acc} \dicDirectTranslationCS{brblat, mručet};  \dicPhraseIS{tauta e‑ð fyrir munni sér} \dicDirectTranslationCS{brblat si (co) pro sebe};  \dicIdiom{tauta}{ \dicPhraseIS{hvað sem tautar og raular}} \dicDirectTranslationCS{ať se děje, co se děje}
\dicEntry[taxi] \dicTerm{tax|i} \dicIPA{{t\smash{\textsuperscript{h}}}{a}{x}{s}{\textsci}} \dicPos{m}[1] \dicFlx{(‑a, ‑ar)}[8] \dicLangCat{hovor.} \dicSynonym{leigubíll} \dicDirectTranslationCS{taxi, taxík} \dicExampleIS{Geturðu kallað á taxa fyrir mig?} \dicExampleCS{Můžeš mi zavolat taxík?}
\dicEntry[taxti] \dicTerm{taxt|i} \dicIPA{{t\smash{\textsuperscript{h}}}{a}{x}{s}{\textsubring{d}}{\textsci}} \dicPos{m}[1] \dicFlx{(‑a, ‑ar)}[8] \dicDirectTranslationCS{sazba, taxa, tarif} \dicExampleIS{Taxtinn er 10 krónur á klukkustund.} \dicExampleCS{Sazba je 10 korun za hodinu.}
\dicEntry[tá] \dicTerm{tá} \dicsymFrequent\  \dicIPA{{t\smash{\textsuperscript{h}}}{au}{\textlengthmark}} \dicPos{f}[9] \dicFlx{(‑ar, tær)}[1] \textbf{1.} \dicDirectTranslationCS{prst (na noze)} \dicExampleIS{Á hvorum fæti eru fimm tær.} \dicExampleCS{Na každé noze je pět prstů.}  \textbf{2.} \dicDirectTranslationCS{špička (boty ap.)};  \dicIdiom{tá}{ \dicPhraseIS{troða e‑m um tær}} \dicLangCat{přen.} \dicDirectTranslationCS{lepit\,/\addthin věšet se (komu) na paty}
\dicEntry[táfýla] \dicTerm{tá··fýl|a} \dicIPA{{t\smash{\textsuperscript{h}}}{au}{\textlengthmark}{f}{i}{l}{a}} \dicPos{f}[1] \dicFlx{(‑u)}[5] \dicDirectTranslationCS{zápach nohou, zápach z~chodidel}
\dicEntry[tág] \dicTerm{tág} \dicIPA{{t\smash{\textsuperscript{h}}}{au}{\textlengthmark}{x}} \dicPos{f}[4] \dicFlx{(‑ar, ‑ar)}[1] \dicDirectTranslationCS{(dlouhý) prut, proutí} \dicIndirectTranslationCS{(zvláště v~pl)}
\dicEntry[tágakarfa] \dicTerm{tága··|karfa} \dicIPA{{t\smash{\textsuperscript{h}}}{au}{\textlengthmark}{a}{k\smash{\textsuperscript{h}}}{a}{r}{v}{a}} \dicPos{f}[1] \dicFlx{(‑körfu, ‑körfur)}[20] \dicDirectTranslationCS{proutěný koš\,/\addthin košík}
\dicEntry[tágamura] \dicTerm{tága··mur|a} \dicIPA{{t\smash{\textsuperscript{h}}}{au}{\textlengthmark}{a}{m}{\textscy}{r}{a}} \dicPos{f}[1] \dicFlx{(‑u, ‑ur)}[7] \dicFieldCat{bot.} \dicDirectTranslationCS{mochna husí} \textit{(l.~{\textLA{Potentilla anserina}})}  \dicsymPhoto\ 
\dicFigure{120518.jpg}{Tágamura}{Tágamura - Horová Tamara, Biolib, Copyright/CC-BY-NC}
\dicEntry[tákn] \dicTerm{tákn} \dicsymFrequent\  \dicIPA{{t\smash{\textsuperscript{h}}}{au}{h}{\r{g}}{\textsubring{n}}} \dicPos{n}[2] \dicFlx{(‑s, ‑)}[5] \textbf{1.} \dicSynonym{merki} \dicDirectTranslationCS{symbol, znak} \dicExampleIS{Krossinn er heilagt tákn.} \dicExampleCS{Kříž je posvátný symbol.}  \textbf{2.} \dicSynonym{einkenni} \dicDirectTranslationCS{znamení, znak}
\dicEntry[tákna] \dicTerm{tákn|a} \dicsymFrequent\  \dicIPA{{t\smash{\textsuperscript{h}}}{au}{h}{\r{g}}{n}{a}} \dicPos{v}[1] \dicFlx{(‑aði)}[1] \dicFlx{acc} \dicSynonym{þýða} \dicDirectTranslationCS{znamenat,  představovat, symbolizovat, značit} \dicExampleIS{Í draumum táknar fugl á tré hamingju.} \dicExampleCS{Ve snu symbolizuje pták na stromě štěstí.}
\dicEntry[táknfræði] \dicTerm{tákn··fræð|i} \dicIPA{{t\smash{\textsuperscript{h}}}{au}{h}{\r{g}}{\textsubring{n}}{f}{r}{a}{i}{ð}{\textsci}} \dicPos{f}[3] \dicFlx{(‑i)}[3] \dicDirectTranslationCS{sémiotika}
\dicEntry[tákngervingur] \dicTerm{tákn··gerv·ing|ur} \dicIPA{{t\smash{\textsuperscript{h}}}{au}{h}{\r{g}}{\textsubring{n}}{\r{\textObardotlessj}}{\textepsilon}{r}{v}{i}{\ng}{\r{g}}{\textscy}{\textsubring{r}}} \dicPos{m}[6] \dicFlx{(‑s, ‑ar)}[8] \dicDirectTranslationCS{symbol, znak} \dicExampleIS{tákngervingur e‑s} \dicExampleCS{projev (čeho)}
\dicEntry[tákngildi] \dicTerm{tákn··gildi} \dicIPA{{t\smash{\textsuperscript{h}}}{au}{h}{\r{g}}{\textsubring{n}}{\r{\textObardotlessj}}{\textsci}{l}{\textsubring{d}}{\textsci}} \dicPos{n}[2] \dicFlx{(‑s, ‑)}[14] \dicDirectTranslationCS{symbolika, symbolismus}
\dicEntry[táknkerfi] \dicTerm{tákn··kerfi} \dicIPA{{t\smash{\textsuperscript{h}}}{au}{h}{\r{g}}{\textsubring{n}}{c\smash{\textsuperscript{h}}}{\textepsilon}{r}{v}{\textsci}} \dicPos{n}[2] \dicFlx{(‑s, ‑)}[14] \dicDirectTranslationCS{znakový systém}
\dicEntry[táknmál] \dicTerm{tákn··mál} \dicIPA{{t\smash{\textsuperscript{h}}}{au}{h}{\r{g}}{\textsubring{n}}{m}{au}{\textsubring{l}}} \dicPos{n}[2] \dicFlx{(‑s, ‑)}[5] \textbf{1.} \dicSynonym{merkjamál} \dicDirectTranslationCS{znaková řeč}  \textbf{2.} \dicSynonym*{safn tákna} \dicDirectTranslationCS{symbolika}
\dicEntry[táknmynd] \dicTerm{tákn··mynd} \dicIPA{{t\smash{\textsuperscript{h}}}{au}{h}{\r{g}}{\textsubring{n}}{m}{\textsci}{n}{\textsubring{d}}} \dicPos{f}[7] \dicFlx{(‑ar, ‑ir)}[1] \dicSynonym{ímynd} \dicDirectTranslationCS{symbol}
\dicEntry[táknrænn] \dicTerm{tákn··rænn} \dicIPA{{t\smash{\textsuperscript{h}}}{au}{h}{\r{g}}{\textsubring{n}}{r}{a}{i}{\textsubring{d}}{\textsubring{n}}} \dicPos{adj}[7]\dicFlx{}[-1] \dicDirectTranslationCS{symbolický} \dicExampleIS{táknræn athöfn} \dicExampleCS{symbolický obřad}
\dicEntry[táknsaga] \dicTerm{tákn··|saga} \dicIPA{{t\smash{\textsuperscript{h}}}{au}{h}{\r{g}}{\textsubring{n}}{s}{a}{\textbabygamma}{a}} \dicPos{f}[1] \dicFlx{(‑sögu, ‑sögur)}[14] \dicFieldCat{lit.} \dicSynonym{allegó\-ría} \dicDirectTranslationCS{alegorie}
\begin{xtolerant}{}{1pt}
\dicEntry[táknsæisstefna] \dicTerm{tákn·sæis··stefn|a}\addthinS\dicIPA{{t\smash{\textsuperscript{h}}}{au}{h}{\r{g}}{\textsubring{n}}{s}{a}{i}{j}{\textsci}{s}{\textsubring{d}}{\textepsilon}{\textsubring{b}}{n}{a}}\addthinS\dicPos{f}[1]\addthinS\dicFlx{(‑u)}[5]\addthinS\dicFieldCat{lit.}\addthinS\dicSynonym{symbólismi} \dicDirectTranslationCS{symbolismus}
\end{xtolerant}
\dicEntry[tál] \dicTerm{tál\smash{\textsuperscript{1}}}\dicTerm{, tál\smash{\textsuperscript{2}}} \dicIPA{{t\smash{\textsuperscript{h}}}\-{au}\-{\textlengthmark}\-{\textsubring{l}}\-} \dicPos{f}[4] \dicFlx{(‑ar, ‑ar)}[1] \dicSynonym{svik} \dicDirectTranslationCS{klam, klamání, podvod};  \dicPhraseIS{draga e‑n á tálar} \dicDirectTranslationCS{oklamat (koho), podvést (koho)}
\dicEntry[tál] \dicTerm{tál\smash{\textsuperscript{2}}} \dicIPA{{t\smash{\textsuperscript{h}}}{au}{\textlengthmark}{\textsubring{l}}} \dicPos{n}[2] \dicFlx{(‑s)}[2] \dicLink{tál\smash{\textsuperscript{1}}}
\dicEntry[táldraga] \dicTerm{tál··|draga} \dicIPA{{t\smash{\textsuperscript{h}}}{au}{l}{\textsubring{d}}{r}{a}{\textbabygamma}{a}} \dicPos{v}[6] \dicFlx{(‑dreg, ‑dró, ‑drógum, ‑drægi, ‑dreg\-ið)}[129] \dicFlx{acc} \dicSynonym{tæla} \dicDirectTranslationCS{obelstít, oklamat, podvést} \dicExampleIS{táldraga e‑n} \dicExampleCS{oklamat (koho)}
\dicEntry[táldreg] \dicTerm{tál··dreg} \dicIPA{{t\smash{\textsuperscript{h}}}{au}{l}{\textsubring{d}}{r}{\textepsilon}{x}} \dicPos{v} \dicFlx{ind praes sg 1 pers} \dicLink{táldraga}
\dicEntry[táldregið] \dicTerm{tál··dregið} \dicIPA{{t\smash{\textsuperscript{h}}}{au}{l}{\textsubring{d}}{r}{ei}{\textsci}{\texttheta}} \dicPos{v} \dicFlx{supin} \dicLink{táldraga}
\dicEntry[táldró] \dicTerm{tál··dró} \dicIPA{{t\smash{\textsuperscript{h}}}{au}{l}{\textsubring{d}}{r}{ou}} \dicPos{v} \dicFlx{ind pf sg 1 pers} \dicLink{táldraga}
\dicEntry[táldrógum] \dicTerm{tál··drógum} \dicIPA{{t\smash{\textsuperscript{h}}}{au}{l}{\textsubring{d}}{r}{ou}{\textscy}{\textsubring{m}}} \dicPos{v} \dicFlx{ind pf pl 1 pers} \dicLink{táldraga}
\dicEntry[táldrægi] \dicTerm{tál··drægi} \dicIPA{{t\smash{\textsuperscript{h}}}{au}{l}{\textsubring{d}}{r}{a}{i}{j}{\textsci}} \dicPos{v} \dicFlx{con pf sg 1 pers} \dicLink{táldraga}
\dicEntry[tálga] \dicTerm{tálg|a} \dicIPA{{t\smash{\textsuperscript{h}}}{au}{l}{\r{g}}{a}} \dicPos{v}[1] \dicFlx{(‑aði)}[1] \dicFlx{acc} \dicDirectTranslationCS{ořezat, ořezávat, vyřezat, vyřezávat} \dicExampleIS{tálga blýantinn til} \dicExampleCS{ořezat tužku}
\dicEntry[tálkn] \dicTerm{tálkn} \dicIPA{{t\smash{\textsuperscript{h}}}{au}{\textsubring{l}}{\r{g}}{\textsubring{n}}} \dicPos{n}[2] \dicFlx{(‑s, ‑)}[5] \dicDirectTranslationCS{žábry}
\dicEntry[tálma] \dicTerm{tálm|a} \dicIPA{{t\smash{\textsuperscript{h}}}{au}{l}{m}{a}} \dicPos{v}[1] \dicFlx{(‑aði)}[1] \dicFlx{acc\,/\addthin dat} \dicPhraseIS{tálma e‑ð} \dicDirectTranslationCS{(za)bránit (čemu), (za)blokovat (co), zatarasit\,/\addthin zatarasovat (co)} \dicExampleIS{tálma umferðina} \dicExampleCS{blokovat dopravní provoz};  \dicPhraseIS{tálma e‑u} \dicDirectTranslationCS{(za)bránit (čemu), zabraňovat (čemu), (za)blokovat (co)} \dicExampleIS{tálma eðlilegri þróun} \dicExampleCS{bránit přirozenému vývoji};  \dicPhraseIS{tálma e‑m e‑ð} \dicDirectTranslationCS{(za)bránit (komu) v~(čem), zamezit\,/\addthin zamezovat (komu co)} \dicExampleIS{tálma honum inngöngu} \dicExampleCS{bránit mu ve vstupu}
\dicEntry[tálmi] \dicTerm{tálm|i} \dicIPA{{t\smash{\textsuperscript{h}}}{au}{l}{m}{\textsci}} \dicPos{m}[1] \dicFlx{(‑a, ‑ar)}[1] \dicSynonym{hindrun} \dicDirectTranslationCS{překážka, zábrana, bariéra, zátaras}
\dicEntry[tálsýn] \dicTerm{tál··sýn} \dicIPA{{t\smash{\textsuperscript{h}}}{au}{l}{s}{i}{\textsubring{n}}} \dicPos{f}[7] \dicFlx{(‑ar, ‑ir)}[1] \dicDirectTranslationCS{iluze, falešná představa}
\dicEntry[tálvon] \dicTerm{tál··von} \dicIPA{{t\smash{\textsuperscript{h}}}{au}{l}{v}{\textopeno}{\textsubring{n}}} \dicPos{f}[7] \dicFlx{(‑ar, ‑ir)}[1] \dicDirectTranslationCS{iluze, falešná naděje}
\dicEntry[táningur] \dicTerm{tán··ing|ur} \dicIPA{{t\smash{\textsuperscript{h}}}{au}{\textlengthmark}{n}{i}{\ng}{\r{g}}{\textscy}{\textsubring{r}}} \dicPos{m}[6] \dicFlx{(‑s, ‑ar)}[8] \dicDirectTranslationCS{náctiletý, náctiletá, teenager(ka)}
\dicEntry[táp] \dicTerm{táp} \dicIPA{{t\smash{\textsuperscript{h}}}{au}{\textlengthmark}{\textsubring{b}}} \dicPos{n}[2] \dicFlx{(‑s)}[2] \dicSynonym{dugur} \dicDirectTranslationCS{verva, elán, vitalita} \dicExampleIS{táp og fjör} \dicExampleCS{elán a~nadšení}
\dicEntry[tápmikill] \dicTerm{táp··|mikill} \dicIPA{{t\smash{\textsuperscript{h}}}{au}{\textlengthmark}{\textsubring{b}}{m}{\textsci}{\r{\textObardotlessj}}{\textsci}{\textsubring{d}}{\textsubring{l}}} \dicPos{adj}[11] \dicFlx{(comp ‑meiri, sup ‑mestur)}[12] \dicDirectTranslationCS{vitální, plný energie\,/\addthin síly}
\dicEntry[tár] \dicTerm{tár} \dicsymFrequent\  \dicIPA{{t\smash{\textsuperscript{h}}}{au}{\textlengthmark}{\textsubring{r}}} \dicPos{n}[2] \dicFlx{(‑s, ‑)}[5] \textbf{1.} \dicDirectTranslationCS{slza} \dicExampleIS{fá tár í augun} \dicExampleCS{rozplakat se};  \dicPhraseIS{fella tár} \dicDirectTranslationCS{brečet, plakat, ronit slzy};  \dicPhraseIS{með tárin í augunum} \dicFlx{adv} \dicDirectTranslationCS{se slzami v~očích}  \textbf{2.} \dicDirectTranslationCS{slza, kapka} \dicIndirectTranslationCS{(malé množství tekutiny)};  \dicIdiom{tár}{ \dicPhraseIS{e‑að er þyngra en tárum taki}} \dicDirectTranslationCS{(co) se nedá vydržet (smutek ap.)};  \dicPhraseIS{gráta þurrum tárum} \dicLangCat{přen.} \dicDirectTranslationCS{ronit krokodýlí slzy}
\dicEntry[táradalur] \dicTerm{tára··dal|ur} \dicIPA{{t\smash{\textsuperscript{h}}}{au}{\textlengthmark}{r}{a}{\textsubring{d}}{a}{l}{\textscy}{\textsubring{r}}} \dicPos{m}[9] \dicFlx{(‑s, ‑ir)}[9] \dicLangCat{přen.} \dicDirectTranslationCS{slzavé údolí}
\dicEntry[táragas] \dicTerm{tára··gas} \dicIPA{{t\smash{\textsuperscript{h}}}{au}{\textlengthmark}{r}{a}{\r{g}}{a}{s}} \dicPos{n}[2] \dicFlx{(‑s)}[2] \dicDirectTranslationCS{slzný plyn}
\dicEntry[tárast] \dicTerm{tár|ast} \dicIPA{{t\smash{\textsuperscript{h}}}{au}{\textlengthmark}{r}{a}{s}{\textsubring{d}}} \dicPos{v}[1] \dicFlx{(‑aðist)}[91] \dicFlx{refl} \dicDirectTranslationCS{plakat, poplakávat, slzet, ronit slzy}
\dicEntry[tárvotur] \dicTerm{tár··votur} \dicIPA{{t\smash{\textsuperscript{h}}}{au}{r}{v}{\textopeno}{\textsubring{d}}{\textscy}{\textsubring{r}}} \dicPos{adj}[1]\dicFlx{}[-1] \dicDirectTranslationCS{uplakaný, ubrečený, uslzený}
\dicEntry[táta] \dicTerm{tát|a} \dicIPA{{t\smash{\textsuperscript{h}}}{au}{\textlengthmark}{\textsubring{d}}{a}} \dicPos{f}[1] \dicFlx{(‑u, ‑ur)}[7] \dicDirectTranslationCS{dívenka, děvčátko}
\dicEntry[Tbílísí] \dicTerm{Tbílísí}\dicTerm{, Tíblisi}\dicTerm{, Tíblísí} \dicIPA{{t\smash{\textsuperscript{h}}}\-{\textsubring{b}}\-{i}\-{\textlengthmark}\-{l}\-{i}\-{s}\-{i}\-} \dicPos{subs} \dicFlx{indecl} \dicFieldCat{geog.} \dicDirectTranslationCS{Tbilisi} \dicIndirectTranslationCS{(hlavní město Gruzie)}
\dicEntry[tbl.] \dicTerm{tbl.} \dicPos{zkr} \dicPhraseIS{tölublað} \dicDirectTranslationCS{číslo (časopisu ap.)}
\dicEntry[te] \dicTerm{te} \dicIPA{{t\smash{\textsuperscript{h}}}{\textepsilon}{\textlengthmark}} \dicPos{n}[2] \dicFlx{(‑s)}[2] \dicDirectTranslationCS{čaj} \dicExampleIS{grænt og svart te} \dicExampleCS{zelený a~černý čaj}
\dicEntry[tebolli] \dicTerm{te··boll|i} \dicIPA{{t\smash{\textsuperscript{h}}}{\textepsilon}{\textlengthmark}{\textsubring{b}}{\textopeno}{\textsubring{d}}{l}{\textsci}} \dicPos{m}[1] \dicFlx{(‑a, ‑ar)}[1] \textbf{1.} \dicSynonym*{bolli með te í} \dicDirectTranslationCS{hrnek\,/\addthin hrníček čaje}  \textbf{2.} \dicSynonym{bolli} \dicDirectTranslationCS{hrnek, hrníček (na čaj ap.)}
\dicEntry[tef] \dicTerm{tef} \dicIPA{{t\smash{\textsuperscript{h}}}{\textepsilon}{\textlengthmark}{f}} \dicPos{v} \dicFlx{ind praes sg 1 pers} \dicLink{tefja}
\dicEntry[tefði] \dicTerm{tefði} \dicIPA{{t\smash{\textsuperscript{h}}}{\textepsilon}{v}{ð}{\textsci}} \dicPos{v} \dicFlx{con pf sg 1 pers} \dicLink{tefja}
\dicEntry[tefja] \dicTerm{tefja} \dicIPA{{t\smash{\textsuperscript{h}}}{\textepsilon}{v}{j}{a}} \dicPos{v}[4] \dicFlx{(tef, tafði, töfðum, tefði, tafið)}[11] \dicFlx{acc} \dicSynonym{hindra} \dicDirectTranslationCS{zdržet, zdržovat, pozdržet} \dicExampleIS{tefja e‑n} \dicExampleCS{zdržovat (koho)};  \dicIdiom{tefjast}{ \dicPhraseIS{tefjast}} \dicFlx{refl} \dicDirectTranslationCS{(po)zdržet se, (po)zastavit se}
\dicEntry[tefla] \dicTerm{tefl|a} \dicsymFrequent\  \dicIPA{{t\smash{\textsuperscript{h}}}{\textepsilon}{\textsubring{b}}{l}{a}} \dicPos{v}[2] \dicFlx{(‑di, ‑t)}[141] \dicFlx{acc} \dicDirectTranslationCS{hrát šachy} \dicExampleIS{tefla skák við hann} \dicExampleCS{hrát s~ním šachy};  \dicIdiom{tefla}[á]{ \dicPhraseIS{tefla á tæpasta vað}} \dicLangCat{přen.} \dicDirectTranslationCS{hrát vabank}; { \dicPhraseIS{tefla á tvær hættur}} \dicLangCat{přen.} \dicDirectTranslationCS{hazardovat, hrát vabank};  \dicIdiom{tefla}[fram]{ \dicPhraseIS{tefla e‑u fram}} \dicDirectTranslationCS{spustit (co), zahájit (co)};  \dicIdiom{tefla}[í]{ \dicPhraseIS{tefla e‑u í tvísýnu}} \dicDirectTranslationCS{hazardovat s~(čím)};  \dicIdiom{tefla}{ \dicPhraseIS{það er um líf eða dauða að tefla}} \dicLangCat{přen.} \dicDirectTranslationCS{je to otázka života a~smrti}
\dicEntry[teflon] \dicTerm{teflon} \dicIPA{{t\smash{\textsuperscript{h}}}{\textepsilon}{f}{l}{\textopeno}{\textsubring{n}}} \dicPos{n}[2] \dicFlx{(‑s)}[33] \dicDirectTranslationCS{teflon}
\dicEntry[teg.] \dicTerm{teg.} \dicPos{zkr} \dicPhraseIS{tegund} \dicDirectTranslationCS{druh, typ, třída}
\dicEntry[tegra] \dicTerm{tegr|a} \dicIPA{{t\smash{\textsuperscript{h}}}{\textepsilon}{\textbabygamma}{r}{a}} \dicPos{v}[1] \dicFlx{(‑aði)}[1] \dicFieldCat{mat.} \dicSynonym{heilda} \dicDirectTranslationCS{integrovat}
\dicEntry[tegrun] \dicTerm{tegr|un} \dicIPA{{t\smash{\textsuperscript{h}}}{\textepsilon}{\textbabygamma}{r}{\textscy}{\textsubring{n}}} \dicPos{f}[7] \dicFlx{(‑unar)}[9] \dicFieldCat{mat.} \dicDirectTranslationCS{integrace}
\dicEntry[tegund] \dicTerm{tegund} \dicsymFrequent\  \dicIPA{{t\smash{\textsuperscript{h}}}{\textepsilon}{\textlengthmark}{\textbabygamma}{\textscy}{n}{\textsubring{d}}} \dicPos{f}[7] \dicFlx{(‑ar, ‑ir)}[1] \textbf{1.} \dicSynonym{flokkur} \dicDirectTranslationCS{druh, typ, skupina, kategorie}  \textbf{2.} \dicFieldCat{biol.} \dicDirectTranslationCS{druh, odrůda} \dicExampleIS{ýmsar tegundir af þangi} \dicExampleCS{různé druhy chaluh}
\dicEntry[tegur] \dicTerm{tegur} \dicIPA{{t\smash{\textsuperscript{h}}}{\textepsilon}{\textlengthmark}{\textbabygamma}{\textscy}{\textsubring{r}}} \dicPos{n}[2] \dicFlx{(‑s, ‑)}[25] \dicFieldCat{mat.} \dicSynonym{heildi} \dicDirectTranslationCS{integrál}
\dicEntry[tehetta] \dicTerm{te··hett|a} \dicIPA{{t\smash{\textsuperscript{h}}}{\textepsilon}{\textlengthmark}{h}{\textepsilon}{h}{\textsubring{d}}{a}} \dicPos{f}[1] \dicFlx{(‑u, ‑ur)}[7] \dicDirectTranslationCS{čajová panenka}
\dicEntry[tehús] \dicTerm{te··hús} \dicIPA{{t\smash{\textsuperscript{h}}}{\textepsilon}{\textlengthmark}{h}{u}{s}} \dicPos{n}[2] \dicFlx{(‑s, ‑)}[5] \dicDirectTranslationCS{čajovna} \dicExampleIS{kínverskt tehús} \dicExampleCS{čínská čajovna}
\dicEntry[teigur] \dicTerm{teig|ur} \dicIPA{{t\smash{\textsuperscript{h}}}{ei}{\textlengthmark}{\textbabygamma}{\textscy}{\textsubring{r}}} \dicPos{m}[6] \dicFlx{(‑s, ‑ir\,/\addthin ‑ar)}[73] \textbf{1.} \dicDirectTranslationCS{pruh\,/\addthin pás louky}  \textbf{2.} \dicDirectTranslationCS{odpaliště} \dicIndirectTranslationCS{(část golfového hřiště)}
\dicEntry[teikn] \dicTerm{teikn} \dicIPA{{t\smash{\textsuperscript{h}}}{ei}{h}{\r{g}}{\textsubring{n}}} \dicPos{n}[2] \dicFlx{(‑s, ‑)}[5] \textbf{1.} \dicSynonym{merki} \dicDirectTranslationCS{znamení, znak} \dicExampleIS{gefa e‑m teikn um e‑ð} \dicExampleCS{dát (komu) znamení o~(čem)}  \textbf{2.} \dicFieldCat{poč.} \dicSynonym{íkon\smash{\textsuperscript{2}}} \dicDirectTranslationCS{ikona}
\dicEntry[teikna] \textls[15]{\dicTerm{teikn|a} \dicsymFrequent\  \dicIPA{{t\smash{\textsuperscript{h}}}{ei}{h}{\r{g}}{n}{a}} \dicPos{v}[1] \dicFlx{(‑aði)}[1] \dicFlx{acc} \dicDirectTranslationCS{(na)kreslit, (na)rýsovat} \dicExampleIS{Hann teiknaði með fingrunum í moldina.} \dicExampleCS{Kreslil \mbox{prsty v~hlíně.}}}
\dicEntry[teiknari] \dicTerm{teikn··ar|i} \dicIPA{{t\smash{\textsuperscript{h}}}{ei}{h}{\r{g}}{n}{a}{r}{\textsci}} \dicPos{m}[1] \dicFlx{(‑a, ‑ar)}[13] \textbf{1.} \dicSynonym*{sá sem teiknar} \dicDirectTranslationCS{kreslíř(ka), kreslič(ka), ilustrátor(ka)}  \textbf{2.} \dicFieldCat{poč.} \dicDirectTranslationCS{souřadnicový zapisovač, plotter}
\dicEntry[teikniborð] \dicTerm{teikni··borð} \dicIPA{{t\smash{\textsuperscript{h}}}{ei}{h}{\r{g}}{n}{\textsci}{\textsubring{b}}{\textopeno}{r}{\texttheta}} \dicPos{n}[2] \dicFlx{(‑s, ‑)}[5] \dicDirectTranslationCS{kreslicí\,/\addthin rýsovací prkno}
\dicEntry[teiknibók] \dicTerm{teikni··|bók} \dicIPA{{t\smash{\textsuperscript{h}}}{ei}{h}{\r{g}}{n}{\textsci}{\textsubring{b}}{ou}{\r{g}}} \dicPos{f}[8] \dicFlx{(‑bókar, ‑bækur)}[5] \dicDirectTranslationCS{skicář, skicák, náčrtník}
\dicEntry[teiknibóla] \dicTerm{teikni··ból|a} \dicIPA{{t\smash{\textsuperscript{h}}}{ei}{h}{\r{g}}{n}{\textsci}{\textsubring{b}}{ou}{l}{a}} \dicPos{f}[1] \dicFlx{(‑u, ‑ur)}[19] \dicDirectTranslationCS{připínáček, napínáček, rýsováček}
\dicEntry[teiknibækur] \dicTerm{teikni··bækur} \dicIPA{{t\smash{\textsuperscript{h}}}{ei}{h}{\r{g}}{n}{\textsci}{\textsubring{b}}{a}{i}{\r{g}}{\textscy}{\textsubring{r}}} \dicPos{f} \dicFlx{pl nom} \dicLink{teiknibók}
\dicEntry[teikniforrit] \dicTerm{teikni··for·rit} \dicIPA{{t\smash{\textsuperscript{h}}}{ei}{h}{\r{g}}{n}{\textsci}{f}{\textopeno}{r}{\textlengthmark}{\textsci}{\textsubring{d}}} \dicPos{n}[2] \dicFlx{(‑s, ‑)}[5] \dicFieldCat{poč.} \dicDirectTranslationCS{kreslicí program\,/\addthin aplikace, editor vektorové grafiky}
\dicEntry[teiknikennsla] \dicTerm{teikni··kennsl|a} \dicIPA{{t\smash{\textsuperscript{h}}}{ei}{h}{\r{g}}{n}{\textsci}{c\smash{\textsuperscript{h}}}{\textepsilon}{n}{s}{\textsubring{d}}{l}{a}} \dicPos{f}[1] \dicFlx{(‑u)}[5] \dicFieldCat{škol.} \dicDirectTranslationCS{(výuka) kreslení\,/\addthin rýsování} \dicIndirectTranslationCS{(vyučovací předmět)}
\dicEntry[teiknimynd] \dicTerm{teikni··mynd} \dicIPA{{t\smash{\textsuperscript{h}}}{ei}{h}{\r{g}}{n}{\textsci}{m}{\textsci}{n}{\textsubring{d}}} \dicPos{f}[7] \dicFlx{(‑ar, ‑ir)}[1] \textbf{1.} \dicSynonym*{hreyfimynd} \dicDirectTranslationCS{kreslený\,/\addthin animovaný film}  \textbf{2.} \dicSynonym*{skopteikning} \dicDirectTranslationCS{satirická kresba, karikatura}
\dicEntry[teiknimyndablað] \dicTerm{teikni·mynda··|blað} \dicIPA{{t\smash{\textsuperscript{h}}}{ei}{h}{\r{g}}{n}{\textsci}{m}{\textsci}{n}{\textsubring{d}}{a}{\textsubring{b}}{l}{a}{\texttheta}} \dicPos{n}[2] \dicFlx{(‑blaðs, ‑blöð)}[8] \dicDirectTranslationCS{komiks, kniha komiksů, komiksový sešit}
\dicEntry[teikning] \dicTerm{teikn··ing} \dicsymFrequent\  \dicIPA{{t\smash{\textsuperscript{h}}}{ei}{h}{\r{g}}{n}{i}{\ng}{\r{g}}} \dicPos{f}[4] \dicFlx{(‑ar, ‑ar)}[5] \textbf{1.} \dicSynonym*{það að teikna} \dicDirectTranslationCS{kreslení, rýsování}  \textbf{2.} \dicSynonym*{teiknuð mynd} \dicDirectTranslationCS{kresba, nákres, náčrt} \dicExampleIS{teikningar barna} \dicExampleCS{kresby dětí}
\dicEntry[teiknistofa] \dicTerm{teikni··stof|a} \dicIPA{{t\smash{\textsuperscript{h}}}{ei}{h}{\r{g}}{n}{\textsci}{s}{\textsubring{d}}{\textopeno}{v}{a}} \dicPos{f}[1] \dicFlx{(‑u, ‑ur)}[7] \textbf{1.} \dicDirectTranslationCS{kreslírna, rýsovna} \dicIndirectTranslationCS{(místnost na kreslení\,/\addthin rýsování ve škole)}  \textbf{2.} \dicSynonym*{vinnustaður arkitekts} \dicDirectTranslationCS{grafické studio}
\dicEntry[teiknun] \dicTerm{teikn|un} \dicIPA{{t\smash{\textsuperscript{h}}}{ei}{h}{\r{g}}{n}{\textscy}{\textsubring{n}}} \dicPos{f}[7] \dicFlx{(‑unar)}[9] \dicDirectTranslationCS{kreslení, rýsování} \dicIndirectTranslationCS{(aktivita nebo školní předmět)}
\dicEntry[teinn] \dicTerm{tein|n} \dicIPA{{t\smash{\textsuperscript{h}}}{ei}{\textsubring{d}}{\textsubring{n}}} \dicPos{m}[6] \dicFlx{(‑s, ‑ar)}[42] \textbf{1.} \dicSynonym*{mjó stöng} \dicDirectTranslationCS{proutek, špejle}  \textbf{2.} \dicSynonym{rák} \dicDirectTranslationCS{proužek, pruh} \dicIndirectTranslationCS{(zvláště v~pl)}  \textbf{3.} \dicSynonym{hjólteinn} \dicDirectTranslationCS{výplet } \dicIndirectTranslationCS{(kola)}  \textbf{4.} \dicPhraseIS{teinar} \dicFlx{pl} \dicDirectTranslationCS{koleje, kolejnice} \dicExampleIS{járnbrautarteinar} \dicExampleCS{železniční koleje};  \dicPhraseIS{fara út af teinunum} \dicDirectTranslationCS{vykolejit}
\dicEntry[teinréttur] \dicTerm{tein··réttur} \dicIPA{{t\smash{\textsuperscript{h}}}{ei}{n}{r}{j}{\textepsilon}{h}{\textsubring{d}}{\textscy}{\textsubring{r}}} \dicPos{adj}[1]\dicFlx{}[-10] \dicDirectTranslationCS{vypnutý, rovný jako prkno}
\dicEntry[teinungur] \dicTerm{tein··ung|ur} \dicIPA{{t\smash{\textsuperscript{h}}}{ei}{\textlengthmark}{n}{u}{\ng}{\r{g}}{\textscy}{\textsubring{r}}} \dicPos{m}[6] \dicFlx{(‑s, ‑ar)}[8] \dicDirectTranslationCS{proutek, haluz, ratolest}
\dicEntry[teista] \dicTerm{teist|a} \dicIPA{{t\smash{\textsuperscript{h}}}{ei}{s}{\textsubring{d}}{a}} \dicPos{f}[1] \dicFlx{(‑u, ‑ur)}[7] \dicFieldCat{zool.} \dicDirectTranslationCS{alkoun obecný} \textit{(l.~{\textLA{Cepphus grylle}})}  \dicsymPhoto\ 
\dicFigure{ds_image_teista_0_2.jpg}{Teista}{Teista - SeanMack, GFDL}
\dicEntry[teitur] \dicTerm{teitur} \dicIPA{{t\smash{\textsuperscript{h}}}{ei}{\textlengthmark}{\textsubring{d}}{\textscy}{\textsubring{r}}} \dicPos{adj}[1]\dicFlx{}[-1] \dicSynonym{glaður} \dicDirectTranslationCS{veselý, radostný}
\dicEntry[tek] \dicTerm{tek} \dicIPA{{t\smash{\textsuperscript{h}}}{\textepsilon}{\textlengthmark}{\r{g}}} \dicPos{v} \dicFlx{ind praes sg 1 pers} \dicLink{taka\smash{\textsuperscript{2}}}
\dicEntry[teketill] \dicTerm{te··|ketill} \dicIPA{{t\smash{\textsuperscript{h}}}{\textepsilon}{\textlengthmark}{c\smash{\textsuperscript{h}}}{\textepsilon}{\textsubring{d}}{\textsci}{\textsubring{d}}{\textsubring{l}}} \dicPos{m}[6] \dicFlx{(‑ketils, ‑katlar)}[79] \dicDirectTranslationCS{čajová konvice}
\dicEntry[tekið] \dicTerm{tekið} \dicIPA{{t\smash{\textsuperscript{h}}}{\textepsilon}{\textlengthmark}{\r{\textObardotlessj}}{\textsci}{\texttheta}} \dicPos{v} \dicFlx{supin} \dicLink{taka\smash{\textsuperscript{2}}}
\dicEntry[tekinn] \dicTerm{tekinn} \dicIPA{{t\smash{\textsuperscript{h}}}{\textepsilon}{\textlengthmark}{\r{\textObardotlessj}}{\textsci}{\textsubring{n}}} \dicPos{adj}[6]\dicFlx{}[-6] \dicSynonym{fölur} \dicDirectTranslationCS{pobledlý, skleslý}
\dicEntry[tekíla] \dicTerm{tekíl|a} \dicIPA{{t\smash{\textsuperscript{h}}}{\textepsilon}{\textlengthmark}{\r{\textObardotlessj}}{i}{l}{a}} \dicPos{n}[4] \dicFlx{(‑a)}[1] \dicDirectTranslationCS{tequila}
\dicEntry[tekja] \dicTerm{tekj|a} \dicsymFrequent\  \dicIPA{{t\smash{\textsuperscript{h}}}{\textepsilon}{\textlengthmark}{\r{\textObardotlessj}}{a}} \dicPos{f}[1] \dicFlx{(‑u, ‑ur)}[17] \textbf{1.} \dicDirectTranslationCS{sběr, sbírka} \dicIndirectTranslationCS{(především ve složených slovech)} \dicExampleIS{eggjatekja} \dicExampleCS{sběr vajec}  \textbf{2.} \dicPhraseIS{tekjur} \dicFlx{pl} \dicLink{tekjur}
\dicEntry[tekjuafgangur] \dicTerm{tekju··af·gang|ur} \dicIPA{{t\smash{\textsuperscript{h}}}{\textepsilon}{\textlengthmark}{\r{\textObardotlessj}}{\textscy}{a}{v}{\r{g}}{au}{\ng}{\r{g}}{\textscy}{\textsubring{r}}} \dicPos{m}[6] \dicFlx{(‑s)}[7] \dicFieldCat{ekon.} \dicDirectTranslationCS{příjmový přebytek}
\dicEntry[tekjuhalli] \dicTerm{tekju··hall|i} \dicIPA{{t\smash{\textsuperscript{h}}}{\textepsilon}{\textlengthmark}{\r{\textObardotlessj}}{\textscy}{h}{a}{\textsubring{d}}{l}{\textsci}} \dicPos{m}[1] \dicFlx{(‑a)}[3] \dicFieldCat{ekon.} \dicDirectTranslationCS{příjmový deficit}
\dicEntry[tekjulaus] \dicTerm{tekju··laus} \dicIPA{{t\smash{\textsuperscript{h}}}{\textepsilon}{\textlengthmark}{\r{\textObardotlessj}}{\textscy}{l}{\oe i}{s}} \dicPos{adj}[5]\dicFlx{}[-1] \dicDirectTranslationCS{(jsoucí) bez příjmů}
\dicEntry[tekjuliður] \dicTerm{tekju··lið|ur} \dicIPA{{t\smash{\textsuperscript{h}}}{\textepsilon}{\textlengthmark}{\r{\textObardotlessj}}{\textscy}{l}{\textsci}{ð}{\textscy}{\textsubring{r}}} \dicPos{m}[10] \dicFlx{(‑ar\,/\addthin ‑s, ‑ir)}[33] \dicFieldCat{ekon.} \dicDirectTranslationCS{příjmová položka}
\dicEntry[tekjulind] \dicTerm{tekju··lind} \dicIPA{{t\smash{\textsuperscript{h}}}{\textepsilon}{\textlengthmark}{\r{\textObardotlessj}}{\textscy}{l}{\textsci}{n}{\textsubring{d}}} \dicPos{f}[7] \dicFlx{(‑ar, ‑ir)}[1] \dicFieldCat{ekon.} \dicDirectTranslationCS{zdroj příjmů}
\dicEntry[tekjumark] \dicTerm{tekju··|mark} \dicIPA{{t\smash{\textsuperscript{h}}}{\textepsilon}{\textlengthmark}{\r{\textObardotlessj}}{\textscy}{m}{a}{\textsubring{r}}{\r{g}}} \dicPos{n}[2] \dicFlx{(‑marks, ‑mörk)}[8] \dicFieldCat{ekon.} \dicDirectTranslationCS{příjmová hranice}
\dicEntry[tekjur] \dicTerm{tekjur} \dicIPA{{t\smash{\textsuperscript{h}}}{\textepsilon}{\textlengthmark}{\r{\textObardotlessj}}{\textscy}{\textsubring{r}}} \dicPos{f}[1] \dicFlx{pl}[18] \dicDirectTranslationCS{příjmy}
\dicEntry[tekjuskattur] \dicTerm{tekju··skatt|ur} \dicIPA{{t\smash{\textsuperscript{h}}}{\textepsilon}{\textlengthmark}{\r{\textObardotlessj}}{\textscy}{s}{\r{g}}{a}{h}{\textsubring{d}}{\textscy}{\textsubring{r}}} \dicPos{m}[6] \dicFlx{(‑s, ‑ar)}[5] \dicFieldCat{ekon.} \dicDirectTranslationCS{daň z~příjmů}
\dicEntry[tekjuskipting] \dicTerm{tekju··skipt·ing} \dicIPA{{t\smash{\textsuperscript{h}}}{\textepsilon}{\textlengthmark}{\r{\textObardotlessj}}{\textscy}{s}{\r{\textObardotlessj}}{\textsci}{f}{\textsubring{d}}{i}{\ng}{\r{g}}} \dicPos{f}[4] \dicFlx{(‑ar)}[7] \dicFieldCat{ekon.} \dicDirectTranslationCS{rozdělení příjmů}
\dicEntry[tekjustofn] \dicTerm{tekju··stofn} \dicIPA{{t\smash{\textsuperscript{h}}}{\textepsilon}{\textlengthmark}{\r{\textObardotlessj}}{\textscy}{s}{\textsubring{d}}{\textopeno}{\textsubring{b}}{\textsubring{n}}} \dicPos{m}[4] \dicFlx{(‑s, ‑ar)}[3] \dicFieldCat{ekon.} \dicSynonym{tekjulind} \dicDirectTranslationCS{zdroj příjmů}
\dicEntry[tekjutengdur] \dicTerm{tekju··tengdur} \dicIPA{{t\smash{\textsuperscript{h}}}{\textepsilon}{\textlengthmark}{\r{\textObardotlessj}}{\textscy}{t\smash{\textsuperscript{h}}}{\textepsilon}{\ng}{\textsubring{d}}{\textscy}{\textsubring{r}}} \dicPos{adj}[2]\dicFlx{}[-17] \dicDirectTranslationCS{závislý na výši příjmu}
\dicEntry[tekjutrygging] \dicTerm{tekju··trygg·ing} \dicIPA{{t\smash{\textsuperscript{h}}}{\textepsilon}{\textlengthmark}{\r{\textObardotlessj}}{\textscy}{t\smash{\textsuperscript{h}}}{r}{\textsci}{\r{\textObardotlessj}}{i}{\ng}{\r{g}}} \dicPos{f}[4] \dicFlx{(‑ar, ‑ar)}[5] \dicFieldCat{ekon.} \dicDirectTranslationCS{pojištění příjmů}
\dicEntry[tekjutryggingarauki] \dicTerm{tekju·tryggingar··auk|i} \dicIPA{{t\smash{\textsuperscript{h}}}\-{\textepsilon}\-{\textlengthmark}\-{\r{\textObardotlessj}}\-{\textscy}\-{t\smash{\textsuperscript{h}}}\-{r}\-{\textsci}\-{\r{\textObardotlessj}}\-{i}\-{\ng}\-{\r{g}}\-{a}\-{r}\-{\oe i}\-{\r{\textObardotlessj}}\-{\textsci}\-} \dicPos{m}[1] \dicFlx{(‑a, ‑ar)}[1] \dicFieldCat{ekon.} \dicDirectTranslationCS{připojištění příjmů}
\dicEntry[tekk] \dicTerm{tekk} \dicIPA{{t\smash{\textsuperscript{h}}}{\textepsilon}{h}{\r{g}}} \dicPos{n}[2] \dicFlx{(‑s)}[2] \dicDirectTranslationCS{teak, týk, teakové dřevo}
\dicEntry[teknetín] \dicTerm{teknetín} \dicIPA{{t\smash{\textsuperscript{h}}}{\textepsilon}{h}{\r{g}}{n}{\textepsilon}{\textsubring{d}}{i}{\textsubring{n}}} \dicPos{n}[2] \dicFlx{(‑s)}[2] \dicFieldCat{chem.} \dicDirectTranslationCS{technecium} \textit{(l.~{\textLA{Tc, Technetium}})}
\dicEntry[tel] \dicTerm{tel} \dicIPA{{t\smash{\textsuperscript{h}}}{\textepsilon}{\textlengthmark}{\textsubring{l}}} \dicPos{v} \dicFlx{ind praes sg 1 pers} \dicLink{telja}
\dicEntry[Tel Avív] \dicTerm{Tel Avív} \dicIPA{{t\smash{\textsuperscript{h}}}{\textepsilon}{l}{a}{v}{i}{v}} \dicPos{subs} \dicFlx{indecl} \dicFieldCat{geog.} \dicDirectTranslationCS{Tel Aviv} \dicIndirectTranslationCS{(město v~Izraeli)}
\dicEntry[telauf] \dicTerm{te··lauf} \dicIPA{{t\smash{\textsuperscript{h}}}{\textepsilon}{\textlengthmark}{l}{\oe i}{f}} \dicPos{n}[2] \dicFlx{(‑s, ‑)}[5] \dicDirectTranslationCS{čajový list}
\dicEntry[teldi] \dicTerm{teldi} \dicIPA{{t\smash{\textsuperscript{h}}}{\textepsilon}{l}{\textsubring{d}}{\textsci}} \dicPos{v} \dicFlx{con pf sg 1 pers} \dicLink{telja}
\dicEntry[telja] \dicTerm{telja} \dicsymFrequent\  \dicIPA{{t\smash{\textsuperscript{h}}}{\textepsilon}{l}{j}{a}} \dicPos{v}[4] \dicFlx{(tel, taldi, töldum, teldi, talið)}[3] \dicFlx{acc} \textbf{1.} \dicDirectTranslationCS{(s)počítat, vyčíslit} \dicExampleIS{Ég tel peningana.} \dicExampleCS{Počítám peníze.}  \textbf{2.} \dicSynonym{álíta} \dicDirectTranslationCS{myslet si, domnívat se, soudit, usuzovat} \dicExampleIS{Ég tel þetta gott.} \dicExampleCS{Myslím si, že je to dobré.};  \dicIdiom{telja}[af]{ \dicPhraseIS{telja e‑n af}} \dicDirectTranslationCS{pokládat (koho) za mrtvého};  \dicIdiom{telja}[á]{ \dicPhraseIS{telja e‑n á e‑ð}} \dicDirectTranslationCS{namluvit (koho) k~(čemu), přemluvit (koho) k~(čemu)} \dicExampleIS{Ég vona að ég telji barnið á heimaverkefnið.} \dicExampleCS{Doufám, že namluvím dítě k~domácímu úkolu.};  \dicIdiom{telja}[eftir]{ \dicPhraseIS{telja e‑ð ekki eftir e‑m}} \dicDirectTranslationCS{nezávidět (co komu) (těžkou práci ap.)};  \dicIdiom{telja}[fram]{ \dicPhraseIS{telja fram (til skatts)}} \dicDirectTranslationCS{připravit a~podat daňové přiznání};  \dicIdiom{telja}[með]{ \dicPhraseIS{telja e‑ð með}} \dicDirectTranslationCS{započítat\,/\addthin započítávat (co), zahrnout\,/\addthin zahrnovat (co)};  \dicIdiom{telja}[niður]{ \dicPhraseIS{telja niður}} \dicDirectTranslationCS{odpočítat, odpočítávat (čas ap.)};  \dicIdiom{telja}[ofan af]{ \dicPhraseIS{telja e‑n ofan af e‑u}} \dicDirectTranslationCS{rozmluvit\,/\addthin rozmlouvat (komu co)};  \dicIdiom{telja}[saman]{ \dicPhraseIS{telja saman e‑ð}} \dicDirectTranslationCS{sečíst (co), sčítat (co), spočítat (co) dohromady};  \dicIdiom{telja}[upp]{ \dicPhraseIS{telja e‑ð upp}} \dicDirectTranslationCS{vypočítat\,/\addthin vypočítávat (co), vyjmenovat\,/\addthin vyjmenovávat (co)};  \dicIdiom{teljast}{ \dicPhraseIS{teljast}} \dicFlx{refl} \dicDirectTranslationCS{platit, být považován} \dicExampleIS{Hann telst vera góður læknir.} \dicExampleCS{Platí za dobrého lékaře.};  \dicIdiom{teljast}[með]{ \dicPhraseIS{teljast með}} \dicFlx{refl} \dicDirectTranslationCS{být zahrnut, počítat se (k)};  \dicIdiom{teljast}[til]{ \dicPhraseIS{teljast til e‑s}} \dicFlx{refl} \dicDirectTranslationCS{náležet k~(čemu), patřit k~(čemu)}; { \dicPhraseIS{e‑m telst (svo) til \dots{}}} \dicFlx{refl impers} \dicDirectTranslationCS{(kdo) odhaduje, že\dots{}}
\dicEntry[teljandi] \dicTerm{telj··andi} \dicIPA{{t\smash{\textsuperscript{h}}}{\textepsilon}{l}{j}{a}{n}{\textsubring{d}}{\textsci}} \dicPos{adj}[13] \dicFlx{indecl}[1] \dicDirectTranslationCS{značný, podstatný} \dicIndirectTranslationCS{(nejčastěji v~záporu)}
\dicEntry[teljari] \dicTerm{telj··ar|i} \dicIPA{{t\smash{\textsuperscript{h}}}{\textepsilon}{l}{j}{a}{r}{\textsci}} \dicPos{m}[1] \dicFlx{(‑a, ‑ar)}[13] \textbf{1.} \dicSynonym*{áhald sem telur} \dicDirectTranslationCS{počítadlo, čítač}  \textbf{2.} \dicFieldCat{mat.} \dicDirectTranslationCS{čitatel}
\dicEntry[tellúr] \dicTerm{tellúr} \dicIPA{{t\smash{\textsuperscript{h}}}{\textepsilon}{l}{\textlengthmark}{u}{\textsubring{r}}} \dicPos{n}[2] \dicFlx{(‑s)}[31] \dicFieldCat{chem.} \dicDirectTranslationCS{telur} \textit{(l.~{\textLA{Te, Tellurium}})}
\dicEntry[telpa] \dicTerm{telp|a} \dicsymFrequent\  \dicIPA{{t\smash{\textsuperscript{h}}}{\textepsilon}{\textsubring{l}}{\textsubring{b}}{a}} \dicPos{f}[1] \dicFlx{(‑u, ‑ur)}[13] \dicSynonym{stúlkubarn} \dicDirectTranslationCS{holčička, děvčátko} \dicExampleIS{Hún reyndi að hugga telpuna.} \dicExampleCS{Pokusila se utěšit děvčátko.}
\dicEntry[tem] \dicTerm{tem} \dicIPA{{t\smash{\textsuperscript{h}}}{\textepsilon}{\textlengthmark}{\textsubring{m}}} \dicPos{v} \dicFlx{ind praes sg 1 pers} \dicLink{temja}
\dicEntry[tema] \dicTerm{tem|a} \dicIPA{{t\smash{\textsuperscript{h}}}{\textepsilon}{\textlengthmark}{m}{a}} \dicPos{n}[1] \dicFlx{(‑a, ‑u)}[1] \dicLink{þema}
\dicEntry[temdi] \dicTerm{temdi} \dicIPA{{t\smash{\textsuperscript{h}}}{\textepsilon}{m}{\textsubring{d}}{\textsci}} \dicPos{v} \dicFlx{con pf sg 1 pers} \dicLink{temja}
\dicEntry[temja] \dicTerm{temja} \dicIPA{{t\smash{\textsuperscript{h}}}{\textepsilon}{m}{j}{a}} \dicPos{v}[4] \dicFlx{(tem, tamdi, tömdum, temdi, tamið)}[3] \dicFlx{acc} \textbf{1.} \dicDirectTranslationCS{ochočit, ochočovat, (z)krotit} \dicExampleIS{temja hest} \dicExampleCS{ochočit koně}  \textbf{2.} \dicPhraseIS{temja sér e‑ð} \dicDirectTranslationCS{osvojit si (co), přijmout za své (co)};  \dicProverb\  \dicPhraseIS{Hvað ungur nemur gamall temur.} \dicLangCat{přís.} \dicDirectTranslationCS{Co se v~mládí naučíš, ve stáří jako když najdeš.}
\dicEntry[templari] \dicTerm{templ··ar|i} \dicIPA{{t\smash{\textsuperscript{h}}}{\textepsilon}{\textsubring{m}}{\textsubring{b}}{l}{a}{r}{\textsci}} \dicPos{m}[1] \dicFlx{(‑a, ‑ar)}[13] \dicFieldCat{hist.} \dicDirectTranslationCS{templář}
\dicEntry[tempó] \dicTerm{tempó} \dicIPA{{t\smash{\textsuperscript{h}}}{\textepsilon}{\textsubring{m}}{\textsubring{b}}{ou}} \dicPos{n}[2] \dicFlx{(‑s, ‑)}[5] \dicLangCat{hovor.} \dicDirectTranslationCS{tempo}
\dicEntry[tempra] \dicTerm{tempr|a} \dicIPA{{t\smash{\textsuperscript{h}}}{\textepsilon}{\textsubring{m}}{\textsubring{b}}{r}{a}} \dicPos{v}[1] \dicFlx{(‑aði)}[1] \dicFlx{acc} \dicSynonym{stilla} \dicDirectTranslationCS{zmírnit, regulovat (teplotu ap.)} \dicExampleIS{tempra hita í herbergi} \dicExampleCS{zmírnit teplotu v~pokoji}
\dicEntry[tempraður] \dicTerm{tempr··|aður} \dicIPA{{t\smash{\textsuperscript{h}}}{\textepsilon}{\textsubring{m}}{\textsubring{b}}{r}{a}{ð}{\textscy}{\textsubring{r}}} \dicPos{adj}[3] \dicFlx{(f ‑uð)}[3] \dicDirectTranslationCS{mírný} \dicIndirectTranslationCS{(o~podnebí)};  \dicPhraseIS{tempraða beltið} \dicFieldCat{geog.} \dicDirectTranslationCS{mírný (podnebný) pás}
\dicEntry[temprun] \dicTerm{tempr|un} \dicIPA{{t\smash{\textsuperscript{h}}}{\textepsilon}{\textsubring{m}}{\textsubring{b}}{r}{\textscy}{\textsubring{n}}} \dicPos{f}[7] \dicFlx{(‑unar)}[9] \dicDirectTranslationCS{zmírnění, regulování (teploty ap.)}
\dicEntry[Tempsá] \dicTerm{Temps··á} \dicIPA{{t\smash{\textsuperscript{h}}}{\textepsilon}{\textsubring{m}}{\textsubring{b}}{s}{au}} \dicPos{f}[4] \dicFlx{(‑r)}[20] \dicFieldCat{geog.} \dicDirectTranslationCS{Temže} \dicIndirectTranslationCS{(řeka v~Anglii)}
\dicEntry[tendra] \dicTerm{tendr|a} \dicIPA{{t\smash{\textsuperscript{h}}}{\textepsilon}{n}{\textsubring{d}}{r}{a}} \dicPos{v}[1] \dicFlx{(‑aði)}[1] \dicFlx{acc} \textbf{1.} \dicSynonym{kveikja\smash{\textsuperscript{2}}} \dicDirectTranslationCS{zapálit, zažehnout} \dicExampleIS{tendra ljós} \dicExampleCS{zapálit světlo}  \textbf{2.} \dicSynonym{örva} \dicDirectTranslationCS{podnítit, vznítit, zapálit (zájem ap.)} \dicExampleIS{tendra áhuga manna} \dicExampleCS{podnítit zájem lidí}
\dicEntry[tengdadóttir] \dicTerm{tengda··|dóttir} \dicIPA{{t\smash{\textsuperscript{h}}}{\textepsilon}{\ng}{\textsubring{d}}{a}{\textsubring{d}}{ou}{h}{\textsubring{d}}{\textsci}{\textsubring{r}}} \dicPos{f}[11] \dicFlx{(‑dóttur, ‑dætur)}[2] \dicDirectTranslationCS{snacha} \dicIndirectTranslationCS{(manželka syna)}
\dicEntry[tengdafaðir] \dicTerm{tengda··|faðir} \dicIPA{{t\smash{\textsuperscript{h}}}{\textepsilon}{\ng}{\textsubring{d}}{a}{f}{a}{ð}{\textsci}{\textsubring{r}}} \dicPos{m}[13] \dicFlx{(‑föður, ‑feður)}[3] \dicDirectTranslationCS{tchán} \dicIndirectTranslationCS{(otec manželky nebo manžela)}
\dicEntry[tengdamamma] \dicTerm{tengda··|mamma} \dicIPA{{t\smash{\textsuperscript{h}}}{\textepsilon}{\ng}{\textsubring{d}}{a}{m}{a}{m}{a}} \dicPos{f}[1] \dicFlx{(‑mömmu, ‑mömmur)}[8] \dicSynonym{tengdamóðir} \dicDirectTranslationCS{tchyně} \dicIndirectTranslationCS{(matka manželky nebo manžela)}
\begin{xtolerant}{}{1pt}
\dicEntry[tengdamóðir] \dicTerm{tengda··|móðir} \dicIPA{{t\smash{\textsuperscript{h}}}{\textepsilon}{\ng}{\textsubring{d}}{a}{m}{ou}{ð}{\textsci}{\textsubring{r}}} \dicPos{f}[11] \dicFlx{(‑móður, ‑mæður)}[1] \dicDirectTranslationCS{tchy\-ně} \dicIndirectTranslationCS{(matka manželky nebo manžela)}
\end{xtolerant}
\dicEntry[tengdapabbi] \dicTerm{tengda··pabb|i} \dicIPA{{t\smash{\textsuperscript{h}}}{\textepsilon}{\ng}{\textsubring{d}}{a}{p\smash{\textsuperscript{h}}}{a}{\textsubring{b}}{\textsci}} \dicPos{m}[1] \dicFlx{(‑a, ‑ar)}[8] \dicSynonym{tengdafaðir} \dicDirectTranslationCS{tchán} \dicIndirectTranslationCS{(otec manželky nebo manžela)}
\dicEntry[tengdasonur] \dicTerm{tengda··|sonur} \dicIPA{{t\smash{\textsuperscript{h}}}{\textepsilon}{\ng}{\textsubring{d}}{a}{s}{\textopeno}{n}{\textscy}{\textsubring{r}}} \dicPos{m}[12] \dicFlx{(‑sonar, ‑synir)}[9] \dicDirectTranslationCS{zeť} \dicIndirectTranslationCS{(manžel dcery)}
\dicEntry[tengdir] \dicTerm{tengdir} \dicIPA{{t\smash{\textsuperscript{h}}}{\textepsilon}{\ng}{\textsubring{d}}{\textsci}{\textsubring{r}}} \dicPos{f}[12] \dicFlx{pl}[10] \dicSynonym{sifjar} \dicDirectTranslationCS{příbuzenské vztahy} \dicIndirectTranslationCS{(vzniklé manželstvím)}
\dicEntry[tengdur] \dicTerm{tengdur} \dicIPA{{t\smash{\textsuperscript{h}}}{\textepsilon}{\ng}{\textsubring{d}}{\textscy}{\textsubring{r}}} \dicPos{adj}[2]\dicFlx{}[-14] \textbf{1.} \dicDirectTranslationCS{svázaný, související}  \textbf{2.} \dicPhraseIS{vera tengdur e‑m} \dicDirectTranslationCS{být s~(kým) příbuzný\,/\addthin spřízněný (svatbou)}
\dicEntry[tengi] \dicTerm{tengi} \dicIPA{{t\smash{\textsuperscript{h}}}{ei}{\textltailn}{\r{\textObardotlessj}}{\textsci}} \dicPos{n}[2] \dicFlx{(‑s, ‑)}[16] \textbf{1.} \dicDirectTranslationCS{spojka, přípojka, konektor}  \textbf{2.} \dicFieldCat{chem.} \dicSynonym*{efnatengi} \dicDirectTranslationCS{vazba}
\dicEntry[tengibraut] \dicTerm{tengi··braut} \dicIPA{{t\smash{\textsuperscript{h}}}{ei}{\textltailn}{\r{\textObardotlessj}}{\textsci}{\textsubring{b}}{r}{\oe i}{\textsubring{d}}} \dicPos{f}[7] \dicFlx{(‑ar, ‑ir)}[1] \dicFieldCat{poč.} \dicDirectTranslationCS{sběrnice}
\dicEntry[tengikló] \dicTerm{tengi··|kló} \dicIPA{{t\smash{\textsuperscript{h}}}{ei}{\textltailn}{\r{\textObardotlessj}}{\textsci}{k\smash{\textsuperscript{h}}}{l}{ou}} \dicPos{f}[9] \dicFlx{(‑klóar, ‑klær)}[1] \dicFieldCat{elek.} \dicDirectTranslationCS{zástrčka}
\dicEntry[tengiliður] \dicTerm{tengi··lið|ur} \dicIPA{{t\smash{\textsuperscript{h}}}{ei}{\textltailn}{\r{\textObardotlessj}}{\textsci}{l}{\textsci}{ð}{\textscy}{\textsubring{r}}} \dicPos{m}[10] \dicFlx{(‑ar\,/\addthin ‑s, ‑ir)}[33] \textbf{1.} \dicSynonym{milliliður} \dicDirectTranslationCS{spojovací článek, pojítko}  \textbf{2.} \dicFieldCat{sport.} \dicDirectTranslationCS{spojka, středový hráč, středová hráčka}  \textbf{3.} \dicDirectTranslationCS{koordinátor(ka), kontaktní osoba}
\dicEntry[tengill] \dicTerm{teng|ill} \dicIPA{{t\smash{\textsuperscript{h}}}{ei}{\textltailn}{\r{\textObardotlessj}}{\textsci}{\textsubring{d}}{\textsubring{l}}} \dicPos{m}[6] \dicFlx{(‑ils, ‑lar)}[35] \textbf{1.} \dicSynonym{tengiliður} \dicDirectTranslationCS{kontakt, kontaktní osoba}  \textbf{2.} \dicFieldCat{poč.} \dicSynonym{krækja\smash{\textsuperscript{1}}} \dicDirectTranslationCS{(internetový) odkaz, link}  \textbf{3.} \dicFieldCat{elek.} \dicSynonym*{rafmagnstengi} \dicDirectTranslationCS{zásuvka}
\dicEntry[tenging] \dicTerm{teng··ing} \dicIPA{{t\smash{\textsuperscript{h}}}{ei}{\textltailn}{\r{\textObardotlessj}}{i}{\ng}{\r{g}}} \dicPos{f}[4] \dicFlx{(‑ar, ‑ar)}[5] \dicSynonym{samband} \dicDirectTranslationCS{spojení, připojení, zapojení}
\dicEntry[tengiorð] \dicTerm{tengi··orð} \dicIPA{{t\smash{\textsuperscript{h}}}{ei}{\textltailn}{\r{\textObardotlessj}}{\textsci}{\textopeno}{r}{\texttheta}} \dicPos{n}[2] \dicFlx{(‑s, ‑)}[5] \dicFieldCat{jaz.} \dicDirectTranslationCS{spojovací výraz}
\dicEntry[tengivagn] \dicTerm{tengi··vagn} \dicIPA{{t\smash{\textsuperscript{h}}}{ei}{\textltailn}{\r{\textObardotlessj}}{\textsci}{v}{a}{\r{g}}{\textsubring{n}}} \dicPos{m}[4] \dicFlx{(‑s, ‑ar)}[4] \dicDirectTranslationCS{přívěs, návěs}
\dicEntry[tengja] \dicTerm{teng|ja} \dicsymFrequent\  \dicIPA{{t\smash{\textsuperscript{h}}}{ei}{\textltailn}{\r{\textObardotlessj}}{a}} \dicPos{v}[2] \dicFlx{(‑di, ‑t)}[129] \dicFlx{acc} \dicDirectTranslationCS{spojit, spojovat, připojit, připojovat, napojit, napojovat};  \dicPhraseIS{tengja e‑ð við e‑ð} \dicDirectTranslationCS{připojit (co) k~(čemu), propojit (co) s~(čím)} \dicExampleIS{tengja iPodinn við tölvuna} \dicExampleCS{připojit iPod k~počítači};  \dicIdiom{tengjast}{ \dicPhraseIS{tengjast e‑u}} \dicFlx{refl} {\textbf{a.}} \dicDirectTranslationCS{spojit se s~(čím), připojit se k~(čemu)};  {\textbf{b.}} \dicDirectTranslationCS{souviset s~(čím), mít spojitost s~(čím), týkat se (čeho)}
\dicEntry[tengsl] \dicTerm{tengsl} \dicsymFrequent\  \dicIPA{{t\smash{\textsuperscript{h}}}{\textepsilon}{\ng}{s}{\textsubring{d}}{\textsubring{l}}} \dicPos{n}[2] \dicFlx{pl}[1] \textbf{1.} \dicDirectTranslationCS{spojení, kontakt} \dicExampleIS{tengsl við ættingja} \dicExampleCS{kontakt s~příbuznými}  \textbf{2.} \dicDirectTranslationCS{spojitost, souvislost}
\dicEntry[tengsli] \dicTerm{tengsli} \dicIPA{{t\smash{\textsuperscript{h}}}{\textepsilon}{\ng}{s}{\textsubring{d}}{l}{\textsci}} \dicPos{n}[2] \dicFlx{(‑s, ‑)}[14] \textbf{1.} \dicSynonym{tengsl} \dicDirectTranslationCS{spojitost, propojení}  \textbf{2.} \dicFieldCat{techn.} \dicSynonym{kúpling} \dicDirectTranslationCS{spojka}
\dicEntry[tengur] \dicTerm{tengur} \dicIPA{{t\smash{\textsuperscript{h}}}{ei}{\ng}{\r{g}}{\textscy}{\textsubring{r}}} \dicPos{f} \dicFlx{pl nom} \dicLink{töng}
\dicEntry[teningur] \dicTerm{ten··ing|ur} \dicIPA{{t\smash{\textsuperscript{h}}}{\textepsilon}{\textlengthmark}{n}{i}{\ng}{\r{g}}{\textscy}{\textsubring{r}}} \dicPos{m}[6] \dicFlx{(‑s, ‑ar)}[8] \textbf{1.} \dicSynonym{kubbur} \dicDirectTranslationCS{kostka, krychle}  \textbf{2.} \dicSynonym*{spilateningur} \dicDirectTranslationCS{kostka} \dicIndirectTranslationCS{(k~hraní označená čísly od 1 do 6)};  \dicIdiom{teningur}{ \dicPhraseIS{teningunum er kastað}} \dicLangCat{přen.} \dicDirectTranslationCS{kostky jsou vrženy}
\dicEntry[tennis] \dicTerm{tennis} \dicIPA{{t\smash{\textsuperscript{h}}}{\textepsilon}{n}{\textlengthmark}{\textsci}{s}} \dicPos{m}[4] \dicFlx{(‑\,/\addthin ‑s)}[38] \dicFieldCat{sport.} \dicDirectTranslationCS{tenis}
\dicEntry[tennisbolti] \dicTerm{tennis··bolt|i} \dicIPA{{t\smash{\textsuperscript{h}}}{\textepsilon}{n}{\textlengthmark}{\textsci}{s}{\textsubring{b}}{\textopeno}{\textsubring{l}}{\textsubring{d}}{\textsci}} \dicPos{m}[1] \dicFlx{(‑a, ‑ar)}[1] \dicDirectTranslationCS{tenisový míček, tenisák}
\dicEntry[tennisspaði] \dicTerm{tennis··spað|i} \dicIPA{{t\smash{\textsuperscript{h}}}{\textepsilon}{n}{\textlengthmark}{\textsci}{s}{\textsubring{b}}{a}{ð}{\textsci}} \dicPos{m}[1] \dicFlx{(‑a, ‑ar)}[8] \dicDirectTranslationCS{tenisová raketa}
\dicEntry[tennisvöllur] \dicTerm{tennis··|völlur} \dicIPA{{t\smash{\textsuperscript{h}}}{\textepsilon}{n}{\textlengthmark}{\textsci}{s}{v}{\oe}{\textsubring{d}}{l}{\textscy}{\textsubring{r}}} \dicPos{m}[11] \dicFlx{(‑vallar, ‑vellir)}[5] \dicDirectTranslationCS{tenisový kurt\,/\addthin dvorec}
\dicEntry[tenntur] \dicTerm{tenntur} \dicIPA{{t\smash{\textsuperscript{h}}}{\textepsilon}{\textsubring{n}}{\textsubring{d}}{\textscy}{\textsubring{r}}} \dicPos{adj}[1]\dicFlx{}[-10] \textbf{1.} \dicDirectTranslationCS{zubatý (nůž ap.)}  \textbf{2.} \dicDirectTranslationCS{ozubený, zoubkovaný (list ap.)}  \textbf{3.} \dicPhraseIS{vera vel\,/\addthin illa tenntur} \dicDirectTranslationCS{mít silné\,/\addthin slabé zuby}
\dicEntry[tennur] \dicTerm{tennur} \dicIPA{{t\smash{\textsuperscript{h}}}{\textepsilon}{n}{\textlengthmark}{\textscy}{\textsubring{r}}} \dicPos{f} \dicFlx{pl nom} \dicLink{tönn}
\dicEntry[tenór] \dicTerm{tenór} \dicIPA{{t\smash{\textsuperscript{h}}}{\textepsilon}{\textlengthmark}{n}{ou}{\textsubring{r}}} \dicPos{m}[4] \dicFlx{(‑s, ‑ar)}[14] \textbf{1.} \dicFieldCat{hud.} \dicDirectTranslationCS{tenor} \dicIndirectTranslationCS{(vysoký mužský zpěvný hlas)}  \textbf{2.} \dicFieldCat{hud.} \dicDirectTranslationCS{tenor} \dicIndirectTranslationCS{(pěvec zpívající tenorem)}
\dicEntry[teoretískur] \dicTerm{teoretískur} \dicIPA{{t\smash{\textsuperscript{h}}}{\textepsilon}{\textlengthmark}{\textopeno}{r}{\textepsilon}{\textsubring{d}}{i}{s}{\r{g}}{\textscy}{\textsubring{r}}} \dicPos{adj}[1]\dicFlx{}[-1] \dicDirectTranslationCS{teoretický}
\dicEntry[teoría] \dicTerm{teorí|a} \dicIPA{{t\smash{\textsuperscript{h}}}{\textepsilon}{\textlengthmark}{\textopeno}{r}{i}{j}{a}} \dicPos{f}[1] \dicFlx{(‑u, ‑ur)}[7] \dicSynonym{kenning} \dicDirectTranslationCS{teorie}
\dicEntry[tepoki] \dicTerm{te··pok|i} \dicIPA{{t\smash{\textsuperscript{h}}}{\textepsilon}{\textlengthmark}{p\smash{\textsuperscript{h}}}{\textopeno}{\r{\textObardotlessj}}{\textsci}} \dicPos{m}[1] \dicFlx{(‑a, ‑ar)}[1] \dicDirectTranslationCS{čajový sáček\,/\addthin pytlík}
\dicEntry[tepottur] \dicTerm{te··pott|ur} \dicIPA{{t\smash{\textsuperscript{h}}}{\textepsilon}{\textlengthmark}{p\smash{\textsuperscript{h}}}{\textopeno}{h}{\textsubring{d}}{\textscy}{\textsubring{r}}} \dicPos{m}[6] \dicFlx{(‑s, ‑ar)}[4] \dicSynonym{teketill} \dicDirectTranslationCS{čajová konvice\,/\addthin konývka}
\dicEntry[teppa] \dicTerm{tepp|a\smash{\textsuperscript{1}}} \dicIPA{{t\smash{\textsuperscript{h}}}{\textepsilon}{h}{\textsubring{b}}{a}} \dicPos{f}[1] \dicFlx{(‑u, ‑ur)}[19] \dicDirectTranslationCS{zácpa, zablokování, zatarasení}
\dicEntry[teppa] \dicTerm{tepp|a\smash{\textsuperscript{2}}} \dicIPA{{t\smash{\textsuperscript{h}}}{\textepsilon}{h}{\textsubring{b}}{a}} \dicPos{v}[2] \dicFlx{(‑ti, ‑t)}[63] \dicFlx{acc} \dicSynonym{stífla\smash{\textsuperscript{2}}} \dicDirectTranslationCS{(za)blokovat, zatarasit, ucpat, ucpávat} \dicExampleIS{teppa umferð} \dicExampleCS{zablokovat dopravní ruch};  \dicIdiom{teppast}{ \dicPhraseIS{teppast}} \dicFlx{refl} \dicDirectTranslationCS{(za)blokovat se, zatarasit se, ucpat se} \dicExampleIS{Hæstu fjallvegir teppast strax í fyrstu snjóum.} \dicExampleCS{Nejvyšší horské cesty se zablokují hned s~prvním sněhem.}
\dicEntry[teppi] \dicTerm{teppi} \dicsymFrequent\  \dicIPA{{t\smash{\textsuperscript{h}}}{\textepsilon}{h}{\textsubring{b}}{\textsci}} \dicPos{n}[2] \dicFlx{(‑s, ‑)}[14] \textbf{1.} \dicDirectTranslationCS{koberec} \dicExampleIS{stofa með teppum á gólfi} \dicExampleCS{obývací pokoj s~koberci na podlaze}  \textbf{2.} \dicSynonym{ábreiða} \dicDirectTranslationCS{přehoz, potah};  \dicIdiom{teppi}{ \dicPhraseIS{taka e‑n á teppið}} \dicLangCat{přen.} \dicDirectTranslationCS{zavolat si (koho) na kobereček}
\dicEntry[tepra] \dicTerm{tepr|a} \dicIPA{{t\smash{\textsuperscript{h}}}{\textepsilon}{\textlengthmark}{\textsubring{b}}{r}{a}} \dicPos{f}[1] \dicFlx{(‑u, ‑ur)}[7] \dicDirectTranslationCS{puritán(ka), mravokárce, mravokárkyně, moralista, moralistka}
\dicEntry[teprulegur] \dicTerm{tepru··legur} \dicIPA{{t\smash{\textsuperscript{h}}}{\textepsilon}{\textlengthmark}{\textsubring{b}}{r}{\textscy}{l}{\textepsilon}{\textbabygamma}{\textscy}{\textsubring{r}}} \dicPos{adj}[1]\dicFlx{}[-8] \dicDirectTranslationCS{puritánský, prudérní, upjatý}
\dicEntry[tepruskapur] \dicTerm{tepru··skap|ur} \dicIPA{{t\smash{\textsuperscript{h}}}{\textepsilon}{\textlengthmark}{\textsubring{b}}{r}{\textscy}{s}{\r{g}}{a}{\textsubring{b}}{\textscy}{\textsubring{r}}} \dicPos{m}[10] \dicFlx{(‑ar)}[15] \dicDirectTranslationCS{puritánství, pruderie, upjatost}
\dicEntry[termíti] \dicTerm{termít|i} \dicIPA{{t\smash{\textsuperscript{h}}}{\textepsilon}{r}{m}{i}{\textsubring{d}}{\textsci}} \dicPos{m}[1] \dicFlx{(‑a, ‑ar)}[1] \dicFieldCat{zool.} \dicDirectTranslationCS{termit, všekaz} \textit{(l.~{\textLA{Isoptera}})}  \dicsymPhoto\ 
\dicFigure{ds_image_termiti_0_1.jpg}{Termíti}{Termíti - Neznámý, PD}
\dicEntry[termostat] \dicTerm{termost|at} \dicIPA{{t\smash{\textsuperscript{h}}}{\textepsilon}{r}{m}{\textopeno}{s}{\textsubring{d}}{a}{\textsubring{d}}} \dicPos{n}[2] \dicFlx{(‑ats, ‑öt)}[8] \dicDirectTranslationCS{termostat}
\dicEntry[terpentína] \dicTerm{terpentín|a} \dicIPA{{t\smash{\textsuperscript{h}}}{\textepsilon}{\textsubring{r}}{\textsubring{b}}{\textepsilon}{\textsubring{n}}{\textsubring{d}}{i}{n}{a}} \dicPos{f}[1] \dicFlx{(‑u)}[5] \dicDirectTranslationCS{terpentýn}
\dicEntry[terror] \dicTerm{terror} \dicIPA{{t\smash{\textsuperscript{h}}}{\textepsilon}{r}{\textlengthmark}{\textopeno}{\textsubring{r}}} \dicPos{m}[4] \dicFlx{(‑s)}[18] \dicLangCat{hovor.} \dicSynonym{grimmd} \dicDirectTranslationCS{teror}
\dicEntry[terroristi] \dicTerm{terror··ist|i} \dicIPA{{t\smash{\textsuperscript{h}}}{\textepsilon}{r}{\textlengthmark}{\textopeno}{r}{\textsci}{s}{\textsubring{d}}{\textsci}} \dicPos{m}[1] \dicFlx{(‑a, ‑ar)}[1] \dicLangCat{hovor.} \dicSynonym{hermdarverkamaður} \dicDirectTranslationCS{terorista, teroristka}
\dicEntry[terta] \dicTerm{tert|a} \dicIPA{{t\smash{\textsuperscript{h}}}{\textepsilon}{\textsubring{r}}{\textsubring{d}}{a}} \dicPos{f}[1] \dicFlx{(‑u, ‑ur)}[19] \dicFieldCat{kulin.} \dicDirectTranslationCS{dort, koláč}
\dicEntry[tertuhnífur] \dicTerm{tertu··hníf|ur} \dicIPA{{t\smash{\textsuperscript{h}}}{\textepsilon}{\textsubring{r}}{\textsubring{d}}{\textscy}{\textsubring{n}}{i}{v}{\textscy}{\textsubring{r}}} \dicPos{m}[6] \dicFlx{(‑s, ‑ar)}[22] \dicDirectTranslationCS{dortový nůž}
\dicEntry[tertusneið] \dicTerm{tertu··sneið} \dicIPA{{t\smash{\textsuperscript{h}}}{\textepsilon}{\textsubring{r}}{\textsubring{d}}{\textscy}{s}{\textsubring{d}}{n}{ei}{\texttheta}} \dicPos{f}[4] \dicFlx{(‑ar, ‑ar)}[1] \dicDirectTranslationCS{kus dortu\,/\addthin koláče}
\dicEntry[tesía] \dicTerm{te··sí|a} \dicIPA{{t\smash{\textsuperscript{h}}}{\textepsilon}{\textlengthmark}{s}{i}{j}{a}} \dicPos{f}[1] \dicFlx{(‑u, ‑ur)}[7] \dicDirectTranslationCS{sítko\,/\addthin cedníček na čaj}
\dicEntry[teskeið] \dicTerm{te··skeið} \dicIPA{{t\smash{\textsuperscript{h}}}{\textepsilon}{\textlengthmark}{s}{\r{\textObardotlessj}}{ei}{\texttheta}} \dicPos{f}[4] \dicFlx{(‑ar, ‑ar)}[1] \dicDirectTranslationCS{čajová\,/\addthin kávová lžička}
\dicEntry[testa] \dicTerm{test|a} \dicIPA{{t\smash{\textsuperscript{h}}}{\textepsilon}{s}{\textsubring{d}}{a}} \dicPos{v}[1] \dicFlx{(‑aði)}[1] \dicFlx{acc} \dicLangCat{hovor.} \dicDirectTranslationCS{(o)testovat}
\dicEntry[testamenti] \dicTerm{testamenti} \dicIPA{{t\smash{\textsuperscript{h}}}{\textepsilon}{s}{\textsubring{d}}{a}{m}{\textepsilon}{\textsubring{n}}{\textsubring{d}}{\textsci}} \dicPos{n}[2] \dicFlx{(‑s, ‑)}[14] \dicPhraseIS{Gamla testamentið} \dicFieldCat{náb.} \dicDirectTranslationCS{Starý zákon};  \dicPhraseIS{Nýja testamentið} \dicFieldCat{náb.} \dicDirectTranslationCS{Nový zákon}
\dicEntry[tetur] \dicTerm{tetur} \dicIPA{{t\smash{\textsuperscript{h}}}{\textepsilon}{\textlengthmark}{\textsubring{d}}{\textscy}{\textsubring{r}}} \dicPos{n}[2] \dicFlx{(‑s, ‑)}[25] \dicSynonym{vesalingur} \dicDirectTranslationCS{ubožák, ubožačka, chudák, chudačka};  \dicPhraseIS{hvorki tangur né tetur af e‑u\,/\addthin e‑m} \dicLangCat{přen.} \dicDirectTranslationCS{ani stopy po (čem\,/\addthin kom), ani památky po (čem\,/\addthin kom)}
\dicEntry[texta] \dicTerm{text|a} \dicIPA{{t\smash{\textsuperscript{h}}}{\textepsilon}{x}{s}{\textsubring{d}}{a}} \dicPos{v}[1] \dicFlx{(‑aði)}[1] \dicFlx{acc} \dicDirectTranslationCS{(o)textovat}
\dicEntry[textafræði] \dicTerm{texta··fræð|i} \dicIPA{{t\smash{\textsuperscript{h}}}{\textepsilon}{x}{s}{\textsubring{d}}{a}{f}{r}{a}{i}{ð}{\textsci}} \dicPos{f}[3] \dicFlx{(‑i)}[3] \dicDirectTranslationCS{filologie}
\dicEntry[textafræðingur] \dicTerm{texta·fræð··ing|ur} \dicIPA{{t\smash{\textsuperscript{h}}}{\textepsilon}{x}{s}{\textsubring{d}}{a}{f}{r}{a}{i}{ð}{i}{\ng}{\r{g}}{\textscy}{\textsubring{r}}} \dicPos{m}[6] \dicFlx{(‑s, ‑ar)}[8] \dicDirectTranslationCS{filolog, filoložka}
\dicEntry[textarýni] \dicTerm{texta··rýn|i} \dicIPA{{t\smash{\textsuperscript{h}}}{\textepsilon}{x}{s}{\textsubring{d}}{a}{r}{i}{n}{\textsci}} \dicPos{f}[3] \dicFlx{(‑i)}[3] \dicDirectTranslationCS{textová kritika}
\dicEntry[textaskrá] \dicTerm{texta··skrá} \dicIPA{{t\smash{\textsuperscript{h}}}{\textepsilon}{x}{s}{\textsubring{d}}{a}{s}{\r{g}}{r}{au}} \dicPos{f}[4] \dicFlx{(‑r\,/\addthin ‑ar, ‑r)}[21] \dicFieldCat{poč.} \dicDirectTranslationCS{textový soubor}
\dicEntry[textavarp] \dicTerm{texta··|varp} \dicIPA{{t\smash{\textsuperscript{h}}}{\textepsilon}{x}{s}{\textsubring{d}}{a}{v}{a}{\textsubring{r}}{\textsubring{b}}} \dicPos{n}[2] \dicFlx{(‑varps, ‑vörp)}[8] \dicDirectTranslationCS{teletext}
\dicEntry[textavinnsla] \dicTerm{texta··vinnsl|a} \dicIPA{{t\smash{\textsuperscript{h}}}{\textepsilon}{x}{s}{\textsubring{d}}{a}{v}{\textsci}{n}{s}{\textsubring{d}}{l}{a}} \dicPos{f}[1] \dicFlx{(‑u)}[5] \dicFieldCat{poč.} \dicDirectTranslationCS{(elektronické) zpracování textu}
\dicEntry[texti] \dicTerm{text|i} \dicsymFrequent\  \dicIPA{{t\smash{\textsuperscript{h}}}{\textepsilon}{x}{s}{\textsubring{d}}{\textsci}} \dicPos{m}[1] \dicFlx{(‑a, ‑ar)}[1] \textbf{1.} \dicDirectTranslationCS{text} \dicExampleIS{lesa textann} \dicExampleCS{číst text}  \textbf{2.} \dicDirectTranslationCS{titulek} \dicIndirectTranslationCS{(psaný text na filmu)}
\dicEntry[textíll] \dicTerm{textíl|l} \dicIPA{{t\smash{\textsuperscript{h}}}{\textepsilon}{x}{s}{\textsubring{d}}{i}{\textsubring{d}}{\textsubring{l}}} \dicPos{m}[6] \dicFlx{(‑s, ‑ar)}[48] \dicDirectTranslationCS{textil}
\dicEntry[textílmennt] \dicTerm{textíl··mennt} \dicIPA{{t\smash{\textsuperscript{h}}}{\textepsilon}{x}{s}{\textsubring{d}}{i}{l}{m}{\textepsilon}{\textsubring{n}}{\textsubring{d}}} \dicPos{f}[7] \dicFlx{(‑ar, ‑ir)}[1] \dicDirectTranslationCS{ruční práce} \dicIndirectTranslationCS{(školní předmět)}
\dicEntry[teyga] \dicTerm{teyg|a} \dicIPA{{t\smash{\textsuperscript{h}}}{ei}{\textlengthmark}{\textbabygamma}{a}} \dicPos{v}[1] \dicFlx{(‑aði)}[34] \dicFlx{acc} \textbf{1.} \dicDirectTranslationCS{hltavě pít, chlemtat} \dicExampleIS{teyga vatnið} \dicExampleCS{chlemtat vodu}  \textbf{2.} \dicDirectTranslationCS{nasát, nasávat, vdechnout, vdechovat}
\begin{xtolerant}{}{1pt}
\dicEntry[teygilegur] \dicTerm{teygi··legur} \dicIPA{{t\smash{\textsuperscript{h}}}{ei}{j}{\textlengthmark}{\textsci}{l}{\textepsilon}{\textbabygamma}{\textscy}{\textsubring{r}}} \dicPos{adj}[1]\dicFlx{}[-8] \textbf{1.} \dicSynonym{teygjanlegur} \dicDirectTranslationCS{pruž\-ný, natažitelný}  \textbf{2.} \dicLangCat{zast.} \dicSynonym{freistandi} \dicDirectTranslationCS{svůdný, okouzlující}
\end{xtolerant}
\dicEntry[teygja] \dicTerm{teygj|a\smash{\textsuperscript{1}}} \dicIPA{{t\smash{\textsuperscript{h}}}{ei}{j}{\textlengthmark}{a}} \dicPos{f}[1] \dicFlx{(‑u, ‑ur)}[7] \textbf{1.} \dicSynonym*{mjór gúmhringur} \dicDirectTranslationCS{guma, gumička}  \textbf{2.} \dicSynonym*{hárteygja} \dicDirectTranslationCS{gumička (do vlasů)}  \textbf{3.} \dicSynonym*{teygjuhreyfing} \dicDirectTranslationCS{natažení, protažení, protahování} \dicIndirectTranslationCS{(zvláště v~pl)} \dicExampleIS{teygjuæfing} \dicExampleCS{protahovací cvičení}  \textbf{4.} \dicSynonym{teygjanleiki} \dicDirectTranslationCS{tažnost, natažitelnost, elastičnost} \dicExampleIS{teygja í efni} \dicExampleCS{tažnost materiálu}
\dicEntry[teygja] \dicTerm{teyg|ja\smash{\textsuperscript{2}}} \dicsymFrequent\  \dicIPA{{t\smash{\textsuperscript{h}}}{ei}{j}{\textlengthmark}{a}} \dicPos{v}[2] \dicFlx{(‑ði, ‑t)}[89] \textbf{1.} \dicDirectTranslationCS{natáhnout, natahovat, vytáhnout, vytahovat} \dicExampleIS{teygja út höndina} \dicExampleCS{natáhnout ruku};  \dicPhraseIS{teygja sig eftir e‑u} \dicDirectTranslationCS{natáhnout se po (čem)}  \textbf{2.} \dicDirectTranslationCS{protáhnout, protahovat (svaly ap.)} \dicExampleIS{teygja sig eftir hlaupið} \dicExampleCS{protáhnout se po běhání};  \dicPhraseIS{teygja á e‑u} \dicDirectTranslationCS{protáhnout (co) (svaly ap.)};  \dicIdiom{teygjast}{ \dicPhraseIS{teygjast}} \dicFlx{refl} {\textbf{a.}} \dicSynonym*{lengjast} \dicDirectTranslationCS{protáhnout se, natáhnout se, prodloužit se (dovolená ap.)} \dicExampleIS{Ferðin teygðist til morguns.} \dicExampleCS{Cesta se protáhla do rána.};  {\textbf{b.}} \dicDirectTranslationCS{být pružný, natáhnout se, vytáhnout se}
\dicEntry[teygjanlegur] \dicTerm{teygjan··legur} \dicIPA{{t\smash{\textsuperscript{h}}}{ei}{j}{a}{n}{l}{\textepsilon}{\textbabygamma}{\textscy}{\textsubring{r}}} \dicPos{adj}[1]\dicFlx{}[-8] \dicDirectTranslationCS{pružný, natažitelný, elastický}
\dicEntry[teygjanleiki] \dicTerm{teygjan··leik|i} \dicIPA{{t\smash{\textsuperscript{h}}}{ei}{j}{a}{n}{l}{ei}{\r{\textObardotlessj}}{\textsci}} \dicPos{m}[1] \dicFlx{(‑a)}[3] \dicDirectTranslationCS{pružnost, elasticita}
\dicEntry[teygjuband] \dicTerm{teygju··|band} \dicIPA{{t\smash{\textsuperscript{h}}}{ei}{j}{\textlengthmark}{\textscy}{\textsubring{b}}{a}{n}{\textsubring{d}}} \dicPos{n}[2] \dicFlx{(‑bands, ‑bönd)}[8] \dicDirectTranslationCS{elastická páska, gumička}
\dicEntry[teygjudýr] \dicTerm{teygju··dýr} \dicIPA{{t\smash{\textsuperscript{h}}}{ei}{j}{\textlengthmark}{\textscy}{\textsubring{d}}{i}{\textsubring{r}}} \dicPos{n}[2] \dicFlx{(‑s, ‑)}[5] \dicFieldCat{zool.} \dicDirectTranslationCS{měňavka, améba} \textit{(l.~{\textLA{Amoebida}})}
\dicEntry[teygjustökk] \dicTerm{teygju··stökk} \dicIPA{{t\smash{\textsuperscript{h}}}{ei}{j}{\textlengthmark}{\textscy}{s}{\textsubring{d}}{\oe}{h}{\r{g}}} \dicPos{n}[2] \dicFlx{(‑s, ‑)}[5] \dicDirectTranslationCS{bungee jumping}
\dicEntry[teygur] \dicTerm{teyg|ur} \dicIPA{{t\smash{\textsuperscript{h}}}{ei}{\textlengthmark}{\textbabygamma}{\textscy}{\textsubring{r}}} \dicPos{m}[6] \dicFlx{(‑s, ‑ar)}[15] \dicSynonym{sopi} \dicDirectTranslationCS{hlt, lok};  \dicPhraseIS{drekka e‑ð í einum teyg} \dicDirectTranslationCS{vypít (co) jedním lokem}
\dicEntry[teyma] \dicTerm{teym|a} \dicIPA{{t\smash{\textsuperscript{h}}}{ei}{\textlengthmark}{m}{a}} \dicPos{v}[2] \dicFlx{(‑di, ‑t)}[141] \dicFlx{acc} \textbf{1.} \dicSynonym*{leiða í taumi} \dicDirectTranslationCS{vést za uzdu\,/\addthin otěže (koně ap.)}  \textbf{2.} \dicDirectTranslationCS{tahat, vláčet (z~obchodu do obchodu ap.)}
\dicEntry[teymi] \dicTerm{teymi} \dicIPA{{t\smash{\textsuperscript{h}}}{ei}{\textlengthmark}{m}{\textsci}} \dicPos{n}[2] \dicFlx{(‑s, ‑)}[14] \dicDirectTranslationCS{tým, kolektiv}
\dicEntry[té] \dicTerm{té} \dicIPA{{t\smash{\textsuperscript{h}}}{j}{\textepsilon}{\textlengthmark}} \dicPos{n}[2] \dicFlx{(‑s)}[2] \textbf{1.} \dicDirectTranslationCS{téčko, písmeno \clqq t\crqq }  \textbf{2.} \dicDirectTranslationCS{téčko} \dicIndirectTranslationCS{(co připomíná tvarem písmeno \clqq t\crqq )};  \dicIdiom{té}{ \dicPhraseIS{láta e‑m e‑ð í té}} \dicDirectTranslationCS{poskytnout (komu co), předat (komu co)} \dicExampleIS{láta honum upplýsingar í té} \dicExampleCS{předat mu informace}
\dicEntry[tékk] \dicTerm{tékk} \dicIPA{{t\smash{\textsuperscript{h}}}{j}{\textepsilon}{h}{\r{g}}} \dicPos{n}[2] \dicFlx{(‑s, ‑)}[5] \dicLangCat{hovor.} \dicDirectTranslationCS{kontrola, ověření}
\dicEntry[tékka] \dicTerm{tékk|a} \dicIPA{{t\smash{\textsuperscript{h}}}{j}{\textepsilon}{h}{\r{g}}{a}} \dicPos{v}[1] \dicFlx{(‑aði)}[1] \dicLangCat{hovor.} \dicSynonym{athuga} \dicDirectTranslationCS{(z)kontrolovat, ověřit, ověřovat};  \dicPhraseIS{tékka á e‑u} \dicDirectTranslationCS{zkontrolovat (co)};  \dicIdiom{tékka}[inn]{ \dicPhraseIS{tékka e‑n inn}} \dicDirectTranslationCS{odbavit (koho) (k~odletu ap.)} \dicExampleIS{tékka sig inn} \dicExampleCS{odbavit se};  \dicIdiom{tékka}[út]{ \dicPhraseIS{tékka e‑n út}} \dicDirectTranslationCS{odhlásit (koho), odregistrovat (koho) (ze seznamu ap.)}
\dicEntry[tékkareikningur] \dicTerm{tékka··reikn·ing|ur} \dicIPA{{t\smash{\textsuperscript{h}}}{j}{\textepsilon}{h}{\r{g}}{a}{r}{ei}{h}{\r{g}}{n}{i}{\ng}{\r{g}}{\textscy}{\textsubring{r}}} \dicPos{m}[6] \dicFlx{(‑s, ‑ar)}[8] \dicFieldCat{ekon.} \dicDirectTranslationCS{běžný účet}
\dicEntry[tékki] \dicTerm{tékk|i\smash{\textsuperscript{1}}} \dicIPA{{t\smash{\textsuperscript{h}}}{j}{\textepsilon}{h}{\r{\textObardotlessj}}{\textsci}} \dicPos{m}[1] \dicFlx{(‑a, ‑ar)}[1] \dicFieldCat{ekon.} \dicSynonym{ávísun} \dicDirectTranslationCS{šek}
\dicEntry[Tékki] \dicTerm{Tékk|i\smash{\textsuperscript{2}}} \dicIPA{{t\smash{\textsuperscript{h}}}{j}{\textepsilon}{h}{\r{\textObardotlessj}}{\textsci}} \dicPos{m}[1] \dicFlx{(‑a, ‑ar)}[1] \dicDirectTranslationCS{Čech, Češka}
\dicEntry[Tékkland] \dicTerm{Tékk··land} \dicIPA{{t\smash{\textsuperscript{h}}}{j}{\textepsilon}{h}{\r{g}}{l}{a}{n}{\textsubring{d}}} \dicPos{n}[2] \dicFlx{(‑s)}[4] \dicFieldCat{geog.} \dicDirectTranslationCS{Čechy, Česko, Česká republika}
\dicEntry[tékkn.] \dicTerm{tékkn.} \dicPos{zkr} \dicPhraseIS{tékkneska, tékkneskur} \dicDirectTranslationCS{čeština, český}
\dicEntry[tékkneska] \dicTerm{tékk··nesk|a} \dicIPA{{t\smash{\textsuperscript{h}}}{j}{\textepsilon}{h}{\r{g}}{n}{\textepsilon}{s}{\r{g}}{a}} \dicPos{f}[1] \dicFlx{(‑u)}[5] \dicDirectTranslationCS{čeština}
\dicEntry[tékkneskur] \dicTerm{tékk··neskur} \dicIPA{{t\smash{\textsuperscript{h}}}{j}{\textepsilon}{h}{\r{g}}{n}{\textepsilon}{s}{\r{g}}{\textscy}{\textsubring{r}}} \dicPos{adj}[1]\dicFlx{}[-1] \dicDirectTranslationCS{český}
\dicEntry[Tékkóslóvakía] \dicTerm{Tékkó··slóvakí|a} \dicIPA{{t\smash{\textsuperscript{h}}}{j}{\textepsilon}{h}{\r{g}}{ou}{s}{\textsubring{d}}{l}{ou}{v}{a}{\r{\textObardotlessj}}{i}{j}{a}} \dicPos{f}[1] \dicFlx{(‑u)}[6] \dicFieldCat{geog. hist.} \dicDirectTranslationCS{Československo}
\dicEntry[Téténía] \dicTerm{Téténí|a} \dicIPA{{t\smash{\textsuperscript{h}}}{j}{\textepsilon}{\textlengthmark}{\textsubring{d}}{j}{\textepsilon}{n}{i}{j}{a}} \dicPos{f}[1] \dicFlx{(‑u)}[6] \dicFieldCat{geog.} \dicDirectTranslationCS{Čečensko}
\dicEntry[tfn.] \dicTerm{tfn.} \dicPos{zkr} \dicPhraseIS{tilvísunarfornafn} \dicFieldCat{jaz.} \dicDirectTranslationCS{vztažné zájmeno}
\dicEntry[tif] \dicTerm{tif} \dicIPA{{t\smash{\textsuperscript{h}}}{\textsci}{\textlengthmark}{f}} \dicPos{n}[2] \dicFlx{(‑s)}[2] \textbf{1.} \dicSynonym*{hröð hreyfing} \dicDirectTranslationCS{cupitání, cupání}  \textbf{2.} \dicSynonym*{lágt og títt ganghljóð} \dicDirectTranslationCS{tikot, tikání} \dicExampleIS{tif í klukku} \dicExampleCS{tikání hodin}
\dicEntry[tifa] \dicTerm{tif|a} \dicIPA{{t\smash{\textsuperscript{h}}}{\textsci}{\textlengthmark}{v}{a}} \dicPos{v}[1] \dicFlx{(‑aði)}[1] \textbf{1.} \dicSynonym*{ganga stuttum skrefum} \dicDirectTranslationCS{cupitat, cupat}  \textbf{2.} \dicDirectTranslationCS{tikat (hodiny ap.)} \dicExampleIS{Klukkan tifar.} \dicExampleCS{Hodiny tikají.}
\dicEntry[tifandi] \dicTerm{tif··andi} \dicIPA{{t\smash{\textsuperscript{h}}}{\textsci}{\textlengthmark}{v}{a}{n}{\textsubring{d}}{\textsci}} \dicPos{adj}[13] \dicFlx{indecl}[1] \dicDirectTranslationCS{tikající}
\dicEntry[tifstjarna] \dicTerm{tif··|stjarna} \dicIPA{{t\smash{\textsuperscript{h}}}{\textsci}{f}{s}{\textsubring{d}}{j}{a}{r}{\textsubring{d}}{n}{a}} \dicPos{f}[1] \dicFlx{(‑stjörnu, ‑stjörnur)}[8] \dicFieldCat{astro.} \dicDirectTranslationCS{pulzar, pulsar}
\dicEntry[tiginn] \dicTerm{tiginn} \dicIPA{{t\smash{\textsuperscript{h}}}{i}{j}{\textlengthmark}{\textsci}{\textsubring{n}}} \dicPos{adj}[6]\dicFlx{}[-2] \dicDirectTranslationCS{urozený, vznešený} \dicExampleIS{tignir gestir} \dicExampleCS{urození hosté}
\dicEntry[tign] \dicTerm{tign} \dicIPA{{t\smash{\textsuperscript{h}}}{\textsci}{\r{g}}{\textsubring{n}}} \dicPos{f}[7] \dicFlx{(‑ar, ‑ir)}[1] \textbf{1.} \dicSynonym*{ættgöfgi} \dicDirectTranslationCS{urozenost, vznešenost};  \dicPhraseIS{hækka í tign} \dicDirectTranslationCS{povýšit, postoupit (ve funkci ap.)} \dicExampleIS{Hann var hækkaður í tign.} \dicExampleCS{Byl povýšen.};  \dicPhraseIS{lækka í tign} \dicDirectTranslationCS{být sesazen\,/\addthin degradován (ve funkci ap.)}  \textbf{2.} \dicDirectTranslationCS{Výsost} \dicExampleIS{yðar tign} \dicExampleCS{Vaše Výsosti}
\dicEntry[tigna] \dicTerm{tign|a} \dicIPA{{t\smash{\textsuperscript{h}}}{\textsci}{\r{g}}{n}{a}} \dicPos{v}[1] \dicFlx{(‑aði)}[1] \dicFlx{acc} \dicSynonym{heiðra} \dicDirectTranslationCS{ctít, uctívat, velebit} \dicExampleIS{tigna guð} \dicExampleCS{velebit Boha}
\dicEntry[tignarlegur] \dicTerm{tignar··legur} \dicIPA{{t\smash{\textsuperscript{h}}}{\textsci}{\r{g}}{n}{a}{r}{l}{\textepsilon}{\textbabygamma}{\textscy}{\textsubring{r}}} \dicPos{adj}[1]\dicFlx{}[-8] \textbf{1.} \dicSynonym{tígulegur} \dicDirectTranslationCS{urozený, vznešený}  \textbf{2.} \dicSynonym*{mikilúðlegur} \dicDirectTranslationCS{majestátní, vznešený, velkolepý} \dicExampleIS{tignarlegt fjall} \dicExampleCS{majestátní hora}
\dicEntry[tignarmaður] \dicTerm{tignar··|maður} \dicIPA{{t\smash{\textsuperscript{h}}}{\textsci}{\r{g}}{n}{a}{r}{m}{a}{ð}{\textscy}{\textsubring{r}}} \dicPos{m}[13] \dicFlx{(‑manns, ‑menn)}[2] \dicDirectTranslationCS{hodnostář(ka)}
\dicEntry[tigulegur] \dicTerm{tigu··legur} \dicIPA{{t\smash{\textsuperscript{h}}}{\textsci}{\textlengthmark}{\textbabygamma}{\textscy}{l}{\textepsilon}{\textbabygamma}{\textscy}{\textsubring{r}}} \dicPos{adj}[1]\dicFlx{}[-8] \dicLink{tígulegur}
\dicEntry[tigull] \dicTerm{tig|ull} \dicIPA{{t\smash{\textsuperscript{h}}}{\textsci}{\textlengthmark}{\textbabygamma}{\textscy}{\textsubring{d}}{\textsubring{l}}} \dicPos{m}[6] \dicFlx{(‑uls, ‑lar)}[31] \dicLink{tígull}
\dicEntry[tigulsteinn] \dicTerm{tigul··stein|n} \dicIPA{{t\smash{\textsuperscript{h}}}{\textsci}{\textlengthmark}{\textbabygamma}{\textscy}{l}{s}{\textsubring{d}}{ei}{\textsubring{d}}{\textsubring{n}}} \dicPos{m}[6] \dicFlx{(‑s, ‑ar)}[42] \dicLink{tígulsteinn}
\dicEntry[tikk] \dicTerm{tikk} \dicIPA{{t\smash{\textsuperscript{h}}}{\textsci}{h}{\r{g}}} \dicPos{n}[2] \dicFlx{(‑s)}[2] \dicDirectTranslationCS{tik, tikání, tikot}
\dicEntry[tikka] \dicTerm{tikk|a} \dicIPA{{t\smash{\textsuperscript{h}}}{\textsci}{h}{\r{g}}{a}} \dicPos{v}[1] \dicFlx{(‑aði)}[1] \dicDirectTranslationCS{tikat} \dicExampleIS{Klukkan tikkar.} \dicExampleCS{Hodiny tikají.}
\dicEntry[til] \dicTerm{til} \dicsymFrequent\  \dicIPA{{t\smash{\textsuperscript{h}}}{\textsci}{\textlengthmark}{\textsubring{l}}} \dicPos{prep\,/\addthin adv} \dicFlx{gen} \textbf{1.} \dicDirectTranslationCS{k(e), do} \dicIndirectTranslationCS{(o~pohybu k~něčemu)} \dicExampleIS{Hún fór til kirkju.} \dicExampleCS{Šla do kostela.}  \textbf{2.} \dicDirectTranslationCS{do} \dicIndirectTranslationCS{(o~čase)} \dicExampleIS{til jóla} \dicExampleCS{do Vánoc}  \textbf{3.} \dicDirectTranslationCS{pro, k(e)} \dicIndirectTranslationCS{(o~účelu, smyslu něčeho)} \dicExampleIS{bréf til hans} \dicExampleCS{dopis pro něj}  \textbf{4.} \dicDirectTranslationCS{až, mezi} \dicIndirectTranslationCS{(vyjadřuje rozpětí údaje)} \dicExampleIS{tvö til þrjú þúsund} \dicExampleCS{dva až tři tisíce};  \dicIdiom{að}[til]{ \dicPhraseIS{að e‑u til}} \dicFlx{prep} {\textbf{a.}} \dicDirectTranslationCS{během (čeho)} \dicExampleIS{að degi til} \dicExampleCS{během dne};  {\textbf{b.}} \dicDirectTranslationCS{s~přihlédnutím k~(čemu)} \dicExampleIS{að magni til} \dicExampleCS{objemově};  \dicIdiom{til}[að]{ \dicPhraseIS{til (þess) að}} \dicFlx{conj} \dicDirectTranslationCS{aby; za účelem, aby; kvůli tomu, aby} \dicExampleIS{Hann fór til Íslands til að læra íslensku.} \dicExampleCS{Odjel na Island, aby se naučil islandštinu.};  \dicIdiom{til}{ \dicPhraseIS{til baka}} \dicFlx{adv} \dicDirectTranslationCS{nazpátek, nazpět, zpátky}; { \dicPhraseIS{til dæmis}} \dicFlx{adv} \dicDirectTranslationCS{například, kupříkladu};  \dicPhraseIS{til hvers} \dicDirectTranslationCS{k~čemu, na co}; { \dicPhraseIS{til og frá}} \dicFlx{adv} \dicSynonym*{hér og hvar} \dicDirectTranslationCS{tam a~zpátky};  \dicPhraseIS{þangað til} \dicFlx{conj} \dicDirectTranslationCS{dokud, než};  \dicPhraseIS{þar til} \dicFlx{conj} \dicDirectTranslationCS{dokud}
\dicEntry[tilbeiðsla] \dicTerm{til··beiðsl|a} \dicIPA{{t\smash{\textsuperscript{h}}}{\textsci}{l}{\textsubring{b}}{ei}{ð}{s}{\textsubring{d}}{l}{a}} \dicPos{f}[1] \dicFlx{(‑u)}[5] \dicDirectTranslationCS{uctívání, zbožňování, adorace}
\dicEntry[tilberi] \dicTerm{til··ber|i} \dicIPA{{t\smash{\textsuperscript{h}}}{\textsci}{l}{\textsubring{b}}{\textepsilon}{r}{\textsci}} \dicPos{m}[1] \dicFlx{(‑a, ‑ar)}[1] \dicFieldCat{pov.} \dicIndirectTranslationCS{stvoření, které dojí krávy a~ovce hospodářů}
\dicEntry[tilboð] \dicTerm{til··boð} \dicIPA{{t\smash{\textsuperscript{h}}}{\textsci}{l}{\textsubring{b}}{\textopeno}{\texttheta}} \dicPos{n}[2] \dicFlx{(‑s, ‑)}[5] \textbf{1.} \dicSynonym*{gott boð} \dicDirectTranslationCS{nabídka, návrh} \dicExampleIS{Stendur tilboðið enn?} \dicExampleCS{Platí ještě ta nabídka?}  \textbf{2.} \dicSynonym*{viðskiptaloforð} \dicDirectTranslationCS{(výhodná) nabídka} \dicExampleIS{tilboð á vöru} \dicExampleCS{nabídka na zboží};  \dicPhraseIS{gera tilboð í e‑u} \dicDirectTranslationCS{nabídnout (výhodně) (co)};  \dicPhraseIS{leggja fram tilboð} \dicDirectTranslationCS{učinit\,/\addthin předložit nabídku}
\dicEntry[tilbót] \dicTerm{til··|bót} \dicIPA{{t\smash{\textsuperscript{h}}}{\textsci}{l}{\textsubring{b}}{ou}{\textsubring{d}}} \dicPos{f}[8] \dicFlx{(‑bótar, ‑bætur)}[5] \dicSynonym{viðbót} \dicDirectTranslationCS{dodatek, přidání} \dicExampleIS{í tilbót} \dicExampleCS{navrch}
\dicEntry[tilbreyting] \textls[15]{\dicTerm{til··breyt·ing} \dicIPA{{t\smash{\textsuperscript{h}}}{\textsci}{l}{\textsubring{b}}{r}{ei}{\textsubring{d}}{i}{\ng}{\r{g}}} \dicPos{f}[4] \dicFlx{(‑ar)}[7] \dicDirectTranslationCS{změna, obměna}} \dicExampleIS{hjóla í vinnuna til tilbreytingar} \dicExampleCS{jet pro změnu do práce na kole}
\dicEntry[tilbreytingalaus] \dicTerm{til·breytinga··laus} \dicIPA{{t\smash{\textsuperscript{h}}}{\textsci}{l}{\textsubring{b}}{r}{ei}{\textsubring{d}}{i}{\ng}{\r{g}}{a}{l}{\oe i}{s}} \dicPos{adj}[5]\dicFlx{}[-1] \dicLink{tilbreytingarlaus}
\dicEntry[tilbreytingaleysi] \dicTerm{til·breytinga··leysi} \dicIPA{{t\smash{\textsuperscript{h}}}{\textsci}{l}{\textsubring{b}}{r}{ei}{\textsubring{d}}{i}{\ng}{\r{g}}{a}{l}{ei}{s}{\textsci}} \dicPos{n}[2] \dicFlx{(‑s)}[20] \dicLink{tilbreytingarleysi}
\dicEntry[tilbreytingarlaus] \dicTerm{til·breytingar··laus}\dicTerm{, tilbreytingalaus} \dicIPA{{t\smash{\textsuperscript{h}}}\-{\textsci}\-{l}\-{\textsubring{b}}\-{r}\-{ei}\-{\textsubring{d}}\-{i}\-{\ng}\-{\r{g}}\-{a}\-{r}\-{l}\-{\oe i}\-{s}\-} \dicPos{adj}[5]\dicFlx{}[-1] \dicDirectTranslationCS{neměnný, monotónní, jednotvárný}
\dicEntry[tilbreytingarleysi] \dicTerm{til·breytingar··leysi}\textls[-10]{\dicTerm{, tilbreytingaleysi} \dicIPA{{t\smash{\textsuperscript{h}}}\-{\textsci}\-{l}\-{\textsubring{b}}\-{r}\-{ei}\-{\textsubring{d}}\-{i}\-{\ng}\-{\r{g}}\-{a}\-{r}\-{l}\-{ei}\-{s}\-{\textsci}}} \dicPos{n}[2] \dicFlx{(‑s)}[20] \dicDirectTranslationCS{neměnnost, monotónnost, jednotvárnost}
\dicEntry[tilbrigði] \dicTerm{til··brigði} \dicIPA{{t\smash{\textsuperscript{h}}}{\textsci}{l}{\textsubring{b}}{r}{\textsci}{\textbabygamma}{ð}{\textsci}} \dicPos{n}[2] \dicFlx{(‑s, ‑)}[14] \textbf{1.} \dicSynonym{umskipti} \dicDirectTranslationCS{obměna, varianta, variace}  \textbf{2.} \dicFieldCat{hud.} \dicDirectTranslationCS{variace}
\dicEntry[tilbúinn] \dicTerm{til··búinn} \dicsymFrequent\  \dicIPA{{t\smash{\textsuperscript{h}}}{\textsci}{l}{\textsubring{b}}{u}{\textsci}{\textsubring{n}}} \dicPos{adj}[6]\dicFlx{}[-6] \textbf{1.} \dicSynonym{viðbúinn} \dicDirectTranslationCS{připravený, přichystaný, nachystaný, hotový} \dicExampleIS{Ertu tilbúinn?} \dicExampleCS{Jsi připraven?}  \textbf{2.} \dicSynonym*{búinn til} \dicDirectTranslationCS{připravený (k~použití ap.)} \dicIndirectTranslationCS{(jistým způsobem předem zpracovaný)};  \dicPhraseIS{tilbúinn réttur} \dicDirectTranslationCS{polotovar}
\dicEntry[tilbúningur] \dicTerm{til··bú·ning|ur} \dicIPA{{t\smash{\textsuperscript{h}}}{\textsci}{l}{\textsubring{b}}{u}{n}{i}{\ng}{\r{g}}{\textscy}{\textsubring{r}}} \dicPos{m}[6] \dicFlx{(‑s)}[9] \dicSynonym{ósannindi} \dicDirectTranslationCS{smyšlenka, výmysl}
\dicEntry[tildra] \dicTerm{tildr|a} \dicIPA{{t\smash{\textsuperscript{h}}}{\textsci}{l}{\textsubring{d}}{r}{a}} \dicPos{f}[1] \dicFlx{(‑u, ‑ur)}[7] \dicFieldCat{zool.} \dicDirectTranslationCS{kameňáček pestrý} \textit{(l.~{\textLA{Arenaria interpres}})}  \dicsymPhoto\ 
\dicFigure{21073.jpg}{Tildra}{Tildra - Bowman, Tim, Biolib, PD}
\dicEntry[tildrög] \dicTerm{til··drög} \dicIPA{{t\smash{\textsuperscript{h}}}{\textsci}{l}{\textsubring{d}}{r}{\oe}{x}} \dicPos{n}[2] \dicFlx{pl}[9] \dicSynonym{aðdragandi} \dicDirectTranslationCS{příčina, důvod} \dicExampleIS{tildrög slyssins} \dicExampleCS{příčina nehody}
\dicEntry[tildur] \dicTerm{tildur} \dicIPA{{t\smash{\textsuperscript{h}}}{\textsci}{l}{\textsubring{d}}{\textscy}{\textsubring{r}}} \dicPos{n}[2] \dicFlx{(‑s)}[28] \dicDirectTranslationCS{parádivost, fintivost}
\dicEntry[tilefni] \dicTerm{til··efni} \dicsymFrequent\  \dicIPA{{t\smash{\textsuperscript{h}}}{\textsci}{\textlengthmark}{l}{\textepsilon}{\textsubring{b}}{n}{\textsci}} \dicPos{n}[2] \dicFlx{(‑s, ‑)}[14] \dicDirectTranslationCS{příležitost, okolnost};  \dicPhraseIS{að gefnu tilefni} \dicFlx{adv} \dicDirectTranslationCS{při příležitosti};  \dicPhraseIS{af sérstöku tilefni} \dicFlx{adv} \dicDirectTranslationCS{při zvláštní příležitosti};  \dicPhraseIS{í tilefni af e‑u} \dicDirectTranslationCS{u~příležitosti (čeho)} \dicExampleIS{Fólkið hittist í tilefni af afmælinu.} \dicExampleCS{Lidé se setkali u~příležitosti oslav narozenin.}
\dicEntry[tilefnislaus] \dicTerm{til·efnis··laus} \dicIPA{{t\smash{\textsuperscript{h}}}{\textsci}{\textlengthmark}{l}{\textepsilon}{\textsubring{b}}{n}{\textsci}{s}{l}{\oe i}{s}} \dicPos{adj}[5]\dicFlx{}[-1] \dicDirectTranslationCS{bezdůvodný, neopodstatněný} \dicExampleIS{tilefnislaus mótmæli} \dicExampleCS{neopodstatněný protest}
\dicEntry[tileinka] \dicTerm{til··eink|a} \dicIPA{{t\smash{\textsuperscript{h}}}{\textsci}{\textlengthmark}{l}{ei}{\r{\ng}}{\r{g}}{a}} \dicPos{v}[1] \dicFlx{(‑aði)}[1] \dicFlx{dat + acc} \textbf{1.} \dicSynonym{helga} \dicDirectTranslationCS{věnovat, dedikovat (knihu ap.)};  \dicPhraseIS{tileinka e‑m e‑ð} \dicDirectTranslationCS{věnovat (komu co), dedikovat (komu co)} \dicExampleIS{tileinka henni bók} \dicExampleCS{věnovat jí knihu}  \textbf{2.} \dicSynonym{læra} \dicDirectTranslationCS{osvojit si, naučit se};  \dicPhraseIS{tileinka sér e‑ð} \dicDirectTranslationCS{osvojit si (co)} \dicExampleIS{tileinka sér málfræðina} \dicExampleCS{osvojit si gramatiku}
\dicEntry[tileinkun] \dicTerm{til··eink|un} \dicIPA{{t\smash{\textsuperscript{h}}}{\textsci}{\textlengthmark}{l}{ei}{\r{\ng}}{\r{g}}{\textscy}{\textsubring{n}}} \dicPos{f}[7] \dicFlx{(‑unar, ‑anir)}[8] \dicDirectTranslationCS{věnování, dedikace} \dicExampleIS{tileinkun til e‑rs} \dicExampleCS{věnování (komu)}
\dicEntry[tileygð] \dicTerm{til··eygð} \dicIPA{{t\smash{\textsuperscript{h}}}{\textsci}{\textlengthmark}{l}{ei}{\textbabygamma}{\texttheta}} \dicPos{f}[4] \dicFlx{(‑ar)}[3] \dicFieldCat{med. } \dicDirectTranslationCS{šilhavost, strabismus}
\dicEntry[tileygður] \dicTerm{til··eygður}\dicTerm{, tileygur} \dicIPA{{t\smash{\textsuperscript{h}}}\-{\textsci}\-{\textlengthmark}\-{l}\-{ei}\-{\textbabygamma}\-{ð}\-{\textscy}\-{\textsubring{r}}\-} \dicPos{adj}[2]\dicFlx{}[-1] \dicFieldCat{med.} \dicDirectTranslationCS{šilhavý};  \dicPhraseIS{vera tileygður til stúlku\,/\addthin pilts} \dicLangCat{přen.} \dicDirectTranslationCS{pokukovat po dívce\,/\addthin klukovi, šilhat za dívkou\,/\addthin klukem}
\dicEntry[tileygur] \dicTerm{til··eygur} \dicIPA{{t\smash{\textsuperscript{h}}}{\textsci}{\textlengthmark}{l}{ei}{\textbabygamma}{\textscy}{\textsubring{r}}} \dicPos{adj}[1]\dicFlx{}[-1] \dicLink{tileygður}
\dicEntry[tilfallandi] \dicTerm{til··fall·andi} \dicIPA{{t\smash{\textsuperscript{h}}}{\textsci}{l}{f}{a}{\textsubring{d}}{l}{a}{n}{\textsubring{d}}{\textsci}} \dicPos{adj}[13] \dicFlx{indecl}[1] \dicDirectTranslationCS{nahodilý, občasný} \dicExampleIS{tilfallandi kýli} \dicExampleCS{občasný vřed}
\dicEntry[tilfelli] \dicTerm{til··felli} \dicIPA{{t\smash{\textsuperscript{h}}}{\textsci}{l}{f}{\textepsilon}{\textsubring{d}}{l}{\textsci}} \dicPos{n}[2] \dicFlx{(‑s, ‑)}[14] \dicSynonym{atvik} \dicDirectTranslationCS{případ, výskyt} \dicExampleIS{Þetta mistekst í mörgum tilfellum.} \dicExampleCS{V~mnoha případech se to nepodaří.}
\dicEntry[tilfinnanlegur] \dicTerm{til·finnan··legur} \dicIPA{{t\smash{\textsuperscript{h}}}{\textsci}{l}{f}{\textsci}{n}{a}{n}{l}{\textepsilon}{\textbabygamma}{\textscy}{\textsubring{r}}} \dicPos{adj}[1]\dicFlx{}[-8] \dicSynonym{verulegur} \dicDirectTranslationCS{citelný, znatelný, značný} \dicExampleIS{tilfinnanlegt tjón} \dicExampleCS{citelná ztráta}
\dicEntry[tilfinning] \dicTerm{til··finn·ing} \dicsymFrequent\  \dicIPA{{t\smash{\textsuperscript{h}}}{\textsci}{l}{f}{\textsci}{n}{i}{\ng}{\r{g}}} \dicPos{f}[4] \dicFlx{(‑ar, ‑ar)}[5] \textbf{1.} \dicDirectTranslationCS{cit} \dicIndirectTranslationCS{(tělesná schopnost vnímat smysly, hmatem)} \dicExampleIS{hafa enga tilfinningu í fætinum} \dicExampleCS{ztratit cit v~noze}  \textbf{2.} \dicSynonym{hugboð} \dicDirectTranslationCS{pocit, tušení};  \dicPhraseIS{fá á tilfinninguna að (e‑að gerist)} \dicDirectTranslationCS{mít pocit, že (se (co) stane)};  \dicPhraseIS{hafa e‑ð á tilfinningunni} \dicDirectTranslationCS{cítit (co), tušit (co)}  \textbf{3.} \dicSynonym{kennd} \dicDirectTranslationCS{cit, pocit};  \dicPhraseIS{hafa tilfinningu fyrir e‑ð} \dicDirectTranslationCS{mít cit pro (co) (hudbu ap.)};  \dicPhraseIS{særa tilfinningar e‑rs} \dicDirectTranslationCS{zranit (čí) city}
\dicEntry[tilfinningagildi] \dicTerm{til·finninga··gildi} \dicIPA{{t\smash{\textsuperscript{h}}}{\textsci}{l}{f}{\textsci}{n}{i}{\ng}{\r{g}}{a}{\r{\textObardotlessj}}{\textsci}{l}{\textsubring{d}}{\textsci}} \dicPos{n}[2] \dicFlx{(‑s, ‑)}[14] \dicDirectTranslationCS{sentiment, citová hodnota}
\dicEntry[tilfinningalaus] \dicTerm{til·finninga··laus} \dicIPA{{t\smash{\textsuperscript{h}}}{\textsci}{l}{f}{\textsci}{n}{i}{\ng}{\r{g}}{a}{l}{\oe i}{s}} \dicPos{adj}[5]\dicFlx{}[-1] \textbf{1.} \dicSynonym{dofinn} \dicDirectTranslationCS{necitlivý, znecitlivělý}  \textbf{2.} \dicSynonym{harðbrjósta} \dicDirectTranslationCS{bezcitný, necitelný, chladný}
\dicEntry[tilfinningalega] \dicTerm{til·finninga··lega} \dicIPA{{t\smash{\textsuperscript{h}}}{\textsci}{l}{f}{\textsci}{n}{i}{\ng}{\r{g}}{a}{l}{\textepsilon}{\textbabygamma}{a}} \dicPos{adv} \dicDirectTranslationCS{citově, emocionálně}
\dicEntry[tilfinningalegur] \textls[15]{\dicTerm{til·finninga··legur} \dicIPA{{t\smash{\textsuperscript{h}}}{\textsci}{l}{f}{\textsci}{n}{i}{\ng}{\r{g}}{a}{l}{\textepsilon}{\textbabygamma}{\textscy}{\textsubring{r}}} \dicPos{adj}[1]\dicFlx{}[-8] \textbf{1.} \dicSynonym*{kenndrænn} \dicDirectTranslationCS{(po)citový, (jsoucí) podle citu} \dicExampleIS{tilfinningalegt mat} \dicExampleCS{ohodnocení podle citu} \textbf{2.} \dicSynonym*{sem varðar tilfinningar} \dicDirectTranslationCS{citový, emocionální} \dicExampleIS{tilfinningalegur þroski} \dicExampleCS{citové dospívání}}
\dicEntry[tilfinningaleysi] \dicTerm{til·finninga··leysi} \dicIPA{{t\smash{\textsuperscript{h}}}{\textsci}{l}{f}{\textsci}{n}{i}{\ng}{\r{g}}{a}{l}{ei}{s}{\textsci}} \dicPos{n}[2] \dicFlx{(‑s)}[20] \textbf{1.} \dicSynonym*{skynjunarleysi} \dicDirectTranslationCS{necitlivost, znecitlivělost}  \textbf{2.} \dicSynonym*{kaldlyndi} \dicDirectTranslationCS{bezcitnost, necitelnost}
\dicEntry[tilfinningalíf] \dicTerm{til·finninga··líf} \dicIPA{{t\smash{\textsuperscript{h}}}{\textsci}{l}{f}{\textsci}{n}{i}{\ng}{\r{g}}{a}{l}{i}{f}} \dicPos{n}[2] \dicFlx{(‑s)}[2] \dicDirectTranslationCS{citový život}
\begin{xtolerant}{}{1pt}
\dicEntry[tilfinningamaður] \dicTerm{til·finninga··|maður}\addthinS\textls[-10]{\dicIPA{{t\smash{\textsuperscript{h}}}{\textsci}{l}{f}{\textsci}{n}{i}{\ng}{\r{g}}{a}{m}{a}{ð}{\textscy}{\textsubring{r}}}} \dicPos{m}[13]\addthinS\dicFlx{(‑manns, ‑menn)}[2] \dicDirectTranslationCS{citlivý\,/\addthin vnímavý člověk}
\end{xtolerant}
\dicEntry[tilfinninganæmi] \dicTerm{til·finninga··næmi} \dicIPA{{t\smash{\textsuperscript{h}}}{\textsci}{l}{f}{\textsci}{n}{i}{\ng}{\r{g}}{a}{n}{a}{i}{m}{\textsci}} \dicPos{n}[2] \dicFlx{(‑s)}[20] \dicDirectTranslationCS{citlivost, vnímavost, senzitivita}
\dicEntry[tilfinninganæmur] \dicTerm{til·finninga··næmur} \dicIPA{{t\smash{\textsuperscript{h}}}{\textsci}{l}{f}{\textsci}{n}{i}{\ng}{\r{g}}{a}{n}{a}{i}{m}{\textscy}{\textsubring{r}}} \dicPos{adj}[1]\dicFlx{}[-1] \dicDirectTranslationCS{citlivý, vnímavý, senzitivní}
\begin{xtolerant}{}{1pt}
\dicEntry[tilfinningasamur] \dicTerm{til·finninga··|samur}\addthinS\dicIPA{{t\smash{\textsuperscript{h}}}{\textsci}{l}{f}{\textsci}{n}{i}{\ng}{\r{g}}{a}{s}{a}{m}{\textscy}{\textsubring{r}}}\addthinS\dicPos{adj}[1]\addthinS\dicFlx{(f ‑söm)}[2]\addthinS\dicSynonym{viðkvæmur} \dicDirectTranslationCS{citový, sentimentální} \dicExampleIS{tilfinningasöm ljóð} \dicExampleCS{sentimentální básně}
\end{xtolerant}
\dicEntry[tilfinningasemi] \dicTerm{til·finninga··sem|i} \dicIPA{{t\smash{\textsuperscript{h}}}{\textsci}{l}{f}{\textsci}{n}{i}{\ng}{\r{g}}{a}{s}{\textepsilon}{m}{\textsci}} \dicPos{f}[2] \dicFlx{(‑i)}[2] \dicDirectTranslationCS{citovost, sentiment}
\dicEntry[tilfinningatengsl] \textls[15]{\dicTerm{til·finninga··tengsl} \dicIPA{{t\smash{\textsuperscript{h}}}{\textsci}{l}{f}{\textsci}{n}{i}{\ng}{\r{g}}{a}{t\smash{\textsuperscript{h}}}{\textepsilon}{\ng}{s}{\textsubring{d}}{\textsubring{l}}} \dicPos{n}[2] \dicFlx{pl}[1] \dicDirectTranslationCS{citová vazba}}
\dicEntry[tilflutningur] \dicTerm{til··flut·ning|ur} \dicIPA{{t\smash{\textsuperscript{h}}}{\textsci}{l}{f}{l}{\textscy}{h}{\textsubring{d}}{n}{i}{\ng}{\r{g}}{\textscy}{\textsubring{r}}} \dicPos{m}[6] \dicFlx{(‑s)}[9] \dicSynonym{flutningur} \dicDirectTranslationCS{přesun, přenos, přemístění} \dicExampleIS{tilflutningur á orku} \dicExampleCS{přesun energie}
\dicEntry[tilfæra] \dicTerm{til··fær|a} \dicIPA{{t\smash{\textsuperscript{h}}}{\textsci}{l}{f}{a}{i}{r}{a}} \dicPos{v}[2] \dicFlx{(‑ði, ‑t)}[99] \dicFlx{acc} \dicSynonym{tilgreina} \dicDirectTranslationCS{uvést, zmínit} \dicExampleIS{tilfæra dæmi} \dicExampleCS{uvést příklad}
\dicEntry[tilfærsla] \dicTerm{til··færsl|a} \dicIPA{{t\smash{\textsuperscript{h}}}{\textsci}{l}{f}{a}{i}{\textsubring{r}}{s}{\textsubring{d}}{l}{a}} \dicPos{f}[1] \dicFlx{(‑u, ‑ur)}[13] \dicDirectTranslationCS{přemístění, přesun, přeskupení} \dicExampleIS{tilfærsla á fjármagni} \dicExampleCS{přesun kapitálu}
\dicEntry[tilgangsfall] \dicTerm{til·gangs··|fall} \dicIPA{{t\smash{\textsuperscript{h}}}{\textsci}{l}{\r{g}}{au}{\ng}{s}{f}{a}{\textsubring{d}}{\textsubring{l}}} \dicPos{n}[2] \dicFlx{(‑falls, ‑föll)}[8] \dicFieldCat{jaz.} \dicDirectTranslationCS{allativ}
\dicEntry[tilgangslaus] \dicTerm{til·gangs··laus} \dicIPA{{t\smash{\textsuperscript{h}}}{\textsci}{l}{\r{g}}{au}{\ng}{s}{l}{\oe i}{s}} \dicPos{adj}[5]\dicFlx{}[-1] \dicSynonym{gagnslaus} \dicDirectTranslationCS{zbytečný, bezúčelný} \dicExampleIS{tilgangslaus leit} \dicExampleCS{bezúčelné pátrání}
\dicEntry[tilgangsleysi] \dicTerm{til·gangs··leysi} \dicIPA{{t\smash{\textsuperscript{h}}}{\textsci}{l}{\r{g}}{au}{\ng}{s}{l}{ei}{s}{\textsci}} \dicPos{n}[2] \dicFlx{(‑s)}[20] \dicDirectTranslationCS{zbytečnost, bezúčelnost}
\dicEntry[tilgangur] \dicTerm{til··gang|ur} \dicsymFrequent\  \dicIPA{{t\smash{\textsuperscript{h}}}{\textsci}{l}{\r{g}}{au}{\ng}{\r{g}}{\textscy}{\textsubring{r}}} \dicPos{m}[6] \dicFlx{(‑s)}[7] \dicSynonym{stefnumark} \dicDirectTranslationCS{účel, cíl, smysl} \dicExampleIS{gera e‑ð í góðum tilgangi} \dicExampleCS{dělat (co) za dobrým účelem};  \dicProverb\  \dicPhraseIS{Tilgangurinn helgar meðalið.} \dicLangCat{přís.} \dicDirectTranslationCS{Účel světí prostředky.}
\dicEntry[tilgáta] \dicTerm{til··gát|a} \dicIPA{{t\smash{\textsuperscript{h}}}{\textsci}{l}{\r{g}}{au}{\textsubring{d}}{a}} \dicPos{f}[1] \dicFlx{(‑u, ‑ur)}[19] \dicSynonym{ágiskun} \dicDirectTranslationCS{domněnka, dohad, hypotéza} \dicExampleIS{sanna tilgátuna} \dicExampleCS{potvrdit domněnku}
\dicEntry[tilgerð] \dicTerm{til··gerð} \dicIPA{{t\smash{\textsuperscript{h}}}{\textsci}{l}{\r{\textObardotlessj}}{\textepsilon}{r}{\texttheta}} \dicPos{f}[7] \dicFlx{(‑ar)}[3] \dicSynonym{fordild} \dicDirectTranslationCS{přetvářka, předstírání, faleš}
\dicEntry[tilgerðarlaus] \dicTerm{til·gerðar··laus} \dicIPA{{t\smash{\textsuperscript{h}}}{\textsci}{l}{\r{\textObardotlessj}}{\textepsilon}{r}{ð}{a}{r}{l}{\oe i}{s}} \dicPos{adj}[5]\dicFlx{}[-1] \dicDirectTranslationCS{prostý, skromný}
\dicEntry[tilgerðarlega] \dicTerm{til·gerðar··lega} \dicIPA{{t\smash{\textsuperscript{h}}}{\textsci}{l}{\r{\textObardotlessj}}{\textepsilon}{r}{ð}{a}{r}{l}{\textepsilon}{\textbabygamma}{a}} \dicPos{adv} \dicDirectTranslationCS{předstíraně, hraně}
\begin{xtolerant}{}{1pt}
\dicEntry[tilgerðarlegur] \dicTerm{til·gerðar··legur}\addthinS\dicIPA{{t\smash{\textsuperscript{h}}}{\textsci}{l}{\r{\textObardotlessj}}{\textepsilon}{r}{ð}{a}{r}{l}{\textepsilon}{\textbabygamma}{\textscy}{\textsubring{r}}}\addthinS\dicPos{adj}[1]\dicFlx{}[-8]\addthinS\dicDirectTranslationCS{předstíraný, hra\-ný, líčený} \dicExampleIS{tilgerðarlegt bros} \dicExampleCS{předstíraný úsměv}
\end{xtolerant}
\dicEntry[tilgreina] \dicTerm{til··grein|a} \dicIPA{{t\smash{\textsuperscript{h}}}{\textsci}{l}{\r{g}}{r}{ei}{n}{a}} \dicPos{v}[2] \dicFlx{(‑di, ‑t)}[141] \dicFlx{acc} \dicDirectTranslationCS{uvést, uvádět, udat, specifikovat} \dicExampleIS{tilgreina eiganda að e‑u} \dicExampleCS{uvést vlastníka (čeho)}
\dicEntry[tilhafður] \dicTerm{til··|hafður} \dicIPA{{t\smash{\textsuperscript{h}}}{\textsci}{\textlengthmark}{\textsubring{l}}{h}{a}{v}{ð}{\textscy}{\textsubring{r}}} \dicPos{adj}[2] \dicFlx{(f ‑höfð)}[3] \dicPhraseIS{vera vel\,/\addthin illa tilhafður} \dicDirectTranslationCS{být upravený\,/\addthin neupravený, být dobře\,/\addthin špatně oblečený}
\dicEntry[tilhald] \dicTerm{til··hald} \dicIPA{{t\smash{\textsuperscript{h}}}{\textsci}{\textlengthmark}{\textsubring{l}}{h}{a}{l}{\textsubring{d}}} \dicPos{n}[2] \dicFlx{(‑s)}[2] \textbf{1.} \dicDirectTranslationCS{paráda} \dicIndirectTranslationCS{(okázale pěkná vnější úprava)}  \textbf{2.} \dicSynonym{viðhöfn} \dicDirectTranslationCS{paráda} \dicIndirectTranslationCS{(okázalá slavnost)}
\dicEntry[tilhaldssamur] \dicTerm{til·halds··|samur} \dicIPA{{t\smash{\textsuperscript{h}}}{\textsci}{\textlengthmark}{\textsubring{l}}{h}{a}{l}{\textsubring{d}}{s}{a}{m}{\textscy}{\textsubring{r}}} \dicPos{adj}[1] \dicFlx{(f ‑söm)}[2] \dicDirectTranslationCS{parádní, okázalý} \dicIndirectTranslationCS{(určený pro parádu)}
\dicEntry[tilheyra] \dicTerm{til··heyr|a} \dicsymFrequent\  \dicIPA{{t\smash{\textsuperscript{h}}}{\textsci}{\textlengthmark}{\textsubring{l}}{h}{ei}{r}{a}} \dicPos{v}[2] \dicFlx{(‑ði, ‑t)}[111] \dicFlx{dat} \textbf{1.} \dicFlx{dat} \dicSynonym*{vera í eigu e‑rs} \dicDirectTranslationCS{patřit, náležet, příslušet} \dicExampleIS{lönd sem tilheyra Evrópu landfræðilega} \dicExampleCS{země, které patří geograficky k~Evropě}  \textbf{2.} \dicSynonym*{vera viðeigandi} \dicDirectTranslationCS{být vhodný\,/\addthin přiměřený, náležet}
\dicEntry[tilheyrandi] \dicTerm{til··heyr·|andi\smash{\textsuperscript{1}}} \dicIPA{{t\smash{\textsuperscript{h}}}{\textsci}{\textlengthmark}{\textsubring{l}}{h}{ei}{r}{a}{n}{\textsubring{d}}{\textsci}} \dicPos{m}[2] \dicFlx{(‑anda, ‑endur)}[1] \dicSynonym{hlustandi} \dicDirectTranslationCS{posluchač(ka)} \dicExampleIS{hafa fáa tilheyrendur} \dicExampleCS{mít málo posluchačů}
\dicEntry[tilheyrandi] \dicTerm{til··heyr·andi\smash{\textsuperscript{2}}} \dicIPA{{t\smash{\textsuperscript{h}}}{\textsci}{\textlengthmark}{\textsubring{l}}{h}{ei}{r}{a}{n}{\textsubring{d}}{\textsci}} \dicPos{adj}[13] \dicFlx{indecl}[1] \dicDirectTranslationCS{patřičný, příslušný, náležitý};  \dicPhraseIS{með öllu tilheyrandi} \dicFlx{adv} \dicDirectTranslationCS{se vším všudy, jak se patří}
\dicEntry[tilhlaup] \dicTerm{til··hlaup} \dicIPA{{t\smash{\textsuperscript{h}}}{\textsci}{l}{\textsubring{l}}{\oe i}{\textsubring{b}}} \dicPos{n}[2] \dicFlx{(‑s, ‑)}[5] \dicDirectTranslationCS{rozběh} \dicIndirectTranslationCS{(začátek běhu)}
\dicEntry[tilhliðrun] \dicTerm{til··hliðr|un} \dicIPA{{t\smash{\textsuperscript{h}}}{\textsci}{l}{\textsubring{l}}{\textsci}{ð}{r}{\textscy}{\textsubring{n}}} \dicPos{f}[7] \dicFlx{(‑unar)}[9] \textbf{1.} \dicSynonym{eftirgjöf} \dicDirectTranslationCS{ústupek, kompromis}  \textbf{2.} \dicSynonym*{lipurleiki} \dicDirectTranslationCS{shovívavost, mírnost, vstřícnost}
\dicEntry[tilhliðrunarsemi] \dicTerm{til·hliðrunar··sem|i} \dicIPA{{t\smash{\textsuperscript{h}}}{\textsci}{l}{\textsubring{l}}{\textsci}{ð}{r}{\textscy}{n}{a}{\textsubring{r}}{s}{\textepsilon}{m}{\textsci}} \dicPos{f}[2] \dicFlx{(‑i)}[2] \dicDirectTranslationCS{shovívavost, ochota ke kompromisu}
\dicEntry[tilhlutan] \dicTerm{til··hlutan} \dicIPA{{t\smash{\textsuperscript{h}}}{\textsci}{l}{\textsubring{l}}{\textscy}{\textsubring{d}}{a}{\textsubring{n}}} \dicPos{f}[4] \dicFlx{(‑ar)}[3] \dicLink{tilhlutun}
\dicEntry[tilhlutun] \dicTerm{til··hlut|un}\dicTerm{, tilhlutan} \dicIPA{{t\smash{\textsuperscript{h}}}\-{\textsci}\-{l}\-{\textsubring{l}}\-{\textscy}\-{\textsubring{d}}\-{\textscy}\-{\textsubring{n}}\-} \dicPos{f}[7] \dicFlx{(‑unar, ‑anir)}[8] \dicDirectTranslationCS{přispění, prostřednictví};  \dicPhraseIS{að tilhlutan e‑rs} \dicDirectTranslationCS{za přispění (koho)} \dicExampleIS{Samningar tókust að tilhlutan stjórnarinnar.} \dicExampleCS{Dohody se podařily za přispění vlády.}
\dicEntry[tilhlýðilegur] \dicTerm{til··hlýði·legur} \dicIPA{{t\smash{\textsuperscript{h}}}{\textsci}{l}{\textsubring{l}}{i}{ð}{\textsci}{l}{\textepsilon}{\textbabygamma}{\textscy}{\textsubring{r}}} \dicPos{adj}[1]\dicFlx{}[-8] \dicSynonym{hæfilegur} \dicDirectTranslationCS{vhodný, patřičný, příhodný}
\dicEntry[tilhlökkun] \dicTerm{til··hlökk|un} \dicIPA{{t\smash{\textsuperscript{h}}}{\textsci}{l}{\textsubring{l}}{\oe}{h}{\r{g}}{\textscy}{\textsubring{n}}} \dicPos{f}[7] \dicFlx{(‑unar)}[12] \dicDirectTranslationCS{očekávání, těšení se} \dicExampleIS{gleði og tilhlökkun} \dicExampleCS{radost a~očekávání}
\dicEntry[tilhneiging] \dicTerm{til··hneig·ing} \dicsymFrequent\  \dicIPA{{t\smash{\textsuperscript{h}}}{\textsci}{l}{\textsubring{n}}{ei}{j}{i}{\ng}{\r{g}}} \dicPos{f}[4] \dicFlx{(‑ar, ‑ar)}[5] \dicDirectTranslationCS{tendence, trend, sklon, tíhnutí, inklinace} \dicExampleIS{tilhneiging til einangrunar} \dicExampleCS{sklon k~izolovanosti};  \dicPhraseIS{hafa tilhneigingu til að (gera e‑ð)} \dicDirectTranslationCS{mít tendenci (dělat (co))}
\dicEntry[tilhugalíf] \dicTerm{til·huga··líf} \dicIPA{{t\smash{\textsuperscript{h}}}{\textsci}{\textlengthmark}{\textsubring{l}}{h}{\textscy}{\textbabygamma}{a}{l}{i}{f}} \dicPos{n}[2] \dicFlx{(‑s)}[2] \dicDirectTranslationCS{námluvy, dvoření}
\dicEntry[tilhugsun] \dicTerm{til··hugs|un} \dicsymFrequent\  \dicIPA{{t\smash{\textsuperscript{h}}}{\textsci}{\textlengthmark}{\textsubring{l}}{h}{\textscy}{x}{s}{\textscy}{\textsubring{n}}} \dicPos{f}[7] \dicFlx{(‑unar)}[9] \dicDirectTranslationCS{pomyšlení, myšlenka, představa};  \dicPhraseIS{e‑n hryllir við tilhugsuninni um e‑ð} \dicFlx{impers} \dicDirectTranslationCS{(kdo) se třese při pomyšlení na (co)}
\dicEntry[tilhæfulaus] \dicTerm{til·hæfu··laus} \dicIPA{{t\smash{\textsuperscript{h}}}{\textsci}{\textlengthmark}{\textsubring{l}}{h}{a}{i}{v}{\textscy}{l}{\oe i}{s}} \dicPos{adj}[5]\dicFlx{}[-1] \dicDirectTranslationCS{nepodložený, neodůvodněný} \dicExampleIS{tilhæfulausar grunsemdir} \dicExampleCS{nepodložená podezření}
\dicEntry[tilhögun] \dicTerm{til··hög|un} \dicIPA{{t\smash{\textsuperscript{h}}}{\textsci}{\textlengthmark}{\textsubring{l}}{h}{\oe}{\textbabygamma}{\textscy}{\textsubring{n}}} \dicPos{f}[7] \dicFlx{(‑unar)}[12] \dicSynonym{fyrirkomulag} \dicDirectTranslationCS{rozvržení, uspořádání} \dicExampleIS{Tilhögun kennslunnar hefur verið óbreytt um árabil.} \dicExampleCS{Rozvržení výuky se po mnoho let nezměnilo.}
\dicEntry[tilkeyra] \dicTerm{til··keyr|a} \dicIPA{{t\smash{\textsuperscript{h}}}{\textsci}{l}{c\smash{\textsuperscript{h}}}{ei}{r}{a}} \dicPos{v}[2] \dicFlx{(‑ði, ‑t)}[99] \dicFlx{acc} \dicDirectTranslationCS{zaběhnout, zajet (nový stroj ap.)} \dicExampleIS{tilkeyra bíl} \dicExampleCS{zaběhnout auto}
\dicEntry[tilkippilegur] \dicTerm{til··kippi·legur} \dicIPA{{t\smash{\textsuperscript{h}}}{\textsci}{l}{c\smash{\textsuperscript{h}}}{\textsci}{h}{\textsubring{b}}{\textsci}{l}{\textepsilon}{\textbabygamma}{\textscy}{\textsubring{r}}} \dicPos{adj}[1]\dicFlx{}[-8] \dicSynonym{reiðubúinn} \dicDirectTranslationCS{připravený, hotový, přichystaný} \dicExampleIS{Ég er ekki tilkippilegur til þess.} \dicExampleCS{Nejsem na to připravený.}
\dicEntry[tilkoma] \dicTerm{til··kom|a} \dicsymFrequent\  \dicIPA{{t\smash{\textsuperscript{h}}}{\textsci}{l}{k\smash{\textsuperscript{h}}}{\textopeno}{m}{a}} \dicPos{f}[1] \dicFlx{(‑u)}[5] \dicDirectTranslationCS{příchod, nástup, objevení se} \dicExampleIS{vænta tilkomu frelsarans} \dicExampleCS{očekávat příchod Spasitele}
\dicEntry[tilkomulaus] \dicTerm{til·komu··laus} \dicIPA{{t\smash{\textsuperscript{h}}}{\textsci}{l}{k\smash{\textsuperscript{h}}}{\textopeno}{m}{\textscy}{l}{\oe i}{s}} \dicPos{adj}[5]\dicFlx{}[-1] \dicLink{tilkomulítill}
\dicEntry[tilkomulítill] \dicTerm{til·komu··|lítill}\dicTerm{, tilkomulaus} \dicIPA{{t\smash{\textsuperscript{h}}}\-{\textsci}\-{l}\-{k\smash{\textsuperscript{h}}}\-{\textopeno}\-{m}\-{\textscy}\-{l}\-{i}\-{\textsubring{d}}\-{\textsci}\-{\textsubring{d}}\-{\textsubring{l}}\-} \dicPos{adj}[11] \dicFlx{(comp ‑minni, sup ‑minn\-st\-ur)}[10] \dicSynonym{óásjálegur} \dicDirectTranslationCS{nevýrazný, nepůsobivý}
\dicEntry[tilkomumikill] \dicTerm{til·komu··|mikill} \dicIPA{{t\smash{\textsuperscript{h}}}{\textsci}{l}{k\smash{\textsuperscript{h}}}{\textopeno}{m}{\textscy}{m}{\textsci}{\r{\textObardotlessj}}{\textsci}{\textsubring{d}}{\textsubring{l}}} \dicPos{adj}[11] \dicFlx{(comp ‑meiri, sup ‑mestur)}[12] \dicSynonym{myndarlegur} \dicDirectTranslationCS{působivý, úchvatný} \dicExampleIS{tilkomumikil útsjón} \dicExampleCS{úchvatný výhled}
\dicEntry[tilkostnaður] \dicTerm{til··kost·nað|ur} \dicIPA{{t\smash{\textsuperscript{h}}}{\textsci}{l}{k\smash{\textsuperscript{h}}}{\textopeno}{s}{\textsubring{d}}{n}{a}{ð}{\textscy}{\textsubring{r}}} \dicPos{m}[10] \dicFlx{(‑ar)}[9] \dicSynonym{kostnaður} \dicDirectTranslationCS{výdaje, vydání, výlohy}
\dicEntry[tilkynna] \dicTerm{til··kynn|a} \dicsymFrequent\  \dicIPA{{t\smash{\textsuperscript{h}}}{\textsci}{l}{c\smash{\textsuperscript{h}}}{\textsci}{n}{a}} \dicPos{v}[2] \dicFlx{(‑ti, ‑t)}[63] \dicFlx{dat + acc} \dicDirectTranslationCS{oznámit, oznamovat, sdělit, sdělovat, uvědomit, uvědomovat, ohlásit, ohlašovat};  \dicPhraseIS{tilkynna e‑m e‑ð} \dicDirectTranslationCS{oznámit (komu co)} \dicExampleIS{tilkynna honum þátttöku sína skriflega} \dicExampleCS{oznámit mu písemně svoji účast};  \dicIdiom{tilkynnast}{ \dicPhraseIS{tilkynnast}} \dicFlx{refl} \dicDirectTranslationCS{oznámit se, sdělit se, být oznámen\,/\addthin sdělen}
\dicEntry[tilkynning] \dicTerm{til··kynn·ing} \dicIPA{{t\smash{\textsuperscript{h}}}{\textsci}{l}{c\smash{\textsuperscript{h}}}{\textsci}{n}{i}{\ng}{\r{g}}} \dicPos{f}[4] \dicFlx{(‑ar, ‑ar)}[5] \dicSynonym{yfirlýsing} \dicDirectTranslationCS{oznámení, sdělení, prohlášení, ohlášení, uvědomění} \dicExampleIS{senda tilkynningu til almennings} \dicExampleCS{vyslat oznámení veřejnosti}
\dicEntry[till.] \dicTerm{till.} \dicPos{zkr} \dicPhraseIS{tillaga} \dicDirectTranslationCS{návrh}
\dicEntry[tillag] \dicTerm{til··|lag} \dicIPA{{t\smash{\textsuperscript{h}}}{\textsci}{l}{\textlengthmark}{a}{x}} \dicPos{n}[2] \dicFlx{(‑lags, ‑lög)}[8] \textbf{1.} \dicSynonym*{árgjald í félagi} \dicDirectTranslationCS{příspěvek (členský ap.)} \dicExampleIS{tillag til félagsins} \dicExampleCS{příspěvek organizaci}  \textbf{2.} \dicSynonym{framlag} \dicDirectTranslationCS{přispění, podpora}
\dicEntry[tillaga] \dicTerm{til··|laga} \dicsymFrequent\  \dicIPA{{t\smash{\textsuperscript{h}}}{\textsci}{l}{\textlengthmark}{a}{\textbabygamma}{a}} \dicPos{f}[1] \dicFlx{(‑lögu, ‑lögur)}[20] \dicSynonym{uppástunga} \dicDirectTranslationCS{(předložený) návrh, propozice} \textls[15]{\dicExampleIS{tillaga um lagabreytingu} \dicExampleCS{návrh na změnu zákona};  \dicPhraseIS{bera fram\,/\addthin leggja fram\,/\addthin flytja tillögu um e‑ð} \dicDirectTranslationCS{předložit návrh na (co), navrhnout\,/\addthin navrhovat (co)}}
\dicEntry[tilleiðanlegur] \dicTerm{til·leiðan··legur} \dicIPA{{t\smash{\textsuperscript{h}}}{\textsci}{l}{\textlengthmark}{ei}{ð}{a}{n}{l}{\textepsilon}{\textbabygamma}{\textscy}{\textsubring{r}}} \dicPos{adj}[1]\dicFlx{}[-8] \dicDirectTranslationCS{přístupný, otevřený, nakloněný (návrhu ap.)} \dicExampleIS{vera tilleiðanlegur til e‑s} \dicExampleCS{být nakloněný (čemu)}
\dicEntry[tilleiðast] \dicTerm{til··lei|ðast} \dicIPA{{t\smash{\textsuperscript{h}}}{\textsci}{l}{\textlengthmark}{ei}{ð}{a}{s}{\textsubring{d}}} \dicPos{v}[2] \dicFlx{(‑ddist, ‑ðst)}[212] \dicFlx{refl} \dicPhraseIS{láta tilleiðast} \dicDirectTranslationCS{nechat se přesvědčit\,/\addthin zviklat}
\dicEntry[tilli] \dicTerm{till|i} \dicIPA{{t\smash{\textsuperscript{h}}}{\textsci}{l}{\textlengthmark}{\textsci}} \dicPos{m}[1] \dicFlx{(‑a, ‑ar)}[1] \dicLangCat{dět.} \dicSynonym{typpi} \dicDirectTranslationCS{frantík, šulínek} \dicIndirectTranslationCS{(penis)}
\dicEntry[tillit] \dicTerm{til··lit} \dicsymFrequent\  \dicIPA{{t\smash{\textsuperscript{h}}}{\textsci}{l}{\textlengthmark}{\textsci}{\textsubring{d}}} \dicPos{n}[2] \dicFlx{(‑s)}[2] \textbf{1.} \dicSynonym{viðmið} \dicDirectTranslationCS{ohled, zřetel, přihlédnutí};  \dicPhraseIS{án tillits til e‑s} \dicDirectTranslationCS{bez ohledu na (co)};  \dicPhraseIS{í (þessu) tilliti} \dicFlx{adv} \dicDirectTranslationCS{v~(tomto) ohledu};  \dicPhraseIS{með tilliti til e‑s} \dicDirectTranslationCS{s~přihlédnutím k~(čemu)};  \dicPhraseIS{taka tillit til e‑rs} \dicDirectTranslationCS{brát ohled\,/\addthin zřetel na (koho)}  \textbf{2.} \dicSynonym{augnaráð} \dicDirectTranslationCS{pohled, podívání}
\dicEntry[tillitslaus] \dicTerm{til·lits··laus} \dicIPA{{t\smash{\textsuperscript{h}}}{\textsci}{l}{\textlengthmark}{\textsci}{\textsubring{d}}{s}{l}{\oe i}{s}} \dicPos{adj}[5]\dicFlx{}[-1] \dicDirectTranslationCS{neohleduplný, necitelný, neberoucí ohled\,/\addthin zřetel}
\dicEntry[tillitsleysi] \dicTerm{til·lits··leysi} \dicIPA{{t\smash{\textsuperscript{h}}}{\textsci}{l}{\textlengthmark}{\textsci}{\textsubring{d}}{s}{l}{ei}{s}{\textsci}} \dicPos{n}[2] \dicFlx{(‑s)}[20] \dicDirectTranslationCS{neohleduplnost, necitelnost}
\dicEntry[tillitssamur] \dicTerm{til·lits··|samur} \dicIPA{{t\smash{\textsuperscript{h}}}{\textsci}{l}{\textlengthmark}{\textsci}{\textsubring{d}}{s}{a}{m}{\textscy}{\textsubring{r}}} \dicPos{adj}[1] \dicFlx{(f ‑söm)}[2] \dicSynonym{nærgætinn} \dicDirectTranslationCS{ohleduplný, taktní, šetrný, diskrétní}
\dicEntry[tillitssemi] \dicTerm{til·lits··sem|i} \dicIPA{{t\smash{\textsuperscript{h}}}{\textsci}{l}{\textlengthmark}{\textsci}{\textsubring{d}}{s}{\textepsilon}{m}{\textsci}} \dicPos{f}[2] \dicFlx{(‑i)}[2] \dicDirectTranslationCS{ohleduplnost, takt, šetrnost, diskrétnost}
\dicEntry[tilmæli] \dicTerm{til··mæli} \dicIPA{{t\smash{\textsuperscript{h}}}{\textsci}{l}{m}{a}{i}{l}{\textsci}} \dicPos{n}[2] \dicFlx{pl}[19] \dicSynonym{beiðni} \dicDirectTranslationCS{žádost, požadavek}
\dicEntry[tilnefna] \dicTerm{til··nefn|a} \dicsymFrequent\  \dicIPA{{t\smash{\textsuperscript{h}}}{\textsci}{l}{n}{\textepsilon}{\textsubring{b}}{n}{a}} \dicPos{v}[2] \dicFlx{(‑di, ‑t)}[141] \dicFlx{acc} \textbf{1.} \dicSynonym*{nefna til} \dicDirectTranslationCS{jmenovat, nominovat (do funkce ap.)} \dicExampleIS{tilnefna hana í stjórn} \dicExampleCS{nominovat ji do vlády}  \textbf{2.} \dicDirectTranslationCS{nominovat} \dicIndirectTranslationCS{(zařadit do soutěže)}
\dicEntry[tilnefning] \dicTerm{til··nefn·ing} \dicIPA{{t\smash{\textsuperscript{h}}}{\textsci}{l}{n}{\textepsilon}{\textsubring{b}}{n}{i}{\ng}{\r{g}}} \dicPos{f}[4] \dicFlx{(‑ar)}[7] \textbf{1.} \dicDirectTranslationCS{jmenování, nominace, nominování (do funkce ap.)}  \textbf{2.} \dicDirectTranslationCS{nominace, nominování} \dicIndirectTranslationCS{(zařazení do soutěže)}
\dicEntry[tilneyddur] \dicTerm{til··neyddur} \dicIPA{{t\smash{\textsuperscript{h}}}{\textsci}{l}{n}{ei}{\textsubring{d}}{\textscy}{\textsubring{r}}} \dicPos{adj}[2]\dicFlx{}[-21] \dicDirectTranslationCS{donucený, přinucený}
\dicEntry[tilraun] \dicTerm{til··raun} \dicsymFrequent\  \dicIPA{{t\smash{\textsuperscript{h}}}{\textsci}{l}{r}{\oe i}{\textsubring{n}}} \dicPos{f}[7] \dicFlx{(‑ar, ‑ir)}[1] \textbf{1.} \dicSynonym*{það að prófa} \dicDirectTranslationCS{pokus, zkouška, test}  \textbf{2.} \dicSynonym*{verkleg rannsókn} \dicDirectTranslationCS{experiment, pokus} \dicExampleIS{efnafræðilegar tilraunir} \dicExampleCS{chemické pokusy}
\dicEntry[tilraunadýr] \dicTerm{til·rauna··dýr} \dicIPA{{t\smash{\textsuperscript{h}}}{\textsci}{l}{r}{\oe i}{n}{a}{\textsubring{d}}{i}{\textsubring{r}}} \dicPos{n}[2] \dicFlx{(‑s, ‑)}[5] \dicDirectTranslationCS{pokusné zvíře}
\dicEntry[tilraunaglas] \dicTerm{til·rauna··|glas} \dicIPA{{t\smash{\textsuperscript{h}}}{\textsci}{l}{r}{\oe i}{n}{a}{\r{g}}{l}{a}{s}} \dicPos{n}[2] \dicFlx{(‑glass, ‑glös)}[8] \dicDirectTranslationCS{zkumavka}
\dicEntry[tilraunastig] \dicTerm{til·rauna··stig} \dicIPA{{t\smash{\textsuperscript{h}}}{\textsci}{l}{r}{\oe i}{n}{a}{s}{\textsubring{d}}{\textsci}{x}} \dicPos{n}[2] \dicFlx{(‑s)}[2] \dicPhraseIS{e‑að er á tilraunastigi} \dicDirectTranslationCS{(co) je ve zkušebním stádiu}
\dicEntry[tilraunastofa] \dicTerm{til·rauna··stof|a} \dicIPA{{t\smash{\textsuperscript{h}}}{\textsci}{l}{r}{\oe i}{n}{a}{s}{\textsubring{d}}{\textopeno}{v}{a}} \dicPos{f}[1] \dicFlx{(‑u, ‑ur)}[7] \dicDirectTranslationCS{laboratoř}
\dicEntry[tilreykja] \dicTerm{til··reyk|ja} \dicIPA{{t\smash{\textsuperscript{h}}}{\textsci}{l}{r}{ei}{\r{\textObardotlessj}}{a}} \dicPos{v}[2] \dicFlx{(‑ti, ‑t)}[32] \dicFlx{acc} \dicDirectTranslationCS{rozpálit, rozpalovat} \dicExampleIS{tilreykja pípuna} \dicExampleCS{rozpálit dýmku}
\dicEntry[tilræði] \dicTerm{til··ræði} \dicIPA{{t\smash{\textsuperscript{h}}}{\textsci}{l}{r}{a}{i}{ð}{\textsci}} \dicPos{n}[2] \dicFlx{(‑s, ‑)}[14] \textbf{1.} \dicSynonym{árás} \dicDirectTranslationCS{atentát, vražedný útok};  \dicPhraseIS{tilræði við e‑n} \dicDirectTranslationCS{atentát na (koho)} \dicExampleIS{tilræði við konunginn} \dicExampleCS{atentát na krále}  \textbf{2.} \dicLangCat{přen.} \dicDirectTranslationCS{útok, atentát (na svobodu slova ap.)}
\dicEntry[tilræðismaður] \dicTerm{til·ræðis··|maður} \dicIPA{{t\smash{\textsuperscript{h}}}{\textsci}{l}{r}{a}{i}{ð}{\textsci}{s}{m}{a}{ð}{\textscy}{\textsubring{r}}} \dicPos{m}[13] \dicFlx{(‑manns, ‑menn)}[2] \dicDirectTranslationCS{atentátník, atentátnice}
\dicEntry[tilsettur] \dicTerm{til··settur} \dicIPA{{t\smash{\textsuperscript{h}}}{\textsci}{l}{s}{\textepsilon}{h}{\textsubring{d}}{\textscy}{\textsubring{r}}} \dicPos{adj}[1]\dicFlx{}[-10] \dicDirectTranslationCS{stanovený, ujednaný} \dicExampleIS{á tilsettum tíma} \dicExampleCS{ve stanovený čas}
\dicEntry[tilsjón] \dicTerm{til··sjón} \dicIPA{{t\smash{\textsuperscript{h}}}{\textsci}{l}{s}{j}{ou}{\textsubring{n}}} \dicPos{f}[7] \dicFlx{(‑ar)}[3] \dicDirectTranslationCS{dohled, dozor};  \dicPhraseIS{hafa tilsjón með e‑u} \dicDirectTranslationCS{dohlížet na (co)}
\dicEntry[tilskilja] \dicTerm{til··skil|ja} \dicIPA{{t\smash{\textsuperscript{h}}}{\textsci}{l}{s}{\r{\textObardotlessj}}{\textsci}{l}{j}{a}} \dicPos{v}[4] \dicFlx{(‑di, ‑ið)}[29] \dicFlx{acc} \dicSynonym{áskilja} \dicDirectTranslationCS{zajistit si, vyhradit si, vymínit si} \dicExampleIS{tilskilja ákveðin tímamörk} \dicExampleCS{vyhradit si určité časové hranice}
\dicEntry[tilskipun] \dicTerm{til··skip|un} \dicIPA{{t\smash{\textsuperscript{h}}}{\textsci}{l}{s}{\r{\textObardotlessj}}{\textsci}{\textsubring{b}}{\textscy}{\textsubring{n}}} \dicPos{f}[7] \dicFlx{(‑unar, ‑anir)}[8] \dicSynonym{ákvæði\smash{\textsuperscript{1}}} \dicDirectTranslationCS{nařízení, ustanovení, výnos, směrnice, dekret} \dicExampleIS{tilskipun um veiðibann} \dicExampleCS{nařízení o~zákazu lovu}
\dicEntry[tilslökun] \dicTerm{til··slök|un} \dicIPA{{t\smash{\textsuperscript{h}}}{\textsci}{l}{s}{\textsubring{d}}{l}{\oe}{\r{g}}{\textscy}{\textsubring{n}}} \dicPos{f}[7] \dicFlx{(‑unar)}[12] \dicDirectTranslationCS{kompromis, ústupek}
\dicEntry[tilstilli] \dicTerm{til··stilli} \dicIPA{{t\smash{\textsuperscript{h}}}{\textsci}{l}{s}{\textsubring{d}}{\textsci}{\textsubring{d}}{l}{\textsci}} \dicPos{n}[2] \dicFlx{(‑s)}[20] \dicSynonym{atbeini} \dicDirectTranslationCS{prostřednictví, přispění, pomoc};  \dicPhraseIS{fyrir\,/\addthin með tilstilli e‑rs} \dicDirectTranslationCS{přispěním (koho)} \dicExampleIS{Þetta tókst með tilstilli hans.} \dicExampleCS{Podařilo se to díky jeho pomoci.}
\begin{xtolerant}{}{1pt}
\dicEntry[tilstuðlan] \dicTerm{til··stuðlan}\dicTerm{,\addthinM tilstuðlun}\addthinS\dicIPA{{t\smash{\textsuperscript{h}}}\-{\textsci}\-{l}\-{s}\-{\textsubring{d}}\-{\textscy}\-{ð}\-{l}\-{a}\-{\textsubring{n}}\-}\addthinS\dicPos{f}[4]\addthinS\dicFlx{(‑ar)}[3]\addthinS\dicSynonym{tilstilli} \dicDirectTranslationCS{podpora, asistence, přispění};  \dicPhraseIS{að\,/\addthin fyrir tilstuðlan e‑rs} \dicDirectTranslationCS{s~přispěním (koho)} \dicExampleIS{að tilstuðlan ráðuneytisins} \dicExampleCS{s~přispěním ministerstva}
\end{xtolerant}
\dicEntry[tilstuðlun] \dicTerm{til··stuðl|un} \dicIPA{{t\smash{\textsuperscript{h}}}{\textsci}{l}{s}{\textsubring{d}}{\textscy}{ð}{l}{\textscy}{\textsubring{n}}} \dicPos{f}[7] \dicFlx{(‑unar)}[9] \dicLink{tilstuðlan}
\dicEntry[tilstyrkur] \dicTerm{til··styrk|ur} \dicIPA{{t\smash{\textsuperscript{h}}}{\textsci}{l}{s}{\textsubring{d}}{\textsci}{\textsubring{r}}{\r{g}}{\textscy}{\textsubring{r}}} \dicPos{m}[9] \dicFlx{(‑s, ‑ir)}[15] \dicSynonym{aðstoð} \dicDirectTranslationCS{přispění, podpora};  \dicPhraseIS{með tilstyrk e‑rs\,/\addthin e‑s} \dicDirectTranslationCS{za přispění (koho\,/\addthin čeho)} \dicExampleIS{með tilstyrk ríkis og sveitarfélaga} \dicExampleCS{za přispění státu a~obcí}
\dicEntry[tilsvar] \dicTerm{til··|svar} \dicIPA{{t\smash{\textsuperscript{h}}}{\textsci}{l}{s}{v}{a}{\textsubring{r}}} \dicPos{n}[2] \dicFlx{(‑svars, ‑svör)}[8] \dicSynonym{svar} \dicDirectTranslationCS{odpověď, odpovídání};  \dicPhraseIS{vera (neyðarlegur) í tilsvörum} \dicDirectTranslationCS{být (rozpačitý) při odpovídání}
\dicEntry[tilsvarandi] \dicTerm{til··svar·andi} \dicIPA{{t\smash{\textsuperscript{h}}}{\textsci}{l}{s}{v}{a}{r}{a}{n}{\textsubring{d}}{\textsci}} \dicPos{adj}[13] \dicFlx{indecl}[1] \dicSynonym{samsvarandi} \dicDirectTranslationCS{odpovídající, adekvátní, korespondující}
\dicEntry[tilsögn] \dicTerm{til··|sögn} \dicIPA{{t\smash{\textsuperscript{h}}}{\textsci}{l}{s}{\oe}{\r{g}}{\textsubring{n}}} \dicPos{f}[7] \dicFlx{(‑sagnar)}[19] \dicDirectTranslationCS{instrukce, pokyny} \dicExampleIS{tilsögn í lestri} \dicExampleCS{pokyny ke čtení}
\dicEntry[tiltaka] \dicTerm{til··|taka} \dicIPA{{t\smash{\textsuperscript{h}}}{\textsci}{l}{t\smash{\textsuperscript{h}}}{a}{\r{g}}{a}} \dicPos{v}[6] \dicFlx{(‑tek, ‑tók, ‑tókum, ‑tæki, ‑tekið)}[127] \dicFlx{acc} \dicSynonym{tilgreina} \dicDirectTranslationCS{udat, uvést, uvádět, sdělit, sdělovat} \dicExampleIS{tiltaka stað og  stund} \dicExampleCS{udat místo a~čas}
\dicEntry[tiltala] \dicTerm{til··|tala} \dicIPA{{t\smash{\textsuperscript{h}}}{\textsci}{l}{t\smash{\textsuperscript{h}}}{a}{l}{a}} \dicPos{f}[1] \dicFlx{(‑tölu)}[2] \dicSynonym{hlutfall} \dicDirectTranslationCS{úměra, proporce} \dicExampleIS{að tiltölu við e‑ð} \dicExampleCS{úměrně k~(čemu), proporcionálně k~(čemu)}
\dicEntry[tiltek] \dicTerm{til··tek} \dicIPA{{t\smash{\textsuperscript{h}}}{\textsci}{l}{t\smash{\textsuperscript{h}}}{\textepsilon}{\r{g}}} \dicPos{v} \dicFlx{ind praes sg 1 pers} \dicLink{tiltaka}
\dicEntry[tiltekið] \dicTerm{til··tekið} \dicIPA{{t\smash{\textsuperscript{h}}}{\textsci}{l}{t\smash{\textsuperscript{h}}}{\textepsilon}{\r{\textObardotlessj}}{\textsci}{\texttheta}} \dicPos{v} \dicFlx{supin} \dicLink{tiltaka}
\dicEntry[tiltekinn] \dicTerm{til··tekinn} \dicsymFrequent\  \dicIPA{{t\smash{\textsuperscript{h}}}{\textsci}{l}{t\smash{\textsuperscript{h}}}{\textepsilon}{\r{\textObardotlessj}}{\textsci}{\textsubring{n}}} \dicPos{adj}[6]\dicFlx{}[-6] \textbf{1.} \dicDirectTranslationCS{ustálený, určený, stanovený} \dicExampleIS{á tilteknum tíma} \dicExampleCS{ve stanovený čas}  \textbf{2.} \dicDirectTranslationCS{konkrétní, určitý};  \dicPhraseIS{nánar tiltekið} \dicFlx{adv} \dicDirectTranslationCS{konkrétně, přesně}
\dicEntry[tiltekt] \dicTerm{til··tekt} \dicIPA{{t\smash{\textsuperscript{h}}}{\textsci}{l}{t\smash{\textsuperscript{h}}}{\textepsilon}{x}{\textsubring{d}}} \dicPos{f}[7] \dicFlx{(‑ar, ‑ir)}[1] \dicDirectTranslationCS{úklid, uklízení}
\dicEntry[tiltók] \dicTerm{til··tók} \dicIPA{{t\smash{\textsuperscript{h}}}{\textsci}{l}{t\smash{\textsuperscript{h}}}{ou}{\r{g}}} \dicPos{v} \dicFlx{ind pf sg 1 pers} \dicLink{tiltaka}
\dicEntry[tiltókum] \dicTerm{til··tókum} \dicIPA{{t\smash{\textsuperscript{h}}}{\textsci}{l}{t\smash{\textsuperscript{h}}}{ou}{\r{g}}{\textscy}{\textsubring{m}}} \dicPos{v} \dicFlx{ind pf pl 1 pers} \dicLink{tiltaka}
\dicEntry[tiltrú] \dicTerm{til··trú} \dicIPA{{t\smash{\textsuperscript{h}}}{\textsci}{l}{t\smash{\textsuperscript{h}}}{r}{u}} \dicPos{f}[4] \dicFlx{(‑ar)}[27] \dicSynonym{traust} \dicDirectTranslationCS{důvěra} \dicExampleIS{hafa fulla tiltrú til hennar} \dicExampleCS{mít k~ní plnou důvěru}
\dicEntry[tiltæki] \dicTerm{til··tæki\smash{\textsuperscript{1}}} \dicIPA{{t\smash{\textsuperscript{h}}}{\textsci}{l}{t\smash{\textsuperscript{h}}}{a}{i}{\r{\textObardotlessj}}{\textsci}} \dicPos{n}[2] \dicFlx{(‑s, ‑)}[16] \dicDirectTranslationCS{výstřelek, výstřednost} \dicExampleIS{Þetta er furðulegt tiltæki.} \dicExampleCS{To je podivný výstřelek.}
\dicEntry[tiltæki] \dicTerm{til··tæki\smash{\textsuperscript{2}}} \dicIPA{{t\smash{\textsuperscript{h}}}{\textsci}{l}{t\smash{\textsuperscript{h}}}{a}{i}{\r{\textObardotlessj}}{\textsci}} \dicPos{v} \dicFlx{con pf sg 1 pers} \dicLink{tiltaka}

\dicEntry[tiltækilegur] \dicTerm{til·tæki··legur} \textls[-5]{\dicIPA{{t\smash{\textsuperscript{h}}}{\textsci}{l}{t\smash{\textsuperscript{h}}}{a}{i}{\r{\textObardotlessj}}{\textsci}{l}{\textepsilon}{\textbabygamma}{\textscy}{\textsubring{r}}}} \dicPos{adj}[1]\dicFlx{}[-8] \dicSynonym{hentugur} \dicDirectTranslationCS{vhod\-ný, příhodný, žádoucí}

\dicEntry[tiltækur] \dicTerm{til··tækur} \dicsymFrequent\  \dicIPA{{t\smash{\textsuperscript{h}}}{\textsci}{l}{t\smash{\textsuperscript{h}}}{a}{i}{\r{g}}{\textscy}{\textsubring{r}}} \dicPos{adj}[1]\dicFlx{}[-1] \textbf{1.} \dicSynonym{tilbúinn} \dicDirectTranslationCS{(jsoucí) po ruce, (jsoucí) při ruce, připravený}  \textbf{2.} \dicSynonym{aðgengilegur} \dicDirectTranslationCS{dostupný, (jsoucí) k~dispozici} \dicExampleIS{Þau börðust með öllum þeim vopnum sem þeim voru tiltæk.} \dicExampleCS{Bojovali všemi zbraněmi, které měli po ruce.}
\dicEntry[tiltölulega] \dicTerm{til·tölu··lega} \dicsymFrequent\  \dicIPA{{t\smash{\textsuperscript{h}}}{\textsci}{l}{t\smash{\textsuperscript{h}}}{\oe}{l}{\textscy}{l}{\textepsilon}{\textbabygamma}{a}} \dicPos{adv} \dicSynonym*{hlutfallslega} \dicDirectTranslationCS{poměrně, relativně} \dicExampleIS{Það getur gerst á tiltölulega skömmum tíma.} \dicExampleCS{Může se to stát během relativně krátké doby.}
\dicEntry[tilurð] \dicTerm{til··urð} \dicIPA{{t\smash{\textsuperscript{h}}}{\textsci}{\textlengthmark}{l}{\textscy}{r}{\texttheta}} \dicPos{f}[7] \dicFlx{(‑ar)}[3] \dicDirectTranslationCS{vznik, zrod, stvoření, geneze} \dicExampleIS{gera grein fyrir tilurð bókarinnar} \dicExampleCS{vysvětlit vznik knihy}
\dicEntry[tilvalinn] \dicTerm{til··valinn} \dicIPA{{t\smash{\textsuperscript{h}}}{\textsci}{l}{v}{a}{l}{\textsci}{\textsubring{n}}} \dicPos{adj}[6]\dicFlx{}[-8] \dicSynonym{ágætur} \dicDirectTranslationCS{dokonalý, ideální} \dicExampleIS{Þetta er tilvalinn staður.} \dicExampleCS{Toto je dokonalé místo.}
\dicEntry[tilvera] \textls[15]{\dicTerm{til··ver|a} \dicsymFrequent\  \dicIPA{{t\smash{\textsuperscript{h}}}{\textsci}{l}{v}{\textepsilon}{r}{a}} \dicPos{f}[1] \dicFlx{(‑u)}[5] \dicDirectTranslationCS{existence, bytí} \dicExampleIS{vera fyllilega sáttur við lífið og tilveruna} \dicExampleCS{být úplně spokojený \mbox{se životem a~bytím}}}
\dicEntry[tilverknaður] \dicTerm{til··verk·nað|ur} \dicIPA{{t\smash{\textsuperscript{h}}}{\textsci}{l}{v}{\textepsilon}{\textsubring{r}}{\r{g}}{n}{a}{ð}{\textscy}{\textsubring{r}}} \dicPos{m}[10] \dicFlx{(‑ar, ‑ir)}[8] \dicSynonym{atbeini} \dicDirectTranslationCS{přičinění, iniciativa, přispění};  \dicPhraseIS{fyrir tilverknað e‑rs} \dicDirectTranslationCS{s~(čím) přičiněním}
\dicEntry[tilveruréttur] \dicTerm{til·veru··rétt|ur} \dicIPA{{t\smash{\textsuperscript{h}}}{\textsci}{l}{v}{\textepsilon}{r}{\textscy}{r}{j}{\textepsilon}{h}{\textsubring{d}}{\textscy}{\textsubring{r}}} \dicPos{m}[10] \dicFlx{(‑ar)}[7] \dicDirectTranslationCS{právo na existenci}
\dicEntry[tilvik] \dicTerm{til··vik} \dicsymFrequent\  \dicIPA{{t\smash{\textsuperscript{h}}}{\textsci}{l}{v}{\textsci}{\r{g}}} \dicPos{n}[2] \dicFlx{(‑s, ‑)}[5] \dicSynonym{tilfelli} \dicDirectTranslationCS{případ, okolnost, situace} \dicExampleIS{í þessu tilviki} \dicExampleCS{v~tomto případě}
\dicEntry[tilviljanakenndur] \dicTerm{til·viljana··kenndur}\textls[-10]{\dicTerm{, tilviljunarkenndur} \dicIPA{{t\smash{\textsuperscript{h}}}\-{\textsci}\-{l}\-{v}\-{\textsci}\-{l}\-{j}\-{a}\-{n}\-{a}\-{c\smash{\textsuperscript{h}}}\-{\textepsilon}\-{n}\-{\textsubring{d}}\-{\textscy}\-{\textsubring{r}}}} \dicPos{adj}[2]\dicFlx{}[-14] \dicDirectTranslationCS{náhodný, nahodilý}
\dicEntry[tilviljun] \dicTerm{til··vilj|un} \dicsymFrequent\  \dicIPA{{t\smash{\textsuperscript{h}}}{\textsci}{l}{v}{\textsci}{l}{j}{\textscy}{\textsubring{n}}} \dicPos{f}[7] \dicFlx{(‑unar, ‑anir)}[8] \dicDirectTranslationCS{náhoda, shoda náhod, souhra okolností};  \dicPhraseIS{af tilviljun} \dicFlx{adv} \dicDirectTranslationCS{náhodou, náhodně, namátkou} \dicExampleIS{hitta hana af tilviljun} \dicExampleCS{potkat ji náhodou}
\dicEntry[tilviljunarkenndur] \dicTerm{til·viljunar··kenndur} \dicIPA{{t\smash{\textsuperscript{h}}}\-{\textsci}\-{l}\-{v}\-{\textsci}\-{l}\-{j}\-{\textscy}\-{n}\-{a}\-{\textsubring{r}}\-{c\smash{\textsuperscript{h}}}\-{\textepsilon}\-{n}\-{\textsubring{d}}\-{\textscy}\-{\textsubring{r}}\-} \dicPos{adj}[2]\dicFlx{}[-14] \dicLink{tilviljanakenndur}
\dicEntry[tilvinnandi] \dicTerm{til··vinn·andi} \dicIPA{{t\smash{\textsuperscript{h}}}{\textsci}{l}{v}{\textsci}{n}{a}{n}{\textsubring{d}}{\textsci}} \dicPos{adj}[13] \dicFlx{indecl}[1] \dicDirectTranslationCS{stojící za to, vyplácející se} \dicExampleIS{Það er ekki tilvinnandi.} \dicExampleCS{Nestojí to za to.}
\dicEntry[tilvist] \dicTerm{til··vist} \dicIPA{{t\smash{\textsuperscript{h}}}{\textsci}{l}{v}{\textsci}{s}{\textsubring{d}}} \dicPos{f}[7] \dicFlx{(‑ar)}[3] \dicSynonym{tilvera} \dicDirectTranslationCS{existence} \dicExampleIS{efast um tilvist drauga} \dicExampleCS{pochybovat o~existenci duchů}
\dicEntry[tilvistarstefna] \dicTerm{til·vistar··stefn|a} \dicIPA{{t\smash{\textsuperscript{h}}}{\textsci}{l}{v}{\textsci}{s}{\textsubring{d}}{a}{\textsubring{r}}{s}{\textsubring{d}}{\textepsilon}{\textsubring{b}}{n}{a}} \dicPos{f}[1] \dicFlx{(‑u)}[5] \dicFieldCat{filos.} \dicDirectTranslationCS{existencionalismus}
\dicEntry[tilvitnun] \dicTerm{til··vitn|un} \dicIPA{{t\smash{\textsuperscript{h}}}{\textsci}{l}{v}{\textsci}{h}{\textsubring{d}}{n}{\textscy}{\textsubring{n}}} \dicPos{f}[7] \dicFlx{(‑unar, ‑anir)}[8] \dicDirectTranslationCS{citát, citace}
\dicEntry[tilvitnunarmerki] \dicTerm{til·vitnunar··merki} \dicIPA{{t\smash{\textsuperscript{h}}}{\textsci}{l}{v}{\textsci}{h}{\textsubring{d}}{n}{\textscy}{n}{a}{r}{m}{\textepsilon}{\textsubring{r}}{\r{\textObardotlessj}}{\textsci}} \dicPos{n}[2] \dicFlx{(‑s, ‑)}[16] \dicFieldCat{jaz.} \dicDirectTranslationCS{uvozovka} \dicIndirectTranslationCS{(interpunkční znaménko)}
\dicEntry[tilvísun] \dicTerm{til··vís|un} \dicIPA{{t\smash{\textsuperscript{h}}}{\textsci}{l}{v}{i}{s}{\textscy}{\textsubring{n}}} \dicPos{f}[7] \dicFlx{(‑unar, ‑anir)}[8] \textbf{1.} \dicSynonym*{það að vísa til} \dicDirectTranslationCS{odkaz, reference} \dicExampleIS{tilvísun til heimilda} \dicExampleCS{odkaz na použitou literaturu}  \textbf{2.} \dicSynonym{leiðbeining} \dicDirectTranslationCS{návod, instrukce, rada};  \dicPhraseIS{eftir\,/\addthin að tilvísun e‑rs} \dicDirectTranslationCS{podle instrukcí (koho)}  \textbf{3.} \dicFieldCat{med.} \dicIndirectTranslationCS{doporučení (obvodního lékaře) ke specialistovi}
\begin{xtolerant}{}{1pt}
\dicEntry[tilvísunarfornafn] \dicTerm{til·vísunar··for·|nafn}\addthinS\dicIPA{{t\smash{\textsuperscript{h}}}\-{\textsci}\-{l}\-{v}\-{i}\-{s}\-{\textscy}\-{n}\-{a}\-{\textsubring{r}}\-{f}\-{\textopeno}\-{r}\-{n}\-{a}\-{\textsubring{b}}\-{\textsubring{n}}}\addthinS\dicPos{n}[2]\addthinS\dicFlx{(‑nafns, ‑nöfn)}[8] \dicFieldCat{jaz.} \dicDirectTranslationCS{vztažné zájmeno}
\end{xtolerant}
\dicEntry[tilvísunarmerki] \dicTerm{til·vísunar··merki} \dicIPA{{t\smash{\textsuperscript{h}}}{\textsci}{l}{v}{i}{s}{\textscy}{n}{a}{r}{m}{\textepsilon}{\textsubring{r}}{\r{\textObardotlessj}}{\textsci}} \dicPos{n}[2] \dicFlx{(‑s, ‑)}[16] \dicFieldCat{jaz.} \dicDirectTranslationCS{uvozovka} \dicIndirectTranslationCS{(interpunkční znaménko)}
\dicEntry[tilvísunarsetning] \dicTerm{til·vísunar··set·ning} \dicIPA{{t\smash{\textsuperscript{h}}}\-{\textsci}\-{l}\-{v}\-{i}\-{s}\-{\textscy}\-{n}\-{a}\-{\textsubring{r}}\-{s}\-{\textepsilon}\-{h}\-{\textsubring{d}}\-{n}\-{i}\-{\ng}\-{\r{g}}\-} \dicPos{f}[4] \dicFlx{(‑ar, ‑ar)}[5] \dicFieldCat{jaz.} \dicDirectTranslationCS{vztažná věta}
\dicEntry[tilvísunartenging] \dicTerm{til·vísunar··teng·ing} \dicIPA{{t\smash{\textsuperscript{h}}}\-{\textsci}\-{l}\-{v}\-{i}\-{s}\-{\textscy}\-{n}\-{a}\-{\textsubring{r}}\-{t\smash{\textsuperscript{h}}}\-{ei}\-{\textltailn}\-{\r{\textObardotlessj}}\-{i}\-{\ng}\-{\r{g}}\-} \dicPos{f}[4] \dicFlx{(‑ar, ‑ar)}[5] \dicFieldCat{jaz.} \dicDirectTranslationCS{vztažná spojka}
\dicEntry[tilvonandi] \dicTerm{til··von·andi} \dicIPA{{t\smash{\textsuperscript{h}}}{\textsci}{l}{v}{\textopeno}{n}{a}{n}{\textsubring{d}}{\textsci}} \dicPos{adj}[13] \dicFlx{indecl}[1] \dicSynonym{væntanlegur} \dicDirectTranslationCS{nastávající, budoucí} \dicExampleIS{tilvonandi formaður} \dicExampleCS{nastávající předseda}
\dicEntry[tilþrif] \dicTerm{til··þrif} \dicIPA{{t\smash{\textsuperscript{h}}}{\textsci}{l}{\texttheta}{r}{\textsci}{f}} \dicPos{n}[2] \dicFlx{pl}[1] \dicSynonym{þróttur} \dicDirectTranslationCS{zaujetí, nasazení} \dicExampleIS{sýna góð tilþrif} \dicExampleCS{vykázat dobré nasazení}
\dicEntry[tilþrifamikill] \dicTerm{til·þrifa··|mikill} \dicIPA{{t\smash{\textsuperscript{h}}}{\textsci}{l}{\texttheta}{r}{\textsci}{v}{a}{m}{\textsci}{\r{\textObardotlessj}}{\textsci}{\textsubring{d}}{\textsubring{l}}} \dicPos{adj}[11] \dicFlx{(comp ‑meiri, sup ‑mest\-ur)}[12] \dicDirectTranslationCS{energický, plný zaujetí\,/\addthin nasazení} \dicExampleIS{vera tilþrifamikill stjórnandi} \dicExampleCS{být energickým velitelem}

\dicEntry[tilætlaður] \textls[15]{\dicTerm{til··ætl·|aður} \dicIPA{{t\smash{\textsuperscript{h}}}{\textsci}{\textlengthmark}{l}{a}{i}{h}{\textsubring{d}}{l}{a}{ð}{\textscy}{\textsubring{r}}} \dicPos{adj}[3] \dicFlx{(f ‑uð)}[4] \dicDirectTranslationCS{zamýšlený, chys\-ta\-ný, plánovaný}}

\dicEntry[tilætlun] \dicTerm{til··ætl|un} \dicIPA{{t\smash{\textsuperscript{h}}}{\textsci}{\textlengthmark}{l}{a}{i}{h}{\textsubring{d}}{l}{\textscy}{\textsubring{n}}} \dicPos{f}[7] \dicFlx{(‑unar, ‑anir)}[8] \dicSynonym{ætlun} \dicDirectTranslationCS{záměr, úmysl, intence}
\dicEntry[tilætlunarsamur] \dicTerm{til·ætlunar··|samur} \dicIPA{{t\smash{\textsuperscript{h}}}{\textsci}{\textlengthmark}{l}{a}{i}{h}{\textsubring{d}}{l}{\textscy}{n}{a}{\textsubring{r}}{s}{a}{m}{\textscy}{\textsubring{r}}} \dicPos{adj}[1] \dicFlx{(f ‑söm)}[2] \dicDirectTranslationCS{náročný, (mnoho) vyžadující} \dicExampleIS{Hann er tilætlunarsamur.} \dicExampleCS{Hodně toho vyžaduje.}
\dicEntry[timbraður] \dicTerm{timbr··|aður} \dicIPA{{t\smash{\textsuperscript{h}}}{\textsci}{m}{\textsubring{b}}{r}{a}{ð}{\textscy}{\textsubring{r}}} \dicPos{adj}[3] \dicFlx{(f ‑uð)}[3] \dicDirectTranslationCS{mající kocovinu}
\dicEntry[timbur] \dicTerm{timbur} \dicsymFrequent\  \dicIPA{{t\smash{\textsuperscript{h}}}{\textsci}{m}{\textsubring{b}}{\textscy}{\textsubring{r}}} \dicPos{n}[2] \dicFlx{(‑s)}[28] \dicDirectTranslationCS{(stavební) dřevo\,/\addthin dříví, řezivo} \dicExampleIS{byggja hús úr timbri} \dicExampleCS{postavit dům ze dřeva}
\dicEntry[timburhús] \dicTerm{timbur··hús} \dicsymFrequent\  \dicIPA{{t\smash{\textsuperscript{h}}}{\textsci}{m}{\textsubring{b}}{\textscy}{\textsubring{r}}{h}{u}{s}} \dicPos{n}[2] \dicFlx{(‑s, ‑)}[5] \dicDirectTranslationCS{dřevěný dům} \dicExampleIS{Timburhúsið er í góðu standi.} \dicExampleCS{Dřevěný dům je v~dobrém stavu.}
\dicEntry[timburmenn] \dicTerm{timbur··menn} \dicIPA{{t\smash{\textsuperscript{h}}}{\textsci}{m}{\textsubring{b}}{\textscy}{r}{m}{\textepsilon}{\textsubring{n}}} \dicPos{m}[13] \dicFlx{pl}[1] \dicDirectTranslationCS{kocovina}
\dicEntry[tin] \dicTerm{tin} \dicIPA{{t\smash{\textsuperscript{h}}}{\textsci}{\textlengthmark}{\textsubring{n}}} \dicPos{n}[2] \dicFlx{(‑s)}[2] \dicFieldCat{chem.} \dicDirectTranslationCS{cín} \textit{(l.~{\textLA{Sn, Stannum}})}
\dicEntry[tinbrúðkaup] \dicTerm{tin··brúð·kaup} \dicIPA{{t\smash{\textsuperscript{h}}}{\textsci}{n}{\textsubring{b}}{r}{u}{ð}{k\smash{\textsuperscript{h}}}{\oe i}{\textsubring{b}}} \dicPos{n}[2] \dicFlx{(‑s, ‑)}[5] \dicDirectTranslationCS{cínová svatba} \dicIndirectTranslationCS{(10. výročí svatby)}
\dicEntry[tindáti] \dicTerm{tin··dát|i} \dicIPA{{t\smash{\textsuperscript{h}}}{\textsci}{n}{\textsubring{d}}{au}{\textsubring{d}}{\textsci}} \dicPos{m}[1] \dicFlx{(‑a, ‑ar)}[1] \dicDirectTranslationCS{cínový vojáček}
\dicEntry[tindra] \dicTerm{tindr|a} \dicIPA{{t\smash{\textsuperscript{h}}}{\textsci}{n}{\textsubring{d}}{r}{a}} \dicPos{v}[1] \dicFlx{(‑aði)}[44] \dicDirectTranslationCS{třpytit se, jiskřit} \dicExampleIS{Stjörnurnar tindra á himninum.} \dicExampleCS{Hvězdy se třpytí na nebi.}
\dicEntry[tindur] \dicTerm{tind|ur} \dicIPA{{t\smash{\textsuperscript{h}}}{\textsci}{n}{\textsubring{d}}{\textscy}{\textsubring{r}}} \dicPos{m}[6] \dicFlx{(‑s, ‑ar)}[4] \textbf{1.} \dicSynonym*{fjallstindur} \dicDirectTranslationCS{vrchol(ek) (hory ap.)}  \textbf{2.} \dicSynonym{gaddur} \dicDirectTranslationCS{špička, hrot (vidličky ap.)} \dicExampleIS{tindur  í gaffli} \dicExampleCS{hrot vidličky}
\dicEntry[tinktúra] \dicTerm{tinktúr|a} \dicIPA{{t\smash{\textsuperscript{h}}}{i}{\r{\ng}}{\textsubring{d}}{u}{r}{a}} \dicPos{f}[1] \dicFlx{(‑u, ‑ur)}[7] \dicDirectTranslationCS{tinktura}
\dicEntry[tinna] \dicTerm{tinn|a} \dicIPA{{t\smash{\textsuperscript{h}}}{\textsci}{n}{\textlengthmark}{a}} \dicPos{f}[1] \dicFlx{(‑u, ‑ur)}[7] \dicFieldCat{geol.} \dicDirectTranslationCS{pazourek, flint}
\dicEntry[tipla] \dicTerm{tipl|a} \dicIPA{{t\smash{\textsuperscript{h}}}{\textsci}{h}{\textsubring{b}}{l}{a}} \dicPos{v}[1] \dicFlx{(‑aði)}[44] \dicDirectTranslationCS{jít po špičkách, našlapovat měkce} \dicExampleIS{tipla á tánum} \dicExampleCS{jít po špičkách}
\dicEntry[tippa] \dicTerm{tipp|a} \dicIPA{{t\smash{\textsuperscript{h}}}{\textsci}{h}{\textsubring{b}}{a}} \dicPos{v}[1] \dicFlx{(‑aði)}[1] \dicDirectTranslationCS{tipovat (na zápasy ap.)} \dicExampleIS{Hann tippar í hverri viku.} \dicExampleCS{Tipuje každý týden.};  \dicPhraseIS{tippa á e‑ð} \dicLangCat{hovor.} \dicDirectTranslationCS{tipovat na (co)}
\dicEntry[tipptopp] \dicTerm{tipp··topp} \dicIPA{{t\smash{\textsuperscript{h}}}{\textsci}{h}{\textsubring{b}}{t\smash{\textsuperscript{h}}}{\textopeno}{h}{\textsubring{b}}} \dicPos{adj}[13] \dicFlx{indecl}[1] \dicDirectTranslationCS{tiptop, dokonalý, perfektní}
\dicEntry[tissjú] \dicTerm{tissjú} \dicIPA{{t\smash{\textsuperscript{h}}}{\textsci}{s}{j}{u}} \dicPos{n}[2] \dicFlx{(‑s, ‑)}[5] \dicDirectTranslationCS{papírový kapesník}
\dicEntry[Titikakavatn] \dicTerm{Titikaka··vatn} \dicIPA{{t\smash{\textsuperscript{h}}}{\textsci}{\textlengthmark}{\textsubring{d}}{\textsci}{\r{g}}{a}{\r{g}}{a}{v}{a}{h}{\textsubring{d}}{\textsubring{n}}} \dicPos{n}[2] \dicFlx{(‑s)}[4] \dicFieldCat{geog.} \dicDirectTranslationCS{Titicaca} \dicIndirectTranslationCS{(jezero v~Centrálních Andách)}
\dicEntry[titill] \dicTerm{tit|ill} \dicIPA{{t\smash{\textsuperscript{h}}}{\textsci}{\textlengthmark}{\textsubring{d}}{\textsci}{\textsubring{d}}{\textsubring{l}}} \dicPos{m}[6] \dicFlx{(‑ils, ‑lar)}[35] \textbf{1.} \dicSynonym{nafn} \dicDirectTranslationCS{titul, název, nadpis} \dicExampleIS{titill bókar} \dicExampleCS{název knihy}  \textbf{2.} \dicSynonym{nafnbót} \dicDirectTranslationCS{titul} \dicIndirectTranslationCS{(vyjádření hodnosti)} \dicExampleIS{doktorstitill} \dicExampleCS{doktorský titul}
\dicEntry[titilsíða] \dicTerm{titil··síð|a} \dicIPA{{t\smash{\textsuperscript{h}}}{\textsci}{\textsubring{d}}{\textsci}{\textsubring{l}}{s}{i}{ð}{a}} \dicPos{f}[1] \dicFlx{(‑u, ‑ur)}[13] \dicDirectTranslationCS{titulní strana\,/\addthin stránka}
\dicEntry[titla] \dicTerm{titl|a} \dicIPA{{t\smash{\textsuperscript{h}}}{\textsci}{h}{\textsubring{d}}{l}{a}} \dicPos{v}[1] \dicFlx{(‑aði)}[1] \dicFlx{acc} \dicDirectTranslationCS{titulovat, oslovovat (hodností)} \dicExampleIS{titla e‑n} \dicExampleCS{titulovat (koho)}
\dicEntry[titra] \dicTerm{titr|a} \dicsymFrequent\  \dicIPA{{t\smash{\textsuperscript{h}}}{\textsci}{\textlengthmark}{\textsubring{d}}{r}{a}} \dicPos{v}[1] \dicFlx{(‑aði)}[1] \dicSynonym{skjálfa} \dicDirectTranslationCS{třást se, chvět se, vibrovat} \dicExampleIS{titra eins og laufblað} \dicExampleCS{chvět se jako list}
\dicEntry[titrandi] \dicTerm{titr··andi} \dicsymFrequent\  \dicIPA{{t\smash{\textsuperscript{h}}}{\textsci}{\textlengthmark}{\textsubring{d}}{r}{a}{n}{\textsubring{d}}{\textsci}} \dicPos{adj}[13] \dicFlx{indecl}[1] \dicDirectTranslationCS{třesoucí se, chvějící se, vibrující} \dicExampleIS{með titrandi höndum} \dicExampleCS{třesoucíma se rukama}
\dicEntry[titrari] \dicTerm{titr··ar|i} \dicIPA{{t\smash{\textsuperscript{h}}}{\textsci}{\textlengthmark}{\textsubring{d}}{r}{a}{r}{\textsci}} \dicPos{m}[1] \dicFlx{(‑a, ‑ar)}[13] \textbf{1.} \dicDirectTranslationCS{vibrátor} \dicIndirectTranslationCS{(zařízení budící vibrací)}  \textbf{2.} \dicDirectTranslationCS{vibrátor} \dicIndirectTranslationCS{(erotická pomůcka)}
\dicEntry[titringur] \dicTerm{titr··ing|ur} \dicIPA{{t\smash{\textsuperscript{h}}}{\textsci}{\textlengthmark}{\textsubring{d}}{r}{i}{\ng}{\r{g}}{\textscy}{\textsubring{r}}} \dicPos{m}[6] \dicFlx{(‑s)}[9] \textbf{1.} \dicSynonym{hristingur} \dicDirectTranslationCS{(za)chvění, třesení} \dicExampleIS{finna fyrir titringi} \dicExampleCS{cítit chvění}  \textbf{2.} \dicSynonym{óróleiki} \dicDirectTranslationCS{chvění, neklid}
\dicEntry[tittlingaskítur] \dicTerm{tittlinga··skít|ur} \dicIPA{{t\smash{\textsuperscript{h}}}{\textsci}{h}{\textsubring{d}}{l}{i}{\ng}{\r{g}}{a}{s}{\r{\textObardotlessj}}{i}{\textsubring{d}}{\textscy}{\textsubring{r}}} \dicPos{m}[9] \dicFlx{(‑s, ‑ir)}[8] \dicDirectTranslationCS{prkotina, drobnost, hloupost}
\dicEntry[tittlingur] \dicTerm{titt··ling|ur} \dicIPA{{t\smash{\textsuperscript{h}}}{\textsci}{h}{\textsubring{d}}{l}{i}{\ng}{\r{g}}{\textscy}{\textsubring{r}}} \dicPos{m}[6] \dicFlx{(‑s, ‑ar)}[8] \textbf{1.} \dicFieldCat{zool.} \dicSynonym{spörfugl} \dicDirectTranslationCS{pěvec} \textit{(l.~{\textLA{Passeriformes}})}  \dicsymPhoto\   \textbf{2.} \dicSynonym{getnaðarlimur} \dicDirectTranslationCS{pták, šulín} \dicIndirectTranslationCS{(penis)};  \dicIdiom{tittlingur}{ \dicPhraseIS{drepa tittlinga}} \dicDirectTranslationCS{(za)mrkat, mrknout} \dicIndirectTranslationCS{(dorozumívat se mžiknutím očí)}
\dicFigure{102360.jpg}{Tittlingur}{Tittlingur - Koleček Jaroslav, Biolib, Copyright/CC-BY}
\dicEntry[tittur] \dicTerm{titt|ur} \dicIPA{{t\smash{\textsuperscript{h}}}{\textsci}{h}{\textsubring{d}}{\textscy}{\textsubring{r}}} \dicPos{m}[9] \dicFlx{(‑s, ‑ir)}[4] \textbf{1.} \dicSynonym{pinni} \dicDirectTranslationCS{kolík, kolíček}  \textbf{2.} \dicSynonym*{lítill nagli} \dicDirectTranslationCS{hřebíček, cvoček, připínáček}  \textbf{3.} \dicSynonym*{lítill maður} \dicDirectTranslationCS{skrček, prcek}
\dicEntry[tía] \dicTerm{tí|a} \dicIPA{{t\smash{\textsuperscript{h}}}{i}{\textlengthmark}{j}{a}} \dicPos{f}[1] \dicFlx{(‑u, ‑ur)}[7] \textbf{1.} \dicDirectTranslationCS{desítka} \dicIndirectTranslationCS{(v~kartách)} \dicExampleIS{hjartatía} \dicExampleCS{srdcová desítka}  \textbf{2.} \dicDirectTranslationCS{desítka (autobus číslo 10 ap.)}
\dicEntry[tíbeskur] \dicTerm{tíbeskur} \dicIPA{{t\smash{\textsuperscript{h}}}{i}{\textlengthmark}{\textsubring{b}}{\textepsilon}{s}{\r{g}}{\textscy}{\textsubring{r}}} \dicPos{adj}[1]\dicFlx{}[-1] \dicDirectTranslationCS{tibetský}
\dicEntry[Tíbet] \dicTerm{Tíbet} \dicIPA{{t\smash{\textsuperscript{h}}}{i}{\textlengthmark}{\textsubring{b}}{\textepsilon}{\textsubring{d}}} \dicPos{n}[2] \dicFlx{(‑\,/\addthin ‑s)}[35] \dicFieldCat{geog.} \dicDirectTranslationCS{Tibet}
\dicEntry[tíbetanska] \dicTerm{tíbetanska} \dicIPA{{t\smash{\textsuperscript{h}}}{i}{\textlengthmark}{\textsubring{b}}{\textepsilon}{\textsubring{d}}{a}{n}{s}{\r{g}}{a}} \dicPos{f}[1] \dicFlx{(tíbetönsku)}[2] \dicDirectTranslationCS{tibetština}
\dicEntry[Tíbeti] \dicTerm{Tíbet|i} \dicIPA{{t\smash{\textsuperscript{h}}}{i}{\textlengthmark}{\textsubring{b}}{\textepsilon}{\textsubring{d}}{\textsci}} \dicPos{m}[1] \dicFlx{(‑a, ‑ar)}[1] \dicDirectTranslationCS{Tibeťan(ka)}
\dicEntry[Tíblisi] \dicTerm{Tíblisi} \dicIPA{{t\smash{\textsuperscript{h}}}{i}{\textsubring{b}}{l}{\textsci}{s}{\textsci}} \dicPos{subs} \dicFlx{indecl} \dicLink{Tbílísí}
\dicEntry[Tíblísí] \dicTerm{Tíblísí} \dicIPA{{t\smash{\textsuperscript{h}}}{i}{\textsubring{b}}{l}{i}{s}{i}} \dicPos{subs} \dicFlx{indecl} \dicLink{Tbílísí}
\dicEntry[tíbrá] \dicTerm{tí··brá} \dicIPA{{t\smash{\textsuperscript{h}}}{i}{\textlengthmark}{\textsubring{b}}{r}{au}} \dicPos{f}[4] \dicFlx{(‑r)}[19] \dicDirectTranslationCS{fata morgána}
\dicEntry[tíð] \dicTerm{tíð} \dicsymFrequent\  \dicIPA{{t\smash{\textsuperscript{h}}}{i}{\textlengthmark}{\texttheta}} \dicPos{f}[7] \dicFlx{(‑ar, ‑ir)}[1] \textbf{1.} \dicSynonym{tími} \dicDirectTranslationCS{doba, období, údobí, čas};  \dicPhraseIS{alla tíð} \dicFlx{adv} \dicDirectTranslationCS{vždy, po celý čas};  \dicPhraseIS{á þeirri tíð} \dicFlx{adv} \dicSynonym{þá\smash{\textsuperscript{2}}} \dicDirectTranslationCS{ve své době};  \dicPhraseIS{innan tíðar} \dicFlx{adv} \dicSynonym{bráðum} \dicDirectTranslationCS{brzy, krátce, co nevidět};  \dicPhraseIS{í háa herrans tíð} \dicFlx{adv} \dicLangCat{přen.} \dicDirectTranslationCS{strašně dlouho};  \dicPhraseIS{í seinni tíð} \dicFlx{adv} \dicDirectTranslationCS{v~poslední době, poslední dobou};  \dicPhraseIS{í tæka tíð} \dicFlx{adv} \dicDirectTranslationCS{včas};  \dicPhraseIS{með tíð og tíma} \dicFlx{adv} \dicDirectTranslationCS{(s) časem};  \dicPhraseIS{vera rólegur í tíðinni} \dicDirectTranslationCS{brát to s~klidem}  \textbf{2.} \dicFieldCat{jaz.} \dicDirectTranslationCS{čas} \dicExampleIS{nútíð} \dicExampleCS{přítomný čas};  \dicPhraseIS{núliðin tíð} \dicFieldCat{jaz.} \dicDirectTranslationCS{předpřítomný čas};  \dicPhraseIS{þáliðin tíð} \dicFieldCat{jaz.} \dicDirectTranslationCS{předminulý čas}  \textbf{3.} \dicSynonym{tíðarfar} \dicDirectTranslationCS{počasí}  \textbf{4.} \dicPhraseIS{tíðir} \dicFlx{pl} \dicSynonym{messa\smash{\textsuperscript{1}}} \dicDirectTranslationCS{bohoslužba, mše} \dicExampleIS{hlýða tíðum} \dicExampleCS{naslouchat mši}  \textbf{5.} \dicPhraseIS{tíðir} \dicFlx{pl} \dicDirectTranslationCS{menstruace, perioda, měsíčky}
\dicEntry[tíðabindi] \dicTerm{tíða··bindi} \dicIPA{{t\smash{\textsuperscript{h}}}{i}{\textlengthmark}{ð}{a}{\textsubring{b}}{\textsci}{n}{\textsubring{d}}{\textsci}} \dicPos{n}[2] \dicFlx{(‑s, ‑)}[14] \dicSynonym{dömubindi} \dicDirectTranslationCS{(menstruační) vložka}
\dicEntry[tíðablæðing] \dicTerm{tíða··blæð·ing} \dicIPA{{t\smash{\textsuperscript{h}}}{i}{\textlengthmark}{ð}{a}{\textsubring{b}}{l}{a}{i}{ð}{i}{\ng}{\r{g}}} \dicPos{f}[4] \dicFlx{(‑ar, ‑ar)}[5] \dicDirectTranslationCS{menstruace}
\dicEntry[tíðahringur] \dicTerm{tíða··hring|ur} \dicIPA{{t\smash{\textsuperscript{h}}}{i}{\textlengthmark}{ð}{a}{\textsubring{r}}{i}{\ng}{\r{g}}{\textscy}{\textsubring{r}}} \dicPos{m}[6] \dicFlx{(‑s)}[17] \dicDirectTranslationCS{menstruační cyklus}
\dicEntry[tíðahvörf] \dicTerm{tíða··hvörf} \dicIPA{{t\smash{\textsuperscript{h}}}{i}{\textlengthmark}{ð}{a}{k\smash{\textsuperscript{h}}}{v}{\oe}{r}{f}} \dicPos{n}[2] \dicFlx{pl}[9] \dicFieldCat{med.} \dicDirectTranslationCS{menopauza, přechod, klimakterium}
\dicEntry[tíðar] \dicTerm{tíðar} \dicIPA{{t\smash{\textsuperscript{h}}}{i}{\textlengthmark}{ð}{a}{\textsubring{r}}} \dicPos{adv} \dicFlx{comp} \dicLink{títt}
\dicEntry[tíðara] \dicTerm{tíðara} \dicIPA{{t\smash{\textsuperscript{h}}}{i}{\textlengthmark}{ð}{a}{r}{a}} \dicPos{adv} \dicFlx{comp} \dicLink{títt}
\dicEntry[tíðarandi] \dicTerm{tíðar··and|i} \dicIPA{{t\smash{\textsuperscript{h}}}{i}{\textlengthmark}{ð}{a}{r}{a}{n}{\textsubring{d}}{\textsci}} \dicPos{m}[1] \dicFlx{(‑a)}[3] \dicLangCat{přen.} \dicDirectTranslationCS{duch doby}
\dicEntry[tíðaratviksorð] \dicTerm{tíðar··at·viks·orð} \dicIPA{{t\smash{\textsuperscript{h}}}{i}{\textlengthmark}{ð}{a}{r}{a}{\textsubring{d}}{v}{\textsci}{\r{g}}{s}{\textopeno}{r}{\texttheta}} \dicPos{n}[2] \dicFlx{(‑s, ‑)}[5] \dicFieldCat{jaz.} \dicDirectTranslationCS{příslovce času}
\dicEntry[tíðarfar] \dicTerm{tíðar··far} \dicIPA{{t\smash{\textsuperscript{h}}}{i}{ð}{a}{\textsubring{r}}{f}{a}{\textsubring{r}}} \dicPos{n}[2] \dicFlx{(‑s)}[2] \dicSynonym{veðrátta} \dicDirectTranslationCS{počasí, klima}
\dicEntry[tíðast] \dicTerm{tíðast} \dicIPA{{t\smash{\textsuperscript{h}}}{i}{\textlengthmark}{ð}{a}{s}{\textsubring{d}}} \dicPos{adv} \dicFlx{sup} \dicLink{títt}
\dicEntry[tíðatappi] \dicTerm{tíða··tapp|i} \dicIPA{{t\smash{\textsuperscript{h}}}{i}{\textlengthmark}{ð}{a}{t\smash{\textsuperscript{h}}}{a}{h}{\textsubring{b}}{\textsci}} \dicPos{m}[1] \dicFlx{(‑a, ‑ar)}[8] \dicDirectTranslationCS{(menstruační) tampon}
\dicEntry[tíðaverkir] \dicTerm{tíða··verkir} \dicIPA{{t\smash{\textsuperscript{h}}}{i}{\textlengthmark}{ð}{a}{v}{\textepsilon}{\textsubring{r}}{\r{\textObardotlessj}}{\textsci}{\textsubring{r}}} \dicPos{m}[10] \dicFlx{pl}[22] \dicDirectTranslationCS{menstruační bolesti}
\dicEntry[tíðbeyging] \dicTerm{tíð··beyg·ing} \dicIPA{{t\smash{\textsuperscript{h}}}{i}{ð}{\textsubring{b}}{ei}{j}{i}{\ng}{\r{g}}} \dicPos{f}[4] \dicFlx{(‑ar)}[7] \dicFieldCat{jaz.} \dicDirectTranslationCS{časování, konjugace}
\dicEntry[tíðindi] \dicTerm{tíð··indi} \dicsymFrequent\  \dicIPA{{t\smash{\textsuperscript{h}}}{i}{\textlengthmark}{ð}{\textsci}{n}{\textsubring{d}}{\textsci}} \dicPos{n}[2] \dicFlx{pl}[19] \dicDirectTranslationCS{novinky, noviny, zprávy, zvěsti} \dicExampleIS{segja tíðindi} \dicExampleCS{vyprávět novinky}
\dicEntry[tíðka] \dicTerm{tíðk|a} \dicsymFrequent\  \dicIPA{{t\smash{\textsuperscript{h}}}{i}{\texttheta}{\r{g}}{a}} \dicPos{v}[1] \dicFlx{(‑aði)}[1] \dicFlx{acc} \dicDirectTranslationCS{být zvyklý, mít ve zvyku, praktikovat};  \dicIdiom{tíðkast}{ \dicPhraseIS{tíðkast}} \dicFlx{refl} \dicDirectTranslationCS{být běžný\,/\addthin obvyklý} \dicExampleIS{Það tíðkast ekki lengur að lesa ljóðin upphátt.} \dicExampleCS{Není už běžné číst básně nahlas.}
\dicEntry[tíðkanlegur] \dicTerm{tíðkan··legur} \dicIPA{{t\smash{\textsuperscript{h}}}{i}{\texttheta}{\r{g}}{a}{n}{l}{\textepsilon}{\textbabygamma}{\textscy}{\textsubring{r}}} \dicPos{adj}[1]\dicFlx{}[-8] \dicSynonym{venjulegur} \dicDirectTranslationCS{zvyklý, obvyklý, běžný}
\dicEntry[tíðni] \dicTerm{tíðn|i} \dicIPA{{t\smash{\textsuperscript{h}}}{i}{ð}{n}{\textsci}} \dicPos{f}[3] \dicFlx{(‑i, ‑ir)}[1] \textbf{1.} \dicSynonym*{fjöldi í tíma eða rúmi} \dicDirectTranslationCS{četnost, frekvence}  \textbf{2.} \dicFieldCat{fyz.} \dicSynonym*{sveiflufjöldi} \dicDirectTranslationCS{kmitočet, frekvence}
\dicEntry[tíðum] \dicTerm{tíðum} \dicIPA{{t\smash{\textsuperscript{h}}}{i}{\textlengthmark}{ð}{\textscy}{\textsubring{m}}} \dicPos{adv} \dicDirectTranslationCS{často, mnohdy};  \dicPhraseIS{oft og tíðum, oft á tíðum} \dicFlx{adv} \dicDirectTranslationCS{častokrát, mnohokrát}
\dicEntry[tíður] \dicTerm{tíður} \dicsymFrequent\  \dicIPA{{t\smash{\textsuperscript{h}}}{i}{\textlengthmark}{ð}{\textscy}{\textsubring{r}}} \dicPos{adj}[2]\dicFlx{}[-6] \dicSynonym{algengur} \dicDirectTranslationCS{častý, běžný, četný, frekventovaný} \dicExampleIS{tíðar heimsóknir} \dicExampleCS{časté návštěvy};  \dicPhraseIS{vera tíður gestur hjá e‑m} \dicDirectTranslationCS{být u~(koho) častým hostem}
\dicEntry[tífaldur] \dicTerm{tí··|faldur} \dicIPA{{t\smash{\textsuperscript{h}}}{i}{\textlengthmark}{f}{a}{l}{\textsubring{d}}{\textscy}{\textsubring{r}}} \dicPos{adj}[2] \dicFlx{(f ‑föld)}[16] \textbf{1.} \dicDirectTranslationCS{desetinásobný, desítinásobný}  \textbf{2.} \dicDirectTranslationCS{desetidílný}
\dicEntry[tífalt] \dicTerm{tí··falt} \dicIPA{{t\smash{\textsuperscript{h}}}{i}{\textlengthmark}{f}{a}{\textsubring{l}}{\textsubring{d}}} \dicPos{adv} \dicDirectTranslationCS{desetinásobně}
\dicEntry[tíglóttur] \dicTerm{tígl··óttur} \dicIPA{{t\smash{\textsuperscript{h}}}{i}{\r{g}}{l}{ou}{h}{\textsubring{d}}{\textscy}{\textsubring{r}}} \dicPos{adj}[1]\dicFlx{}[-10] \dicDirectTranslationCS{kosočtverečný, kosočtvercový}
\dicEntry[Tígris] \dicTerm{Tígris} \dicIPA{{t\smash{\textsuperscript{h}}}{i}{\textbabygamma}{r}{\textsci}{s}} \dicPos{n}[2] \dicFlx{(‑s)}[32] \dicFieldCat{geog.} \dicDirectTranslationCS{Tigris} \dicIndirectTranslationCS{(řeka v~Asii)}
\dicEntry[tígrisdýr] \dicTerm{tígris··dýr}\dicTerm{, tígur} \dicIPA{{t\smash{\textsuperscript{h}}}\-{i}\-{\textbabygamma}\-{r}\-{\textsci}\-{s}\-{\textsubring{d}}\-{i}\-{\textsubring{r}}\-} \dicPos{n}[2] \dicFlx{(‑s, ‑)}[5] \dicFieldCat{zool.} \dicDirectTranslationCS{tygr} \textit{(l.~{\textLA{Panthera tigris}})}  \dicsymPhoto\ 
\dicFigure{26787.jpg}{Tígrisdýr}{Tígrisdýr - Velechovský Marek, Biolib, Copyright/CC-BY-NC}
\dicEntry[tígulegur] \dicTerm{tígu··legur}\dicTerm{, tigulegur} \dicIPA{{t\smash{\textsuperscript{h}}}\-{i}\-{\textlengthmark}\-{\textbabygamma}\-{\textscy}\-{l}\-{\textepsilon}\-{\textbabygamma}\-{\textscy}\-{\textsubring{r}}\-} \dicPos{adj}[1]\dicFlx{}[-8] \dicSynonym{myndarlegur} \dicDirectTranslationCS{elegantní, vznešený, spanilý} \dicExampleIS{vera tígulegur að sjá} \dicExampleCS{být elegantní od pohledu}
\dicEntry[tígull] \dicTerm{tíg|ull}\dicTerm{, tigull} \dicIPA{{t\smash{\textsuperscript{h}}}\-{i}\-{\textlengthmark}\-{\textbabygamma}\-{\textscy}\-{\textsubring{d}}\-{\textsubring{l}}\-} \dicPos{m}[6] \dicFlx{(‑uls, ‑lar)}[31] \textbf{1.} \dicSynonym*{litur í spilum} \dicDirectTranslationCS{kára} \dicIndirectTranslationCS{(barva v~kartách)}  \textbf{2.} \dicFieldCat{mat.} \dicDirectTranslationCS{kosočtverec}  \textbf{3.} \dicSynonym{tígulsteinn} \dicDirectTranslationCS{(okrasná) cihla, dlaždice}
\dicEntry[tígullaga] \dicTerm{tígul··laga} \dicIPA{{t\smash{\textsuperscript{h}}}{i}{\textlengthmark}{\textbabygamma}{\textscy}{l}{a}{\textbabygamma}{a}} \dicPos{adj}[13] \dicFlx{indecl}[1] \dicDirectTranslationCS{kosočtverečný, kosočtvercový}
\dicEntry[tígulsteinn] \textls[10]{\dicTerm{tígul··stein|n}\dicTerm{, tigulsteinn} \dicIPA{{t\smash{\textsuperscript{h}}}\-{i}\-{\textlengthmark}\-{\textbabygamma}\-{\textscy}\-{\textsubring{l}}\-{s}\-{\textsubring{d}}\-{ei}\-{\textsubring{d}}\-{\textsubring{n}}\-} \dicPos{m}[6] \dicFlx{(‑s, ‑ar)}[42]} \dicSynonym{tíg\-ull} \dicDirectTranslationCS{(okrasná) cihla, dlaždice}
\dicEntry[tígur] \dicTerm{tíg|ur} \dicIPA{{t\smash{\textsuperscript{h}}}{i}{\textlengthmark}{\textbabygamma}{\textscy}{\textsubring{r}}} \dicPos{m}[5] \dicFlx{(‑urs, ‑rar)}[1] \dicLink{tígrisdýr}
\dicEntry[tíhyrningur] \dicTerm{tí··hyrn·ing|ur} \dicIPA{{t\smash{\textsuperscript{h}}}{i}{\textlengthmark}{h}{\textsci}{r}{\textsubring{d}}{n}{i}{\ng}{\r{g}}{\textscy}{\textsubring{r}}} \dicPos{m}[6] \dicFlx{(‑s, ‑ar)}[8] \dicFieldCat{mat.} \dicDirectTranslationCS{desetiúhelník}
\dicEntry[tík] \dicTerm{tík} \dicIPA{{t\smash{\textsuperscript{h}}}{i}{\textlengthmark}{\r{g}}} \dicPos{f}[10] \dicFlx{(‑ur\,/\addthin ‑ar, ‑ur)}[7] \textbf{1.} \dicDirectTranslationCS{psice, fena, čubka}  \textbf{2.} \dicLangCat{han.} \dicDirectTranslationCS{čubka, mrcha}
\dicEntry[tíkall] \dicTerm{tí··kall} \dicIPA{{t\smash{\textsuperscript{h}}}{i}{\textlengthmark}{k\smash{\textsuperscript{h}}}{a}{\textsubring{d}}{\textsubring{l}}} \dicPos{m}[4] \dicFlx{(‑s, ‑ar)}[4] \textbf{1.} \dicLangCat{hovor.} \dicDirectTranslationCS{deset (korun)} \dicIndirectTranslationCS{(částka)}  \textbf{2.} \dicLangCat{hovor.} \dicDirectTranslationCS{desetikoruna, pětka} \dicIndirectTranslationCS{(mince)}
\dicEntry[tíkarlegur] \dicTerm{tíkar··legur} \dicIPA{{t\smash{\textsuperscript{h}}}{i}{\textlengthmark}{\r{g}}{a}{r}{l}{\textepsilon}{\textbabygamma}{\textscy}{\textsubring{r}}} \dicPos{adj}[1]\dicFlx{}[-8] \dicSynonym{skammarlegur} \dicDirectTranslationCS{nízký, hanebný, ostudný} \dicExampleIS{vera tíkarlegur við hana} \dicExampleCS{hanebně se k~ní chovat}
\dicEntry[tíkarspeni] \dicTerm{tíkar··spen|i} \dicIPA{{t\smash{\textsuperscript{h}}}{i}{\textlengthmark}{\r{g}}{a}{\textsubring{r}}{s}{\textsubring{b}}{\textepsilon}{n}{\textsci}} \dicPos{m}[1] \dicFlx{(‑a, ‑ar)}[1] \textbf{1.} \dicSynonym*{speni á tík} \dicDirectTranslationCS{struk (feny)}  \textbf{2.} \dicSynonym*{hárflétta} \dicDirectTranslationCS{cop, copánek} \dicIndirectTranslationCS{(pletenec vlasů)}
\dicEntry[tíma] \dicTerm{tím|a} \dicIPA{{t\smash{\textsuperscript{h}}}{i}{\textlengthmark}{m}{a}} \dicPos{v}[2] \dicFlx{(‑di, ‑t)}[155] \dicDirectTranslationCS{mít ochotu\,/\addthin chuť, chtít} \dicIndirectTranslationCS{(nejčastěji se záporem)} \dicExampleIS{Ég tími ekki að vekja hann.} \dicExampleCS{Nemám chuť ho budit.}
\dicEntry[tímaáætlun] \dicTerm{tíma··á·ætl|un} \dicIPA{{t\smash{\textsuperscript{h}}}{i}{\textlengthmark}{m}{a}{au}{a}{i}{h}{\textsubring{d}}{l}{\textscy}{\textsubring{n}}} \dicPos{f}[7] \dicFlx{(‑unar, ‑anir)}[8] \textbf{1.} \dicDirectTranslationCS{časový plán, časové rozvržení, harmonogram}  \textbf{2.} \dicDirectTranslationCS{jízdní řád}
\dicEntry[tímabelti] \dicTerm{tíma··belti} \dicIPA{{t\smash{\textsuperscript{h}}}{i}{\textlengthmark}{m}{a}{\textsubring{b}}{\textepsilon}{\textsubring{l}}{\textsubring{d}}{\textsci}} \dicPos{n}[2] \dicFlx{(‑s, ‑)}[14] \dicFieldCat{geog.} \dicDirectTranslationCS{časové pásmo}
\dicEntry[tímabil] \dicTerm{tíma··bil} \dicsymFrequent\  \dicIPA{{t\smash{\textsuperscript{h}}}{i}{\textlengthmark}{m}{a}{\textsubring{b}}{\textsci}{\textsubring{l}}} \dicPos{n}[2] \dicFlx{(‑s, ‑)}[5] \dicDirectTranslationCS{(časové) období, epocha} \dicExampleIS{á þriggja mánaða tímabili} \dicExampleCS{v~rozpětí tří měsíců}
\dicEntry[tímabundinn] \dicTerm{tíma··bundinn} \dicsymFrequent\  \dicIPA{{t\smash{\textsuperscript{h}}}{i}{\textlengthmark}{m}{a}{\textsubring{b}}{\textscy}{n}{\textsubring{d}}{\textsci}{\textsubring{n}}} \dicPos{adj}[6]\dicFlx{}[-2] \textbf{1.} \dicDirectTranslationCS{dočasný, prozatímní, provizorní} \dicExampleIS{e‑að er tímabundið} \dicExampleCS{(co) je dočasné}  \textbf{2.} \dicDirectTranslationCS{zaneprázdněný, vázaný časem}
\dicEntry[tímabær] \dicTerm{tíma··bær} \dicIPA{{t\smash{\textsuperscript{h}}}{i}{\textlengthmark}{m}{a}{\textsubring{b}}{a}{i}{\textsubring{r}}} \dicPos{adj}[5] \dicFlx{(f ‑)}[8] \dicDirectTranslationCS{příhodný, pravý (okamžik ap.)};  \dicPhraseIS{það er tímabært að (gera e‑ð)} \dicDirectTranslationCS{je (nejvyšší) čas (udělat (co))}
\dicEntry[tímaeyðsla] \dicTerm{tíma··eyðsl|a} \dicIPA{{t\smash{\textsuperscript{h}}}{i}{\textlengthmark}{m}{a}{ei}{ð}{s}{\textsubring{d}}{l}{a}} \dicPos{f}[1] \dicFlx{(‑u)}[5] \dicDirectTranslationCS{plýtvání časem, ztráta času}
\dicEntry[tímafrekur] \dicTerm{tíma··frekur} \dicIPA{{t\smash{\textsuperscript{h}}}{i}{\textlengthmark}{m}{a}{f}{r}{\textepsilon}{\r{g}}{\textscy}{\textsubring{r}}} \dicPos{adj}[1]\dicFlx{}[-1] \dicSynonym{tafsamur} \dicDirectTranslationCS{časově náročný, zdlouhavý}
\dicEntry[tímalaus] \dicTerm{tíma··laus} \dicIPA{{t\smash{\textsuperscript{h}}}{i}{\textlengthmark}{m}{a}{l}{\oe i}{s}} \dicPos{adj}[5]\dicFlx{}[-1] \textbf{1.} \dicSynonym{tímabundinn} \dicDirectTranslationCS{zaneprázdněný, nemající čas} \dicExampleIS{Ég er alveg tímalaus.} \dicExampleCS{Nemám čas.}  \textbf{2.} \dicSynonym*{óháður tíma} \dicDirectTranslationCS{nadčasový, nestárnoucí}
\dicEntry[tímamót] \dicTerm{tíma··mót} \dicIPA{{t\smash{\textsuperscript{h}}}{i}{\textlengthmark}{m}{a}{m}{ou}{\textsubring{d}}} \dicPos{n}[2] \dicFlx{pl}[1] \dicDirectTranslationCS{(časový) předěl, mezník, přelom} \dicExampleIS{Þessi ákvörðun markaði tímamót í lífi mínu.} \dicExampleCS{Toto rozhodnutí bylo přelomem v~mém životě.}
\dicEntry[tímamörk] \dicTerm{tíma··mörk} \dicIPA{{t\smash{\textsuperscript{h}}}{i}{\textlengthmark}{m}{a}{m}{\oe}{\textsubring{r}}{\r{g}}} \dicPos{n}[2] \dicFlx{pl}[9] \dicDirectTranslationCS{(časová) lhůta, termín}
\dicEntry[tímanlega] \dicTerm{tíman··lega} \dicIPA{{t\smash{\textsuperscript{h}}}{i}{\textlengthmark}{m}{a}{n}{l}{\textepsilon}{\textbabygamma}{a}} \dicPos{adv} \dicDirectTranslationCS{(za)včas}
\dicEntry[tímanlegur] \dicTerm{tíman··legur} \dicIPA{{t\smash{\textsuperscript{h}}}{i}{\textlengthmark}{m}{a}{n}{l}{\textepsilon}{\textbabygamma}{\textscy}{\textsubring{r}}} \dicPos{adj}[1]\dicFlx{}[-8] \dicDirectTranslationCS{včasný, (jsoucí) včas}
\dicEntry[tímapressa] \dicTerm{tíma··press|a} \dicIPA{{t\smash{\textsuperscript{h}}}{i}{\textlengthmark}{m}{a}{p\smash{\textsuperscript{h}}}{r}{\textepsilon}{s}{a}} \dicPos{f}[1] \dicFlx{(‑u)}[5] \dicDirectTranslationCS{časová tíseň, časový tlak}
\dicEntry[tímarit] \dicTerm{tíma··rit} \dicsymFrequent\  \dicIPA{{t\smash{\textsuperscript{h}}}{i}{\textlengthmark}{m}{a}{r}{\textsci}{\textsubring{d}}} \dicPos{n}[2] \dicFlx{(‑s, ‑)}[5] \dicDirectTranslationCS{časopis, magazín, periodikum} \dicExampleIS{grein í íslensku tímariti} \dicExampleCS{článek v~islandském časopise}
\dicEntry[tímaritsgrein] \dicTerm{tíma·rits··grein} \dicIPA{{t\smash{\textsuperscript{h}}}{i}{\textlengthmark}{m}{a}{r}{\textsci}{\textsubring{d}}{s}{\r{g}}{r}{ei}{\textsubring{n}}} \dicPos{f}[4] \dicFlx{(‑ar, ‑ar)}[1] \dicDirectTranslationCS{článek v~časopise}
\dicEntry[tímarúm] \dicTerm{tíma··rúm} \dicIPA{{t\smash{\textsuperscript{h}}}{i}{\textlengthmark}{m}{a}{r}{u}{\textsubring{m}}} \dicPos{n}[2] \dicFlx{(‑s)}[2] \dicFieldCat{fyz.} \dicDirectTranslationCS{časoprostor}
\dicEntry[tímaröð] \dicTerm{tíma··|röð} \dicIPA{{t\smash{\textsuperscript{h}}}{i}{\textlengthmark}{m}{a}{r}{\oe}{\texttheta}} \dicPos{f}[7] \dicFlx{(‑raðar)}[19] \dicDirectTranslationCS{časové\,/\addthin chronologické pořadí, časový sled, chronologie} \dicExampleIS{í tímaröð} \dicExampleCS{chronologicky}
\dicEntry[tímaset] \dicTerm{tíma··set} \dicIPA{{t\smash{\textsuperscript{h}}}{i}{\textlengthmark}{m}{a}{s}{\textepsilon}{\textsubring{d}}} \dicPos{v} \dicFlx{ind praes sg 1 pers} \dicLink{tímasetja}
\dicEntry[tímasetja] \dicTerm{tíma··|setja} \dicIPA{{t\smash{\textsuperscript{h}}}{i}{\textlengthmark}{m}{a}{s}{\textepsilon}{\textsubring{d}}{j}{a}} \dicPos{v}[4] \dicFlx{(‑set, ‑setti, ‑settum, ‑setti, ‑sett)}[26] \dicFlx{acc} \textbf{1.} \dicDirectTranslationCS{datovat, zařadit časově} \dicExampleIS{tímasetja gamalt handrit} \dicExampleCS{datovat starý rukopis}  \textbf{2.} \dicDirectTranslationCS{načasovat, (časově) rozvrhnout} \dicExampleIS{tímasetja fund} \dicExampleCS{načasovat schůzi}
\dicEntry[tímasetning] \dicTerm{tíma··set·ning} \dicIPA{{t\smash{\textsuperscript{h}}}{i}{\textlengthmark}{m}{a}{s}{\textepsilon}{h}{\textsubring{d}}{n}{i}{\ng}{\r{g}}} \dicPos{f}[4] \dicFlx{(‑ar, ‑ar)}[5] \textbf{1.} \dicSynonym{dagsetning} \dicDirectTranslationCS{datování, datum}  \textbf{2.} \dicSynonym*{valinn tími} \dicDirectTranslationCS{(na)časování, časový rozvrh}
\dicEntry[tímasett] \dicTerm{tíma··sett} \dicIPA{{t\smash{\textsuperscript{h}}}{i}{\textlengthmark}{m}{a}{s}{\textepsilon}{h}{\textsubring{d}}} \dicPos{v} \dicFlx{supin} \dicLink{tímasetja}
\dicEntry[tímasetti] \dicTerm{tíma··setti} \dicIPA{{t\smash{\textsuperscript{h}}}{i}{\textlengthmark}{m}{a}{s}{\textepsilon}{h}{\textsubring{d}}{\textsci}} \dicPos{v} \dicFlx{ind\,/\addthin con pf sg 1 pers} \dicLink{tímasetja}
\dicEntry[tímasettum] \dicTerm{tíma··settum} \dicIPA{{t\smash{\textsuperscript{h}}}{i}{\textlengthmark}{m}{a}{s}{\textepsilon}{h}{\textsubring{d}}{\textscy}{\textsubring{m}}} \dicPos{v} \dicFlx{ind pf pl 1 pers} \dicLink{tímasetja}
\dicEntry[tímaskekkja] \dicTerm{tíma··skekkj|a} \dicIPA{{t\smash{\textsuperscript{h}}}{i}{\textlengthmark}{m}{a}{s}{\r{\textObardotlessj}}{\textepsilon}{h}{\r{\textObardotlessj}}{a}} \dicPos{f}[1] \dicFlx{(‑u, ‑ur)}[17] \dicDirectTranslationCS{anachronismus}
\dicEntry[tímaskyn] \dicTerm{tíma··skyn}\dicTerm{, tímaskynjun} \dicIPA{{t\smash{\textsuperscript{h}}}\-{i}\-{\textlengthmark}\-{m}\-{a}\-{s}\-{\r{\textObardotlessj}}\-{\textsci}\-{\textsubring{n}}\-} \dicPos{n}[2] \dicFlx{(‑s)}[2] \dicDirectTranslationCS{smysl pro čas, vnímání času}
\dicEntry[tímaskynjun] \dicTerm{tíma··skynj|un} \dicIPA{{t\smash{\textsuperscript{h}}}{i}{\textlengthmark}{m}{a}{s}{\r{\textObardotlessj}}{\textsci}{n}{j}{\textscy}{\textsubring{n}}} \dicPos{f}[7] \dicFlx{(‑unar)}[9] \dicLink{tímaskyn}
\dicEntry[tímasókn] \dicTerm{tíma··sókn} \dicIPA{{t\smash{\textsuperscript{h}}}{i}{\textlengthmark}{m}{a}{s}{ou}{h}{\r{g}}{\textsubring{n}}} \dicPos{f}[7] \dicFlx{(‑ar, ‑ir)}[1] \dicDirectTranslationCS{docházka, prezence}
\dicEntry[tímasprengja] \dicTerm{tíma··sprengj|a} \dicIPA{{t\smash{\textsuperscript{h}}}{i}{\textlengthmark}{m}{a}{s}{\textsubring{b}}{r}{ei}{\textltailn}{\r{\textObardotlessj}}{a}} \dicPos{f}[1] \dicFlx{(‑u, ‑ur)}[25] \textbf{1.} \dicDirectTranslationCS{časovaná bomba}  \textbf{2.} \dicLangCat{přen.} \dicDirectTranslationCS{časovaná\,/\addthin tikající bomba}
\dicEntry[tímaspursmál] \dicTerm{tíma··spurs·mál} \dicIPA{{t\smash{\textsuperscript{h}}}{i}{\textlengthmark}{m}{a}{s}{\textsubring{b}}{\textscy}{\textsubring{r}}{s}{m}{au}{\textsubring{l}}} \dicPos{n}[2] \dicFlx{(‑s, ‑)}[5] \dicDirectTranslationCS{otázka času};  \dicPhraseIS{það er tímaspursmál hvenær (e‑að gerist)} \dicDirectTranslationCS{je to otázka času, kdy (se (co) stane)}
\dicEntry[tímatal] \dicTerm{tíma··|tal} \dicIPA{{t\smash{\textsuperscript{h}}}{i}{\textlengthmark}{m}{a}{t\smash{\textsuperscript{h}}}{a}{\textsubring{l}}} \dicPos{n}[2] \dicFlx{(‑tals, ‑töl)}[8] \dicSynonym{almanak} \dicDirectTranslationCS{kalendář}
\dicEntry[tímavél] \dicTerm{tíma··vél} \dicIPA{{t\smash{\textsuperscript{h}}}{i}{\textlengthmark}{m}{a}{v}{j}{\textepsilon}{\textsubring{l}}} \dicPos{f}[4] \dicFlx{(‑ar, ‑ar)}[1] \dicDirectTranslationCS{stroj času}
\dicEntry[tímavinna] \dicTerm{tíma··vinn|a} \dicIPA{{t\smash{\textsuperscript{h}}}{i}{\textlengthmark}{m}{a}{v}{\textsci}{n}{a}} \dicPos{f}[1] \dicFlx{(‑u)}[5] \dicDirectTranslationCS{práce placená od hodiny}
\dicEntry[tímaþröng] \dicTerm{tíma··þröng} \dicIPA{{t\smash{\textsuperscript{h}}}{i}{\textlengthmark}{m}{a}{\texttheta}{r}{\oe i}{\ng}{\r{g}}} \dicPos{f}[4] \dicFlx{(‑ar)}[3] \dicDirectTranslationCS{časová tíseň};  \dicPhraseIS{vera í tímaþröng} \dicDirectTranslationCS{být v~časové tísni}
\dicEntry[tími] \dicTerm{tím|i} \dicsymFrequent\  \dicIPA{{t\smash{\textsuperscript{h}}}{i}{\textlengthmark}{m}{\textsci}} \dicPos{m}[1] \dicFlx{(‑a, ‑ar)}[1] \textbf{1.} \dicSynonym{tíð} \dicDirectTranslationCS{čas, doba};  \dicPhraseIS{á óhentugum tíma} \dicFlx{adv} \dicDirectTranslationCS{v~nevhodnou dobu};  \dicPhraseIS{drepa tímann} \dicLangCat{přen.} \dicDirectTranslationCS{trávit\,/\addthin zabíjet čas};  \dicPhraseIS{einhvern tíma(nn)} \dicFlx{adv} \dicDirectTranslationCS{někdy, jednou};  \dicPhraseIS{fylgjast með tímanum} \dicFlx{refl} \dicLangCat{přen.} \dicDirectTranslationCS{jít s~dobou};  \dicPhraseIS{gefa sér tíma til að (gera e‑ð)} \dicDirectTranslationCS{nechat si čas na (udělání (čeho))};  \dicPhraseIS{hafa tímann fyrir sér} \dicDirectTranslationCS{mít čas pro sebe};  \dicPhraseIS{í tíma} \dicFlx{adv} \dicDirectTranslationCS{včas};  \dicPhraseIS{í tíma og ótíma} \dicFlx{adv} \dicDirectTranslationCS{neustále, nepřetržitě, bez přestání};  \dicPhraseIS{í þann tíma} \dicFlx{adv} \dicSynonym{þá\smash{\textsuperscript{2}}} \dicDirectTranslationCS{v~té době, tou dobou};  \dicPhraseIS{með tímanum} \dicFlx{adv} \dicDirectTranslationCS{časem};  \dicPhraseIS{nokkurn tíma(nn)} \dicFlx{adv} \dicDirectTranslationCS{vůbec};  \dicPhraseIS{um tíma} \dicFlx{adv} \dicDirectTranslationCS{občas, chvílemi};  \dicPhraseIS{það er tími til kominn að (gera e‑ð)} \dicDirectTranslationCS{je (nejvyšší) čas (udělat (co))}  \textbf{2.} \dicSynonym*{rétt stund} \dicDirectTranslationCS{(správný) čas\,/\addthin moment};  \dicPhraseIS{á réttum tíma} \dicFlx{adv} \dicDirectTranslationCS{v~pravý čas};  \dicPhraseIS{fyrir tímann} \dicFlx{adv} \dicSynonym*{of snemma} \dicDirectTranslationCS{předčasně}  \textbf{3.} \dicSynonym{klukkustund} \dicDirectTranslationCS{hodina} \dicIndirectTranslationCS{(šedesát minut)} \dicExampleIS{tveggja tíma gangur} \dicExampleCS{dvouhodinový pochod};  \dicPhraseIS{á tólfta tímanum} \dicFlx{adv} \dicLangCat{přen.} \dicDirectTranslationCS{v~hodině dvanácté, v\,/\addthin na poslední chvíli};  \dicPhraseIS{tímunum saman} \dicFlx{adv} \dicDirectTranslationCS{(celými) hodinami}  \textbf{4.} \dicSynonym{kennslustund} \dicDirectTranslationCS{(vyučovací) hodina,  lekce} \dicExampleIS{íslenskutími} \dicExampleCS{hodina islandštiny};  \dicPhraseIS{mæta í tíma} \dicDirectTranslationCS{přijít na hodinu}
\dicEntry[tína] \dicTerm{tín|a} \dicsymFrequent\  \dicIPA{{t\smash{\textsuperscript{h}}}{i}{\textlengthmark}{n}{a}} \dicPos{v}[2] \dicFlx{(‑di, ‑t)}[141] \dicFlx{acc} \dicDirectTranslationCS{sbírat, trhat, česat (maliny ap.)} \dicExampleIS{tína blóm} \dicExampleCS{sbírat květiny};  \dicIdiom{tína}[saman]{ \dicPhraseIS{tína saman e‑ð}} \dicDirectTranslationCS{(po)sbírat (co), sesbírat (co) (smetí ap.)};  \dicIdiom{tína}[upp]{ \dicPhraseIS{tína upp e‑ð}} \dicDirectTranslationCS{vysbírat (co), sesbírat (co)};  \dicIdiom{tínast}[burt]{ \dicPhraseIS{tínast burt}} \dicFlx{refl} \dicDirectTranslationCS{sebrat se, sbírat se, odcházet jeden za druhým (z~večírku ap.)} \dicExampleIS{Fólk tínist burt úr veislunni.} \dicExampleCS{Lidé se sbírají z~oslavy.};  \dicIdiom{tínast}[inn]{ \dicPhraseIS{tínast inn}} \dicFlx{refl} \dicDirectTranslationCS{začít se objevovat, přicházet (nové zboží ap.)}
\dicEntry[tíningur] \dicTerm{tín··ing|ur} \dicIPA{{t\smash{\textsuperscript{h}}}{i}{\textlengthmark}{n}{i}{\ng}{\r{g}}{\textscy}{\textsubring{r}}} \dicPos{m}[6] \dicFlx{(‑s, ‑ar)}[8] \dicSynonym*{slæðingur} \dicDirectTranslationCS{střípek, paběrek}
\dicEntry[tínsla] \dicTerm{tínsl|a} \dicIPA{{t\smash{\textsuperscript{h}}}{i}{n}{s}{\textsubring{d}}{l}{a}} \dicPos{f}[1] \dicFlx{(‑u)}[5] \dicDirectTranslationCS{sběr, kolekce}
\dicEntry[tíra] \dicTerm{tír|a\smash{\textsuperscript{1}}}\dicTerm{, týra\smash{\textsuperscript{1}}} \dicIPA{{t\smash{\textsuperscript{h}}}\-{i}\-{\textlengthmark}\-{r}\-{a}\-} \dicPos{f}[1] \dicFlx{(‑u, ‑ur)}[7] \dicSynonym*{dauft ljós} \dicDirectTranslationCS{slabé\,/\addthin mdlé světlo}
\dicEntry[tíra] \dicTerm{tír|a\smash{\textsuperscript{2}}}\dicTerm{, týra\smash{\textsuperscript{2}}} \dicIPA{{t\smash{\textsuperscript{h}}}\-{i}\-{\textlengthmark}\-{r}\-{a}\-} \dicPos{v}[2] \dicFlx{(‑ði, ‑t)}[111] \dicSynonym{loga} \dicDirectTranslationCS{slabě\,/\addthin mdle svítit};  \dicPhraseIS{það tírir á e‑u} \dicFlx{impers} \dicDirectTranslationCS{(co) slabě svítí}
\dicEntry[tíræðisaldur] \dicTerm{tí·ræðis··ald|ur} \dicIPA{{t\smash{\textsuperscript{h}}}{i}{\textlengthmark}{r}{a}{i}{ð}{\textsci}{s}{a}{l}{\textsubring{d}}{\textscy}{\textsubring{r}}} \dicPos{m}[5] \dicFlx{(‑urs, ‑rar)}[3] \dicDirectTranslationCS{věk po devadesátce, věk mezi 90. a~99. rokem}
\dicEntry[tíræður] \dicTerm{tí··ræður} \dicIPA{{t\smash{\textsuperscript{h}}}{i}{\textlengthmark}{r}{a}{i}{ð}{\textscy}{\textsubring{r}}} \dicPos{adj}[2]\dicFlx{}[-10] \textbf{1.} \dicSynonym*{100 ára} \dicDirectTranslationCS{stoletý} \dicExampleIS{tíræður að aldri} \dicExampleCS{sto let starý}  \textbf{2.} \dicSynonym*{100 faðma djúpur} \dicDirectTranslationCS{měřící sto sáhů} \dicExampleIS{tírætt bjarg} \dicExampleCS{sto sáhů vysoká skála}
\dicEntry[tíska] \dicTerm{tísk|a} \dicIPA{{t\smash{\textsuperscript{h}}}{i}{s}{\r{g}}{a}} \dicPos{f}[1] \dicFlx{(‑u, ‑ur)}[7] \dicDirectTranslationCS{móda};  \dicPhraseIS{e‑að er í tísku} \dicDirectTranslationCS{(co) je v~módě};  \dicPhraseIS{e‑að er komið úr tísku} \dicDirectTranslationCS{(co) vychází z~módy}
\dicEntry[tískublað] \dicTerm{tísku··|blað} \dicIPA{{t\smash{\textsuperscript{h}}}{i}{s}{\r{g}}{\textscy}{\textsubring{b}}{l}{a}{\texttheta}} \dicPos{n}[2] \dicFlx{(‑blaðs, ‑blöð)}[8] \dicDirectTranslationCS{módní časopis}
\dicEntry[tískubúð] \dicTerm{tísku··búð} \dicIPA{{t\smash{\textsuperscript{h}}}{i}{s}{\r{g}}{\textscy}{\textsubring{b}}{u}{\texttheta}} \dicPos{f}[7] \dicFlx{(‑ar, ‑ir)}[1] \dicDirectTranslationCS{módní obchod, butik}
\dicEntry[tískufatnaður] \dicTerm{tísku··fat·nað|ur} \dicIPA{{t\smash{\textsuperscript{h}}}{i}{s}{\r{g}}{\textscy}{f}{a}{\textsubring{d}}{n}{a}{ð}{\textscy}{\textsubring{r}}} \dicPos{m}[10] \dicFlx{(‑ar)}[9] \dicDirectTranslationCS{módní oblečení}
\dicEntry[tískuhús] \dicTerm{tísku··hús} \dicIPA{{t\smash{\textsuperscript{h}}}{i}{s}{\r{g}}{\textscy}{h}{u}{s}} \dicPos{n}[2] \dicFlx{(‑s, ‑)}[5] \dicDirectTranslationCS{módní dům} \dicExampleIS{franskt tískuhús} \dicExampleCS{francouzský módní dům}
\dicEntry[tískuhönnuður] \dicTerm{tísku··hönn·uð|ur} \dicIPA{{t\smash{\textsuperscript{h}}}{i}{s}{\r{g}}{\textscy}{h}{\oe}{n}{\textscy}{ð}{\textscy}{\textsubring{r}}} \dicPos{m}[10] \dicFlx{(‑ar, ‑ir)}[4] \dicDirectTranslationCS{módní návrhář(ka)}
\dicEntry[tískusýning] \dicTerm{tísku··sýn·ing} \dicIPA{{t\smash{\textsuperscript{h}}}{i}{s}{\r{g}}{\textscy}{s}{i}{n}{i}{\ng}{\r{g}}} \dicPos{f}[4] \dicFlx{(‑ar, ‑ar)}[5] \dicDirectTranslationCS{módní přehlídka}
\begin{xtolerant}{}{1pt}
\dicEntry[tískusýningarstúlka] \dicTerm{tísku·sýningar··stúlk|a}\addthinS\dicIPA{{t\smash{\textsuperscript{h}}}\-{i}\-{s}\-{\r{g}}\-{\textscy}\-{s}\-{i}\-{n}\-{i}\-{\ng}\-{\r{g}}\-{a}\-{\textsubring{r}}\-{s}\-{\textsubring{d}}\-{u}\-{\textsubring{l}}\-{\r{g}}\-{a}}\addthinS\dicPos{f}[1]\addthinS\dicFlx{(‑u, ‑ur)}[13] \dicDirectTranslationCS{modelka, manekýnka}
\end{xtolerant}
\dicEntry[tískutímarit] \dicTerm{tísku··tíma·rit} \dicIPA{{t\smash{\textsuperscript{h}}}{i}{s}{\r{g}}{\textscy}{t\smash{\textsuperscript{h}}}{i}{m}{a}{r}{\textsci}{\textsubring{d}}} \dicPos{n}[2] \dicFlx{(‑s, ‑)}[5] \dicDirectTranslationCS{módní časopis}
\dicEntry[tískuverslun] \dicTerm{tísku··versl|un} \dicIPA{{t\smash{\textsuperscript{h}}}{i}{s}{\r{g}}{\textscy}{v}{\textepsilon}{\textsubring{r}}{s}{\textsubring{d}}{l}{\textscy}{\textsubring{n}}} \dicPos{f}[7] \dicFlx{(‑unar, ‑anir)}[8] \dicDirectTranslationCS{módní obchod, butik}
\dicEntry[tíst] \dicTerm{tíst} \dicIPA{{t\smash{\textsuperscript{h}}}{i}{s}{\textsubring{d}}} \dicPos{n}[2] \dicFlx{(‑s)}[2] \textbf{1.} \dicSynonym*{veikt hljóð fugla} \dicDirectTranslationCS{pípání (ptáčka ap.)} \dicExampleIS{heyra tíst í fugli} \dicExampleCS{slyšet pípání ptáčka}  \textbf{2.} \dicSynonym*{veikt hljóð músa} \dicDirectTranslationCS{pištění (myši ap.)}
\dicEntry[tísta] \dicTerm{tíst|a} \dicIPA{{t\smash{\textsuperscript{h}}}{i}{s}{\textsubring{d}}{a}} \dicPos{v}[2] \dicFlx{(‑i, ‑)}[17] \textbf{1.} \dicDirectTranslationCS{(za)pískat, (za)pípat, (za)švitořit (pták ap.)} \dicExampleIS{Fuglinn tístir.} \dicExampleCS{Ptáček pípá.}  \textbf{2.} \dicDirectTranslationCS{(za)pištět (myš ap.)}
\dicEntry[títan] \dicTerm{títan} \dicIPA{{t\smash{\textsuperscript{h}}}{i}{\textlengthmark}{\textsubring{d}}{a}{\textsubring{n}}} \dicPos{n}[2] \dicFlx{(‑s)}[2] \dicFieldCat{chem.} \dicDirectTranslationCS{titan} \textit{(l.~{\textLA{Ti, Titanium}})}
\dicEntry[títt] \dicTerm{títt} \dicIPA{{t\smash{\textsuperscript{h}}}{i}{h}{\textsubring{d}}} \dicPos{adv} \dicFlx{(comp tíðar\,/\addthin tíðara, sup tíðast)} \textbf{1.} \dicSynonym*{algengt} \dicDirectTranslationCS{běžně, normálně};  \dicPhraseIS{eins og títt er} \dicFlx{adv} \dicDirectTranslationCS{jak je běžné, jak tomu bývá}  \textbf{2.} \dicSynonym*{fréttnæmt} \dicDirectTranslationCS{nově, čerstvě};  \dicPhraseIS{hvað er títt?} \dicDirectTranslationCS{co je nového?};  \dicIdiom{títt}{ \dicPhraseIS{láta sér títt um e‑n}} \dicDirectTranslationCS{mít (koho) velmi rád}
\dicEntry[títuprjónn] \dicTerm{títu··prjón|n} \dicIPA{{t\smash{\textsuperscript{h}}}{i}{\textlengthmark}{\textsubring{d}}{\textscy}{p\smash{\textsuperscript{h}}}{r}{j}{ou}{\textsubring{d}}{\textsubring{n}}} \dicPos{m}[6] \dicFlx{(‑s, ‑ar)}[42] \dicDirectTranslationCS{špendlík}
\begin{xtolerant}{}{1pt}
\dicEntry[títuprjónshaus] \dicTerm{títu·prjóns··haus}\addthinS\dicIPA{{t\smash{\textsuperscript{h}}}{i}{\textlengthmark}{\textsubring{d}}{\textscy}{p\smash{\textsuperscript{h}}}{r}{j}{ou}{n}{s}{h}{\oe i}{s}}\addthinS\dicPos{m}[4]\addthinS\dicFlx{(‑s, ‑ar)}[14]\addthinS\dicDirectTranslationCS{špendlíková hlavička}
\end{xtolerant}
\dicEntry[tíu] \dicTerm{tíu} \dicsymFrequent\  \dicIPA{{t\smash{\textsuperscript{h}}}{i}{\textlengthmark}{j}{\textscy}} \dicPos{num} \dicFlx{indecl} \dicDirectTranslationCS{deset} \dicExampleIS{tíu metra breiður} \dicExampleCS{deset metrů široký};  \dicPhraseIS{kunna e‑ð upp á sína tíu fingur} \dicLangCat{přen.} \dicDirectTranslationCS{znát (co) jako své boty}
\dicEntry[tíund] \dicTerm{tí··und} \dicIPA{{t\smash{\textsuperscript{h}}}{i}{\textlengthmark}{j}{\textscy}{n}{\textsubring{d}}} \dicPos{f}[7] \dicFlx{(‑ar, ‑ir)}[1] \dicFieldCat{hist.} \dicDirectTranslationCS{desátek}
\dicEntry[tíunda] \dicTerm{tíund|a} \dicIPA{{t\smash{\textsuperscript{h}}}{i}{\textlengthmark}{j}{\textscy}{n}{\textsubring{d}}{a}} \dicPos{v}[1] \dicFlx{(‑aði)}[1] \dicFlx{acc} \dicSynonym*{skýra frá} \dicDirectTranslationCS{vylíčit, vykládat, popsat} \dicExampleIS{Hér er ekki rúm til að tíunda þau verk.} \dicExampleCS{Nedostatek prostoru nedovoluje popsat tu práci.}
\dicEntry[tíundi] \dicTerm{tíundi} \dicsymFrequent\  \dicIPA{{t\smash{\textsuperscript{h}}}{i}{\textlengthmark}{j}{\textscy}{n}{\textsubring{d}}{\textsci}} \dicPos{num}[5] \dicFlx{ord} \dicDirectTranslationCS{desátý} \dicExampleIS{í lok tíundu aldar} \dicExampleCS{na konci desátého století}
\dicEntry[tja] \dicTerm{tja} \dicIPA{{t\smash{\textsuperscript{h}}}{j}{a}{\textlengthmark}} \dicPos{inter} \dicDirectTranslationCS{hm, no}
\dicEntry[tjakka] \dicTerm{tjakk|a} \dicIPA{{t\smash{\textsuperscript{h}}}{j}{a}{h}{\r{g}}{a}} \dicPos{v}[1] \dicFlx{(‑aði)}[13] \dicFlx{acc} \dicDirectTranslationCS{zvednout\,/\addthin zvedat heverem} \dicExampleIS{tjakka e‑ð upp} \dicExampleCS{zvednout (co) heverem}
\dicEntry[tjakkur] \dicTerm{tjakk|ur} \dicIPA{{t\smash{\textsuperscript{h}}}{j}{a}{h}{\r{g}}{\textscy}{\textsubring{r}}} \dicPos{m}[6] \dicFlx{(‑s, ‑ar)}[25] \dicFieldCat{techn.} \dicDirectTranslationCS{hever, zvedák}
\dicEntry[tjald] \dicTerm{tjald} \dicsymFrequent\  \dicIPA{{t\smash{\textsuperscript{h}}}{j}{a}{l}{\textsubring{d}}} \dicPos{n}[2] \dicFlx{(‑s, tjöld)}[8] \textbf{1.} \dicSynonym*{skýli úr dúki} \dicDirectTranslationCS{stan};  \dicPhraseIS{reisa tjald} \dicDirectTranslationCS{postavit stan}  \textbf{2.} \dicSynonym*{forhengi} \dicDirectTranslationCS{opona};  \dicPhraseIS{tjaldið fellur} \dicDirectTranslationCS{opona padá}  \textbf{3.} \dicSynonym{gluggatjald} \dicDirectTranslationCS{závěs, záclona};  \dicPhraseIS{draga tjaldið\,/\addthin tjöldin frá} \dicDirectTranslationCS{roztáhnout záclony};  \dicPhraseIS{draga tjaldið\,/\addthin tjöldin fyrir} \dicDirectTranslationCS{zatáhnout záclony}  \textbf{4.} \dicSynonym*{sýningartjald} \dicDirectTranslationCS{(promítací) plátno}  \textbf{5.} \dicFieldCat{hist.} \dicSynonym{járntjald} \dicDirectTranslationCS{(železná) opona};  \dicIdiom{tjald}{ \dicPhraseIS{(e‑að gerist) á bak við tjöldin}} \dicFlx{refl} \dicLangCat{přen.} \dicDirectTranslationCS{((co) se děje) v~zákulisí, ((co) se děje) za oponou};  \dicPhraseIS{fyrir opnum tjöldum} \dicFlx{adv} \dicDirectTranslationCS{veřejně, na veřejnosti}
\dicEntry[tjalda] \dicTerm{tjald|a} \dicIPA{{t\smash{\textsuperscript{h}}}{j}{a}{l}{\textsubring{d}}{a}} \dicPos{v}[1] \dicFlx{(‑aði)}[13] \dicFlx{dat} \textbf{1.} \dicSynonym*{gista í tjaldi} \dicDirectTranslationCS{stanovat, kempovat, tábořit}  \textbf{2.} \dicFlx{dat} \dicSynonym*{klæða veggi tjöldum} \dicDirectTranslationCS{pověsit látku, ověsit látkou (zeď ap.)}
\dicEntry[tjaldborg] \dicTerm{tjald··borg} \dicIPA{{t\smash{\textsuperscript{h}}}{j}{a}{l}{\textsubring{d}}{\textsubring{b}}{\textopeno}{r}{\r{g}}} \dicPos{f}[7] \dicFlx{(‑ar, ‑ir)}[1] \dicDirectTranslationCS{stanový tábor, stanové městečko}
\dicEntry[tjalddýna] \dicTerm{tjald··dýn|a} \dicIPA{{t\smash{\textsuperscript{h}}}{j}{a}{l}{\textsubring{d}}{\textlengthmark}{i}{n}{a}} \dicPos{f}[1] \dicFlx{(‑u, ‑ur)}[7] \dicDirectTranslationCS{karimatka}
\dicEntry[tjaldhæll] \dicTerm{tjald··hæl|l} \dicIPA{{t\smash{\textsuperscript{h}}}{j}{a}{l}{\textsubring{d}}{h}{a}{i}{\textsubring{d}}{\textsubring{l}}} \dicPos{m}[6] \dicFlx{(‑s, ‑ar)}[48] \dicDirectTranslationCS{(stanový) kolík}
\dicEntry[tjaldstæði] \dicTerm{tjald··stæði} \dicIPA{{t\smash{\textsuperscript{h}}}{j}{a}{l}{\textsubring{d}}{s}{\textsubring{d}}{a}{i}{ð}{\textsci}} \dicPos{n}[2] \dicFlx{(‑s, ‑)}[14] \dicDirectTranslationCS{tábořiště, kemp, kempink}
\dicEntry[tjaldsvæði] \dicTerm{tjald··svæði} \dicIPA{{t\smash{\textsuperscript{h}}}{j}{a}{l}{\textsubring{d}}{s}{v}{a}{i}{ð}{\textsci}} \dicPos{n}[2] \dicFlx{(‑s, ‑)}[14] \dicDirectTranslationCS{kempink, kemp, tábořiště}
\dicEntry[tjaldur] \dicTerm{tjald|ur} \dicIPA{{t\smash{\textsuperscript{h}}}{j}{a}{l}{\textsubring{d}}{\textscy}{\textsubring{r}}} \dicPos{m}[6] \dicFlx{(‑s, ‑ar)}[5] \dicFieldCat{zool.} \dicDirectTranslationCS{ústřičník velký} \textit{(l.~{\textLA{Haematopus ostralegus}})}  \dicsymPhoto\ 
\dicFigure{61170.jpg}{Tjaldur}{Tjaldur - Baum Petr, Biolib, CC-BY-SA}
\dicEntry[tjaldvagn] \dicTerm{tjald··vagn} \dicIPA{{t\smash{\textsuperscript{h}}}{j}{a}{l}{\textsubring{d}}{v}{a}{\r{g}}{\textsubring{n}}} \dicPos{m}[4] \dicFlx{(‑s, ‑ar)}[4] \dicDirectTranslationCS{stanový přívěs}
\dicEntry[tjara] \dicTerm{tjara} \dicIPA{{t\smash{\textsuperscript{h}}}{j}{a}{\textlengthmark}{r}{a}} \dicPos{f}[1] \dicFlx{(tjöru)}[2] \dicDirectTranslationCS{dehet};  \dicPhraseIS{þetta er tóm tjara} \dicDirectTranslationCS{to je naprostý nesmysl}
\dicEntry[tjarga] \dicTerm{tjarg|a} \dicIPA{{t\smash{\textsuperscript{h}}}{j}{a}{r}{\r{g}}{a}} \dicPos{v}[1] \dicFlx{(‑aði)}[13] \dicFlx{acc} \dicDirectTranslationCS{(na)dehtovat}
\dicEntry[tjargaður] \dicTerm{tjargaður} \dicIPA{{t\smash{\textsuperscript{h}}}{j}{a}{r}{\r{g}}{a}{ð}{\textscy}{\textsubring{r}}} \dicPos{adj}[3] \dicFlx{(f tjörguð)}[1] \dicDirectTranslationCS{dehtový, dehtovaný}
\dicEntry[tjarnar] \dicTerm{tjarnar} \dicIPA{{t\smash{\textsuperscript{h}}}{j}{a}{r}{\textsubring{d}}{n}{a}{\textsubring{r}}} \dicPos{f} \dicFlx{sg gen} \dicLink{tjörn}
\dicEntry[tjarnarblaðka] \dicTerm{tjarnar··|blaðka} \dicIPA{{t\smash{\textsuperscript{h}}}{j}{a}{r}{\textsubring{d}}{n}{a}{r}{\textsubring{b}}{l}{a}{\texttheta}{\r{g}}{a}} \dicPos{f}[1] \dicFlx{(‑blöðku, ‑blöðkur)}[20] \dicFieldCat{bot.} \dicDirectTranslationCS{rdesno obojživelné} \textit{(l.~{\textLA{Persicaria amphibia}})}  \dicsymPhoto\ 
\dicFigure{ds_image_tjarnarbladka_0_1.jpg}{Tjarnarblaðka}{Tjarnarblaðka - Svdmolen, CC-BY-2.5}
\dicEntry[tjarnastör] \dicTerm{tjarna··|stör} \dicIPA{{t\smash{\textsuperscript{h}}}{j}{a}{r}{\textsubring{d}}{n}{a}{s}{\textsubring{d}}{\oe}{\textsubring{r}}} \dicPos{f}[7] \dicFlx{(‑starar, ‑starir)}[16] \dicFieldCat{bot.} \dicDirectTranslationCS{ostřice zobánkatá} \textit{(l.~{\textLA{Carex rostrata}})}  \dicsymPhoto\ 
\dicFigure{18508.jpg}{Tjarnastör}{Tjarnastör - Zicha Ondřej, Biolib, Copyright/CC-BY-NC}
\dicEntry[tjarnir] \dicTerm{tjarnir} \dicIPA{{t\smash{\textsuperscript{h}}}{j}{a}{r}{\textsubring{d}}{n}{\textsci}{\textsubring{r}}} \dicPos{f} \dicFlx{pl nom} \dicLink{tjörn}
\dicEntry[tjasla] \dicTerm{tjasl|a} \dicIPA{{t\smash{\textsuperscript{h}}}{j}{a}{s}{\textsubring{d}}{l}{a}} \dicPos{v}[1] \dicFlx{(‑aði)}[13] \dicPhraseIS{tjasla saman e‑u} \dicDirectTranslationCS{dát (co) dohromady, smontovat (co)} \dicExampleIS{tjasla saman gamalli myndavél} \dicExampleCS{dát dohromady starý foťák};  \dicPhraseIS{tjasla upp á e‑ð} \dicDirectTranslationCS{smontovat (co), sestavit (co)}
\dicEntry[tjá] \dicTerm{tjá} \dicsymFrequent\  \dicIPA{{t\smash{\textsuperscript{h}}}{j}{au}{\textlengthmark}} \dicPos{v}[2] \dicFlx{(‑ði, ‑ð)}[115] \dicFlx{(dat +) acc} \textbf{1.} \dicFlx{acc} \dicSynonym*{skýra frá} \dicDirectTranslationCS{vyjádřit, vyjadřovat, projevit, projevovat};  \dicPhraseIS{tjá sig um e‑ð} \dicDirectTranslationCS{vyjádřit se k~(čemu)}  \textbf{2.} \dicFlx{dat + acc} \dicDirectTranslationCS{sdělit, sdělovat, informovat} \dicExampleIS{tjá henni tilfinningar sínar} \dicExampleCS{sdělit jí své pocity}  \textbf{3.} \dicPhraseIS{það tjáir ekki að (gera e‑ð)} \dicFlx{impers} \dicDirectTranslationCS{nemá cenu (dělat (co)) (rozčilovat se ap.)}
\dicEntry[tjáning] \dicTerm{tjá··ning} \dicIPA{{t\smash{\textsuperscript{h}}}{j}{au}{\textlengthmark}{n}{i}{\ng}{\r{g}}} \dicPos{f}[4] \dicFlx{(‑ar)}[7] \textbf{1.} \dicDirectTranslationCS{vyjádření, projev (pocitů ap.)}  \textbf{2.} \dicDirectTranslationCS{vyjadřování, vyjadřovací schopnosti}
\dicEntry[tjáningarfrelsi] \dicTerm{tjáningar··frelsi} \dicIPA{{t\smash{\textsuperscript{h}}}{j}{au}{\textlengthmark}{n}{i}{\ng}{\r{g}}{a}{\textsubring{r}}{f}{r}{\textepsilon}{l}{s}{\textsci}} \dicPos{n}[2] \dicFlx{(‑s)}[20] \dicSynonym{málfrelsi} \dicDirectTranslationCS{svoboda vyjadřování\,/\addthin projevu}
\dicEntry[tjáskipti] \dicTerm{tjá··skipti} \dicIPA{{t\smash{\textsuperscript{h}}}{j}{au}{\textlengthmark}{s}{\r{\textObardotlessj}}{\textsci}{f}{\textsubring{d}}{\textsci}} \dicPos{n}[2] \dicFlx{pl}[19] \dicSynonym{boðmiðlun} \dicDirectTranslationCS{komunikace, komunikování}
\dicEntry[tjástefna] \dicTerm{tjá··stefn|a} \dicIPA{{t\smash{\textsuperscript{h}}}{j}{au}{\textlengthmark}{s}{\textsubring{d}}{\textepsilon}{\textsubring{b}}{n}{a}} \dicPos{f}[1] \dicFlx{(‑u)}[5] \dicDirectTranslationCS{expresionismus}
\dicEntry[tjóðra] \dicTerm{tjóðr|a} \dicIPA{{t\smash{\textsuperscript{h}}}{j}{ou}{ð}{r}{a}} \dicPos{v}[1] \dicFlx{(‑aði)}[1] \dicFlx{acc} \dicDirectTranslationCS{uvázat, přivázat (na provaz ap.)} \dicExampleIS{tjóðra hestinn við staurinn} \dicExampleCS{přivázat koně ke sloupku}
\dicEntry[tjóður] \dicTerm{tjóður} \dicIPA{{t\smash{\textsuperscript{h}}}{j}{ou}{\textlengthmark}{ð}{\textscy}{\textsubring{r}}} \dicPos{n}[2] \dicFlx{(‑s, ‑)}[25] \dicDirectTranslationCS{provaz} \dicIndirectTranslationCS{(na uvázání domácích zvířat)}
\dicEntry[tjón] \dicTerm{tjón} \dicIPA{{t\smash{\textsuperscript{h}}}{j}{ou}{\textlengthmark}{\textsubring{n}}} \dicPos{n}[2] \dicFlx{(‑s, ‑)}[5] \dicDirectTranslationCS{škoda, újma, ztráta} \dicExampleIS{tjón á eignum} \dicExampleCS{škoda na majetku};  \dicPhraseIS{verða fyrir tjóni} \dicDirectTranslationCS{utrpět škodu\,/\addthin újmu}
\dicEntry[tjónka] \dicTerm{tjónk|a} \dicIPA{{t\smash{\textsuperscript{h}}}{j}{ou}{\r{\ng}}{\r{g}}{a}} \dicPos{v}[1] \dicFlx{(‑aði)}[1] \dicPhraseIS{tjónka við e‑n\,/\addthin e‑ð} \dicDirectTranslationCS{dostat (koho\,/\addthin co) pod kontrolu, řídit (koho\,/\addthin co)}
\dicEntry[tjónvaldur] \dicTerm{tjón··vald|ur} \dicIPA{{t\smash{\textsuperscript{h}}}{j}{ou}{n}{v}{a}{l}{\textsubring{d}}{\textscy}{\textsubring{r}}} \dicPos{m}[6] \dicFlx{(‑s, ‑ar)}[5] \dicDirectTranslationCS{pachatel(ka), viník, vinice}
\dicEntry[tjónþoli] \dicTerm{tjón··þol|i} \dicIPA{{t\smash{\textsuperscript{h}}}{j}{ou}{n}{\texttheta}{\textopeno}{l}{\textsci}} \dicPos{m}[1] \dicFlx{(‑a, ‑ar)}[1] \dicDirectTranslationCS{poškozený, poškozená, poškozená osoba}
\dicEntry[tjöld] \dicTerm{tjöld} \dicIPA{{t\smash{\textsuperscript{h}}}{j}{\oe}{l}{\textsubring{d}}} \dicPos{n} \dicFlx{pl nom} \dicLink{tjald}
\dicEntry[tjörn] \dicTerm{tjörn} \dicIPA{{t\smash{\textsuperscript{h}}}{j}{\oe}{r}{\textsubring{d}}{\textsubring{n}}} \dicPos{f}[7] \dicFlx{(tjarnar, tjarnir)}[16] \dicDirectTranslationCS{rybník, jezero}
\dicEntry[to.] \dicTerm{to.} \dicPos{zkr} \dicPhraseIS{töluorð} \dicFieldCat{jaz.} \dicDirectTranslationCS{číslovka}
\dicEntry[tog] \dicTerm{tog} \dicIPA{{t\smash{\textsuperscript{h}}}{\textopeno}{\textlengthmark}{x}} \dicPos{n}[2] \dicFlx{(‑s)}[2] \textbf{1.} \dicSynonym*{það að toga} \dicDirectTranslationCS{tažení, táhnutí, vlečení}  \textbf{2.} \dicSynonym{kaðall} \dicDirectTranslationCS{vlečná šňůra, vlečné lano};  \dicPhraseIS{hafa e‑ð í togi} \dicDirectTranslationCS{mít (co) ve vleku, vléct (co)}  \textbf{3.} \dicSynonym{ádráttur} \dicDirectTranslationCS{zatažení\,/\addthin zátah sítě};  \dicIdiom{tog}{ \dicPhraseIS{það skiptir engum togum að (e‑að gerist)}} \dicFlx{impers} \dicLangCat{přen.} \dicDirectTranslationCS{během okamžiku (se (co) stane)}
\dicEntry[toga] \dicTerm{tog|a} \dicsymFrequent\  \dicIPA{{t\smash{\textsuperscript{h}}}{\textopeno}{\textlengthmark}{\textbabygamma}{a}} \dicPos{v}[1] \dicFlx{(‑aði)}[1] \dicFlx{acc} \textbf{1.} \dicSynonym{draga} \dicDirectTranslationCS{(vy)tahat, (vy)táhnout} \dicExampleIS{toga hana með sér út úr bænum} \dicExampleCS{vytáhnout ji sebou z~města};  \dicPhraseIS{toga í e‑ð} \dicDirectTranslationCS{tahat za (co)} \dicExampleIS{toga af öllu afli í bandið} \dicExampleCS{tahat za provaz ze všech sil}  \textbf{2.} \dicDirectTranslationCS{rybařit\,/\addthin lovit vlečnými sítěmi};  \dicIdiom{toga}[upp úr]{ \dicPhraseIS{toga e‑ð upp úr e‑m}} \dicLangCat{přen.} \dicDirectTranslationCS{vytáhnout (co) z~(koho), dostat (co) z~(koho) (informaci ap.)} \dicExampleIS{toga þetta upp úr henni} \dicExampleCS{dostat to z~ní};  \dicIdiom{togast}[á]{ \dicPhraseIS{togast á}} \dicFlx{refl} {\textbf{a.}} \dicDirectTranslationCS{přetahovat se, soutěžit v~přetahování lanem};  {\textbf{b.}} \dicLangCat{přen.} \dicDirectTranslationCS{přetahovat se, bojovat, válčit} \dicExampleIS{Þær togast á um völd.} \dicExampleCS{Přetahují se o~vládu.};  \dicIdiom{toga}{ \dicPhraseIS{hlaupa eins og fætur toga}} \dicDirectTranslationCS{utíkat jako o~život}
\dicEntry[togari] \dicTerm{tog··ar|i} \dicsymFrequent\  \dicIPA{{t\smash{\textsuperscript{h}}}{\textopeno}{\textlengthmark}{\textbabygamma}{a}{r}{\textsci}} \dicPos{m}[1] \dicFlx{(‑a, ‑ar)}[13] \dicFieldCat{nám.} \dicDirectTranslationCS{trauler, trawler} \dicExampleIS{Togarinn fiskar vel.} \dicExampleCS{Trauler dobře rybaří.}
\dicEntry[togbátur] \dicTerm{tog··bát|ur} \dicIPA{{t\smash{\textsuperscript{h}}}{\textopeno}{\textbabygamma}{\textsubring{b}}{au}{\textsubring{d}}{\textscy}{\textsubring{r}}} \dicPos{m}[6] \dicFlx{(‑s, ‑ar)}[22] \dicFieldCat{nám.} \dicDirectTranslationCS{(menší) trauler\,/\addthin trawler}
\dicEntry[togbraut] \dicTerm{tog··braut} \dicIPA{{t\smash{\textsuperscript{h}}}{\textopeno}{\textbabygamma}{\textsubring{b}}{r}{\oe i}{\textsubring{d}}} \dicPos{f}[7] \dicFlx{(‑ar, ‑ir)}[1] \dicDirectTranslationCS{lanovka, lanová dráha}
\dicEntry[togi] \dicTerm{tog|i} \dicIPA{{t\smash{\textsuperscript{h}}}{\textopeno i}{j}{\textlengthmark}{\textsci}} \dicPos{m}[1] \dicFlx{(‑a, ‑ar)}[1] \dicPhraseIS{e‑að er af ólíkum toga} \dicDirectTranslationCS{(co) je různého\,/\addthin rozdílného druhu};  \dicPhraseIS{e‑að er af sama toga} \dicDirectTranslationCS{(co) je stejného typu\,/\addthin ražení}
\dicEntry[toglyfta] \dicTerm{tog··lyft|a} \dicIPA{{t\smash{\textsuperscript{h}}}{\textopeno}{\textbabygamma}{l}{\textsci}{f}{\textsubring{d}}{a}} \dicPos{f}[1] \dicFlx{(‑u, ‑ur)}[7] \dicDirectTranslationCS{lyžařský vlek}
\dicEntry[togna] \dicTerm{togn|a} \dicIPA{{t\smash{\textsuperscript{h}}}{\textopeno}{\r{g}}{n}{a}} \dicPos{v}[1] \dicFlx{(‑aði)}[1] \textbf{1.} \dicSynonym*{teygjast} \dicDirectTranslationCS{natáhnout si, namoct si (sval ap.)}  \textbf{2.} \dicDirectTranslationCS{vymknout si, vyvrtnout si (kotník ap.)} \dicExampleIS{togna á e‑u} \dicExampleCS{vymknout si (co)}  \textbf{3.} \dicPhraseIS{það tognar úr e‑m} \dicFlx{impers} \dicDirectTranslationCS{(kdo) povyroste, (kdo) se vytáhne} \dicExampleIS{Það hefur tognað úr honum.} \dicExampleCS{Povyrostl.}
\dicEntry[tognun] \dicTerm{togn|un} \dicIPA{{t\smash{\textsuperscript{h}}}{\textopeno}{\r{g}}{n}{\textscy}{\textsubring{n}}} \dicPos{f}[7] \dicFlx{(‑unar, ‑anir)}[8] \textbf{1.} \dicDirectTranslationCS{natáhnutí, namožení (svalu ap.)}  \textbf{2.} \dicSynonym*{liðtognun} \dicDirectTranslationCS{vymknutí, vyvrtnutí (kotníku ap.)}
\dicEntry[togstreita] \dicTerm{tog··streit|a} \dicIPA{{t\smash{\textsuperscript{h}}}{\textopeno}{x}{s}{\textsubring{d}}{r}{ei}{\textsubring{d}}{a}} \dicPos{f}[1] \dicFlx{(‑u)}[5] \dicSynonym{spenna\smash{\textsuperscript{1}}} \dicDirectTranslationCS{tahanice, spor};  \dicPhraseIS{togstreita um e‑ð} \dicDirectTranslationCS{tahanice o~(co)} \dicExampleIS{togstreita um völd} \dicExampleCS{tahanice o~moc}
\dicEntry[togveiði] \dicTerm{tog··veið|i} \dicIPA{{t\smash{\textsuperscript{h}}}{\textopeno}{\textbabygamma}{v}{ei}{ð}{\textsci}} \dicPos{f}[2] \dicFlx{(‑i\,/\addthin ‑ar, ‑ar)}[3] \dicDirectTranslationCS{rybaření vlečnými sítěmi}
\dicEntry[tolla] \dicTerm{toll|a\smash{\textsuperscript{1}}} \dicIPA{{t\smash{\textsuperscript{h}}}{\textopeno}{\textsubring{d}}{l}{a}} \dicPos{v}[1] \dicFlx{(‑aði)}[1] \dicFlx{acc} \dicDirectTranslationCS{(pro)clít, proclívat, vyclít, vyclívat}
\dicEntry[tolla] \dicTerm{toll|a\smash{\textsuperscript{2}}} \dicIPA{{t\smash{\textsuperscript{h}}}{\textopeno}{\textsubring{d}}{l}{a}} \dicPos{v}[2] \dicFlx{(‑di, ‑að)}[166] \dicDirectTranslationCS{(vy)držet, (při)lnout} \dicExampleIS{Myndin tollir ekki á veggnum.} \dicExampleCS{Obrázek nedrží na stěně.};  \dicPhraseIS{e‑r tollir hvergi} \dicLangCat{přen.} \dicDirectTranslationCS{(kdo) nikde nevydrží (v~práci ap.)};  \dicPhraseIS{tolla í tískunni} \dicDirectTranslationCS{držet se módy}
\dicEntry[tollafgreiða] \dicTerm{toll··af·grei|ða} \dicIPA{{t\smash{\textsuperscript{h}}}{\textopeno}{\textsubring{d}}{l}{a}{v}{\r{g}}{r}{ei}{ð}{a}} \dicPos{v}[2] \dicFlx{(‑ddi, ‑tt)}[167] \dicFlx{acc} \dicDirectTranslationCS{proclít, proclívat} \dicExampleIS{tollafgreiða vörur} \dicExampleCS{proclít zboží}
\dicEntry[tollafgreiðsla] \dicTerm{toll··af·greiðsl|a} \dicIPA{{t\smash{\textsuperscript{h}}}{\textopeno}{\textsubring{d}}{l}{a}{v}{\r{g}}{r}{ei}{ð}{s}{\textsubring{d}}{l}{a}} \dicPos{f}[1] \dicFlx{(‑u, ‑ur)}[13] \textbf{1.} \dicSynonym*{tollmeðferð} \dicDirectTranslationCS{proclení, celní odbavení}  \textbf{2.} \dicSynonym*{tollskrifstofa} \dicDirectTranslationCS{celní úřad}
\dicEntry[tollari] \dicTerm{toll··ar|i} \dicIPA{{t\smash{\textsuperscript{h}}}{\textopeno}{\textsubring{d}}{l}{a}{r}{\textsci}} \dicPos{m}[1] \dicFlx{(‑a, ‑ar)}[13] \dicLangCat{hovor.} \dicDirectTranslationCS{celník}
\dicEntry[tolleftirlit] \dicTerm{toll··eftir·lit} \dicIPA{{t\smash{\textsuperscript{h}}}{\textopeno}{\textsubring{d}}{l}{\textepsilon}{f}{\textsubring{d}}{\textsci}{r}{l}{\textsci}{\textsubring{d}}} \dicPos{n}[2] \dicFlx{(‑s)}[2] \dicSynonym{tollgæsla} \dicDirectTranslationCS{celní kontrola}
\dicEntry[tollfrjáls] \dicTerm{toll··frjáls} \dicIPA{{t\smash{\textsuperscript{h}}}{\textopeno}{\textsubring{d}}{\textsubring{l}}{f}{r}{j}{au}{l}{s}} \dicPos{adj}[5]\dicFlx{}[-1] \dicDirectTranslationCS{bezcelní, prostý cla} \dicExampleIS{tollfrjáls vara} \dicExampleCS{výrobek prostý cla}
\dicEntry[tollgæsla] \dicTerm{toll··gæsl|a} \dicIPA{{t\smash{\textsuperscript{h}}}{\textopeno}{\textsubring{d}}{\textsubring{l}}{\r{\textObardotlessj}}{a}{i}{s}{\textsubring{d}}{l}{a}} \dicPos{f}[1] \dicFlx{(‑u)}[5] \dicDirectTranslationCS{celní kontrola}
\dicEntry[tollleit] \dicTerm{toll··leit} \dicIPA{{t\smash{\textsuperscript{h}}}{\textopeno}{\textsubring{d}}{l}{ei}{\textsubring{d}}} \dicPos{f}[7] \dicFlx{(‑ar)}[3] \dicDirectTranslationCS{celní prohlídka}
\dicEntry[tollmúr] \dicTerm{toll··múr} \dicIPA{{t\smash{\textsuperscript{h}}}{\textopeno}{\textsubring{d}}{l}{m}{u}{\textsubring{r}}} \dicPos{m}[4] \dicFlx{(‑s, ‑ar)}[14] \dicFieldCat{ekon.} \dicDirectTranslationCS{celní bariéra}
\dicEntry[tollsamningur] \dicTerm{toll··samn·ing|ur} \dicIPA{{t\smash{\textsuperscript{h}}}{\textopeno}{\textsubring{d}}{\textsubring{l}}{s}{a}{m}{n}{i}{\ng}{\r{g}}{\textscy}{\textsubring{r}}} \dicPos{m}[6] \dicFlx{(‑s, ‑ar)}[8] \dicDirectTranslationCS{celní dohoda}
\dicEntry[tollskoða] \dicTerm{toll··skoð|a} \dicIPA{{t\smash{\textsuperscript{h}}}{\textopeno}{\textsubring{d}}{\textsubring{l}}{s}{\r{g}}{\textopeno}{ð}{a}} \dicPos{v}[1] \dicFlx{(‑aði)}[1] \dicFlx{acc} \dicDirectTranslationCS{provést\,/\addthin provádět celní kontrolu}
\dicEntry[tollskoðun] \dicTerm{toll··skoð|un} \dicIPA{{t\smash{\textsuperscript{h}}}{\textopeno}{\textsubring{d}}{\textsubring{l}}{s}{\r{g}}{\textopeno}{ð}{\textscy}{\textsubring{n}}} \dicPos{f}[7] \dicFlx{(‑unar)}[9] \dicDirectTranslationCS{celní prohlídka\,/\addthin kontrola}
\dicEntry[tollskrá] \dicTerm{toll··skrá} \dicIPA{{t\smash{\textsuperscript{h}}}{\textopeno}{\textsubring{d}}{\textsubring{l}}{s}{\r{g}}{r}{au}} \dicPos{f}[4] \dicFlx{(‑r\,/\addthin ‑ar, ‑r)}[21] \dicDirectTranslationCS{celní sazebník}
\dicEntry[tollskyldur] \dicTerm{toll··skyldur} \dicIPA{{t\smash{\textsuperscript{h}}}{\textopeno}{\textsubring{d}}{\textsubring{l}}{s}{\r{\textObardotlessj}}{\textsci}{l}{\textsubring{d}}{\textscy}{\textsubring{r}}} \dicPos{adj}[2]\dicFlx{}[-17] \dicDirectTranslationCS{podléhající clu} \dicExampleIS{tollskyld vara} \dicExampleCS{výrobek podléhající clu}
\dicEntry[tollskýrsla] \dicTerm{toll··skýrsl|a} \dicIPA{{t\smash{\textsuperscript{h}}}{\textopeno}{\textsubring{d}}{\textsubring{l}}{s}{\r{\textObardotlessj}}{i}{\textsubring{r}}{s}{\textsubring{d}}{l}{a}} \dicPos{f}[1] \dicFlx{(‑u, ‑ur)}[13] \dicDirectTranslationCS{celní prohlášení\,/\addthin deklarace}
\dicEntry[tollstöð] \dicTerm{toll··stöð} \dicIPA{{t\smash{\textsuperscript{h}}}{\textopeno}{\textsubring{d}}{\textsubring{l}}{s}{\textsubring{d}}{\oe}{\texttheta}} \dicPos{f}[6] \dicFlx{(‑var, ‑var)}[1] \dicDirectTranslationCS{celnice}
\dicEntry[tollsvik] \dicTerm{toll··svik} \dicIPA{{t\smash{\textsuperscript{h}}}{\textopeno}{\textsubring{d}}{\textsubring{l}}{s}{v}{\textsci}{\r{g}}} \dicPos{n}[2] \dicFlx{pl}[1] \dicDirectTranslationCS{celní podvod}
\dicEntry[tollsvæði] \dicTerm{toll··svæði} \dicIPA{{t\smash{\textsuperscript{h}}}{\textopeno}{\textsubring{d}}{\textsubring{l}}{s}{v}{a}{i}{ð}{\textsci}} \dicPos{n}[2] \dicFlx{(‑s, ‑)}[14] \dicDirectTranslationCS{celní prostor}
\dicEntry[tollur] \dicTerm{toll|ur} \dicIPA{{t\smash{\textsuperscript{h}}}{\textopeno}{\textsubring{d}}{l}{\textscy}{\textsubring{r}}} \dicPos{m}[6] \dicFlx{(‑s, ‑ar)}[4] \textbf{1.} \dicSynonym{innflutningsgjald} \dicDirectTranslationCS{clo, celní poplatek} \dicExampleIS{greiða toll af e‑u} \dicExampleCS{zaplatit z~(čeho) clo};  \dicPhraseIS{leggja toll á e‑u} \dicSynonym{tolla\smash{\textsuperscript{1}}} \dicDirectTranslationCS{uvalit na (co) clo}  \textbf{2.} \dicSynonym{afnotagjald} \dicDirectTranslationCS{poplatek za užívání};  \dicPhraseIS{e‑að tekur toll af e‑u} \dicLangCat{přen.} \dicDirectTranslationCS{(co) si vybírá daň na (čem)}
\dicEntry[tollvara] \dicTerm{toll··|vara} \dicIPA{{t\smash{\textsuperscript{h}}}{\textopeno}{\textsubring{d}}{l}{v}{a}{r}{a}} \dicPos{f}[1] \dicFlx{(‑vöru, ‑vörur)}[8] \dicDirectTranslationCS{zboží podléhající clu}
\dicEntry[tollvörður] \dicTerm{toll··|vörður}\dicTerm{, tollþjónn} \dicIPA{{t\smash{\textsuperscript{h}}}\-{\textopeno}\-{\textsubring{d}}\-{l}\-{v}\-{\oe}\-{r}\-{ð}\-{\textscy}\-{\textsubring{r}}\-} \dicPos{m}[11] \dicFlx{(‑varðar, ‑verðir)}[5] \dicDirectTranslationCS{celník, celnice}
\dicEntry[tollvörugeymsla] \dicTerm{toll·vöru··geymsl|a} \dicIPA{{t\smash{\textsuperscript{h}}}{\textopeno}{\textsubring{d}}{l}{v}{\oe}{r}{\textscy}{\r{\textObardotlessj}}{ei}{m}{s}{\textsubring{d}}{l}{a}} \dicPos{f}[1] \dicFlx{(‑u, ‑ur)}[13] \dicDirectTranslationCS{celní skladiště}
\dicEntry[tollþjónn] \dicTerm{toll··þjón|n} \dicIPA{{t\smash{\textsuperscript{h}}}{\textopeno}{\textsubring{d}}{\textsubring{l}}{\texttheta}{j}{ou}{\textsubring{d}}{\textsubring{n}}} \dicPos{m}[6] \dicFlx{(‑s, ‑ar)}[42] \dicLink{tollvörður}
\dicEntry[tombóla] \dicTerm{tomból|a} \dicIPA{{t\smash{\textsuperscript{h}}}{\textopeno}{m}{\textsubring{b}}{ou}{l}{a}} \dicPos{f}[1] \dicFlx{(‑u, ‑ur)}[19] \dicDirectTranslationCS{tombola}
\dicEntry[tomma] \dicTerm{tomm|a} \dicIPA{{t\smash{\textsuperscript{h}}}{\textopeno}{m}{\textlengthmark}{a}} \dicPos{f}[1] \dicFlx{(‑u, ‑ur)}[7] \dicDirectTranslationCS{palec} \dicIndirectTranslationCS{(délková míra)}
\dicEntry[tommustokkur] \dicTerm{tommu··stokk|ur} \dicIPA{{t\smash{\textsuperscript{h}}}{\textopeno}{m}{\textlengthmark}{\textscy}{s}{\textsubring{d}}{\textopeno}{h}{\r{g}}{\textscy}{\textsubring{r}}} \dicPos{m}[6] \dicFlx{(‑s, ‑ar)}[22] \dicDirectTranslationCS{skládací metr}
\dicEntry[Tonga] \dicTerm{Tonga} \dicIPA{{t\smash{\textsuperscript{h}}}{\textopeno}{\ng}{\r{g}}{a}} \dicPos{n}[4] \dicFlx{indecl}[2] \dicFieldCat{geog.} \dicDirectTranslationCS{Tonga}
\dicEntry[tongskur] \dicTerm{tongskur}\dicTerm{, tongverskur} \dicIPA{{t\smash{\textsuperscript{h}}}\-{\textopeno}\-{\ng}\-{s}\-{\r{g}}\-{\textscy}\-{\textsubring{r}}\-} \dicPos{adj}[1]\dicFlx{}[-6] \dicDirectTranslationCS{tonžský}
\dicEntry[Tongverji] \dicTerm{Tong··verj|i} \dicIPA{{t\smash{\textsuperscript{h}}}{\textopeno}{\ng}{\r{g}}{v}{\textepsilon}{r}{j}{\textsci}} \dicPos{m}[1] \dicFlx{(‑a, ‑ar)}[1] \dicDirectTranslationCS{Tonžan(ka)}
\dicEntry[tongverskur] \dicTerm{tong··verskur} \dicIPA{{t\smash{\textsuperscript{h}}}{\textopeno}{\ng}{\r{g}}{v}{\textepsilon}{\textsubring{r}}{s}{\r{g}}{\textscy}{\textsubring{r}}} \dicPos{adj}[1]\dicFlx{}[-6] \dicLink{tongskur}
\dicEntry[tonn] \dicTerm{tonn} \dicsymFrequent\  \dicIPA{{t\smash{\textsuperscript{h}}}{\textopeno}{n}{\textlengthmark}} \dicPos{n}[2] \dicFlx{(‑s, ‑)}[5] \dicDirectTranslationCS{tuna} \dicIndirectTranslationCS{(jednotka hmotnosti)} \dicExampleIS{tveggja tonna ballest} \dicExampleCS{dvoutunový balast}
\dicEntry[toppa] \dicTerm{topp|a} \dicIPA{{t\smash{\textsuperscript{h}}}{\textopeno}{h}{\textsubring{b}}{a}} \dicPos{v}[1] \dicFlx{(‑aði)}[1] \dicFlx{acc} \dicLangCat{hovor.} \dicDirectTranslationCS{překonat, trumfnout}
\dicEntry[toppfígúra] \dicTerm{topp··fígúr|a} \dicIPA{{t\smash{\textsuperscript{h}}}{\textopeno}{h}{\textsubring{b}}{f}{i}{\textbabygamma}{u}{r}{a}} \dicPos{f}[1] \dicFlx{(‑u, ‑ur)}[7] \dicDirectTranslationCS{loutka, figurka (nastrčená do politiky ap.)}
\dicEntry[toppfundur] \dicTerm{topp··fund|ur} \dicIPA{{t\smash{\textsuperscript{h}}}{\textopeno}{h}{\textsubring{b}}{f}{\textscy}{n}{\textsubring{d}}{\textscy}{\textsubring{r}}} \dicPos{m}[10] \dicFlx{(‑ar, ‑ir)}[4] \dicFieldCat{pol.} \dicDirectTranslationCS{schůzka na nejvyšší úrovni, summit}
\dicEntry[toppgrind] \dicTerm{topp··grind} \dicIPA{{t\smash{\textsuperscript{h}}}{\textopeno}{h}{\textsubring{b}}{\r{g}}{r}{\textsci}{n}{\textsubring{d}}} \dicPos{f}[8] \dicFlx{(‑ar, ‑ur)}[1] \dicDirectTranslationCS{zahrádka} \dicIndirectTranslationCS{(auta)}
\dicEntry[toppkál] \dicTerm{topp··kál} \dicIPA{{t\smash{\textsuperscript{h}}}{\textopeno}{h}{\textsubring{b}}{k\smash{\textsuperscript{h}}}{au}{\textsubring{l}}} \dicPos{n}[2] \dicFlx{(‑s)}[2] \dicFieldCat{bot.} \dicDirectTranslationCS{hlávkové zelí} \textit{(l.~{\textLA{Brassica oleracea var. capitata}})}
\dicEntry[toppskarfur] \dicTerm{topp··skarf|ur} \dicIPA{{t\smash{\textsuperscript{h}}}{\textopeno}{h}{\textsubring{b}}{s}{\r{g}}{a}{r}{v}{\textscy}{\textsubring{r}}} \dicPos{m}[6] \dicFlx{(‑s, ‑ar)}[5] \dicFieldCat{zool.} \dicDirectTranslationCS{kormorán chocholatý} \textit{(l.~{\textLA{Phalacrocorax aristotelis}})}  \dicsymPhoto\ 
\dicFigure{ds_image_toppskarfur_0_1.jpg}{Toppskarfur}{Toppskarfur - MPF, CC BY-SA 3.0}
\dicEntry[toppur] \dicTerm{topp|ur} \dicsymFrequent\  \dicIPA{{t\smash{\textsuperscript{h}}}{\textopeno}{h}{\textsubring{b}}{\textscy}{\textsubring{r}}} \dicPos{m}[6] \dicFlx{(‑s, ‑ar)}[10] \textbf{1.} \dicSynonym{hártoppur} \dicDirectTranslationCS{ofina}  \textbf{2.} \dicSynonym{tindur} \dicDirectTranslationCS{vrchol, špice (stromu ap.)} \dicExampleIS{toppurinn á trénu} \dicExampleCS{vrchol stromu}  \textbf{3.} \dicDirectTranslationCS{střecha} \dicIndirectTranslationCS{(auta)}  \textbf{4.} \dicDirectTranslationCS{top} \dicIndirectTranslationCS{(horní část ženského oblečení)};  \dicIdiom{toppur}{ \dicPhraseIS{vera á toppnum}} \dicLangCat{přen.} \dicDirectTranslationCS{být na vrcholu\,/\addthin špici}
\dicEntry[toppönd] \dicTerm{topp··|önd} \dicIPA{{t\smash{\textsuperscript{h}}}{\textopeno}{h}{\textsubring{b}}{\oe}{n}{\textsubring{d}}} \dicPos{f}[8] \dicFlx{(‑andar, ‑endur)}[2] \dicFieldCat{zool.} \dicDirectTranslationCS{morčák prostřední} \textit{(l.~{\textLA{Mergus serrator}})}  \dicsymPhoto\ 
\dicFigure{ds_image_toppond_0_2.jpg}{Toppönd}{Toppönd - Peter Massas, CC BY-SA 2.0}
\dicEntry[tor] \dicTerm{tor-} \dicIPA{{t\smash{\textsuperscript{h}}}{\textopeno}{r}} \dicPos{predp} \dicDirectTranslationCS{obtížně, těžko} \dicExampleIS{torskilinn} \dicExampleCS{obtížně srozumitelný}
\dicEntry[torf] \dicTerm{torf} \dicIPA{{t\smash{\textsuperscript{h}}}{\textopeno}{r}{f}} \dicPos{n}[2] \dicFlx{(‑s)}[2] \textbf{1.} \dicFieldCat{geol.} \dicDirectTranslationCS{rašelina}  \textbf{2.} \dicSynonym*{túnþaka} \dicDirectTranslationCS{drn}  \textbf{3.} \dicSynonym*{erfiður texti} \dicDirectTranslationCS{obtížný text}
\dicEntry[torfa] \dicTerm{torf|a} \dicIPA{{t\smash{\textsuperscript{h}}}{\textopeno}{r}{v}{a}} \dicPos{f}[1] \dicFlx{(‑u, ‑ur)}[19] \textbf{1.} \dicSynonym*{grastorfa} \dicDirectTranslationCS{travnatý pruh, drnový koberec}  \textbf{2.} \dicSynonym*{fiskatorfa} \dicDirectTranslationCS{hejno\,/\addthin houf (ryb)}  \textbf{3.} \dicSynonym*{húsaþyrping} \dicDirectTranslationCS{shluk stavení};  \dicIdiom{torfa}{ \dicPhraseIS{vera gróinn við torfuna}} \dicLangCat{přen.} \dicDirectTranslationCS{být přirostlý k~rodné hroudě}; { \dicPhraseIS{vera kominn undir græna torfu}} \dicLangCat{přen.} \dicDirectTranslationCS{ležet\,/\addthin být pod drnem}
\dicEntry[torfbær] \dicTerm{torf··bæ|r} \dicIPA{{t\smash{\textsuperscript{h}}}{\textopeno}{r}{v}{\textsubring{b}}{a}{i}{\textsubring{r}}} \dicPos{m}[10] \dicFlx{(‑jar, ‑ir)}[25] \dicDirectTranslationCS{drnový dům, dům s~drnovou střechou} \dicIndirectTranslationCS{(tradiční islandský dům)}
\dicEntry[torfenginn] \dicTerm{tor··fenginn} \dicIPA{{t\smash{\textsuperscript{h}}}{\textopeno}{\textsubring{r}}{f}{ei}{\textltailn}{\r{\textObardotlessj}}{\textsci}{\textsubring{n}}} \dicPos{adj}[6]\dicFlx{}[-2] \dicDirectTranslationCS{těžce dosažitelný, obtížně získatelný}
\dicEntry[torfær] \dicTerm{tor··fær} \dicIPA{{t\smash{\textsuperscript{h}}}{\textopeno}{\textsubring{r}}{f}{a}{i}{\textsubring{r}}} \dicPos{adj}[5] \dicFlx{(f ‑)}[8] \dicDirectTranslationCS{obtížně schůdný\,/\addthin sjízdný}
\dicEntry[torfæra] \dicTerm{tor··fær|a}\dicTerm{, torfæri} \dicIPA{{t\smash{\textsuperscript{h}}}\-{\textopeno}\-{\textsubring{r}}\-{f}\-{a}\-{i}\-{r}\-{a}\-} \dicPos{f}[1] \dicFlx{(‑u, ‑ur)}[19] \dicDirectTranslationCS{obtížně sjízdný terén}
\dicEntry[torfæri] \dicTerm{tor··færi} \dicIPA{{t\smash{\textsuperscript{h}}}{\textopeno}{\textsubring{r}}{f}{a}{i}{r}{\textsci}} \dicPos{n}[2] \dicFlx{(‑s, ‑)}[14] \dicLink{torfæra}
\dicEntry[torfærubíll] \dicTerm{tor·færu··bíl|l} \dicIPA{{t\smash{\textsuperscript{h}}}{\textopeno}{\textsubring{r}}{f}{a}{i}{r}{\textscy}{\textsubring{b}}{i}{\textsubring{d}}{\textsubring{l}}} \dicPos{m}[6] \dicFlx{(‑s, ‑ar)}[48] \dicDirectTranslationCS{terénní vozidlo}
\dicEntry[torfæruhjól] \dicTerm{tor·færu··hjól} \dicIPA{{t\smash{\textsuperscript{h}}}{\textopeno}{\textsubring{r}}{f}{a}{i}{r}{\textscy}{\c{c}}{ou}{\textsubring{l}}} \dicPos{n}[2] \dicFlx{(‑s, ‑)}[5] \dicDirectTranslationCS{terénní motocykl}
\dicEntry[torg] \dicTerm{torg} \dicsymFrequent\  \dicIPA{{t\smash{\textsuperscript{h}}}{\textopeno}{r}{\r{g}}} \dicPos{n}[2] \dicFlx{(‑s, ‑)}[5] \dicDirectTranslationCS{náměstí, plac, rynek} \dicExampleIS{Minnisvarði var reistur á torginu.} \dicExampleCS{Pomník byl vztyčen na náměstí.}
\dicEntry[torga] \dicTerm{torg|a} \dicIPA{{t\smash{\textsuperscript{h}}}{\textopeno}{r}{\r{g}}{a}} \dicPos{v}[1] \dicFlx{(‑aði)}[1] \dicFlx{dat} \dicSynonym*{éta upp} \dicDirectTranslationCS{sníst, spořádat, pozřít}
\dicEntry[torkennilegur] \dicTerm{tor·kenni··legur} \dicIPA{{t\smash{\textsuperscript{h}}}{\textopeno}{\textsubring{r}}{c\smash{\textsuperscript{h}}}{\textepsilon}{n}{\textsci}{l}{\textepsilon}{\textbabygamma}{\textscy}{\textsubring{r}}} \dicPos{adj}[1]\dicFlx{}[-8] \dicDirectTranslationCS{těžko\,/\addthin obtížně (roz)poznatelný}
\dicEntry[torlæsi] \dicTerm{tor··læsi} \dicIPA{{t\smash{\textsuperscript{h}}}{\textopeno}{r}{l}{a}{i}{s}{\textsci}} \dicPos{n}[2] \dicFlx{(‑s)}[20] \dicDirectTranslationCS{obtížná\,/\addthin slabá čitelnost}
\dicEntry[tormeltur] \dicTerm{tor··meltur} \dicIPA{{t\smash{\textsuperscript{h}}}{\textopeno}{r}{m}{\textepsilon}{\textsubring{l}}{\textsubring{d}}{\textscy}{\textsubring{r}}} \dicPos{adj}[1]\dicFlx{}[-10] \dicDirectTranslationCS{těžce\,/\addthin obtížně stravitelný}
\dicEntry[tormerki] \dicTerm{tor··merki} \dicIPA{{t\smash{\textsuperscript{h}}}{\textopeno}{r}{m}{\textepsilon}{\textsubring{r}}{\r{\textObardotlessj}}{\textsci}} \dicPos{n}[2] \dicFlx{(‑s, ‑)}[16] \dicSynonym*{erfiðleikar} \dicDirectTranslationCS{těžkosti, nesnáze};  \dicPhraseIS{telja tormerki á e‑u} \dicDirectTranslationCS{vypočítávat záporné stránky (čeho)};  \dicPhraseIS{það eru engin tormerki á e‑u} \dicDirectTranslationCS{nejsou žádné překážky v~(čem) (uzavření manželství ap.)}
\dicEntry[tornæmi] \dicTerm{tor··næmi} \dicIPA{{t\smash{\textsuperscript{h}}}{\textopeno}{r}{n}{a}{i}{m}{\textsci}} \dicPos{n}[2] \dicFlx{(‑s)}[20] \dicDirectTranslationCS{nechápavost, tupost}
\dicEntry[tornæmur] \dicTerm{tor··næmur} \dicIPA{{t\smash{\textsuperscript{h}}}{\textopeno}{r}{n}{a}{i}{m}{\textscy}{\textsubring{r}}} \dicPos{adj}[1]\dicFlx{}[-1] \dicDirectTranslationCS{mající potíže s~učením, nechápavý}
\dicEntry[torráðinn] \dicTerm{tor··ráðinn} \dicIPA{{t\smash{\textsuperscript{h}}}{\textopeno}{r}{\textlengthmark}{au}{ð}{\textsci}{\textsubring{n}}} \dicPos{adj}[6]\dicFlx{}[-2] \dicDirectTranslationCS{obtížně vyřešitelný, obtížný, těžký} \dicExampleIS{torráðin gáta} \dicExampleCS{obtížná hádanka}
\dicEntry[torskilinn] \dicTerm{tor··skilinn} \dicIPA{{t\smash{\textsuperscript{h}}}{\textopeno}{\textsubring{r}}{s}{\r{\textObardotlessj}}{\textsci}{l}{\textsci}{\textsubring{n}}} \dicPos{adj}[6]\dicFlx{}[-7] \dicDirectTranslationCS{těžko\,/\addthin obtížně pochopitelný, nesrozumitelný} \dicExampleIS{torskilið orð} \dicExampleCS{nesrozumitelné slovo}
\dicEntry[torsóttur] \dicTerm{tor··sóttur} \dicIPA{{t\smash{\textsuperscript{h}}}{\textopeno}{\textsubring{r}}{s}{ou}{h}{\textsubring{d}}{\textscy}{\textsubring{r}}} \dicPos{adj}[1]\dicFlx{}[-10] \dicDirectTranslationCS{svízelný, nesnadný} \dicExampleIS{Námið er mér torsótt.} \dicExampleCS{Kurz je pro mě svízelný.}
\dicEntry[tortíma] \dicTerm{tor··tím|a} \dicIPA{{t\smash{\textsuperscript{h}}}{\textopeno}{\textsubring{r}}{t\smash{\textsuperscript{h}}}{i}{m}{a}} \dicPos{v}[2] \dicFlx{(‑di, ‑t)}[148] \dicFlx{dat} \dicDirectTranslationCS{(z)ničit, (z)likvidovat, zneškodnit, zneškodňovat} \dicExampleIS{tortíma e‑m\,/\addthin e‑u} \dicExampleCS{zlikvidovat (koho\,/\addthin co)};  \dicIdiom{tortímast}{ \dicPhraseIS{tortímast}} \dicFlx{refl} \dicDirectTranslationCS{zničit se, zlikvidovat se}
\dicEntry[tortíming] \dicTerm{tor··tím·ing} \dicIPA{{t\smash{\textsuperscript{h}}}{\textopeno}{\textsubring{r}}{t\smash{\textsuperscript{h}}}{i}{m}{i}{\ng}{\r{g}}} \dicPos{f}[4] \dicFlx{(‑ar)}[7] \dicSynonym{eyðilegging} \dicDirectTranslationCS{zničení, zlikvidování, zneškodnění}
\dicEntry[tortryggilegur] \dicTerm{tor·tryggi··legur} \dicIPA{{t\smash{\textsuperscript{h}}}{\textopeno}{\textsubring{r}}{t\smash{\textsuperscript{h}}}{r}{\textsci}{\r{\textObardotlessj}}{\textsci}{l}{\textepsilon}{\textbabygamma}{\textscy}{\textsubring{r}}} \dicPos{adj}[1]\dicFlx{}[-8] \dicSynonym{grunsamlegur} \dicDirectTranslationCS{podezřelý, budící podezření, pofidérní} \dicExampleIS{Útlitið gerir hann tortryggilegan.} \dicExampleCS{Jeho vzhled budí podezření.}
\dicEntry[tortrygginn] \dicTerm{tor··trygginn} \dicIPA{{t\smash{\textsuperscript{h}}}{\textopeno}{\textsubring{r}}{t\smash{\textsuperscript{h}}}{r}{\textsci}{\r{\textObardotlessj}}{\textsci}{\textsubring{n}}} \dicPos{adj}[6]\dicFlx{}[-2] \dicDirectTranslationCS{podezíravý, podezřívavý, nedůvěřivý} \dicExampleIS{vera tortrygginn gagnvart nýjum félögum} \dicExampleCS{být podezřívavý vůči novým členům}
\dicEntry[tortryggja] \dicTerm{tor··trygg|ja} \dicIPA{{t\smash{\textsuperscript{h}}}{\textopeno}{\textsubring{r}}{t\smash{\textsuperscript{h}}}{r}{\textsci}{\r{\textObardotlessj}}{a}} \dicPos{v}[2] \dicFlx{(‑ði, ‑t)}[89] \dicFlx{acc} \dicDirectTranslationCS{podezírat, podezřívat, nedůvěřovat} \dicExampleIS{tortryggja e‑n} \dicExampleCS{podezřívat (koho)}
\dicEntry[tortryggni] \dicTerm{tor··tryggn|i} \dicIPA{{t\smash{\textsuperscript{h}}}{\textopeno}{\textsubring{r}}{t\smash{\textsuperscript{h}}}{r}{\textsci}{\r{g}}{n}{\textsci}} \dicPos{f}[3] \dicFlx{(‑i)}[3] \dicDirectTranslationCS{podezíravost, podezřívavost, nedůvěra, skepse} \dicExampleIS{tortryggni í garð hans} \dicExampleCS{podezřívavost vůči němu}
\dicEntry[torvelda] \dicTerm{tor··veld|a} \dicIPA{{t\smash{\textsuperscript{h}}}{\textopeno}{r}{v}{\textepsilon}{l}{\textsubring{d}}{a}} \dicPos{v}[1] \dicFlx{(‑aði)}[1] \dicFlx{(dat +) acc} \dicDirectTranslationCS{ztížit, ztěžovat, znesnadnit, znesnadňovat, (z)komplikovat} \dicExampleIS{torvelda honum verkið} \dicExampleCS{ztížit mu práci}
\dicEntry[torveldur] \dicTerm{tor··veldur} \dicIPA{{t\smash{\textsuperscript{h}}}{\textopeno}{r}{v}{\textepsilon}{l}{\textsubring{d}}{\textscy}{\textsubring{r}}} \dicPos{adj}[2]\dicFlx{}[-14] \dicSynonym{örðugur} \dicDirectTranslationCS{ztížený, nesnadný} \dicExampleIS{torveldur í framkvæmd} \dicExampleCS{ztížený v~realizaci}
\dicEntry[tosa] \dicTerm{tos|a} \dicIPA{{t\smash{\textsuperscript{h}}}{\textopeno}{\textlengthmark}{s}{a}} \dicPos{v}[1] \dicFlx{(‑aði)}[1] \dicFlx{acc} \dicSynonym{draga} \dicDirectTranslationCS{vléct, vláčet};  \dicIdiom{tosast}{ \dicPhraseIS{tosast}} \dicFlx{refl} \dicDirectTranslationCS{vléct se, prodlužovat se}
\dicEntry[tossi] \dicTerm{toss|i} \dicIPA{{t\smash{\textsuperscript{h}}}{\textopeno}{s}{\textlengthmark}{\textsci}} \dicPos{m}[1] \dicFlx{(‑a, ‑ar)}[1] \dicDirectTranslationCS{nechápavý\,/\addthin tupý žák, nechápavá\,/\addthin tupá žačka\,/\addthin žákyně}
\dicEntry[tota] \dicTerm{tot|a} \dicIPA{{t\smash{\textsuperscript{h}}}{\textopeno}{\textlengthmark}{\textsubring{d}}{a}} \dicPos{f}[1] \dicFlx{(‑u, ‑ur)}[19] \dicDirectTranslationCS{špička} \dicExampleIS{tota á sokk} \dicExampleCS{špička ponožky}
\dicEntry[tott] \dicTerm{tott} \dicIPA{{t\smash{\textsuperscript{h}}}{\textopeno}{h}{\textsubring{d}}} \dicPos{n}[2] \dicFlx{(‑s)}[2] \textbf{1.} \dicDirectTranslationCS{(na)sání}  \textbf{2.} \dicDirectTranslationCS{bafání, dýmání}  \textbf{3.} \dicLangCat{hrub.} \dicDirectTranslationCS{(vy)kouření} \dicIndirectTranslationCS{(orální sex)}
\dicEntry[totta] \dicTerm{tott|a} \dicIPA{{t\smash{\textsuperscript{h}}}{\textopeno}{h}{\textsubring{d}}{a}} \dicPos{v}[1] \dicFlx{(‑aði)}[1] \dicFlx{acc} \textbf{1.} \dicSynonym{sjúga} \dicDirectTranslationCS{sát, nasávat} \dicExampleIS{totta spena} \dicExampleCS{sát cecík}  \textbf{2.} \dicSynonym{reykja} \dicDirectTranslationCS{bafat, dýmat} \dicExampleIS{totta pípu sína} \dicExampleCS{bafat dýmku}  \textbf{3.} \dicLangCat{hrub.} \dicDirectTranslationCS{(vy)kouřit} \dicIndirectTranslationCS{(orálně uspokojovat muže)}
\dicEntry[tó] \dicTerm{tó\smash{\textsuperscript{1}}} \dicIPA{{t\smash{\textsuperscript{h}}}{ou}{\textlengthmark}} \dicPos{f}[4] \dicFlx{(‑r\,/\addthin ‑ar, ‑r\,/\addthin tær)}[28] \dicSynonym*{grastoppur} \dicDirectTranslationCS{travnaté místo, trávník} \dicExampleIS{mosató} \dicExampleCS{trávník z~mechu}
\dicEntry[tó] \dicTerm{tó\smash{\textsuperscript{2}}} \dicIPA{{t\smash{\textsuperscript{h}}}{ou}{\textlengthmark}} \dicPos{n}[2] \dicFlx{(‑s)}[2] \textbf{1.} \dicSynonym{ull} \dicDirectTranslationCS{vlna} \dicExampleIS{tóvinna} \dicExampleCS{práce s~vlnou}  \textbf{2.} \dicSynonym{taug} \dicDirectTranslationCS{lano, provaz}
\dicEntry[tóbak] \dicTerm{tóbak} \dicIPA{{t\smash{\textsuperscript{h}}}{ou}{\textlengthmark}{\textsubring{b}}{a}{\r{g}}} \dicPos{n}[2] \dicFlx{(‑s)}[2] \dicDirectTranslationCS{tabák} \dicExampleIS{neftóbak} \dicExampleCS{šňupací tabák}
\dicEntry[tóbakshorn] \dicTerm{tóbaks··horn} \dicIPA{{t\smash{\textsuperscript{h}}}{ou}{\textlengthmark}{\textsubring{b}}{a}{\r{g}}{s}{h}{\textopeno}{r}{\textsubring{d}}{\textsubring{n}}} \dicPos{n}[2] \dicFlx{(‑s, ‑)}[5] \textbf{1.} \dicDirectTranslationCS{krabička na šňupací tabák}  \textbf{2.} \dicDirectTranslationCS{petúnie, petúnie zahradní} \textit{(l.~{\textLA{Petunia hybrida}})}  \dicsymPhoto\ 
\dicFigure{ds_image_tobakshorn_0_1.jpg}{Tóbakshorn}{Tóbakshorn - Fir0002/Flagstaffotos, CC-BY-NC}
\dicEntry[tóbakspípa] \dicTerm{tóbaks··píp|a} \dicIPA{{t\smash{\textsuperscript{h}}}{ou}{\textlengthmark}{\textsubring{b}}{a}{\r{g}}{s}{p\smash{\textsuperscript{h}}}{i}{\textsubring{b}}{a}} \dicPos{f}[1] \dicFlx{(‑u, ‑ur)}[19] \dicSynonym*{reykjarpípa} \dicDirectTranslationCS{(tabáková) dýmka}
\dicEntry[tófa] \dicTerm{tóf|a} \dicIPA{{t\smash{\textsuperscript{h}}}{ou}{\textlengthmark}{a}} \dicPos{f}[1] \dicFlx{(‑u, ‑ur)}[19] \dicFieldCat{zool.} \dicDirectTranslationCS{liška polární} \textit{(l.~{\textLA{Alopex lagopus}})}  \dicsymPhoto\ 
\dicFigure{20787.jpg}{Tófa}{Tófa - Morehouse, Keith, Biolib, PD}
\dicEntry[tófugras] \dicTerm{tófu··|gras} \dicIPA{{t\smash{\textsuperscript{h}}}{ou}{\textlengthmark}{\textscy}{\r{g}}{r}{a}{s}} \dicPos{n}[2] \dicFlx{(‑grass, ‑grös)}[8] \dicFieldCat{bot.} \dicDirectTranslationCS{puchýřník křehký} \textit{(l.~{\textLA{Cystopteris fragilis}})}  \dicsymPhoto\ 
\dicFigure{91395.jpg}{Tófugras}{Tófugras - Ševčík Jan, Biolib, Copyright/PD}
\dicEntry[Tógó] \dicTerm{Tógó} \dicIPA{{t\smash{\textsuperscript{h}}}{ou}{\textlengthmark}{\textbabygamma}{ou}} \dicPos{n}[4] \dicFlx{indecl}[2] \dicFieldCat{geog.} \dicDirectTranslationCS{Togo}
\dicEntry[Tógói] \dicTerm{Tógó|i} \dicIPA{{t\smash{\textsuperscript{h}}}{ou}{\textlengthmark}{\textbabygamma}{ou}{\textsci}} \dicPos{m}[1] \dicFlx{(‑a, ‑ar)}[1] \dicLink{Tógómaður}
\begin{xtolerant}{}{1pt}
\dicEntry[Tógómaður] \dicTerm{Tógó··|maður}\dicTerm{,\addthinM Tógói}\addthinS\dicIPA{{t\smash{\textsuperscript{h}}}\-{ou}\-{\textlengthmark}\-{\textbabygamma}\-{ou}\-{m}\-{a}\-{ð}\-{\textscy}\-{\textsubring{r}}}\addthinS\dicPos{m}[13]\addthinS\dicFlx{(‑manns, ‑menn)}[2] \dicDirectTranslationCS{Tožan(ka), Togan(ka)}
\end{xtolerant}
\dicEntry[tógóskur] \dicTerm{tógóskur} \dicIPA{{t\smash{\textsuperscript{h}}}{ou}{\textlengthmark}{\textbabygamma}{ou}{s}{\r{g}}{\textscy}{\textsubring{r}}} \dicPos{adj}[1]\dicFlx{}[-6] \dicDirectTranslationCS{tožský, togský}
\dicEntry[tók] \dicTerm{tók} \dicIPA{{t\smash{\textsuperscript{h}}}{ou}{\textlengthmark}{\r{g}}} \dicPos{v} \dicFlx{ind pf sg 1 pers} \dicLink{taka\smash{\textsuperscript{2}}}
\dicEntry[tókum] \dicTerm{tókum} \dicIPA{{t\smash{\textsuperscript{h}}}{ou}{\textlengthmark}{\r{g}}{\textscy}{\textsubring{m}}} \dicPos{v} \dicFlx{ind pf pl 1 pers} \dicLink{taka\smash{\textsuperscript{2}}}
\dicEntry[Tókýó] \dicTerm{Tókýó} \dicIPA{{t\smash{\textsuperscript{h}}}{ou}{\textlengthmark}{\r{\textObardotlessj}}{i}{j}{ou}} \dicPos{subs} \dicFlx{indecl} \dicFieldCat{geog.} \dicDirectTranslationCS{Tokio} \dicIndirectTranslationCS{(hlavní město Japonska)}
\dicEntry[tól] \dicTerm{tól} \dicsymFrequent\  \dicIPA{{t\smash{\textsuperscript{h}}}{ou}{\textlengthmark}{\textsubring{l}}} \dicPos{n}[2] \dicFlx{(‑s, ‑)}[5] \textbf{1.} \dicSynonym{verkfæri} \dicDirectTranslationCS{nástroj, nářadí, náčiní} \dicExampleIS{Ritillinn er auðvelt tól.} \dicExampleCS{Textový editor je jednoduchý nástroj.}  \textbf{2.} \dicLangCat{slang.} \dicSynonym{getnaðarlimur} \dicDirectTranslationCS{žihadlo, klacek} \dicIndirectTranslationCS{(penis)}
\dicEntry[tólf] \dicTerm{tólf} \dicsymFrequent\  \dicIPA{{t\smash{\textsuperscript{h}}}{ou}{l}{f}} \dicPos{num} \dicFlx{indecl} \dicDirectTranslationCS{dvanáct} \dicExampleIS{eftir tólf mánuði} \dicExampleCS{po dvanácti měsících}
\dicEntry[tólffótungur] \dicTerm{tólf··fót·ung|ur} \dicIPA{{t\smash{\textsuperscript{h}}}{ou}{l}{f}{ou}{\textsubring{d}}{u}{\ng}{\r{g}}{\textscy}{\textsubring{r}}} \dicPos{m}[6] \dicFlx{(‑s, ‑ar)}[8] \dicDirectTranslationCS{housenka}
\dicEntry[tólfti] \dicTerm{tólfti} \dicIPA{{t\smash{\textsuperscript{h}}}{ou}{\textsubring{l}}{\textsubring{d}}{\textsci}} \dicPos{num}[5] \dicFlx{ord} \dicDirectTranslationCS{dvanáctý}
\dicEntry[tólg] \dicTerm{tólg} \dicIPA{{t\smash{\textsuperscript{h}}}{ou}{l}{\r{g}}} \dicPos{f}[4] \dicFlx{(‑ar)}[3] \dicDirectTranslationCS{lůj}
\dicEntry[tóm] \dicTerm{tóm} \dicIPA{{t\smash{\textsuperscript{h}}}{ou}{\textlengthmark}{\textsubring{m}}} \dicPos{n}[2] \dicFlx{(‑s)}[2] \textbf{1.} \dicSynonym*{tómt rými} \dicDirectTranslationCS{prázdno, prázdnota, vakuum}  \textbf{2.} \dicSynonym*{næði} \dicDirectTranslationCS{volná chvíle, klid};  \dicPhraseIS{í góðu tómi} \dicFlx{adv} \dicDirectTranslationCS{v~klidu} \dicIndirectTranslationCS{(nerušeně)}
\dicEntry[tómarúm] \dicTerm{tóma··rúm} \dicIPA{{t\smash{\textsuperscript{h}}}{ou}{\textlengthmark}{m}{a}{r}{u}{\textsubring{m}}} \dicPos{n}[2] \dicFlx{(‑s)}[2] \textbf{1.} \dicSynonym{lofttóm} \textls[15]{\dicDirectTranslationCS{vakuum, vzduchoprázdno} \textbf{2.} \dicSynonym{tómleikatilfinning} \dicDirectTranslationCS{prázdno, pusto}}
\dicEntry[tómatmauk] \dicTerm{tómat··mauk} \dicIPA{{t\smash{\textsuperscript{h}}}{ou}{\textlengthmark}{m}{a}{\textsubring{d}}{m}{\oe i}{\r{g}}} \dicPos{n}[2] \dicFlx{(‑s)}[2] \dicDirectTranslationCS{rajčatový\,/\addthin rajský protlak}
\dicEntry[tómatsósa] \dicTerm{tómat··sós|a} \dicIPA{{t\smash{\textsuperscript{h}}}{ou}{\textlengthmark}{m}{a}{\textsubring{d}}{s}{ou}{s}{a}} \dicPos{f}[1] \dicFlx{(‑u, ‑ur)}[7] \dicDirectTranslationCS{kečup}
\dicEntry[tómatsúpa] \dicTerm{tómat··súp|a} \dicIPA{{t\smash{\textsuperscript{h}}}{ou}{\textlengthmark}{m}{a}{\textsubring{d}}{s}{u}{\textsubring{b}}{a}} \dicPos{f}[1] \dicFlx{(‑u, ‑ur)}[19] \dicFieldCat{kulin.} \dicDirectTranslationCS{rajská polévka}
\dicEntry[tómatur] \dicTerm{tómat|ur} \dicIPA{{t\smash{\textsuperscript{h}}}{ou}{\textlengthmark}{m}{a}{\textsubring{d}}{\textscy}{\textsubring{r}}} \dicPos{m}[6] \dicFlx{(‑s, ‑ar)}[20] \dicFieldCat{bot.} \dicDirectTranslationCS{rajské jablko, rajče, rajče jedlé} \textit{(l.~{\textLA{Lycopersicon esculentum}})}
\dicEntry[tómhentur] \dicTerm{tóm··hentur} \dicIPA{{t\smash{\textsuperscript{h}}}{ou}{\textlengthmark}{\textsubring{m}}{h}{\textepsilon}{\textsubring{n}}{\textsubring{d}}{\textscy}{\textsubring{r}}} \dicPos{adj}[1]\dicFlx{}[-10] \dicSynonym{slyppur} \dicDirectTranslationCS{mající prázdné ruce, (jsoucí) s~prázdnýma rukama};  \dicPhraseIS{koma tómhentur} \dicLangCat{přen.} \dicDirectTranslationCS{vrátit se s~prázdnýma rukama}
\dicEntry[tómlátur] \dicTerm{tóm··látur} \dicIPA{{t\smash{\textsuperscript{h}}}{ou}{m}{l}{au}{\textsubring{d}}{\textscy}{\textsubring{r}}} \dicPos{adj}[1]\dicFlx{}[-1] \dicSynonym*{hirðulítill} \dicDirectTranslationCS{lhostejný, apatický, netečný} \dicExampleIS{tómlátur um klæðaburð sinn} \dicExampleCS{lhostejný ke svému oblečení}
\dicEntry[tómlegur] \dicTerm{tóm··legur} \dicIPA{{t\smash{\textsuperscript{h}}}{ou}{m}{l}{\textepsilon}{\textbabygamma}{\textscy}{\textsubring{r}}} \dicPos{adj}[1]\dicFlx{}[-8] \textbf{1.} \dicSynonym{tómur} \dicDirectTranslationCS{prázdný} \dicExampleIS{tómlegt herbergi} \dicExampleCS{prázdný pokoj}  \textbf{2.} \dicSynonym{dapurlegur} \dicDirectTranslationCS{prázdný, smutný}
\dicEntry[tómleikatilfinning] \dicTerm{tóm·leika··til·finn·ing} \dicIPA{{t\smash{\textsuperscript{h}}}\-{ou}\-{m}\-{l}\-{ei}\-{\r{g}}\-{a}\-{t\smash{\textsuperscript{h}}}\-{\textsci}\-{l}\-{f}\-{\textsci}\-{n}\-{i}\-{\ng}\-{\r{g}}\-} \dicPos{f}[4] \dicFlx{(‑ar, ‑ar)}[5] \dicDirectTranslationCS{pocit prázdnoty}
\dicEntry[tómleiki] \dicTerm{tóm··leik|i} \dicIPA{{t\smash{\textsuperscript{h}}}{ou}{m}{l}{ei}{\r{\textObardotlessj}}{\textsci}} \dicPos{m}[1] \dicFlx{(‑a)}[3] \dicDirectTranslationCS{prázdnota, nicota} \dicExampleIS{finna fyrir tómleika} \dicExampleCS{cítit prázdnotu}
\dicEntry[tómlæti] \dicTerm{tóm··læti} \dicIPA{{t\smash{\textsuperscript{h}}}{ou}{m}{l}{a}{i}{\textsubring{d}}{\textsci}} \dicPos{n}[2] \dicFlx{(‑s)}[20] \dicSynonym{deyfð} \dicDirectTranslationCS{lhostejnost, apatie, netečnost} \dicExampleIS{sýna honum tómlæti} \dicExampleCS{dávat mu najevo lhostejnost}
\dicEntry[tómstund] \textls[15]{\dicTerm{tóm··stund} \dicIPA{{t\smash{\textsuperscript{h}}}{ou}{m}{s}{\textsubring{d}}{\textscy}{n}{\textsubring{d}}} \dicPos{f}[7] \dicFlx{(‑ar, ‑ir)}[1] \dicPhraseIS{tómstundir}} \dicFlx{pl} \dicDirectTranslationCS{volný čas, volno} \dicIndirectTranslationCS{(čas věnovaný zálibám)}
\dicEntry[tómstundagaman] \dicTerm{tóm·stunda··gaman} \dicIPA{{t\smash{\textsuperscript{h}}}{ou}{m}{s}{\textsubring{d}}{\textscy}{n}{\textsubring{d}}{a}{\r{g}}{a}{m}{a}{\textsubring{n}}} \dicPos{n}[2] \dicFlx{(‑s)}[2] \dicDirectTranslationCS{koníček, záliba, hobby}

\dicEntry[tómstundaiðja] \dicTerm{tóm·stunda··iðj|a} \dicIPA{{t\smash{\textsuperscript{h}}}{ou}{m}{s}{\textsubring{d}}{\textscy}{n}{\textsubring{d}}{a}{\textsci}{ð}{j}{a}} \dicPos{f}[1] \dicFlx{(‑u)}[5] \dicDirectTranslationCS{oddecho\-vá\,/\addthin zájmová činnost, záliba, hobby}

\dicEntry[tómstundanámskeið] \dicTerm{tóm·stunda··nám·skeið} \dicIPA{{t\smash{\textsuperscript{h}}}\-{ou}\-{m}\-{s}\-{\textsubring{d}}\-{\textscy}\-{n}\-{\textsubring{d}}\-{a}\-{n}\-{au}\-{m}\-{s}\-{\r{\textObardotlessj}}\-{ei}\-{\texttheta}\-} \dicPos{n}[2] \dicFlx{(‑s, ‑)}[5] \dicFieldCat{škol.} \dicDirectTranslationCS{zájmový kurz}
\dicEntry[tómstundaskóli] \dicTerm{tóm·stunda··skól|i} \dicIPA{{t\smash{\textsuperscript{h}}}{ou}{m}{s}{\textsubring{d}}{\textscy}{n}{\textsubring{d}}{a}{s}{\r{g}}{ou}{l}{\textsci}} \dicPos{m}[1] \dicFlx{(‑a, ‑ar)}[1] \dicFieldCat{škol.} \dicDirectTranslationCS{škola nabízející zájmové kurzy}
\dicEntry[tómstundastarf] \dicTerm{tóm·stunda··|starf} \dicIPA{{t\smash{\textsuperscript{h}}}{ou}{m}{s}{\textsubring{d}}{\textscy}{n}{\textsubring{d}}{a}{s}{\textsubring{d}}{a}{r}{f}} \dicPos{n}[2] \dicFlx{(‑starfs, ‑störf)}[8] \dicDirectTranslationCS{volnočasová aktivita}
\dicEntry[tómur] \dicTerm{tómur} \dicsymFrequent\  \dicIPA{{t\smash{\textsuperscript{h}}}{ou}{\textlengthmark}{m}{\textscy}{\textsubring{r}}} \dicPos{adj}[1]\dicFlx{}[-1] \textbf{1.} \dicSynonym*{innihaldslaus} \dicDirectTranslationCS{prázdný, nevyplněný, pustý} \dicExampleIS{tómt ílát} \dicExampleCS{prázdná nádoba}  \textbf{2.} \dicSynonym{merkingarlaus} \dicDirectTranslationCS{prázdný, jalový, plytký} \dicExampleIS{tóm orð} \dicExampleCS{prázdná slova};  \dicIdiom{tómur}{ \dicPhraseIS{koma ekki að tómum kofanum hjá e‑m}} \dicLangCat{přen.} \dicDirectTranslationCS{neodejít s~prázdnou od (koho)}
\dicEntry[tóna] \dicTerm{tón|a} \dicIPA{{t\smash{\textsuperscript{h}}}{ou}{\textlengthmark}{n}{a}} \dicPos{v}[1] \dicFlx{(‑aði)}[1] \dicFlx{acc} \dicDirectTranslationCS{přednést (melodicky)} \dicIndirectTranslationCS{(o~knězi při mši)}
\dicEntry[tónband] \dicTerm{tón··|band} \dicIPA{{t\smash{\textsuperscript{h}}}{ou}{n}{\textsubring{b}}{a}{n}{\textsubring{d}}} \dicPos{n}[2] \dicFlx{(‑bands, ‑bönd)}[8] \dicDirectTranslationCS{kazeta} \dicIndirectTranslationCS{(zvukový nosič)}
\dicEntry[tónbil] \dicTerm{tón··bil} \dicIPA{{t\smash{\textsuperscript{h}}}{ou}{n}{\textsubring{b}}{\textsci}{\textsubring{l}}} \dicPos{n}[2] \dicFlx{(‑s, ‑)}[5] \dicFieldCat{hud.} \dicDirectTranslationCS{interval}
\dicEntry[tónblær] \dicTerm{tón··blæ|r} \dicIPA{{t\smash{\textsuperscript{h}}}{ou}{n}{\textsubring{b}}{l}{a}{i}{\textsubring{r}}} \dicPos{m}[9] \dicFlx{(‑s\,/\addthin ‑jar)}[30] \dicFieldCat{hud.} \dicDirectTranslationCS{témbr, zabarvení hlasu\,/\addthin tónu}
\dicEntry[tónfall] \dicTerm{tón··fall} \dicIPA{{t\smash{\textsuperscript{h}}}{ou}{n}{f}{a}{\textsubring{d}}{\textsubring{l}}} \dicPos{n}[2] \dicFlx{(‑s)}[2] \dicFieldCat{jaz.} \dicDirectTranslationCS{intonace}
\dicEntry[tónfræði] \dicTerm{tón··fræð|i} \dicIPA{{t\smash{\textsuperscript{h}}}{ou}{n}{f}{r}{a}{i}{ð}{\textsci}} \dicPos{f}[3] \dicFlx{(‑i)}[3] \dicFieldCat{hud.} \dicDirectTranslationCS{hudební věda, muzikologie}
\dicEntry[tónheyrn] \dicTerm{tón··heyrn} \dicIPA{{t\smash{\textsuperscript{h}}}{ou}{\textlengthmark}{\textsubring{n}}{h}{ei}{r}{\textsubring{d}}{\textsubring{n}}} \dicPos{f}[7] \dicFlx{(‑ar)}[3] \dicDirectTranslationCS{(hudební) sluch}
\dicEntry[tónhæð] \dicTerm{tón··hæð} \dicIPA{{t\smash{\textsuperscript{h}}}{ou}{\textlengthmark}{\textsubring{n}}{h}{a}{i}{\texttheta}} \dicPos{f}[7] \dicFlx{(‑ar)}[3] \dicFieldCat{hud.} \dicDirectTranslationCS{výška tónu}
\dicEntry[tónik] \dicTerm{tónik} \dicIPA{{t\smash{\textsuperscript{h}}}{ou}{\textlengthmark}{n}{\textsci}{\r{g}}} \dicPos{n}[2] \dicFlx{(‑s)}[31] \dicDirectTranslationCS{tonik}
\dicEntry[tónkvísl] \dicTerm{tón··kvísl} \dicIPA{{t\smash{\textsuperscript{h}}}{ou}{n}{k\smash{\textsuperscript{h}}}{v}{i}{s}{\textsubring{d}}{\textsubring{l}}} \dicPos{f}[4] \dicFlx{(‑ar, ‑ar\,/\addthin ‑ir)}[30] \dicDirectTranslationCS{ladička}
\dicEntry[tónleikar] \dicTerm{tón··leikar} \dicIPA{{t\smash{\textsuperscript{h}}}{ou}{n}{l}{ei}{\r{g}}{a}{\textsubring{r}}} \dicPos{m}[6] \dicFlx{pl}[2] \dicDirectTranslationCS{koncert}
\dicEntry[tónleikasalur] \dicTerm{tón·leika··sal|ur} \dicIPA{{t\smash{\textsuperscript{h}}}{ou}{n}{l}{ei}{\r{g}}{a}{s}{a}{l}{\textscy}{\textsubring{r}}} \dicPos{m}[10] \dicFlx{(‑ar, ‑ir)}[14] \dicDirectTranslationCS{koncert\-ní sál}
\dicEntry[tónlist] \dicTerm{tón··list} \dicsymFrequent\  \dicIPA{{t\smash{\textsuperscript{h}}}{ou}{n}{l}{\textsci}{s}{\textsubring{d}}} \dicPos{f}[7] \dicFlx{(‑ar)}[3] \dicDirectTranslationCS{hudba, muzika} \dicExampleIS{íslensk tónlist} \dicExampleCS{islandská hudba}
\dicEntry[tónlistargáfa] \dicTerm{tón·listar··gáf|a} \dicIPA{{t\smash{\textsuperscript{h}}}{ou}{n}{l}{\textsci}{s}{\textsubring{d}}{a}{r}{\r{g}}{au}{v}{a}} \dicPos{f}[1] \dicFlx{(‑u)}[5] \dicDirectTranslationCS{hudební nadání}
\dicEntry[tónlistarháskóli] \dicTerm{tón·listar··há·skól|i} \dicIPA{{t\smash{\textsuperscript{h}}}\-{ou}\-{n}\-{l}\-{\textsci}\-{s}\-{\textsubring{d}}\-{a}\-{\textsubring{r}}\-{h}\-{au}\-{s}\-{\r{g}}\-{ou}\-{l}\-{\textsci}\-} \dicPos{m}[1] \dicFlx{(‑a, ‑ar)}[1] \dicFieldCat{škol.} \dicDirectTranslationCS{(hudební) konzervatoř}
\dicEntry[tónlistarhátíð] \dicTerm{tón·listar··há·tíð} \dicIPA{{t\smash{\textsuperscript{h}}}{ou}{n}{l}{\textsci}{s}{\textsubring{d}}{a}{\textsubring{r}}{h}{au}{t\smash{\textsuperscript{h}}}{i}{\texttheta}} \dicPos{f}[7] \dicFlx{(‑ar, ‑ir)}[1] \dicDirectTranslationCS{hudební festival}
\begin{xtolerant}{}{1pt}
\dicEntry[tónlistarmaður] \dicTerm{tón·listar··|maður}\addthinS\dicIPA{{t\smash{\textsuperscript{h}}}{ou}{n}{l}{\textsci}{s}{\textsubring{d}}{a}{r}{m}{a}{ð}{\textscy}{\textsubring{r}}}\addthinS\dicPos{m}[13]\addthinS\dicFlx{(‑manns, ‑menn)}[2] \dicDirectTranslationCS{hudebník, hudebnice, muzikant(ka)}
\end{xtolerant}
\begin{xtolerant}{}{1pt}
\dicEntry[tónlistarmyndband] \dicTerm{tón·listar··mynd·|band}\addthinS\textls[-10]{\dicIPA{{t\smash{\textsuperscript{h}}}\-{ou}\-{n}\-{l}\-{\textsci}\-{s}\-{\textsubring{d}}\-{a}\-{r}\-{m}\-{\textsci}\-{n}\-{\textsubring{d}}\-{\textsubring{b}}\-{a}\-{n}\-{\textsubring{d}}}} \dicPos{n}[2]\addthinS\dicFlx{(‑bands, ‑bönd)}[8] \dicDirectTranslationCS{(hudební) videoklip\,/\addthin klip}
\end{xtolerant}
\dicEntry[tónlistarskóli] \dicTerm{tón·listar··skól|i} \dicIPA{{t\smash{\textsuperscript{h}}}{ou}{n}{l}{\textsci}{s}{\textsubring{d}}{a}{\textsubring{r}}{s}{\r{g}}{ou}{l}{\textsci}} \dicPos{m}[1] \dicFlx{(‑a, ‑ar)}[1] \dicDirectTranslationCS{hudební škola, konzervatoř}
\begin{xtolerant}{}{1pt}
\dicEntry[tónlistarunnandi] \dicTerm{tón·listar··unn·|andi}\addthinS\textls[-10]{\dicIPA{{t\smash{\textsuperscript{h}}}\-{ou}\-{n}\-{l}\-{\textsci}\-{s}\-{\textsubring{d}}\-{a}\-{r}\-{\textscy}\-{n}\-{a}\-{n}\-{\textsubring{d}}\-{\textsci}}} \dicPos{m}[2]\addthinSS\dicFlx{(‑anda,\addthinSS‑endur)}[1] \dicDirectTranslationCS{milovník\,/\addthin milovnice hudby}
\end{xtolerant}
\dicEntry[tónmennt] \dicTerm{tón··mennt} \dicIPA{{t\smash{\textsuperscript{h}}}{ou}{n}{m}{\textepsilon}{\textsubring{n}}{\textsubring{d}}} \dicPos{f}[7] \dicFlx{(‑ar, ‑ir)}[1] \dicFieldCat{škol.} \dicDirectTranslationCS{hudební výchova} \dicIndirectTranslationCS{(předmět ve škole)}
\dicEntry[tónn] \dicTerm{tón|n} \dicsymFrequent\  \dicIPA{{t\smash{\textsuperscript{h}}}{ou}{\textsubring{d}}{\textsubring{n}}} \dicPos{m}[6] \dicFlx{(‑s, ‑ar)}[44] \textbf{1.} \dicFieldCat{hud.} \dicDirectTranslationCS{tón} \dicIndirectTranslationCS{(zvuk určité výšky)} \dicExampleIS{falskur tónn} \dicExampleCS{falešný tón}  \textbf{2.} \dicSynonym*{raddblær} \dicDirectTranslationCS{tón} \dicIndirectTranslationCS{(zabarvení řeči)};  \dicIdiom{tónn}{ \dicPhraseIS{senda e‑m tóninn}} \dicDirectTranslationCS{vynadat (komu), dát (komu) najevo svou nelibost, slovně zaútočit na (koho)};  \dicPhraseIS{svara í sama tón} \dicLangCat{přen.} \dicDirectTranslationCS{odpovědět stejným tónem}
\dicEntry[tónset] \dicTerm{tón··set} \dicIPA{{t\smash{\textsuperscript{h}}}{ou}{n}{s}{\textepsilon}{\textsubring{d}}} \dicPos{v} \dicFlx{ind praes sg 1 pers} \dicLink{tónsetja}
\dicEntry[tónsetja] \dicTerm{tón··|setja} \dicIPA{{t\smash{\textsuperscript{h}}}{ou}{n}{s}{\textepsilon}{\textsubring{d}}{j}{a}} \dicPos{v}[4] \dicFlx{(‑set, ‑setti, ‑settum, ‑setti, ‑sett)}[26] \dicFlx{acc} \dicDirectTranslationCS{zhudebnit} \dicExampleIS{tónsetja ljóðið} \dicExampleCS{zhudebnit báseň}
\dicEntry[tónsetning] \dicTerm{tón··set·ning} \dicIPA{{t\smash{\textsuperscript{h}}}{ou}{n}{s}{\textepsilon}{\textsubring{d}}{n}{i}{\ng}{\r{g}}} \dicPos{f}[4] \dicFlx{(‑ar, ‑ar)}[5] \dicDirectTranslationCS{zhudebnění}
\dicEntry[tónsett] \dicTerm{tón··sett} \dicIPA{{t\smash{\textsuperscript{h}}}{ou}{n}{s}{\textepsilon}{h}{\textsubring{d}}} \dicPos{v} \dicFlx{supin} \dicLink{tónsetja}
\dicEntry[tónsetti] \dicTerm{tón··setti} \dicIPA{{t\smash{\textsuperscript{h}}}{ou}{n}{s}{\textepsilon}{h}{\textsubring{d}}{\textsci}} \dicPos{v} \dicFlx{ind\,/\addthin con pf sg 1 pers} \dicLink{tónsetja}
\dicEntry[tónsettum] \dicTerm{tón··settum} \dicIPA{{t\smash{\textsuperscript{h}}}{ou}{n}{s}{\textepsilon}{h}{\textsubring{d}}{\textscy}{\textsubring{m}}} \dicPos{v} \dicFlx{ind pf pl 1 pers} \dicLink{tónsetja}
\dicEntry[tónskáld] \dicTerm{tón··skáld} \dicsymFrequent\  \dicIPA{{t\smash{\textsuperscript{h}}}{ou}{n}{s}{\r{g}}{au}{l}{\textsubring{d}}} \dicPos{n}[2] \dicFlx{(‑s, ‑)}[5] \dicDirectTranslationCS{skladatel(ka), komponista, komponistka} \dicExampleIS{Tónskáldið stjórnaði hljómsveitinni sjálft.} \dicExampleCS{Skladatel osobně dirigoval orchestr.}
\dicEntry[tónsmíð] \dicTerm{tón··smíð}\dicTerm{, tónsmíði} \dicIPA{{t\smash{\textsuperscript{h}}}\-{ou}\-{n}\-{s}\-{m}\-{i}\-{\texttheta}\-} \dicPos{f}[4] \dicFlx{(‑ar, ‑ar)}[1] \textbf{1.} \dicSynonym{tónverk} \dicDirectTranslationCS{kompozice, (hudební) skladba}  \textbf{2.} \dicPhraseIS{tónsmíðar} \dicFlx{pl} \dicSynonym*{samning tónverks} \dicDirectTranslationCS{komponování, skládání (hudby)}
\dicEntry[tónsmíði] \dicTerm{tón··smíð|i} \dicIPA{{t\smash{\textsuperscript{h}}}{ou}{n}{s}{m}{i}{ð}{\textsci}} \dicPos{f}[2] \dicFlx{(‑i, ‑ar)}[1] \dicLink{tónsmíð}
\dicEntry[tónstigi] \dicTerm{tón··stig|i} \dicIPA{{t\smash{\textsuperscript{h}}}{ou}{n}{s}{\textsubring{d}}{i}{j}{\textsci}} \dicPos{m}[1] \dicFlx{(‑a, ‑ar)}[1] \dicFieldCat{hud.} \dicDirectTranslationCS{(tónová) stupnice}
\dicEntry[tóntegund] \dicTerm{tón··tegund} \dicIPA{{t\smash{\textsuperscript{h}}}{ou}{n}{t\smash{\textsuperscript{h}}}{\textepsilon}{\textbabygamma}{\textscy}{n}{\textsubring{d}}} \dicPos{f}[7] \dicFlx{(‑ar, ‑ir)}[1] \dicFieldCat{hud.} \dicDirectTranslationCS{tónina}
\dicEntry[tónverk] \dicTerm{tón··verk} \dicIPA{{t\smash{\textsuperscript{h}}}{ou}{n}{v}{\textepsilon}{\textsubring{r}}{\r{g}}} \dicPos{n}[2] \dicFlx{(‑s, ‑)}[5] \dicDirectTranslationCS{(hudební) skladba, kompozice}
\dicEntry[tóra] \dicTerm{tór|a} \dicIPA{{t\smash{\textsuperscript{h}}}{ou}{\textlengthmark}{r}{a}} \dicPos{v}[2] \dicFlx{(‑ði, ‑að\,/\addthin ‑t)}[185] \textbf{1.} \dicDirectTranslationCS{přežít, přežívat}  \textbf{2.} \dicDirectTranslationCS{plápolat, slabě svítit}
\dicEntry[tótem] \dicTerm{tótem} \dicIPA{{t\smash{\textsuperscript{h}}}{ou}{\textlengthmark}{\textsubring{d}}{\textepsilon}{\textsubring{m}}} \dicPos{n}[2] \dicFlx{(‑s, ‑)}[5] \dicDirectTranslationCS{totem}
\dicEntry[traðka] \dicTerm{traðk|a} \dicIPA{{t\smash{\textsuperscript{h}}}{r}{a}{\texttheta}{\r{g}}{a}} \dicPos{v}[1] \dicFlx{(‑aði)}[13] \textbf{1.} \dicSynonym{troða} \dicDirectTranslationCS{pošlapat, podupat (trávu ap.)} \dicExampleIS{traðka niður blóm} \dicExampleCS{pošlapat květiny}  \textbf{2.} \dicPhraseIS{traðka á e‑m} \dicLangCat{přen.} \dicSynonym{misbjóða} \dicDirectTranslationCS{šlapat po (kom)}
\dicEntry[traf] \dicTerm{traf} \dicIPA{{t\smash{\textsuperscript{h}}}{r}{a}{\textlengthmark}{f}} \dicPos{n}[2] \dicFlx{(‑s, tröf)}[8] \textbf{1.} \dicSynonym{kögur} \dicDirectTranslationCS{střapec, třáseň}  \textbf{2.} \dicSynonym{vaf} \dicDirectTranslationCS{obvaz, obinadlo, gáza}
\dicEntry[trafali] \dicTerm{trafal|i} \dicIPA{{t\smash{\textsuperscript{h}}}{r}{a}{\textlengthmark}{v}{a}{l}{\textsci}} \dicPos{m}[1] \dicFlx{(‑a)}[3] \dicPhraseIS{e‑að er e‑m til trafala} \dicDirectTranslationCS{(co) je (komu) na obtíž}
\dicEntry[traffík] \dicTerm{traffík} \dicIPA{{t\smash{\textsuperscript{h}}}{r}{a}{f}{\textlengthmark}{i}{\r{g}}} \dicPos{f}[10] \dicFlx{(‑ur)}[2] \dicDirectTranslationCS{ruch, shon, provoz}
\dicEntry[tragedía] \dicTerm{tragedí|a} \dicIPA{{t\smash{\textsuperscript{h}}}{r}{a}{\textlengthmark}{\textbabygamma}{\textepsilon}{\textsubring{d}}{i}{j}{a}} \dicPos{f}[1] \dicFlx{(‑u, ‑ur)}[7] \textbf{1.} \dicSynonym{harmleikur} \dicDirectTranslationCS{tragédie} \dicIndirectTranslationCS{(drama s~tragickým závěrem)}  \textbf{2.} \dicSynonym*{sorglegur atburður} \dicDirectTranslationCS{tragická událost, tragédie}
\dicEntry[traktor] \dicTerm{traktor} \dicIPA{{t\smash{\textsuperscript{h}}}{r}{a}{x}{\textsubring{d}}{\textopeno}{\textsubring{r}}} \dicPos{m}[4] \dicFlx{(‑s, ‑ar)}[14] \dicDirectTranslationCS{traktor}
\dicEntry[tralla] \dicTerm{trall|a} \dicIPA{{t\smash{\textsuperscript{h}}}{r}{a}{l}{\textlengthmark}{a}} \dicPos{v}[1] \dicFlx{(‑aði)}[13] \dicDirectTranslationCS{zpívat si tra-la-la-la}
\dicEntry[trampa] \dicTerm{tramp|a} \dicIPA{{t\smash{\textsuperscript{h}}}{r}{a}{\textsubring{m}}{\textsubring{b}}{a}} \dicPos{v}[1] \dicFlx{(‑aði)}[13] \dicDirectTranslationCS{dupnout, dupat} \dicExampleIS{trampa á e‑u} \dicExampleCS{dupnout na (co)}
\dicEntry[trampólín] \dicTerm{trampólín} \dicIPA{{t\smash{\textsuperscript{h}}}{r}{a}{\textsubring{m}}{\textsubring{b}}{ou}{l}{i}{\textsubring{n}}} \dicPos{n}[2] \dicFlx{(‑s, ‑)}[5] \dicDirectTranslationCS{trampolína}
\dicEntry[trana] \dicTerm{trana\smash{\textsuperscript{1}}} \dicIPA{{t\smash{\textsuperscript{h}}}{r}{a}{\textlengthmark}{n}{a}} \dicPos{f}[1] \dicFlx{(trönu, trönur)}[8] \textbf{1.} \dicFieldCat{zool.} \dicDirectTranslationCS{jeřáb} \textit{(l.~{\textLA{Grus}})}  \dicsymPhoto\   \textbf{2.} \dicPhraseIS{trönur} \dicFlx{pl} {\textbf{a.}} \dicDirectTranslationCS{stojan na sušení ryb};  {\textbf{b.}} \dicDirectTranslationCS{malířský stojan}
\dicFigure{20773.jpg}{Trana}{Trana - Hillebrand, Steve, Biolib, PD}
\dicEntry[trana] \dicTerm{tran|a\smash{\textsuperscript{2}}} \dicIPA{{t\smash{\textsuperscript{h}}}{r}{a}{\textlengthmark}{n}{a}} \dicPos{v}[1] \dicFlx{(‑aði)}[13] \dicFlx{dat} \dicPhraseIS{trana sér fram} \dicDirectTranslationCS{cpát se dopředu, vnucovat se}
\dicEntry[trans] \dicTerm{trans} \dicIPA{{t\smash{\textsuperscript{h}}}{r}{a}{n}{s}} \dicPos{m}[4] \dicFlx{(‑, ‑ar)}[36] \dicDirectTranslationCS{trans, vytržení} \dicExampleIS{vera í transi} \dicExampleCS{být v~transu}
\dicEntry[transistor] \dicTerm{transistor} \dicIPA{{t\smash{\textsuperscript{h}}}{r}{a}{n}{s}{\textsci}{s}{\textsubring{d}}{\textopeno}{\textsubring{r}}} \dicPos{m}[4] \dicFlx{(‑s, ‑ar)}[14] \dicFieldCat{elek.} \dicDirectTranslationCS{tranzistor}
\dicEntry[trantur] \dicTerm{trant|ur} \dicIPA{{t\smash{\textsuperscript{h}}}{r}{a}{\textsubring{n}}{\textsubring{d}}{\textscy}{\textsubring{r}}} \dicPos{m}[6] \dicFlx{(‑s, ‑ar)}[5] \dicLangCat{han.} \dicSynonym{kjaftur} \dicDirectTranslationCS{huba, papula} \dicExampleIS{loka á sér trantinum} \dicExampleCS{zavřít papulu}
\dicEntry[trappa] \dicTerm{trappa\smash{\textsuperscript{1}}} \dicsymFrequent\  \dicIPA{{t\smash{\textsuperscript{h}}}{r}{a}{h}{\textsubring{b}}{a}} \dicPos{f}[1] \dicFlx{(tröppu, tröppur)}[8] \textbf{1.} \dicSynonym{þrep} \dicDirectTranslationCS{stupeň, schod} \dicExampleIS{fyrsta trappa} \dicExampleCS{první schod}  \textbf{2.} \dicSynonym*{lítill stigi} \dicDirectTranslationCS{štafle}  \textbf{3.} \dicPhraseIS{tröppur} \dicFlx{pl} \dicSynonym{stigi\smash{\textsuperscript{1}}} \dicDirectTranslationCS{schody} \dicExampleIS{detta í tröppunum} \dicExampleCS{spadnout ze schodů}
\dicEntry[trappa] \dicTerm{trapp|a\smash{\textsuperscript{2}}} \dicIPA{{t\smash{\textsuperscript{h}}}{r}{a}{h}{\textsubring{b}}{a}} \dicPos{v}[1] \dicFlx{(‑aði)}[13] \dicFlx{acc} \dicPhraseIS{trappa niður e‑ð} \dicLangCat{hovor.} \dicDirectTranslationCS{ohraničit (co), omezit (co) (kouření ap.)};  \dicPhraseIS{trappa sig niður} \dicLangCat{hovor.} \dicDirectTranslationCS{ohraničit se, omezit se}
\dicEntry[trassa] \dicTerm{trass|a} \dicIPA{{t\smash{\textsuperscript{h}}}{r}{a}{s}{\textlengthmark}{a}} \dicPos{v}[1] \dicFlx{(‑aði)}[13] \dicFlx{acc} \dicSynonym{vanrækja} \dicDirectTranslationCS{zanedbat, zanedbávat, opomenout, opomíjet} \dicExampleIS{Hann hefur lengi trassað að skila ritgerðum.} \dicExampleCS{Dlouho opomíjel odevzdání eseje.}
\dicEntry[trassafenginn] \dicTerm{trassa··fenginn} \dicIPA{{t\smash{\textsuperscript{h}}}{r}{a}{s}{\textlengthmark}{a}{f}{ei}{\textltailn}{\r{\textObardotlessj}}{\textsci}{\textsubring{n}}} \dicPos{adj}[6]\dicFlx{}[-2] \dicSynonym{hirðulaus} \dicDirectTranslationCS{nedbalý, ledabylý}
\dicEntry[trassi] \dicTerm{trass|i} \dicIPA{{t\smash{\textsuperscript{h}}}{r}{a}{s}{\textlengthmark}{\textsci}} \dicPos{m}[1] \dicFlx{(‑a, ‑ar)}[8] \dicDirectTranslationCS{lajdák, lajdačka} \dicExampleIS{letingjar og trassar} \dicExampleCS{lenoši a~lajdáci}
\dicEntry[trauður] \dicTerm{trauður} \dicIPA{{t\smash{\textsuperscript{h}}}{r}{\oe i}{\textlengthmark}{ð}{\textscy}{\textsubring{r}}} \dicPos{adj}[2]\dicFlx{}[-6] \dicSynonym{ófús} \dicDirectTranslationCS{zdráhavý, neochotný}
\dicEntry[traust] \dicTerm{traust} \dicsymFrequent\  \dicIPA{{t\smash{\textsuperscript{h}}}{r}{\oe i}{s}{\textsubring{d}}} \dicPos{n}[2] \dicFlx{(‑s)}[2] \dicDirectTranslationCS{důvěra, důvěřování};  \dicPhraseIS{bregðast trausti e‑rs} \dicFlx{refl} \dicDirectTranslationCS{zklamat (čí) důvěru};  \dicPhraseIS{njóta trausts e‑rs} \dicDirectTranslationCS{mít (čí) důvěru};  \dicPhraseIS{sýna e‑m traust} \dicDirectTranslationCS{důvěřovat (komu)};  \dicPhraseIS{vera e‑m til halds og trausts} \dicDirectTranslationCS{být (komu) oporou}
\dicEntry[traustlega] \dicTerm{traust··lega} \dicIPA{{t\smash{\textsuperscript{h}}}{r}{\oe i}{s}{\textsubring{d}}{l}{\textepsilon}{\textbabygamma}{a}} \dicPos{adv} \dicDirectTranslationCS{důvěryhodně, spolehlivě, solidně}
\dicEntry[traustlegur] \dicTerm{traust··legur} \dicIPA{{t\smash{\textsuperscript{h}}}{r}{\oe i}{s}{\textsubring{d}}{l}{\textepsilon}{\textbabygamma}{\textscy}{\textsubring{r}}} \dicPos{adj}[1]\dicFlx{}[-8] \dicDirectTranslationCS{důvěryhodný, spolehlivý, solidní} \dicExampleIS{traustleg bygging} \dicExampleCS{solidní budova}
\dicEntry[traustleiki] \dicTerm{traust··leik|i} \dicIPA{{t\smash{\textsuperscript{h}}}{r}{\oe i}{s}{\textsubring{d}}{l}{ei}{\r{\textObardotlessj}}{\textsci}} \dicPos{m}[1] \dicFlx{(‑a)}[3] \dicSynonym{áreiðanleiki} \dicDirectTranslationCS{důvěryhodnost, spolehlivost, solidnost}
\dicEntry[traustsyfirlýsing] \dicTerm{trausts··yfir·lýs·ing} \dicIPA{{t\smash{\textsuperscript{h}}}\-{r}\-{\oe i}\-{s}\-{\textsubring{d}}\-{s}\-{\textsci}\-{v}\-{\textsci}\-{r}\-{l}\-{i}\-{s}\-{i}\-{\ng}\-{\r{g}}\-} \dicPos{f}[4] \dicFlx{(‑ar, ‑ar)}[5] \dicFieldCat{pol.} \dicDirectTranslationCS{vyslovení důvěry}
\dicEntry[traustur] \dicTerm{traustur} \dicsymFrequent\  \dicIPA{{t\smash{\textsuperscript{h}}}{r}{\oe i}{s}{\textsubring{d}}{\textscy}{\textsubring{r}}} \dicPos{adj}[1]\dicFlx{}[-10] \textbf{1.} \dicSynonym{áreiðanlegur} \dicDirectTranslationCS{důvěryhodný, spolehlivý, jistý} \dicExampleIS{traustur vinur} \dicExampleCS{spolehlivý přítel}  \textbf{2.} \dicSynonym{sterkur} \dicDirectTranslationCS{solidní, pevný};  \dicPhraseIS{e‑að stendur á traustum grunni} \dicLangCat{přen.} \dicDirectTranslationCS{(co) stojí na solidních základech};  \dicPhraseIS{e‑að stendur traustum fótum} \dicDirectTranslationCS{(co) je solidní}
\dicEntry[tráss] \dicTerm{tráss} \dicIPA{{t\smash{\textsuperscript{h}}}{r}{au}{s}{\textlengthmark}} \dicPos{n}[2] \dicFlx{(‑)}[24] \dicSynonym*{óvilji} \dicDirectTranslationCS{truc};  \dicPhraseIS{(gera e‑ð) í trássi við e‑n} \dicDirectTranslationCS{(udělat (co)) (komu) na truc}
\dicEntry[treð] \dicTerm{treð} \dicIPA{{t\smash{\textsuperscript{h}}}{r}{\textepsilon}{\textlengthmark}{\texttheta}} \dicPos{v} \dicFlx{ind praes sg 1 pers} \dicLink{troða}
\dicEntry[trefill] \dicTerm{tref|ill} \dicIPA{{t\smash{\textsuperscript{h}}}{r}{\textepsilon}{\textlengthmark}{v}{\textsci}{\textsubring{d}}{\textsubring{l}}} \dicPos{m}[6] \dicFlx{(‑ils, ‑lar)}[35] \dicDirectTranslationCS{šála}
\dicEntry[trefjaefni] \dicTerm{trefja··efni} \dicIPA{{t\smash{\textsuperscript{h}}}{r}{\textepsilon}{v}{j}{a}{\textepsilon}{\textsubring{b}}{n}{\textsci}} \dicPos{n}[2] \dicFlx{(‑s, ‑)}[14] \dicFieldCat{biol.} \dicDirectTranslationCS{fibrin}
\dicEntry[trefjagler] \dicTerm{trefja··gler} \dicIPA{{t\smash{\textsuperscript{h}}}{r}{\textepsilon}{v}{j}{a}{\r{g}}{l}{\textepsilon}{\textsubring{r}}} \dicPos{n}[2] \dicFlx{(‑s, ‑)}[6] \dicFieldCat{techn.} \dicDirectTranslationCS{sklolaminát}
\dicEntry[trefjar] \dicTerm{trefjar} \dicIPA{{t\smash{\textsuperscript{h}}}{r}{\textepsilon}{v}{j}{a}{\textsubring{r}}} \dicPos{f}[12] \dicFlx{pl}[1] \textbf{1.} \dicSynonym*{plöntutrefjar} \dicDirectTranslationCS{vláknina} \dicIndirectTranslationCS{(významná složka potravy)}  \textbf{2.} \dicSynonym*{efnisþræðir} \dicDirectTranslationCS{vláknina} \dicIndirectTranslationCS{(materiál složený z~vláken)}
\dicEntry[trefjóttur] \dicTerm{trefj··óttur} \dicIPA{{t\smash{\textsuperscript{h}}}{r}{\textepsilon}{v}{j}{ou}{h}{\textsubring{d}}{\textscy}{\textsubring{r}}} \dicPos{adj}[1]\dicFlx{}[-10] \textbf{1.} \dicSynonym*{þráðóttur} \dicDirectTranslationCS{vláknitý}  \textbf{2.} \dicSynonym{trosnaður} \dicDirectTranslationCS{roztřepený}
\dicEntry[trega] \dicTerm{treg|a} \dicIPA{{t\smash{\textsuperscript{h}}}{r}{\textepsilon}{\textlengthmark}{\textbabygamma}{a}} \dicPos{v}[1] \dicFlx{(‑aði)}[1] \dicFlx{acc} \dicSynonym{syrgja} \dicDirectTranslationCS{truchlit, rmoutit se} \dicExampleIS{trega dauða hennar sárt} \dicExampleCS{truchlit hluboce nad její smrtí}
\dicEntry[tregða] \dicTerm{tregð|a} \dicIPA{{t\smash{\textsuperscript{h}}}{r}{\textepsilon}{\textbabygamma}{ð}{a}} \dicPos{f}[1] \dicFlx{(‑u, ‑ur)}[7] \textbf{1.} \dicDirectTranslationCS{neochota, liknavost, zdráhavost} \dicExampleIS{sýna tregðu} \dicExampleCS{dávat najevo neochotu}  \textbf{2.} \dicFieldCat{fyz.} \dicDirectTranslationCS{setrvačnost}
\dicEntry[tregðulögmál] \dicTerm{tregðu··lög·mál} \dicIPA{{t\smash{\textsuperscript{h}}}{r}{\textepsilon}{\textbabygamma}{ð}{\textscy}{l}{\oe}{\textbabygamma}{m}{au}{\textsubring{l}}} \dicPos{n}[2] \dicFlx{(‑s)}[2] \dicFieldCat{fyz.} \dicDirectTranslationCS{zákon setrvačnosti}
\dicEntry[treggáfaður] \dicTerm{treg··gáf·|aður} \dicIPA{{t\smash{\textsuperscript{h}}}{r}{\textepsilon}{\textbabygamma}{\r{g}}{au}{v}{a}{ð}{\textscy}{\textsubring{r}}} \dicPos{adj}[3] \dicFlx{(f ‑uð)}[3] \dicSynonym{tornæmur} \dicDirectTranslationCS{nechápavý, pomalý (v~učení se)}
\dicEntry[tregi] \dicTerm{treg|i} \dicIPA{{t\smash{\textsuperscript{h}}}{r}{ei}{\textsci}} \dicPos{m}[1] \dicFlx{(‑a, ‑ar)}[1] \dicSynonym{harmur} \dicDirectTranslationCS{žal, zármutek} \dicExampleIS{Minningin um hana vekur trega.} \dicExampleCS{Vzpomínka na ni budí zármutek.}
\dicEntry[tregur] \dicTerm{tregur} \dicIPA{{t\smash{\textsuperscript{h}}}{r}{\textepsilon}{\textlengthmark}{\textbabygamma}{\textscy}{\textsubring{r}}} \dicPos{adj}[1]\dicFlx{}[-1] \textbf{1.} \dicSynonym{ófús} \dicDirectTranslationCS{neochotný, liknavý, zdráhavý} \dicExampleIS{Hann er tregur til að viðurkenna mistök sín.} \dicExampleCS{Není ochotný přiznat svoje chyby.}  \textbf{2.} \dicSynonym{tornæmur} \dicDirectTranslationCS{nechápavý, tupý, pomalý (v~učení ap.)}  \textbf{3.} \dicSynonym{torveldur} \dicDirectTranslationCS{obtížný, těžký}
\dicEntry[treina] \dicTerm{trein|a} \dicIPA{{t\smash{\textsuperscript{h}}}{r}{ei}{\textlengthmark}{n}{a}} \dicPos{v}[2] \dicFlx{(‑di, ‑t)}[155] \dicFlx{dat + acc} \dicPhraseIS{treina sér e‑ð} \dicDirectTranslationCS{šetřit si (co), šetřit (čím)} \dicExampleIS{treina sér matinn} \dicExampleCS{šetřit si jídlo}
\dicEntry[trekkja] \dicTerm{trekk|ja} \dicIPA{{t\smash{\textsuperscript{h}}}{r}{\textepsilon}{h}{\r{\textObardotlessj}}{a}} \dicPos{v}[2] \dicFlx{(‑ti, ‑t)}[27] \dicFlx{acc} \textbf{1.} \dicSynonym*{draga upp} \dicDirectTranslationCS{natáhnout, natahovat (hodiny ap.)} \dicExampleIS{trekkja úr} \dicExampleCS{natáhnout hodinky};  \dicPhraseIS{trekkja upp e‑ð} \dicDirectTranslationCS{natáhnout (co) (hodiny ap.)}  \textbf{2.} \dicPhraseIS{það trekkir} \dicFlx{impers} \dicDirectTranslationCS{táhne, je průvan} \dicExampleIS{Það trekkir um gættina.} \dicExampleCS{V~průchodu táhne.}  \textbf{3.} \dicPhraseIS{e‑að trekkir að e‑n} \dicLangCat{hovor.} \dicDirectTranslationCS{(co) táhne (koho), (co) láká (koho), (co) přitahuje (čí) pozornost} \dicExampleIS{Hljómsveitin trekkir að ferðamenn.} \dicExampleCS{Kapela přitahuje pozornost turistů.}
\dicEntry[trekktur] \dicTerm{trekktur} \dicIPA{{t\smash{\textsuperscript{h}}}{r}{\textepsilon}{x}{\textsubring{d}}{\textscy}{\textsubring{r}}} \dicPos{adj}[1]\dicFlx{}[-10] \textbf{1.} \dicSynonym*{dreginn upp} \dicDirectTranslationCS{natažený (hodiny ap.)}  \textbf{2.} \dicDirectTranslationCS{napnutý, napjatý, vynervovaný};  \dicPhraseIS{vera trekktur á taugum} \dicDirectTranslationCS{být vynervovaný, mít napnuté nervy}
\dicEntry[trekkur] \dicTerm{trekk|ur} \dicIPA{{t\smash{\textsuperscript{h}}}{r}{\textepsilon}{h}{\r{g}}{\textscy}{\textsubring{r}}} \dicPos{m}[6] \dicFlx{(‑s)}[23] \dicSynonym{súgur} \dicDirectTranslationCS{tah} \dicIndirectTranslationCS{(proudění vzduchu zpravidla do místnosti)}
\dicEntry[trekt] \dicTerm{trekt} \dicIPA{{t\smash{\textsuperscript{h}}}{r}{\textepsilon}{x}{\textsubring{d}}} \dicPos{f}[7] \dicFlx{(‑ar, ‑ir)}[1] \dicDirectTranslationCS{trychtýř, nálevka}
\dicEntry[treyja] \dicTerm{treyj|a} \dicIPA{{t\smash{\textsuperscript{h}}}{r}{ei}{j}{\textlengthmark}{a}} \dicPos{f}[1] \dicFlx{(‑u, ‑ur)}[7] \dicDirectTranslationCS{tričko, triko, mikina}
\dicEntry[treysta] \dicTerm{treyst|a} \dicsymFrequent\  \dicIPA{{t\smash{\textsuperscript{h}}}{r}{ei}{s}{\textsubring{d}}{a}} \dicPos{v}[2] \dicFlx{(‑i, ‑)}[14] \dicFlx{dat\,/\addthin acc} \textbf{1.} \dicFlx{dat} \dicDirectTranslationCS{důvěřovat, spolehnout se, spoléhat se};  \dicPhraseIS{treysta e‑m} \dicDirectTranslationCS{důvěřovat (komu)} \dicExampleIS{Ég treysti henni líka.} \dicExampleCS{Také jí důvěřuji.};  \dicPhraseIS{treysta á e‑ð\,/\addthin e‑n} \dicDirectTranslationCS{spoléhat se na (co\,/\addthin koho)} \dicExampleIS{treysta á sjálfan sig} \dicExampleCS{spoléhat se sám na sebe};  \dicPhraseIS{treysta sér til að (gera e‑ð)} \dicDirectTranslationCS{věřit si na (udělání (čeho))};  \dicPhraseIS{treysta því að (e‑að gerist)} \dicDirectTranslationCS{spolehnout se na to, že (se (co) stane)};  \dicPhraseIS{það er engu að treysta} \dicLangCat{přen.} \dicDirectTranslationCS{jeden nikdy neví, člověk si nikdy nemůže být jistý}  \textbf{2.} \dicFlx{acc} \dicSynonym{styrkja} \dicDirectTranslationCS{posílit, posilovat, zesílit, zesilovat, zpevnit, zpevňovat} \dicExampleIS{treysta vegginn} \dicExampleCS{zpevnit zeď};  \dicIdiom{treystast}{ \dicPhraseIS{treystast til að (gera e‑ð)}} \dicFlx{refl} \dicDirectTranslationCS{věřit si na (udělání (čeho)), odvážit se na (udělání (čeho))} \dicIndirectTranslationCS{(nejčastěji se záporem)}
\dicEntry[tré] \dicTerm{tré} \dicsymFrequent\  \dicIPA{{t\smash{\textsuperscript{h}}}{r}{j}{\textepsilon}{\textlengthmark}} \dicPos{n}[2] \dicFlx{(‑s, ‑)}[36] \textbf{1.} \dicDirectTranslationCS{strom(ek)} \dicExampleIS{milli blómstrandi trjáa} \dicExampleCS{mezi kvetoucími stromy}  \textbf{2.} \dicSynonym{timbur} \dicDirectTranslationCS{dřevo} \dicIndirectTranslationCS{(materiál)} \dicExampleIS{smíða úr tré} \dicExampleCS{vyrábět ze dřeva};  \dicIdiom{tré}{ \dicPhraseIS{hafa\,/\addthin eiga í fullu tré við e‑n}} \dicLangCat{přen.} \dicDirectTranslationCS{vyrovnat se (komu), být s~(kým) na stejné úrovni}
\dicEntry[tréblásturshljóðfæri] \dicTerm{tré·blásturs··hljóð·færi} \dicIPA{{t\smash{\textsuperscript{h}}}\-{r}\-{j}\-{\textepsilon}\-{\textlengthmark}\-{\textsubring{b}}\-{l}\-{au}\-{s}\-{\textsubring{d}}\-{\textscy}\-{\textsubring{r}}\-{s}\-{\textsubring{l}}\-{j}\-{ou}\-{ð}\-{f}\-{a}\-{i}\-{r}\-{\textsci}\-} \dicPos{n}[2] \dicFlx{(‑s, ‑)}[14] \dicFieldCat{hud.} \dicDirectTranslationCS{dřevěný dechový nástroj}
\dicEntry[trébútur] \dicTerm{tré··bút|ur} \dicIPA{{t\smash{\textsuperscript{h}}}{r}{j}{\textepsilon}{\textlengthmark}{\textsubring{b}}{u}{\textsubring{d}}{\textscy}{\textsubring{r}}} \dicPos{m}[6] \dicFlx{(‑s, ‑ar)}[15] \dicSynonym{bútur} \dicDirectTranslationCS{pařez}
\dicEntry[tréklossi] \dicTerm{tré··kloss|i} \dicIPA{{t\smash{\textsuperscript{h}}}{r}{j}{\textepsilon}{\textlengthmark}{k\smash{\textsuperscript{h}}}{l}{\textopeno}{s}{\textsci}} \dicPos{m}[1] \dicFlx{(‑a, ‑ar)}[1] \dicDirectTranslationCS{dřevák}
\dicEntry[trélitur] \dicTerm{tré··lit|ur} \dicIPA{{t\smash{\textsuperscript{h}}}{r}{j}{\textepsilon}{\textlengthmark}{l}{\textsci}{\textsubring{d}}{\textscy}{\textsubring{r}}} \dicPos{m}[10] \dicFlx{(‑ar, ‑ir)}[13] \dicDirectTranslationCS{pastelka}
\dicEntry[tréna] \dicTerm{trén|a} \dicIPA{{t\smash{\textsuperscript{h}}}{r}{j}{\textepsilon}{\textlengthmark}{n}{a}} \dicPos{v}[1] \dicFlx{(‑aði)}[1] \dicDirectTranslationCS{(z)dřevnatět}
\dicEntry[tréni] \dicTerm{tréni} \dicIPA{{t\smash{\textsuperscript{h}}}{r}{j}{\textepsilon}{\textlengthmark}{n}{\textsci}} \dicPos{n}[2] \dicFlx{(‑s)}[20] \dicSynonym{sellulósi} \dicDirectTranslationCS{buničina, celulóza} \dicIndirectTranslationCS{(hmota získaná ze dřeva)} \dicAntonym{beðmi}
\dicEntry[trérista] \dicTerm{tré··rist|a} \dicIPA{{t\smash{\textsuperscript{h}}}{r}{j}{\textepsilon}{\textlengthmark}{r}{\textsci}{s}{\textsubring{d}}{a}} \dicPos{f}[1] \dicFlx{(‑u, ‑ur)}[7] \dicSynonym{tréskurður} \dicDirectTranslationCS{dřevoryt, dřevořezba}
\dicEntry[tréskór] \dicTerm{tré··skó|r} \dicIPA{{t\smash{\textsuperscript{h}}}{r}{j}{\textepsilon}{\textlengthmark}{s}{\r{g}}{ou}{\textsubring{r}}} \dicPos{m}[6] \dicFlx{(‑s, ‑r)}[52] \dicDirectTranslationCS{dřevák} \dicIndirectTranslationCS{(zvláště v~pl)}
\dicEntry[tréskurðarmynd] \dicTerm{tré·skurðar··mynd} \dicIPA{{t\smash{\textsuperscript{h}}}{r}{j}{\textepsilon}{\textlengthmark}{s}{\r{g}}{\textscy}{r}{ð}{a}{r}{m}{\textsci}{n}{\textsubring{d}}} \dicPos{f}[7] \dicFlx{(‑ar, ‑ir)}[1] \dicDirectTranslationCS{dřevoryt, dřevořezba}
\dicEntry[tréskurður] \dicTerm{tré··skurð|ur} \dicIPA{{t\smash{\textsuperscript{h}}}{r}{j}{\textepsilon}{\textlengthmark}{s}{\r{g}}{\textscy}{r}{ð}{\textscy}{\textsubring{r}}} \dicPos{m}[10] \dicFlx{(‑ar)}[7] \textbf{1.} \dicDirectTranslationCS{řezbářství, dřevorytectví}  \textbf{2.} \dicDirectTranslationCS{dřevořezba, dřevoryt}
\dicEntry[trésmiður] \dicTerm{tré··smið|ur} \dicIPA{{t\smash{\textsuperscript{h}}}{r}{j}{\textepsilon}{\textlengthmark}{s}{m}{\textsci}{ð}{\textscy}{\textsubring{r}}} \dicPos{m}[9] \dicFlx{(‑s, ‑ir)}[8] \textbf{1.} \dicDirectTranslationCS{tesař(ka)}  \textbf{2.} \dicDirectTranslationCS{truhlář(ka)}
\dicEntry[trésmíðavinnustofa] \textls[10]{\dicTerm{tré·smíða··vinnu·stof|a} \dicIPA{{t\smash{\textsuperscript{h}}}\-{r}\-{j}\-{\textepsilon}\-{\textlengthmark}\-{s}\-{m}\-{i}\-{ð}\-{a}\-{v}\-{\textsci}\-{n}\-{\textscy}\-{s}\-{\textsubring{d}}\-{\textopeno}\-{v}\-{a}\-} \dicPos{f}[1] \dicFlx{(‑u, ‑ur)}[7]} \dicDirectTranslationCS{truhlárna}
\dicEntry[trésmíði] \dicTerm{tré··smíð|i} \dicIPA{{t\smash{\textsuperscript{h}}}{r}{j}{\textepsilon}{\textlengthmark}{s}{m}{i}{ð}{\textsci}} \dicPos{f}[2] \dicFlx{(‑i, ‑ar)}[1] \textbf{1.} \dicDirectTranslationCS{tesařství}  \textbf{2.} \dicDirectTranslationCS{truhlářství}
\dicEntry[tréspíritus] \dicTerm{tré··spíritus} \dicIPA{{t\smash{\textsuperscript{h}}}{r}{j}{\textepsilon}{\textlengthmark}{s}{\textsubring{b}}{i}{r}{\textsci}{\textsubring{d}}{\textscy}{s}} \dicPos{m}[4] \dicFlx{(‑s)}[18] \dicFieldCat{chem.} \dicDirectTranslationCS{methanol, dřevný\,/\addthin dřevitý líh}
\dicEntry[tréverk] \dicTerm{tré··verk} \dicIPA{{t\smash{\textsuperscript{h}}}{r}{j}{\textepsilon}{\textlengthmark}{v}{\textepsilon}{\textsubring{r}}{\r{g}}} \dicPos{n}[2] \dicFlx{(‑s, ‑)}[5] \dicDirectTranslationCS{tesařská\,/\addthin truhlářská prá\-ce} \dicExampleIS{Íbúðin er tilbúin undir tréverk.} \dicExampleCS{Byt je připraven na tesařské práce.}
\dicEntry[trikk] \dicTerm{trikk} \dicIPA{{t\smash{\textsuperscript{h}}}{r}{\textsci}{h}{\r{g}}} \dicPos{n}[2] \dicFlx{(‑s, ‑)}[5] \dicLangCat{hovor.} \dicSynonym{bragð} \dicDirectTranslationCS{trik, finta}
\dicEntry[trilla] \dicTerm{trill|a\smash{\textsuperscript{1}}} \dicIPA{{t\smash{\textsuperscript{h}}}{r}{\textsci}{l}{\textlengthmark}{a}} \dicPos{f}[1] \dicFlx{(‑u, ‑ur)}[19] \textbf{1.} \dicSynonym{fiskibátur} \dicDirectTranslationCS{malý motorový člun} \dicExampleIS{trillubátur} \dicExampleCS{motorový člun}  \textbf{2.} \dicSynonym{börur} \dicDirectTranslationCS{rudl, rudlík} \dicIndirectTranslationCS{(speciální vozík)}
\dicEntry[trilla] \dicTerm{trill|a\smash{\textsuperscript{2}}} \dicIPA{{t\smash{\textsuperscript{h}}}{r}{\textsci}{l}{\textlengthmark}{a}} \dicPos{f}[1] \dicFlx{(‑u, ‑ur)}[19] \dicFieldCat{hud.} \dicDirectTranslationCS{trylek}
\dicEntry[trilla] \dicTerm{trill|a\smash{\textsuperscript{3}}} \dicIPA{{t\smash{\textsuperscript{h}}}{r}{\textsci}{l}{\textlengthmark}{a}} \dicPos{v}[1] \dicFlx{(‑aði)}[1] \dicFlx{dat} \dicDirectTranslationCS{vézt, vozit (na vozíku ap.)}
\dicEntry[trilljón] \dicTerm{trilljón} \dicIPA{{t\smash{\textsuperscript{h}}}{r}{\textsci}{l}{j}{ou}{\textsubring{n}}} \dicPos{f}[7] \dicFlx{(‑ar, ‑ir)}[1] \dicDirectTranslationCS{trilion, trilión}
\dicEntry[trimm] \dicTerm{trimm} \dicIPA{{t\smash{\textsuperscript{h}}}{r}{\textsci}{m}{\textlengthmark}} \dicPos{n}[2] \dicFlx{(‑s)}[2] \dicDirectTranslationCS{cvičení, trénink} \dicIndirectTranslationCS{(chůze, běhání, posilování ap.)}
\dicEntry[trimma] \dicTerm{trimm|a} \dicIPA{{t\smash{\textsuperscript{h}}}{r}{\textsci}{m}{\textlengthmark}{a}} \dicPos{v}[1] \dicFlx{(‑aði)}[1] \dicDirectTranslationCS{cvičit, trénovat} \dicIndirectTranslationCS{(běhat pro zdraví, posilovat ap.)}
\dicEntry[trippi] \dicTerm{trippi} \dicIPA{{t\smash{\textsuperscript{h}}}{r}{\textsci}{h}{\textsubring{b}}{\textsci}} \dicPos{n}[2] \dicFlx{(‑s, ‑)}[14] \dicDirectTranslationCS{mladý kůň} \dicIndirectTranslationCS{(od jednoho roku do tří let)}
\dicEntry[trissa] \dicTerm{triss|a} \dicIPA{{t\smash{\textsuperscript{h}}}{r}{\textsci}{s}{\textlengthmark}{a}} \dicPos{f}[1] \dicFlx{(‑u, ‑ur)}[7] \dicDirectTranslationCS{kladka};  \dicPhraseIS{vera út um allar trissur} \dicLangCat{přen.} \dicDirectTranslationCS{být všude kolem}
\dicEntry[Trínidad og Tóbagó] \dicTerm{Trínidad og Tóbagó} \dicIPA{{t\smash{\textsuperscript{h}}}{r}{i}{\textlengthmark}{n}{\textsci}{\textsubring{d}}{a}{\textsubring{d}}{}{\textopeno}{\textlengthmark}{x}{t\smash{\textsuperscript{h}}}{ou}{\textlengthmark}{\textsubring{b}}{a}{\textbabygamma}{ou}} \dicPos{n}[4] \dicFlx{indecl}[2] \dicFieldCat{geog.} \dicDirectTranslationCS{Trinidad a~Tobago}
\dicEntry[Trínidadi] \dicTerm{Trínidad|i} \dicIPA{{t\smash{\textsuperscript{h}}}{r}{i}{\textlengthmark}{n}{\textsci}{\textsubring{d}}{a}{\textsubring{d}}{\textsci}} \dicPos{m}[1] \dicFlx{(‑a, ‑ar)}[8] \dicDirectTranslationCS{Trinidaďan(ka)}
\dicEntry[trínidadískur] \dicTerm{trínidadískur} \dicIPA{{t\smash{\textsuperscript{h}}}{r}{i}{\textlengthmark}{n}{\textsci}{\textsubring{d}}{a}{\textsubring{d}}{i}{s}{\r{g}}{\textscy}{\textsubring{r}}} \dicPos{adj}[1]\dicFlx{}[-6] \dicDirectTranslationCS{trinidadský}
\dicEntry[tríó] \dicTerm{tríó} \dicIPA{{t\smash{\textsuperscript{h}}}{r}{i}{\textlengthmark}{j}{ou}} \dicPos{n}[2] \dicFlx{(‑s, ‑)}[5] \textbf{1.} \dicFieldCat{hud.} \dicDirectTranslationCS{trio} \dicIndirectTranslationCS{(komorní hudební skladba pro tři nástroje nebo hlasy)}  \textbf{2.} \dicFieldCat{hud.} \dicDirectTranslationCS{trio} \dicIndirectTranslationCS{(tři umělci tuto skladbu provádějící)}
\dicEntry[trítla] \dicTerm{trítl|a} \dicIPA{{t\smash{\textsuperscript{h}}}{r}{i}{h}{\textsubring{d}}{l}{a}} \dicPos{v}[1] \dicFlx{(‑aði)}[44] \dicDirectTranslationCS{cupat, cupitat}
\dicEntry[trjábolur] \dicTerm{trjá··bol|ur} \dicIPA{{t\smash{\textsuperscript{h}}}{r}{j}{au}{\textlengthmark}{\textsubring{b}}{\textopeno}{l}{\textscy}{\textsubring{r}}} \dicPos{m}[9] \dicFlx{(‑s, ‑ir)}[8] \dicDirectTranslationCS{kmen stromu}
\dicEntry[trjábútur] \dicTerm{trjá··bút|ur} \dicIPA{{t\smash{\textsuperscript{h}}}{r}{j}{au}{\textlengthmark}{\textsubring{b}}{u}{\textsubring{d}}{\textscy}{\textsubring{r}}} \dicPos{m}[6] \dicFlx{(‑s, ‑ar)}[15] \dicSynonym{trébútur} \dicDirectTranslationCS{pařez}
\dicEntry[trjágarður] \dicTerm{trjá··garð|ur} \dicIPA{{t\smash{\textsuperscript{h}}}{r}{j}{au}{\textlengthmark}{\r{g}}{a}{r}{ð}{\textscy}{\textsubring{r}}} \dicPos{m}[6] \dicFlx{(‑s, ‑ar)}[5] \dicDirectTranslationCS{park, arboretum}
\dicEntry[trjágrein] \dicTerm{trjá··grein} \dicIPA{{t\smash{\textsuperscript{h}}}{r}{j}{au}{\textlengthmark}{\r{g}}{r}{ei}{\textsubring{n}}} \dicPos{f}[4] \dicFlx{(‑ar, ‑ar)}[1] \dicDirectTranslationCS{větev (stromu)}
\dicEntry[trjágöng] \dicTerm{trjá··göng} \dicIPA{{t\smash{\textsuperscript{h}}}{r}{j}{au}{\textlengthmark}{\r{g}}{\oe i}{\ng}{\r{g}}} \dicPos{n}[2] \dicFlx{pl}[9] \dicDirectTranslationCS{alej, stromořadí}
\dicEntry[trjákróna] \dicTerm{trjá··krón|a} \dicIPA{{t\smash{\textsuperscript{h}}}{r}{j}{au}{\textlengthmark}{k\smash{\textsuperscript{h}}}{r}{ou}{n}{a}} \dicPos{f}[1] \dicFlx{(‑u, ‑ur)}[7] \dicDirectTranslationCS{koruna stromu}
\dicEntry[trjákvoða] \dicTerm{trjá··kvoð|a} \dicIPA{{t\smash{\textsuperscript{h}}}{r}{j}{au}{\textlengthmark}{k\smash{\textsuperscript{h}}}{v}{\textopeno}{ð}{a}} \dicPos{f}[1] \dicFlx{(‑u, ‑ur)}[19] \dicDirectTranslationCS{smůla, pryskyřice}
\dicEntry[trjálundur] \dicTerm{trjá··lund|ur} \dicIPA{{t\smash{\textsuperscript{h}}}{r}{j}{au}{\textlengthmark}{l}{\textscy}{n}{\textsubring{d}}{\textscy}{\textsubring{r}}} \dicPos{m}[10] \dicFlx{(‑ar, ‑ir)}[4] \dicDirectTranslationCS{háj}
\dicEntry[trjárækt] \dicTerm{trjá··rækt} \dicIPA{{t\smash{\textsuperscript{h}}}{r}{j}{au}{\textlengthmark}{r}{a}{i}{x}{\textsubring{d}}} \dicPos{f}[4] \dicFlx{(‑ar)}[3] \dicDirectTranslationCS{lesnictví, pěstování stromů}
\dicEntry[trjátegund] \dicTerm{trjá··tegund} \dicIPA{{t\smash{\textsuperscript{h}}}{r}{j}{au}{\textlengthmark}{t\smash{\textsuperscript{h}}}{\textepsilon}{\textbabygamma}{\textscy}{n}{\textsubring{d}}} \dicPos{f}[7] \dicFlx{(‑ar, ‑ir)}[1] \dicDirectTranslationCS{druh stromu}
\dicEntry[trjáviður] \dicTerm{trjá··við|ur} \dicIPA{{t\smash{\textsuperscript{h}}}{r}{j}{au}{\textlengthmark}{v}{\textsci}{ð}{\textscy}{\textsubring{r}}} \dicPos{m}[10] \dicFlx{(‑ar, ‑ir)}[19] \dicLangCat{zast.} \dicSynonym{timbur} \dicDirectTranslationCS{dřevo, dříví}
\dicEntry[trjóna] \dicTerm{trjón|a} \dicIPA{{t\smash{\textsuperscript{h}}}{r}{j}{ou}{\textlengthmark}{n}{a}} \dicPos{f}[1] \dicFlx{(‑u, ‑ur)}[7] \textbf{1.} \dicSynonym{trýni} \dicDirectTranslationCS{čumák, rypák, chobot}  \textbf{2.} \dicFieldCat{nám.} \dicIndirectTranslationCS{řezba na přídi lodi}
\dicEntry[troða] \dicTerm{troða} \dicsymFrequent\  \dicIPA{{t\smash{\textsuperscript{h}}}{r}{\textopeno}{\textlengthmark}{ð}{a}} \dicPos{v}[6] \dicFlx{(treð, tróð, tróðum, træði, troðið)}[122] \dicFlx{acc\,/\addthin dat} \textbf{1.} \dicFlx{acc} \dicSynonym{traðka} \dicDirectTranslationCS{(u)šlapat, pošlapat, udusat} \dicExampleIS{troða snjóinn undir fótum} \dicExampleCS{udusat sníh pod nohama}  \textbf{2.} \dicFlx{dat} \dicSynonym{þjappa\smash{\textsuperscript{2}}} \dicDirectTranslationCS{nacpat, cpát, (na)pěchovat} \dicExampleIS{troða heyinu í poka} \dicExampleCS{napěchovat seno do pytle};  \dicIdiom{troða}[á]{ \dicPhraseIS{troða á e‑m}} \dicLangCat{přen.} \dicDirectTranslationCS{šlapat po (kom), utiskovat (koho)}; { \dicPhraseIS{láta ekki troða á sér}} \dicLangCat{přen.} \dicDirectTranslationCS{nenechat po sobě šlapat};  \dicIdiom{troða}[í]{ \dicPhraseIS{troða e‑u í e‑n}} \dicDirectTranslationCS{nacpat (co) do (koho)};  \dicIdiom{troða}[sér]{ \dicPhraseIS{troða sér}} \dicDirectTranslationCS{nacpat  se, cpát se, (na)tlačit se} \dicExampleIS{troða sér áfram} \dicExampleCS{cpát se dopředu};  \dicIdiom{troða}[um]{ \dicPhraseIS{troða e‑m um tær}} \dicLangCat{přen.} \dicDirectTranslationCS{překážet (komu), šlapat (komu) na paty};  \dicIdiom{troða}[út]{ \dicPhraseIS{troða sig út af e‑u}} \dicDirectTranslationCS{nacpat se (čím)};  \dicPhraseIS{troða e‑ð út} \dicDirectTranslationCS{napěchovat (co) (po okraj)};  \dicIdiom{troðast}{ \dicPhraseIS{troðast}} \dicFlx{refl} \dicSynonym*{ryðjast} \dicDirectTranslationCS{tlačit se, cpát se};  \dicIdiom{troðast}[undir]{ \dicPhraseIS{troðast undir}} \dicFlx{refl} \dicDirectTranslationCS{být ušlapán k~smrti}
\dicEntry[troðfullur] \dicTerm{troð··fullur} \dicIPA{{t\smash{\textsuperscript{h}}}{r}{\textopeno}{ð}{f}{\textscy}{\textsubring{d}}{l}{\textscy}{\textsubring{r}}} \dicPos{adj}[10]\dicFlx{}[-8] \dicDirectTranslationCS{naplněný, zaplněný, napěchovaný}
\dicEntry[troðfylla] \dicTerm{troð··fyll|a} \dicIPA{{t\smash{\textsuperscript{h}}}{r}{\textopeno}{ð}{f}{\textsci}{\textsubring{d}}{l}{a}} \dicPos{v}[2] \dicFlx{(‑ti, ‑t)}[63] \dicFlx{acc} \dicDirectTranslationCS{napěchovat, namačkat, natlačit} \dicExampleIS{troðfylla ferðatöskuna} \dicExampleCS{napěchovat cestovní tašku}
\dicEntry[troðið] \dicTerm{troðið} \dicIPA{{t\smash{\textsuperscript{h}}}{r}{\textopeno}{\textlengthmark}{ð}{\textsci}{\texttheta}} \dicPos{v} \dicFlx{supin} \dicLink{troða}
\begin{xtolerant}{}{1pt}
\dicEntry[troðinn] \dicTerm{troðinn} \dicIPA{{t\smash{\textsuperscript{h}}}{r}{\textopeno}{\textlengthmark}{ð}{\textsci}{\textsubring{n}}} \dicPos{adj}[6]\dicFlx{}[-2] \textbf{1.} \dicDirectTranslationCS{ušlapaný, vyšlapaný, udusa\-ný} \dicExampleIS{troðin slóð} \dicExampleCS{vyšlapaná stezka}  \textbf{2.} \dicDirectTranslationCS{napěchovaný, nacpaný};  \dicIdiom{troðinn}{ \dicPhraseIS{feta ekki troðnar slóðir}} \dicLangCat{přen.} \dicDirectTranslationCS{chodit po nevyšlapaných stezkách}
\end{xtolerant}
\dicEntry[troðningur] \dicTerm{troðn··ing|ur} \dicIPA{{t\smash{\textsuperscript{h}}}{r}{\textopeno}{ð}{n}{i}{\ng}{\r{g}}{\textscy}{\textsubring{r}}} \dicPos{m}[6] \dicFlx{(‑s, ‑ar)}[8] \textbf{1.} \dicSynonym*{það að troðast} \dicDirectTranslationCS{tlačení (se), mačkání (se)}  \textbf{2.} \dicSynonym{þröng} \dicDirectTranslationCS{tlačenice, mačkanice} \dicExampleIS{troðningur við innganginn} \dicExampleCS{tlačenice u~vchodu}  \textbf{3.} \dicSynonym*{troðinn stigur} \dicDirectTranslationCS{vyšlapaná stezka} \dicExampleIS{Við fórum gamla troðninga upp brekkurnar.} \dicExampleCS{Šli jsme starou vyšlapanou stezkou vzhůru do svahu.}  \textbf{4.} \dicSynonym{marvaði} \dicDirectTranslationCS{šlapání vody}
\dicEntry[trog] \dicTerm{trog} \dicIPA{{t\smash{\textsuperscript{h}}}{r}{\textopeno}{\textlengthmark}{x}} \dicPos{n}[2] \dicFlx{(‑s, ‑)}[5] \dicDirectTranslationCS{díž(e), káď};  \dicPhraseIS{skera e‑ð niður við trog} \dicLangCat{přen.} \dicDirectTranslationCS{omezit\,/\addthin snížit (co) na minimum}
\dicEntry[troll] \dicTerm{troll} \dicIPA{{t\smash{\textsuperscript{h}}}{r}{\textopeno}{l}{\textlengthmark}} \dicPos{n}[2] \dicFlx{(‑s, ‑)}[5] \dicFieldCat{nám.} \dicDirectTranslationCS{lovení vlečnou sítí}
\dicEntry[tromma] \dicTerm{tromm|a\smash{\textsuperscript{1}}} \dicIPA{{t\smash{\textsuperscript{h}}}{r}{\textopeno}{m}{\textlengthmark}{a}} \dicPos{f}[1] \dicFlx{(‑u, ‑ur)}[7] \dicFieldCat{hud.} \dicDirectTranslationCS{buben}
\dicEntry[tromma] \dicTerm{tromm|a\smash{\textsuperscript{2}}} \dicIPA{{t\smash{\textsuperscript{h}}}{r}{\textopeno}{m}{\textlengthmark}{a}} \dicPos{v}[1] \dicFlx{(‑aði)}[1] \dicDirectTranslationCS{(za)bubnovat}
\dicEntry[trommari] \dicTerm{tromm··ar|i} \dicIPA{{t\smash{\textsuperscript{h}}}{r}{\textopeno}{m}{\textlengthmark}{a}{r}{\textsci}} \dicPos{m}[1] \dicFlx{(‑a, ‑ar)}[13] \dicLink{trommuleikari}
\dicEntry[trommuleikari] \dicTerm{trommu··leik·ar|i}\dicTerm{, trommari} \dicIPA{{t\smash{\textsuperscript{h}}}\-{r}\-{\textopeno}\-{m}\-{\textlengthmark}\-{\textscy}\-{l}\-{ei}\-{\r{g}}\-{a}\-{r}\-{\textsci}\-} \dicPos{m}[1] \dicFlx{(‑a, ‑ar)}[13] \dicDirectTranslationCS{bubeník, bubenice}
\dicEntry[trommusett] \dicTerm{trommu··sett} \dicIPA{{t\smash{\textsuperscript{h}}}{r}{\textopeno}{m}{\textlengthmark}{\textscy}{s}{\textepsilon}{h}{\textsubring{d}}} \dicPos{n}[2] \dicFlx{(‑s, ‑)}[5] \dicDirectTranslationCS{bicí (souprava)}
\dicEntry[tromp] \dicTerm{tromp} \dicIPA{{t\smash{\textsuperscript{h}}}{r}{\textopeno}{\textsubring{m}}{\textsubring{b}}} \dicPos{n}[2] \dicFlx{(‑s, ‑)}[5] \dicDirectTranslationCS{trumf} \dicIndirectTranslationCS{(v~kartách)};  \dicPhraseIS{hafa tromp á hendinni} \dicLangCat{přen.} \dicDirectTranslationCS{mít trumfy v~rukávě};  \dicPhraseIS{taka e‑n með trompi} \dicLangCat{přen.} \dicDirectTranslationCS{trumfnout (koho)}
\dicEntry[trompa] \dicTerm{tromp|a} \dicIPA{{t\smash{\textsuperscript{h}}}{r}{\textopeno}{\textsubring{m}}{\textsubring{b}}{a}} \dicPos{v}[1] \dicFlx{(‑aði)}[1] \textbf{1.} \dicDirectTranslationCS{trumfovat, vynést trumf} \dicIndirectTranslationCS{(v~kartách)}  \textbf{2.} \dicLangCat{přen.} \dicDirectTranslationCS{trumfnout, trumfovat}
\dicEntry[trompet] \dicTerm{trompet} \dicIPA{{t\smash{\textsuperscript{h}}}{r}{\textopeno}{\textsubring{m}}{\textsubring{b}}{\textepsilon}{\textsubring{d}}} \dicPos{n}[2] \dicFlx{(‑s, ‑)}[5] \dicFieldCat{hud.} \dicDirectTranslationCS{trumpeta}
\dicEntry[trompetleikari] \dicTerm{trompet··leik·ar|i} \dicIPA{{t\smash{\textsuperscript{h}}}{r}{\textopeno}{\textsubring{m}}{\textsubring{b}}{\textepsilon}{\textsubring{d}}{l}{ei}{\r{g}}{a}{r}{\textsci}} \dicPos{m}[1] \dicFlx{(‑a, ‑ar)}[13] \dicDirectTranslationCS{trumpetista, trumpetistka}
\dicEntry[trosna] \dicTerm{trosn|a} \dicIPA{{t\smash{\textsuperscript{h}}}{r}{\textopeno}{s}{\textsubring{d}}{n}{a}} \dicPos{v}[1] \dicFlx{(‑aði)}[1] \dicDirectTranslationCS{(roz)třepit se, rozplést se} \dicExampleIS{Dúkurinn er farinn að trosna á endunum.} \dicExampleCS{Ubrus se začíná třepit na okrajích.}
\dicEntry[trosnaður] \dicTerm{trosn··|aður} \dicIPA{{t\smash{\textsuperscript{h}}}{r}{\textopeno}{s}{\textsubring{d}}{n}{a}{ð}{\textscy}{\textsubring{r}}} \dicPos{adj}[3] \dicFlx{(f ‑uð)}[3] \dicDirectTranslationCS{roztřepený, třepící se}
\dicEntry[tróð] \dicTerm{tróð\smash{\textsuperscript{1}}} \dicIPA{{t\smash{\textsuperscript{h}}}{r}{ou}{\textlengthmark}{\texttheta}} \dicPos{n}[2] \dicFlx{(‑s)}[2] \dicSynonym{stopp\smash{\textsuperscript{1}}} \dicDirectTranslationCS{vycpávka, polstrování}
\dicEntry[tróð] \dicTerm{tróð\smash{\textsuperscript{2}}} \dicIPA{{t\smash{\textsuperscript{h}}}{r}{ou}{\textlengthmark}{\texttheta}} \dicPos{v} \dicFlx{ind pf sg 1 pers} \dicLink{troða}
\dicEntry[tróðum] \dicTerm{tróðum} \dicIPA{{t\smash{\textsuperscript{h}}}{r}{ou}{\textlengthmark}{ð}{\textscy}{\textsubring{m}}} \dicPos{v} \dicFlx{ind pf pl 1 pers} \dicLink{troða}
\dicEntry[trójuhestur] \dicTerm{tróju··hest|ur} \dicIPA{{t\smash{\textsuperscript{h}}}{r}{ou}{\textlengthmark}{j}{\textscy}{h}{\textepsilon}{s}{\textsubring{d}}{\textscy}{\textsubring{r}}} \dicPos{m}[6] \dicFlx{(‑s, ‑ar)}[4] \textbf{1.} \dicDirectTranslationCS{trojský kůň} \dicIndirectTranslationCS{(kůň z~řecké mytologie)}  \textbf{2.} \dicFieldCat{poč.} \dicDirectTranslationCS{trojský kůň, trojan}
\dicEntry[trufla] \dicTerm{trufl|a} \dicsymFrequent\  \dicIPA{{t\smash{\textsuperscript{h}}}{r}{\textscy}{\textsubring{b}}{l}{a}} \dicPos{v}[1] \dicFlx{(‑aði)}[1] \dicFlx{acc} \textbf{1.} \dicSynonym{ónáða} \dicDirectTranslationCS{(vy)rušit, vyrušovat} \dicExampleIS{Má ég trufla þig augnablik?} \dicExampleCS{Můžu tě na chvilku vyrušit?}  \textbf{2.} \dicSynonym{rugla} \dicDirectTranslationCS{(na)rušit, pokazit}  \textbf{3.} \dicDirectTranslationCS{přerušit, způsobit poruchu};  \dicIdiom{truflast}{ \dicPhraseIS{truflast}} \dicFlx{refl} {\textbf{a.}} \dicDirectTranslationCS{narušit se, porušit se, oslabit se};  {\textbf{b.}} \dicDirectTranslationCS{pomást se, zbláznit se}
\dicEntry[truflaður] \dicTerm{trufl··|aður} \dicIPA{{t\smash{\textsuperscript{h}}}{r}{\textscy}{\textsubring{b}}{l}{a}{ð}{\textscy}{\textsubring{r}}} \dicPos{adj}[3] \dicFlx{(f ‑uð)}[3] \dicSynonym{geggjaður} \dicDirectTranslationCS{pomatený, šílený}
\dicEntry[truflandi] \dicTerm{trufl··andi} \dicIPA{{t\smash{\textsuperscript{h}}}{r}{\textscy}{\textsubring{b}}{l}{a}{n}{\textsubring{d}}{\textsci}} \dicPos{adj}[13] \dicFlx{indecl}[1] \dicDirectTranslationCS{rušivý, vyrušující, obtěžující}
\dicEntry[truflun] \dicTerm{trufl|un} \dicIPA{{t\smash{\textsuperscript{h}}}{r}{\textscy}{\textsubring{b}}{l}{\textscy}{\textsubring{n}}} \dicPos{f}[7] \dicFlx{(‑unar, ‑anir)}[8] \textbf{1.} \dicSynonym{ónæði} \dicDirectTranslationCS{(vy)rušení} \dicExampleIS{truflun á starfseminni} \dicExampleCS{vyrušení v~činnosti}  \textbf{2.} \dicSynonym{geggjun} \dicDirectTranslationCS{(duševní) pomatení} \dicExampleIS{Mér hélt við truflun.} \dicExampleCS{Pomátl jsem se.}  \textbf{3.} \dicSynonym*{truflun í útsendingu} \dicDirectTranslationCS{(pře)rušení (signálu ap.)}
\dicEntry[trukkur] \dicTerm{trukk|ur} \dicIPA{{t\smash{\textsuperscript{h}}}{r}{\textscy}{h}{\r{g}}{\textscy}{\textsubring{r}}} \dicPos{m}[6] \dicFlx{(‑s, ‑ar)}[15] \dicSynonym{vörubíll} \dicDirectTranslationCS{kamion, nákladní auto, náklaďák}
\dicEntry[trumba] \dicTerm{trumb|a} \dicIPA{{t\smash{\textsuperscript{h}}}{r}{\textscy}{m}{\textsubring{b}}{a}} \dicPos{f}[1] \dicFlx{(‑u, ‑ur)}[13] \dicFieldCat{hud.} \dicDirectTranslationCS{buben}
\dicEntry[trumbuslagari] \dicTerm{trumbu··slag·ar|i} \dicIPA{{t\smash{\textsuperscript{h}}}{r}{\textscy}{m}{\textsubring{b}}{\textscy}{s}{\textsubring{d}}{l}{a}{\textbabygamma}{a}{r}{\textsci}} \dicPos{m}[1] \dicFlx{(‑a, ‑ar)}[10] \dicDirectTranslationCS{bubeník, bubenice}
\dicEntry[trunta] \dicTerm{trunt|a} \dicIPA{{t\smash{\textsuperscript{h}}}{r}{\textscy}{\textsubring{n}}{\textsubring{d}}{a}} \dicPos{f}[1] \dicFlx{(‑u, ‑ur)}[19] \textbf{1.} \dicSynonym{bikkja} \dicDirectTranslationCS{herka} \dicIndirectTranslationCS{(sešlý kůň)}  \textbf{2.} \dicLangCat{han.} \dicSynonym*{óvingjarnleg kona} \dicDirectTranslationCS{mrcha, lítice}
\dicEntry[trú] \dicTerm{trú} \dicsymFrequent\  \dicIPA{{t\smash{\textsuperscript{h}}}{r}{u}{\textlengthmark}} \dicPos{f}[4] \dicFlx{(‑ar, ‑r)}[26] \textbf{1.} \dicSynonym{traust} \dicDirectTranslationCS{důvěra, víra};  \dicPhraseIS{hafa trú á e‑m} \dicDirectTranslationCS{mít ke (komu) důvěru, důvěřovat (komu)};  \dicPhraseIS{í góðri trú} \dicFlx{adv} \dicDirectTranslationCS{v~dobré víře};  \dicPhraseIS{missa trúna á e‑ð} \dicDirectTranslationCS{ztratit důvěru\,/\addthin víru v~(co) (smysl života ap.)}  \textbf{2.} \dicSynonym{átrúnaður} \dicDirectTranslationCS{víra, vyznání, přesvědčení} \dicExampleIS{trú á guð} \dicExampleCS{víra v~Boha};  \dicPhraseIS{ganga af trúnni} \dicDirectTranslationCS{zřeknout se víry};  \dicPhraseIS{kasta trúnni} \dicDirectTranslationCS{ztratit víru};  \dicPhraseIS{kristin trú} \dicFieldCat{náb.} \dicDirectTranslationCS{křesťanství};  \dicPhraseIS{snúast til (kristinnar) trúar} \dicFlx{refl} \dicDirectTranslationCS{obrátit se na (křesťanskou) víru};  \dicPhraseIS{vera blendinn í trúnni} \dicDirectTranslationCS{pochybovat (ve víře)}  \textbf{3.} \dicSynonym{skoðun} \dicDirectTranslationCS{názor, přesvědčení};  \dicPhraseIS{telja e‑m trú um e‑ð} \dicSynonym{sannfæra} \dicDirectTranslationCS{přesvědčovat (koho) o~(čem)};  \dicPhraseIS{það er trú e‑rs að (e‑að gerist)} \dicDirectTranslationCS{je to (čí) přesvědčení, že (se (co) stane)}
\dicEntry[trúa] \dicTerm{trúa} \dicsymFrequent\  \dicIPA{{t\smash{\textsuperscript{h}}}{r}{u}{\textlengthmark}{a}} \dicPos{v}[4] \dicFlx{(trúi, trúði, trúðum, tryði, trúað)}[44] \dicFlx{dat} \textbf{1.} \dicFieldCat{náb.} \dicDirectTranslationCS{věřit, mít víru};  \dicPhraseIS{trúa á guð} \dicDirectTranslationCS{věřit v~Boha}  \textbf{2.} \dicFlx{dat} \dicSynonym*{taka sem sannleik} \dicDirectTranslationCS{věřit, důvěřovat};  \dicPhraseIS{trúa e‑u\,/\addthin e‑m} \dicDirectTranslationCS{věřit (čemu\,/\addthin komu)};  \dicIdiom{trúa}[fyrir]{ \dicPhraseIS{trúa e‑m fyrir e‑u}} {\textbf{a.}} \dicDirectTranslationCS{svěřit\,/\addthin svěřovat (komu co) (úkol ap.)};  {\textbf{b.}} \dicDirectTranslationCS{svěřit\,/\addthin svěřovat (komu co) (tajemství ap.)};  \dicIdiom{trúa}[til]{ \dicPhraseIS{trúa e‑m til að (gera e‑ð)}} \dicDirectTranslationCS{věřit (komu), že (udělá (co))}
\dicEntry[trúað] \dicTerm{trúað} \dicIPA{{t\smash{\textsuperscript{h}}}{r}{u}{\textlengthmark}{a}{\texttheta}} \dicPos{v} \dicFlx{supin} \dicLink{trúa}
\dicEntry[trúaður] \dicTerm{trú··|aður} \dicIPA{{t\smash{\textsuperscript{h}}}{r}{u}{\textlengthmark}{a}{ð}{\textscy}{\textsubring{r}}} \dicPos{adj}[3] \dicFlx{(f ‑uð)}[3] \dicDirectTranslationCS{věřící, pobožný} \dicExampleIS{Hún er trúuð og fer oft í kirkju.} \dicExampleCS{Je věřící a~chodí často do kostela.};  \dicPhraseIS{vera trúaður á að (e‑að gerist)} \dicDirectTranslationCS{být přesvědčen, že se ((co) stane)}
\dicEntry[trúanlegur] \dicTerm{trúan··legur} \dicIPA{{t\smash{\textsuperscript{h}}}{r}{u}{\textlengthmark}{a}{n}{l}{\textepsilon}{\textbabygamma}{\textscy}{\textsubring{r}}} \dicPos{adj}[1]\dicFlx{}[-8] \dicDirectTranslationCS{věrohodný, důvěryhodný};  \dicPhraseIS{taka e‑n trúanlegan} \dicDirectTranslationCS{(u)věřit (komu)};  \dicPhraseIS{taka e‑ð trúanlegt} \dicDirectTranslationCS{(u)věřit (čemu)}
\dicEntry[trúaratriði] \dicTerm{trúar··at·riði} \dicIPA{{t\smash{\textsuperscript{h}}}{r}{u}{\textlengthmark}{a}{r}{a}{\textsubring{d}}{r}{\textsci}{ð}{\textsci}} \dicPos{n}[2] \dicFlx{(‑s, ‑)}[14] \dicDirectTranslationCS{věc\,/\addthin otázka víry}
\dicEntry[trúarbrögð] \dicTerm{trúar··brögð} \dicIPA{{t\smash{\textsuperscript{h}}}{r}{u}{\textlengthmark}{a}{r}{\textsubring{b}}{r}{\oe}{\textbabygamma}{\texttheta}} \dicPos{n}[2] \dicFlx{pl}[9] \dicDirectTranslationCS{náboženství, víra}
\dicEntry[trúarjátning] \dicTerm{trúar··ját·ning} \dicIPA{{t\smash{\textsuperscript{h}}}{r}{u}{\textlengthmark}{a}{r}{j}{au}{h}{\textsubring{d}}{n}{i}{\ng}{\r{g}}} \dicPos{f}[4] \dicFlx{(‑ar, ‑ar)}[5] \dicDirectTranslationCS{vyznání víry, konfese}
\dicEntry[trúarkenning] \dicTerm{trúar··kenn·ing} \dicIPA{{t\smash{\textsuperscript{h}}}{r}{u}{\textlengthmark}{a}{\textsubring{r}}{c\smash{\textsuperscript{h}}}{\textepsilon}{n}{i}{\ng}{\r{g}}} \dicPos{f}[4] \dicFlx{(‑ar, ‑ar)}[5] \dicDirectTranslationCS{náboženské učení}
\dicEntry[trúarlegur] \dicTerm{trúar··legur} \dicIPA{{t\smash{\textsuperscript{h}}}{r}{u}{\textlengthmark}{a}{r}{l}{\textepsilon}{\textbabygamma}{\textscy}{\textsubring{r}}} \dicPos{adj}[1]\dicFlx{}[-8] \dicDirectTranslationCS{duchovní, náboženský}
\dicEntry[trúarlíf] \dicTerm{trúar··líf} \dicIPA{{t\smash{\textsuperscript{h}}}{r}{u}{\textlengthmark}{a}{r}{l}{i}{f}} \dicPos{n}[2] \dicFlx{(‑s)}[2] \dicDirectTranslationCS{duchovní\,/\addthin náboženský život}
\dicEntry[trúarofstæki] \dicTerm{trúar··of·stæki} \dicIPA{{t\smash{\textsuperscript{h}}}{r}{u}{\textlengthmark}{a}{r}{\textopeno}{f}{s}{\textsubring{d}}{a}{i}{\r{\textObardotlessj}}{\textsci}} \dicPos{n}[2] \dicFlx{(‑s)}[20] \dicDirectTranslationCS{náboženský fanatismus}
\dicEntry[trúarsetning] \dicTerm{trúar··set·ning} \dicIPA{{t\smash{\textsuperscript{h}}}{r}{u}{\textlengthmark}{a}{\textsubring{r}}{s}{\textepsilon}{h}{\textsubring{d}}{n}{i}{\ng}{\r{g}}} \dicPos{f}[4] \dicFlx{(‑ar, ‑ar)}[5] \dicDirectTranslationCS{doktrína, dogma}
\dicEntry[trúbador] \dicTerm{trúbador} \dicIPA{{t\smash{\textsuperscript{h}}}{r}{u}{\textlengthmark}{\textsubring{b}}{a}{\textsubring{d}}{\textopeno}{\textsubring{r}}} \dicPos{m}[4] \dicFlx{(‑s, ‑ar)}[14] \dicLink{trúbadúr}
\dicEntry[trúbadúr] \dicTerm{trúbadúr}\dicTerm{, trúbador} \dicIPA{{t\smash{\textsuperscript{h}}}\-{r}\-{u}\-{\textlengthmark}\-{\textsubring{b}}\-{a}\-{\textsubring{d}}\-{u}\-{\textsubring{r}}\-} \dicPos{m}[4] \dicFlx{(‑s, ‑ar)}[14] \dicFieldCat{hist.} \dicDirectTranslationCS{trubadúr}
\dicEntry[trúboð] \dicTerm{trú··boð} \dicIPA{{t\smash{\textsuperscript{h}}}{r}{u}{\textlengthmark}{\textsubring{b}}{\textopeno}{\texttheta}} \dicPos{n}[2] \dicFlx{(‑s, ‑)}[5] \dicDirectTranslationCS{misie, církevní mise}
\dicEntry[trúboði] \dicTerm{trú··boð|i} \dicIPA{{t\smash{\textsuperscript{h}}}{r}{u}{\textlengthmark}{\textsubring{b}}{\textopeno}{ð}{\textsci}} \dicPos{m}[1] \dicFlx{(‑a, ‑ar)}[1] \dicDirectTranslationCS{misionář(ka)}
\dicEntry[trúði] \dicTerm{trúði} \dicIPA{{t\smash{\textsuperscript{h}}}{r}{u}{\textlengthmark}{ð}{\textsci}} \dicPos{v} \dicFlx{ind pf sg 1 pers} \dicLink{trúa}
\dicEntry[trúðum] \dicTerm{trúðum} \dicIPA{{t\smash{\textsuperscript{h}}}{r}{u}{\textlengthmark}{ð}{\textscy}{\textsubring{m}}} \dicPos{v} \dicFlx{ind pf pl 1 pers} \dicLink{trúa}
\dicEntry[trúður] \dicTerm{trúð|ur} \dicIPA{{t\smash{\textsuperscript{h}}}{r}{u}{\textlengthmark}{ð}{\textscy}{\textsubring{r}}} \dicPos{m}[6] \dicFlx{(‑s, ‑ar)}[24] \textbf{1.} \dicSynonym*{sirkuslistamaður} \dicDirectTranslationCS{šašek, klaun(ka)} \dicIndirectTranslationCS{(cirkusový umělec)}  \textbf{2.} \dicSynonym{bjáni} \dicDirectTranslationCS{šašek, kašpar, tatrman(ka)} \dicIndirectTranslationCS{(kdo ze sebe dělá hlupáka)}
\dicEntry[trúfastur] \dicTerm{trú··|fastur} \dicIPA{{t\smash{\textsuperscript{h}}}{r}{u}{\textlengthmark}{f}{a}{s}{\textsubring{d}}{\textscy}{\textsubring{r}}} \dicPos{adj}[1] \dicFlx{(f ‑föst)}[11] \dicSynonym{tryggur} \dicDirectTranslationCS{věrný, loajální}
\dicEntry[trúfesti] \dicTerm{trú··fest|i} \dicIPA{{t\smash{\textsuperscript{h}}}{r}{u}{\textlengthmark}{f}{\textepsilon}{s}{\textsubring{d}}{\textsci}} \dicPos{f}[5] \dicFlx{(‑ar)}[4] \dicDirectTranslationCS{věrnost, loajalita}
\dicEntry[trúfélag] \dicTerm{trú··fé·|lag} \dicIPA{{t\smash{\textsuperscript{h}}}{r}{u}{\textlengthmark}{f}{j}{\textepsilon}{l}{a}{x}} \dicPos{n}[2] \dicFlx{(‑lags, ‑lög)}[8] \dicDirectTranslationCS{náboženská obec\,/\addthin komunita}
\dicEntry[trúflokkur] \dicTerm{trú··flokk|ur} \dicIPA{{t\smash{\textsuperscript{h}}}{r}{u}{\textlengthmark}{f}{l}{\textopeno}{h}{\r{g}}{\textscy}{\textsubring{r}}} \dicPos{m}[6] \dicFlx{(‑s, ‑ar)}[8] \dicDirectTranslationCS{církev, náboženská skupina}
\dicEntry[trúfrelsi] \dicTerm{trú··frelsi} \dicIPA{{t\smash{\textsuperscript{h}}}{r}{u}{\textlengthmark}{f}{r}{\textepsilon}{l}{s}{\textsci}} \dicPos{n}[2] \dicFlx{(‑s)}[20] \dicDirectTranslationCS{náboženská svoboda, svoboda (náboženského) vyznání}
\dicEntry[trúfræði] \dicTerm{trú··fræð|i} \dicIPA{{t\smash{\textsuperscript{h}}}{r}{u}{\textlengthmark}{f}{r}{a}{i}{ð}{\textsci}} \dicPos{f}[3] \dicFlx{(‑i)}[3] \dicDirectTranslationCS{věrouka, dogmatika}
\dicEntry[trúgirni] \dicTerm{trú··girn|i} \dicIPA{{t\smash{\textsuperscript{h}}}{r}{u}{\textlengthmark}{\r{\textObardotlessj}}{\textsci}{r}{\textsubring{d}}{n}{\textsci}} \dicPos{f}[3] \dicFlx{(‑i)}[3] \dicDirectTranslationCS{důvěřivost, lehkověrnost}
\dicEntry[trúgjarn] \dicTerm{trú··|gjarn} \dicIPA{{t\smash{\textsuperscript{h}}}{r}{u}{\textlengthmark}{\r{\textObardotlessj}}{a}{r}{\textsubring{d}}{\textsubring{n}}} \dicPos{adj}[5] \dicFlx{(f ‑gjörn)}[6] \dicDirectTranslationCS{důvěřivý, lehkověrný}
\dicEntry[trúhneigð] \dicTerm{trú··hneigð} \dicIPA{{t\smash{\textsuperscript{h}}}{r}{u}{\textlengthmark}{\textsubring{n}}{ei}{\textbabygamma}{\texttheta}} \dicPos{f}[7] \dicFlx{(‑ar)}[3] \dicDirectTranslationCS{pobožnost, zbožnost, bohabojnost}
\dicEntry[trúhneigður] \dicTerm{trú··hneigður} \dicIPA{{t\smash{\textsuperscript{h}}}{r}{u}{\textlengthmark}{\textsubring{n}}{ei}{\textbabygamma}{ð}{\textscy}{\textsubring{r}}} \dicPos{adj}[2]\dicFlx{}[-1] \dicDirectTranslationCS{pobožný, zbožný, bohabojný}
\dicEntry[trúi] \dicTerm{trúi} \dicIPA{{t\smash{\textsuperscript{h}}}{r}{u}{\textlengthmark}{\textsci}} \dicPos{v} \dicFlx{ind praes sg 1 pers} \dicLink{trúa}
\dicEntry[trúlaus] \dicTerm{trú··laus} \dicIPA{{t\smash{\textsuperscript{h}}}{r}{u}{\textlengthmark}{l}{\oe i}{s}} \dicPos{adj}[5]\dicFlx{}[-1] \dicDirectTranslationCS{nevěřící, bezbožný, ateistický}
\dicEntry[trúlega] \dicTerm{trú··lega} \dicsymFrequent\  \dicIPA{{t\smash{\textsuperscript{h}}}{r}{u}{\textlengthmark}{l}{\textepsilon}{\textbabygamma}{a}} \dicPos{adv} \dicSynonym{líklega} \dicDirectTranslationCS{patrně, pravděpodobně} \dicExampleIS{skemmtilegur og trúlega einstæður viðburður} \dicExampleCS{zábavná a~pravděpodobně unikátní událost}
\dicEntry[trúlegur] \dicTerm{trú··legur} \dicIPA{{t\smash{\textsuperscript{h}}}{r}{u}{\textlengthmark}{l}{\textepsilon}{\textbabygamma}{\textscy}{\textsubring{r}}} \dicPos{adj}[1]\dicFlx{}[-8] \textbf{1.} \dicSynonym{líklegur} \dicDirectTranslationCS{patrný, pravděpodobný}  \textbf{2.} \dicSynonym{öruggur} \dicDirectTranslationCS{důvěryhodný, věrohodný} \dicExampleIS{trúleg saga} \dicExampleCS{věrohodný příběh}
\dicEntry[trúleysi] \dicTerm{trú··leysi} \dicIPA{{t\smash{\textsuperscript{h}}}{r}{u}{\textlengthmark}{l}{ei}{s}{\textsci}} \dicPos{n}[2] \dicFlx{(‑s)}[20] \dicDirectTranslationCS{ateismus}
\dicEntry[trúleysingi] \dicTerm{trú··leys·ing|i} \dicIPA{{t\smash{\textsuperscript{h}}}{r}{u}{\textlengthmark}{l}{ei}{s}{i}{\textltailn}{\r{\textObardotlessj}}{\textsci}} \dicPos{m}[1] \dicFlx{(‑ja, ‑jar)}[14] \dicDirectTranslationCS{ateista, ateistka, nevěřící}
\dicEntry[trúlofa] \dicTerm{trú··lof|a} \dicIPA{{t\smash{\textsuperscript{h}}}{r}{u}{\textlengthmark}{l}{\textopeno}{v}{a}} \dicPos{v}[1] \dicFlx{(‑aði)}[1] \dicFlx{acc} \dicDirectTranslationCS{zasnoubit se};  \dicPhraseIS{e‑ir trúlofa sig} \dicDirectTranslationCS{(kdo) se zasnoubí};  \dicIdiom{trúlofast}{ \dicPhraseIS{trúlofast}} \dicFlx{refl} \dicDirectTranslationCS{zasnoubit se, být zasnouben}; { \dicPhraseIS{trúlofast e‑m}} \dicFlx{refl} \dicDirectTranslationCS{zasnoubit se s~(kým), být zasnouben s~(kým)}
\dicEntry[trúlofaður] \textls[15]{\dicTerm{trú··lof·|aður} \dicIPA{{t\smash{\textsuperscript{h}}}{r}{u}{\textlengthmark}{l}{\textopeno}{v}{a}{ð}{\textscy}{\textsubring{r}}} \dicPos{adj}[3] \dicFlx{(f ‑uð)}[4] \dicDirectTranslationCS{zasnoubený}} \dicExampleIS{vera trúlofaður henni} \dicExampleCS{být s~ní zasnoubený}
\dicEntry[trúlofun] \dicTerm{trú··lof|un} \dicIPA{{t\smash{\textsuperscript{h}}}{r}{u}{\textlengthmark}{l}{\textopeno}{v}{\textscy}{\textsubring{n}}} \dicPos{f}[7] \dicFlx{(‑unar)}[9] \dicDirectTranslationCS{zasnoubení, zásnuby} \dicExampleIS{Þau hafa opinberað trúlofun sína.} \dicExampleCS{Zveřejnili svoje zásnuby.}
\dicEntry[trúlofunarhringur] \dicTerm{trú·lofunar··hring|ur} \dicIPA{{t\smash{\textsuperscript{h}}}\-{r}\-{u}\-{\textlengthmark}\-{l}\-{\textopeno}\-{v}\-{\textscy}\-{n}\-{a}\-{\textsubring{r}}\-{i}\-{\ng}\-{\r{g}}\-{\textscy}\-{\textsubring{r}}\-} \dicPos{m}[6] \dicFlx{(‑s, ‑ir\,/\addthin ‑ar)}[67] \dicDirectTranslationCS{zásnubní prsten}
\dicEntry[trúmaður] \dicTerm{trú··|maður} \dicIPA{{t\smash{\textsuperscript{h}}}{r}{u}{\textlengthmark}{m}{a}{ð}{\textscy}{\textsubring{r}}} \dicPos{m}[13] \dicFlx{(‑manns, ‑menn)}[2] \dicDirectTranslationCS{věřící} \dicIndirectTranslationCS{(člověk věřící v~Boha)}
\dicEntry[trúmál] \dicTerm{trú··mál} \dicIPA{{t\smash{\textsuperscript{h}}}{r}{u}{\textlengthmark}{m}{au}{\textsubring{l}}} \dicPos{n}[2] \dicFlx{pl}[1] \dicDirectTranslationCS{náboženské otázky\,/\addthin záležitosti}
\dicEntry[trúmennska] \dicTerm{trú··mennsk|a} \dicIPA{{t\smash{\textsuperscript{h}}}{r}{u}{\textlengthmark}{m}{\textepsilon}{n}{s}{\r{g}}{a}} \dicPos{f}[1] \dicFlx{(‑u)}[5] \dicSynonym{dyggð} \dicDirectTranslationCS{věrnost, oddanost, loajalita} \dicExampleIS{vinna af trúmennsku} \dicExampleCS{pracovat s~oddaností}
\dicEntry[trúnaðarbrot] \dicTerm{trúnaðar··brot} \dicIPA{{t\smash{\textsuperscript{h}}}{r}{u}{\textlengthmark}{n}{a}{ð}{a}{r}{\textsubring{b}}{r}{\textopeno}{\textsubring{d}}} \dicPos{n}[2] \dicFlx{(‑s, ‑)}[5] \dicDirectTranslationCS{porušení důvěrnosti}
\dicEntry[trúnaðarkona] \dicTerm{trúnaðar··kon|a} \dicIPA{{t\smash{\textsuperscript{h}}}{r}{u}{\textlengthmark}{n}{a}{ð}{a}{\textsubring{r}}{k\smash{\textsuperscript{h}}}{\textopeno}{n}{a}} \dicPos{f}[1] \dicFlx{(‑u, ‑ur)}[27] \textbf{1.} \dicDirectTranslationCS{důvěrnice}  \textbf{2.} \dicDirectTranslationCS{důvěrná přítelkyně}
\dicEntry[trúnaðarlæknir] \dicTerm{trúnaðar··lækn|ir} \dicIPA{{t\smash{\textsuperscript{h}}}{r}{u}{\textlengthmark}{n}{a}{ð}{a}{r}{l}{a}{i}{h}{\r{g}}{n}{\textsci}{\textsubring{r}}} \dicPos{m}[7] \dicFlx{(‑is, ‑ar)}[1] \dicDirectTranslationCS{firem\-ní lékař(ka)}
\dicEntry[trúnaðarmaður] \dicTerm{trúnaðar··|maður} \dicIPA{{t\smash{\textsuperscript{h}}}{r}{u}{\textlengthmark}{n}{a}{ð}{a}{r}{m}{a}{ð}{\textscy}{\textsubring{r}}} \dicPos{m}[13] \dicFlx{(‑manns, ‑menn)}[2] \textbf{1.} \dicSynonym{fulltrúi} \dicDirectTranslationCS{důvěrník}  \textbf{2.} \dicSynonym*{náinn vinur} \dicDirectTranslationCS{důvěrná osoba, důvěrný přítel}
\dicEntry[trúnaðarmál] \dicTerm{trúnaðar··mál} \dicIPA{{t\smash{\textsuperscript{h}}}{r}{u}{\textlengthmark}{n}{a}{ð}{a}{r}{m}{au}{\textsubring{l}}} \dicPos{n}[2] \dicFlx{(‑s, ‑)}[5] \dicDirectTranslationCS{důvěrnost, důvěrná záležitost} \dicExampleIS{Efni þessarar skýrslu er trúnaðarmál.} \dicExampleCS{Obsah této zprávy je důvěrnou záležitostí.}
\dicEntry[trúnaðartraust] \dicTerm{trúnaðar··traust} \dicIPA{{t\smash{\textsuperscript{h}}}{r}{u}{\textlengthmark}{n}{a}{ð}{a}{\textsubring{r}}{t\smash{\textsuperscript{h}}}{r}{\oe i}{s}{\textsubring{d}}} \dicPos{n}[2] \dicFlx{(‑s)}[2] \dicDirectTranslationCS{důvěra}
\dicEntry[trúnaður] \dicTerm{trú··nað|ur} \dicIPA{{t\smash{\textsuperscript{h}}}{r}{u}{\textlengthmark}{n}{a}{ð}{\textscy}{\textsubring{r}}} \dicPos{m}[10] \dicFlx{(‑ar)}[9] \textbf{1.} \dicSynonym{tryggð} \dicDirectTranslationCS{věrnost, oddanost, loajalita};  \dicPhraseIS{halda trúnað við e‑n} \dicDirectTranslationCS{být věrný\,/\addthin oddaný (komu)}  \textbf{2.} \dicSynonym{traust} \dicDirectTranslationCS{důvěra};  \dicPhraseIS{leggja trúnað á e‑ð} \dicDirectTranslationCS{důvěřovat (čemu)} \dicExampleIS{leggja engan trúnað á orð hans} \dicExampleCS{nedůvěřovat jeho slovům}  \textbf{3.} \dicSynonym{trúnaðarmál} \dicDirectTranslationCS{důvěrnost};  \dicPhraseIS{segja e‑m e‑ð í trúnaði} \dicDirectTranslationCS{říci (komu co) v~důvěrnosti}
\dicEntry[trúr] \dicTerm{trúr} \dicIPA{{t\smash{\textsuperscript{h}}}{r}{u}{\textlengthmark}{\textsubring{r}}} \dicPos{adj}[4]\dicFlx{}[-1] \textbf{1.} \dicSynonym{dyggur} \dicDirectTranslationCS{věrný, oddaný};  \dicPhraseIS{vera trúr e‑m} \dicDirectTranslationCS{být věrný (komu)} \dicExampleIS{vera trúr maka sínum} \dicExampleCS{být věrný svému partnerovi}  \textbf{2.} \dicSynonym{öruggur} \dicDirectTranslationCS{věrný, spolehlivý (vzpomínka ap.)} \dicExampleIS{trútt minni} \dicExampleCS{věrná vzpomínka}
\dicEntry[trúrækinn] \dicTerm{trú··rækinn} \dicIPA{{t\smash{\textsuperscript{h}}}{r}{u}{\textlengthmark}{r}{a}{i}{\r{\textObardotlessj}}{\textsci}{\textsubring{n}}} \dicPos{adj}[6]\dicFlx{}[-2] \dicDirectTranslationCS{zbožný, pobožný}
\dicEntry[trúrækni] \dicTerm{trú··rækn|i} \dicIPA{{t\smash{\textsuperscript{h}}}{r}{u}{\textlengthmark}{r}{a}{i}{h}{\r{g}}{n}{\textsci}} \dicPos{f}[3] \dicFlx{(‑i)}[3] \dicDirectTranslationCS{zbožnost, pobožnost}
\dicEntry[trúskipti] \dicTerm{trú··skipti} \dicIPA{{t\smash{\textsuperscript{h}}}{r}{u}{\textlengthmark}{s}{\r{\textObardotlessj}}{\textsci}{f}{\textsubring{d}}{\textsci}} \dicPos{n}[2] \dicFlx{pl}[19] \dicFieldCat{náb.} \dicDirectTranslationCS{obrácení víry}
\dicEntry[trúskiptingur] \dicTerm{trú··skipt·ing|ur} \dicIPA{{t\smash{\textsuperscript{h}}}{r}{u}{\textlengthmark}{s}{\r{\textObardotlessj}}{\textsci}{f}{\textsubring{d}}{i}{\ng}{\r{g}}{\textscy}{\textsubring{r}}} \dicPos{m}[6] \dicFlx{(‑s, ‑ar)}[8] \dicFieldCat{náb.} \dicDirectTranslationCS{novověrec, novověrkyně, konvertita, konvertitka}
\dicEntry[trúss] \dicTerm{trúss} \dicIPA{{t\smash{\textsuperscript{h}}}{r}{u}{s}{\textlengthmark}} \dicPos{n}[2] \dicFlx{(‑, ‑)}[22] \dicSynonym{farangur} \dicDirectTranslationCS{zavazadlo, bagáž};  \dicPhraseIS{binda trúss sitt við e‑n} \dicDirectTranslationCS{přidružit se ke (komu)}
\dicEntry[trúverðugur] \dicTerm{trú··verð·ugur}\dicTerm{, trúverður} \dicIPA{{t\smash{\textsuperscript{h}}}\-{r}\-{u}\-{\textlengthmark}\-{v}\-{\textepsilon}\-{r}\-{ð}\-{\textscy}\-{\textbabygamma}\-{\textscy}\-{\textsubring{r}}\-} \dicPos{adj}[1]\dicFlx{}[-8] \dicSynonym{traustur} \dicDirectTranslationCS{věrohodný, důvěryhodný}
\dicEntry[trúverður] \dicTerm{trú··verður} \dicIPA{{t\smash{\textsuperscript{h}}}{r}{u}{\textlengthmark}{v}{\textepsilon}{r}{ð}{\textscy}{\textsubring{r}}} \dicPos{adj}[2]\dicFlx{}[-1] \dicLink{trúverðugur}
\dicEntry[trúvilla] \dicTerm{trú··vill|a} \dicIPA{{t\smash{\textsuperscript{h}}}{r}{u}{\textlengthmark}{v}{\textsci}{\textsubring{d}}{l}{a}} \dicPos{f}[1] \dicFlx{(‑u, ‑ur)}[13] \dicFieldCat{náb.} \dicDirectTranslationCS{hereze, kacířství, bludařství}
\dicEntry[trúvillingur] \dicTerm{trú··vill·ing|ur} \dicIPA{{t\smash{\textsuperscript{h}}}{r}{u}{\textlengthmark}{v}{\textsci}{\textsubring{d}}{l}{i}{\ng}{\r{g}}{\textscy}{\textsubring{r}}} \dicPos{m}[6] \dicFlx{(‑s, ‑ar)}[8] \dicFieldCat{náb.} \dicDirectTranslationCS{heretik, heretička, kacíř(ka), bludař(ka)}
\dicEntry[tryði] \dicTerm{tryði} \dicIPA{{t\smash{\textsuperscript{h}}}{r}{\textsci}{\textlengthmark}{ð}{\textsci}} \dicPos{v} \dicFlx{con pf sg 1 pers} \dicLink{trúa}
\dicEntry[tryggð] \dicTerm{tryggð} \dicIPA{{t\smash{\textsuperscript{h}}}{r}{\textsci}{\textbabygamma}{\texttheta}} \dicPos{f}[7] \dicFlx{(‑ar, ‑ir)}[1] \dicDirectTranslationCS{věrnost, oddanost, loajalita};  \dicPhraseIS{halda tryggð við e‑n} \dicDirectTranslationCS{být věrný\,/\addthin oddaný (komu)};  \dicPhraseIS{svíkja e‑n í tryggðum} \dicDirectTranslationCS{zradit (koho)};  \dicPhraseIS{sýna e‑m tryggð} \dicDirectTranslationCS{být věrný\,/\addthin oddaný (komu)}
\dicEntry[tryggður] \dicTerm{tryggður} \dicIPA{{t\smash{\textsuperscript{h}}}{r}{\textsci}{\textbabygamma}{ð}{\textscy}{\textsubring{r}}} \dicPos{adj}[2]\dicFlx{}[-4] \dicDirectTranslationCS{pojištěný};  \dicPhraseIS{e‑að er tryggt gegn e‑u} \dicDirectTranslationCS{(co) je pojištěné proti (čemu)}
\dicEntry[tryggilega] \dicTerm{tryggi··lega} \dicIPA{{t\smash{\textsuperscript{h}}}{r}{\textsci}{\r{\textObardotlessj}}{\textlengthmark}{\textsci}{l}{\textepsilon}{\textbabygamma}{a}} \dicPos{adv} \dicDirectTranslationCS{opatrně}
\dicEntry[tryggilegur] \dicTerm{tryggi··legur} \dicIPA{{t\smash{\textsuperscript{h}}}{r}{\textsci}{\r{\textObardotlessj}}{\textlengthmark}{\textsci}{l}{\textepsilon}{\textbabygamma}{\textscy}{\textsubring{r}}} \dicPos{adj}[1]\dicFlx{}[-8] \dicSynonym{áreiðanlegur} \dicDirectTranslationCS{spolehlivý, solidní} \dicExampleIS{tryggilegt veður} \dicExampleCS{spolehlivé počasí}
\dicEntry[trygging] \dicTerm{trygg··ing} \dicIPA{{t\smash{\textsuperscript{h}}}{r}{\textsci}{\r{\textObardotlessj}}{\textlengthmark}{i}{\ng}{\r{g}}} \dicPos{f}[4] \dicFlx{(‑ar, ‑ar)}[5] \textbf{1.} \dicSynonym{öryggi} \dicDirectTranslationCS{záruka, garance}  \textbf{2.} \dicFieldCat{ekon.} \dicSynonym{vátrygging} \dicDirectTranslationCS{pojištění, pojistka}
\dicEntry[tryggingafélag] \dicTerm{trygginga··fé·|lag}\dicTerm{, tryggingarfélag} \dicIPA{{t\smash{\textsuperscript{h}}}\-{r}\-{\textsci}\-{\r{\textObardotlessj}}\-{\textlengthmark}\-{i}\-{\ng}\-{\r{g}}\-{a}\-{f}\-{j}\-{\textepsilon}\-{l}\-{a}\-{x}\-} \dicPos{n}[2] \dicFlx{(‑lags, ‑lög)}[8] \dicDirectTranslationCS{pojišťovna}
\dicEntry[tryggingafræðingur] \dicTerm{trygginga·fræð··ing|ur} \dicIPA{{t\smash{\textsuperscript{h}}}\-{r}\-{\textsci}\-{\r{\textObardotlessj}}\-{\textlengthmark}\-{i}\-{\ng}\-{\r{g}}\-{a}\-{f}\-{r}\-{a}\-{i}\-{ð}\-{i}\-{\ng}\-{\r{g}}\-{\textscy}\-{\textsubring{r}}\-} \dicPos{m}[6] \dicFlx{(‑s, ‑ar)}[8] \dicDirectTranslationCS{pojistný matematik, pojistná matematička}
\dicEntry[tryggingakerfi] \dicTerm{trygginga··kerfi} \dicIPA{{t\smash{\textsuperscript{h}}}{r}{\textsci}{\r{\textObardotlessj}}{\textlengthmark}{i}{\ng}{\r{g}}{a}{c\smash{\textsuperscript{h}}}{\textepsilon}{r}{v}{\textsci}} \dicPos{n}[2] \dicFlx{(‑s, ‑)}[14] \dicDirectTranslationCS{sociální systém, systém sociálního zabezpečení}
\dicEntry[tryggingalæknir] \dicTerm{trygginga··lækn|ir} \dicIPA{{t\smash{\textsuperscript{h}}}{r}{\textsci}{\r{\textObardotlessj}}{\textlengthmark}{i}{\ng}{\r{g}}{a}{l}{a}{i}{h}{\r{g}}{n}{\textsci}{\textsubring{r}}} \dicPos{m}[7] \dicFlx{(‑is, ‑ar)}[1] \dicIndirectTranslationCS{lékař, který posuzuje nároky na platby sociálního zabezpečení}
\dicEntry[tryggingarfé] \dicTerm{tryggingar··|fé} \dicIPA{{t\smash{\textsuperscript{h}}}{r}{\textsci}{\r{\textObardotlessj}}{\textlengthmark}{i}{\ng}{\r{g}}{a}{\textsubring{r}}{f}{j}{\textepsilon}} \dicPos{n}[3] \dicFlx{(‑fjár)}[5] \textbf{1.} \dicSynonym*{innborgun á samning} \dicDirectTranslationCS{záloha, zajišťovací poplatek}  \textbf{2.} \dicSynonym*{greiðsla tryggingabóta} \dicDirectTranslationCS{pojistné}
\begin{xtolerant}{}{1pt}
\dicEntry[tryggingarfélag] \dicTerm{tryggingar··fé·\hspace{-0.06667em}|lag}\addthinS\dicIPA{{t\smash{\textsuperscript{h}}}{r}{\textsci}{\r{\textObardotlessj}}{\textlengthmark}{i}{\ng}{\r{g}}{a}{\textsubring{r}}{f}{j}{\textepsilon}{l}{a}{x}}\addthinS\dicPos{n}[2]\addthinS\dicFlx{(‑lags,\addthinSS‑lög)}[8]\addthinS\dicLink{tryggingafélag}
\end{xtolerant}
\dicEntry[tryggingarhæfur] \dicTerm{tryggingar··hæfur} \dicIPA{{t\smash{\textsuperscript{h}}}{r}{\textsci}{\r{\textObardotlessj}}{\textlengthmark}{i}{\ng}{\r{g}}{a}{\textsubring{r}}{h}{a}{i}{v}{\textscy}{\textsubring{r}}} \dicPos{adj}[1]\dicFlx{}[-1] \dicDirectTranslationCS{pojistitelný}
\dicEntry[tryggingartaki] \dicTerm{tryggingar··tak|i} \dicIPA{{t\smash{\textsuperscript{h}}}{r}{\textsci}{\r{\textObardotlessj}}{\textlengthmark}{i}{\ng}{\r{g}}{a}{\textsubring{r}}{t\smash{\textsuperscript{h}}}{a}{\r{\textObardotlessj}}{\textsci}} \dicPos{m}[1] \dicFlx{(‑a, ‑ar)}[8] \dicDirectTranslationCS{pojištěnec, pojištěnka}
\dicEntry[tryggingaskírteini] \dicTerm{trygginga··skír·teini} \dicIPA{{t\smash{\textsuperscript{h}}}\-{r}\-{\textsci}\-{\r{\textObardotlessj}}\-{\textlengthmark}\-{i}\-{\ng}\-{\r{g}}\-{a}\-{s}\-{\r{\textObardotlessj}}\-{i}\-{\textsubring{r}}\-{t\smash{\textsuperscript{h}}}\-{ei}\-{n}\-{\textsci}\-} \dicPos{n}[2] \dicFlx{(‑s, ‑)}[14] \dicDirectTranslationCS{průkaz pojištěnce}
\dicEntry[tryggingastofnun] \dicTerm{trygginga··stofn|un} \dicIPA{{t\smash{\textsuperscript{h}}}{r}{\textsci}{\r{\textObardotlessj}}{\textlengthmark}{i}{\ng}{\r{g}}{a}{s}{\textsubring{d}}{\textopeno}{\textsubring{b}}{n}{\textscy}{\textsubring{n}}} \dicPos{f}[7] \dicFlx{(‑unar, ‑anir)}[8] \dicDirectTranslationCS{správa sociálního zabezpečení}
\dicEntry[tryggja] \dicTerm{trygg|ja} \dicsymFrequent\  \dicIPA{{t\smash{\textsuperscript{h}}}{r}{\textsci}{\r{\textObardotlessj}}{\textlengthmark}{a}} \dicPos{v}[2] \dicFlx{(‑ði, ‑t)}[89] \dicFlx{(dat +) acc} \textbf{1.} \dicFlx{acc} \dicSynonym*{setja í tryggingu} \dicDirectTranslationCS{pojistit, pojišťovat} \dicExampleIS{tryggja húsið hjá tryggingafélagi} \dicExampleCS{pojistit dům v~pojišťovně}  \textbf{2.} \dicFlx{(dat +) acc} \dicSynonym*{gera öruggt} \dicDirectTranslationCS{(za)ručit, zajistit, zajišťovat, garantovat} \dicExampleIS{tryggja mannréttindi í landinu} \dicExampleCS{zajistit lidská práva v~zemi};  \dicPhraseIS{tryggja e‑m e‑ð} \dicDirectTranslationCS{zaručit (komu co)};  \dicPhraseIS{tryggja sér e‑ð} \dicDirectTranslationCS{zajistit si (co)} \dicExampleIS{tryggja sér eintak af bókinni} \dicExampleCS{zajistit si jeden exemplář knihy}
\dicEntry[trygglyndi] \dicTerm{trygg··lyndi} \dicIPA{{t\smash{\textsuperscript{h}}}{r}{\textsci}{\r{g}}{l}{\textsci}{n}{\textsubring{d}}{\textsci}} \dicPos{n}[2] \dicFlx{(‑s)}[20] \dicDirectTranslationCS{věrnost, loajalita}
\dicEntry[trygglyndur] \dicTerm{trygg··lyndur} \dicIPA{{t\smash{\textsuperscript{h}}}{r}{\textsci}{\r{g}}{l}{\textsci}{n}{\textsubring{d}}{\textscy}{\textsubring{r}}} \dicPos{adj}[2]\dicFlx{}[-14] \dicSynonym{trúr} \dicDirectTranslationCS{věrný, loajální}
\dicEntry[tryggur] \dicTerm{tryggur} \dicIPA{{t\smash{\textsuperscript{h}}}{r}{\textsci}{\r{g}}{\textlengthmark}{\textscy}{\textsubring{r}}} \dicPos{adj}[1]\dicFlx{}[-1] \textbf{1.} \dicSynonym{trúfastur} \dicDirectTranslationCS{spolehlivý, věrný} \dicExampleIS{vera vini sínum tryggur} \dicExampleCS{být věrný svému příteli}  \textbf{2.} \dicSynonym{áreiðanlegur} \dicDirectTranslationCS{spolehlivý, solidní}  \textbf{3.} \dicSynonym{öruggur} \dicDirectTranslationCS{spolehlivý, bezpečný}
\dicEntry[trylla] \dicTerm{tryll|a} \dicIPA{{t\smash{\textsuperscript{h}}}{r}{\textsci}{\textsubring{d}}{l}{a}} \dicPos{v}[2] \dicFlx{(‑ti, ‑t)}[63] \dicFlx{acc} \dicLangCat{hovor.} \dicSynonym{æra\smash{\textsuperscript{2}}} \dicDirectTranslationCS{rozrušit, vytočit, dostat do varu};  \dicIdiom{tryllast}{ \dicPhraseIS{tryllast}} \dicFlx{refl} \dicDirectTranslationCS{vytočit se, dostat se do varu}
\dicEntry[tryllingslega] \dicTerm{tryllings··lega} \dicIPA{{t\smash{\textsuperscript{h}}}{r}{\textsci}{\textsubring{d}}{l}{i}{\ng}{\r{g}}{s}{l}{\textepsilon}{\textbabygamma}{a}} \dicPos{adv} \dicDirectTranslationCS{divoce, nevázaně, šíleně}
\dicEntry[tryllingslegur] \dicTerm{tryllings··legur} \dicIPA{{t\smash{\textsuperscript{h}}}{r}{\textsci}{\textsubring{d}}{l}{i}{\ng}{\r{g}}{s}{l}{\textepsilon}{\textbabygamma}{\textscy}{\textsubring{r}}} \dicPos{adj}[1]\dicFlx{}[-8] \dicDirectTranslationCS{divoký, nevázaný, šílený} \dicExampleIS{tryllingslegur hlátur} \dicExampleCS{divoký smích}
\dicEntry[tryllitæki] \dicTerm{trylli··tæki} \dicIPA{{t\smash{\textsuperscript{h}}}{r}{\textsci}{\textsubring{d}}{l}{\textsci}{t\smash{\textsuperscript{h}}}{a}{i}{\r{\textObardotlessj}}{\textsci}} \dicPos{n}[2] \dicFlx{(‑s, ‑)}[16] \dicLangCat{hovor.} \dicDirectTranslationCS{mašina, kára } \dicIndirectTranslationCS{(o~autě nebo motorce)}
\dicEntry[trylltur] \dicTerm{trylltur} \dicsymFrequent\  \dicIPA{{t\smash{\textsuperscript{h}}}{r}{\textsci}{\textsubring{l}}{\textsubring{d}}{\textscy}{\textsubring{r}}} \dicPos{adj}[1]\dicFlx{}[-10] \dicSynonym{æstur} \dicDirectTranslationCS{divoký, šílený, zuřivý} \dicExampleIS{trylltur hávaði} \dicExampleCS{šílený kravál}
\dicEntry[trýni] \dicTerm{trýni} \dicIPA{{t\smash{\textsuperscript{h}}}{r}{i}{\textlengthmark}{n}{\textsci}} \dicPos{n}[2] \dicFlx{(‑s, ‑)}[14] \dicDirectTranslationCS{čenich, čumák} \dicExampleIS{trýni á hundinum} \dicExampleCS{psí čenich};  \dicPhraseIS{fitja upp á trýnið} \dicLangCat{přen.} \dicDirectTranslationCS{ohrnout nos}
\dicEntry[træði] \dicTerm{træði} \dicIPA{{t\smash{\textsuperscript{h}}}{r}{a}{i}{\textlengthmark}{ð}{\textsci}} \dicPos{v} \dicFlx{con pf sg 1 pers} \dicLink{troða}
\dicEntry[tröf] \dicTerm{tröf} \dicIPA{{t\smash{\textsuperscript{h}}}{r}{\oe}{\textlengthmark}{f}} \dicPos{n} \dicFlx{pl nom} \dicLink{traf}
\dicEntry[tröll] \dicTerm{tröll} \dicIPA{{t\smash{\textsuperscript{h}}}{r}{\oe}{\textsubring{d}}{\textsubring{l}}} \dicPos{n}[2] \dicFlx{(‑s, ‑)}[5] \textbf{1.} \dicFieldCat{pov.} \dicSynonym{risi\smash{\textsuperscript{1}}} \dicDirectTranslationCS{obr, troll}  \textbf{2.} \dicSynonym*{mjög stór maður} \dicDirectTranslationCS{obr, velikán, goliáš}
\dicEntry[tröllasaga] \dicTerm{trölla··|saga} \dicIPA{{t\smash{\textsuperscript{h}}}{r}{\oe}{\textsubring{d}}{l}{a}{s}{a}{\textbabygamma}{a}} \dicPos{f}[1] \dicFlx{(‑sögu, ‑sögur)}[14] \textbf{1.} \dicSynonym*{saga um tröll} \dicDirectTranslationCS{příběh o~obrech}  \textbf{2.} \dicSynonym{lygasaga} \dicDirectTranslationCS{pohádka, historka, báchorka, povídačka} \dicIndirectTranslationCS{(nevěrohodný příběh)}
\dicEntry[tröllastakkur] \dicTerm{trölla··stakk|ur} \dicIPA{{t\smash{\textsuperscript{h}}}{r}{\oe}{\textsubring{d}}{l}{a}{s}{\textsubring{d}}{a}{h}{\r{g}}{\textscy}{\textsubring{r}}} \dicPos{m}[6] \dicFlx{(‑s, ‑ar)}[25] \dicFieldCat{bot.} \dicIndirectTranslationCS{druh všivce} \textit{(l.~{\textLA{Pedicularis flammea}})}  \dicsymPhoto\ 
\dicFigure{ds_image_trollastakkur_0_1.jpg}{Tröllastakkur}{Tröllastakkur - Kim Hansen, GFDL}
\dicEntry[tröllasúra] \dicTerm{trölla··súr|a} \dicIPA{{t\smash{\textsuperscript{h}}}{r}{\oe}{\textsubring{d}}{l}{a}{s}{u}{r}{a}} \dicPos{f}[1] \dicFlx{(‑u, ‑ur)}[7] \dicFieldCat{bot.} \dicDirectTranslationCS{reveň} \textit{(l.~{\textLA{Rheum}})}
\dicEntry[tröllatrú] \dicTerm{trölla··trú} \dicIPA{{t\smash{\textsuperscript{h}}}{r}{\oe}{\textsubring{d}}{l}{a}{t\smash{\textsuperscript{h}}}{r}{u}} \dicPos{f}[4] \dicFlx{(‑ar)}[27] \dicDirectTranslationCS{absolutní víra\,/\addthin důvěra};  \dicPhraseIS{hafa tröllatrú á e‑m} \dicDirectTranslationCS{mít v~(koho) absolutní důvěru, absolutně (komu) důvěřovat}
\dicEntry[tröllaukinn] \dicTerm{tröll··aukinn} \dicIPA{{t\smash{\textsuperscript{h}}}{r}{\oe}{\textsubring{d}}{l}{\oe i}{\r{\textObardotlessj}}{\textsci}{\textsubring{n}}} \dicPos{adj}[6]\dicFlx{}[-2] \dicSynonym{risavaxinn} \dicDirectTranslationCS{obrovský, obrovitý, gigantický} \dicExampleIS{tröllaukin skepna} \dicExampleCS{obrovité zvíře}
\dicEntry[tröllkarl] \dicTerm{tröll··karl} \dicIPA{{t\smash{\textsuperscript{h}}}{r}{\oe}{\textsubring{d}}{\textsubring{l}}{k\smash{\textsuperscript{h}}}{a}{r}{\textsubring{d}}{\textsubring{l}}} \dicPos{m}[4] \dicFlx{(‑s, ‑ar)}[4] \dicFieldCat{pov.} \dicDirectTranslationCS{obr, troll}
\dicEntry[tröllkona] \dicTerm{tröll··kon|a} \dicIPA{{t\smash{\textsuperscript{h}}}{r}{\oe}{\textsubring{d}}{\textsubring{l}}{k\smash{\textsuperscript{h}}}{\textopeno}{n}{a}} \dicPos{f}[1] \dicFlx{(‑u, ‑ur)}[27] \dicFieldCat{pov.} \dicSynonym{skessa} \dicDirectTranslationCS{obryně}
\dicEntry[trönu] \dicTerm{trönu} \dicIPA{{t\smash{\textsuperscript{h}}}{r}{\oe}{\textlengthmark}{n}{\textscy}} \dicPos{f} \dicFlx{sg gen} \dicLink{trana\smash{\textsuperscript{1}}}
\dicEntry[trönuber] \dicTerm{trönu··ber} \dicIPA{{t\smash{\textsuperscript{h}}}{r}{\oe}{\textlengthmark}{n}{\textscy}{\textsubring{b}}{\textepsilon}{\textsubring{r}}} \dicPos{n}[2] \dicFlx{(‑s, ‑)}[13] \dicDirectTranslationCS{klikva} \dicIndirectTranslationCS{(plod)}
\dicEntry[trönuberjalyng] \dicTerm{trönu·berja··lyng} \dicIPA{{t\smash{\textsuperscript{h}}}{r}{\oe}{\textlengthmark}{n}{\textscy}{\textsubring{b}}{\textepsilon}{r}{j}{a}{l}{i}{\ng}{\r{g}}} \dicPos{n}[2] \dicFlx{(‑s)}[2] \dicFieldCat{bot.} \dicDirectTranslationCS{klikva, klikva bahenní} \textit{(l.~{\textLA{Vaccinium oxycoccus}})}  \dicsymPhoto\ 
\dicFigure{ds_image_tronuberjalyng_0_1.jpg}{Trönuberjalyng}{Trönuberjalyng - Christian Fischer, CC BY-SA 3.0}
\dicEntry[tröppu] \dicTerm{tröppu} \dicIPA{{t\smash{\textsuperscript{h}}}{r}{\oe}{h}{\textsubring{b}}{\textscy}} \dicPos{f} \dicFlx{sg gen} \dicLink{trappa\smash{\textsuperscript{1}}}
\dicEntry[Tsjad] \dicTerm{Tsjad} \dicIPA{{\textsubring{d}}{s}{j}{a}{\textlengthmark}{\textsubring{d}}} \dicPos{n}[4] \dicFlx{indecl}[2] \dicFieldCat{geog.} \dicDirectTranslationCS{Čad}
\dicEntry[Tsjadi] \dicTerm{Tsjad|i} \dicIPA{{\textsubring{d}}{s}{j}{a}{\textlengthmark}{\textsubring{d}}{\textsci}} \dicPos{m}[1] \dicFlx{(‑a, ‑ar)}[8] \dicDirectTranslationCS{Čaďan(ka)}
\dicEntry[tsjadneskur] \dicTerm{tsjad··neskur} \dicIPA{{\textsubring{d}}{s}{j}{a}{\textsubring{d}}{n}{\textepsilon}{s}{\r{g}}{\textscy}{\textsubring{r}}} \dicPos{adj}[1]\dicFlx{}[-6] \dicDirectTranslationCS{čadský}
\dicEntry[tuddi] \dicTerm{tudd|i} \dicIPA{{t\smash{\textsuperscript{h}}}{\textscy}{\textsubring{d}}{\textlengthmark}{\textsci}} \dicPos{m}[1] \dicFlx{(‑a, ‑ar)}[1] \textbf{1.} \dicSynonym{boli} \dicDirectTranslationCS{býk}  \textbf{2.} \dicSynonym{fantur} \dicDirectTranslationCS{gauner, ničema}
\dicEntry[tuða] \dicTerm{tuð|a} \dicIPA{{t\smash{\textsuperscript{h}}}{\textscy}{\textlengthmark}{ð}{a}} \dicPos{v}[1] \dicFlx{(‑aði)}[1] \dicSynonym{nöldra} \dicDirectTranslationCS{mrmlat, reptat} \dicExampleIS{tuða í e‑u} \dicExampleCS{(neustále) mluvit do (čeho)}
\dicEntry[tugabrot] \dicTerm{tuga··brot} \dicIPA{{t\smash{\textsuperscript{h}}}{\textscy}{\textlengthmark}{\textbabygamma}{a}{\textsubring{b}}{r}{\textopeno}{\textsubring{d}}} \dicPos{n}[2] \dicFlx{(‑s, ‑)}[5] \dicFieldCat{mat.} \dicDirectTranslationCS{desetinný zlomek}
\dicEntry[tugakerfi] \dicTerm{tuga··kerfi} \dicIPA{{t\smash{\textsuperscript{h}}}{\textscy}{\textlengthmark}{\textbabygamma}{a}{c\smash{\textsuperscript{h}}}{\textepsilon}{r}{v}{\textsci}} \dicPos{n}[2] \dicFlx{(‑s, ‑)}[14] \dicFieldCat{mat.} \dicDirectTranslationCS{desítková\,/\addthin dekadická soustava}
\dicEntry[tugakomma] \dicTerm{tuga··komm|a} \dicIPA{{t\smash{\textsuperscript{h}}}{\textscy}{\textlengthmark}{\textbabygamma}{a}{k\smash{\textsuperscript{h}}}{\textopeno}{m}{a}} \dicPos{f}[1] \dicFlx{(‑u, ‑ur)}[7] \dicFieldCat{mat.} \dicDirectTranslationCS{desetinná čárka}
\dicEntry[tugga] \dicTerm{tugg|a} \dicIPA{{t\smash{\textsuperscript{h}}}{\textscy}{\r{g}}{\textlengthmark}{a}} \dicPos{f}[1] \dicFlx{(‑u, ‑ur)}[7] \textbf{1.} \dicSynonym*{munnfylli} \dicDirectTranslationCS{žvanec}  \textbf{2.} \dicSynonym*{upptugga} \dicDirectTranslationCS{odvar} \dicIndirectTranslationCS{(špatná napodobenina)}
\dicEntry[tuggði] \dicTerm{tuggði} \dicIPA{{t\smash{\textsuperscript{h}}}{\textscy}{\textbabygamma}{ð}{\textsci}} \dicPos{v} \dicFlx{ind pf sg 1 pers} \dicLink{tyggja}
\dicEntry[tuggðum] \dicTerm{tuggðum} \dicIPA{{t\smash{\textsuperscript{h}}}{\textscy}{\textbabygamma}{ð}{\textscy}{\textsubring{m}}} \dicPos{v} \dicFlx{ind pf pl 1 pers} \dicLink{tyggja}
\dicEntry[tuggið] \dicTerm{tuggið} \dicIPA{{t\smash{\textsuperscript{h}}}{\textscy}{\r{\textObardotlessj}}{\textlengthmark}{\textsci}{\texttheta}} \dicPos{v} \dicFlx{supin} \dicLink{tyggja}
\dicEntry[tugur] \dicTerm{tug|ur} \dicsymFrequent\  \dicIPA{{t\smash{\textsuperscript{h}}}{\textscy}{\textlengthmark}{\textbabygamma}{\textscy}{\textsubring{r}}} \dicPos{m}[10] \dicFlx{(‑ar, ‑ir)}[18] \dicDirectTranslationCS{desítka, dekáda} \dicExampleIS{tugir manna} \dicExampleCS{desítky lidí}
\dicEntry[tugþraut] \dicTerm{tug··þraut} \dicIPA{{t\smash{\textsuperscript{h}}}{\textscy}{x}{\texttheta}{r}{\oe i}{\textsubring{d}}} \dicPos{f}[7] \dicFlx{(‑ar)}[3] \dicFieldCat{sport.} \dicDirectTranslationCS{desetiboj}
\dicEntry[tukta] \dicTerm{tukt|a} \dicIPA{{t\smash{\textsuperscript{h}}}{\textscy}{x}{\textsubring{d}}{a}} \dicPos{v}[1] \dicFlx{(‑aði)}[1] \dicFlx{acc} \dicSynonym{refsa} \dicDirectTranslationCS{ukáznit, potrestat (disciplinárně)};  \dicPhraseIS{tukta e‑n til} \dicDirectTranslationCS{ukáznit (koho)}
\dicEntry[tukthús] \dicTerm{tukt··hús} \dicIPA{{t\smash{\textsuperscript{h}}}{\textscy}{x}{\textsubring{d}}{h}{u}{s}} \dicPos{n}[2] \dicFlx{(‑s, ‑)}[5] \dicSynonym{fangelsi} \dicDirectTranslationCS{káznice, vězení} \dicExampleIS{dæma hann í tukthús} \dicExampleCS{odsoudit ho do káznice}
\dicEntry[tuldra] \dicTerm{tuldr|a} \dicsymFrequent\  \dicIPA{{t\smash{\textsuperscript{h}}}{\textscy}{l}{\textsubring{d}}{r}{a}} \dicPos{v}[1] \dicFlx{(‑aði)}[1] \dicFlx{acc} \dicSynonym{tauta} \dicDirectTranslationCS{(za)huhňat, (za)mumlat} \dicExampleIS{tuldra e‑ð hálfum hljóðum fyrir munni sér} \dicExampleCS{zamumlat si (co) potichu pro sebe};  \dicPhraseIS{tuldra í barm sér} \dicLangCat{přen.} \dicDirectTranslationCS{mumlat si pod vousy}
\dicEntry[tuldur] \dicTerm{tuldur} \dicIPA{{t\smash{\textsuperscript{h}}}{\textscy}{l}{\textsubring{d}}{\textscy}{\textsubring{r}}} \dicPos{n}[2] \dicFlx{(‑s)}[28] \dicSynonym{taut} \dicDirectTranslationCS{(za)huhňání, (za)mumlání}
\dicEntry[tundur] \dicTerm{tundur} \dicIPA{{t\smash{\textsuperscript{h}}}{\textscy}{n}{\textsubring{d}}{\textscy}{\textsubring{r}}} \dicPos{n}[2] \dicFlx{(‑s)}[28] \textbf{1.} \dicSynonym*{kveikiefni} \dicDirectTranslationCS{troud, podpal}  \textbf{2.} \dicSynonym{sprengiefni} \dicDirectTranslationCS{třaskavina, výbušnina};  \dicIdiom{tundur}{ \dicPhraseIS{það er allt á tjá og tundri}} \dicLangCat{přen.} \dicDirectTranslationCS{všechno je vzhůru nohama}
\dicEntry[tundurdufl] \dicTerm{tundur··dufl} \dicIPA{{t\smash{\textsuperscript{h}}}{\textscy}{n}{\textsubring{d}}{\textscy}{r}{\textsubring{d}}{\textscy}{\textsubring{b}}{\textsubring{l}}} \dicPos{n}[2] \dicFlx{(‑s, ‑)}[5] \dicDirectTranslationCS{mořská mina}
\dicEntry[tundurduflaslæðari] \textls[19]{\dicTerm{tundur·dufla··slæð·ar|i} \dicIPA{{t\smash{\textsuperscript{h}}}\-{\textscy}\-{n}\-{\textsubring{d}}\-{\textscy}\-{r}\-{\textsubring{d}}\-{\textscy}\-{\textsubring{b}}\-{l}\-{a}\-{s}\-{\textsubring{d}}\-{l}\-{a}\-{i}\-{ð}\-{a}\-{r}\-{\textsci}\-} \dicPos{m}[1] \dicFlx{(‑a, ‑ar)}[13]} \dicFieldCat{nám.} \dicDirectTranslationCS{minolovka}
\dicEntry[tundurskeyti] \dicTerm{tundur··skeyti} \dicIPA{{t\smash{\textsuperscript{h}}}{\textscy}{n}{\textsubring{d}}{\textscy}{\textsubring{r}}{s}{\r{\textObardotlessj}}{ei}{\textsubring{d}}{\textsci}} \dicPos{n}[2] \dicFlx{(‑s, ‑)}[14] \dicFieldCat{nám.} \dicDirectTranslationCS{torpédo}
\dicEntry[tunga] \dicTerm{tung|a} \dicsymFrequent\  \dicIPA{{t\smash{\textsuperscript{h}}}{u}{\ng}{\r{g}}{a}} \dicPos{f}[1] \dicFlx{(‑u, ‑ur)}[13] \textbf{1.} \dicFieldCat{anat.} \dicDirectTranslationCS{jazyk} \dicIndirectTranslationCS{(v~ústech)}  \textbf{2.} \dicSynonym{tungumál} \dicDirectTranslationCS{jazyk, řeč, mluva} \dicExampleIS{þróun íslenskrar tungu fram um aldir} \dicExampleCS{vývoj islandského jazyka napříč stoletími}  \textbf{3.} \dicSynonym*{landspilda} \dicDirectTranslationCS{šíje, úžina}  \textbf{4.} \dicSynonym*{tunga á skó} \dicDirectTranslationCS{jazyk} \dicIndirectTranslationCS{(boty)};  \dicIdiom{tunga}{ \dicPhraseIS{e‑m vefst tunga um tönn}} \dicFlx{impers} \dicDirectTranslationCS{(kdo) si šlape na jazyk}; { \dicPhraseIS{gæta tungu sinnar}} \dicDirectTranslationCS{dávat si pozor na jazyk};  \dicPhraseIS{skæðar tungur} \dicLangCat{přen.} \dicDirectTranslationCS{zlé jazyky};  \dicPhraseIS{það leikur ekki á tveim tungum að (e‑að gerist)} \dicFlx{impers} \dicDirectTranslationCS{není pochyb o~tom, že se ((co) stane)}
\dicEntry[tungl] \dicTerm{tungl} \dicsymFrequent\  \dicIPA{{t\smash{\textsuperscript{h}}}{u}{\ng}{\textsubring{l}}} \dicPos{n}[2] \dicFlx{(‑s, ‑)}[5] \textbf{1.} \dicFieldCat{astro.} \dicDirectTranslationCS{měsíc, Měsíc} \dicIndirectTranslationCS{(nebeské těleso obíhající kolem Země)};  \dicPhraseIS{fullt tungl} \dicDirectTranslationCS{úplněk};  \dicPhraseIS{minnkandi tungl} \dicDirectTranslationCS{ubývající měsíc};  \dicPhraseIS{nýtt tungl} \dicDirectTranslationCS{nov};  \dicPhraseIS{vaxandi tungl} \dicDirectTranslationCS{dorůstající měsíc}  \textbf{2.} \dicFieldCat{astro.} \dicSynonym{fylgihnöttur} \dicDirectTranslationCS{měsíc} \dicIndirectTranslationCS{(oběžnice jiných těles)}
\dicEntry[tunglferð] \dicTerm{tungl··ferð} \dicIPA{{t\smash{\textsuperscript{h}}}{u}{\ng}{l}{f}{\textepsilon}{r}{\texttheta}} \dicPos{f}[7] \dicFlx{(‑ar, ‑ir)}[1] \dicDirectTranslationCS{cesta\,/\addthin let na měsíc}
\dicEntry[tungljurt] \dicTerm{tungl··jurt} \dicIPA{{t\smash{\textsuperscript{h}}}{u}{\ng}{l}{j}{\textscy}{\textsubring{r}}{\textsubring{d}}} \dicPos{f}[7] \dicFlx{(‑ar, ‑ir)}[1] \dicFieldCat{bot.} \dicDirectTranslationCS{vratička měsíční} \textit{(l.~{\textLA{Botrychium lunaria}})}  \dicsymPhoto\ 
\dicFigure{ds_image_tungljurt_0_1.jpg}{Tungljurt}{Tungljurt - Abalg, GFDL}
\dicEntry[tunglkoma] \dicTerm{tungl··kom|a} \dicIPA{{t\smash{\textsuperscript{h}}}{u}{\ng}{l}{k\smash{\textsuperscript{h}}}{\textopeno}{m}{a}} \dicPos{f}[1] \dicFlx{(‑u)}[5] \dicSynonym*{nýtt tungl} \dicDirectTranslationCS{nov}
\dicEntry[tunglmánuður] \dicTerm{tungl··mán·|uður} \dicIPA{{t\smash{\textsuperscript{h}}}{u}{\ng}{l}{m}{au}{n}{\textscy}{ð}{\textscy}{\textsubring{r}}} \dicPos{m}[10] \dicFlx{(‑aðar, ‑uðir)}[42] \dicFieldCat{astro.} \dicDirectTranslationCS{lunární měsíc}
\dicEntry[tunglmyrkvi] \dicTerm{tungl··myrkv|i} \dicIPA{{t\smash{\textsuperscript{h}}}{u}{\ng}{l}{m}{\textsci}{\textsubring{r}}{\r{g}}{v}{\textsci}} \dicPos{m}[1] \dicFlx{(‑a, ‑ar)}[1] \dicFieldCat{astro.} \dicDirectTranslationCS{zatmění měsíce}
\dicEntry[tunglsbirta] \dicTerm{tungls··birt|a} \dicIPA{{t\smash{\textsuperscript{h}}}{u}{\ng}{l}{s}{\textsubring{b}}{\textsci}{\textsubring{r}}{\textsubring{d}}{a}} \dicPos{f}[1] \dicFlx{(‑u)}[5] \dicDirectTranslationCS{měsíční svit}
\dicEntry[tunglskin] \dicTerm{tungl··skin} \dicIPA{{t\smash{\textsuperscript{h}}}{u}{\ng}{l}{s}{\r{\textObardotlessj}}{\textsci}{\textsubring{n}}} \dicPos{n}[2] \dicFlx{(‑s)}[2] \dicSynonym{tunglsbirta} \dicDirectTranslationCS{měsíční svit}
\dicEntry[tungubrjótur] \dicTerm{tungu··brjót|ur} \dicIPA{{t\smash{\textsuperscript{h}}}{u}{\ng}{\r{g}}{\textscy}{\textsubring{b}}{r}{j}{ou}{\textsubring{d}}{\textscy}{\textsubring{r}}} \dicPos{m}[6] \dicFlx{(‑s, ‑ar)}[10] \dicDirectTranslationCS{jazykolam}
\dicEntry[tungubroddur] \dicTerm{tungu··brodd|ur} \dicIPA{{t\smash{\textsuperscript{h}}}{u}{\ng}{\r{g}}{\textscy}{\textsubring{b}}{r}{\textopeno}{\textsubring{d}}{\textscy}{\textsubring{r}}} \dicPos{m}[6] \dicFlx{(‑s, ‑ar)}[4] \dicDirectTranslationCS{špička\,/\addthin hrot jazyka}
\dicEntry[tungulangur] \dicTerm{tungu··|langur} \dicIPA{{t\smash{\textsuperscript{h}}}{u}{\ng}{\r{g}}{\textscy}{l}{au}{\ng}{\r{g}}{\textscy}{\textsubring{r}}} \dicPos{adj}[10] \dicFlx{(f ‑löng, comp ‑lengri, sup ‑lengstur)}[2] \textbf{1.} \dicSynonym*{með langa tungu} \dicDirectTranslationCS{mající dlouhý jazyk}  \textbf{2.} \dicSynonym*{djarfmæltur} \dicDirectTranslationCS{prostořeký}
\dicEntry[tungulipur] \dicTerm{tungu··lipur} \dicIPA{{t\smash{\textsuperscript{h}}}{u}{\ng}{\r{g}}{\textscy}{l}{\textsci}{\textsubring{b}}{\textscy}{\textsubring{r}}} \dicPos{adj}[9]\dicFlx{}[-5] \dicDirectTranslationCS{výřečný, mající obratný jazyk}
\dicEntry[tungumál] \dicTerm{tungu··mál} \dicsymFrequent\  \dicIPA{{t\smash{\textsuperscript{h}}}{u}{\ng}{\r{g}}{\textscy}{m}{au}{\textsubring{l}}} \dicPos{n}[2] \dicFlx{(‑s, ‑)}[5] \dicSynonym{mál\smash{\textsuperscript{1}}} \dicDirectTranslationCS{jazyk, řeč} \dicExampleIS{tala mörg tungumál} \dicExampleCS{hovořit mnoha jazyky}
\begin{xtolerant}{}{1pt}
\dicEntry[tungumálakennari] \dicTerm{tungu·mála··kenn·ar|i}\addthinM\textls[-10]{\dicIPA{{t\smash{\textsuperscript{h}}}\-{u}\-{\ng}\-{\r{g}}\-{\textscy}\-{m}\-{au}\-{l}\-{a}\-{c\smash{\textsuperscript{h}}}\-{\textepsilon}\-{n}\-{a}\-{r}\-{\textsci}}} \dicPos{m}[1]\addthinS\dicFlx{(‑a, ‑ar)}[13] \dicDirectTranslationCS{učitel(ka) jazyků}
\end{xtolerant}
\dicEntry[tungumálakennsla] \textls[10]{\dicTerm{tungu·mála··kennsl|a} \dicIPA{{t\smash{\textsuperscript{h}}}{u}{\ng}{\r{g}}{\textscy}{m}{au}{l}{a}{c\smash{\textsuperscript{h}}}{\textepsilon}{n}{s}{\textsubring{d}}{l}{a}} \dicPos{f}[1] \dicFlx{(‑u)}[5]} \dicFieldCat{škol.} \dicDirectTranslationCS{výuka jazyků}
\begin{xtolerant}{}{1pt}
\dicEntry[tungumálakunnátta] \dicTerm{tungu·mála··kunn·átt|a}\addthinS\textls[-13]{\dicIPA{{t\smash{\textsuperscript{h}}}\-{u}\-{\ng}\-{\r{g}}\-{\textscy}\-{m}\-{au}\-{l}\-{a}\-{k\smash{\textsuperscript{h}}}\-{\textscy}\-{n}\-{au}\-{h}\-{\textsubring{d}}\-{a}}} \dicPos{f}[1]\addthinSS\dicFlx{(‑u)}[5] \dicDirectTranslationCS{jazykové znalosti}
\end{xtolerant}
\dicEntry[tungumjúkur] \dicTerm{tungu··|mjúkur} \dicIPA{{t\smash{\textsuperscript{h}}}{u}{\ng}{\r{g}}{\textscy}{m}{j}{u}{\r{g}}{\textscy}{\textsubring{r}}} \dicPos{adj}[10] \dicFlx{(comp ‑mýkri, sup ‑mýkstur)}[11] \dicDirectTranslationCS{vemlouvavý, lichotný, úlisný}
\dicEntry[tungutak] \dicTerm{tungu··tak} \dicIPA{{t\smash{\textsuperscript{h}}}{u}{\ng}{\r{g}}{\textscy}{t\smash{\textsuperscript{h}}}{a}{\r{g}}} \dicPos{n}[2] \dicFlx{(‑s)}[2] \dicDirectTranslationCS{mluva, vyjadřování} \dicExampleIS{liðugt tungutak} \dicExampleCS{volné vyjadřování}
\dicEntry[tungutækni] \dicTerm{tungu··tækn|i} \dicIPA{{t\smash{\textsuperscript{h}}}{u}{\ng}{\r{g}}{\textscy}{t\smash{\textsuperscript{h}}}{a}{i}{h}{\r{g}}{n}{\textsci}} \dicPos{f}[3] \dicFlx{(‑i)}[3] \dicFieldCat{poč.} \dicDirectTranslationCS{jazykové technologie}
\dicEntry[tunna] \dicTerm{tunn|a} \dicIPA{{t\smash{\textsuperscript{h}}}{\textscy}{n}{\textlengthmark}{a}} \dicPos{f}[1] \dicFlx{(‑u, ‑ur)}[7] \textbf{1.} \dicSynonym*{sívalt ílát} \dicDirectTranslationCS{sud, bečka, barel}  \textbf{2.} \dicSynonym{sorptunna} \dicDirectTranslationCS{popelnice}
\dicEntry[turn] \dicTerm{turn} \dicsymFrequent\  \dicIPA{{t\smash{\textsuperscript{h}}}{\textscy}{r}{\textsubring{d}}{\textsubring{n}}} \dicPos{m}[4] \dicFlx{(‑s, ‑ar)}[3] \textbf{1.} \dicDirectTranslationCS{věž} \dicExampleIS{litast um uppi í turninum} \dicExampleCS{rozhlížet se z~věže}  \textbf{2.} \dicDirectTranslationCS{rozhledna, pozorovatelna}
\dicEntry[turnfálki] \dicTerm{turn··fálk|i} \dicIPA{{t\smash{\textsuperscript{h}}}{\textscy}{r}{\textsubring{d}}{\textsubring{n}}{f}{au}{\textsubring{l}}{\r{\textObardotlessj}}{\textsci}} \dicPos{m}[1] \dicFlx{(‑a, ‑ar)}[1] \dicFieldCat{zool.} \dicDirectTranslationCS{poštolka, poštolka obecná} \textit{(l.~{\textLA{Falco tinnunculus}})}  \dicsymPhoto\ 
\dicFigure{6176.jpg}{Turnfálki}{Turnfálki - Deml Miroslav, Biolib, Copyright/CC-BY}
\dicEntry[turnspíra] \dicTerm{turn··spír|a} \dicIPA{{t\smash{\textsuperscript{h}}}{\textscy}{r}{\textsubring{d}}{\textsubring{n}}{s}{\textsubring{b}}{i}{r}{a}} \dicPos{f}[1] \dicFlx{(‑u, ‑ur)}[7] \dicDirectTranslationCS{špička věže, věžička}
\dicEntry[turtildúfa] \dicTerm{turtil··dúf|a} \dicIPA{{t\smash{\textsuperscript{h}}}{\textscy}{\textsubring{r}}{\textsubring{d}}{\textsci}{l}{\textsubring{d}}{u}{v}{a}} \dicPos{f}[1] \dicFlx{(‑u, ‑ur)}[13] \dicFieldCat{zool.} \dicDirectTranslationCS{hrdlička, hrdlička divoká} \textit{(l.~{\textLA{Streptopelia turtur}})}  \dicsymPhoto\ 
\dicFigure{ds_image_turtildufa_0_2.jpg}{Turtildúfa}{Turtildúfa - David King, CC BY 2.0}
\dicEntry[tusk] \dicTerm{tusk} \dicIPA{{t\smash{\textsuperscript{h}}}{\textscy}{s}{\r{g}}} \dicPos{n}[2] \dicFlx{(‑s)}[2] \dicDirectTranslationCS{pranice, přetahovaná};  \dicPhraseIS{vera til í tuskið} \dicLangCat{přen.} \dicDirectTranslationCS{být připravený k~boji}
\dicEntry[tuska] \dicTerm{tusk|a} \dicsymFrequent\  \dicIPA{{t\smash{\textsuperscript{h}}}{\textscy}{s}{\r{g}}{a}} \dicPos{f}[1] \dicFlx{(‑u, ‑ur)}[13] \textbf{1.} \dicDirectTranslationCS{hadr, hadřík, prachovka} \dicExampleIS{gólftuska} \dicExampleCS{hadr na podlahu}  \textbf{2.} \dicSynonym{ræfill} \dicDirectTranslationCS{trhan(ka)} \dicIndirectTranslationCS{(zvláště jako poslední část složeného slova)} \dicExampleIS{ræfilstuska} \dicExampleCS{chudák}  \textbf{3.} \dicSynonym*{ónýt flík} \dicDirectTranslationCS{hadr, cár (oblečení ap.)};  \dicIdiom{tuska}{ \dicPhraseIS{vera eins og undin tuska}} \dicLangCat{přen.} \dicDirectTranslationCS{být úplně vyždímaný\,/\addthin vyšťavený};  \dicPhraseIS{það kjaftar á e‑m hver tuska} \dicFlx{impers} \dicLangCat{přen.} \dicDirectTranslationCS{pusa se (komu) nezastaví};  \dicPhraseIS{það er líf í tuskunum} \dicLangCat{přen.} \dicDirectTranslationCS{žije to, pulzuje to životem, má to šmrnc}
\dicEntry[tuskulegur] \dicTerm{tusku··legur} \dicIPA{{t\smash{\textsuperscript{h}}}{\textscy}{s}{\r{g}}{\textscy}{l}{\textepsilon}{\textbabygamma}{\textscy}{\textsubring{r}}} \dicPos{adj}[1]\dicFlx{}[-8] \textbf{1.} \dicSynonym{druslulegur} \dicDirectTranslationCS{otrhaný, ošoupaný} \dicIndirectTranslationCS{(o~oblečení)} \dicExampleIS{tuskuleg flík} \dicExampleCS{otrhaný kus oblečení}  \textbf{2.} \dicSynonym*{linjulegur} \dicDirectTranslationCS{strhaný, schlíplý} \dicExampleIS{vera tuskulegur í framan} \dicExampleCS{být strhaný ve tváři}
\dicEntry[tussa] \dicTerm{tuss|a} \dicIPA{{t\smash{\textsuperscript{h}}}{\textscy}{s}{\textlengthmark}{a}} \dicPos{f}[1] \dicFlx{(‑u, ‑ur)}[7] \dicLangCat{hrub.} \dicDirectTranslationCS{píča, kunda}
\dicEntry[tussulegur] \dicTerm{tussu··legur} \dicIPA{{t\smash{\textsuperscript{h}}}{\textscy}{s}{\textlengthmark}{\textscy}{l}{\textepsilon}{\textbabygamma}{\textscy}{\textsubring{r}}} \dicPos{adj}[1]\dicFlx{}[-8] \textbf{1.} \dicSynonym{lítilfjörlegur} \dicDirectTranslationCS{nanicovatý}  \textbf{2.} \dicSynonym{leiðinlegur} \dicDirectTranslationCS{otravný} \dicIndirectTranslationCS{(o~počasí)}
\dicEntry[tuttugasti] \dicTerm{tut·tug··asti} \dicIPA{{t\smash{\textsuperscript{h}}}{\textscy}{h}{\textsubring{d}}{\textscy}{\textbabygamma}{a}{s}{\textsubring{d}}{\textsci}} \dicPos{num}[6] \dicFlx{ord} \dicDirectTranslationCS{dvacátý}
\dicEntry[tuttugu] \dicTerm{tut··tugu} \dicsymFrequent\  \dicIPA{{t\smash{\textsuperscript{h}}}{\textscy}{h}{\textsubring{d}}{\textscy}{\textbabygamma}{\textscy}} \dicPos{num} \dicFlx{indecl} \dicDirectTranslationCS{dvacet} \dicExampleIS{um tuttugu mínútum síðar} \dicExampleCS{přibližně po dvaceti minutách}
\dicEntry[tuttugukall] \dicTerm{tut·tugu··kall} \dicIPA{{t\smash{\textsuperscript{h}}}{\textscy}{h}{\textsubring{d}}{\textscy}{\textbabygamma}{\textscy}{k\smash{\textsuperscript{h}}}{a}{\textsubring{d}}{\textsubring{l}}} \dicPos{m}[4] \dicFlx{(‑s, ‑ar)}[4] \textbf{1.} \dicLangCat{hovor.} \dicDirectTranslationCS{dvacet (korun)} \dicIndirectTranslationCS{(částka)}  \textbf{2.} \dicLangCat{hovor.} \dicDirectTranslationCS{dvacka} \dicIndirectTranslationCS{(mince\,/\addthin bankovka)}
\dicEntry[túba] \dicTerm{túb|a}\dicTerm{, túpa} \dicIPA{{t\smash{\textsuperscript{h}}}\-{u}\-{\textlengthmark}\-{\textsubring{b}}\-{a}\-} \dicPos{f}[1] \dicFlx{(‑u, ‑ur)}[7] \textbf{1.} \dicFieldCat{hud.} \dicDirectTranslationCS{tuba}  \textbf{2.} \dicDirectTranslationCS{tuba} \dicIndirectTranslationCS{(válcovitý obal se šroubovacím závěrem)} \dicExampleIS{olíulitir í túbum} \dicExampleCS{olejové barvy v~tubách}
\dicEntry[túff] \dicTerm{túff} \dicIPA{{t\smash{\textsuperscript{h}}}{u}{f}{\textlengthmark}} \dicPos{n}[2] \dicFlx{(‑s)}[2] \dicFieldCat{geol.} \dicDirectTranslationCS{tuf}
\dicEntry[túkall] \dicTerm{tú··kall} \dicIPA{{t\smash{\textsuperscript{h}}}{u}{\textlengthmark}{k\smash{\textsuperscript{h}}}{a}{\textsubring{d}}{\textsubring{l}}} \dicPos{m}[4] \dicFlx{(‑s, ‑ar)}[4] \dicLangCat{zast.} \dicDirectTranslationCS{dvě (koruny)} \dicIndirectTranslationCS{(částka)}
\dicEntry[túkani] \dicTerm{túkan|i} \dicIPA{{t\smash{\textsuperscript{h}}}{u}{\textlengthmark}{\r{g}}{a}{n}{\textsci}} \dicPos{m}[1] \dicFlx{(‑a, ‑ar)}[8] \dicFieldCat{zool.} \dicSynonym*{piparfugl} \dicDirectTranslationCS{tukan} \textit{(l.~{\textLA{Ramphastos}})}  \dicsymPhoto\ 
\dicFigure{ds_image_tukani_0_2.jpg}{Túkani}{Túkani - Donar Reiskoffer, GFDL}
\dicEntry[túli] \dicTerm{túl|i} \dicIPA{{t\smash{\textsuperscript{h}}}{u}{\textlengthmark}{l}{\textsci}} \dicPos{m}[1] \dicFlx{(‑a, ‑ar)}[1] \dicLangCat{han.} \dicSynonym{munnur} \dicDirectTranslationCS{papula, huba, držka};  \dicPhraseIS{halda á sér túlanum} \dicDirectTranslationCS{sklapnout, zmlknout}
\dicEntry[túlipani] \dicTerm{túlipan|i} \dicIPA{{t\smash{\textsuperscript{h}}}{u}{\textlengthmark}{l}{\textsci}{\textsubring{b}}{a}{n}{\textsci}} \dicPos{m}[1] \dicFlx{(‑a, ‑ar)}[8] \dicFieldCat{bot.} \dicDirectTranslationCS{tulipán} \textit{(l.~{\textLA{Tulipa}})}  \rlap{\dicsymPhoto\ }
\dicFigure{ds_image_tulipani_0_1.jpg}{Túlipani}{Túlipani - John O'Neill, CC BY-SA 3.0}
\dicEntry[túlka] \dicTerm{túlk|a} \dicIPA{{t\smash{\textsuperscript{h}}}{u}{\textsubring{l}}{\r{g}}{a}} \dicPos{v}[1] \dicFlx{(‑aði)}[1] \dicFlx{acc} \textbf{1.} \dicSynonym{skýra} \dicDirectTranslationCS{vysvětlit, vysvětlovat, objasnit, objasňovat, vyjasnit, vyjasňovat} \dicExampleIS{túlka mál e‑rs} \dicExampleCS{vysvětlit (čí) záležitost}  \textbf{2.} \dicSynonym*{birta eðli} \dicDirectTranslationCS{(obsahově) vyložit\,/\addthin vykládat, interpretovat (dílo ap.)}  \textbf{3.} \dicSynonym{þýða} \dicDirectTranslationCS{(pře)tlumočit} \dicExampleIS{túlka ræðuna á íslensku} \dicExampleCS{přetlumočit řeč do islandštiny}  \textbf{4.} \dicFieldCat{poč.} \dicDirectTranslationCS{interpretovat};  \dicPhraseIS{túlkað forritunarmál} \dicFieldCat{poč.} \dicDirectTranslationCS{interpretovaný jazyk}
\dicEntry[túlkafræði] \dicTerm{túlka··fræð|i} \dicIPA{{t\smash{\textsuperscript{h}}}{u}{\textsubring{l}}{\r{g}}{a}{f}{r}{a}{i}{ð}{\textsci}} \dicPos{f}[3] \dicFlx{(‑i)}[3] \dicDirectTranslationCS{tlumočnictví} \dicExampleIS{nám í túlkafræði} \dicExampleCS{studium tlumočnictví}
\dicEntry[túlkaklefi] \dicTerm{túlka··klef|i} \dicIPA{{t\smash{\textsuperscript{h}}}{u}{\textsubring{l}}{\r{g}}{a}{k\smash{\textsuperscript{h}}}{l}{\textepsilon}{v}{\textsci}} \dicPos{m}[1] \dicFlx{(‑a, ‑ar)}[1] \dicDirectTranslationCS{tlumočnická kabina}
\dicEntry[túlkandi] \dicTerm{túlk··|andi} \dicIPA{{t\smash{\textsuperscript{h}}}{u}{\textsubring{l}}{\r{g}}{a}{n}{\textsubring{d}}{\textsci}} \dicPos{m}[2] \dicFlx{(‑anda, ‑endur)}[1] \textbf{1.} \dicDirectTranslationCS{tlumočník, tlumočnice}  \textbf{2.} \dicDirectTranslationCS{interpret(ka)}
\dicEntry[túlkun] \dicTerm{túlk|un} \dicIPA{{t\smash{\textsuperscript{h}}}{u}{\textsubring{l}}{\r{g}}{\textscy}{\textsubring{n}}} \dicPos{f}[7] \dicFlx{(‑unar, ‑anir)}[8] \textbf{1.} \dicSynonym{framsetning} \dicDirectTranslationCS{interpretace, vyjasnění, výklad} \dicExampleIS{túlkun á orðum hans} \dicExampleCS{interpretace jeho slov}  \textbf{2.} \dicDirectTranslationCS{tlumočení} \dicExampleIS{túlkun á ensku} \dicExampleCS{tlumočení do angličtiny};  \dicPhraseIS{samtímis túlkun} \dicDirectTranslationCS{simultánní tlumočení}  \textbf{3.} \dicDirectTranslationCS{interpretace, provedení, ztvárnění (role ap.)}
\dicEntry[túlkunarfræði] \dicTerm{túlkunar··fræð|i} \dicIPA{{t\smash{\textsuperscript{h}}}{u}{\textsubring{l}}{\r{g}}{\textscy}{n}{a}{\textsubring{r}}{f}{r}{a}{i}{ð}{\textsci}} \dicPos{f}[3] \dicFlx{(‑i)}[3] \dicFieldCat{filos.} \dicDirectTranslationCS{hermeneutika}
\dicEntry[túlkur] \dicTerm{túlk|ur} \dicIPA{{t\smash{\textsuperscript{h}}}{u}{\textsubring{l}}{\r{g}}{\textscy}{\textsubring{r}}} \dicPos{m}[6] \dicFlx{(‑s, ‑ar)}[24] \textbf{1.} \dicDirectTranslationCS{tlumočník, tlumočnice}  \textbf{2.} \dicFieldCat{poč.} \dicDirectTranslationCS{interpret}
\dicEntry[tún] \dicTerm{tún} \dicsymFrequent\  \dicIPA{{t\smash{\textsuperscript{h}}}{u}{\textlengthmark}{\textsubring{n}}} \dicPos{n}[2] \dicFlx{(‑s, ‑)}[5] \dicDirectTranslationCS{louka, palouk, humno} \dicExampleIS{Ég var alein úti á grænu túni.} \dicExampleCS{Byla jsem sama na zelené louce.}
\dicEntry[túndra] \dicTerm{túndr|a} \dicIPA{{t\smash{\textsuperscript{h}}}{u}{n}{\textsubring{d}}{r}{a}} \dicPos{f}[1] \dicFlx{(‑u, ‑ur)}[7] \dicFieldCat{geog.} \dicDirectTranslationCS{tundra}
\dicEntry[túnfiskur] \dicTerm{tún··fisk|ur} \dicIPA{{t\smash{\textsuperscript{h}}}{u}{n}{f}{\textsci}{s}{\r{g}}{\textscy}{\textsubring{r}}} \dicPos{m}[6] \dicFlx{(‑s, ‑ar)}[10] \dicFieldCat{zool.} \dicDirectTranslationCS{tuňák obecný} \textit{(l.~{\textLA{Thunnus thynnus}})}
\dicEntry[túnfífill] \dicTerm{tún··fíf|ill} \dicIPA{{t\smash{\textsuperscript{h}}}{u}{n}{f}{i}{v}{\textsci}{\textsubring{d}}{\textsubring{l}}} \dicPos{m}[6] \dicFlx{(‑ils, ‑lar)}[35] \dicFieldCat{bot.} \dicDirectTranslationCS{pampeliška, smetanka} \textit{(l.~{\textLA{Taraxacum}})}  \dicsymPhoto\ 
\dicFigure{6912.jpg}{Túnfífill}{Túnfífill - Deml Miroslav, Biolib, Copyright/CC-BY}
\dicEntry[túngarður] \dicTerm{tún··garð|ur} \dicIPA{{t\smash{\textsuperscript{h}}}{u}{n}{\r{g}}{a}{r}{ð}{\textscy}{\textsubring{r}}} \dicPos{m}[6] \dicFlx{(‑s, ‑ar)}[5] \dicDirectTranslationCS{ohrada louky}
\dicEntry[Túnis] \dicTerm{Túnis} \dicIPA{{t\smash{\textsuperscript{h}}}{u}{\textlengthmark}{n}{\textsci}{s}} \dicPos{n}[4] \dicFlx{indecl}[2] \dicFieldCat{geog.} \dicDirectTranslationCS{Tunisko}
\dicEntry[Túnisbúi] \dicTerm{Túnis··bú|i} \dicIPA{{t\smash{\textsuperscript{h}}}{u}{\textlengthmark}{n}{\textsci}{s}{\textsubring{b}}{u}{\textsci}} \dicPos{m}[1] \dicFlx{(‑a, ‑ar)}[1] \dicLink{Túnisi}
\dicEntry[Túnisi] \dicTerm{Túnis|i}\dicTerm{, Túnisbúi} \dicIPA{{t\smash{\textsuperscript{h}}}\-{u}\-{\textlengthmark}\-{n}\-{\textsci}\-{s}\-{\textsci}\-} \dicPos{m}[1] \dicFlx{(‑a, ‑ar)}[1] \dicDirectTranslationCS{Tunisan(ka)}
\dicEntry[túniskur] \dicTerm{túniskur} \dicIPA{{t\smash{\textsuperscript{h}}}{u}{\textlengthmark}{n}{\textsci}{s}{\r{g}}{\textscy}{\textsubring{r}}} \dicPos{adj}[1]\dicFlx{}[-6] \dicDirectTranslationCS{tuniský}
\dicEntry[túnsúra] \dicTerm{tún··súr|a} \dicIPA{{t\smash{\textsuperscript{h}}}{u}{n}{s}{u}{r}{a}} \dicPos{f}[1] \dicFlx{(‑u, ‑ur)}[7] \dicFieldCat{bot.} \dicDirectTranslationCS{šťovík kyselý} \textit{(l.~{\textLA{Rumex acetosa}})}  \dicsymPhoto\ 
\dicFigure{ds_image_tunsura_0_1.jpg}{Túnsúra}{Túnsúra - Christian Hummert, CC BY-SA 2.5}
\dicEntry[túnvingull] \dicTerm{tún··ving|ull} \dicIPA{{t\smash{\textsuperscript{h}}}{u}{n}{v}{i}{\ng}{\r{g}}{\textscy}{\textsubring{d}}{\textsubring{l}}} \dicPos{m}[6] \dicFlx{(‑uls, ‑lar)}[31] \dicFieldCat{bot.} \dicDirectTranslationCS{kostřava červená} \textit{(l.~{\textLA{Festuca rubra}})}
\dicEntry[túpa] \dicTerm{túp|a} \dicIPA{{t\smash{\textsuperscript{h}}}{u}{\textlengthmark}{\textsubring{b}}{a}} \dicPos{f}[1] \dicFlx{(‑u, ‑ur)}[19] \dicLink{túba}
\dicEntry[túr] \dicTerm{túr} \dicIPA{{t\smash{\textsuperscript{h}}}{u}{\textlengthmark}{\textsubring{r}}} \dicPos{m}[4] \dicFlx{(‑s, ‑ar)}[14] \textbf{1.} \dicSynonym{ferð} \dicDirectTranslationCS{výlet, cesta, vyjížďka, výprava}  \textbf{2.} \dicSynonym*{útivist skips} \dicDirectTranslationCS{(rybolovná) výprava};  \dicPhraseIS{vera á túr} {\textbf{a.}} \dicLangCat{přen.} \dicDirectTranslationCS{být na tahu, kalit};  {\textbf{b.}} \dicDirectTranslationCS{mít měsíčky}
\dicEntry[túramaður] \dicTerm{túra··|maður} \dicIPA{{t\smash{\textsuperscript{h}}}{u}{\textlengthmark}{r}{a}{m}{a}{ð}{\textscy}{\textsubring{r}}} \dicPos{m}[13] \dicFlx{(‑manns, ‑menn)}[2] \dicDirectTranslationCS{pařič(ka), kalič(ka)}
\dicEntry[túrban] \dicTerm{túrban} \dicIPA{{t\smash{\textsuperscript{h}}}{u}{r}{\textsubring{b}}{a}{\textsubring{n}}} \dicPos{m}[4] \dicFlx{(‑s, ‑ar)}[11] \dicDirectTranslationCS{turban}
\dicEntry[túrbína] \dicTerm{túrbín|a} \dicIPA{{t\smash{\textsuperscript{h}}}{u}{r}{\textsubring{b}}{i}{n}{a}} \dicPos{f}[1] \dicFlx{(‑u, ‑ur)}[7] \dicDirectTranslationCS{turbína}
\dicEntry[túristi] \dicTerm{túr··ist|i} \dicIPA{{t\smash{\textsuperscript{h}}}{u}{\textlengthmark}{r}{\textsci}{s}{\textsubring{d}}{\textsci}} \dicPos{m}[1] \dicFlx{(‑a, ‑ar)}[1] \dicSynonym{ferðamaður} \dicDirectTranslationCS{turista, turistka}
\dicEntry[túrkis] \dicTerm{túrkis} \dicIPA{{t\smash{\textsuperscript{h}}}{u}{\textsubring{r}}{\r{\textObardotlessj}}{\textsci}{s}} \dicPos{m}[4] \dicFlx{(‑s, ‑ar)}[14] \dicFieldCat{geol.} \dicDirectTranslationCS{tyrkys}
\dicEntry[Túrkmeni] \dicTerm{Túrkmen|i} \dicIPA{{t\smash{\textsuperscript{h}}}{u}{\textsubring{r}}{\r{g}}{m}{\textepsilon}{n}{\textsci}} \dicPos{m}[1] \dicFlx{(‑a, ‑ar)}[1] \dicDirectTranslationCS{Turkmen(ka)}
\dicEntry[Túrkmenistan] \dicTerm{Túrkmenistan} \dicIPA{{t\smash{\textsuperscript{h}}}{u}{\textsubring{r}}{\r{g}}{m}{\textepsilon}{n}{\textsci}{s}{\textsubring{d}}{a}{\textsubring{n}}} \dicPos{n}[2] \dicFlx{(‑\,/\addthin ‑s)}[35] \dicFieldCat{geog.} \dicDirectTranslationCS{Turkmenistán}
\dicEntry[túrkmenska] \dicTerm{túrkmensk|a} \dicIPA{{t\smash{\textsuperscript{h}}}{u}{\textsubring{r}}{\r{g}}{m}{\textepsilon}{n}{s}{\r{g}}{a}} \dicPos{f}[1] \dicFlx{(‑u)}[5] \dicDirectTranslationCS{turkmenština}
\dicEntry[túrkmenskur] \dicTerm{túrkmenskur} \dicIPA{{t\smash{\textsuperscript{h}}}{u}{\textsubring{r}}{\r{g}}{m}{\textepsilon}{n}{s}{\r{g}}{\textscy}{\textsubring{r}}} \dicPos{adj}[1]\dicFlx{}[-6] \dicDirectTranslationCS{turkmenský}
\dicEntry[tússpenni] \dicTerm{túss··penn|i} \dicIPA{{t\smash{\textsuperscript{h}}}{u}{s}{\textlengthmark}{p\smash{\textsuperscript{h}}}{\textepsilon}{n}{\textsci}} \dicPos{m}[1] \dicFlx{(‑a, ‑ar)}[1] \dicDirectTranslationCS{fix, popisovač}
\dicEntry[tútna] \dicTerm{tútn|a} \dicIPA{{t\smash{\textsuperscript{h}}}{u}{h}{\textsubring{d}}{n}{a}} \dicPos{v}[1] \dicFlx{(‑aði)}[1] \dicPhraseIS{tútna út} {\textbf{a.}} \dicDirectTranslationCS{(na)bobtnat, (na)růst};  {\textbf{b.}} \dicDirectTranslationCS{(z)tloustnout}
\dicEntry[tútta] \dicTerm{tútt|a} \dicIPA{{t\smash{\textsuperscript{h}}}{u}{h}{\textsubring{d}}{a}} \dicPos{f}[1] \dicFlx{(‑u, ‑ur)}[7] \textbf{1.} \dicDirectTranslationCS{dudlík, cumel}  \textbf{2.} \dicLangCat{hovor.} \dicDirectTranslationCS{cecík}
\dicEntry[Túvalú] \dicTerm{Túvalú} \dicIPA{{t\smash{\textsuperscript{h}}}{u}{\textlengthmark}{v}{a}{l}{u}} \dicPos{n}[4] \dicFlx{indecl}[2] \dicFieldCat{geog.} \dicDirectTranslationCS{Tuvalu}
\dicEntry[Túvalúi] \dicTerm{Túvalú|i} \dicIPA{{t\smash{\textsuperscript{h}}}{u}{\textlengthmark}{v}{a}{l}{u}{\textsci}} \dicPos{m}[1] \dicFlx{(‑a, ‑ar)}[1] \dicDirectTranslationCS{Tuvalan(ka)}
\dicEntry[túvalúskur] \dicTerm{túvalúskur} \dicIPA{{t\smash{\textsuperscript{h}}}{u}{\textlengthmark}{v}{a}{l}{u}{s}{\r{g}}{\textscy}{\textsubring{r}}} \dicPos{adj}[1]\dicFlx{}[-6] \dicDirectTranslationCS{tuvalský}
\dicEntry[tveggja] \dicTerm{tveggja} \dicIPA{{t\smash{\textsuperscript{h}}}{v}{\textepsilon}{\r{\textObardotlessj}}{\textlengthmark}{a}} \dicPos{num} \dicFlx{gen} \dicLink{tveir}
\dicEntry[tveggjahæða] \dicTerm{tveggja··hæða} \dicIPA{{t\smash{\textsuperscript{h}}}{v}{\textepsilon}{\r{\textObardotlessj}}{\textlengthmark}{a}{h}{a}{i}{ð}{a}} \dicPos{adj}[13] \dicFlx{indecl}[1] \dicDirectTranslationCS{dvoupatrový, dvoupodlažní} \dicExampleIS{tveggjahæða hús} \dicExampleCS{dvoupatrový dům}
\dicEntry[tveir] \dicTerm{tveir} \dicsymFrequent\  \dicIPA{{t\smash{\textsuperscript{h}}}{v}{ei}{\textlengthmark}{\textsubring{r}}} \dicPos{num}[2] \dicFlx{pl} \dicFlx{m} \dicDirectTranslationCS{dva, dvě} \dicExampleIS{strekkja gaddavír milli tveggja staura} \dicExampleCS{natáhnout ostnatý drát mezi dvěma sloupy};  \dicIdiom{tveir}{ \dicPhraseIS{leika tveimur skjöldum}} \dicLangCat{přen.} \dicDirectTranslationCS{hrát dvojí hru}; { \dicPhraseIS{taka e‑m tveim höndum}} \dicDirectTranslationCS{přijmout (koho) s~otevřenou náručí};  \dicPhraseIS{(gera e‑ð) með tvær hendur tómar} \dicLangCat{přen.} \dicDirectTranslationCS{(udělat (co)) s~prázdnýma rukama}
\dicEntry[tvennd] \dicTerm{tvennd} \dicIPA{{t\smash{\textsuperscript{h}}}{v}{\textepsilon}{n}{\textsubring{d}}} \dicPos{f}[7] \dicFlx{(‑ar, ‑ir)}[1] \dicDirectTranslationCS{pár, dvojice} \dicExampleIS{fimm tvenndir e‑s} \dicExampleCS{pět párů (čeho)}
\dicEntry[tvenndarkeppni] \dicTerm{tvenndar··keppn|i} \dicIPA{{t\smash{\textsuperscript{h}}}{v}{\textepsilon}{n}{\textsubring{d}}{a}{\textsubring{r}}{c\smash{\textsuperscript{h}}}{\textepsilon}{h}{\textsubring{b}}{n}{\textsci}} \dicPos{f}[3] \dicFlx{(‑i, ‑ir)}[1] \dicDirectTranslationCS{(smíšená) čtyřhra, soutěž\,/\addthin závod (smíšených) dvojic}
\dicEntry[tvennskonar] \dicTerm{tvenns··konar}\dicTerm{, tvenns konar} \dicIPA{{t\smash{\textsuperscript{h}}}\-{v}\-{\textepsilon}\-{n}\-{s}\-{k\smash{\textsuperscript{h}}}\-{\textopeno}\-{n}\-{a}\-{\textsubring{r}}\-} \dicPos{adj}[13] \dicFlx{indecl}[1] \dicDirectTranslationCS{dvojí, (jsoucí) dvojího druhu} \dicExampleIS{tvennskonar brauð} \dicExampleCS{dva druhy chleba}
\dicEntry[tvennur] \dicTerm{tvennur} \dicIPA{{t\smash{\textsuperscript{h}}}{v}{\textepsilon}{n}{\textlengthmark}{\textscy}{\textsubring{r}}} \dicPos{adj}[1]\dicFlx{}[-7] \dicDirectTranslationCS{dvojitý, dva páry} \dicExampleIS{tvennir sokkar} \dicExampleCS{dvoje ponožky};  \dicPhraseIS{í tvennt} \dicFlx{adv} \dicDirectTranslationCS{ve dví, na dvě části, na půl}
\dicEntry[tvfn.] \dicTerm{tvfn.} \dicPos{zkr} \dicPhraseIS{tilvísunarfornafn} \dicFieldCat{jaz.} \dicDirectTranslationCS{vztažné zájmeno}
\dicEntry[tvinna] \dicTerm{tvinn|a} \dicIPA{{t\smash{\textsuperscript{h}}}{v}{\textsci}{n}{\textlengthmark}{a}} \dicPos{v}[1] \dicFlx{(‑aði)}[1] \dicFlx{acc} \textbf{1.} \dicSynonym*{gera tvöfalt} \dicDirectTranslationCS{zdvojit, zdvojovat} \dicExampleIS{tvinna band} \dicExampleCS{zdvojit provázek}  \textbf{2.} \dicSynonym{víxla} \dicDirectTranslationCS{(z)křížit, sloučit, slučovat} \dicExampleIS{Hesturinn er tvinnaður.} \dicExampleCS{Kůň je křížený.};  \dicIdiom{tvinnast}{ \dicPhraseIS{tvinnast}} \dicFlx{refl} \dicDirectTranslationCS{zkřížit se (osudy ap.)}
\dicEntry[tvinnakefli] \dicTerm{tvinna··kefli} \dicIPA{{t\smash{\textsuperscript{h}}}{v}{\textsci}{n}{\textlengthmark}{a}{c\smash{\textsuperscript{h}}}{\textepsilon}{\textsubring{b}}{l}{\textsci}} \dicPos{n}[2] \dicFlx{(‑s, ‑)}[14] \dicDirectTranslationCS{cívka\,/\addthin špulka (s~nití)}
\dicEntry[tvinni] \dicTerm{tvinn|i} \dicIPA{{t\smash{\textsuperscript{h}}}{v}{\textsci}{n}{\textlengthmark}{\textsci}} \dicPos{m}[1] \dicFlx{(‑a, ‑ar)}[1] \dicSynonym{þráður} \dicDirectTranslationCS{nit} \dicExampleIS{nál og tvinni} \dicExampleCS{jehla a~nit}
\dicEntry[tvinntala] \dicTerm{tvinn··|tala} \dicIPA{{t\smash{\textsuperscript{h}}}{v}{\textsci}{n}{t\smash{\textsuperscript{h}}}{a}{l}{a}} \dicPos{f}[1] \dicFlx{(‑tölu, ‑tölur)}[14] \dicFieldCat{mat.} \dicDirectTranslationCS{komplexní číslo}
\dicEntry[tvistur] \dicTerm{tvist|ur} \dicIPA{{t\smash{\textsuperscript{h}}}{v}{\textsci}{s}{\textsubring{d}}{\textscy}{\textsubring{r}}} \dicPos{m}[6] \dicFlx{(‑s, ‑ar)}[4] \textbf{1.} \dicDirectTranslationCS{dvojka} \dicIndirectTranslationCS{(v~kartách)} \dicExampleIS{tígul\-tvistur} \dicExampleCS{kárová dvojka}  \textbf{2.} \dicDirectTranslationCS{dvojka (autobus číslo 2 ap.)}  \textbf{3.} \dicDirectTranslationCS{bavlněný odpad};  \dicIdiom{tvistur}{ \dicPhraseIS{e‑að fer á tvist og bast}} \dicLangCat{přen.} \dicDirectTranslationCS{(co) leží vzhůru nohama, (co) je zpřeházené}
\dicEntry[tvisvar] \dicTerm{tvisvar} \dicsymFrequent\  \dicIPA{{t\smash{\textsuperscript{h}}}{v}{\textsci}{\textlengthmark}{s}{v}{a}{\textsubring{r}}} \dicPos{adv} \dicSynonym{tvívegis} \dicDirectTranslationCS{dvakrát} \dicExampleIS{tvisvar í viku} \dicExampleCS{dvakrát týdně}
\dicEntry[tví] \dicTerm{tví-} \dicIPA{{t\smash{\textsuperscript{h}}}{v}{i}{\textlengthmark}} \dicPos{predp} \dicDirectTranslationCS{dvou-, dvoj-, bi-} \dicExampleIS{tvíhendur} \dicExampleCS{dvouruký}
\dicEntry[tvíbaka] \dicTerm{tví··|baka} \dicIPA{{t\smash{\textsuperscript{h}}}{v}{i}{\textlengthmark}{\textsubring{b}}{a}{\r{g}}{a}} \dicPos{f}[1] \dicFlx{(‑böku, ‑bökur)}[20] \dicFieldCat{kulin.} \dicDirectTranslationCS{suchar, křupavý chléb}
\dicEntry[tvíbendur] \dicTerm{tví··bendur} \dicIPA{{t\smash{\textsuperscript{h}}}{v}{i}{\textlengthmark}{\textsubring{b}}{\textepsilon}{n}{\textsubring{d}}{\textscy}{\textsubring{r}}} \dicPos{adj}[2]\dicFlx{}[-17] \dicLink{tvíbentur}
\dicEntry[tvíbentur] \dicTerm{tví··bentur}\dicTerm{, tvíbendur} \dicIPA{{t\smash{\textsuperscript{h}}}\-{v}\-{i}\-{\textlengthmark}\-{\textsubring{b}}\-{\textepsilon}\-{\textsubring{n}}\-{\textsubring{d}}\-{\textscy}\-{\textsubring{r}}\-} \dicPos{adj}[1]\dicFlx{}[-13] \textbf{1.} \dicSynonym{efablandinn} \dicDirectTranslationCS{nerozhodný, váhavý} \dicExampleIS{vera tvíbentur í afstöðu sinni} \dicExampleCS{nerozhodný ve svém postoji}  \textbf{2.} \dicSynonym{tvíræður} \dicDirectTranslationCS{dvojznačný, dvojaký} \dicExampleIS{tvíbent svar} \dicExampleCS{dvojznačná odpověď}
\dicEntry[tvíblöðungur] \dicTerm{tví··blöð·ung|ur} \dicIPA{{t\smash{\textsuperscript{h}}}{v}{i}{\textlengthmark}{\textsubring{b}}{l}{\oe}{ð}{u}{\ng}{\r{g}}{\textscy}{\textsubring{r}}} \dicPos{m}[6] \dicFlx{(‑s, ‑ar)}[8] \dicDirectTranslationCS{dvojlist, folio}
\begin{xtolerant}{}{1pt}
\dicEntry[tvíbreiður] \dicTerm{tví··breiður}\addthinS\dicIPA{{t\smash{\textsuperscript{h}}}{v}{i}{\textlengthmark}{\textsubring{b}}{r}{ei}{ð}{\textscy}{\textsubring{r}}}\addthinS\dicPos{adj}[2]\dicFlx{}[-10]\addthinS\textbf{1.}\addthinS\dicDirectTranslationCS{dvouproudový (most ap.)} \dicAntonym{einbreiður}  \textbf{2.} \dicDirectTranslationCS{dvojitý (postel ap.)};  \dicPhraseIS{tvíbreitt rúm} \dicDirectTranslationCS{manželská postel}
\end{xtolerant}
\dicEntry[tvíbroddur] \dicTerm{tví··brodd|ur} \dicIPA{{t\smash{\textsuperscript{h}}}{v}{i}{\textlengthmark}{\textsubring{b}}{r}{\textopeno}{\textsubring{d}}{\textscy}{\textsubring{r}}} \dicPos{m}[6] \dicFlx{(‑s, ‑ar)}[4] \dicFieldCat{jaz.} \dicDirectTranslationCS{stříška, vokáň, cirkumflex} \dicIndirectTranslationCS{(diakritické znaménko)}
\dicEntry[tvíburi] \dicTerm{tví··bur|i} \dicsymFrequent\  \dicIPA{{t\smash{\textsuperscript{h}}}{v}{i}{\textlengthmark}{\textsubring{b}}{\textscy}{r}{\textsci}} \dicPos{m}[1] \dicFlx{(‑a, ‑ar)}[1] \textbf{1.} \dicDirectTranslationCS{dvojče} \dicExampleIS{Tvíburarnir voru eins klæddir.} \dicExampleCS{Dvojčata byla stejně oblečená.}  \textbf{2.} \dicPhraseIS{tvíburar} \dicFlx{pl} \dicDirectTranslationCS{Blíženci} \dicIndirectTranslationCS{(znamení zvěrokruhu)} \textit{(l.~{\textLA{Gemini}})}
\dicEntry[tvíbytna] \dicTerm{tví··bytn|a} \dicIPA{{t\smash{\textsuperscript{h}}}{v}{i}{\textlengthmark}{\textsubring{b}}{\textsci}{h}{\textsubring{d}}{n}{a}} \dicPos{f}[1] \dicFlx{(‑u, ‑ur)}[7] \dicFieldCat{nám.} \dicDirectTranslationCS{katamarán}
\dicEntry[tvíbýli] \dicTerm{tví··býli} \dicIPA{{t\smash{\textsuperscript{h}}}{v}{i}{\textlengthmark}{\textsubring{b}}{i}{l}{\textsci}} \dicPos{n}[2] \dicFlx{(‑s, ‑)}[14] \textbf{1.} \dicSynonym{tvíbýlishús} \dicDirectTranslationCS{dvougenerační dům}  \textbf{2.} \dicSynonym*{tvíbýlisjörð} \dicDirectTranslationCS{dvougenerační usedlost\,/\addthin hospodářství}
\dicEntry[tvíbýlishús] \dicTerm{tví·býlis··hús} \dicIPA{{t\smash{\textsuperscript{h}}}{v}{i}{\textlengthmark}{\textsubring{b}}{i}{l}{\textsci}{s}{h}{u}{s}} \dicPos{n}[2] \dicFlx{(‑s, ‑)}[5] \dicDirectTranslationCS{dvougenerační dům}
\dicEntry[tvídálka] \dicTerm{tví··dálka} \dicIPA{{t\smash{\textsuperscript{h}}}{v}{i}{\textlengthmark}{\textsubring{d}}{au}{\textsubring{l}}{\r{g}}{a}} \dicPos{adj}[13] \dicFlx{indecl}[1] \dicDirectTranslationCS{dvousloupcový}
\dicEntry[tvíeggja] \dicTerm{tví··eggja} \dicIPA{{t\smash{\textsuperscript{h}}}{v}{i}{\textlengthmark}{\textepsilon}{\r{\textObardotlessj}}{a}} \dicPos{adj}[13] \dicFlx{indecl}[1] \dicFieldCat{med.} \dicDirectTranslationCS{dvojvaječný} \dicExampleIS{tvíeggja tvíburar} \dicExampleCS{dvojvaječná dvojčata} \dicAntonym{eineggja}
\dicEntry[tvíeggjaður] \dicTerm{tví··eggj·|aður} \dicIPA{{t\smash{\textsuperscript{h}}}{v}{i}{\textlengthmark}{\textepsilon}{\r{\textObardotlessj}}{a}{ð}{\textscy}{\textsubring{r}}} \dicPos{adj}[3] \dicFlx{(f ‑uð)}[4] \textbf{1.} \dicDirectTranslationCS{dvousečný, dvojsečný, dvoubřitý, dvojbřitý} \dicExampleIS{tvíeggjað sverð} \dicExampleCS{dvousečný meč}  \textbf{2.} \dicLangCat{přen.} \dicDirectTranslationCS{dvousečný, dvojsečný, sporný};  \dicPhraseIS{e‑að er tvíeggjað} \dicLangCat{přen.} \dicDirectTranslationCS{(co) je dvousečné}
\dicEntry[tvífari] \dicTerm{tví··far|i} \dicIPA{{t\smash{\textsuperscript{h}}}{v}{i}{\textlengthmark}{f}{a}{r}{\textsci}} \dicPos{m}[1] \dicFlx{(‑a, ‑ar)}[8] \dicDirectTranslationCS{dvojník, dvojnice} \dicExampleIS{Hann er tvífari frænda síns.} \dicExampleCS{On je dvojníkem svého strýčka.}
\dicEntry[tvífættur] \dicTerm{tví··fættur} \dicIPA{{t\smash{\textsuperscript{h}}}{v}{i}{\textlengthmark}{f}{a}{i}{h}{\textsubring{d}}{\textscy}{\textsubring{r}}} \dicPos{adj}[1]\dicFlx{}[-13] \dicDirectTranslationCS{dvounohý}
\begin{xtolerant}{}{1pt}
\dicEntry[tvígangur] \dicTerm{tví··gang|ur} \dicIPA{{t\smash{\textsuperscript{h}}}{v}{i}{\textlengthmark}{\r{g}}{au}{\ng}{\r{g}}{\textscy}{\textsubring{r}}} \dicPos{m}[6] \dicFlx{(‑s)}[3] \dicPhraseIS{í tvígang} \dicFlx{adv} \dicSynonym{tvis\-var} \dicDirectTranslationCS{(na) dvakrát} \dicExampleIS{reyna þetta í tvígang} \dicExampleCS{zkusit to dvakrát}
\end{xtolerant}
\dicEntry[tvígildur] \dicTerm{tví··gildur} \dicIPA{{t\smash{\textsuperscript{h}}}{v}{i}{\textlengthmark}{\r{\textObardotlessj}}{\textsci}{l}{\textsubring{d}}{\textscy}{\textsubring{r}}} \dicPos{adj}[2]\dicFlx{}[-17] \dicFieldCat{chem.} \dicDirectTranslationCS{dvojmocný, dvoumocný}
\dicEntry[tvíhjól] \dicTerm{tví··hjól} \dicIPA{{t\smash{\textsuperscript{h}}}{v}{i}{\textlengthmark}{\c{c}}{ou}{\textsubring{l}}} \dicPos{n}[2] \dicFlx{(‑s, ‑)}[5] \dicDirectTranslationCS{dvoukolka, dvojkolka}
\dicEntry[tvíhjóla] \dicTerm{tví··hjóla} \dicIPA{{t\smash{\textsuperscript{h}}}{v}{i}{\textlengthmark}{\c{c}}{ou}{l}{a}} \dicPos{adj}[13] \dicFlx{indecl}[1] \dicDirectTranslationCS{dvoukolový}
\dicEntry[tvíhliða] \dicTerm{tví··hliða} \dicIPA{{t\smash{\textsuperscript{h}}}{v}{i}{\textlengthmark}{\textsubring{l}}{\textsci}{ð}{a}} \dicPos{adj}[13] \dicFlx{indecl}[1] \textbf{1.} \dicSynonym*{með tvær hliðar} \dicDirectTranslationCS{dvoustranný, dvojstranný} \dicIndirectTranslationCS{(který má dvě strany)}  \textbf{2.} \dicDirectTranslationCS{dvoustranný, dvojstranný, bilaterální} \dicIndirectTranslationCS{(uskutečňující se mezi dvěma stranami)} \dicExampleIS{tvíhliða samningur} \dicExampleCS{bilaterální smlouva}
\dicEntry[tvíhljóð] \dicTerm{tví··hljóð} \dicIPA{{t\smash{\textsuperscript{h}}}{v}{i}{\textlengthmark}{\textsubring{l}}{j}{ou}{\texttheta}} \dicPos{n}[2] \dicFlx{(‑s, ‑)}[5] \dicFieldCat{jaz.} \dicDirectTranslationCS{dvojhláska, diftong}
\dicEntry[tvíhnepptur] \dicTerm{tví··hnepptur} \dicIPA{{t\smash{\textsuperscript{h}}}{v}{i}{\textlengthmark}{\textsubring{n}}{\textepsilon}{f}{\textsubring{d}}{\textscy}{\textsubring{r}}} \dicPos{adj}[1]\dicFlx{}[-13] \dicDirectTranslationCS{dvouřadový (kabát ap.)} \dicExampleIS{tvíhnepptur jakki} \dicExampleCS{dvouřadový kabát}
\dicEntry[tvíhyggja] \dicTerm{tví··hyggj|a} \dicIPA{{t\smash{\textsuperscript{h}}}{v}{i}{\textlengthmark}{h}{\textsci}{\r{\textObardotlessj}}{a}} \dicPos{f}[1] \dicFlx{(‑u)}[5] \dicFieldCat{filos.} \dicDirectTranslationCS{dualismus}
\dicEntry[tvíhyrndur] \dicTerm{tví··hyrndur} \dicIPA{{t\smash{\textsuperscript{h}}}{v}{i}{\textlengthmark}{h}{\textsci}{r}{n}{\textsubring{d}}{\textscy}{\textsubring{r}}} \dicPos{adj}[2]\dicFlx{}[-17] \dicDirectTranslationCS{dvourohý} \dicExampleIS{tvíhyrnt dýr} \dicExampleCS{dvourohé zvíře}
\dicEntry[tvíhöfði] \dicTerm{tví··höfð|i} \dicIPA{{t\smash{\textsuperscript{h}}}{v}{i}{\textlengthmark}{h}{\oe}{v}{ð}{\textsci}} \dicPos{m}[1] \dicFlx{(‑a, ‑ar)}[1] \dicFieldCat{anat.} \dicDirectTranslationCS{biceps, dvojhlavý sval}
\dicEntry[tvíkvæni] \dicTerm{tví··kvæni} \dicIPA{{t\smash{\textsuperscript{h}}}{v}{i}{\textlengthmark}{k\smash{\textsuperscript{h}}}{v}{a}{i}{n}{\textsci}} \dicPos{n}[2] \dicFlx{(‑s)}[20] \dicDirectTranslationCS{bigamie}
\dicEntry[tvíkynja] \dicTerm{tví··kynja} \dicIPA{{t\smash{\textsuperscript{h}}}{v}{i}{\textlengthmark}{c\smash{\textsuperscript{h}}}{\textsci}{n}{j}{a}} \dicPos{adj}[13] \dicFlx{indecl}[1] \dicFieldCat{biol.} \dicDirectTranslationCS{oboupohlavní, obojpohlavní} \dicAntonym{einkynja}
\dicEntry[tvílemba] \dicTerm{tví··lemb|a} \dicIPA{{t\smash{\textsuperscript{h}}}{v}{i}{\textlengthmark}{l}{\textepsilon}{m}{\textsubring{b}}{a}} \dicPos{f}[1] \dicFlx{(‑u, ‑ur)}[13] \dicDirectTranslationCS{ovce mající dvě jehňata} \dicAntonym{einlemba}
\dicEntry[tvílembdur] \dicTerm{tví··lembdur} \dicIPA{{t\smash{\textsuperscript{h}}}{v}{i}{\textlengthmark}{l}{\textepsilon}{m}{\textsubring{d}}{\textscy}{\textsubring{r}}} \dicPos{adj}[2]\dicFlx{}[-17] \dicDirectTranslationCS{mající dvě jehňata} \dicIndirectTranslationCS{(o~ovci)}
\dicEntry[tvíliður] \dicTerm{tví··lið|ur} \dicIPA{{t\smash{\textsuperscript{h}}}{v}{i}{\textlengthmark}{l}{\textsci}{ð}{\textscy}{\textsubring{r}}} \dicPos{m}[10] \dicFlx{(‑ar\,/\addthin ‑s, ‑ir)}[33] \dicPhraseIS{réttur tvíliður} \dicFieldCat{lit.} \dicDirectTranslationCS{trochej};  \dicPhraseIS{öfugur tvíliður} \dicFieldCat{lit.} \dicDirectTranslationCS{jamb, jambický verš}
\dicEntry[tvímála] \dicTerm{tví··mála} \dicIPA{{t\smash{\textsuperscript{h}}}{v}{i}{\textlengthmark}{m}{au}{l}{a}} \dicPos{adj}[13] \dicFlx{indecl}[1] \dicDirectTranslationCS{dvojjazyčný, dvoujazyčný} \dicExampleIS{tvímála orðabækur} \dicExampleCS{dvojjazyčné slovníky}
\dicEntry[tvímenna] \dicTerm{tví··menn|a} \dicIPA{{t\smash{\textsuperscript{h}}}{v}{i}{\textlengthmark}{m}{\textepsilon}{n}{a}} \dicPos{v}[2] \dicFlx{(‑ti, ‑t)}[77] \dicDirectTranslationCS{jet na koni\,/\addthin kole ve dvou}
\dicEntry[tvímenningur] \dicTerm{tví··menn·ing|ur} \dicIPA{{t\smash{\textsuperscript{h}}}{v}{i}{\textlengthmark}{m}{\textepsilon}{n}{i}{\ng}{\r{g}}{\textscy}{\textsubring{r}}} \dicPos{m}[6] \dicFlx{(‑s, ‑ar)}[8] \textbf{1.} \dicSynonym{skyldmenni} \dicDirectTranslationCS{příbuzný\,/\addthin příbuzná z~prvního kolena}  \textbf{2.} \dicPhraseIS{tvímenningar} \dicFlx{pl} \dicDirectTranslationCS{pár\,/\addthin dvojice (lidí)}
\dicEntry[tvímælalaus] \dicTerm{tví·mæla··laus} \dicIPA{{t\smash{\textsuperscript{h}}}{v}{i}{\textlengthmark}{m}{a}{i}{l}{a}{l}{\oe i}{s}} \dicPos{adj}[5]\dicFlx{}[-1] \dicSynonym{greinilegur} \dicDirectTranslationCS{nepochybný, nesporný}
\dicEntry[tvímælalaust] \dicTerm{tví·mæla··laust} \dicsymFrequent\  \dicIPA{{t\smash{\textsuperscript{h}}}{v}{i}{\textlengthmark}{m}{a}{i}{l}{a}{l}{\oe i}{s}{\textsubring{d}}} \dicPos{adv} \dicDirectTranslationCS{nepochybně, nesporně, bezpochyby} \dicExampleIS{Hún lifir tvímælalaust þægilegu lífi.} \dicExampleCS{Žije nepochybně pohodlným životem.}
\dicEntry[tvímæli] \dicTerm{tví··mæli} \dicIPA{{t\smash{\textsuperscript{h}}}{v}{i}{\textlengthmark}{m}{a}{i}{l}{\textsci}} \dicPos{n}[2] \dicFlx{(‑s, ‑)}[14] \dicSynonym{efi} \dicDirectTranslationCS{spornost, pochybnost};  \dicPhraseIS{e‑að orkar tvímælis} \dicDirectTranslationCS{(co) budí pochybnosti};  \dicPhraseIS{taka af öll tvímæli um e‑ð} \dicDirectTranslationCS{rozptýlit všechny pochybnosti o~(čem)}
\dicEntry[tvínóna] \dicTerm{tví··nón|a} \dicIPA{{t\smash{\textsuperscript{h}}}{v}{i}{\textlengthmark}{n}{ou}{n}{a}} \dicPos{v}[1] \dicFlx{(‑aði)}[1] \dicDirectTranslationCS{rozpakovat se, rozmýšlet se} \dicExampleIS{vera ekki að tvínóna við e‑u} \dicExampleCS{nerozpakovat se nad (čím)}
\dicEntry[tvípunktur] \dicTerm{tví··punkt|ur} \dicIPA{{t\smash{\textsuperscript{h}}}{v}{i}{\textlengthmark}{p\smash{\textsuperscript{h}}}{u}{\textsubring{n}}{\textsubring{d}}{\textscy}{\textsubring{r}}} \dicPos{m}[6] \dicFlx{(‑s, ‑ar)}[4] \textbf{1.} \dicFieldCat{jaz.} \dicDirectTranslationCS{dvojtečka} \dicIndirectTranslationCS{(interpunkční znaménko)}  \textbf{2.} \dicFieldCat{jaz.} \dicDirectTranslationCS{tréma} \dicIndirectTranslationCS{(diakritické znaménko)}
\dicEntry[tvíraddað] \dicTerm{tví··raddað} \dicIPA{{t\smash{\textsuperscript{h}}}{v}{i}{\textlengthmark}{r}{a}{\textsubring{d}}{a}{\texttheta}} \dicPos{adv} \dicDirectTranslationCS{dvojhlasně, dvouhlasně}
\dicEntry[tvíraddaður] \dicTerm{tví··|radd·aður} \dicIPA{{t\smash{\textsuperscript{h}}}{v}{i}{\textlengthmark}{r}{a}{\textsubring{d}}{a}{ð}{\textscy}{\textsubring{r}}} \dicPos{adj}[3] \dicFlx{(f ‑rödduð)}[2] \dicDirectTranslationCS{dvojhlasný, dvouhlasný, dvojhlasý} \dicExampleIS{tvíraddaður söngur} \dicExampleCS{dvojhlasná píseň}
\dicEntry[tvírit] \dicTerm{tví··rit} \dicIPA{{t\smash{\textsuperscript{h}}}{v}{i}{\textlengthmark}{r}{\textsci}{\textsubring{d}}} \dicPos{n}[2] \dicFlx{(‑s, ‑)}[5] \dicDirectTranslationCS{dvojitá kopie (originál a~kopie)} \dicExampleIS{taka skjalið í tvíriti} \dicExampleCS{vzít dokument ve dvojité kopii}
\dicEntry[tvírím] \dicTerm{tví··rím} \dicIPA{{t\smash{\textsuperscript{h}}}{v}{i}{\textlengthmark}{r}{i}{\textsubring{m}}} \dicPos{n}[2] \dicFlx{(‑s, ‑)}[5] \dicFieldCat{lit.} \dicSynonym{kvenrím} \dicIndirectTranslationCS{koncový rým, ve kterém se rýmují obě slabiky dvouslabičných slov (např. maður : staður)}
\dicEntry[tvíræðni] \dicTerm{tví··ræðn|i} \dicIPA{{t\smash{\textsuperscript{h}}}{v}{i}{\textlengthmark}{r}{a}{i}{ð}{n}{\textsci}} \dicPos{f}[3] \dicFlx{(‑i)}[3] \dicDirectTranslationCS{dvojznačnost, dvouznačnost, dvojsmyslnost}
\dicEntry[tvíræður] \dicTerm{tví··ræður} \dicIPA{{t\smash{\textsuperscript{h}}}{v}{i}{\textlengthmark}{r}{a}{i}{ð}{\textscy}{\textsubring{r}}} \dicPos{adj}[2]\dicFlx{}[-6] \dicDirectTranslationCS{dvojznačný, dvouznačný, dvojsmyslný} \dicExampleIS{tvírætt svar} \dicExampleCS{dvojznačná odpověď}
\dicEntry[tvísaga] \dicTerm{tví··saga} \dicIPA{{t\smash{\textsuperscript{h}}}{v}{i}{\textlengthmark}{s}{a}{\textbabygamma}{a}} \dicPos{adj}[13] \dicFlx{indecl}[1] \dicDirectTranslationCS{protiřečící si, rozporuplný, rozporný};  \dicPhraseIS{verða tvísaga} \dicDirectTranslationCS{protiřečit si}
\dicEntry[tvískinnungur] \dicTerm{tví··skinn·ung|ur} \dicIPA{{t\smash{\textsuperscript{h}}}{v}{i}{\textlengthmark}{s}{\r{\textObardotlessj}}{\textsci}{n}{u}{\ng}{\r{g}}{\textscy}{\textsubring{r}}} \dicPos{m}[6] \dicFlx{(‑s, ‑ar)}[8] \dicDirectTranslationCS{rozporuplnost, rozpornost, protikladnost}
\dicEntry[tvísköttun] \dicTerm{tví··skött|un} \dicIPA{{t\smash{\textsuperscript{h}}}{v}{i}{\textlengthmark}{s}{\r{g}}{\oe}{h}{\textsubring{d}}{\textscy}{\textsubring{n}}} \dicPos{f}[7] \dicFlx{(‑unar)}[12] \dicFieldCat{ekon.} \dicDirectTranslationCS{dvojité zdanění}
\dicEntry[tvíslá] \dicTerm{tví··slá} \dicIPA{{t\smash{\textsuperscript{h}}}{v}{i}{\textlengthmark}{s}{\textsubring{d}}{l}{au}} \dicPos{f}[4] \dicFlx{(‑r, ‑r)}[18] \dicFieldCat{sport.} \dicDirectTranslationCS{bradla}
\dicEntry[tvísmella] \dicTerm{tví··smell|a} \dicIPA{{t\smash{\textsuperscript{h}}}{v}{i}{\textlengthmark}{s}{m}{\textepsilon}{\textsubring{d}}{l}{a}} \dicPos{v}[2] \dicFlx{(‑ti, ‑t)}[78] \dicFieldCat{poč.} \dicDirectTranslationCS{dvakrát kliknout\,/\addthin ťuknout};  \dicPhraseIS{tvísmella á e‑ð} \dicDirectTranslationCS{dvakrát kliknout na (co)}
\dicEntry[tvísteig] \dicTerm{tví··steig} \dicIPA{{t\smash{\textsuperscript{h}}}{v}{i}{\textlengthmark}{s}{\textsubring{d}}{ei}{x}} \dicPos{v} \dicFlx{ind pf sg 1 pers} \dicLink{tvístíga}
\dicEntry[tvístigi] \dicTerm{tví··stigi} \dicIPA{{t\smash{\textsuperscript{h}}}{v}{i}{\textlengthmark}{s}{\textsubring{d}}{i}{j}{\textsci}} \dicPos{v} \dicFlx{con pf sg 1 pers} \dicLink{tvístíga}
\dicEntry[tvístigið] \dicTerm{tví··stigið} \dicIPA{{t\smash{\textsuperscript{h}}}{v}{i}{\textlengthmark}{s}{\textsubring{d}}{i}{j}{\textsci}{\texttheta}} \dicPos{v} \dicFlx{supin} \dicLink{tvístíga}
\dicEntry[tvístigum] \dicTerm{tví··stigum} \dicIPA{{t\smash{\textsuperscript{h}}}{v}{i}{\textlengthmark}{s}{\textsubring{d}}{\textsci}{\textbabygamma}{\textscy}{\textsubring{m}}} \dicPos{v} \dicFlx{ind pf pl 1 pers} \dicLink{tvístíga}
\dicEntry[tvístíg] \dicTerm{tví··stíg} \dicIPA{{t\smash{\textsuperscript{h}}}{v}{i}{\textlengthmark}{s}{\textsubring{d}}{i}{x}} \dicPos{v} \dicFlx{ind praes sg 1 pers} \dicLink{tvístíga}
\dicEntry[tvístíga] \dicTerm{tví··|stíga} \dicIPA{{t\smash{\textsuperscript{h}}}{v}{i}{\textlengthmark}{s}{\textsubring{d}}{i}{\textbabygamma}{a}} \dicPos{v}[6] \dicFlx{(‑stíg, ‑steig, ‑stigum, ‑stigi, ‑stigið)}[79] \textbf{1.} \dicSynonym*{stjákla} \dicDirectTranslationCS{přešlapovat na místě}  \textbf{2.} \dicSynonym{hika} \dicDirectTranslationCS{rozmýšlet se, rozpakovat se}
\dicEntry[tvístígandi] \dicTerm{tví··stíg·andi} \dicIPA{{t\smash{\textsuperscript{h}}}{v}{i}{\textlengthmark}{s}{\textsubring{d}}{i}{\textbabygamma}{a}{n}{\textsubring{d}}{\textsci}} \dicPos{adj}[13] \dicFlx{indecl}[1] \dicDirectTranslationCS{váhavý, nerozhodný, kolísavý}
\dicEntry[tvístra] \dicTerm{tvístr|a} \dicIPA{{t\smash{\textsuperscript{h}}}{v}{i}{s}{\textsubring{d}}{r}{a}} \dicPos{v}[1] \dicFlx{(‑aði)}[1] \dicFlx{dat} \dicSynonym{dreifa} \dicDirectTranslationCS{rozptýlit, rozptylovat, rozehnat, rozhánět, rozprášit, rozprašovat} \dicExampleIS{tvístra hópnum} \dicExampleCS{rozehnat dav};  \dicIdiom{tvístrast}{ \dicPhraseIS{tvístrast}} \dicFlx{refl} \dicSynonym*{dreifast} \dicDirectTranslationCS{rozptýlit se, rozprášit se, rozprchnout se} \dicExampleIS{Þegar skólanum lauk tvístraðist nemendahópurinn um allt land.} \dicExampleCS{Když skončila škola, rozprchli se žáci po celé zemi.}
\dicEntry[tvístrun] \dicTerm{tvístr|un} \dicIPA{{t\smash{\textsuperscript{h}}}{v}{i}{s}{\textsubring{d}}{r}{\textscy}{\textsubring{n}}} \dicPos{f}[7] \dicFlx{(‑unar)}[9] \textbf{1.} \dicSynonym{dreifing} \dicDirectTranslationCS{rozptýlení, rozprášení}  \textbf{2.} \dicFieldCat{fyz.} \dicDirectTranslationCS{rozptyl}
\dicEntry[tvísýna] \dicTerm{tví··sýn|a} \dicIPA{{t\smash{\textsuperscript{h}}}{v}{i}{\textlengthmark}{s}{i}{n}{a}} \dicPos{f}[1] \dicFlx{(‑u, ‑ur)}[7] \dicSynonym{óvissa} \dicDirectTranslationCS{nejistota, pochybnost};  \dicPhraseIS{tefla í tvísýnu} \dicDirectTranslationCS{riskovat, hazardovat, zahrávat si};  \dicPhraseIS{tefla e‑u í tvísýnu} \dicDirectTranslationCS{hazardovat s~(čím), zahrávat si s~(čím)}
\dicEntry[tvísýnn] \dicTerm{tví··sýnn} \dicIPA{{t\smash{\textsuperscript{h}}}{v}{i}{\textlengthmark}{s}{i}{\textsubring{d}}{\textsubring{n}}} \dicPos{adj}[7]\dicFlx{}[-1] \dicSynonym{óviss} \dicDirectTranslationCS{nejistý, nejednoznač\-ný, nepředvídatelný} \dicExampleIS{tvísýn keppni} \dicExampleCS{utkání s~nejistým výsledkem}
\dicEntry[tvísýringur] \dicTerm{tví··sýr·ing|ur} \dicIPA{{t\smash{\textsuperscript{h}}}{v}{i}{\textlengthmark}{s}{i}{r}{i}{\ng}{\r{g}}{\textscy}{\textsubring{r}}} \dicPos{m}[6] \dicFlx{(‑s, ‑ar)}[8] \dicFieldCat{chem.} \dicDirectTranslationCS{dioxid}
\dicEntry[tvísögli] \dicTerm{tví··sögl|i} \dicIPA{{t\smash{\textsuperscript{h}}}{v}{i}{\textlengthmark}{s}{\oe}{\r{g}}{l}{\textsci}} \dicPos{f}[3] \dicFlx{(‑i)}[3] \dicDirectTranslationCS{protiřečení, protimluv, kontradikce}
\dicEntry[tvísöngur] \dicTerm{tví··söng|ur} \dicIPA{{t\smash{\textsuperscript{h}}}{v}{i}{\textlengthmark}{s}{\oe i}{\ng}{\r{g}}{\textscy}{\textsubring{r}}} \dicPos{m}[8] \dicFlx{(‑s, ‑var)}[2] \dicFieldCat{hud.} \dicDirectTranslationCS{duo, dueto, duet}
\dicEntry[tvítak] \dicTerm{tví··|tak} \dicIPA{{t\smash{\textsuperscript{h}}}{v}{i}{\textlengthmark}{t\smash{\textsuperscript{h}}}{a}{\r{g}}} \dicPos{n}[2] \dicFlx{(‑taks, ‑tök)}[8] \dicDirectTranslationCS{duplikát}
\dicEntry[tvítaka] \dicTerm{tví··|taka} \dicIPA{{t\smash{\textsuperscript{h}}}{v}{i}{\textlengthmark}{t\smash{\textsuperscript{h}}}{a}{\r{g}}{a}} \dicPos{v}[6] \dicFlx{(‑tek, ‑tók, ‑tókum, ‑tæki, ‑tekið)}[128] \dicFlx{acc} \dicSynonym{endurtaka} \dicDirectTranslationCS{dvakrát (z)opakovat}
\dicEntry[tvítala] \dicTerm{tví··|tala} \dicIPA{{t\smash{\textsuperscript{h}}}{v}{i}{\textlengthmark}{t\smash{\textsuperscript{h}}}{a}{l}{a}} \dicPos{f}[1] \dicFlx{(‑tölu)}[2] \dicFieldCat{jaz.} \dicDirectTranslationCS{dvojné číslo, duál}
\dicEntry[tvítek] \dicTerm{tví··tek} \dicIPA{{t\smash{\textsuperscript{h}}}{v}{i}{\textlengthmark}{t\smash{\textsuperscript{h}}}{\textepsilon}{\r{g}}} \dicPos{v} \dicFlx{ind praes sg 1 pers} \dicLink{tvítaka}
\dicEntry[tvítekið] \dicTerm{tví··tekið} \dicIPA{{t\smash{\textsuperscript{h}}}{v}{i}{\textlengthmark}{t\smash{\textsuperscript{h}}}{\textepsilon}{\r{\textObardotlessj}}{\textsci}{\texttheta}} \dicPos{v} \dicFlx{supin} \dicLink{tvítaka}
\dicEntry[tvítekning] \dicTerm{tví··tek·ning} \dicIPA{{t\smash{\textsuperscript{h}}}{v}{i}{\textlengthmark}{t\smash{\textsuperscript{h}}}{\textepsilon}{h}{\r{g}}{n}{i}{\ng}{\r{g}}} \dicPos{f}[4] \dicFlx{(‑ar, ‑ar)}[5] \dicDirectTranslationCS{(dvojité) (z)opakování}
\dicEntry[tvítók] \dicTerm{tví··tók} \dicIPA{{t\smash{\textsuperscript{h}}}{v}{i}{\textlengthmark}{t\smash{\textsuperscript{h}}}{ou}{\r{g}}} \dicPos{v} \dicFlx{ind pf sg 1 pers} \dicLink{tvítaka}
\dicEntry[tvítugfaldur] \dicTerm{tví·tug··|faldur} \dicIPA{{t\smash{\textsuperscript{h}}}{v}{i}{\textlengthmark}{t\smash{\textsuperscript{h}}}{\textscy}{x}{f}{a}{l}{\textsubring{d}}{\textscy}{\textsubring{r}}} \dicPos{adj}[2] \dicFlx{(f ‑föld)}[16] \textbf{1.} \dicDirectTranslationCS{dvacetinásobný}  \textbf{2.} \dicDirectTranslationCS{dvacetidílný}
\dicEntry[tvítugsaldur] \dicTerm{tví·tugs··ald|ur} \dicIPA{{t\smash{\textsuperscript{h}}}{v}{i}{\textlengthmark}{t\smash{\textsuperscript{h}}}{\textscy}{x}{s}{a}{l}{\textsubring{d}}{\textscy}{\textsubring{r}}} \dicPos{m}[5] \dicFlx{(‑urs, ‑rar)}[3] \dicDirectTranslationCS{věk pozdního dospívání} \dicIndirectTranslationCS{(věk mezi 16. a~20. rokem)}
\dicEntry[tvítugur] \dicTerm{tví··tugur} \dicsymFrequent\  \dicIPA{{t\smash{\textsuperscript{h}}}{v}{i}{\textlengthmark}{t\smash{\textsuperscript{h}}}{\textscy}{\textbabygamma}{\textscy}{\textsubring{r}}} \dicPos{adj}[1]\dicFlx{}[-6] \textbf{1.} \dicSynonym*{tuttugu ára} \dicDirectTranslationCS{dvacetiletý} \dicExampleIS{Hún er tvítug.} \dicExampleCS{Je jí dvacet let.}  \textbf{2.} \dicSynonym*{tuttugu faðma langur} \dicDirectTranslationCS{měřící dvacet sáhů}
\dicEntry[tvítyngdur] \dicTerm{tví··tyngdur} \dicIPA{{t\smash{\textsuperscript{h}}}{v}{i}{\textlengthmark}{t\smash{\textsuperscript{h}}}{\textsci}{\ng}{\textsubring{d}}{\textscy}{\textsubring{r}}} \dicPos{adj}[2]\dicFlx{}[-17] \dicDirectTranslationCS{dvojjazyčný, bilingvní, bilingvální}
\dicEntry[tvítæki] \dicTerm{tví··tæki} \dicIPA{{t\smash{\textsuperscript{h}}}{v}{i}{\textlengthmark}{t\smash{\textsuperscript{h}}}{a}{i}{\r{\textObardotlessj}}{\textsci}} \dicPos{v} \dicFlx{con pf sg 1 pers} \dicLink{tvítaka}
\dicEntry[tvíund] \dicTerm{tví··und} \dicIPA{{t\smash{\textsuperscript{h}}}{v}{i}{\textlengthmark}{j}{\textscy}{n}{\textsubring{d}}} \dicPos{f}[7] \dicFlx{(‑ar, ‑ir)}[1] \dicFieldCat{hud.} \dicDirectTranslationCS{sekunda} \dicIndirectTranslationCS{(hudební interval mezi prvním a~druhým tónem)}
\dicEntry[tvíundakerfi] \dicTerm{tví·unda··kerfi} \dicIPA{{t\smash{\textsuperscript{h}}}{v}{i}{\textlengthmark}{j}{\textscy}{n}{\textsubring{d}}{a}{c\smash{\textsuperscript{h}}}{\textepsilon}{r}{v}{\textsci}} \dicPos{n}[2] \dicFlx{(‑s)}[20] \dicFieldCat{mat.\,/\addthin poč.} \dicDirectTranslationCS{dvojková\,/\addthin binární soustava}
\dicEntry[tvívarahljóð] \dicTerm{tví·vara··hljóð} \dicIPA{{t\smash{\textsuperscript{h}}}{v}{i}{\textlengthmark}{v}{a}{r}{a}{\textsubring{l}}{j}{ou}{\texttheta}} \dicPos{n}[2] \dicFlx{(‑s, ‑)}[5] \dicFieldCat{jaz.} \dicDirectTranslationCS{obouretná souhláska, bilabiála}
\dicEntry[tvívaramæltur] \dicTerm{tví·vara··mæltur} \dicIPA{{t\smash{\textsuperscript{h}}}{v}{i}{\textlengthmark}{v}{a}{r}{a}{m}{a}{i}{\textsubring{l}}{\textsubring{d}}{\textscy}{\textsubring{r}}} \dicPos{adj}[1]\dicFlx{}[-13] \dicFieldCat{jaz.} \dicDirectTranslationCS{obouret\-ný, bilabiální}
\dicEntry[tvívegis] \dicTerm{tví··vegis} \dicIPA{{t\smash{\textsuperscript{h}}}{v}{i}{\textlengthmark}{v}{ei}{\textsci}{s}} \dicPos{adv} \dicDirectTranslationCS{dvakrát}
\dicEntry[tvíverknaður] \dicTerm{tví··verk·nað|ur} \dicIPA{{t\smash{\textsuperscript{h}}}{v}{i}{\textlengthmark}{v}{\textepsilon}{\textsubring{r}}{\r{g}}{n}{a}{ð}{\textscy}{\textsubring{r}}} \dicPos{m}[10] \dicFlx{(‑ar, ‑ir)}[8] \dicDirectTranslationCS{(zbytečná) dublovaná práce}
\dicEntry[tvívetni] \dicTerm{tví··vetni} \dicIPA{{t\smash{\textsuperscript{h}}}{v}{i}{\textlengthmark}{v}{\textepsilon}{h}{\textsubring{d}}{n}{\textsci}} \dicPos{n}[2] \dicFlx{(‑s)}[20] \dicFieldCat{chem.} \dicDirectTranslationCS{deuterium, těžký vodík}
\dicEntry[tvívíður] \dicTerm{tví··víður} \dicIPA{{t\smash{\textsuperscript{h}}}{v}{i}{\textlengthmark}{v}{i}{ð}{\textscy}{\textsubring{r}}} \dicPos{adj}[2]\dicFlx{}[-10] \dicDirectTranslationCS{dvojrozměrný} \dicAntonym{þrívíður}
\dicEntry[tvíþættur] \dicTerm{tví··þættur} \dicIPA{{t\smash{\textsuperscript{h}}}{v}{i}{\textlengthmark}{\texttheta}{a}{i}{h}{\textsubring{d}}{\textscy}{\textsubring{r}}} \dicPos{adj}[1]\dicFlx{}[-13] \textbf{1.} \dicDirectTranslationCS{dvoudílný} \dicExampleIS{tvíþætt verkefni} \dicExampleCS{dvoudílná úloha}  \textbf{2.} \dicDirectTranslationCS{dvojí} \dicExampleIS{hafa tvíþætta merkingu} \dicExampleCS{mít dvojí význam}
\dicEntry[tvíær] \dicTerm{tví··ær} \dicIPA{{t\smash{\textsuperscript{h}}}{v}{i}{\textlengthmark}{a}{i}{\textsubring{r}}} \dicPos{adj}[5] \dicFlx{(f ‑)}[10] \dicDirectTranslationCS{dvouletý} \dicIndirectTranslationCS{(o~rostlině)} \dicExampleIS{tvíær jurt} \dicExampleCS{dvouletá rostlina}
\dicEntry[tvt.] \dicTerm{tvt.} \dicPos{zkr} \dicPhraseIS{tvítala} \dicFieldCat{jaz.} \dicDirectTranslationCS{duál, dvojné číslo}
\dicEntry[tvær] \dicTerm{tvær} \dicIPA{{t\smash{\textsuperscript{h}}}{v}{a}{i}{\textlengthmark}{\textsubring{r}}} \dicPos{num} \dicFlx{pl} \dicFlx{f nom\,/\addthin acc} \dicLink{tveir}
\dicEntry[tvö] \dicTerm{tvö} \dicIPA{{t\smash{\textsuperscript{h}}}{v}{\oe}{\textlengthmark}} \dicPos{num} \dicFlx{pl} \dicFlx{n nom\,/\addthin acc} \dicLink{tveir}
\dicEntry[tvöfalda] \dicTerm{tvö··fald|a} \dicIPA{{t\smash{\textsuperscript{h}}}{v}{\oe}{\textlengthmark}{f}{a}{l}{\textsubring{d}}{a}} \dicPos{v}[1] \dicFlx{(‑aði)}[13] \dicFlx{acc} \dicDirectTranslationCS{zdvojnásobit, zdvojnásobovat} \dicExampleIS{tvöfalda stærðina} \dicExampleCS{zdvojnásobit velikost};  \dicIdiom{tvöfaldast}{ \dicPhraseIS{tvöfaldast}} \dicFlx{refl} \dicDirectTranslationCS{zdvojnásobit se, zdvojnásobovat se} \dicExampleIS{Skuldirnar hafa tvöfaldast á skömmum tíma.} \dicExampleCS{Dluhy se zdvojnásobily během krátkého období.}
\dicEntry[tvöfaldur] \dicTerm{tvö··|faldur} \dicIPA{{t\smash{\textsuperscript{h}}}{v}{\oe}{\textlengthmark}{f}{a}{l}{\textsubring{d}}{\textscy}{\textsubring{r}}} \dicPos{adj}[2] \dicFlx{(f ‑föld)}[16] \textbf{1.} \dicDirectTranslationCS{dvojitý, zdvojený} \dicExampleIS{tvöfaldir skattar} \dicExampleCS{dvojitá daň};  \dicPhraseIS{tvöfaldur (vodka)} \dicDirectTranslationCS{dvojitá (vodka)};  \dicPhraseIS{tvöfalt vaff} \dicFieldCat{jaz.} \dicDirectTranslationCS{dvojité \clqq v\crqq } \dicIndirectTranslationCS{(název písmene \clqq w\crqq  v~islandském pravopise)};  \dicPhraseIS{tvöfalt i\,/\addthin í} \dicFieldCat{jaz.} \dicDirectTranslationCS{krátký\,/\addthin dlouhý ypsilon} \dicIndirectTranslationCS{(název písmene \clqq y\crqq \,/\addthin \clqq ý\crqq  v~islandském pravopise)}  \textbf{2.} \dicDirectTranslationCS{dvojnásobný}
\dicEntry[tvöfalt] \dicTerm{tvö··falt} \dicIPA{{t\smash{\textsuperscript{h}}}{v}{\oe}{\textlengthmark}{f}{a}{\textsubring{l}}{\textsubring{d}}} \dicPos{adv} \dicDirectTranslationCS{dvojnásobně}
\dicEntry[tvöfeldni] \dicTerm{tvö··feldn|i} \dicIPA{{t\smash{\textsuperscript{h}}}{v}{\oe}{\textlengthmark}{f}{\textepsilon}{l}{\textsubring{d}}{n}{\textsci}} \dicPos{f}[3] \dicFlx{(‑i)}[3] \dicDirectTranslationCS{dvojakost, falešnost, licoměrnost}
\dicEntry[tvöföldun] \dicTerm{tvö··föld|un} \dicIPA{{t\smash{\textsuperscript{h}}}{v}{\oe}{\textlengthmark}{f}{\oe}{l}{\textsubring{d}}{\textscy}{\textsubring{n}}} \dicPos{f}[7] \dicFlx{(‑unar)}[12] \dicDirectTranslationCS{zdvojnásobení}
\dicEntry[tvöþúsundkall] \dicTerm{tvö·þúsund··kall} \dicIPA{{t\smash{\textsuperscript{h}}}{v}{\oe}{\textlengthmark}{\texttheta}{u}{s}{\textscy}{n}{\textsubring{d}}{k\smash{\textsuperscript{h}}}{a}{\textsubring{d}}{\textsubring{l}}} \dicPos{m}[4] \dicFlx{(‑s, ‑ar)}[4] \textbf{1.} \dicLangCat{hovor.} \dicDirectTranslationCS{dva tisíce (korun)} \dicIndirectTranslationCS{(částka)}  \textbf{2.} \dicLangCat{hovor.} \dicDirectTranslationCS{dvoutisícikoruna, dvoutisícovka} \dicIndirectTranslationCS{(bankovka)}
\dicEntry[tygg] \dicTerm{tygg} \dicIPA{{t\smash{\textsuperscript{h}}}{\textsci}{\r{g}}{\textlengthmark}} \dicPos{v} \dicFlx{ind praes sg 1 pers} \dicLink{tyggja}
\dicEntry[tyggði] \dicTerm{tyggði} \dicIPA{{t\smash{\textsuperscript{h}}}{\textsci}{\textbabygamma}{ð}{\textsci}} \dicPos{v} \dicFlx{con pf sg 1 pers} \dicLink{tyggja}
\dicEntry[tyggigúmí] \dicTerm{tyggi··gúmí} \dicIPA{{t\smash{\textsuperscript{h}}}{\textsci}{\r{\textObardotlessj}}{\textlengthmark}{\textsci}{\r{g}}{u}{m}{i}} \dicPos{n}[2] \dicFlx{(‑s)}[2] \dicLink{tyggigúmmí}
\dicEntry[tyggigúmmí] \dicTerm{tyggi··gúmmí}\dicTerm{, tyggigúmí} \dicIPA{{t\smash{\textsuperscript{h}}}\-{\textsci}\-{\r{\textObardotlessj}}\-{\textlengthmark}\-{\textsci}\-{\r{g}}\-{u}\-{m}\-{i}\-} \dicPos{n}[2] \dicFlx{(‑s)}[2] \dicDirectTranslationCS{žvýkačka, žvýkací guma}
\dicEntry[tyggja] \dicTerm{tyggja} \dicsymFrequent\  \dicIPA{{t\smash{\textsuperscript{h}}}{\textsci}{\r{\textObardotlessj}}{\textlengthmark}{a}} \dicPos{v}[4] \dicFlx{(tygg, tuggði, tuggðum, tyggði, tuggið)}[51] \dicFlx{acc} \dicDirectTranslationCS{(roz)žvýkat, (roz)kousat} \dicExampleIS{tyggja fæðuna vandlega} \dicExampleCS{důkladně rozžvýkat stravu};  \dicIdiom{tyggja}[í]{ \dicPhraseIS{tyggja e‑ð í e‑n}} \dicLangCat{přen.} \dicDirectTranslationCS{namluvit (komu co), dostat (komu co) do hlavy};  \dicIdiom{tyggja}[upp]{ \dicPhraseIS{tyggja e‑ð upp eftir e‑m}} \dicLangCat{přen.} \dicDirectTranslationCS{(nekriticky) opakovat (co) po (kom), přežvykovat (co) po (kom)}
\dicEntry[tyggjó] \dicTerm{tyggjó} \dicIPA{{t\smash{\textsuperscript{h}}}{\textsci}{\r{\textObardotlessj}}{\textlengthmark}{ou}} \dicPos{n}[2] \dicFlx{(‑s, ‑)}[5] \dicLangCat{hovor.} \dicSynonym{tyggigúmmí} \dicDirectTranslationCS{žvýkačka, žvejka}
\dicEntry[tygi] \dicTerm{tygi} \dicIPA{{t\smash{\textsuperscript{h}}}{i}{j}{\textlengthmark}{\textsci}} \dicPos{n}[2] \dicFlx{(‑s, ‑)}[16] \dicSynonym{ástarsamband} \dicDirectTranslationCS{(milostný) vztah, poměr};  \dicPhraseIS{vera í tygjum við e‑n} \dicDirectTranslationCS{mít s~(kým) poměr}
\dicEntry[tygja] \dicTerm{tygj|a} \dicIPA{{t\smash{\textsuperscript{h}}}{i}{j}{\textlengthmark}{a}} \dicPos{v}[1] \dicFlx{(‑aði)}[1] \dicFlx{acc} \dicPhraseIS{tygja sig} \dicDirectTranslationCS{nachystat\,/\addthin připravit se k~odchodu}
\dicEntry[tylft] \dicTerm{tylft} \dicIPA{{t\smash{\textsuperscript{h}}}{\textsci}{\textsubring{l}}{\textsubring{d}}} \dicPos{f}[7] \dicFlx{(‑ar, ‑ir)}[1] \dicDirectTranslationCS{tucet}
\dicEntry[tylla] \dicTerm{tyll|a} \dicsymFrequent\  \dicIPA{{t\smash{\textsuperscript{h}}}{\textsci}{\textsubring{d}}{l}{a}} \dicPos{v}[2] \dicFlx{(‑ti, ‑t)}[63] \dicFlx{dat} \textbf{1.} \dicPhraseIS{tylla sér} \dicDirectTranslationCS{usednout (na chvíli), usadit se (na chvíli), sednout si (na chvíli)}  \textbf{2.} \dicDirectTranslationCS{umístit, umísťovat, usadit, usazovat}  \textbf{3.} \dicSynonym*{festa lauslega} \dicDirectTranslationCS{přichytit, přichycovat, (volně) upevnit\,/\addthin upevňovat} \dicExampleIS{tylla bót á flíkina} \dicExampleCS{přichytit záplatu na oblečení}  \textbf{4.} \dicDirectTranslationCS{stoupnout (si), stoupat (si)} \dicExampleIS{tylla fætinum á e‑u} \dicExampleCS{stoupnout nohou na (co)};  \dicPhraseIS{tylla sér á tá} {\textbf{a.}} \dicDirectTranslationCS{postavit se na špičky, stát na špičkách};  {\textbf{b.}} \dicLangCat{přen.} \dicDirectTranslationCS{stavět se na špičky} \dicIndirectTranslationCS{(snažit se na sebe upozornit)}
\dicEntry[tylliástæða] \textls[15]{\dicTerm{tylli··á·stæð|a} \dicIPA{{t\smash{\textsuperscript{h}}}{\textsci}{\textsubring{d}}{l}{\textsci}{au}{s}{\textsubring{d}}{a}{i}{ð}{a}} \dicPos{f}[1] \dicFlx{(‑u, ‑ur)}[13] \dicDirectTranslationCS{záminka, předstíraný důvod}}
\dicEntry[tyllidagur] \dicTerm{tylli··dag|ur} \dicIPA{{t\smash{\textsuperscript{h}}}{\textsci}{\textsubring{d}}{l}{\textsci}{\textsubring{d}}{a}{\textbabygamma}{\textscy}{\textsubring{r}}} \dicPos{m}[6] \dicFlx{(‑s, ‑ar)}[62] \dicDirectTranslationCS{sváteční\,/\addthin slavnostní den};  \dicPhraseIS{á tyllidögum} \dicFlx{adv} \dicDirectTranslationCS{o~svátečních dnech, během svátečních dnů}
\dicEntry[typpi] \dicTerm{typpi} \dicIPA{{t\smash{\textsuperscript{h}}}{\textsci}{h}{\textsubring{b}}{\textsci}} \dicPos{n}[2] \dicFlx{(‑s, ‑)}[14] \textbf{1.} \dicSynonym{hnappur} \dicDirectTranslationCS{tlačítko, knoflík, čudlík}  \textbf{2.} \dicSynonym{getnaðarlimur} \dicDirectTranslationCS{šulín} \dicIndirectTranslationCS{(penis)};  \dicIdiom{typpi}{ \dicPhraseIS{það er uppi á e‑m typpið}} \dicFlx{impers} \dicLangCat{přen.} \dicDirectTranslationCS{(kdo) je ve svém živlu}
\dicEntry[tyrfa] \dicTerm{tyrf|a} \dicIPA{{t\smash{\textsuperscript{h}}}{\textsci}{r}{v}{a}} \dicPos{v}[2] \dicFlx{(‑ði, ‑t)}[105] \dicFlx{acc} \dicDirectTranslationCS{položit\,/\addthin pokládat (drnový) trávník} \dicExampleIS{tyrfa þakið} \dicExampleCS{položit drnový trávník na střechu}
\dicEntry[tyrfinn] \dicTerm{tyrfinn} \dicIPA{{t\smash{\textsuperscript{h}}}{\textsci}{r}{v}{\textsci}{\textsubring{n}}} \dicPos{adj}[6]\dicFlx{}[-2] \dicSynonym{torskilinn} \dicDirectTranslationCS{těžko srozumitelný, nesrozumitelný} \dicExampleIS{tyrfið orðalag} \dicExampleCS{nesrozumitelné vyjádření}
\dicEntry[Tyrki] \dicTerm{Tyrk|i} \dicIPA{{t\smash{\textsuperscript{h}}}{\textsci}{\textsubring{r}}{\r{\textObardotlessj}}{\textsci}} \dicPos{m}[3] \dicFlx{(‑ja, ‑ir)}[2] \dicDirectTranslationCS{Turek, Turkyně}
\begin{xtolerant}{}{10pt}
\dicEntry[tyrkjadúfa] \dicTerm{tyrkja··dúf|a} \dicIPA{{t\smash{\textsuperscript{h}}}{\textsci}{\textsubring{r}}{\r{\textObardotlessj}}{a}{\textsubring{d}}{u}{v}{a}} \dicPos{f}[1] \dicFlx{(‑u, ‑ur)}[13] \dicFieldCat{zool.} \dicDirectTranslationCS{hrdlička za\-hradní} \textit{(l.~{\textLA{Streptopelia decaocto}})}
\end{xtolerant}
\dicEntry[tyrkjasól] \dicTerm{tyrkja··sól} \dicIPA{{t\smash{\textsuperscript{h}}}{\textsci}{\textsubring{r}}{\r{\textObardotlessj}}{a}{s}{ou}{\textsubring{l}}} \dicPos{f}[7] \dicFlx{(‑ar, ‑ir)}[1] \dicFieldCat{bot.} \dicDirectTranslationCS{mák východní} \textit{(l.~{\textLA{Papaver orientale}})}  \dicsymPhoto\ 
\dicFigure{ds_image_tyrkjasol_0_1.jpg}{Tyrkjasól}{Tyrkjasól - James Steakley, CC BY-SA 3.0}
\dicEntry[Tyrkland] \dicTerm{Tyrk··land} \dicIPA{{t\smash{\textsuperscript{h}}}{\textsci}{\textsubring{r}}{\r{g}}{l}{a}{n}{\textsubring{d}}} \dicPos{n}[2] \dicFlx{(‑s)}[4] \dicFieldCat{geog.} \dicDirectTranslationCS{Turecko}
\dicEntry[tyrkneska] \dicTerm{tyrk··nesk|a} \dicIPA{{t\smash{\textsuperscript{h}}}{\textsci}{\textsubring{r}}{\r{g}}{n}{\textepsilon}{s}{\r{g}}{a}} \dicPos{f}[1] \dicFlx{(‑u)}[5] \dicDirectTranslationCS{turečtina}
\dicEntry[tyrkneskur] \dicTerm{tyrk··neskur} \dicIPA{{t\smash{\textsuperscript{h}}}{\textsci}{\textsubring{r}}{\r{g}}{n}{\textepsilon}{s}{\r{g}}{\textscy}{\textsubring{r}}} \dicPos{adj}[1]\dicFlx{}[-1] \dicDirectTranslationCS{turecký}
\dicEntry[týna] \dicTerm{týn|a} \dicsymFrequent\  \dicIPA{{t\smash{\textsuperscript{h}}}{i}{\textlengthmark}{n}{a}} \dicPos{v}[2] \dicFlx{(‑di, ‑t)}[141] \dicFlx{dat} \dicDirectTranslationCS{ztratit, (po)ztrácet, přijít\,/\addthin přicházet o} \dicExampleIS{týna bókinni} \dicExampleCS{ztratit knížku};  \dicPhraseIS{týna lífinu} \dicDirectTranslationCS{přijít o~život};  \dicIdiom{týna}[niður]{ \dicPhraseIS{týna e‑u niður}} \dicDirectTranslationCS{zapomenout (co)} \dicIndirectTranslationCS{(něco, co se člověk dříve naučil)};  \dicIdiom{týna}[sér]{ \dicPhraseIS{týna sjálfum sér}} \dicDirectTranslationCS{zabít se, spáchat sebevraždu};  \dicIdiom{týnast}{ \dicPhraseIS{týnast}} \dicFlx{refl} \dicDirectTranslationCS{ztratit se, ztrácet se, (z)mizet} \dicExampleIS{Skipið týndist í hafi.} \dicExampleCS{Loď se ztratila na moři.}
\dicEntry[týndur] \dicTerm{týndur} \dicIPA{{t\smash{\textsuperscript{h}}}{i}{n}{\textsubring{d}}{\textscy}{\textsubring{r}}} \dicPos{adj}[2]\dicFlx{}[-14] \dicDirectTranslationCS{ztracený, nezvěstný, pohřešovaný}
\dicEntry[týpa] \dicTerm{týp|a} \dicIPA{{t\smash{\textsuperscript{h}}}{i}{\textlengthmark}{\textsubring{b}}{a}} \dicPos{f}[1] \dicFlx{(‑u, ‑ur)}[19] \textbf{1.} \dicSynonym{tegund} \dicDirectTranslationCS{typ, model}  \textbf{2.} \dicSynonym{manngerð} \dicDirectTranslationCS{typ (člověka)}
\dicEntry[týpískur] \dicTerm{týpískur} \dicIPA{{t\smash{\textsuperscript{h}}}{i}{\textlengthmark}{\textsubring{b}}{i}{s}{\r{g}}{\textscy}{\textsubring{r}}} \dicPos{adj}[1]\dicFlx{}[-1] \dicLangCat{hovor.} \dicSynonym{dæmigerður} \dicDirectTranslationCS{typický}
\dicEntry[Týr] \dicTerm{Tý|r} \dicIPA{{t\smash{\textsuperscript{h}}}{i}{\textlengthmark}{\textsubring{r}}} \dicPos{m}[6] \dicFlx{(‑s)}[53] \dicFlx{prop} \dicFieldCat{myt.} \dicDirectTranslationCS{Týr} \dicIndirectTranslationCS{(bůh války)}
\dicEntry[týra] \dicTerm{týr|a\smash{\textsuperscript{1}}} \dicIPA{{t\smash{\textsuperscript{h}}}{i}{\textlengthmark}{r}{a}} \dicPos{f}[1] \dicFlx{(‑u, ‑ur)}[7] \dicLink{tíra\smash{\textsuperscript{1}}}
\dicEntry[týra] \dicTerm{týr|a\smash{\textsuperscript{2}}} \dicIPA{{t\smash{\textsuperscript{h}}}{i}{\textlengthmark}{r}{a}} \dicPos{v}[2] \dicFlx{(‑ði, ‑t)}[111] \dicLink{tíra\smash{\textsuperscript{2}}}
\dicEntry[týsfjóla] \dicTerm{týs··fjól|a} \dicIPA{{t\smash{\textsuperscript{h}}}{i}{\textlengthmark}{s}{f}{j}{ou}{l}{a}} \dicPos{f}[1] \dicFlx{(‑u, ‑ur)}[19] \dicFieldCat{bot.} \dicDirectTranslationCS{violka psí} \textit{(l.~{\textLA{Viola canina}})}  \dicsymPhoto\ 
\dicFigure{69690.jpg}{Týsfjóla}{Týsfjóla - Krejčík Stanislav, Biolib, Copyright/CC-BY}
\dicEntry[týtuber] \dicTerm{týtu··ber} \dicIPA{{t\smash{\textsuperscript{h}}}{i}{\textlengthmark}{\textsubring{d}}{\textscy}{\textsubring{b}}{\textepsilon}{\textsubring{r}}} \dicPos{n}[2] \dicFlx{(‑s, ‑)}[13] \dicSynonym{rauðber} \dicDirectTranslationCS{brusinka} \dicIndirectTranslationCS{(plod)}
\dicEntry[tæfa] \dicTerm{tæf|a} \dicIPA{{t\smash{\textsuperscript{h}}}{a}{i}{\textlengthmark}{v}{a}} \dicPos{f}[1] \dicFlx{(‑u, ‑ur)}[19] \textbf{1.} \dicSynonym{tófa} \dicDirectTranslationCS{liška} \dicIndirectTranslationCS{(samice)}  \textbf{2.} \dicLangCat{han.} \dicSynonym*{ómerkileg kona} \dicDirectTranslationCS{mrcha, potvora, čůza}
\dicEntry[tæki] \dicTerm{tæki\smash{\textsuperscript{1}}} \dicsymFrequent\  \dicIPA{{t\smash{\textsuperscript{h}}}{a}{i}{\textlengthmark}{\r{\textObardotlessj}}{\textsci}} \dicPos{n}[2] \dicFlx{(‑s, ‑)}[16] \textbf{1.} \dicSynonym{verkfæri} \dicDirectTranslationCS{nářadí, náčiní, nástroj}  \textbf{2.} \dicSynonym{vél\smash{\textsuperscript{1}}} \dicDirectTranslationCS{přístroj, zařízení, aparát} \dicExampleIS{Tækið virkar eðlilega.} \dicExampleCS{Zařízení funguje normálně.}
\dicEntry[tæki] \dicTerm{tæki\smash{\textsuperscript{2}}} \dicIPA{{t\smash{\textsuperscript{h}}}{a}{i}{\textlengthmark}{\r{\textObardotlessj}}{\textsci}} \dicPos{v} \dicFlx{con pf sg 1 pers} \dicLink{taka\smash{\textsuperscript{2}}}
\dicEntry[tækifæri] \dicTerm{tæki··færi} \dicsymFrequent\  \dicIPA{{t\smash{\textsuperscript{h}}}{a}{i}{\textlengthmark}{\r{\textObardotlessj}}{\textsci}{f}{a}{i}{r}{\textsci}} \dicPos{n}[2] \dicFlx{(‑s, ‑)}[14] \textbf{1.} \dicSynonym{möguleiki} \dicDirectTranslationCS{příležitost, možnost, šance};  \dicPhraseIS{grípa\,/\addthin nota tækifærið} \dicDirectTranslationCS{využít příležitosti};  \dicPhraseIS{við fyrsta tækifæri} \dicFlx{adv} \dicDirectTranslationCS{při první příležitosti};  \dicPhraseIS{þegar tækifæri gefst\,/\addthin býðst} \dicFlx{refl} \dicDirectTranslationCS{až se naskytne příležitost};  \dicPhraseIS{við tækifæri} \dicFlx{adv} \dicDirectTranslationCS{při příležitosti, příležitostně}  \textbf{2.} \dicSynonym{tilefni} \dicDirectTranslationCS{příležitost, vhodná okolnost};  \dicPhraseIS{við hátíðleg tækifæri} \dicFlx{adv} \dicDirectTranslationCS{při slavnostních příležitostech}
\dicEntry[tækifæriskjóll] \dicTerm{tæki·færis··kjól|l} \dicIPA{{t\smash{\textsuperscript{h}}}{a}{i}{\textlengthmark}{\r{\textObardotlessj}}{\textsci}{f}{a}{i}{r}{\textsci}{s}{c\smash{\textsuperscript{h}}}{ou}{\textsubring{d}}{\textsubring{l}}} \dicPos{m}[6] \dicFlx{(‑s, ‑ar)}[48] \dicDirectTranslationCS{těhotenské šaty, šaty pro těhotné ženy}
\dicEntry[tækifærissinni] \dicTerm{tæki·færis··sinn|i} \dicIPA{{t\smash{\textsuperscript{h}}}{a}{i}{\textlengthmark}{\r{\textObardotlessj}}{\textsci}{f}{a}{i}{r}{\textsci}{s}{\textsci}{n}{\textsci}} \dicPos{m}[1] \dicFlx{(‑a, ‑ar)}[1] \dicDirectTranslationCS{prospěchář(ka), oportunista, oportunistka}
\dicEntry[tækifærisstefna] \dicTerm{tæki·færis··stefn|a} \dicIPA{{t\smash{\textsuperscript{h}}}{a}{i}{\textlengthmark}{\r{\textObardotlessj}}{\textsci}{f}{a}{i}{r}{\textsci}{s}{\textsubring{d}}{\textepsilon}{\textsubring{b}}{n}{a}} \dicPos{f}[1] \dicFlx{(‑u)}[5] \dicSynonym{hentistefna} \dicDirectTranslationCS{prospěchářství, oportunismus}
\dicEntry[tækisfall] \dicTerm{tækis··|fall} \dicIPA{{t\smash{\textsuperscript{h}}}{a}{i}{\textlengthmark}{\r{\textObardotlessj}}{\textsci}{s}{f}{a}{\textsubring{d}}{\textsubring{l}}} \dicPos{n}[2] \dicFlx{(‑falls, ‑föll)}[8] \dicFieldCat{jaz.} \dicSynonym{verkfærisfall} \dicDirectTranslationCS{7. pád, instrumentál}
\dicEntry[tækjabúnaður] \dicTerm{tækja··bú·nað|ur} \dicIPA{{t\smash{\textsuperscript{h}}}{a}{i}{\textlengthmark}{\r{\textObardotlessj}}{a}{\textsubring{b}}{u}{n}{a}{ð}{\textscy}{\textsubring{r}}} \dicPos{m}[10] \dicFlx{(‑ar)}[9] \dicDirectTranslationCS{technické vybavení\,/\addthin zařízení}
\dicEntry[tækni] \dicTerm{tækn|i} \dicsymFrequent\  \dicIPA{{t\smash{\textsuperscript{h}}}{a}{i}{h}{\r{g}}{n}{\textsci}} \dicPos{f}[3] \dicFlx{(‑i)}[3] \textbf{1.} \dicSynonym*{véla- og verkkunnátta} \dicDirectTranslationCS{technologie} \dicExampleIS{skoða nýtingu nýrrar tækni} \dicExampleCS{zkoumat využití nové technologie}  \textbf{2.} \dicSynonym{aðferð} \dicDirectTranslationCS{technika} \dicIndirectTranslationCS{(způsob provádění činnosti)}
\dicEntry[tæknibrella] \dicTerm{tækni··brell|a} \dicIPA{{t\smash{\textsuperscript{h}}}{a}{i}{h}{\r{g}}{n}{\textsci}{\textsubring{b}}{r}{\textepsilon}{\textsubring{d}}{l}{a}} \dicPos{f}[1] \dicFlx{(‑u, ‑ur)}[13] \dicDirectTranslationCS{speciální\,/\addthin zvláštní efekt}
\dicEntry[tæknifrjóvgun] \dicTerm{tækni··frjóvg|un} \dicIPA{{t\smash{\textsuperscript{h}}}{a}{i}{h}{\r{g}}{n}{\textsci}{f}{r}{j}{ou}{v}{\textbabygamma}{\textscy}{\textsubring{n}}} \dicPos{f}[7] \dicFlx{(‑unar, ‑anir)}[8] \textbf{1.} \dicFieldCat{med.} \dicSynonym{glasafrjóvgun} \dicDirectTranslationCS{umělé oplodnění}  \textbf{2.} \dicSynonym{sæðing} \dicDirectTranslationCS{inseminace}
\dicEntry[tæknifræði] \dicTerm{tækni··fræð|i} \dicIPA{{t\smash{\textsuperscript{h}}}{a}{i}{h}{\r{g}}{n}{\textsci}{f}{r}{a}{i}{ð}{\textsci}} \dicPos{f}[3] \dicFlx{(‑i)}[3] \dicDirectTranslationCS{technika, technologie} \dicIndirectTranslationCS{(nauka)}
\dicEntry[tæknifræðingur] \dicTerm{tækni·fræð··ing|ur} \dicIPA{{t\smash{\textsuperscript{h}}}{a}{i}{h}{\r{g}}{n}{\textsci}{f}{r}{a}{i}{ð}{i}{\ng}{\r{g}}{\textscy}{\textsubring{r}}} \dicPos{m}[6] \dicFlx{(‑s, ‑ar)}[8] \dicDirectTranslationCS{technik, technička, technolog, technoložka}
\dicEntry[tækniháskóli] \dicTerm{tækni··há·skól|i} \dicIPA{{t\smash{\textsuperscript{h}}}{a}{i}{h}{\r{g}}{n}{\textsci}{h}{au}{s}{\r{g}}{ou}{l}{\textsci}} \dicPos{m}[1] \dicFlx{(‑a, ‑ar)}[1] \dicFieldCat{škol.} \dicDirectTranslationCS{polytechnika, technická univerzita}
\dicEntry[tæknilega] \dicTerm{tækni··lega} \dicIPA{{t\smash{\textsuperscript{h}}}{a}{i}{h}{\r{g}}{n}{\textsci}{l}{\textepsilon}{\textbabygamma}{a}} \dicPos{adv} \dicDirectTranslationCS{technicky}
\dicEntry[tæknilegur] \dicTerm{tækni··legur} \dicIPA{{t\smash{\textsuperscript{h}}}{a}{i}{h}{\r{g}}{n}{\textsci}{l}{\textepsilon}{\textbabygamma}{\textscy}{\textsubring{r}}} \dicPos{adj}[1]\dicFlx{}[-8] \textbf{1.} \dicSynonym*{sem varðar tækni} \dicDirectTranslationCS{technický}  \textbf{2.} \dicSynonym*{sem varðar formið} \dicDirectTranslationCS{formální, technický}
\dicEntry[tæknimaður] \dicTerm{tækni··|maður} \dicIPA{{t\smash{\textsuperscript{h}}}{a}{i}{h}{\r{g}}{n}{\textsci}{m}{a}{ð}{\textscy}{\textsubring{r}}} \dicPos{m}[13] \dicFlx{(‑manns, ‑menn)}[2] \dicDirectTranslationCS{technik, technička}
\begin{xtolerant}{}{4pt}
\dicEntry[tæknimenntaður] \dicTerm{tækni··mennt·|aður} \dicIPA{{t\smash{\textsuperscript{h}}}{a}{i}{h}{\r{g}}{n}{\textsci}{m}{\textepsilon}{\textsubring{n}}{\textsubring{d}}{a}{ð}{\textscy}{\textsubring{r}}} \dicPos{adj}[3] \dicFlx{(f ‑uð)}[3] \dicDirectTranslationCS{tech\-nicky vzdělaný}
\end{xtolerant}
\dicEntry[tæknimenntun] \dicTerm{tækni··mennt|un} \dicIPA{{t\smash{\textsuperscript{h}}}{a}{i}{h}{\r{g}}{n}{\textsci}{m}{\textepsilon}{\textsubring{n}}{\textsubring{d}}{\textscy}{\textsubring{n}}} \dicPos{f}[7] \dicFlx{(‑unar)}[9] \dicDirectTranslationCS{technické vzdělání}
\dicEntry[tækninám] \dicTerm{tækni··nám} \dicIPA{{t\smash{\textsuperscript{h}}}{a}{i}{h}{\r{g}}{n}{\textsci}{n}{au}{\textsubring{m}}} \dicPos{n}[2] \dicFlx{(‑s)}[2] \dicDirectTranslationCS{technické učení}
\dicEntry[tækninýjung] \dicTerm{tækni··nýj·ung} \dicIPA{{t\smash{\textsuperscript{h}}}{a}{i}{h}{\r{g}}{n}{\textsci}{n}{i}{j}{u}{\ng}{\r{g}}} \dicPos{f}[4] \dicFlx{(‑ar, ‑ar)}[1] \dicDirectTranslationCS{technická\,/\addthin technologická novinka}
\dicEntry[tækniskóli] \dicTerm{tækni··skól|i} \dicIPA{{t\smash{\textsuperscript{h}}}{a}{i}{h}{\r{g}}{n}{\textsci}{s}{\r{g}}{ou}{l}{\textsci}} \dicPos{m}[1] \dicFlx{(‑a, ‑ar)}[1] \dicFieldCat{škol.} \dicDirectTranslationCS{střední škola technická}
\dicEntry[tækniteiknari] \dicTerm{tækni··teikn·ar|i} \dicIPA{{t\smash{\textsuperscript{h}}}{a}{i}{h}{\r{g}}{n}{\textsci}{t\smash{\textsuperscript{h}}}{ei}{h}{\r{g}}{n}{a}{r}{\textsci}} \dicPos{m}[1] \dicFlx{(‑a, ‑ar)}[13] \dicDirectTranslationCS{kreslič(ka), projektant(ka)}
\dicEntry[tækniteiknun] \dicTerm{tækni··teikn|un} \dicIPA{{t\smash{\textsuperscript{h}}}{a}{i}{h}{\r{g}}{n}{\textsci}{t\smash{\textsuperscript{h}}}{ei}{h}{\r{g}}{n}{\textscy}{\textsubring{n}}} \dicPos{f}[7] \dicFlx{(‑unar)}[9] \dicDirectTranslationCS{technické kreslení}
\dicEntry[tæknivæddur] \dicTerm{tækni··væddur} \dicIPA{{t\smash{\textsuperscript{h}}}{a}{i}{h}{\r{g}}{n}{\textsci}{v}{a}{i}{\textsubring{d}}{\textscy}{\textsubring{r}}} \dicPos{adj}[2]\dicFlx{}[-21] \dicDirectTranslationCS{(jsoucí) s~vyspělou technologií}
\dicEntry[tæknivæða] \dicTerm{tækni··væ|ða} \dicIPA{{t\smash{\textsuperscript{h}}}{a}{i}{h}{\r{g}}{n}{\textsci}{v}{a}{i}{ð}{a}} \dicPos{v}[2] \dicFlx{(‑ddi, ‑tt)}[171] \dicFlx{acc} \dicDirectTranslationCS{zavádět nové technologie}
\dicEntry[tækniþróun] \dicTerm{tækni··þró|un} \dicIPA{{t\smash{\textsuperscript{h}}}{a}{i}{h}{\r{g}}{n}{\textsci}{\texttheta}{r}{ou}{\textscy}{\textsubring{n}}} \dicPos{f}[7] \dicFlx{(‑unar)}[9] \dicDirectTranslationCS{technický\,/\addthin technologický rozvoj}
\dicEntry[tækniöld] \dicTerm{tækni··|öld} \dicIPA{{t\smash{\textsuperscript{h}}}{a}{i}{h}{\r{g}}{n}{\textsci}{\oe}{l}{\textsubring{d}}} \dicPos{f}[7] \dicFlx{(‑aldar)}[19] \dicDirectTranslationCS{století\,/\addthin věk technologie}
\dicEntry[tækur] \dicTerm{tækur} \dicIPA{{t\smash{\textsuperscript{h}}}{a}{i}{\textlengthmark}{\r{g}}{\textscy}{\textsubring{r}}} \dicPos{adj}[1]\dicFlx{}[-1] \dicSynonym{hæfilegur} \dicDirectTranslationCS{vhodný, příhodný, přijatelný, použitelný}
\dicEntry[tæla] \dicTerm{tæl|a} \dicIPA{{t\smash{\textsuperscript{h}}}{a}{i}{\textlengthmark}{l}{a}} \dicPos{v}[2] \dicFlx{(‑di, ‑t)}[141] \dicFlx{acc} \dicSynonym{ginna} \dicDirectTranslationCS{svést, svádět, (z)lákat, (z)vábit} \dicExampleIS{tæla stúlku} \dicExampleCS{svést dívku}
\dicEntry[Tæland] \dicTerm{Tæ··land} \dicIPA{{t\smash{\textsuperscript{h}}}{a}{i}{\textlengthmark}{l}{a}{n}{\textsubring{d}}} \dicPos{n}[2] \dicFlx{(‑s)}[4] \dicLink{Taíland}
\dicEntry[tælandi] \dicTerm{tæl··andi} \dicIPA{{t\smash{\textsuperscript{h}}}{a}{i}{\textlengthmark}{l}{a}{n}{\textsubring{d}}{\textsci}} \dicPos{adj}[13] \dicFlx{indecl}[1] \dicDirectTranslationCS{svůdný, svádivý, lákavý, vábivý}
\dicEntry[Tælendingur] \dicTerm{Tæ·lend··ing|ur} \dicIPA{{t\smash{\textsuperscript{h}}}{a}{i}{\textlengthmark}{l}{\textepsilon}{n}{\textsubring{d}}{i}{\ng}{\r{g}}{\textscy}{\textsubring{r}}} \dicPos{m}[6] \dicFlx{(‑s, ‑ar)}[8] \dicLink{Taílendingur}
\dicEntry[tælenska] \dicTerm{tæ··lensk|a} \dicIPA{{t\smash{\textsuperscript{h}}}{a}{i}{\textlengthmark}{l}{\textepsilon}{n}{s}{\r{g}}{a}} \dicPos{f}[1] \dicFlx{(‑u)}[5] \dicLink{taílenska}
\dicEntry[tælenskur] \dicTerm{tæ··lenskur} \dicIPA{{t\smash{\textsuperscript{h}}}{a}{i}{\textlengthmark}{l}{\textepsilon}{n}{s}{\r{g}}{\textscy}{\textsubring{r}}} \dicPos{adj}[1]\dicFlx{}[-6] \dicLink{taílenskur}
\dicEntry[tæma] \dicTerm{tæm|a} \dicsymFrequent\  \dicIPA{{t\smash{\textsuperscript{h}}}{a}{i}{\textlengthmark}{m}{a}} \dicPos{v}[2] \dicFlx{(‑di, ‑t)}[144] \dicFlx{acc} \dicDirectTranslationCS{vyprázdnit, vyprazdňovat} \dicExampleIS{tæma ílátið} \dicExampleCS{vyprázdnit nádobu};  \dicIdiom{tæmast}{ \dicPhraseIS{tæmast}} \dicFlx{refl} \dicDirectTranslationCS{vyprázdnit se, vyprazdňovat se}; { \dicPhraseIS{e‑m tæmist arfur}} \dicFlx{refl impers} \dicDirectTranslationCS{(komu) připadá dědictví, (kdo) zdědil dědictví}
\dicEntry[tæmandi] \dicTerm{tæm··andi} \dicIPA{{t\smash{\textsuperscript{h}}}{a}{i}{\textlengthmark}{m}{a}{n}{\textsubring{d}}{\textsci}} \dicPos{adj}[13] \dicFlx{indecl}[1] \dicSynonym{endanlegur} \dicDirectTranslationCS{vyčerpávající, úplný (seznam ap.)} \dicIndirectTranslationCS{(nejčastěji se záporem)} \dicExampleIS{tæmandi listi} \dicExampleCS{vyčerpávající seznam}
\dicEntry[tæming] \dicTerm{tæm··ing} \dicIPA{{t\smash{\textsuperscript{h}}}{a}{i}{\textlengthmark}{m}{i}{\ng}{\r{g}}} \dicPos{f}[4] \dicFlx{(‑ar)}[7] \dicDirectTranslationCS{vyprázdnění, vyprazdňování}
\dicEntry[tæpa] \dicTerm{tæp|a} \dicIPA{{t\smash{\textsuperscript{h}}}{a}{i}{\textlengthmark}{\textsubring{b}}{a}} \dicPos{v}[2] \dicFlx{(‑ti, ‑t)}[73] \dicPhraseIS{tæpa á e‑u} \dicDirectTranslationCS{narážet na (co) (slovem), dělat slovní narážky na (co), připomenout (co), připomínat (co)}
\dicEntry[tæpitunga] \dicTerm{tæpi··tung|a} \dicIPA{{t\smash{\textsuperscript{h}}}{a}{i}{\textlengthmark}{\textsubring{b}}{\textsci}{t\smash{\textsuperscript{h}}}{u}{\ng}{\r{g}}{a}} \dicPos{f}[1] \dicFlx{(‑u, ‑ur)}[13] \dicSynonym{smámæli} \dicDirectTranslationCS{šišlání, žvatlání};  \dicPhraseIS{tala enga tæpitungu} \dicLangCat{přen.} \dicDirectTranslationCS{mluvit k~věci, mluvit přímo}
\dicEntry[tæplega] \dicTerm{tæp··lega} \dicsymFrequent\  \dicIPA{{t\smash{\textsuperscript{h}}}{a}{i}{\textlengthmark}{\textsubring{b}}{l}{\textepsilon}{\textbabygamma}{a}} \dicPos{adv} \textbf{1.} \dicSynonym*{knappt} \dicDirectTranslationCS{téměř, skoro} \dicExampleIS{tæplega fjögur prósent} \dicExampleCS{téměř čtyři procenta}  \textbf{2.} \dicSynonym{varla} \dicDirectTranslationCS{sotva, stěží}
\dicEntry[tæpt] \dicTerm{tæp|t} \dicsymFrequent\  \dicIPA{{t\smash{\textsuperscript{h}}}{a}{i}{f}{\textsubring{d}}} \dicPos{adv} \dicFlx{(comp ‑ar\,/\addthin ‑ara, sup ‑ast)} \dicDirectTranslationCS{těsně, o~vlásek};  \dicPhraseIS{standa tæpt} \dicLangCat{přen.} \dicDirectTranslationCS{viset na vlásku}
\dicEntry[tæpur] \dicTerm{tæpur} \dicsymFrequent\  \dicIPA{{t\smash{\textsuperscript{h}}}{a}{i}{\textlengthmark}{\textsubring{b}}{\textscy}{\textsubring{r}}} \dicPos{adj}[1]\dicFlx{}[-1] \textbf{1.} \dicSynonym*{aðeins minni en e‑ð} \dicDirectTranslationCS{necelý} \dicExampleIS{tæpar 100 krónur} \dicExampleCS{necelých 100 korun}  \textbf{2.} \dicSynonym*{ekki nógur} \dicDirectTranslationCS{těsný, nedostatečný, (jsoucí) na hraně};  \dicPhraseIS{vera tæpur með e‑ð} \dicDirectTranslationCS{mít jen těsné zásoby (čeho), být na hraně s~(čím)};  \dicPhraseIS{það var tæpt að (e‑að gerðist)} \dicDirectTranslationCS{málem (se (co) stalo), chybělo málo a~((co) by se stalo)};  \dicPhraseIS{það mátti ekki tæpara standa} \dicFlx{impers} \dicLangCat{přen.} \dicDirectTranslationCS{bylo to o~vlásek\,/\addthin fous}
\dicEntry[tær] \dicTerm{tær\smash{\textsuperscript{1}}} \dicIPA{{t\smash{\textsuperscript{h}}}{a}{i}{\textlengthmark}{\textsubring{r}}} \dicPos{f} \dicFlx{pl nom} \dicLink{tá}
\dicEntry[tær] \dicTerm{tær\smash{\textsuperscript{2}}} \dicIPA{{t\smash{\textsuperscript{h}}}{a}{i}{\textlengthmark}{\textsubring{r}}} \dicPos{adj}[5] \dicFlx{(f ‑)}[8] \dicDirectTranslationCS{průzračný, čirý} \dicExampleIS{tært vatn} \dicExampleCS{průzračná voda}
\dicEntry[tæra] \dicTerm{tær|a} \dicIPA{{t\smash{\textsuperscript{h}}}{a}{i}{\textlengthmark}{r}{a}} \dicPos{v}[2] \dicFlx{(‑ði, ‑t)}[99] \dicFlx{acc} \dicDirectTranslationCS{prožrat, prožírat, rozežrat, rozežírat (rzí ap.)} \dicExampleIS{Ryðið tærir málminn.} \dicExampleCS{Rez rozežírá kov.};  \dicIdiom{tærast}{ \dicPhraseIS{tærast}} \dicFlx{refl} \dicDirectTranslationCS{korodovat};  \dicIdiom{tærast}[upp]{ \dicPhraseIS{tærast upp}} \dicFlx{refl} \dicDirectTranslationCS{chřadnout, ztrácet se před očima}
\dicEntry[tærður] \dicTerm{tærður} \dicIPA{{t\smash{\textsuperscript{h}}}{a}{i}{r}{ð}{\textscy}{\textsubring{r}}} \dicPos{adj}[2]\dicFlx{}[-1] \dicDirectTranslationCS{zdecimovaný, vychrtlý}
\dicEntry[tæri] \dicTerm{tæri} \dicIPA{{t\smash{\textsuperscript{h}}}{a}{i}{\textlengthmark}{r}{\textsci}} \dicPos{n}[2] \dicFlx{(‑s)}[20] \dicPhraseIS{komast í tæri við e‑ð} \dicFlx{refl} \dicDirectTranslationCS{seznámit se s~(čím), přijít s~(čím) do styku}
\dicEntry[tæring] \dicTerm{tær··ing} \dicIPA{{t\smash{\textsuperscript{h}}}{a}{i}{\textlengthmark}{r}{i}{\ng}{\r{g}}} \dicPos{f}[4] \dicFlx{(‑ar)}[7] \textbf{1.} \dicSynonym*{eyðing málms} \dicDirectTranslationCS{koroze, rozežírání (kovu ap.)}  \textbf{2.} \dicFieldCat{med.} \dicSynonym{berklar} \dicDirectTranslationCS{tuberkulóza, souchotiny, úbytě}
\dicEntry[tæta] \dicTerm{tæt|a} \dicIPA{{t\smash{\textsuperscript{h}}}{a}{i}{\textlengthmark}{\textsubring{d}}{a}} \dicPos{v}[2] \dicFlx{(‑ti, ‑t)}[54] \dicFlx{acc} \dicDirectTranslationCS{roztrhnout, roztrhat, rozsápat, rozervat} \dicExampleIS{tæta e‑ð í sundur} \dicExampleCS{rozervat (co) na kusy};  \dicIdiom{tætast}{ \dicPhraseIS{tætast}} \dicFlx{refl} \dicDirectTranslationCS{být roztrhán\,/\addthin rozerván (na kusy ap.)}
\dicEntry[tætari] \dicTerm{tæt··ar|i} \dicIPA{{t\smash{\textsuperscript{h}}}{a}{i}{\textlengthmark}{\textsubring{d}}{a}{r}{\textsci}} \dicPos{m}[1] \dicFlx{(‑a, ‑ar)}[13] \dicDirectTranslationCS{drtička} \dicIndirectTranslationCS{(stroj)}
\dicEntry[tætingslegur] \dicTerm{tætings··legur} \dicIPA{{t\smash{\textsuperscript{h}}}{a}{i}{\textlengthmark}{\textsubring{d}}{i}{\ng}{\r{g}}{s}{l}{\textepsilon}{\textbabygamma}{\textscy}{\textsubring{r}}} \dicPos{adj}[1]\dicFlx{}[-8] \textbf{1.} \dicDirectTranslationCS{roztroušený, rozptýlený} \dicExampleIS{tætingslegt safn} \dicExampleCS{roztroušená sbírka}  \mbox{\textbf{2.} \dicDirectTranslationCS{rozcuchaný, pocuchaný}} \dicExampleIS{tætingslegt hár} \dicExampleCS{rozcuchané vlasy}
\dicEntry[tætla] \dicTerm{tætl|a} \dicIPA{{t\smash{\textsuperscript{h}}}{a}{i}{h}{\textsubring{d}}{l}{a}} \dicPos{f}[1] \dicFlx{(‑u, ‑ur)}[13] \dicSynonym{flygsa} \dicDirectTranslationCS{cár, útržek, kus (látky ap.)};  \dicPhraseIS{í tætlur} \dicFlx{adv} \dicDirectTranslationCS{na cáry\,/\addthin útržky} \dicExampleIS{rífa blaðið í tætlur} \dicExampleCS{roztrhat papír na cáry}
\dicEntry[tættur] \dicTerm{tættur} \dicIPA{{t\smash{\textsuperscript{h}}}{a}{i}{h}{\textsubring{d}}{\textscy}{\textsubring{r}}} \dicPos{adj}[1]\dicFlx{}[-10] \textbf{1.} \dicSynonym*{rifinn} \dicDirectTranslationCS{roztrhaný, rozervaný}  \textbf{2.} \dicSynonym*{sem líður illa} \dicDirectTranslationCS{rozpolcený, rozervaný} \dicIndirectTranslationCS{(o~člověku)}
\dicEntry[töð] \dicTerm{töð} \dicIPA{{t\smash{\textsuperscript{h}}}{\oe}{\textlengthmark}{\texttheta}} \dicPos{n} \dicFlx{pl nom} \dicLink{tað}
\dicEntry[töðu] \dicTerm{töðu} \dicIPA{{t\smash{\textsuperscript{h}}}{\oe}{\textlengthmark}{ð}{\textscy}} \dicPos{f} \dicFlx{sg gen} \dicLink{taða}
\dicEntry[töf] \dicTerm{töf} \dicIPA{{t\smash{\textsuperscript{h}}}{\oe}{\textlengthmark}{f}} \dicPos{f}[7] \dicFlx{(tafar, tafir)}[16] \dicSynonym{seinkun} \dicDirectTranslationCS{zpoždění, zdržení, prodleva} \dicExampleIS{Það verður töf á afgreiðslu.} \dicExampleCS{Servis bude mít zpoždění.}
\dicEntry[töfðum] \dicTerm{töfðum} \dicIPA{{t\smash{\textsuperscript{h}}}{\oe}{v}{ð}{\textscy}{\textsubring{m}}} \dicPos{v} \dicFlx{ind pf pl 1 pers} \dicLink{tefja}
\dicEntry[töff] \dicTerm{töff} \dicIPA{{t\smash{\textsuperscript{h}}}{\oe}{f}{\textlengthmark}} \dicPos{adj}[13] \dicFlx{indecl}[1] \dicLangCat{hovor.} \dicDirectTranslationCS{drsný, hustý, mazec (boty ap.)}
\dicEntry[töffari] \dicTerm{töff··ar|i} \dicIPA{{t\smash{\textsuperscript{h}}}{\oe}{f}{\textlengthmark}{a}{r}{\textsci}} \dicPos{m}[1] \dicFlx{(‑a, ‑ar)}[13] \dicLangCat{hovor.} \dicDirectTranslationCS{drsňák, drsňačka, tvrďák}
\dicEntry[töfl] \dicTerm{töfl} \dicIPA{{t\smash{\textsuperscript{h}}}{\oe}{\textsubring{b}}{\textsubring{l}}} \dicPos{n} \dicFlx{pl nom} \dicLink{tafl}
\dicEntry[töflu] \dicTerm{töflu} \dicIPA{{t\smash{\textsuperscript{h}}}{\oe}{\textsubring{b}}{l}{\textscy}} \dicPos{f} \dicFlx{sg gen} \dicLink{tafla}
\dicEntry[töflureiknir] \dicTerm{töflu··reikn|ir} \dicIPA{{t\smash{\textsuperscript{h}}}{\oe}{\textsubring{b}}{l}{\textscy}{r}{ei}{h}{\r{g}}{n}{\textsci}{\textsubring{r}}} \dicPos{m}[7] \dicFlx{(‑is)}[2] \dicFieldCat{poč.} \dicDirectTranslationCS{tabulkový procesor}
\dicEntry[töfra] \dicTerm{töfr|a} \dicIPA{{t\smash{\textsuperscript{h}}}{\oe}{v}{r}{a}} \dicPos{v}[1] \dicFlx{(‑aði)}[1] \textbf{1.} \dicSynonym{galdra} \dicDirectTranslationCS{očarovat, začarovat, zaklít}  \textbf{2.} \dicLangCat{přen.} \dicSynonym*{hrifa} \dicDirectTranslationCS{okouzlit, okouzlovat, očarovat (krása ap.)}
\dicEntry[töfrabrögð] \dicTerm{töfra··brögð} \dicIPA{{t\smash{\textsuperscript{h}}}{\oe}{v}{r}{a}{\textsubring{b}}{r}{\oe}{\textbabygamma}{\texttheta}} \dicPos{n}[2] \dicFlx{pl}[9] \dicDirectTranslationCS{kouzla, čáry}
\dicEntry[töfradrykkur] \dicTerm{töfra··drykk|ur} \dicIPA{{t\smash{\textsuperscript{h}}}{\oe}{v}{r}{a}{\textsubring{d}}{r}{\textsci}{h}{\r{g}}{\textscy}{\textsubring{r}}} \dicPos{m}[10] \dicFlx{(‑jar, ‑ir)}[23] \dicDirectTranslationCS{kouzelný\,/\addthin čarovný nápoj}
\dicEntry[töfragripur] \dicTerm{töfra··grip|ur} \dicIPA{{t\smash{\textsuperscript{h}}}{\oe}{v}{r}{a}{\r{g}}{r}{\textsci}{\textsubring{b}}{\textscy}{\textsubring{r}}} \dicPos{m}[9] \dicFlx{(‑s, ‑ir)}[8] \dicDirectTranslationCS{amulet, talisman}
\dicEntry[töfralæknir] \dicTerm{töfra··lækn|ir} \dicIPA{{t\smash{\textsuperscript{h}}}{\oe}{v}{r}{a}{l}{a}{i}{h}{\r{g}}{n}{\textsci}{\textsubring{r}}} \dicPos{m}[7] \dicFlx{(‑is, ‑ar)}[1] \dicDirectTranslationCS{medicinman}
\begin{xtolerant}{}{1pt}
\dicEntry[töframaður] \dicTerm{töfra··|maður} \dicIPA{{t\smash{\textsuperscript{h}}}{\oe}{v}{r}{a}{m}{a}{ð}{\textscy}{\textsubring{r}}} \dicPos{m}[13] \dicFlx{(‑manns, ‑menn)}[2] \dicSynonym{gal\-dramaður} \dicDirectTranslationCS{čaroděj(ka), kouzelník, kouzelnice}
\end{xtolerant}
\dicEntry[töframáttur] \dicTerm{töfra··mátt|ur} \dicIPA{{t\smash{\textsuperscript{h}}}{\oe}{v}{r}{a}{m}{au}{h}{\textsubring{d}}{\textscy}{\textsubring{r}}} \dicPos{m}[12] \dicFlx{(‑ar)}[8] \dicDirectTranslationCS{kouzelná\,/\addthin magická moc}
\dicEntry[töfrandi] \dicTerm{töfr··andi} \dicIPA{{t\smash{\textsuperscript{h}}}{\oe}{v}{r}{a}{n}{\textsubring{d}}{\textsci}} \dicPos{adj}[13] \dicFlx{indecl}[1] \dicDirectTranslationCS{okouzlující, čarovný, úchvatný}
\dicEntry[töfrar] \dicTerm{töfrar} \dicIPA{{t\smash{\textsuperscript{h}}}{\oe}{v}{r}{a}{\textsubring{r}}} \dicPos{m}[14] \dicFlx{pl}[1] \textbf{1.} \dicSynonym{galdur} \dicDirectTranslationCS{kouzla, čáry, magie}  \textbf{2.} \dicSynonym{þokki} \dicDirectTranslationCS{šarm, čarovná krása, kouzlo} \dicExampleIS{njóta töfra náttúrunnar} \dicExampleCS{těšit se z~čarovné krásy přírody}
\dicEntry[töfrasproti] \dicTerm{töfra··sprot|i} \dicIPA{{t\smash{\textsuperscript{h}}}{\oe}{v}{r}{a}{s}{\textsubring{b}}{r}{\textopeno}{\textsubring{d}}{\textsci}} \dicPos{m}[1] \dicFlx{(‑a, ‑ar)}[1] \textbf{1.} \dicSynonym*{stafur galdramanns} \dicDirectTranslationCS{kouzelná hůlka}  \textbf{2.} \dicSynonym*{eldhúsáhald} \dicDirectTranslationCS{ruční\,/\addthin tyčový mixér}
\dicEntry[töfraveröld] \dicTerm{töfra··ver·|öld} \dicIPA{{t\smash{\textsuperscript{h}}}{\oe}{v}{r}{a}{v}{\textepsilon}{r}{\oe}{l}{\textsubring{d}}} \dicPos{f}[7] \dicFlx{(‑aldar, ‑aldir)}[16] \dicDirectTranslationCS{magický\,/\addthin kouzelný svět}
\dicEntry[töggur] \dicTerm{tögg|ur} \dicIPA{{t\smash{\textsuperscript{h}}}{\oe}{\r{g}}{\textlengthmark}{\textscy}{\textsubring{r}}} \dicPos{m}[6] \dicFlx{(‑s)}[17] \dicPhraseIS{það er töggur í e‑m} \dicFlx{impers} \dicDirectTranslationCS{(kdo) má kuráž}
\dicEntry[tök] \dicTerm{tök} \dicIPA{{t\smash{\textsuperscript{h}}}{\oe}{\textlengthmark}{\r{g}}} \dicPos{n} \dicFlx{pl nom} \dicLink{tak}
\dicEntry[tökuorð] \dicTerm{töku··orð} \dicIPA{{t\smash{\textsuperscript{h}}}{\oe}{\textlengthmark}{\r{g}}{\textscy}{\textopeno}{r}{\texttheta}} \dicPos{n}[2] \dicFlx{(‑s, ‑)}[5] \dicFieldCat{jaz.} \dicDirectTranslationCS{přejaté slovo, přejímka, výpůjčka}
\dicEntry[tökuvél] \dicTerm{töku··vél} \dicIPA{{t\smash{\textsuperscript{h}}}{\oe}{\textlengthmark}{\r{g}}{\textscy}{v}{j}{\textepsilon}{\textsubring{l}}} \dicPos{f}[4] \dicFlx{(‑ar, ‑ar)}[1] \dicDirectTranslationCS{(video)kamera}
\dicEntry[töl] \dicTerm{töl} \dicIPA{{t\smash{\textsuperscript{h}}}{\oe}{\textlengthmark}{\textsubring{l}}} \dicPos{n} \dicFlx{pl nom} \dicLink{tal}
\dicEntry[töldum] \dicTerm{töldum} \dicIPA{{t\smash{\textsuperscript{h}}}{\oe}{l}{\textsubring{d}}{\textscy}{\textsubring{m}}} \dicPos{v} \dicFlx{ind pf pl 1 pers} \dicLink{telja}
\dicEntry[tölfræði] \dicTerm{töl··fræð|i} \dicIPA{{t\smash{\textsuperscript{h}}}{\oe}{\textsubring{l}}{f}{r}{a}{i}{ð}{\textsci}} \dicPos{f}[3] \dicFlx{(‑i)}[3] \dicDirectTranslationCS{statistika} \dicIndirectTranslationCS{(vědní obor)}
\dicEntry[tölfræðilegur] \dicTerm{töl·fræði··legur} \dicIPA{{t\smash{\textsuperscript{h}}}{\oe}{\textsubring{l}}{f}{r}{a}{i}{ð}{\textsci}{l}{\textepsilon}{\textbabygamma}{\textscy}{\textsubring{r}}} \dicPos{adj}[1]\dicFlx{}[-8] \dicDirectTranslationCS{statistický}
\dicEntry[tölfræðingur] \dicTerm{töl·fræð··ing|ur} \dicIPA{{t\smash{\textsuperscript{h}}}{\oe}{\textsubring{l}}{f}{r}{a}{i}{ð}{i}{\ng}{\r{g}}{\textscy}{\textsubring{r}}} \dicPos{m}[6] \dicFlx{(‑s, ‑ar)}[8] \dicDirectTranslationCS{statistik, statistička}
\dicEntry[tölt] \dicTerm{tölt} \dicIPA{{t\smash{\textsuperscript{h}}}{\oe}{\textsubring{l}}{\textsubring{d}}} \dicPos{n}[2] \dicFlx{(‑s)}[2] \dicDirectTranslationCS{tölt} \dicIndirectTranslationCS{(čtyřdobý rychlý mimochod)}
\dicEntry[tölu] \dicTerm{tölu} \dicIPA{{t\smash{\textsuperscript{h}}}{\oe}{\textlengthmark}{l}{\textscy}} \dicPos{f} \dicFlx{sg gen} \dicLink{tala\smash{\textsuperscript{1}}}
\dicEntry[tölublað] \dicTerm{tölu··|blað} \dicIPA{{t\smash{\textsuperscript{h}}}{\oe}{\textlengthmark}{l}{\textscy}{\textsubring{b}}{l}{a}{\texttheta}} \dicPos{n}[2] \dicFlx{(‑blaðs, ‑blöð)}[8] \dicDirectTranslationCS{číslo (časopisu ap.)}
\dicEntry[töludálkur] \dicTerm{tölu··dálk|ur} \dicIPA{{t\smash{\textsuperscript{h}}}{\oe}{\textlengthmark}{l}{\textscy}{\textsubring{d}}{au}{\textsubring{l}}{\r{g}}{\textscy}{\textsubring{r}}} \dicPos{m}[6] \dicFlx{(‑s, ‑ar)}[24] \dicDirectTranslationCS{sloupec čísel}
\dicEntry[tölugildi] \dicTerm{tölu··gildi} \dicIPA{{t\smash{\textsuperscript{h}}}{\oe}{\textlengthmark}{l}{\textscy}{\r{\textObardotlessj}}{\textsci}{l}{\textsubring{d}}{\textsci}} \dicPos{n}[2] \dicFlx{(‑s, ‑)}[14] \dicFieldCat{mat.} \dicDirectTranslationCS{číselná hodnota}
\dicEntry[tölulegur] \dicTerm{tölu··legur} \dicIPA{{t\smash{\textsuperscript{h}}}{\oe}{\textlengthmark}{l}{\textscy}{l}{\textepsilon}{\textbabygamma}{\textscy}{\textsubring{r}}} \dicPos{adj}[1]\dicFlx{}[-8] \dicDirectTranslationCS{statistický, numerický, číselný}
\dicEntry[töluliður] \dicTerm{tölu··lið|ur} \dicIPA{{t\smash{\textsuperscript{h}}}{\oe}{\textlengthmark}{l}{\textscy}{l}{\textsci}{ð}{\textscy}{\textsubring{r}}} \dicPos{m}[10] \dicFlx{(‑ar\,/\addthin ‑s, ‑ir)}[33] \dicDirectTranslationCS{odstavec, paragraf (trestního zákoníku ap.)}
\dicEntry[töluorð] \dicTerm{tölu··orð} \dicIPA{{t\smash{\textsuperscript{h}}}{\oe}{\textlengthmark}{l}{\textscy}{\textopeno}{r}{\texttheta}} \dicPos{n}[2] \dicFlx{(‑s, ‑)}[5] \dicFieldCat{jaz.} \dicDirectTranslationCS{číslovka} \textit{(l.~{\textLA{nomen numerale}})}
\dicEntry[töluset] \dicTerm{tölu··set} \dicIPA{{t\smash{\textsuperscript{h}}}{\oe}{\textlengthmark}{l}{\textscy}{s}{\textepsilon}{\textsubring{d}}} \dicPos{v} \dicFlx{ind praes sg 1 pers} \dicLink{tölusetja}
\dicEntry[tölusetja] \dicTerm{tölu··|setja} \dicIPA{{t\smash{\textsuperscript{h}}}{\oe}{\textlengthmark}{l}{\textscy}{s}{\textepsilon}{\textsubring{d}}{j}{a}} \dicPos{v}[4] \dicFlx{(‑set, ‑setti, ‑settum, ‑setti, ‑sett)}[26] \dicFlx{acc} \dicDirectTranslationCS{(o)číslovat}
\dicEntry[tölusett] \dicTerm{tölu··sett} \dicIPA{{t\smash{\textsuperscript{h}}}{\oe}{\textlengthmark}{l}{\textscy}{s}{\textepsilon}{h}{\textsubring{d}}} \dicPos{v} \dicFlx{supin} \dicLink{tölusetja}
\dicEntry[tölusetti] \dicTerm{tölu··setti} \dicIPA{{t\smash{\textsuperscript{h}}}{\oe}{\textlengthmark}{l}{\textscy}{s}{\textepsilon}{h}{\textsubring{d}}{\textsci}} \dicPos{v} \dicFlx{ind\,/\addthin con pf sg 1 pers} \dicLink{tölusetja}
\dicEntry[tölusettum] \dicTerm{tölu··settum} \dicIPA{{t\smash{\textsuperscript{h}}}{\oe}{\textlengthmark}{l}{\textscy}{s}{\textepsilon}{h}{\textsubring{d}}{\textscy}{\textsubring{m}}} \dicPos{v} \dicFlx{ind pf pl 1 pers} \dicLink{tölusetja}
\dicEntry[tölusettur] \dicTerm{tölu··settur} \dicIPA{{t\smash{\textsuperscript{h}}}{\oe}{\textlengthmark}{l}{\textscy}{s}{\textepsilon}{h}{\textsubring{d}}{\textscy}{\textsubring{r}}} \dicPos{adj}[1]\dicFlx{}[-13] \dicDirectTranslationCS{očíslovaný} \dicExampleIS{tölusett sæti} \dicExampleCS{očíslované sedadlo}
\dicEntry[tölustafur] \textls[15]{\dicTerm{tölu··staf|ur} \dicIPA{{t\smash{\textsuperscript{h}}}{\oe}{\textlengthmark}{l}{\textscy}{s}{\textsubring{d}}{a}{v}{\textscy}{\textsubring{r}}} \dicPos{m}[9] \dicFlx{(‑s, ‑ir)}[9] \dicDirectTranslationCS{číslice, cifra}} \dicExampleIS{skrifa þrettán með tölustöfum} \dicExampleCS{napsat třináct číslicemi}
\begin{xtolerant}{}{1pt}
\dicEntry[töluverður] \dicTerm{tölu··verður} \dicsymFrequent\addthinS\dicIPA{{t\smash{\textsuperscript{h}}}{\oe}{\textlengthmark}{l}{\textscy}{v}{\textepsilon}{r}{ð}{\textscy}{\textsubring{r}}}\addthinS\dicPos{adj}[2]\dicFlx{}[-4]\addthinS\dicSynonym{talsverður} \dicDirectTranslationCS{znač\-ný, znatelný} \dicExampleIS{töluverður skortur} \dicExampleCS{značný nedostatek}
\end{xtolerant}
\dicEntry[töluvert] \dicTerm{tölu··vert} \dicsymFrequent\  \dicIPA{{t\smash{\textsuperscript{h}}}{\oe}{\textlengthmark}{l}{\textscy}{v}{\textepsilon}{\textsubring{r}}{\textsubring{d}}} \dicPos{adv} \dicSynonym{býsna} \dicDirectTranslationCS{značně, znatelně} \dicExampleIS{Hún er töluvert yngri en hann.} \dicExampleCS{Je značně mladší než on.}
\dicEntry[tölva] \dicTerm{tölv|a} \dicsymFrequent\  \dicIPA{{t\smash{\textsuperscript{h}}}{\oe}{l}{v}{a}} \dicPos{f}[1] \dicFlx{(‑u, ‑ur)}[7] \dicFieldCat{poč.} \dicDirectTranslationCS{počítač, komputer} \dicExampleIS{Ættu ung börn að nota tölvur?} \dicExampleCS{Měly by malé děti používat počítače?}
\dicEntry[tölvubúnaður] \dicTerm{tölvu··bú·nað|ur} \dicIPA{{t\smash{\textsuperscript{h}}}{\oe}{l}{v}{\textscy}{\textsubring{b}}{u}{n}{a}{ð}{\textscy}{\textsubring{r}}} \dicPos{m}[10] \dicFlx{(‑ar)}[9] \dicFieldCat{poč.} \dicDirectTranslationCS{počítačové\,/\addthin IT vybavení}
\dicEntry[tölvuforrit] \dicTerm{tölvu··for·rit} \dicIPA{{t\smash{\textsuperscript{h}}}{\oe}{l}{v}{\textscy}{f}{\textopeno}{r}{\textlengthmark}{\textsci}{\textsubring{d}}} \dicPos{n}[2] \dicFlx{(‑s, ‑)}[5] \dicFieldCat{poč.} \dicSynonym{forrit} \dicDirectTranslationCS{(počítačový) program, (počítačová) aplikace}
\dicEntry[tölvufræði] \dicTerm{tölvu··fræð|i}\dicTerm{, tölvunarfræði} \dicIPA{{t\smash{\textsuperscript{h}}}\-{\oe}\-{l}\-{v}\-{\textscy}\-{f}\-{r}\-{a}\-{i}\-{ð}\-{\textsci}\-} \dicPos{f}[3] \dicFlx{(‑i)}[3] \dicDirectTranslationCS{informatika}
\dicEntry[tölvufræðingur] \dicTerm{tölvu·fræð··ing|ur}\textls[-15]{\dicTerm{, tölvunarfræðingur} \dicIPA{{t\smash{\textsuperscript{h}}}\-{\oe}\-{l}\-{v}\-{\textscy}\-{f}\-{r}\-{a}\-{i}\-{ð}\-{i}\-{\ng}\-{\r{g}}\-{\textscy}\-{\textsubring{r}}}} \dicPos{m}[6] \dicFlx{(‑s, ‑ar)}[8] \dicDirectTranslationCS{informatik, informatička}
\dicEntry[tölvugrafík] \dicTerm{tölvu··grafík} \dicIPA{{t\smash{\textsuperscript{h}}}{\oe}{l}{v}{\textscy}{\r{g}}{r}{a}{f}{i}{\r{g}}} \dicPos{f}[10] \dicFlx{(‑ur)}[2] \dicFieldCat{poč.} \dicDirectTranslationCS{počítačová grafika}
\dicEntry[tölvuhakkari] \dicTerm{tölvu··hakk·ar|i} \dicIPA{{t\smash{\textsuperscript{h}}}{\oe}{l}{v}{\textscy}{h}{a}{h}{\r{g}}{a}{r}{\textsci}} \dicPos{m}[1] \dicFlx{(‑a, ‑ar)}[10] \dicDirectTranslationCS{(počítačový) hacker, (počítačová) hackerka}
\dicEntry[tölvukerfi] \dicTerm{tölvu··kerfi} \dicIPA{{t\smash{\textsuperscript{h}}}{\oe}{l}{v}{\textscy}{c\smash{\textsuperscript{h}}}{\textepsilon}{r}{v}{\textsci}} \dicPos{n}[2] \dicFlx{(‑s, ‑)}[14] \dicFieldCat{poč.} \dicDirectTranslationCS{počítačový\,/\addthin IT systém}
\dicEntry[tölvukunnátta] \dicTerm{tölvu··kunn·átt|a} \dicIPA{{t\smash{\textsuperscript{h}}}{\oe}{l}{v}{\textscy}{k\smash{\textsuperscript{h}}}{\textscy}{n}{au}{h}{\textsubring{d}}{a}} \dicPos{f}[1] \dicFlx{(‑u)}[5] \dicDirectTranslationCS{IT znalosti\,/\addthin dovednosti, IT kompetence}
\dicEntry[tölvuleikur] \dicTerm{tölvu··leik|ur} \dicIPA{{t\smash{\textsuperscript{h}}}{\oe}{l}{v}{\textscy}{l}{ei}{\r{g}}{\textscy}{\textsubring{r}}} \dicPos{m}[9] \dicFlx{(‑s, ‑ir)}[15] \dicFieldCat{poč.} \dicDirectTranslationCS{počítačová hra}
\dicEntry[tölvunarfræði] \dicTerm{tölvunar··fræð|i} \dicIPA{{t\smash{\textsuperscript{h}}}{\oe}{l}{v}{\textscy}{n}{a}{\textsubring{r}}{f}{r}{a}{i}{ð}{\textsci}} \dicPos{f}[3] \dicFlx{(‑i)}[3] \dicLink{tölvufræði}
\dicEntry[tölvunarfræðingur] \dicTerm{tölvunar·fræð··ing|ur} \dicIPA{{t\smash{\textsuperscript{h}}}\-{\oe}\-{l}\-{v}\-{\textscy}\-{n}\-{a}\-{\textsubring{r}}\-{f}\-{r}\-{a}\-{i}\-{ð}\-{i}\-{\ng}\-{\r{g}}\-{\textscy}\-{\textsubring{r}}\-} \dicPos{m}[6] \dicFlx{(‑s, ‑ar)}[8] \dicLink{tölvufræðingur}
\dicEntry[tölvunet] \dicTerm{tölvu··net} \dicIPA{{t\smash{\textsuperscript{h}}}{\oe}{l}{v}{\textscy}{n}{\textepsilon}{\textsubring{d}}} \dicPos{n}[2] \dicFlx{(‑s, ‑)}[5] \dicFieldCat{poč.} \dicDirectTranslationCS{(počítačová) síť}
\dicEntry[tölvuorðabók] \dicTerm{tölvu··orða·|bók} \dicIPA{{t\smash{\textsuperscript{h}}}{\oe}{l}{v}{\textscy}{\textopeno}{r}{ð}{a}{\textsubring{b}}{ou}{\r{g}}} \dicPos{f}[8] \dicFlx{(‑bókar, ‑bækur)}[5] \dicDirectTranslationCS{elektronický slovník}
\dicEntry[tölvupóstfang] \dicTerm{tölvu··póst·|fang} \dicIPA{{t\smash{\textsuperscript{h}}}{\oe}{l}{v}{\textscy}{p\smash{\textsuperscript{h}}}{ou}{s}{\textsubring{d}}{f}{au}{\ng}{\r{g}}} \dicPos{n}[2] \dicFlx{(‑fangs, ‑föng)}[8] \dicFieldCat{poč.} \dicDirectTranslationCS{e‑mailová adresa}
\dicEntry[tölvupóstur] \dicTerm{tölvu··póst|ur} \dicIPA{{t\smash{\textsuperscript{h}}}{\oe}{l}{v}{\textscy}{p\smash{\textsuperscript{h}}}{ou}{s}{\textsubring{d}}{\textscy}{\textsubring{r}}} \dicPos{m}[6] \dicFlx{(‑s)}[7] \textbf{1.} \dicFieldCat{poč.} \dicSynonym{boðskipti} \dicDirectTranslationCS{e‑mail, elektronická pošta} \dicIndirectTranslationCS{(způsob odesílání, doručování a~přijímání zpráv)}  \textbf{2.} \dicFieldCat{poč.} \dicSynonym{tölvuskeyti} \dicDirectTranslationCS{e‑mail, elektronická zpráva} \dicExampleIS{senda tölvupóst} \dicExampleCS{poslat e‑mail}
\dicEntry[tölvuprentari] \dicTerm{tölvu··prent·ar|i} \dicIPA{{t\smash{\textsuperscript{h}}}{\oe}{l}{v}{\textscy}{p\smash{\textsuperscript{h}}}{r}{\textepsilon}{\textsubring{n}}{\textsubring{d}}{a}{r}{\textsci}} \dicPos{m}[1] \dicFlx{(‑a, ‑ar)}[13] \dicFieldCat{poč.} \dicDirectTranslationCS{(počítačová) tiskárna}
\dicEntry[tölvuskeyti] \dicTerm{tölvu··skeyti} \dicIPA{{t\smash{\textsuperscript{h}}}{\oe}{l}{v}{\textscy}{s}{\r{\textObardotlessj}}{ei}{\textsubring{d}}{\textsci}} \dicPos{n}[2] \dicFlx{(‑s, ‑)}[14] \dicFieldCat{poč.} \dicDirectTranslationCS{e‑mail, mail, elektronická zpráva}
\dicEntry[tölvuskjár] \dicTerm{tölvu··skjá|r} \dicIPA{{t\smash{\textsuperscript{h}}}{\oe}{l}{v}{\textscy}{s}{\r{\textObardotlessj}}{au}{\textsubring{r}}} \dicPos{m}[9] \dicFlx{(‑s, ‑ir)}[12] \dicFieldCat{poč.} \dicDirectTranslationCS{monitor, (počítačová) obrazovka}
\dicEntry[tölvuskrá] \dicTerm{tölvu··skrá\smash{\textsuperscript{1}}} \dicIPA{{t\smash{\textsuperscript{h}}}{\oe}{l}{v}{\textscy}{s}{\r{g}}{r}{au}} \dicPos{f}[4] \dicFlx{(‑r\,/\addthin ‑ar, ‑r)}[21] \dicFieldCat{poč.} \dicSynonym{skrá\smash{\textsuperscript{1}}} \dicDirectTranslationCS{(počítačový) soubor}
\dicEntry[tölvuskrá] \dicTerm{tölvu··skrá\smash{\textsuperscript{2}}} \dicIPA{{t\smash{\textsuperscript{h}}}{\oe}{l}{v}{\textscy}{s}{\r{g}}{r}{au}} \dicPos{v}[2] \dicFlx{(‑ði, ‑ð)}[118] \dicFlx{acc} \dicDirectTranslationCS{zadat\,/\addthin psát do počítače} \dicExampleIS{tölvuskrá textann} \dicExampleCS{psát text do počítače}
\dicEntry[tölvusnillingur] \dicTerm{tölvu··snill·ing|ur} \dicIPA{{t\smash{\textsuperscript{h}}}{\oe}{l}{v}{\textscy}{s}{\textsubring{d}}{n}{\textsci}{\textsubring{d}}{\textsubring{l}}{i}{\ng}{\r{g}}{\textscy}{\textsubring{r}}} \dicPos{m}[6] \dicFlx{(‑s, ‑ar)}[8] \dicDirectTranslationCS{počítačový génius}
\dicEntry[tölvustýrður] \dicTerm{tölvu··stýrður} \dicIPA{{t\smash{\textsuperscript{h}}}{\oe}{l}{v}{\textscy}{s}{\textsubring{d}}{i}{r}{ð}{\textscy}{\textsubring{r}}} \dicPos{adj}[2]\dicFlx{}[-4] \dicDirectTranslationCS{řízený počítačem}
\dicEntry[tölvutækur] \dicTerm{tölvu··tækur} \dicIPA{{t\smash{\textsuperscript{h}}}{\oe}{l}{v}{\textscy}{t\smash{\textsuperscript{h}}}{a}{i}{\r{g}}{\textscy}{\textsubring{r}}} \dicPos{adj}[1]\dicFlx{}[-1] \dicFieldCat{poč.} \dicDirectTranslationCS{elektronický, digitální} \dicExampleIS{Bókin er á tölvutæku formi.} \dicExampleCS{Knížka je v~elektronické podobě.}
\dicEntry[tölvuveira] \dicTerm{tölvu··veir|a} \dicIPA{{t\smash{\textsuperscript{h}}}{\oe}{l}{v}{\textscy}{v}{ei}{r}{a}} \dicPos{f}[1] \dicFlx{(‑u, ‑ur)}[7] \dicFieldCat{poč.} \dicDirectTranslationCS{(počítačový) virus\,/\addthin vir}
\dicEntry[tölvuver] \dicTerm{tölvu··ver} \dicIPA{{t\smash{\textsuperscript{h}}}{\oe}{l}{v}{\textscy}{v}{\textepsilon}{\textsubring{r}}} \dicPos{n}[2] \dicFlx{(‑s, ‑)}[5] \dicDirectTranslationCS{počítačová místnost\,/\addthin učebna}
\dicEntry[tölvuverslun] \dicTerm{tölvu··versl|un} \dicIPA{{t\smash{\textsuperscript{h}}}{\oe}{l}{v}{\textscy}{v}{\textepsilon}{\textsubring{r}}{s}{\textsubring{d}}{l}{\textscy}{\textsubring{n}}} \dicPos{f}[7] \dicFlx{(‑unar, ‑anir)}[8] \dicDirectTranslationCS{obchod s~počítači}
\dicEntry[tölvuvæða] \dicTerm{tölvu··væ|ða} \dicIPA{{t\smash{\textsuperscript{h}}}{\oe}{l}{v}{\textscy}{v}{a}{i}{ð}{a}} \dicPos{v}[2] \dicFlx{(‑ddi, ‑tt)}[171] \dicFlx{acc} \dicDirectTranslationCS{zavádět informační technologie}
\dicEntry[tölvuvæðing] \dicTerm{tölvu··væð·ing} \dicIPA{{t\smash{\textsuperscript{h}}}{\oe}{l}{v}{\textscy}{v}{a}{i}{ð}{i}{\ng}{\r{g}}} \dicPos{f}[4] \dicFlx{(‑ar)}[7] \dicDirectTranslationCS{zavádění informačních technologií}
\dicEntry[töm] \dicTerm{töm} \dicIPA{{t\smash{\textsuperscript{h}}}{\oe}{\textlengthmark}{\textsubring{m}}} \dicPos{adj} \dicFlx{f sg nom pos} \dicLink{tamur}
\dicEntry[tömdum] \dicTerm{tömdum} \dicIPA{{t\smash{\textsuperscript{h}}}{\oe}{m}{\textsubring{d}}{\textscy}{\textsubring{m}}} \dicPos{v} \dicFlx{ind pf pl 1 pers} \dicLink{temja}
\dicEntry[töng] \dicTerm{töng} \dicIPA{{t\smash{\textsuperscript{h}}}{\oe i}{\ng}{\r{g}}} \dicPos{f}[8] \dicFlx{(tangar, tengur\,/\addthin tangir)}[9] \dicDirectTranslationCS{kleště} \dicExampleIS{vírtöng} \dicExampleCS{kleště na drát}
\dicEntry[tönn] \dicTerm{tönn} \dicsymFrequent\  \dicIPA{{t\smash{\textsuperscript{h}}}{\oe}{n}{\textlengthmark}} \dicPos{f}[8] \dicFlx{(tannar, tennur\,/\addthin tannir\,/\addthin tönnur)}[14] \textbf{1.} \dicDirectTranslationCS{zub} \dicIndirectTranslationCS{(v~puse)} \dicExampleIS{Tennurnar í honum eru skemmdar.} \dicExampleCS{Má zkažené zuby.};  \dicPhraseIS{e‑að stendur í tönnunum á e‑m} \dicDirectTranslationCS{(kdo) má (co) mezi zuby (zbytky jídla ap.)};  \dicPhraseIS{bursta tennur(nar)} \dicDirectTranslationCS{čistit si zuby};  \dicPhraseIS{falskar tennur} \dicSynonym{gervitennur} \dicDirectTranslationCS{umělý chrup};  \dicPhraseIS{gnísta tönnum} \dicLangCat{přen.} \dicDirectTranslationCS{skřípat zuby}  \textbf{2.} \dicSynonym{teinn} \dicDirectTranslationCS{zub} \dicIndirectTranslationCS{(část nástrojů připomínající špičatým tvarem zub)} \dicExampleIS{tönn í sög} \dicExampleCS{zub pily};  \dicIdiom{tönn}{ \dicPhraseIS{standast tímans tönn}} \dicFlx{refl} \dicLangCat{přen.} \dicDirectTranslationCS{odolat zubu času};  \dicPhraseIS{vera á milli tannanna á e‑m} \dicDirectTranslationCS{dostat se do řečí (koho)}
\dicEntry[tönnlast] \dicTerm{tönnl|ast} \dicIPA{{t\smash{\textsuperscript{h}}}{\oe}{\ng}{l}{a}{s}{\textsubring{d}}} \dicPos{v}[1] \dicFlx{(‑aðist)}[95] \dicFlx{refl} \dicPhraseIS{tönnlast á e‑u} \dicDirectTranslationCS{omílat (co) dokola} \dicExampleIS{tönnlast á því sama} \dicExampleCS{omílat stále to samé dokola}
\dicEntry[töp] \dicTerm{töp} \dicIPA{{t\smash{\textsuperscript{h}}}{\oe}{\textlengthmark}{\textsubring{b}}} \dicPos{n} \dicFlx{pl nom} \dicLink{tap}
\dicEntry[törn] \dicTerm{törn} \dicIPA{{t\smash{\textsuperscript{h}}}{\oe}{r}{\textsubring{d}}{\textsubring{n}}} \dicPos{f}[7] \dicFlx{(tarnar, tarnir)}[16] \dicDirectTranslationCS{zabíračka, dřina}
\dicEntry[tösku] \dicTerm{tösku} \dicIPA{{t\smash{\textsuperscript{h}}}{\oe}{s}{\r{g}}{\textscy}} \dicPos{f} \dicFlx{sg gen} \dicLink{taska}
\dicEntry[tötralega] \dicTerm{tötra··lega} \dicIPA{{t\smash{\textsuperscript{h}}}{\oe}{\textlengthmark}{\textsubring{d}}{r}{a}{l}{\textepsilon}{\textbabygamma}{a}} \dicPos{adv} \dicDirectTranslationCS{otrhaně, odraně}
\dicEntry[tötralegur] \dicTerm{tötra··legur} \dicIPA{{t\smash{\textsuperscript{h}}}{\oe}{\textlengthmark}{\textsubring{d}}{r}{a}{l}{\textepsilon}{\textbabygamma}{\textscy}{\textsubring{r}}} \dicPos{adj}[1]\dicFlx{}[-8] \dicSynonym{druslulegur} \dicDirectTranslationCS{otrhaný, odraný}

\else
\makeatletter
\@for\next:=\alphabetF\do{\input{letters/\next.tex}}
\makeatother 
\fi
\restoregeometry
\pagestyle{plain}

\fi

% ================================================================ PHONETICS =

\ifinputphon

\cleardoublepage

\chapter{Seznam islandských fonémů}              \label{sec:phon_phonems}

\begin{table}[h]
\begin{tabular}{lllll} \toprule

\textbf{foném} & \textbf{vysvětlivka} & \textbf{příklad} & \textbf{fonetický přepis} & \textbf{překlad} \\
\midrule
{\textipa{[{a}]}} & otevřená přední nezaokrouhlená samohl. & rata & {\textipa{[{r}{a}{\textlengthmark}{\textsubring{d}}{a}]}} & trefit \\ 
{\textipa{[{\textepsilon}]}} & polootevřená přední nezaokrouhlená samohl. & meta & {\textipa{[{m}{\textepsilon}{\textlengthmark}{\textsubring{d}}{a}]}} & ocenit \\ 
{\textipa{[{e}]}} & polozavřená přední nezaokrouhlená samohl. & heiður & {\textipa{[{h}{ei}{\textlengthmark}{ð}{\textscy}{\textsubring{r}}]}} & čest \\
{\textipa{[{\textsci}]}} & téměř zavřená téměř přední nezaokrouhlená samohl. & vinna & {\textipa{[{v}{\textsci}{n}{\textlengthmark}{a}]}} & pracovat \\ 
{\textipa{[{i}]}} & zavřená přední nezaokrouhlená samohl. & fínn & {\textipa{[{f}{i}{\textsubring{d}}{\textsubring{n}}]}} & jemný \\ 
{\textipa{[{\textopeno}]}} & polootevřená zadní zaokrouhlená samohl. & lofa & {\textipa{[{l}{\textopeno}{\textlengthmark}{v}{a}]}} & slibovat \\ 
{\textipa{[{o}]}} & polozavřená zadní zaokrouhlená samohl. & bóndi & {\textipa{[{\textsubring{b}}{ou}{n}{\textsubring{d}}{\textsci}]}} & farmář \\ 
{\textipa{[{\textscy}]}} & téměř zavřená téměř přední zaokrouhlená samohl. & hundur & {\textipa{[{h}{\textscy}{n}{\textsubring{d}}{\textscy}{\textsubring{r}}]}} & pes \\ 
{\textipa{[{u}]}} & zavřená zadní zaokrouhlená samohl. & súkkulaði & {\textipa{[{s}{u}{h}{\r{g}}{\textscy}{l}{a}{ð}{\textsci}]}} & čokoláda \\ 
{\textipa{[{\oe}]}} & téměř otevřená přední nezaokrouhlená samohl. & lönd & {\textipa{[{l}{\oe}{n}{\textsubring{d}}]}} & země \\ 
{\textipa{[{m}]}} & bilabiální nazála & mýri & {\textipa{[{m}{i}{\textlengthmark}{r}{\textsci}]}} & mokřina \\ 
{\textipa{[{\textsubring{m}}]}} & neznělá bilabiální nazála & þyrmt & {\textipa{[{\texttheta}{\textsci}{\textsubring{r}}{\textsubring{m}}{\textsubring{d}}]}} & {\textit{$\shortrightarrow$ þyrma}} \\  
{\textipa{[{n}]}} & alveolární nazála & nemandi & {\textipa{[{n}{\textepsilon}{\textlengthmark}{m}{a}{n}{\textsubring{d}}{\textsci}]}} & žák \\ 
{\textipa{[{\textsubring{n}}]}} & neznělá alveolární nazála & mennt & {\textipa{[{m}{\textepsilon}{\textsubring{n}}{\textsubring{d}}]}} & vzdělání \\ 
{\textipa{[{\textltailn}]}} & palatální nazála & angi & {\textipa{[{au}{\textltailn}{\r{\textObardotlessj}}{\textsci}]}} & drobek \\ 
{\textipa{[{\r{\textltailn}}]}} & neznělá palatální nazála & banki & {\textipa{[{\textsubring{b}}{au}{\r{\textltailn}}{\r{\textObardotlessj}}{\textsci}]}} & banka \\
{\textipa{[{\ng}]}} & velární nazála & skyggndur & {\textipa{[{s}{\r{\textObardotlessj}}{\textsci}{\ng}{\textsubring{d}}{\textscy}{\textsubring{r}}]}} & {\textit{$\shortrightarrow$ skyggna}} \\ 
{\textipa{[{\r{\ng}}]}} & neznělá velární nazála & skyggnt & {\textipa{[{s}{\r{\textObardotlessj}}{\textsci}{\r{\ng}}{\textsubring{d}}]}} & {\textit{$\shortrightarrow$ skyggna}} \\ 
{\textipa{[{\textsubring{b}}]}} & neznělá bilabiální ploziva & bara & {\textipa{[{\textsubring{b}}{a}{\textlengthmark}{r}{a}]}} & jenom \\ 
{\textipa{[{p\textsuperscript{h}}]}} & neznělá bilabiální ploziva s přídechem & prestur & {\textipa{[{p\textsuperscript{h}}{r}{\textepsilon}{s}{\textsubring{d}}{\textscy}{\textsubring{r}}]}} & kněz \\ 
{\textipa{[{\textsubring{d}}]}} & neznělá alveolární ploziva & dæmi & {\textipa{[{\textsubring{d}}{a}{i}{\textlengthmark}{m}{\textsci}]}} & příklad \\ 
{\textipa{[{t\textsuperscript{h}}]}} & neznělá alveolární ploziva s přídechem & tómur & {\textipa{[{t\textsuperscript{h}}{ou}{\textlengthmark}{m}{\textscy}{\textsubring{r}}]}} & prázdný \\ 
{\textipa{[{\r{\textObardotlessj}}]}} & neznělá palatální ploziva & gefa & {\textipa{[{\r{\textObardotlessj}}{\textepsilon}{\textlengthmark}{v}{a}]}} & dát \\ 
{\textipa{[{c\textsuperscript{h}}]}} & neznělá palatální ploziva s přídechem & kenna & {\textipa{[{c\textsuperscript{h}}{\textepsilon}{n}{\textlengthmark}{a}]}} & učit \\
{\textipa{[{\r{g}}]}} & neznělá velární ploziva & gata & {\textipa{[{\r{g}}{a}{\textlengthmark}{\textsubring{d}}{a}]}} & ulice \\ 
{\textipa{[{k\textsuperscript{h}}]}} & neznělá velární ploziva s přídechem & koma & {\textipa{[{k\textsuperscript{h}}{\textopeno}{\textlengthmark}{m}{a}]}} & přijít \\
{\textipa{[{v}]}} & znělá labiodentální frikativa & hafa & {\textipa{[{h}{a}{\textlengthmark}{v}{a}]}} & mít \\ 
{\textipa{[{f}]}} & neznělá labiodentální frikativa & finna & {\textipa{[{f}{\textsci}{n}{\textlengthmark}{a}]}} & najít \\ 
{\textipa{[{ð}]}} & znělá dentální frikativa & borða & {\textipa{[{\textsubring{b}}{\textopeno}{r}{ð}{a}]}} & jíst \\ 
{\textipa{[{\texttheta}]}} & neznělá dentální frikativa & þola & {\textipa{[{\texttheta}{\textopeno}{\textlengthmark}{l}{a}]}} & snést \\
{\textipa{[{s}]}} & neznělá alveolární frikativa & systir & {\textipa{[{s}{\textsci}{s}{\textsubring{d}}{\textsci}{\textsubring{r}}]}} & sestra \\ 
{\textipa{[{\c{c}}]}} & neznělá palatální frikativa & hjóla & {\textipa{[{\c{c}}{ou}{\textlengthmark}{l}{a}]}} & jet na kole \\ 
{\textipa{[{\textbabygamma}]}} & znělá velární frikativa & sagði & {\textipa{[{s}{a}{\textbabygamma}{ð}{\textsci}]}} & řekl \\ 
{\textipa{[{x}]}} & neznělá velární frikativa & lag & {\textipa{[{l}{a}{\textlengthmark}{x}]}} & píseň \\
{\textipa{[{h}]}} & neznělá glotální frikativa & heima & {\textipa{[{h}{ei}{\textlengthmark}{m}{a}]}} & doma \\ 
{\textipa{[{l}]}} & alveolární laterální aproximanta & loforð & {\textipa{[{l}{\textopeno}{\textlengthmark}{v}{\textopeno}{r}{\texttheta}]}} & slib \\ 
{\textipa{[{\textsubring{l}}]}} & neznělá alveolární laterální aproximanta & hlutur & {\textipa{[{\textsubring{l}}{\textscy}{\textlengthmark}{\textsubring{d}}{\textscy}{\textsubring{r}}]}} & věc \\  
{\textipa{[{j}]}} & palatální aproximanta & jól & {\textipa{[{j}{ou}{\textlengthmark}{\textsubring{l}}]}} & Vánoce \\ 
{\textipa{[{r}]}} & alveolární vibranta & reynsla & {\textipa{[{r}{ei}{n}{s}{\textsubring{d}}{l}{a}]}} & zkušenost \\ 
{\textipa{[{\textsubring{r}}]}} & neznělá alveolární vibranta & hraun & {\textipa{[{\textsubring{r}}{\oe i}{\textlengthmark}{\textsubring{n}}]}} & láva \\ 
\bottomrule
\end{tabular}
\end{table}

\clearpage

%\section{Výslovnost islandských samohlásek}      \label{sec:phon_vowels}
%\LTXtable{\columnwidth}{phon/vowels}
%\clearpage

%\section{Výslovnost islandských souhlásek}       \label{sec:phon_consonants}
%\LTXtable{\columnwidth}{phon/consonants}
%\clearpage
\fi

% =============================================================== MORPHOLOGY =

\ifinputmorpho

\cleardoublepage

\chapter{Morfologický klíč}                      \label{sec:morpho}

\section{Podstatné jméno, rod mužský}            \label{sec:morpho_m}
{\small\LTXfw{morpho/m}}
\clearpage

\section{Podstatné jméno, rod ženský}            \label{sec:morpho_f}
{\small\LTXfw{morpho/f}}
\clearpage

\section{Podstatné jméno, rod střední}           \label{sec:morpho_n}
{\small\LTXfw{morpho/n}}
\clearpage

\section{Přídavné jméno}                         \label{sec:morpho_adj}
{\small\LTXfw{morpho/adj}}
\clearpage

\section{Zájmeno}                                \label{sec:morpho_pron}
{\small\LTXfw{morpho/pron}}
 \begin{center}
\begin{minipage}[t]{.45\textwidth}
\begin{tabular}{l>{\footnotesize\itshape}lll}
 {\textbf{\textit{třída}}} & {\textit{pád}}   & \textit{sg} & \textit{pl}  \\ 
 \hline
\multirow{3}{*}{{{\textbf{pron} \Large{\textbf{17}}}}}  &  nom & ég & við   \\*
 & acc &  mig  & okkur  \\*
 & dat & mér & okkur   \\*
 & gen & mín  & okkar  \\
\hline

\multirow{3}{*}{{{\textbf{pron} \Large{\textbf{17}}}}}  &  nom & hann & þeir   \\*
 & acc &  hann  & þá  \\*
 & dat & honum & þeim   \\*
 & gen & hans  & þeirra  \\
\hline

\multirow{3}{*}{{{\textbf{pron} \Large{\textbf{17}}}}}  &  nom & hún & þær   \\*
 & acc &  hana  & þær  \\*
 & dat & henni & þeim   \\*
 & gen & hennar  & þeirra  \\
\hline

\multirow{3}{*}{{{\textbf{pron} \Large{\textbf{17}}}}}  &  nom &  & vér   \\*
 & acc &    & oss  \\*
 & dat &  & oss   \\*
 & gen &   & vor  \\
\hline

\multirow{3}{*}{{{\textbf{pron} \Large{\textbf{17}}}}}  &  nom &  & þér   \\*
 & acc &    & yður  \\*
 & dat &  & yður   \\*
 & gen &   & yðar  \\
\hline

\multirow{3}{*}{{{\textbf{pron} \Large{\textbf{17}}}}}  &  nom & þú & þið   \\*
 & acc &  þig  & ykkur  \\*
 & dat & þér & ykkur   \\*
 & gen & þín  & ykkar  \\
\hline

\multirow{3}{*}{{{\textbf{pron} \Large{\textbf{17}}}}}  &  nom & það & þau   \\*
 & acc &  það  & þau  \\*
 & dat & því & þeim   \\*
 & gen & þess  & þeirra  \\
\hline
\end{tabular}
\end{minipage}\hfil
\begin{minipage}[t]{.45\textwidth}
\begin{tabular}{l>{\footnotesize\itshape}lll}
\hline
\multirow{3}{*}{{{\textbf{pron} \Large{\textbf{18}}}}}  &  nom &  &    \\*
 & acc &  sig  &   \\*
 & dat & sér &    \\*
 & gen & sín  &   \\
 \hline
\end{tabular}
\end{minipage}
\end{center}

\clearpage

\section{Číslovka}                               \label{sec:morpho_num}
 {\small\LTXfw{morpho/num}}
\clearpage

\begin{landscape}
\section{Sloveso}                                \label{sec:morpho_v}
{\small\LTXfw{morpho/v}}
\end{landscape}

\clearpage

\fi

% ============================================================== BACK MATTER =

\ifmakebackmatter

% ------------------------------------------------ Photographs authors index -

\cleardoublepage
%\pagestyle{empty}
%\pagestyle{indexstyle}
%\phantomsection
%\addcontentsline{toc}{chapter}
%  {Seznam autorů fotografií a ilustrací}        
\label{sec:photo}
\printindex[figures]

% ------------------------------------------------------------- Bibliography -
\cleardoublepage
\chapter{Bibliografie}
%\twocolumn
%\phantomsection
%\addcontentsline{toc}{chapter}{Bibliografie}
\begin{multicols}{2}
\nocite{*}\printbibliography[heading=none]
\end{multicols}

% ------------------------------------------------------------------ Apendix -
\cleardoublepage
\onecolumn
\ifPDF
\else
\chapter{Apendix: Islandská gramatika}
\clearpage
\covergeometry  
\makeatletter
\includepdf[
  pages={1,2,3,4,5,6,7,8,9,10,11,12,13,14,15,16,17,18,19,20,21,22,23,24,25,26,27,28,29,30,31,32,33,34,35,36,37,38,39,40,41,42,43,44,45,46,47,48,49,50,51,52,53,54,55,56,57,58,59,60,61,62,63,64,65,66,67,68,69,70,71,72,73},
  turn=false,
  scale=0.79,
  keepaspectratio,
  offset={\dimexpr(\paperwidth-\csname Gm@layoutwidth\endcsname)/2-\csname Gm@layouthoffset\endcsname\relax}
  {\dimexpr(\paperheight-\csname  Gm@layoutheight\endcsname)/2-\csname Gm@layoutvoffset\endcsname\relax}
   ]{gram.pdf}
\makeatother
\restoregeometry  
\fi

% ------------------------------------------------------------------ Barcode -
\pagenumbering{gobble}
\ifPDF
\else
\cleardoublepage\null\clearpage % clear to verso
\vspace*\fill
\EANisbn[SC3]
\fi
\fi

% --------------------------------------------------------------- Back cover -

\ifmakecovers
  \cleardoublepage\null\clearpage % inner side of the cover
\covergeometry  
  \makecoverwith\backcoverimages
\restoregeometry  
\fi

\end{document}

